% Combined TeX source bundle

% Title: Classical-modern Chinese translation drafts

% Note: constituent files are preserved below in sources/. Some are standalone

% documents with their own preambles; this combined file is for inspection,

% comparison, and source packet rather than guaranteed one-command compilation.





% ============================================================

% Source: translations/zh/chinese/jiuzhang-suanshu-vols1-3_classical-modern_bilingual.tex

% ============================================================



%% =============================================================================

%% 九章算術 —— 文言文/白话文双语对照版

%% Jiuzhang Suanshu -- Classical Chinese / Modern Chinese Bilingual Edition

%% Volumes 1-3: 方田, 粟米, 衰分

%% =============================================================================

\documentclass[10pt,a4paper]{article}



%% -------- Core Packages --------

\usepackage{fontspec}

\usepackage{polyglossia}

\usepackage{amsmath,amssymb,amsthm}

\usepackage{geometry}

\usepackage{xcolor}

\usepackage{hyperref}

\usepackage{titlesec}

\usepackage{fancyhdr}

\usepackage{xeCJK}

\usepackage{etoolbox}

\usepackage{enumitem}

\usepackage{array,booktabs}

\usepackage{colortbl}



%% -------- Page Geometry --------

\geometry{

  paperwidth=210mm,paperheight=297mm,

  left=12mm,right=12mm,top=18mm,bottom=18mm,

  headheight=14pt,footskip=25pt

}



%% -------- Language Setup --------

\setmainlanguage{chinese}

\setotherlanguage{english}



%% -------- Font Configuration --------

\setmainfont{Microsoft YaHei}

\setsansfont{Microsoft YaHei}

\setmonofont{Microsoft YaHei}



%% -------- xeCJK --------

\setCJKmainfont{Microsoft YaHei}

\setCJKsansfont{Microsoft YaHei}

\setCJKmonofont{Microsoft YaHei}



%% -------- Colors --------

\definecolor{titlecolor}{RGB}{45,55,72}

\definecolor{sectioncolor}{RGB}{55,65,81}

\definecolor{problemcolor}{RGB}{30,60,90}

\definecolor{answercolor}{RGB}{40,80,60}

\definecolor{methodcolor}{RGB}{80,50,30}

\definecolor{commentarycolor}{RGB}{70,90,110}

\definecolor{colrule}{RGB}{200,180,160}

\definecolor{classicalbg}{RGB}{255,250,240}

\definecolor{modernbg}{RGB}{245,250,255}

\definecolor{headcolor}{RGB}{120,100,80}



%% -------- Column widths --------

\newlength{\lw}

\newlength{\rw}

\setlength{\lw}{0.48\textwidth}

\setlength{\rw}{0.48\textwidth}



%% -------- Section Formatting --------

\titleformat{\section}

  {\normalfont\Large\bfseries\color{sectioncolor}}

  {\thesection}{0.5em}{}

\titleformat{\subsection}

  {\normalfont\normalsize\bfseries\color{sectioncolor}}

  {\thesubsection}{0.5em}{}



%% -------- Header/Footer --------

\pagestyle{fancy}

\fancyhf{}

\fancyhead[L]{\small\color{headcolor}九章算術·卷一至卷三}

\fancyhead[C]{\small\color{headcolor}文言文$\circ$白话文对照版}

\fancyhead[R]{\small\color{headcolor}\thepage}

\renewcommand{\headrulewidth}{0.3pt}



%% -------- Bilingual Commands --------

% Classical column box - simple colorbg approach

\newcommand{\biling}[2]{%

  \par\vspace{0.2em}%

  \noindent%

  \colorbox{classicalbg}{\parbox{\dimexpr\lw-2\fboxsep\relax}{#1}}%

  \hfill%

  \colorbox{modernbg}{\parbox{\dimexpr\rw-2\fboxsep\relax}{#2}}%

  \vspace{0.2em}%

}



% Compact bilingual (no colorbg, smaller)

\newcommand{\bilingc}[2]{%

  \par\vspace{0.1em}%

  \noindent\begin{minipage}[t]{\lw}

  {\small #1}

  \end{minipage}%

  \hfill%

  \begin{minipage}[t]{\rw}

  {\small #2}

  \end{minipage}%

  \vspace{0.1em}%

}



% Problem/Answer/Method labels

\newcommand{\wen}{\textbf{\textcolor{problemcolor}{问：}}}

\newcommand{\da}{\textbf{\textcolor{answercolor}{答曰：}}}

\newcommand{\daB}{\textbf{\textcolor{answercolor}{【答】}}}

\newcommand{\shu}{\textbf{\textcolor{methodcolor}{术曰：}}}

\newcommand{\shuB}{\textbf{\textcolor{methodcolor}{【术】}}}

\newcommand{\zhu}{\textcolor{commentarycolor}{\small\textit{（注）}}}

\newcommand{\zhuB}{\textcolor{commentarycolor}{\small\textit{【注】}}}

\newcommand{\pn}[1]{\textbf{[#1]}}

\newcommand{\sk}{\vspace{0.1em}}



%% -------- Hyperref --------

\hypersetup{

  colorlinks=true,

  linkcolor=sectioncolor,

  pdftitle={九章算術 卷一至卷三 · 文言文白话文对照版},

  pdfauthor={译者注},

}



\setlength{\parskip}{0.1em}

\setlength{\parindent}{0pt}

\setlength{\fboxsep}{3pt}



%% =============================================================================

\begin{document}



%% -------- Title Page --------

\begin{titlepage}

\begin{center}

\vspace*{2cm}

{\fontsize{36}{42}\selectfont\bfseries\color{titlecolor} 九章算術}\\[0.5cm]

{\fontsize{16}{20}\selectfont\bfseries\color{sectioncolor} 文言文与白话文对照版}\\[0.8cm]

{\large\textit{Bilingual Edition: Classical \& Modern Chinese}}\\[1.5cm]

\rule{0.6\textwidth}{0.8pt}\\[1cm]

{\Large\bfseries 卷一至卷三}\\[0.5cm]

\begin{tabular}{cl}

\textbf{卷一} & 方田（矩形田地面积计算）\\[3pt]

\textbf{卷二} & 粟米（谷物比例交换）\\[3pt]

\textbf{卷三} & 衰分（按比例分配）\\

\end{tabular}

\vspace{1cm}

\textit{含刘徽注（公元263年）及李淳风等注释（公元656年）}\\[0.3cm]

{\small 原文为文言文（左栏），译文为现代汉语白话文（右栏）}

\vfill

\fbox{\parbox{0.85\textwidth}{\centering\small

《九章算术》是中国最重要的古典数学著作\\[2pt]

约成书于公元前200年至公元200年 \quad$\bullet$\quad 刘徽注于公元263年\\[2pt]

本章包括分数运算、面积计算、比例交换与分配等核心数学内容

}}

\vspace{1cm}

{\small 使用 \LaTeX\ + XeLaTeX + xeCJK 排版}\\[0.3cm]

\today

\end{center}

\end{titlepage}



%% -------- TOC --------

\tableofcontents

\newpage



%% =============================================================================

\section*{翻译凡例}

\addcontentsline{toc}{section}{翻译凡例}



\biling{

\textbf{左栏：文言文原文}\\[0.2em]

左栏为《九章算術》原文，采用汉代数学经典文献的原始文言文表述，保留「问」「答」「术」等原始格式。

}{

\textbf{右栏：白话文译文}\\[0.2em]

右栏为现代汉语白话文译文，力求忠实原文数学内容，以现代语言解释古代数学概念与算法。

}



\sk



\biling{

\textbf{术语对照}\\[0.2em]

\begin{tabular}{@{}l@{：}l@{}}

方田 & 矩形田地面积计算\\

约分 & 分数化简\\

合分 & 分数加法\\

今有术 & 比例计算法（三率法）\\

衰分 & 按比例分配\\

实 & 被除数/分子\\

法 & 除数/分母\\

\end{tabular}

}{

\textbf{Key Terminology}\\[0.2em]

\begin{tabular}{@{}l@{ $\to$ }l@{}}

方田 & 矩形田地面积计算\\

约分 & 分数化简\\

合分 & 分数加法\\

今有术 & 比例计算法/三率法\\

衰分 & 按比例分配\\

实 & 被除数/分子\\

法 & 除数/分母\\

\end{tabular}

}



\newpage



%% =============================================================================

\section*{刘徽九章算術注原序}

\addcontentsline{toc}{section}{刘徽序（Preface by Liu Hui）}



\biling{

\textbf{〔原文〕刘徽序}\\[0.2em]

昔在包牺氏始画八卦，以通神明之德，以类万物之情，作九九之术以合六爻之变。暨于黄帝神而化之，引而伸之，于是建历纪，协律吕，用稽道原，然后两仪四象精微之气可得而效焉。记称隶首作数，其详未之闻也。按周公制礼而有九数，九数之流，则九章是矣。

}{

\textbf{〔译文〕刘徽序}\\[0.2em]

远古伏羲氏始创八卦，用以沟通神明的德性，类比万物的情状，并发明九九算术以配合六爻的变化。到了黄帝时代，进一步神妙地发展它，引申推广，于是建立历法，协调音律，用以考察道之本源，然后天地两仪、四时四象的精微之气才能够被效法。传说隶首创造了数学，但详情已不可考。据周公制礼而有"九数"，九数的传承发展，便是《九章算术》。

}



\sk



\biling{

往者暴秦焚书，经术散坏。自时厥后，汉北平侯张苍、大司农中丞耿寿昌皆以善算命世。苍等因旧文之遗残，各称删补。故校其目则与古或异，而所论者多近语也。

}{

从前暴秦焚书，经典学术散失毁坏。自此以后，汉北平侯张苍、大司农中丞耿寿昌都以擅长算术闻名于世。张苍等人根据残存的旧文，各自进行删补。所以校对其目录与古代或有不同，而所论述的多用当时的语言。

}



\sk



\biling{

徽幼习九章，长再详览。观阴阳之割裂，总算术之根源，探赜之暇，遂悟其意。是以敢竭顽鲁，采其所见，为之作注。事类相推，各有攸归，故枝条虽分而同本干者，知发其一端而已。又所析理以辞，解体用图，庶亦约而能周，通而不黩，览之者思过半矣。

}{

刘徽幼年学习《九章》，成年后再详细研读。观察阴阳的分化，总结数学的根源，在探索深奥之理的空暇，遂领悟其意。因此斗胆竭尽愚钝，采录自己的见解，为之作注。事物按类相推，各有归属，所以枝条虽分而同一根本者，可知其发于同一端绪。又通过言辞分析道理，用图形解释形体，希望可以简约而周全，通达而不繁琐，读者可以事半功倍。

}



\sk



\biling{

且算在六艺，古者以宾兴贤能，教习国子。虽曰九数，其能穷纤入微，探测无方。至于以法相传，亦犹规矩度量可得而共，非特难为也。当今好之者寡，故世虽多通才达学，而未必能综于此耳。

}{

而且算术属于六艺之一，古代用以选拔贤能，教习贵族子弟。虽说只是九数，但它能够穷究纤微，探测无方。至于以算法相传，也如同规矩度量可以共同使用，并非特别困难。如今喜好它的人少了，所以世上虽有很多博学多才之士，却未必能精通此道。

}



\sk



\biling{

周官大司徒职，夏至日中立八尺之表，其景尺有五寸，谓之地中。说云，南戴日下万五千里。夫云尔者，以术推之。按九章立四表望远及因木望山之术，皆端旁互见，无有超邈若斯之类。然则苍等为术犹未足以博尽群数也。徽寻九数有重差之名，原其指趣乃所以施于此也。凡望极高、测绝深而兼知其远者必用重差，句股则必以重差为率，故曰重差也。

}{

《周官·大司徒》记载，夏至日中立八尺之表，其影长一尺五寸，称为地中。传说日下南方一万五千里。这种说法，是用算法推算的。按《九章》中立四表望远及因木望山的方法，都是端旁互见，没有像这般超越邈远的。然而张苍等人的算法还不足以博尽群数。刘徽探究九数中有"重差"之名，推原其旨趣乃是为此而设的。凡测量极高、极深而同时要知道其远近的，必用重差。勾股则必以重差为比率，所以叫重差。

}



\sk



\biling{

立两表于洛阳之城，令高八尺。南北各尽平地，同日度其正中之景。以景差为法，表高乘表间为实，实如法而一，所得加表高，即日去地也。以南表之景乘表间为实，实如法而一，即为从南表至南戴日下也。以南戴日下及日去地为句、股，为之求弦，即日去人也。以径寸之筒南望日，日满筒空，则定筒之长短以为股率，以筒径为句率，日去人之数为大股，大股之句即日径也。

}{

在洛阳城立两根表竿，高八尺。南北都是平地，同一天测量正午的日影。以影差为法（除数），表高乘以表间距离为实（被除数），实除以法，所得加表高，就是日离地面的高度。以南表之影乘以表间距离为实，实除以法，就是从南表到南戴日下的距离。以南戴日下及日去地为勾、股，求其弦，就是日离人的距离。用直径一寸的竹筒向南望日，日光满筒，则确定筒的长短作为股率，以筒径为勾率，日去人的距离为大股，大股的勾就是日的直径。

}



\sk



\biling{

虽天圆穹之象犹曰可度，又况泰山之高与江海之广哉。徽以为今之史籍且略举天地之物，考论厥数，载之于志，以阐世术之美。辄造重差，并为注解，以究古人之意，缀于句股之下。度高者重表，测深者累矩，孤离者三望，离而又旁求者四望。触类而长之，则虽幽遐诡伏，靡所不入。



\hfill\textit{博物君子，详而览焉。}

}{

虽天圆穹隆之象犹可度量，又何况泰山之高与江海之广呢。刘徽认为当今史籍应略举天地之物，考论其数，记载于志，以阐明数学之美。遂造重差，并为之注解，以探究古人之意，附于勾股章之后。测高者用重表，测深者用累矩，孤立者用三望，离而又旁求者用四望。触类推广，则虽幽远诡秘之处，无不涉及。



\hfill\textit{博学君子，请详览焉。}

}



\newpage





%% =============================================================================

\section{卷第一：方田}

\label{sec:fangtian}

\addcontentsline{toc}{section}{卷一：方田（矩形田地面积计算）}



\biling{

\begin{center}\textit{方田以御田畴界域}\end{center}



\textbf{方田}者，田之正也。诸田不等，以方为正，故曰方田。此章以御田畴界域，论各种田地面积之计算法，并及分数四则运算之法。

}{

\begin{center}\textit{方田——用于计算田地的边界与面积}\end{center}



\textbf{【方田章引言】}方田，即矩形的田地。各种田地形状不同，以方形（矩形）为正，所以叫方田。本章用于计算田地的边界与面积，讨论各种田地面积的计算方法，以及分数的四则运算法则。

}



\bigskip



%% ===== Problems 1-2 =====

\bilingc{

\pn{一} 今有田广十五步，从十六步。问为田几何？\quad\da 一亩。



\pn{二} 又有田广十二步，从十四步。问为田几何？\quad\da 一百六十八步。

}{

\pn{一} 现有田地宽十五步，长十六步。问田地面积是多少？\quad\daB 一亩。



\pn{二} 又有田地宽十二步，长十四步。问田地面积是多少？\quad\daB 一百六十八平方步。

}



\sk



\bilingc{

\shu 方田术曰：广从步数相乘得积步。以亩法二百四十步除之，即亩数。百亩为一顷。

}{

\shuB \textbf{方田术} 算法：宽与长的步数相乘得积步。以亩法二百四十步除之，即得亩数。一百亩为一顷。\textit{即：面积 = 宽$\times$长，亩数 = $\frac{\text{积步}}{240}$}

}



\vspace{0.3em}



%% ===== Problems 3-4 =====

\bilingc{

\pn{三} 今有田广一里，从一里。问为田几何？\quad\da 三顷七十五亩。



\pn{四} 又有田广二里，从三里。问为田几何？\quad\da 二十二顷五十亩。



\shu 里田术曰：广从里数相乘得积里。以三百七十五乘之，即亩数。

}{

\pn{三} 现有田地宽一里，长一里。问田地面积是多少？\quad\daB 三顷七十五亩。



\pn{四} 又有田地宽二里，长三里。问田地面积是多少？\quad\daB 二十二顷五十亩。



\shuB \textbf{里田术} 算法：宽与长的里数相乘得积里。以三百七十五乘之，即得亩数。

}



\vspace{0.3em}



%% ===== Problems 5-6 (Fraction Reduction) =====

\bilingc{

\pn{五} 今有十八分之十二。问约之得几何？\quad\da 三分之二。



\pn{六} 又有九十一分之四十九。问约之得几何？\quad\da 十三分之七。

}{

\pn{五} 现有分数$\frac{12}{18}$。问约分后得多少？\quad\daB 三分之二。



\pn{六} 又有分数$\frac{49}{91}$。问约分后得多少？\quad\daB 十三分之七。

}



\sk



\bilingc{

\shu 约分术曰：可半者半之，不可半者，副置分母子之数，以少减多，更相减损，求其等也。以等数约之。

}{

\shuB \textbf{约分术（分数化简）} 算法：分子分母可半者先减半（同除以2），不可半者，另置分子分母之数，以少减多，更相减损，求其等数（最大公约数）。以等数约分子分母。\textit{即欧几里得辗转相除法求GCD。}}

}



\vspace{0.3em}



%% ===== Problems 7-9 (Fraction Addition) =====

\bilingc{

\pn{七} 今有$\frac13$，$\frac25$。问合之得几何？\quad\da $\frac{11}{15}$。



\pn{八} 又有$\frac23$，$\frac47$，$\frac59$。问合之得几何？\quad\da 得一、$\frac{50}{63}$。



\pn{九} 又有$\frac12$，$\frac23$，$\frac34$，$\frac45$。问合之得几何？\quad\da 得二、$\frac{43}{60}$。

}{

\pn{七} 现有$\frac{1}{3}$和$\frac{2}{5}$。问相加得多少？\quad\daB 十五分之十一。



\pn{八} 又有$\frac{2}{3}$、$\frac{4}{7}$和$\frac{5}{9}$。问相加得多少？\quad\daB 得一又六十三分之五十。



\pn{九} 又有$\frac{1}{2}$、$\frac{2}{3}$、$\frac{3}{4}$和$\frac{4}{5}$。问相加得多少？\quad\daB 得二又六十分之四十三。

}



\sk



\bilingc{

\shu 合分术曰：母互乘子，并为实，母相乘为法，实如法而一。不满法者，以法命之。其母同者，直相从之。

}{

\shuB \textbf{合分术（分数加法）} 算法：分母互乘分子，相加为实，分母相乘为法，实除以法得一。\textit{即：$\frac{a}{b}+\frac{c}{d}=\frac{ad+bc}{bd}$}

}



\vspace{0.3em}



%% ===== Problems 10-11 =====

\bilingc{

\pn{一〇} 今有$\frac89$，减其$\frac15$。问余几何？\quad\da $\frac{31}{45}$。



\pn{一一} 又有$\frac34$，减其$\frac13$。问余几何？\quad\da $\frac{5}{12}$。

}{

\pn{一〇} 现有$\frac{8}{9}$，减去其$\frac{1}{5}$。问余多少？\quad\daB 四十五分之三十一。



\pn{一一} 又有$\frac{3}{4}$，减去其$\frac{1}{3}$。问余多少？\quad\daB 十二分之五。

}



\sk



\bilingc{

\shu 减分术曰：母互乘子，以少减多，余为实，母相乘为法，实如法而一。

}{

\shuB \textbf{减分术（分数减法）} 算法：分母互乘分子，以大减小，余数为实，分母相乘为法，实除以法得一。\textit{即：$\frac{a}{b}-\frac{c}{d}=\frac{ad-bc}{bd}$}

}



\vspace{0.3em}



%% ===== Problems 12-14 =====

\bilingc{

\pn{一二} 今有$\frac58$，$\frac{16}{25}$。问孰多？多几何？\quad\da $\frac{16}{25}$多，多$\frac{3}{200}$。



\pn{一三} 又有$\frac89$，$\frac67$。问孰多？多几何？\quad\da $\frac89$多，多$\frac{2}{63}$。



\pn{一四} 又有$\frac{8}{21}$，$\frac{17}{50}$。问孰多？多几何？\quad\da $\frac{8}{21}$多，多$\frac{43}{1050}$。

}{

\pn{一二} 现有$\frac{5}{8}$和$\frac{16}{25}$。问哪个大？大多少？\quad\daB $\frac{16}{25}$大，大$\frac{3}{200}$。



\pn{一三} 又有$\frac{8}{9}$和$\frac{6}{7}$。问哪个大？大多少？\quad\daB $\frac{8}{9}$大，大$\frac{2}{63}$。



\pn{一四} 又有$\frac{8}{21}$和$\frac{17}{50}$。问哪个大？大多少？\quad\daB $\frac{8}{21}$大，大$\frac{43}{1050}$。

}



\sk



\bilingc{

\shu 课分术曰：母互乘子，以少减多，余为实，母相乘为法，实如法而一，即相多也。

}{

\shuB \textbf{课分术（分数比较）} 算法：分母互乘分子，以大减小，余为实，分母相乘为法，实除以法，即得两数之差。

}



\vspace{0.3em}



%% ===== Problems 15-16 =====

\bilingc{

\pn{一五} 今有$\frac13$，$\frac23$，$\frac34$。问减多益少，各几何而平？



\da 减$\frac34$者二，减$\frac23$者一，并以益$\frac13$，而各平于$\frac{7}{12}$。



\pn{一六} 又有$\frac12$，$\frac23$，$\frac34$。问减多益少，各几何而平？



\da 减$\frac23$者一，减$\frac34$者四，并以益$\frac12$，而各平于$\frac{23}{36}$。

}{

\pn{一五} 现有$\frac{1}{3}$、$\frac{2}{3}$、$\frac{3}{4}$。问从多的减去、补给少的，各多少才能相等？



\daB 从$\frac{3}{4}$减$\frac{2}{12}$，从$\frac{2}{3}$减$\frac{1}{12}$，都补给$\frac{1}{3}$，各等于$\frac{7}{12}$。



\pn{一六} 又有$\frac{1}{2}$、$\frac{2}{3}$、$\frac{3}{4}$。问从多的减去、补给少的，各多少才能相等？



\daB 从$\frac{2}{3}$减$\frac{1}{36}$，从$\frac{3}{4}$减$\frac{4}{36}$，都补给$\frac{1}{2}$，各等于$\frac{23}{36}$。

}



\sk



\bilingc{

\shu 平分术曰：母互乘子，副并为平实，母相乘为法。以列数乘未并者各自为列实。亦以列数乘法，以平实减列实，余，约之为所减。并所减以益于少，以法命平实，各得其平。

}{

\shuB \textbf{平分术（分数等分）} 算法：分母互乘分子，另加一起作为平实，分母相乘为法。以个数乘未并者各自为列实。以平实减列实，余数约分后即为应减之数。将所减之数合并补给少的，各得相等之数。

}



\vspace{0.3em}



%% ===== Problems 17-18 =====

\bilingc{

\pn{一七} 今有七人，分八钱三分钱之一。问人得几何？\quad\da 人得一钱、二十一分钱之四。



\pn{一八} 又有三人，三分人之一，分六钱三分钱之一，四分钱之三。问人得几何？\quad\da 人得二钱、八分钱之一。

}{

\pn{一七} 现有七人，分$8\frac{1}{3}$钱。问每人得多少？\quad\daB 每人得一又二十一分之四钱。



\pn{一八} 又有$3\frac{1}{3}$人，分$6\frac{1}{3}+\frac{3}{4}$钱。问每人得多少？\quad\daB 每人得二又八分之一钱。

}



\sk



\bilingc{

\shu 经分术曰：以人数为法，钱数为实，实如法而一。有分者通之，重有分者同而通之。

}{

\shuB \textbf{经分术（分数除法）} 算法：以人数为法（除数），钱数为实（被除数），实除以法得一。有分数的先通分，有多重分数的先统一再通分。

}



\vspace{0.3em}



%% ===== Problems 19-21 =====

\bilingc{

\pn{一九} 今有田广$\frac47$步，从$\frac35$步。问为田几何？\quad\da $\frac{12}{35}$步。



\pn{二〇} 又有田广$\frac79$步，从$\frac9{11}$步。问为田几何？\quad\da $\frac{7}{11}$步。



\pn{二一} 又有田广$\frac45$步，从$\frac59$步。问为田几何？\quad\da $\frac49$步。

}{

\pn{一九} 现有田地宽$\frac{4}{7}$步，长$\frac{3}{5}$步。问面积多少？\quad\daB $\frac{12}{35}$平方步。



\pn{二〇} 又有田地宽$\frac{7}{9}$步，长$\frac{9}{11}$步。问面积多少？\quad\daB $\frac{7}{11}$平方步。



\pn{二一} 又有田地宽$\frac{4}{5}$步，长$\frac{5}{9}$步。问面积多少？\quad\daB $\frac{4}{9}$平方步。

}



\sk



\bilingc{

\shu 乘分术曰：母相乘为法，子相乘为实，实如法而一。

}{

\shuB \textbf{乘分术（分数乘法）} 算法：分母相乘为法，分子相乘为实，实除以法得一。\textit{即：$\frac{a}{b}\times\frac{c}{d}=\frac{ac}{bd}$}

}



\vspace{0.3em}



%% ===== Problems 22-24 =====

\bilingc{

\pn{二二} 今有田广$3\frac13$步，从$5\frac25$步。问为田几何？\quad\da 十八步。



\pn{二三} 又有田广$7\frac34$步，从$15\frac59$步。问为田几何？\quad\da 百二十步、$\frac59$步。



\pn{二四} 又有田广$18\frac57$步，从$23\frac6{11}$步。问为田几何？\quad\da 一亩二百步、$\frac7{11}$步。

}{

\pn{二二} 现有田地宽$3\frac{1}{3}$步，长$5\frac{2}{5}$步。问面积多少？\quad\daB 十八平方步。



\pn{二三} 又有田地宽$7\frac{3}{4}$步，长$15\frac{5}{9}$步。问面积多少？\quad\daB 一百二十又九分之五平方步。



\pn{二四} 又有田地宽$18\frac{5}{7}$步，长$23\frac{6}{11}$步。问面积多少？\quad\daB 一亩二百又十一分之七平方步。

}



\sk



\bilingc{

\shu 大广田术曰：分母各乘其全，分子从之，相乘为实。分母相乘为法。实如法而一。

}{

\shuB \textbf{大广田术（带分数乘法）} 算法：分母各乘其整数部分，加上分子，相乘为实。分母相乘为法。实除以法得一。\textit{即计算带分数的乘法。}

}



\vspace{0.3em}



%% ===== Problems 25-26 =====

\bilingc{

\pn{二五} 今有圭田广十二步，正从二十一步。问为田几何？\quad\da 一百二十六步。



\pn{二六} 又有圭田广$5\frac12$步，从$8\frac23$步。问为田几何？\quad\da $23\frac56$步。

}{

\pn{二五} 现有圭田（三角形田）底宽十二步，正长二十一步。问面积多少？\quad\daB 一百二十六平方步。



\pn{二六} 又有圭田底宽$5\frac{1}{2}$步，长$8\frac{2}{3}$步。问面积多少？\quad\daB 二十三又六分之五平方步。

}



\sk



\bilingc{

\shu 术曰：半广以乘正从。



\zhu 半广者，以盈补虚，为直田也。亦可以半正从以乘广。

}{

\shuB \textbf{圭田术} 算法：以底宽的一半乘以正长。\textit{即：三角形面积 = $\frac12\times$底$\times$高}



\zhuB 半广，即以多余补不足，化为矩形田。也可以用半正长乘以底宽。

}



\vspace{0.3em}



%% ===== Problems 27-28 =====

\bilingc{

\pn{二七} 今有邪田，一头广三十步，一头广四十二步，正从六十四步。问为田几何？



\da 九亩一百四十四步。



\pn{二八} 又有邪田，正广六十五步，一畔从一百步，一畔从七十二步。问为田几何？



\da 二十三亩七十步。

}{

\pn{二七} 现有邪田（梯形田），一头宽三十步，另一头宽四十二步，正长六十四步。问面积多少？



\daB 九亩一百四十四平方步。



\pn{二八} 又有邪田，正宽六十五步，一斜边长一百步，另一斜边长七十二步。问面积多少？



\daB 二十三亩七十平方步。

}



\sk



\bilingc{

\shu 术曰：并两邪而半之，以乘正从若广。又可半正从若广，以乘并，亩法而一。

}{

\shuB \textbf{邪田术（梯形面积）} 算法：将两斜边（两底）相加取半，乘以正长（高）。也可以半正长乘以两底之和，再以亩法换算。\textit{即：梯形面积 = $\frac{\text{上底}+\text{下底}}{2}\times$高}

}



\vspace{0.3em}



%% ===== Problems 29-30 =====

\bilingc{

\pn{二九} 今有箕田，舌广二十步，踵广五步，正从三十步。问为田几何？\quad\da 一亩一百三十五步。



\pn{三〇} 又有箕田，舌广一百一十七步，踵广五十步，正从一百三十五步。问为田几何？



\da 四十六亩二百三十二步半。

}{

\pn{二九} 现有箕田（簸箕形田），舌宽二十步，踵宽五步，正长三十步。问面积多少？\quad\daB 一亩一百三十五平方步。



\pn{三〇} 又有箕田，舌宽一百一十七步，踵宽五十步，正长一百三十五步。问面积多少？



\daB 四十六亩二百三十二又二分之一平方步。

}



\sk



\bilingc{

\shu 术曰：并踵舌而半之，以乘正从。亩法而一。

}{

\shuB \textbf{箕田术} 算法：将踵宽与舌宽相加取半，乘以正长。以亩法换算。\textit{即：梯形面积 = $\frac{\text{上底}+\text{下底}}{2}\times$高}

}



\vspace{0.3em}



%% ===== Problems 31-32 =====

\bilingc{

\pn{三一} 今有圆田，周三十步，径十步。问为田几何？\quad\da 七十五步。



\pn{三二} 又有圆田，周一百八十一步，径六十步、三分步之一。问为田几何？



\da 十一亩九十步、十二分步之一。

}{

\pn{三一} 现有圆田，周长三十步，直径十步。问面积多少？\quad\daB 七十五平方步。



\pn{三二} 又有圆田，周长一百八十一步，直径$60\frac{1}{3}$步。问面积多少？



\daB 十一亩九十又十二分之一平方步。

}



\sk



\bilingc{

\shu 术曰：半周半径相乘得积步。又术曰：周径相乘，四而一。又术曰：径自相乘，三之，四而一。又术曰：周自相乘，十二而一。

}{

\shuB \textbf{圆田术} 四种等价方法：一、半周长与半径相乘得积步。二、周长与直径相乘，除以四。三、直径自乘，乘以三，除以四。四、周长自乘，除以十二。\textit{注：此处采用"周三径一"近似，$\pi\approx 3$}

}



\sk



\bilingc{

\zhu 此于徽术，当为田十亩二百八步三百一十四分步之一百一。臣淳风等谨依密率，为田十亩二百步八十八分步之八十七。按：半周为从，半径为广，故广从相乘为积步也。假令圆径二尺，圆中容六觚之一面，与圆径之半，其数均等。合径率一而外周率三也。



又按：为图，以六觚之一面乘一弧半径，三之，得十二觚之幂。若又割之，次以十二觚之一面乘一弧之半径，六之，则得二十四觚之幂。割之弥细，所失弥少。割之又割，以至于不可割，则与圆周合体而无所失矣。

}{

\zhuB 按刘徽的精确算法，应为十亩二百又三百一十四分之一百一平方步。李淳风等依密率计算，为十亩二百又八十八分之八十七平方步。



注：半周长为长，半径为宽，故宽长相乘为积步。假设圆直径二尺，圆内接正六边形的一边与圆半径相等。所以径率为一，周率为三。



又按：用图形分析，以正六边形一边乘以弧半径，乘以三，得正十二边形的面积。若继续分割，以正十二边形一边乘以弧半径，乘以六，得正二十四边形的面积。分割越细，误差越小。割之又割，以至于不可分割，则与圆完全重合而无误差了。\textit{此即刘徽著名的"割圆术"思想。}

}



\vspace{0.3em}



%% ===== Problems 33-34 =====

\bilingc{

\pn{三三} 今有宛田，下周三十步，径十六步。问为田几何？\quad\da 一百二十步。



\pn{三四} 又有宛田，下周九十九步，径五十一步。问为田几何？\quad\da 五亩六十二步、四分步之一。



\shu 术曰：以径乘周，四而一。

}{

\pn{三三} 现有宛田（丘形田），下周三十步，径十六步。问面积多少？\quad\daB 一百二十平方步。



\pn{三四} 又有宛田，下周九十九步，径五十一步。问面积多少？\quad\daB 五亩六十二又四分之一平方步。



\shuB \textbf{宛田术} 算法：以直径乘以周长，除以四。

}



\vspace{0.3em}



%% ===== Problems 35-36 =====

\bilingc{

\pn{三五} 今有弧田，弦三十步，矢十五步。问为田几何？\quad\da 一亩九十七步半。



\pn{三六} 又有弧田，弦七十八步、二分步之一，矢十三步、九分步之七。问为田几何？



\da 二亩一百五十五步、八十一分步之五十六。



\shu 术曰：以弦乘矢，矢又自乘，并之，二而一。

}{

\pn{三五} 现有弧田（弓形田），弦长三十步，矢高十五步。问面积多少？\quad\daB 一亩九十七又二分之一平方步。



\pn{三六} 又有弧田，弦长$78\frac{1}{2}$步，矢高$13\frac{7}{9}$步。问面积多少？



\daB 二亩一百五十五又八十一分之五十六平方步。



\shuB \textbf{弧田术} 算法：以弦乘以矢高，矢高又自乘，两者相加，除以二。\textit{即：弓形面积$\approx\frac12(\text{弦}\times\text{矢}+\text{矢}^2)$}

}



\vspace{0.3em}



%% ===== Problems 37-38 =====

\bilingc{

\pn{三七} 今有环田，中周九十二步，外周一百二十二步，径五步。问为田几何？\quad\da 二亩五十五步。



\pn{三八} 又有环田，中周$62\frac34$步，外周$113\frac12$步，径$12\frac23$步。问为田几何？



\da 四亩一百五十六步、四分步之一。



\shu 术曰：并中外周而半之，以径乘之为积步。



密率术曰：置中外周步数，分母、子各居其下。母互乘子，通全步，内分子。以中周减外周，余半之，以益中周。径亦通分内子，以乘周为实。分母相乘为法，除之为积步，余积步之分。以亩法除之，即亩数也。

}{

\pn{三七} 现有环田（环形田），中周九十二步，外周一百二十二步，径宽五步。问面积多少？\quad\daB 二亩五十五平方步。



\pn{三八} 又有环田，中周$62\frac{3}{4}$步，外周$113\frac{1}{2}$步，径宽$12\frac{2}{3}$步。问面积多少？



\daB 四亩一百五十六又四分之一平方步。



\shuB \textbf{环田术}\\算法：中外周相加取半，乘以径宽得积步。\\密率算法：置中外周步数，分子分母各列其下。分母互乘分子，通全步，内加分子。以外周减中周，余数取半，加入中周。径也通分内加分子，乘以周为实。分母相乘为法，除之得积步。\\\textit{即：环形面积 = 平均周长$\times$径宽}

}



\newpage





%% =============================================================================

\section{卷第二：粟米}

\label{sec:sumi}

\addcontentsline{toc}{section}{卷二：粟米（谷物比例交换）}



\biling{

\begin{center}\textit{粟米以御交质变易}\end{center}



本章论各种谷物之交换比率，及比例计算之法。\textbf{今有术}（即今之三率法、比例法）为本章之核心算法。

}{

\begin{center}\textit{粟米——用于谷物交换与比例换算}\end{center}



本章讨论各种谷物之间的交换比率，以及比例计算的方法。\textbf{【今有术】}（即现代所说的三率法、比例法）是本章的核心算法。\textit{公式：所求数 = $\frac{\text{所有数}\times\text{所求率}}{\text{所有率}}$}

}



\sk



\biling{

\begin{center}\textbf{粟米之法}\end{center}

\begin{tabular}{@{}llll@{}}

粟率五十 & 糲米三十 & 粺米二十七 & 凿米二十四 \\

御米二十一 & 小$\bullet$十三半 & 大$\bullet$五十四 & 糲饭七十五 \\

粺饭五十四 & 凿饭四十八 & 御饭四十二 & 菽荅麻麦各四十五 \\

稻六十 & 豉六十三 & 飧九十 & 熟菽一百三半 \\

\multicolumn{4}{@{}l@{}}{櫱一百七十五} \\

\end{tabular}

}{

\begin{center}\textbf{【粟米兑换率表】}\end{center}

\begin{tabular}{@{}llll@{}}

粟米50 & 糙粳米30 & 精白米27 & 精米24 \\

御米21 & 小豆$13\frac12$ & 大豆54 & 糙粳饭75 \\

精白米饭54 & 精米饭48 & 御米饭42 & 豆麻麦各45 \\

稻60 & 豉63 & 飧90 & 熟豆$103\frac12$ \\

\multicolumn{4}{@{}l@{}}{麦芽175} \\

\end{tabular}

}



\sk



\bilingc{

\shu 今有术曰：以所有数乘所求率为实，以所有率为法，实如法而一。

}{

\shuB \textbf{今有术（比例法/三率法）} 算法：以所有数乘以所求率为实（被除数），以所有率为法（除数），实除以法得一。\textit{即：所求数 = $\frac{\text{所有数}\times\text{所求率}}{\text{所有率}}$}

}



\bigskip



%% ===== Problems 1-4 =====

\bilingc{

\pn{一} 今有粟一斗，欲为糲米。问得几何？\quad\da 为糲米六升。



\shu 术曰：以粟求糲米，三之，五而一。

}{

\pn{一} 现有粟米一斗，想换成糙粳米。问能得多少？\quad\daB 得糙粳米六升。



\shuB 算法：以粟求糙粳米，乘以三，除以五。\textit{即：$\frac{30}{50}\times 1$斗 = 6升}

}



\sk



\bilingc{

\pn{二} 今有粟二斗一升，欲为粺米。问得几何？\quad\da 为粺米一斗一升、五十分升之十七。



\shu 术曰：以粟求粺米，二十七之，五十而一。

}{

\pn{二} 现有粟米二斗一升，想换成精白米。问能得多少？\quad\daB 得精白米一斗一升又五十分之十七升。



\shuB 算法：以粟求精白米，乘以二十七，除以五十。

}



\sk



\bilingc{

\pn{三} 今有粟四斗五升，欲为凿米。问得几何？\quad\da 为凿米二斗一升、五分升之三。



\shu 术曰：以粟求凿米，十二之，二十五而一。

}{

\pn{三} 现有粟米四斗五升，想换成精米。问能得多少？\quad\daB 得精米二斗一升又五分之三升。



\shuB 算法：以粟求精米，乘以十二，除以二十五。

}



\sk



\bilingc{

\pn{四} 今有粟七斗九升，欲为御米。问得几何？\quad\da 为御米三斗三升、五十分升之九。



\shu 术曰：以粟求御米，二十一之，五十而一。

}{

\pn{四} 现有粟米七斗九升，想换成御米。问能得多少？\quad\daB 得御米三斗三升又五十分之九升。



\shuB 算法：以粟求御米，乘以二十一，除以五十。

}



\vspace{0.3em}



%% ===== Problems 5-6 =====

\bilingc{

\pn{五} 今有粟一斗，欲为小$\bullet$。问得几何？\quad\da 为小$\bullet$二升、一十分升之七。



\shu 术曰：以粟求小$\bullet$，二十七之，百而一。



\pn{六} 今有粟九斗八升，欲为大$\bullet$。问得几何？\quad\da 为大$\bullet$一十斗五升、二十五分升之二十一。



\shu 术曰：以粟求大$\bullet$，二十七之，二十五而一。

}{

\pn{五} 现有粟米一斗，想换成小豆。问能得多少？\quad\daB 得小豆二升又十分之七升。



\shuB 算法：以粟求小豆，乘以二十七，除以一百。



\pn{六} 现有粟米九斗八升，想换成大豆。问能得多少？\quad\daB 得大豆十斗五升又二十五分之二十一升。



\shuB 算法：以粟求大豆，乘以二十七，除以二十五。

}



\vspace{0.3em}



%% ===== Problems 7-10 =====

\bilingc{

\pn{七} 今有粟二斗三升，欲为糲饭。问得几何？\quad\da 为糲饭三斗四升半。



\shu 术曰：以粟求糲饭，三之，二而一。



\pn{八} 今有粟三斗六升，欲为粺饭。问得几何？\quad\da 为粺饭三斗八升、二十五分升之二十二。



\shu 术曰：以粟求粺饭，二十七之，二十五而一。

}{

\pn{七} 现有粟米二斗三升，想换成糙粳饭。问能得多少？\quad\daB 得糙粳饭三斗四升半。



\shuB 算法：以粟求糙粳饭，乘以三，除以二。



\pn{八} 现有粟米三斗六升，想换成精白米饭。问能得多少？\quad\daB 得精白米饭三斗八升又二十五分之二十二升。



\shuB 算法：以粟求精白米饭，乘以二十七，除以二十五。

}



\sk



\bilingc{

\pn{九} 今有粟八斗六升，欲为凿饭。问得几何？\quad\da 为凿饭八斗二升、二十五分升之一十四。



\shu 术曰：以粟求凿饭，二十四之，二十五而一。



\pn{一〇} 今有粟九斗八升，欲为御饭。问得几何？\quad\da 为御饭八斗二升、二十五分升之八。



\shu 术曰：以粟求御饭，二十一之，二十五而一。

}{

\pn{九} 现有粟米八斗六升，想换成精米饭。问能得多少？\quad\daB 得精米饭八斗二升又二十五分之十四升。



\shuB 算法：以粟求精米饭，乘以二十四，除以二十五。



\pn{一〇} 现有粟米九斗八升，想换成御米饭。问能得多少？\quad\daB 得御米饭八斗二升又二十五分之八升。



\shuB 算法：以粟求御米饭，乘以二十一，除以二十五。

}



\vspace{0.3em}



%% ===== Problems 11-14 =====

\bilingc{

\pn{一一} 今有粟三斗少半升，欲为菽。\quad\da 为菽二斗七升、一十分升之三。



\pn{一二} 今有粟四斗一升、太半升，欲为荅。\quad\da 为荅三斗七升半。



\pn{一三} 今有粟五斗、太半升，欲为麻。\quad\da 为麻四斗五升、五分升之三。



\pn{一四} 今有粟一十斗八升、五分升之二，欲为麦。\quad\da 为麦九斗七升、二十五分升之一十四。



\shu 术曰：以粟求菽、荅、麻、麦，皆九之，十而一。

}{

\pn{一一} 现有粟米三斗又$\frac13$升，想换成豆子。\quad\daB 得豆子二斗七升又十分之三升。



\pn{一二} 现有粟米四斗一升又$\frac23$升，想换成荅（小豆）。\quad\daB 得荅三斗七升半。



\pn{一三} 现有粟米五斗又$\frac23$升，想换成麻。\quad\daB 得麻四斗五升又五分之三升。



\pn{一四} 现有粟米十斗八升又$\frac25$升，想换成麦。\quad\daB 得麦九斗七升又二十五分之十四升。



\shuB 算法：以粟求豆、荅、麻、麦，皆乘以九，除以十。\textit{即：所求 = 粟$\times\frac{9}{10}$}

}



\vspace{0.3em}



%% ===== Problems 15-19 =====

\bilingc{

\pn{一五} 今有粟七斗五升、七分升之四，欲为稻。\quad\da 为稻九斗、三十五分升之二十四。



\shu 术曰：以粟求稻，六之，五而一。



\pn{一六} 今有粟七斗八升，欲为豉。\quad\da 为豉九斗八升、二十五分升之七。



\shu 术曰：以粟求豉，六十三之，五十而一。



\pn{一七} 今有粟五斗五升，欲为飧。\quad\da 为飧九斗九升。



\shu 术曰：以粟求飧，九之，五而一。



\pn{一八} 今有粟四斗，欲为熟菽。\quad\da 为熟菽八斗二升、五分升之四。



\shu 术曰：以粟求熟菽，二百七之，百而一。



\pn{一九} 今有粟二斗，欲为櫱。\quad\da 为櫱七斗。



\shu 术曰：以粟求櫱，七之，二而一。

}{

\pn{一五} 现有粟米七斗五升又七分之四升，想换成稻。\quad\daB 得稻九斗又三十五分之二十四升。



\shuB 算法：以粟求稻，乘以六，除以五。



\pn{一六} 现有粟米七斗八升，想换成豉。\quad\daB 得豉九斗八升又二十五分之七升。



\shuB 算法：以粟求豉，乘以六十三，除以五十。



\pn{一七} 现有粟米五斗五升，想换成飧。\quad\daB 得飧九斗九升。



\shuB 算法：以粟求飧，乘以九，除以五。



\pn{一八} 现有粟米四斗，想换成熟豆。\quad\daB 得熟豆八斗二升又五分之四升。



\shuB 算法：以粟求熟豆，乘以二百零七，除以一百。



\pn{一九} 现有粟米二斗，想换成麦芽。\quad\daB 得麦芽七斗。



\shuB 算法：以粟求麦芽，乘以七，除以二。

}



\vspace{0.3em}



%% ===== Problems 20-24 =====

\bilingc{

\pn{二〇} 今有糲米十五斗五升、五分升之二，欲为粟。\quad\da 为粟二十五斗九升。



\shu 术曰：以糲米求粟，五之，三而一。



\pn{二一} 今有粺米二斗，欲为粟。\quad\da 为粟三斗七升、二十七分升之一。



\shu 术曰：以粺米求粟，五十之，二十七而一。



\pn{二二} 今有凿米三斗、少半升，欲为粟。\quad\da 为粟六斗三升、三十六分升之七。



\shu 术曰：以凿米求粟，二十五之，十三而一。



\pn{二三} 今有御米十四斗，欲为粟。\quad\da 为粟三十三斗三升、少半升。



\shu 术曰：以御米求粟，五十之，二十一而一。



\pn{二四} 今有稻一十二斗六升、一十五分升之一十四，欲为粟。\quad\da 为粟一十斗五升、九分升之七。



\shu 术曰：以稻求粟，五之，六而一。

}{

\pn{二〇} 现有糙粳米十五斗五升又五分之二升，想换成粟。\quad\daB 得粟米二十五斗九升。



\shuB 算法：以糙粳米求粟，乘以五，除以三。



\pn{二一} 现有精白米二斗，想换成粟。\quad\daB 得粟米三斗七升又二十七分之一升。



\shuB 算法：以精白米求粟，乘以五十，除以二十七。



\pn{二二} 现有精米三斗又$\frac13$升，想换成粟。\quad\daB 得粟米六斗三升又三十六分之七升。



\shuB 算法：以精米求粟，乘以二十五，除以十三。



\pn{二三} 现有御米十四斗，想换成粟。\quad\daB 得粟米三十三斗三升又三分之一升。



\shuB 算法：以御米求粟，乘以五十，除以二十一。



\pn{二四} 现有稻十二斗六升又十五分之十四升，想换成粟。\quad\daB 得粟米十斗五升又九分之七升。



\shuB 算法：以稻求粟，乘以五，除以六。

}



\vspace{0.3em}



%% ===== Problems 25-27 =====

\bilingc{

\pn{二五} 今有糲米一十九斗二升、七分升之一，欲为粺米。\quad\da 为粺米一十七斗二升、一十四分升之一十三。



\shu 术曰：以糲米求粺米，九之，十而一。



\pn{二六} 今有糲米六斗四升、五分升之三，欲为糲饭。\quad\da 为糲饭一十六斗一升半。



\shu 术曰：以糲米求糲饭，五之，二而一。



\pn{二七} 今有糲饭七斗六升、七分升之四，欲为飧。\quad\da 为飧九斗一升、三十五分升之三十一。



\shu 术曰：以糲饭求飧，六之，五而一。

}{

\pn{二五} 现有糙粳米十九斗二升又七分之一升，想换成精白米。\quad\daB 得精白米十七斗二升又十四分之十三升。



\shuB 算法：以糙粳米求精白米，乘以九，除以十。



\pn{二六} 现有糙粳米六斗四升又五分之三升，想换成糙粳饭。\quad\daB 得糙粳饭十六斗一升半。



\shuB 算法：以糙粳米求糙粳饭，乘以五，除以二。



\pn{二七} 现有糙粳饭七斗六升又七分之四升，想换成飧。\quad\daB 得飧九斗一升又三十五分之三十一升。



\shuB 算法：以糙粳饭求飧，乘以六，除以五。

}



\vspace{0.3em}



%% ===== Problems 28-31 =====

\bilingc{

\pn{二八} 今有菽一斗，欲为熟菽。\quad\da 为熟菽二斗三升。



\shu 术曰：以菽求熟菽，二十三之，十而一。



\pn{二九} 今有菽二斗，欲为豉。\quad\da 为豉二斗八升。



\shu 术曰：以菽求豉，七之，五而一。



\pn{三〇} 今有麦八斗六升、七分升之三，欲为小$\bullet$。\quad\da 为小$\bullet$二斗五升、一十四分升之一十三。



\shu 术曰：以麦求小$\bullet$，三之，十而一。



\pn{三一} 今有麦一斗，欲为大$\bullet$。\quad\da 为大$\bullet$一斗二升。



\shu 术曰：以麦求大$\bullet$，六之，五而一。

}{

\pn{二八} 现有豆子一斗，想换成熟豆。\quad\daB 得熟豆二斗三升。



\shuB 算法：以豆子求熟豆，乘以二十三，除以十。



\pn{二九} 现有豆子二斗，想换成豉。\quad\daB 得豉二斗八升。



\shuB 算法：以豆子求豉，乘以七，除以五。



\pn{三〇} 现有麦八斗六升又七分之三升，想换成小豆。\quad\daB 得小豆二斗五升又十四分之十三升。



\shuB 算法：以麦求小豆，乘以三，除以十。



\pn{三一} 现有麦一斗，想换成大豆。\quad\daB 得大豆一斗二升。



\shuB 算法：以麦求大豆，乘以六，除以五。

}



\vspace{0.3em}



%% ===== Problems 32-33 =====

\bilingc{

\pn{三二} 今有出钱一百六十，买瓴甓十八枚。问枚几何？\quad\da 一枚，八钱、九分钱之八。



\pn{三三} 今有出钱一万三千五百，买竹二千三百五十个。问个几何？\quad\da 一个，五钱、四十七分钱之三十五。



\shu 经率术曰：以所买率为法，所出钱数为实，实如法得一钱。

}{

\pn{三二} 现有出钱一百六十，买瓦砖十八枚。问每枚多少钱？\quad\daB 每枚八又九分之八钱。



\pn{三三} 现有出钱一万三千五百，买竹二千三百五十个。问每个多少钱？\quad\daB 每个五又四十七分之三十五钱。



\shuB \textbf{经率术（单价计算）} 算法：以所买数量为法（除数），所出钱数为实（被除数），实除以法得每单位价格。

}



\vspace{0.3em}



%% ===== Problems 34-37 =====

\bilingc{

\pn{三四} 今有出钱五千七百八十五，买漆一斛六斗七升、太半升。欲斗率之，问斗几何？



\da 一斗，三百四十五钱、五百三分钱之一十五。



\pn{三五} 今有出钱七百二十，买缣一匹二丈一尺。欲丈率之，问丈几何？



\da 一丈，一百一十八钱、六十一分钱之二。

}{

\pn{三四} 现有出钱五千七百八十五，买漆一斛六斗七升又$\frac23$升。想按每斗计单价，问每斗多少钱？



\daB 每斗三百四十五又五百零三分之十五钱。



\pn{三五} 现有出钱七百二十，买缣（细绢）一匹二丈一尺。想按每丈计单价，问每丈多少钱？



\daB 每丈一百一十八又六十一分之二钱。

}



\sk



\bilingc{

\pn{三六} 今有出钱二千三百七十，买布九匹二丈七尺。欲匹率之，问匹几何？



\da 一匹，二百四十四钱、一百二十九分升之一百二十四。



\pn{三七} 今有出钱一万三千六百七十，买丝一石二钧一十七斤。欲石率之，问石几何？



\da 一石，八千三百二十六钱、一百九十七分钱之一百七十八。



\shu 经术术曰：以所求率乘钱数为实，以所买率为法，实如法得一。

}{

\pn{三六} 现有出钱二千三百七十，买布九匹二丈七尺。想按每匹计单价，问每匹多少钱？



\daB 每匹二百四十四又一百二十九分之一百二十四钱。



\pn{三七} 现有出钱一万三千六百七十，买丝一石二钧一十七斤。想按每石计单价，问每石多少钱？



\daB 每石八千三百二十六又一百九十七分之一百七十八钱。



\shuB \textbf{经术术} 算法：以所求率乘以钱数为实，以所买数量为法，实除以法得一。

}



\vspace{0.3em}



%% ===== Problems 38-39 =====

\bilingc{

\pn{三八} 今有出钱五百七十六，买竹七十八个。欲其大小率之，问各几何？



\da 其四十八个，个七钱。其三十个，个八钱。



\pn{三九} 今有出钱一千一百二十，买丝一石二钧一十八斤。欲其贵贱斤率之，问各几何？



\da 其二钧八斤，斤五钱。其一石一十斤，斤六钱。

}{

\pn{三八} 现有出钱五百七十六，买竹七十八个。想按大小不同单价，问各多少钱？



\daB 其中四十八个，每个七钱；三十个，每个八钱。



\pn{三九} 现有出钱一千一百二十，买丝一石二钧一十八斤。想按贵贱两种单价论斤计，问各多少钱一斤？



\daB 其中二钧八斤，每斤五钱；一石十斤，每斤六钱。

}



\sk



\bilingc{

\shu 其率术曰：各置所买石、钧、斤、两以为法，以所率乘钱数为实，实如法而一。不满法者反以实减法，法贱实贵。

}{

\shuB \textbf{其率术（贵贱分配）} 算法：各置所买石、钧、斤、两作为法，以所率乘钱数为实，实除以法得一。余数不满法者，反过来以实减法，法对应贱价，实对应贵价。

}



\vspace{0.3em}



%% ===== Problems 45-46 =====

\bilingc{

\pn{四五} 今有出钱六百一十，买羽二千一百翭。欲其贵贱率之，问各几何？



\da 其一千一百四十翭，三翭一钱。其九百六十翭，四翭一钱。



\pn{四六} 今有出钱九百八十，买矢簳五千八百二十枚。欲其贵贱率之，问各几何？



\da 其三百枚，五枚一钱。其五千五百二十枚，六枚一钱。

}{

\pn{四五} 现有出钱六百一十，买羽二千一百翭（鸟羽单位）。想按贵贱两种单价，问各多少？



\daB 其中一千一百四十翭，三翭一钱；九百六十翭，四翭一钱。



\pn{四六} 现有出钱九百八十，买箭杆五千八百二十枚。想按贵贱两种单价，问各多少？



\daB 其中三百枚，五枚一钱；五千五百二十枚，六枚一钱。

}



\sk



\bilingc{

\shu 反其率术曰：以钱数为法，所率为实，实如法而一。不满法者反以实减法，法少，实多。二物各以所得多少之数乘法实，即物数。

}{

\shuB \textbf{反其率术（反贵贱分配）} 算法：以钱数为法，所率为实，实除以法得一。不满法者反以实减法，法对应少数，实对应多数。两种物品各以所得多少之数乘法实，即得物品数量。

}



\newpage





%% =============================================================================

\section{卷第三：衰分}

\label{sec:cuifen}

\addcontentsline{toc}{section}{卷三：衰分（按比例分配）}



\biling{

\begin{center}\textit{衰分以御贵贱粟税之率}\end{center}



本章论比例分配之法，包括按等级、按数量、按比率之分配问题。\textbf{衰分术}为本章之核心算法。

}{

\begin{center}\textit{衰分——用于按等级、比率分配赋税与物资}\end{center}



本章讨论比例分配的方法，包括按等级、按数量、按比率分配的问题。\textbf{【衰分术】}（即按比例分配）是本章的核心算法。\textit{公式：各份 = $\frac{\text{总数}\times\text{该份比率}}{\text{各比率之和}}$}

}



\sk



\bilingc{

\shu 衰分术曰：各置列衰，副并为法，以所分乘未并者各自为实，实如法而一。不满法者，以法命之。

}{

\shuB \textbf{衰分术（按比例分配）} 算法：各置分配比率（列衰），另加一起作为法（除数），以所分总数乘未并之各比率各自为实（被除数），实除以法得一。余数不满法者，以分数表示。\textit{即：各份 = $\frac{\text{总数}\times\text{该份比率}}{\text{各比率之和}}$}

}



\bigskip



%% ===== Problem 1 =====

\bilingc{

\pn{一} 今有大夫、不更、簪裹、上造、公士，凡五人，共猎得五鹿。欲以爵次分之，问各得几何？



\da 大夫得一鹿、三分鹿之二。不更得一鹿、三分鹿之一。簪裹得一鹿。上造得三分鹿之二。公士得三分鹿之一。



\shu 术曰：列置爵数，各自为衰，副并为法。以五鹿乘未并者，各自为实。实如法得一鹿。

}{

\pn{一} 现有大夫、不更、簪袅、上造、公士共五人，一起打猎得五只鹿。想按爵位等级分配，问各得多少？



\daB 大夫得一又三分之二只鹿。不更得一又三分之一只鹿。簪袅得一只鹿。上造得三分之二只鹿。公士得三分之一只鹿。



\shuB 算法：列置各爵位等级数作为分配比率，另加一起作为法。以五鹿乘各比率作为实。实除以法得各人所分鹿数。\textit{比率：大夫5、不更4、簪袅3、上造2、公士1，总和15。}

}



\vspace{0.3em}



%% ===== Problem 2 =====

\bilingc{

\pn{二} 今有牛、马、羊食人苗。苗主责之粟五斗。羊主曰：「我羊食半马。」马主曰：「我马食半牛。」今欲衰偿之，问各出几何？



\da 牛主出二斗八升、七分升之四。马主出一斗四升、七分升之二。羊主出七升、七分升之一。



\shu 术曰：置牛四、马二、羊一，各自为列衰，副并为法。以五斗乘未并者各自为实。实如法得一斗。

}{

\pn{二} 现有牛、马、羊吃了人家的禾苗。苗主要求赔偿粟五斗。羊主人说：「我的羊吃的是马的一半。」马主人说：「我的马吃的是牛的一半。」现在想按比例赔偿，问各出多少？



\daB 牛主人出二斗八升又七分之四升。马主人出一斗四升又七分之二升。羊主人出七升又七分之一升。



\shuB 算法：置牛四、马二、羊一作为分配比率，另加一起（共七）作为法。以五斗乘各比率作为实。实除以法得各应出粟数。\textit{即：牛$:$马$:$羊 = $4:2:1$}

}



\vspace{0.3em}



%% ===== Problem 3 =====

\bilingc{

\pn{三} 今有甲持钱五百六十，乙持钱三百五十，丙持钱一百八十，凡三人俱出关，关税百钱。欲以钱数多少衰出之，问各几何？



\da 甲出五十一钱、一百九分钱之四十一。乙出三十二钱、一百九分钱之一十二。丙出一十六钱、一百九分钱之五十六。



\shu 术曰：各置钱数为列衰，副并为法，以百钱乘未并者，各自为实，实如法得一钱。

}{

\pn{三} 现有甲带钱五百六十，乙带钱三百五十，丙带钱一百八十，三人一起出关，关税一百钱。想按各人所带钱数比例出税，问各出多少？



\daB 甲出五十一又一百零九分之四十一钱。乙出三十二又一百零九分之十二钱。丙出十六又一百零九分之五十六钱。



\shuB 算法：各置所持钱数作为分配比率，另加一起作为法。以一百钱乘各比率作为实。实除以法得各人应纳税额。\textit{比率：甲560、乙350、丙180，总和1090。}

}



\vspace{0.3em}



%% ===== Problem 4 =====

\bilingc{

\pn{四} 今有女子善织，日自倍，五日织五尺。问日织几何？



\da 初日织一寸、三十一分寸之十九。次日织三寸、三十一分寸之七。次日织六寸、三十一分寸之十四。次日织一尺二寸、三十一分寸之二十八。次日织二尺五寸、三十一分寸之二十五。



\shu 术曰：置一、二、四、八、十六为列衰，副并为法，以五尺乘未并者，各自为实，实如法得一尺。

}{

\pn{四} 现有女子善于织布，每天织的布是前一天的二倍，五天共织五尺。问每天各织多少？



\daB 第一天织一又三十一分之十九寸。第二天织三又三十一分之七寸。第三天织六又三十一分之十四寸。第四天织一尺二又三十一分之二十八寸。第五天织二尺五又三十一分之二十五寸。



\shuB 算法：置1、2、4、8、16作为分配比率（等比数列），另加一起（共31）作为法。以五尺乘各比率作为实。实除以法得每日织布数。\textit{即：每天织布量按$1:2:4:8:16$的比例分配五尺。}

}



\vspace{0.3em}



%% ===== Problem 5 =====

\bilingc{

\pn{五} 今有北乡算八千七百五十八，西乡算七千二百三十六，南乡算八千三百五十六，凡三乡，发傜三百七十八人。欲以算数多少衰出之，问各几何？



\da 北乡遣一百三十五人、一万二千一百七十五分人之一万一千六百三十七。西乡遣一百一十二人、一万二千一百七十五分人之四千四。南乡遣一百二十九人、一万二千一百七十五分之八千七百九。



\shu 术曰：各置算数为列衰，副并为法，以所发傜人数乘未并者，各自为实，实如法得一人。

}{

\pn{五} 现有北乡人口（算数）八千七百五十八，西乡七千二百三十六，南乡八千三百五十六，共三乡，需征发徭役三百七十八人。想按人口比例征发，问各乡出多少人？



\daB 北乡征发一百三十五又一万二千一百七十五分之一万一千六百三十七人。西乡征发一百一十二又一万二千一百七十五分之四千零四人。南乡征发一百二十九又一万二千一百七十五分之八千七百零九人。



\shuB 算法：各置人口数为分配比率，另加一起作为法。以所征发徭役人数乘各乡人口数作为实。实除以法得各乡应出人数。\textit{比率即各乡人口数：北8758、西7236、南8356，总和24350。}

}



\vspace{0.3em}



%% ===== Problem 6 =====

\bilingc{

\pn{六} 今有禀粟，大夫、不更、簪裹、上造、公士，凡五人，一十五斗。今有大夫一人后来，亦当禀五斗。仓无粟，欲以衰出之，问各几何？



\da 大夫出一斗、四分斗之一。不更出一斗。簪褭出四分斗之三。上造出四分斗之二。公士出四分斗之一。



\shu 术曰：各置所禀粟斛斗数，爵次均之，以为列衰，副并而加后来大夫亦五斗，得二十以为法。以五斗乘未并者各自为实。实如法得一斗。

}{

\pn{六} 现有发放给大夫、不更、簪袅、上造、公士五人的粟共十五斗。现在有一位大夫后到，也应发五斗。仓库无粟可取，想按比例从各人之粟中分出补给，问各出多少？



\daB 大夫出一又四分之一斗。不更出一斗。簪袅出四分之三斗。上造出四分之二斗。公士出四分之一斗。



\shuB 算法：各置所发粟数，按爵位等级平均作为列衰，另加一起并加入后来大夫的五斗，共二十作为法。以五斗乘未并之各衰作为实。实除以法得各应出粟数。

}



\vspace{0.3em}



%% ===== Problem 7 =====

\bilingc{

\pn{七} 今有禀粟五斛，五人分之，欲令三人得三，二人得二。问各几何？



\da 三人，人得一斛一斗五升、十三分升之五。二人，人得七斗六升、十三分升之十二。



\shu 术曰：置三人，人三；二人，人二，为列衰。副并为法。以五斛乘未并者，各自为实。实如法得一斛。

}{

\pn{七} 现有分发的粟五斛，分给五人，想让三人各得三份，二人各得二份。问每人各得多少？



\daB 三人中每人得一斛一斗五升又十三分之五升。二人中每人得七斗六升又十三分之十二升。



\shuB 算法：置三人，每人分配比率3；二人，每人分配比率2，作为列衰。另加一起（共13）作为法。以五斛乘未并之各衰作为实。实除以法得各人应得粟数。\textit{总衰 = $3\times3+2\times2=13$}

}



\vspace{0.3em}



%% ===== Reverse Distribution =====

\bilingc{

\shu 返衰术曰：列置衰而令相乘，动者为不动者衰。

}{

\shuB \textbf{返衰术（反比例分配）} 算法：列置各衰而令其相乘，动者作为不动者的衰（比率）。\textit{即：反比例分配，按倒数比率分配，使得等级高者少出，等级低者多出。}

}



\vspace{0.3em}



%% ===== Problem 8 =====

\bilingc{

\pn{八} 今有大夫、不更、簪褭、上造、公士，凡五人，共出百钱。欲令高爵出少，以次渐多，问各几何？



\da 大夫出八钱、一百三十七分钱之一百四。不更出一十钱、一百三十七分钱之一百三十。簪褭出一十四钱、一百三十七分钱之八十二。上造出二十一钱、一百三十七分钱之一百二十三。公士出四十三钱、一百三十七分钱之一百九。



\shu 术曰：置爵数各自为衰，而返衰之，副并为法。以百钱乘未并者各自为实。实如法得一钱。

}{

\pn{八} 现有大夫、不更、簪袅、上造、公士共五人，一起出一百钱。想让爵位高的少出，依次渐多，问各出多少？



\daB 大夫出八又一百三十七分之一百零四钱。不更出十又一百三十七分之一百三十钱。簪袅出十四又一百三十七分之八十二钱。上造出二十一又一百三十七分之一百二十三钱。公士出四十三又一百三十七分之一百零九钱。



\shuB 算法：置各爵位等级数作为衰，取倒数（返衰），另加一起作为法。以一百钱乘各返衰作为实。实除以法得各应出钱数。\textit{正衰为$5,4,3,2,1$，返衰为$\frac15,\frac14,\frac13,\frac12,1$。}

}



\vspace{0.3em}



%% ===== Problem 9 =====

\bilingc{

\pn{九} 今有甲持粟三升，乙持糲米三升，丙持糲饭三升。欲令合而分之，问各几何？



\da 甲二升、一十分升之七。乙四升、一十分升之五。丙一升、一十分升之八。



\shu 术曰：以粟率五十、糲米率三十、糲饭率七十五为衰、而返衰之，副并为法。以九升乘未并者各自为实。实如法得一升。

}{

\pn{九} 现有甲带粟三升，乙带糙粳米三升，丙带糙粳饭三升。想将三者合并后重新按价值分配，问各得多少？



\daB 甲得二又十分之七升。乙得四又十分之五升。丙得一又十分之八升。



\shuB 算法：以粟率50、糙粳米率30、糙粳饭率75作为衰，取倒数（返衰），另加一起作为法。以九升乘各返衰作为实。实除以法得各人应得之升数。\textit{返衰后：粟$\frac{1}{50}$、米$\frac{1}{30}$、饭$\frac{1}{75}$，按此比例重新分配。}

}



\vspace{0.3em}



%% ===== Problem 10 =====

\bilingc{

\pn{一〇} 今有丝一斤，价直二百四十。今有钱一千三百二十八，问得丝几何？



\da 五斤八两一十二铢、三分铢之四。



\shu 术曰：以一斤价数为法，以一斤乘今有钱数为实，实如法得丝数。

}{

\pn{一〇} 现有丝一斤，价值二百四十钱。现有钱一千三百二十八，问能买多少丝？



\daB 五斤八两十二又三分之四铢。



\shuB 算法：以一斤的价格为法（除数），以一斤乘现有钱数为实（被除数），实除以法得可买丝数。\textit{即：丝数 = $\frac{1\times1328}{240}$ 斤}

}



\vspace{0.3em}



%% ===== Problem 11 =====

\bilingc{

\pn{一一} 今有丝一斤价直三百四十三。今有丝七两一十二铢，问得钱几何？



\da 一百六十一钱、三十二分钱之二十三。



\shu 术曰：以一斤铢数为法，以一斤价数，乘七两一十二铢为实。实如法得钱数。

}{

\pn{一一} 现有丝一斤价值三百四十三钱。现有丝七两十二铢，问能卖多少钱？



\daB 一百六十一又三十二分之二十三钱。



\shuB 算法：以一斤的铢数（384铢）为法，以一斤的价格乘七两十二铢为实。实除以法得钱数。

}



\vspace{0.3em}



%% ===== Problem 12 =====

\bilingc{

\pn{一二} 今有缣一丈价直一百二十六。今有缣一匹九尺五寸，问得钱几何？



\da 六百三十三钱、五分钱之三。



\shu 术曰：以一丈寸数为法，以价钱数乘今有缣寸数为实，实如法得钱数。

}{

\pn{一二} 现有缣（细绢）一丈价值一百二十六钱。现有缣一匹九尺五寸，问能卖多少钱？



\daB 六百三十三又五分之三钱。



\shuB 算法：以一丈的寸数为法，以价钱乘现有缣的寸数为实。实除以法得钱数。

}



\vspace{0.3em}



%% ===== Problem 13 =====

\bilingc{

\pn{一三} 今有布一匹，价直一百二十三。今有布二丈七尺，问得钱几何？



\da 八十四钱、八十一分钱之三。



\shu 术曰：以一匹尺数为法，今有布尺数乘价钱为实，实如法得钱数。

}{

\pn{一三} 现有布一匹价值一百二十三钱。现有布二丈七尺，问能卖多少钱？



\daB 八十四又八十一分之三钱。



\shuB 算法：以一匹的尺数为法，现有布尺数乘价钱为实。实除以法得钱数。

}



\vspace{0.3em}



%% ===== Problem 14 =====

\bilingc{

\pn{一四} 今有素一匹一丈，价直六百二十五。今有钱五百，问得素几何？



\shu 术曰：以价直为法，以一匹一丈尺数乘今有钱数为实。实如法得素数。

}{

\pn{一四} 现有素（白绢）一匹一丈价值六百二十五钱。现有钱五百，问能买多少素？



\shuB 算法：以价格为法，以一匹一丈的尺数乘现有钱数为实。实除以法得可买素数。

}



\vspace{0.3em}



%% ===== Problem 15 =====

\bilingc{

\pn{一五} 今有与人丝一十四斤，约得缣一十斤。今与人丝四十五斤八两，问得缣几何？



\da 三十二斤八两。



\shu 术曰：以一十四斤两数为法，以一十斤乘今有丝两数为实，实如法得缣数。

}{

\pn{一五} 现有给人丝十四斤，约定换得缣十斤。现在给人丝四十五斤八两，问能换得多少缣？



\daB 三十二斤八两。



\shuB 算法：以十四斤的铢数为法，以十斤乘现有丝的铢数为实。实除以法得缣数。

}



\vspace{0.3em}



%% ===== Problem 16 =====

\bilingc{

\pn{一六} 今有丝一斤，耗七两。今有丝二十三斤五两，问耗几何？



\da 一百六十三两四铢半。



\shu 术曰：以一斤展十六两为法，以七两乘今有丝两数为实，实如法得耗数。

}{

\pn{一六} 现有丝一斤耗损七两。现有丝二十三斤五两，问耗损多少？



\daB 一百六十三两四又二分之一铢。



\shuB 算法：以一斤折合十六两为法，以七两乘现有丝的两数为实。实除以法得耗损数。

}



\vspace{0.3em}



%% ===== Problem 17 =====

\bilingc{

\pn{一七} 今有生丝三十斤，乾之，耗三斤十二两。今有乾丝一十二斤，问生丝几何？



\da 一十三斤一十一两十铢、七分铢之二。



\shu 术曰：置生丝两数，除耗数，余，以为法。三十斤乘乾丝两数为实。实如法得生丝数。

}{

\pn{一七} 现有生丝三十斤，晒干后耗损三斤十二两。现有干丝十二斤，问原需生丝多少？



\daB 十三斤十一两十又七分之二铢。



\shuB 算法：置生丝两数减去耗数，余数作为法。三十斤乘干丝两数为实。实除以法得所需生丝数。\textit{即：生丝 $:$ 干丝 = $30 : (30-3.75)$，已知干丝求生丝。}

}



\vspace{0.3em}



%% ===== Problem 18 =====

\bilingc{

\pn{一八} 今有田一亩，收粟六升、太半升。今有田一顷二十六亩一百五十九步，问收粟几何？



\da 八斛四斗四升、一十二分升之五。



\shu 术曰：以亩二百四十步为法，以六升、太半升乘今有田积步为实，实如法得粟数。

}{

\pn{一八} 现有田一亩，收粟六升又三分之二升。现有田一顷二十六亩一百五十九步，问收粟多少？



\daB 八斛四斗四升又十二分之五升。



\shuB 算法：以亩二百四十步为法，以六又三分之二升乘现有田的积步为实。实除以法得粟数。

}



\vspace{0.3em}



%% ===== Problem 19 =====

\bilingc{

\pn{一九} 今有取保一岁，价钱二千五百。今先取一千二百，问当作日几何？



\da 一百六十九日、二十五分日之二十三。



\shu 术曰：以价钱为法，以一岁三百五十四日乘先取钱数为实，实如法得日数。

}{

\pn{一九} 现有雇佣一年的工价二千五百钱。现已先取一千二百钱，问应做工多少天？



\daB 一百六十九又二十五分之二十三天。



\shuB 算法：以工价为法，以一年三百五十四天乘先取钱数为实。实除以法得应做工天数。\textit{即：做工天数 = $\frac{354\times1200}{2500}$}

}



\vspace{0.3em}



%% ===== Problem 20 =====

\bilingc{

\pn{二〇} 今有贷人千钱，月息三十。今有贷人七百五十钱，九日归之，问息几何？



\da 六钱、四分钱之三。



\shu 术曰：以月三十日，乘千钱为法。以息三十乘今所贷钱数，又以九日乘之，为实。实如法得一钱。

}{

\pn{二〇} 现有借给人一千钱，月息三十钱。现有借给人七百五十钱，九天归还，问利息多少？



\daB 六又四分之三钱。



\shuB 算法：以月三十天乘一千钱为法。以息三十乘所贷钱数，再以九天乘之，为实。实除以法得利息钱数。\textit{即：利息 = $\frac{30\times750\times9}{30\times1000}$ 钱}

}



\vspace{0.8em}



%% -------- End matter --------

\begin{center}

\rule{0.5\textwidth}{0.4pt}

\end{center}



\vspace{0.3em}



\biling{

\begin{center}\textit{（以上为《九章算術》卷第一至卷第三之全文，含刘徽注及李淳风等注释。）}\end{center}

}{

\begin{center}\textit{（以上为《九章算术》卷一至卷三全文，含刘徽注及李淳风等注释。）}\end{center}

}



\newpage



%% =============================================================================

\section*{译注说明}

\addcontentsline{toc}{section}{译注说明}



\biling{

\textbf{一、术语说明}



\textbf{实}与\textbf{法}为古代算学之核心概念。实即被除数（分子），法即除数（分母）。「实如法而一」意为以法除实，得一整数商；若有余数，则以分数表示。

}{

\textbf{一、术语说明}



\textbf{实}与\textbf{法}是古代数学的核心概念。实即被除数（分子），法即除数（分母）。「实如法而一」的意思是用法去除实，得到一个整数商；如果有余数，则用分数表示。

}



\sk



\biling{

\textbf{今有术}即今日所称之比例法或三率法（Rule of Three）。以所有数乘所求率为实，所有率为法，实如法得所求数。此术为粟米、衰分诸章之根本算法，广泛应用于商业、税收和工程计算。

}{

\textbf{今有术}就是今天所说的比例法或三率法。公式为：所求数 = （所有数 $\times$ 所求率）$\div$ 所有率。这是粟米、衰分各章的根本算法，广泛应用于商业、税收和工程计算。

}



\sk



\biling{

\textbf{衰分术}中「衰」（cu$\bar{\mathrm{i}}$）即比率、等级之意。列衰即排列各比率，用于按比例分配资源。返衰则为反比例分配，用于等级高者少出、等级低者多出之情形。

}{

\textbf{衰分术}中的「衰」（cu$\bar{\mathrm{i}}$）就是比率、等级的意思。列衰即排列各比率，用于按比例分配资源。返衰则为反比例分配，用于等级高者少出、等级低者多出之情形。衰分术是古代分配理论和税收理论的重要数学基础。

}



\sk



\biling{

\textbf{二、度量衡换算}



\begin{tabular}{@{}l@{}}

1顷 = 100亩 \\[1pt]

1亩 = 240平方步 \\[1pt]

1里 = 300步 \\[1pt]

1斤 = 16两 = 384铢 \\[1pt]

1匹 = 4丈 = 40尺 \\[1pt]

1斛 = 10斗 = 100升 \\[1pt]

1石 = 4钧 = 120斤 \\[1pt]

\end{tabular}

}{

\textbf{二、度量衡换算}



\begin{tabular}{@{}l@{}}

1顷 = 100亩（面积单位） \\[1pt]

1亩 = 240平方步 \\[1pt]

1里 = 300步（长度单位） \\[1pt]

1斤 = 16两 = 384铢（重量单位） \\[1pt]

1匹 = 4丈 = 40尺（布帛长度） \\[1pt]

1斛 = 10斗 = 100升（容量单位） \\[1pt]

1石 = 4钧 = 120斤（重量单位） \\[1pt]

\end{tabular}

}



\sk



\biling{

\textbf{三、刘徽割圆术}



刘徽于圆田注中提出著名的「割圆术」。以圆内接正六边形开始，逐次倍增边数，「割之弥细，所失弥少。割之又割，以至于不可割，则与圆周合体而无所失矣」。刘徽以此求得 $\pi \approx \frac{157}{50} = 3.14$，后世称为「徽率」。祖冲之在此基础上进一步求得密率 $\frac{355}{113}$。

}{

\textbf{三、刘徽割圆术}



刘徽在圆田术的注释中提出了著名的「割圆术」。从圆内接正六边形开始，逐次倍增边数，「分割越细，误差越小。割之又割，以至于不可分割，则与圆完全重合而无误差」。刘徽以此求得圆周率 $\pi \approx \frac{157}{50} = 3.14$，后世称为「徽率」。祖冲之在此基础上进一步求得密率 $\frac{355}{113}$。这是中国数学史上对圆周率精确计算的重要贡献。

}



\vspace{0.8em}



\begin{center}

{\small 使用 \LaTeX\ + XeLaTeX + xeCJK 排版 \quad$\bullet$\quad \today}

\end{center}



\end{document}







% ============================================================

% Source: translations/zh/chinese/jiuzhang-suanshu-vols4-6_classical-modern_bilingual.tex

% ============================================================



% =============================================================================

% 九章算術·卷四至卷六 — 文言文↔白话文对照本

% Jiuzhang Suanshu Volumes 4-6: Classical↔Modern Chinese Bilingual Edition

% =============================================================================

\documentclass[11pt,a4paper]{article}



% -------- Core Packages --------

\usepackage{fontspec}

\usepackage{xeCJK}

\usepackage{polyglossia}

\setmainlanguage{english}



% CJK Font Configuration

\setCJKmainfont{Microsoft YaHei}



% -------- Line breaking for Chinese text --------

\tolerance=1200

\emergencystretch=3em

\hyphenpenalty=5000

\exhyphenpenalty=5000



\usepackage{amsmath,amssymb}

\usepackage{geometry}

\usepackage{graphicx}

\usepackage{xcolor}

\usepackage{titlesec}

\usepackage{fancyhdr}

\usepackage{array,booktabs,longtable,multirow}

\usepackage{hyperref}

\usepackage{tocloft}

\usepackage{indentfirst}

\usepackage{ragged2e}

\usepackage{paracol}

\usepackage{tcolorbox}

\tcbuselibrary{skins,breakable}



% -------- Page Geometry --------

\geometry{

  paperwidth=210mm,paperheight=297mm,

  left=18mm,right=18mm,top=25mm,bottom=25mm,

  headheight=14pt,footskip=25pt

}



% -------- Colors --------

\definecolor{titlecolor}{RGB}{45,55,72}

\definecolor{sectioncolor}{RGB}{55,65,81}

\definecolor{volcolor}{RGB}{120,53,15}

\definecolor{commentarycolor}{RGB}{80,60,40}

\definecolor{rulecolor}{RGB}{180,160,140}

\definecolor{problemcolor}{RGB}{25,55,95}

\definecolor{answercolor}{RGB}{35,85,55}

\definecolor{methodcolor}{RGB}{85,35,35}

\definecolor{leftbg}{RGB}{252,250,245}

\definecolor{rightbg}{RGB}{245,250,252}

\definecolor{leftframe}{RGB}{160,140,110}

\definecolor{rightframe}{RGB}{110,140,160}

\definecolor{classicalcolor}{RGB}{90,50,30}

\definecolor{moderncolor}{RGB}{30,60,90}



% -------- Column Ratio --------

\columnratio{0.50,0.50}



% -------- Section Formatting --------

\titleformat{\section}

  {\normalfont\Large\bfseries\color{titlecolor}}

  {\thesection}{1em}{}

\titleformat{\subsection}

  {\normalfont\large\bfseries\color{sectioncolor}}

  {\thesubsection}{1em}{}

\titleformat{\subsubsection}

  {\normalfont\normalsize\bfseries\color{sectioncolor}}

  {\thesubsubsection}{1em}{}



% -------- Header/Footer --------

\pagestyle{fancy}

\fancyhf{}

\fancyhead[L]{\small\color{sectioncolor}九章算術·卷四至卷六 — 文言文↔白话文对照本}

\fancyhead[R]{\small\thepage}

\renewcommand{\headrulewidth}{0.4pt}

\renewcommand{\headrule}{\hbox to\headwidth{\color{rulecolor}\leaders\hrule height \headrulewidth\hfill}}



% -------- Custom Environments --------



% Classical (left column) problem box

\newtcolorbox{classicalproblem}[1]{

  colback=leftbg,colframe=leftframe!80!black,

  fonttitle=\bfseries\color{classicalcolor},title={#1},

  breakable,enhanced,boxrule=0.6pt,arc=2pt,

  before upper={\RaggedRight},

  left=4pt,right=4pt,top=3pt,bottom=3pt

}



% Modern (right column) problem box

\newtcolorbox{modernproblem}[1]{

  colback=rightbg,colframe=rightframe!80!black,

  fonttitle=\bfseries\color{moderncolor},title={#1},

  breakable,enhanced,boxrule=0.6pt,arc=2pt,

  before upper={\RaggedRight},

  left=4pt,right=4pt,top=3pt,bottom=3pt

}



% Classical answer box

\newtcolorbox{classicalanswer}[1]{

  colback=green!2!white,colframe=answercolor!50!black,

  fonttitle=\bfseries\color{classicalcolor},title={#1},

  breakable,enhanced,boxrule=0.5pt,arc=2pt,

  before upper={\RaggedRight},

  left=4pt,right=4pt,top=3pt,bottom=3pt

}



% Modern answer box

\newtcolorbox{modernanswer}[1]{

  colback=green!2!white,colframe=answercolor!50!black,

  fonttitle=\bfseries\color{moderncolor},title={#1},

  breakable,enhanced,boxrule=0.5pt,arc=2pt,

  before upper={\RaggedRight},

  left=4pt,right=4pt,top=3pt,bottom=3pt

}



% Classical method box

\newtcolorbox{classicalmethod}[1]{

  colback=red!2!white,colframe=methodcolor!50!black,

  fonttitle=\bfseries\color{classicalcolor},title={#1},

  breakable,enhanced,boxrule=0.5pt,arc=2pt,

  before upper={\RaggedRight},

  left=4pt,right=4pt,top=3pt,bottom=3pt

}



% Modern method box

\newtcolorbox{modernmethod}[1]{

  colback=red!2!white,colframe=methodcolor!50!black,

  fonttitle=\bfseries\color{moderncolor},title={#1},

  breakable,enhanced,boxrule=0.5pt,arc=2pt,

  before upper={\RaggedRight},

  left=4pt,right=4pt,top=3pt,bottom=3pt

}



% Classical commentary box

\newtcolorbox{classicalcommentary}[1]{

  colback=yellow!3!white,colframe=commentarycolor!50!black,

  fonttitle=\bfseries\itshape\color{classicalcolor},title={#1},

  breakable,enhanced,boxrule=0.4pt,arc=2pt,

  before upper={\RaggedRight},

  left=4pt,right=4pt,top=3pt,bottom=3pt

}



% Modern commentary box

\newtcolorbox{moderncommentary}[1]{

  colback=yellow!3!white,colframe=commentarycolor!50!black,

  fonttitle=\bfseries\itshape\color{moderncolor},title={#1},

  breakable,enhanced,boxrule=0.4pt,arc=2pt,

  before upper={\RaggedRight},

  left=4pt,right=4pt,top=3pt,bottom=3pt

}



% Volume header

\newtcolorbox{volumeheader}[2]{

  colback=volcolor!8!white,colframe=volcolor!70!black,

  fonttitle=\bfseries\Large,title={#2},#1,

  breakable,enhanced,boxrule=1pt,arc=3pt

}



% Bilingual header box

\newtcolorbox{bilingheader}[1]{

  colback=volcolor!5!white,colframe=volcolor!60!black,

  fonttitle=\bfseries\large,#1,

  breakable,enhanced,boxrule=0.8pt,arc=3pt,

  before upper={\RaggedRight}

}



% -------- Custom Commands --------

\newcommand{\longline}{\noindent\rule{\textwidth}{0.4pt}}

\newcommand{\shortline}{\noindent\rule{0.5\textwidth}{0.4pt}\par}



% Synchronized environment markers

\newcommand{\wenleft}{\textcolor{problemcolor}{\textbf{【问】}}}

\newcommand{\wenright}{\textcolor{problemcolor}{\textbf{【问题】}}}

\newcommand{\daleft}{\textcolor{answercolor}{\textbf{【答】}}}

\newcommand{\daright}{\textcolor{answercolor}{\textbf{【答案】}}}

\newcommand{\shuleft}{\textcolor{methodcolor}{\textbf{【术】}}}

\newcommand{\shuright}{\textcolor{methodcolor}{\textbf{【算法】}}}

\newcommand{\zhuleft}{\textcolor{commentarycolor}{\textbf{【注】}}}

\newcommand{\zhuright}{\textcolor{commentarycolor}{\textbf{【注释】}}}



\hypersetup{

  colorlinks=true,linkcolor=sectioncolor,citecolor=sectioncolor,

  urlcolor=sectioncolor,pdfencoding=auto

}



\setcounter{tocdepth}{2}

\setcounter{secnumdepth}{2}



% Paracol synchronization

\footnotelayout{m}



% =============================================================================

\begin{document}



% ---------------------------------------------------------------------------

% TITLE PAGE

% ---------------------------------------------------------------------------

\begin{titlepage}

\begin{center}

\vspace*{1.5cm}



{\Huge\bfseries\color{titlecolor} 九章算術}\[0.8em]

{\LARGE\color{sectioncolor} 卷四至卷六}\[1.5em]



{\Large\color{volcolor}

 文言文 ↔ 白话文 对照本}\[2em]



{\large\itshape Jiuzhang Suanshu / Nine Chapters on the Mathematical Art}\[0.5em]

{\normalsize\itshape Volumes 4--6: Classical ↔ Modern Chinese Bilingual Edition}\[2em]



\shortline



\vspace{1cm}



{\large\bfseries 卷四：少广（开方和体积计算）}\\[0.3em]

{\normalsize Vol.\ 4 --- Shaoguang: Root extraction, sphere volume}\\[0.8em]



{\large\bfseries 卷五：商功（工程体积计算）}\\[0.3em]

{\normalsize Vol.\ 5 --- Shanggong: Volumes of walls, dykes, canals, various solids}\\[0.8em]



{\large\bfseries 卷六：均输（公平分配和运输）}\\[0.3em]

{\normalsize Vol.\ 6 --- Junshu: Weighted distribution, travel problems}\\[1.5em]



\shortline



\vspace{1cm}



\begin{minipage}{0.9\textwidth}

\centering

\textit{左栏：原文（文言文）} \hspace{1em} | \hspace{1em} \textit{右栏：白话文翻译}\\[0.3em]

\textit{Left: Original Classical Chinese} \hspace{0.5em} | \hspace{0.5em} \textit{Right: Modern Chinese Vernacular}

\end{minipage}



\vspace{1.5cm}



{\small （晋）刘徽注 \quad | \quad （唐）李淳风注释}\\[0.5em]

{\footnotesize With Commentary by Liu Hui (263 CE) and sub-commentary by Li Chunfeng ( Tang)}



\vfill



{\small 排版使用 XeLaTeX + Microsoft YaHei}\\

{\footnotesize 来源：钦定四库全书版}



\end{center}

\end{titlepage}



% ---------------------------------------------------------------------------

% TABLE OF CONTENTS

% ---------------------------------------------------------------------------

\tableofcontents

\newpage



% ---------------------------------------------------------------------------

% INTRODUCTION

% ---------------------------------------------------------------------------

\section{简介 / Introduction}



\begin{paracol}{2}



\textbf{\color{classicalcolor}【文言文原文】}



《九章算術》是中國古代最重要的數學經典，編纂於漢代（公元前202年—公元220年）。全書分為九卷，每卷論述一類實際數學問題。劉徽（約225—295年）於263年完成《九章算術注》，不僅解釋了各種算法，還運用「出入相補」原理給出了嚴格的幾何證明。李淳風（602—670年）奉唐詔對本文進一步疏證。



本對照本呈現\textbf{卷四至卷六}：

\begin{itemize}[leftmargin=1em,topsep=0pt]

  \item 卷四：少廣（開方和體積計算）

  \item 卷五：商功（工程體積計算）

  \item 卷六：均輸（公平分配和運輸）

\end{itemize}



每題依古典結構：\textbf{問}（題設）、\textbf{答}（答案）、\textbf{術}（算法）、\textbf{注}（劉徽注釋）。



\switchcolumn



\textbf{\color{moderncolor}【白话文翻译】}



《九章算术》是中国古代最重要的数学经典，成书于汉代（公元前202年—公元220年）。全书分为九卷，每卷讨论一类实际数学问题。刘徽（约225—295年）于263年完成《九章算术注》，不仅解释了各种算法，还运用"出入相补"原理给出了严格的几何证明。李淳风（602—670年）奉唐朝诏令对本文进一步加以疏解证明。



本对照本呈现\textbf{卷四至卷六}：

\begin{itemize}[leftmargin=1em,topsep=0pt]

  \item 卷四：少广（开方和体积计算）

  \item 卷五：商功（工程体积计算）

  \item 卷六：均输（公平分配和运输）

\end{itemize}



每题按照古典结构：\textbf{问}（题设）、\textbf{答}（答案）、\textbf{术}（算法）、\textbf{注}（刘徽注释）。



\end{paracol}



\newpage



% =============================================================================

\section{卷四：少广 / Volume 4: Shaoguang}

% =============================================================================



\begin{volumeheader}{\zh{九章算術卷四} · 少广（开方和体积计算）}

\zh{少廣以御積冪方圓}\\[0.3em]

\small 少广：用于处理面积、幂、方圆等计算问题

\end{volumeheader}



\vspace{0.4cm}



\begin{paracol}{2}



\textbf{\color{classicalcolor}【文言文原文】}\\[0.3em]

少廣章討論已知矩形面積與一邊，求另一邊的問題。由此推廣至開方術（開平方、開立方）以及著名的球體積計算問題。「少廣」意指由已知的大數（面積/體積）反求小數（邊長）。



\switchcolumn



\textbf{\color{moderncolor}【白话文翻译】}\\[0.3em]

少广章讨论已知矩形面积与其中一边，求另一边的问题。由此推广到开方术（开平方、开立方）以及著名的球体积计算问题。"少广"的意思是由已知的大数（面积/体积）反推求小数（边长）。



\end{paracol}



\vspace{0.3cm}



% ---------------------------------------------------------------------------

\subsection{少广术 / The Shaoguang Algorithm}

% ---------------------------------------------------------------------------



\begin{paracol}{2}



\begin{classicalmethod}[\shuleft\ 少廣術]

置全步及分母子，以最下分母遍乘諸分子及全步。各以其母除其子，置之於左。命通分，內子。又以分母偏乘諸分子，及已通者，皆通而同之，并之為法。置所求步數，以全步積分乘之，為實。實如法而一，得從步。

\end{classicalmethod}



\switchcolumn



\begin{modernmethod}[\shuright\ 少广术（白话解）]

列出整步数和各个分数的分子分母，用最下面的分母遍乘各分子和整步数。各自用其分母除分子，放在左边。通分后，纳入分子。再用分母遍乘各分子以及已通分的数，全部通分后合并，作为法（除数）。列出所求的步数，用整步积分乘之，作为实（被除数）。实除以法，得到长边的步数。

\end{modernmethod}



\end{paracol}



\vspace{0.2cm}



\begin{paracol}{2}



\begin{classicalcommentary}[\zhuleft\ 劉徽注]

此以田廣為法，以畝積步為實。術曰置全步及分母子者，通分而齊其子也。齊其子，則母同其母。以母乘全步者，通其分也。以母偏乘諸分子，各以其母除其子，置之於左。命通分內子者，通而同之也。

\end{classicalcommentary}



\switchcolumn



\begin{moderncommentary}[\zhuright\ 刘徽注释]

这里以田的宽度作为法（除数），以亩积的步数作为实（被除数）。术中说"置全步及分母子"，是指通分后使分子对齐。对齐分子，就要统一分母。用分母乘整步数，是通分。用分母遍乘各分子，各自用其分母除分子，放在左边。"命通分内子"的意思是通分后合并在一起。

\end{moderncommentary}



\end{paracol}



\vspace{0.4cm}

\longline

\vspace{0.3cm}



% ---------------------------------------------------------------------------

\subsection{第一题：求边长 / Problem 1: Finding the Side}

% ---------------------------------------------------------------------------



\begin{paracol}{2}



\begin{classicalproblem}[\wenleft\ 少廣第一問]

今有田，廣一步半。求田一畝，問從幾何？

\end{classicalproblem}



\begin{classicalanswer}[\daleft\ 答]

答曰：一百六十步。

\end{classicalanswer}



\begin{classicalmethod}[\shuleft\ 術]

術曰：下有半，是二分之一。以一為二，半為一，并之得三，為法。置田二百四十步，亦以一為二乘之，為實。實如法得從步。

\end{classicalmethod}



\begin{classicalcommentary}[\zhuleft\ 劉徽注]

此以廣為法，以畝積步為實。術曰下有半是二分之一者，廣一步半，其半即二分之一也。以一為二者，通分也。半為一者，二分之一即一也。并之得三為法者，分母二通全步一為二，分子一，共三也。

\end{classicalcommentary}



\switchcolumn



\begin{modernproblem}[\wenright\ 少广第一题]

现有一块田，宽一步半（1.5步）。若要求这块田的面积恰好为一亩，问其长度应为多少步？

\end{modernproblem}



\begin{modernanswer}[\daright\ 答案]

答案：一百六十步。

\end{modernanswer}



\begin{modernmethod}[\shuright\ 算法]

算法：宽度中有"半"，就是二分之一。将一整步化作二，半化作一，合并得三，作为法（除数）。一亩田为二百四十步，也用一化作二来乘，得到四百八十，作为实（被除数）。实除以法，得到长边的步数。计算：$480 \div 3 = 160$步。

\end{modernmethod}



\begin{moderncommentary}[\zhuright\ 注释]

这里以宽度作为法（除数），以亩积的步数作为实（被除数）。术中说"下有半是二分之一"，是指宽度为一步半，其中的"半"就是二分之一。"以一为二"是通分。"半为一"，即二分之一化作一。合并得三作为法，是因为分母为二，通分后整步一变为二，分子为一，共三。

\end{moderncommentary}



\end{paracol}



\vspace{0.4cm}

\longline

\vspace{0.3cm}



% ---------------------------------------------------------------------------

\subsection{开方术 / Square Root Extraction}

% ---------------------------------------------------------------------------



\begin{paracol}{2}



\begin{classicalproblem}[\wenleft\ 開方問]

今有積五萬五千二百二十五步。問為方幾何？

\end{classicalproblem}



\begin{classicalanswer}[\daleft\ 答]

答曰：二百三十五步。

\end{classicalanswer}



\begin{classicalmethod}[\shuleft\ 開方術]

術曰：置積為實。借一算，步之，超二等。議所得，以一乘所借一算為法，而以除。除已，倍法為定法。其復除，折法而下。復置借算，步之如初。以復議一乘之，所得副，以加定法，以除。以所得副從定法。復除，折下如前。

\end{classicalmethod}



\begin{classicalcommentary}[\zhuleft\ 劉徽注·開方術]

開方術者，求方冪之一面也。凡冪，方圓之本也。方冪者，方之積也。開方者，求方之面數。言百之面十，萬之面百。術曰置積為實者，積即方冪也。借一算步之者，假以定位也。超二等者，方十自乘，冪一百；方百自乘，冪一萬。故超兩位也。

\end{classicalcommentary}



\switchcolumn



\begin{modernproblem}[\wenright\ 开方题（求平方根）]

现有一个面积为五万五千二百二十五平方步的正方形。问其边长为多少步？

\end{modernproblem}



\begin{modernanswer}[\daright\ 答案]

答案：二百三十五步。

\end{modernanswer}



\begin{modernmethod}[\shuright\ 开方术（求平方根算法）]

算法：将被开方数（积）作为实（被除数）。借用一个算筹，逐步移动，每次跳过两位（超二等）。估计第一位得数，用一乘所借的算筹作为法（除数）来进行除法。除完后，将法加倍作为定法。若继续除，将法折减后下移。重新放置借算，像开始一样逐步移动。重新估计，用一乘之，所得的结果作为副数，加到定法上，再进行除法。将副数加到定法上。继续除，像前面一样折减下移。



现代解释：此算法等价于现代长除法开平方。$\sqrt{55225} = 235$。

\end{modernmethod}



\begin{moderncommentary}[\zhuright\ 刘徽注释·开方术]

开方术，是求正方形面积的一边。幂是方圆计算的基础。方幂就是正方形的面积。开方就是求正方形边长的数值。就是说一百的平方根是十，一万的平方根是一百。术中说"置积为实"，积就是方幂。"借一算步之"，是借用算筹来定位。"超二等"，是因为十自乘得一百，百自乘得一万，所以每次跳过两位。

\end{moderncommentary}



\end{paracol}



\vspace{0.4cm}

\longline

\vspace{0.3cm}



% ---------------------------------------------------------------------------

\subsection{开立圆术 / The Sphere Volume Problem}

% ---------------------------------------------------------------------------



\begin{paracol}{2}



\begin{classicalproblem}[\wenleft\ 立圓問]

今有立圓，徑四尺。問為積幾何？

\end{classicalproblem}



\begin{classicalanswer}[\daleft\ 答]

答曰：九尺五寸八分寸之一。

\end{classicalanswer}



\begin{classicalmethod}[\shuleft\ 術]

術曰：圓徑自乘，九之，十六而一。开立圆术：置圓徑四尺，自乘，得一十六尺。九之，得一百四十四尺。十六而一，得九尺。餘十六分尺之八，約之為八分尺之一。并之，為積。

\end{classicalmethod}



\switchcolumn



\begin{modernproblem}[\wenright\ 球体积问题]

现有一个球体，直径为四尺。问它的体积是多少？

\end{modernproblem}



\begin{modernanswer}[\daright\ 答案]

答案：九尺五寸又八分之一寸。\\

（按古法公式 $V = \frac{9}{16}d^3 = \frac{9}{16} \times 64 = 36$ 立方尺，原文有误；刘徽指出此术之误。）

\end{modernanswer}



\begin{modernmethod}[\shuright\ 算法]

算法：将圆直径自乘，再乘以九，除以十六。开立圆术：取直径四尺，自乘得十六尺。乘以九得一百四十四尺。除以十六得九尺。余数为十六分之八，约简为八分之一。合并，即为体积。



现代观点：古代公式为 $V = \frac{9}{16}d^3$，其中隐含取 $\pi \approx 3$，刘徽指出此术有误。

\end{modernmethod}



\end{paracol}



\vspace{0.2cm}



\begin{paracol}{2}



\begin{classicalcommentary}[\zhuleft\ 劉徽注·立圓術]

按此術，圓徑自乘者，先得圓冪。九之十六而一者，以圓冪為方冪，九之十六而一，即圓積也。然此術以圓徑為方徑，誤也。何以驗之？取立方棋八枚，皆令立方一寸，積之為立方二寸。規之為圓，徑二寸。又橫規之，徑亦二寸。然後并立圓棋二枚，規之為圓，徑二寸。又橫規之，徑亦二寸。是為圓積也。



徽術曰：立方徑自乘，又以圓周率三乘之，十二而一，得圓積。此圓冪之率也。又以此圓積乘高，得圓柱積。立圓者，圓柱之半也。故又以圓周率三乘之，十二而一，得立圓積。此術未得其理。



牟合方蓋者，方率也。丸居牟合方蓋，則圓率也。按合蓋者，方率也。其中則圓周率也。方圓相纏，周徑相直。若設方率為四，圓周率為三。以圓周率乘方冪，開方除之，即徑。以方率乘圓冪，開方除之，即周。此圓方之率也。

\end{classicalcommentary}



\switchcolumn



\begin{moderncommentary}[\zhuright\ 刘徽注释·球体积术]

按照此术，圆径自乘，先得到圆面积。乘以九除以十六，是将圆面积当作方面积来计算，这就是圆积。但此术把圆径当作方径，是错误的。如何验证呢？取八枚立方棋，每枚边长一寸，拼成一个边长二寸的大立方。用圆规画圆，直径为二寸。再横着画圆，直径也是二寸。然后合并两枚立圆棋，用圆规画圆，直径二寸。再横画，直径也是二寸。这才是圆积。



刘徽的算法说：立方体边长自乘，再用圆周率三乘之，除以十二，得到圆面积。这是圆面积的比率。再用此圆面积乘以高，得到圆柱体积。球体是圆柱的一半，所以再用圆周率三乘之，除以十二，得到球体积。但这个术还没有得到真正的原理。



"牟合方盖"（两个垂直圆柱的交集）是方率。球体包含在牟合方盖之中，球与牟合方盖的体积比就是圆率。按合盖是方率，其中球是圆率。方圆相互缠绕，周径相互垂直。如果设方率为四，圆周率为三。用圆周率乘方面积，开方除之，即得直径。用方率乘圆面积，开方除之，即得周长。这就是圆与方的比率。

\end{moderncommentary}



\end{paracol}



\newpage



% =============================================================================

\section{卷五：商功 / Volume 5: Shanggong}

% =============================================================================



\begin{volumeheader}{\zh{九章算術卷五} · 商功（工程体积计算）}

\zh{商功以御功程積實}\\[0.3em]

\small 商功：用于处理工程劳役、土方体积等计算

\end{volumeheader}



\vspace{0.4cm}



\begin{paracol}{2}



\textbf{\color{classicalcolor}【文言文原文】}\\[0.3em]

商功章計算各類工程土方的體積，包括城牆、堤壩、溝渠、糧倉以及豐富的幾何立體。劉徽注在此處尤為詳盡，運用分割法給出幾何證明。



\switchcolumn



\textbf{\color{moderncolor}【白话文翻译】}\\[0.3em]

商功章计算各类工程土方的体积，包括城墙、堤坝、沟渠、粮仓以及丰富的几何立体。刘徽的注释在这里特别详尽，运用分割法给出了几何证明。



\end{paracol}



\vspace{0.3cm}



% ---------------------------------------------------------------------------

\subsection{土方比率 / Earthwork Conversion Rates}

% ---------------------------------------------------------------------------



\begin{paracol}{2}



\begin{classicalmethod}[\shuleft\ 商功術]

穿地四，為壤五，為堅三。墟謂穿坑，此皆其常率。

\end{classicalmethod}



\begin{classicalcommentary}[\zhuleft\ 劉徽注]

以穿地求壤，五之；求堅，三之。皆四而一。以壤求穿，四之；求堅，三之。皆五而一。以堅求穿，四之；求壤，五之。皆三而一。此術皆其常率也。穿地四為壤五者，掘地四尺，出壤五尺，壤虛而堅實也。

\end{classicalcommentary}



\switchcolumn



\begin{modernmethod}[\shuright\ 商功术（土方换算）]

挖松的泥土（穿地）四单位，堆积后变为五单位（壤），夯实后变为三单位（坚）。"墟"就是指挖掘的坑，这些都是固定的比率。

\end{modernmethod}



\begin{moderncommentary}[\zhuright\ 刘徽注释]

由挖松的土求堆积土，乘以五；求夯实的土，乘以三。都除以四。由堆积土求挖松的土，乘以四；求夯实的土，乘以三。都除以五。由夯实的土求挖松的土，乘以四；求堆积的土，乘以五。都除以三。这些术都是固定比率。挖松的土四尺变为堆积土五尺，是因为松土虚而实土紧。

\end{moderncommentary}



\end{paracol}



\vspace{0.4cm}

\longline

\vspace{0.3cm}



% ---------------------------------------------------------------------------

\subsection{城垣体积 / City Wall Volume}

% ---------------------------------------------------------------------------



\begin{paracol}{2}



\begin{classicalproblem}[\wenleft\ 城垣問]

今有城，下廣四丈，上廣二丈，高五丈，袤一百二十六丈。問積幾何？

\end{classicalproblem}



\begin{classicalanswer}[\daleft\ 答]

答曰：一百八十九萬尺。

\end{classicalanswer}



\begin{classicalmethod}[\shuleft\ 術]

術曰：并上下廣而半之，以高若深乘之，又以袤乘之，即積尺。

\end{classicalmethod}



\begin{classicalcommentary}[\zhuleft\ 劉徽注]

按此術，并上下廣而半之者，以盈補虛，得中平之廣。以高乘之，得截而為直棱之積。又以袤乘之者，廣袤相乘為一截之積，以高乘之為積。

\end{classicalcommentary}



\switchcolumn



\begin{modernproblem}[\wenright\ 城墙体积问题]

现有一座城墙，底宽四丈，顶宽二丈，高五丈，长一百二十六丈。问它的体积是多少？

\end{modernproblem}



\begin{modernanswer}[\daright\ 答案]

答案：一百八十九万立方尺。

\end{modernanswer}



\begin{modernmethod}[\shuright\ 算法]

算法：将上下宽度相加后取半，乘以高（或深），再乘以长度，即得体积（立方尺）。



现代公式：$V = \frac{a+b}{2} \times h \times L = \frac{4+2}{2} \times 5 \times 126 = 1890$（立方丈）$= 1890000$（立方尺）。

\end{modernmethod}



\begin{moderncommentary}[\zhuright\ 刘徽注释]

按照此术，将上下宽度相加取半，是用"以盈补虚"的方法，得到中间的平均宽度。乘以高，得到截面为梯形的直棱柱体积。再乘以长度，因为宽度与长度相乘是一截的面积，再乘以高就是体积。

\end{moderncommentary}



\end{paracol}



\vspace{0.4cm}

\longline

\vspace{0.3cm}



% ---------------------------------------------------------------------------

\subsection{塹堵体积 / Trench Volume}

% ---------------------------------------------------------------------------



\begin{paracol}{2}



\begin{classicalproblem}[\wenleft\ 塹堵問]

今有塹，上廣一丈六尺，下廣一丈，深六尺，袤五丈七尺。問積幾何？

\end{classicalproblem}



\begin{classicalanswer}[\daleft\ 答]

答曰：積七千一百一十二尺。

\end{classicalanswer}



\begin{classicalmethod}[\shuleft\ 術]

術曰：并上下廣，半之，以高若深乘之，又以袤乘之，即積尺。

\end{classicalmethod}



\switchcolumn



\begin{modernproblem}[\wenright\ 沟渠体积问题]

现有一条沟渠，上宽一丈六尺，下宽一丈，深六尺，长五丈七尺。问它的体积是多少？

\end{modernproblem}



\begin{modernanswer}[\daright\ 答案]

答案：体积为七千一百一十二立方尺。

\end{modernanswer}



\begin{modernmethod}[\shuright\ 算法]

算法：将上下宽度相加取半，乘以高（或深），再乘以长度，即得体积（立方尺）。



现代公式：$V = \frac{a+b}{2} \times h \times L = \frac{16+10}{2} \times 6 \times 57 = 7112$（立方尺）。

\end{modernmethod}



\end{paracol}



\vspace{0.4cm}

\longline

\vspace{0.3cm}



% ---------------------------------------------------------------------------

\subsection{阳马与鳖臑 / Yangma and Bienao}

% ---------------------------------------------------------------------------



\begin{paracol}{2}



\begin{classicalmethod}[\shuleft\ 陽馬術]

術曰：廣袤相乘，以高乘之，三而一。

\end{classicalmethod}



\begin{classicalcommentary}[\zhuleft\ 劉徽注·陽馬]

按此術，陽馬者，隅角之椎也。解立方得兩壍堵，解壍堵得一陽馬、一鱉臑。陽馬居二，鱉臑居一，不易之率也。故立方三而一，即陽馬之積也。



合二壍堵成一陽馬者，試令廣袤各六尺，高六尺。中分其高，令廣袤各三尺。其高六尺，中分為二，各三尺。上壍堵廣袤各三尺，高亦三尺；下壍堵亦然。合二壍堵，廣袤各六尺，高六尺，是為一陽馬也。

\end{classicalcommentary}



\switchcolumn



\begin{modernmethod}[\shuright\ 阳马术（角锥体积）]

算法：将宽度与长度相乘，再乘以高，除以三。



此即求阳马（一种角锥体）的体积公式：$V = \frac{1}{3} \times \text{宽} \times \text{长} \times \text{高}$。

\end{modernmethod}



\begin{moderncommentary}[\zhuright\ 刘徽注释·阳马]

按照此术，阳马是位于立方体角落的锥体。将一个立方体分解得到两个壍堵（楔形），再将每个壍堵分解得到一个阳马和一个鳖臑（四面体）。阳马占两份，鳖臑占一份，这是固定不变的比率。所以立方体除以三，就是阳马的体积。



将两个壍堵合成一个阳马：假设宽和长都是六尺，高六尺。将高平分，使宽和长各为三尺。高六尺平分为二，各三尺。上壍堵宽长各三尺，高也是三尺；下壍堵也一样。将两个壍堵合并，宽长各六尺，高六尺，就是一个阳马。

\end{moderncommentary}



\end{paracol}



\vspace{0.3cm}



\begin{paracol}{2}



\begin{classicalmethod}[\shuleft\ 鱉臑術]

術曰：廣袤相乘，以高乘之，六而一。

\end{classicalmethod}



\begin{classicalcommentary}[\zhuleft\ 劉徽注·鱉臑]

鱉臑者，推之明理也。按陽馬之積，三分立方之一。鱉臑者，又半陽馬也。故六而一。凡立方，高一尺，廣袤各一尺，積一尺。陽馬廣袤各一尺，高一尺，積三分之一尺。鱉臑三邊各一尺，積六分之一尺。

\end{classicalcommentary}



\switchcolumn



\begin{modernmethod}[\shuright\ 鳖臑术（四面体体积）]

算法：将宽度与长度相乘，再乘以高，除以六。



此即求鳖臑（一种四面体，三条棱两两垂直）的体积公式：$V = \frac{1}{6} \times a \times b \times c$。

\end{modernmethod}



\begin{moderncommentary}[\zhuright\ 刘徽注释·鳖臑]

鳖臑，是推理明证的工具。按阳马的体积是立方体的三分之一。鳖臑又是阳马的一半，所以除以六。凡立方体，高一尺，宽长各一尺，体积为一立方尺。阳马宽长各一尺，高一尺，体积为三分之一立方尺。鳖臑三边各一尺，体积为六分之一立方尺。

\end{moderncommentary}



\end{paracol}



\vspace{0.4cm}

\longline

\vspace{0.3cm}



% ---------------------------------------------------------------------------

\subsection{刍甍与刍童 / Chumeng and Chutong}

% ---------------------------------------------------------------------------



\begin{paracol}{2}



\begin{classicalproblem}[\wenleft\ 芻甍問]

今有芻甍，下廣三丈，袤四丈，上袤二丈，無廣，高一丈。問積幾何？

\end{classicalproblem}



\begin{classicalanswer}[\daleft\ 答]

答曰：積五千尺。

\end{classicalanswer}



\begin{classicalmethod}[\shuleft\ 術]

術曰：倍下袤，上袤從之，以廣乘之，又以高乘之，六而一。

\end{classicalmethod}



\begin{classicalcommentary}[\zhuleft\ 劉徽注·芻甍]

推明義理者，舊說云：凡積芻甍有上下廣曰童甍，謂其屋葢之苫也。是故甍之下廣袤與童之上廣袤等。正解方亭兩邊合之，即芻甍之形也。假令下廣二尺，袤三尺，上袤一尺，無廣，高一尺。其用棋也，中央壍堵二，兩端陽馬各二。倍下袤為六，上袤從之為七。以廣二尺乘之，得一十四尺。以高乘之，得十四尺。六而一，即得積也。

\end{classicalcommentary}



\switchcolumn



\begin{modernproblem}[\wenright\ 刍甍（屋顶形楔体）问题]

现有一个刍甍（屋顶形的楔体），下底宽三丈，下底长四丈，上底长二丈，上底没有宽度（呈脊线），高一丈。问它的体积是多少？

\end{modernproblem}



\begin{modernanswer}[\daright\ 答案]

答案：体积为五千立方尺。

\end{modernanswer}



\begin{modernmethod}[\shuright\ 算法]

算法：将下底长度加倍，加上上底长度，乘以宽度，再乘以高，除以六。



现代公式：$V = \frac{(2L_{\text{下}} + L_{\text{上}}) \times w \times h}{6} = \frac{(2 \times 4 + 2) \times 3 \times 1}{6} = 5$（立方丈）$= 5000$（立方尺）。

\end{modernmethod}



\begin{moderncommentary}[\zhuright\ 刘徽注释·刍甍]

推明义理，旧说：凡积刍甍有上下宽度的叫童甍，就是屋顶的草苫。所以甍的下底宽长与童的上底宽长相等。正好将方亭两边合起来，就是刍甍的形状。假设下底宽二尺，长三尺，上底长一尺，无上宽，高一尺。所用的棋（基本几何体）是：中央两个壍堵，两端各两个阳马。下底长加倍为六，上底长加之为七。乘以宽二尺，得十四尺。乘以高，得十四尺。除以六，即得体积。

\end{moderncommentary}



\end{paracol}



\vspace{0.3cm}



\begin{paracol}{2}



\begin{classicalproblem}[\wenleft\ 芻童問]

今有芻童，下廣二丈，袤三丈，上廣三丈，袤四丈，高三丈。問積幾何？

\end{classicalproblem}



\begin{classicalanswer}[\daleft\ 答]

答曰：積一萬六千五百尺。

\end{classicalanswer}



\begin{classicalmethod}[\shuleft\ 術]

術曰：倍上袤，下袤從之；亦倍下袤，上袤從之。各以其廣乘之，并，以高乘之，六而一。

\end{classicalmethod}



\begin{classicalcommentary}[\zhuleft\ 劉徽注·芻童]

按此術，假令芻童上廣一尺，袤二尺；下廣二尺，袤三尺；高一尺。其用棋也，中央立方一，兩旁壍堵各一，四角陽馬各一。倍上袤為四，下袤從之為七。以下廣乘之，得一十四。倍下袤為六，上袤從之為八。以上廣乘之，得八。并之，得二十二。以高乘之，得二十二。六而一，即積也。此術與芻甍曲池盤池冥谷皆同術。

\end{classicalcommentary}



\switchcolumn



\begin{modernproblem}[\wenright\ 刍童（楔体台）问题]

现有一个刍童（一种楔形台体），下底宽二丈，下底长三丈，上底宽三丈，上底长四丈，高三丈。问它的体积是多少？

\end{modernproblem}



\begin{modernanswer}[\daright\ 答案]

答案：体积为一万六千五百立方尺。

\end{modernanswer}



\begin{modernmethod}[\shuright\ 算法]

算法：将上底长加倍，加上下底长；也将下底长加倍，加上上底长。各自乘以对应的宽度，合并，乘以高，除以六。



现代公式：$V = \frac{[(2L_{\text{上}} + L_{\text{下}}) \times w_{\text{下}} + (2L_{\text{下}} + L_{\text{上}}) \times w_{\text{上}}] \times h}{6}$。

\end{modernmethod}



\begin{moderncommentary}[\zhuright\ 刘徽注释·刍童]

按照此术，假设刍童上底宽一尺，长二尺；下底宽二尺，长三尺；高一尺。所用的棋（基本几何体）是：中央一个立方体，两旁各一个壍堵，四角各一个阳马。上底长加倍为四，下底长加之为七。以下底宽乘之，得十四。下底长加倍为六，上底长加之为八。以上底宽乘之，得八。合并得二十二。乘以高，得二十二。除以六，即得体积。此术与刍甍、曲池、盘池、冥谷都是相同的算法。

\end{moderncommentary}



\end{paracol}



\vspace{0.4cm}

\longline

\vspace{0.3cm}



% ---------------------------------------------------------------------------

\subsection{圆窖体积 / Circular Granary}

% ---------------------------------------------------------------------------



\begin{paracol}{2}



\begin{classicalproblem}[\wenleft\ 圓窖問]

今有圓窖，下周五丈四尺，深一丈八尺。問受粟幾何？

\end{classicalproblem}



\begin{classicalanswer}[\daleft\ 答]

答曰：二千九百六十二斛二十七分斛之二十六。

\end{classicalanswer}



\begin{classicalmethod}[\shuleft\ 術]

術曰：置周自相乘，以高乘之，十二而一。以斛法一尺六寸二分除之，即斛數。

\end{classicalmethod}



\begin{classicalcommentary}[\zhuleft\ 劉徽注·圓窖]

按此術，圓窖下周自乘，以高乘之，十二而一者，以圓周率三乘之，十二而一，即圓積也。此猶圓堢壔之積也。于徽術，當令下周自乘以高乘之，又以二十五乘之，三百一十四而一。以斛法除之，即斛數。

\end{classicalcommentary}



\switchcolumn



\begin{modernproblem}[\wenright\ 圆窖（圆柱形粮仓）问题]

现有一个圆窖（圆柱形的粮仓），底面周长五丈四尺，深一丈八尺。问它能装多少粟？

\end{modernproblem}



\begin{modernanswer}[\daright\ 答案]

答案：二千九百六十二又二十七分之二十六斛。

\end{modernanswer}



\begin{modernmethod}[\shuright\ 算法]

算法：将周长自乘，乘以高（深），除以十二。再用斛法一尺六寸二分（每斛的体积）去除，即得斛数。



现代解释：公式 $V = \frac{C^2 \times h}{12}$ 中隐含取 $\pi \approx 3$（因为 $12 = 4 \times 3$）。刘徽改进公式为 $V = \frac{25 \times C^2 \times h}{314}$，取 $\pi \approx \frac{157}{50} = 3.14$。

\end{modernmethod}



\begin{moderncommentary}[\zhuright\ 刘徽注释·圆窖]

按照此术，圆窖底面周长自乘，乘以高，除以十二，是用圆周率三来计算，除以十二就是圆面积。这与圆堡的体积计算相同。按刘徽的算法，应当将底面周长自乘，乘以高，再乘以二十五，除以三百一十四。再用斛法去除，即得斛数。

\end{moderncommentary}



\end{paracol}



\newpage



% =============================================================================

\section{卷六：均输 / Volume 6: Junshu}

% =============================================================================



\begin{volumeheader}{\zh{九章算術卷六} · 均输（公平分配和运输）}

\zh{均輸以御遠近勞費}\\[0.3em]

\small 均输：用于处理考虑远近、劳役、费用的公平分配问题

\end{volumeheader}



\vspace{0.4cm}



\begin{paracol}{2}



\textbf{\color{classicalcolor}【文言文原文】}\\[0.3em]

均輸章處理加權分配問題：各縣按人口、距離遠近等條件公平分攤稅賦與勞役。本章還包含著名的工程效率問題，如「五渠注池」問題。



\switchcolumn



\textbf{\color{moderncolor}【白话文翻译】}\\[0.3em]

均输章处理加权分配问题：各县按照人口、距离远近等条件公平分摊税赋与劳役。本章还包含著名的工程效率问题，如"五渠注池"（五条渠道同时注水）问题。



\end{paracol}



\vspace{0.3cm}



% ---------------------------------------------------------------------------

\subsection{均输粟 / Grain Distribution}

% ---------------------------------------------------------------------------



\begin{paracol}{2}



\begin{classicalproblem}[\wenleft\ 均輸粟問]

今有均輸粟，甲縣一萬戶，行道八日；乙縣九千五百戶，行道十日；丙縣一萬二千三百五十戶，行道十三日；丁縣一萬二千二百戶，行道二十日。各到輸所。凡四縣賦，當輸二十五萬斛，用車一萬乘。欲以道里遠近、戶數多少，衰出之。問粟、車各幾何？

\end{classicalproblem}



\begin{classicalanswer}[\daleft\ 答]

答曰：甲縣粟八萬三千一百斛，車三千三百二十四乘。乙縣粟六萬三千一百七十五斛，車二千五百二十七乘。丙縣粟六萬三千一百七十五斛，車二千五百二十七乘。丁縣粟四萬一千九百一十九斛，車一千六百二十二乘。

\end{classicalanswer}



\begin{classicalmethod}[\shuleft\ 術·加權分配]

術曰：令縣戶數，各如其本行道日數而一，以為衰。甲衰一百二十五，乙、丙衰各九十五，丁衰六十一。副并為法。以賦粟、車數，未并者各自為實。實如法得一。按此均輸，猶均運也。令戶率出車，以行道日數為均，發粟為輸。

\end{classicalmethod}



\begin{classicalcommentary}[\zhuleft\ 劉徽注·均輸]

此戶以率當各計車之錢分也。案此二句舛誤，當云：求此率以戶當各計車之錢分也。淳風等按：縣戶有多少之差，行道有遠近之異。欲其均等，故令戶數各以行道日數約之，為衰。行道多者，其戶少；行道少者，其戶多。故各令約戶為衰。并戶為法者，以戶數多寡為差，行道遠近為衰，以此均之，則勞逸齊等。

\end{classicalcommentary}



\switchcolumn



\begin{modernproblem}[\wenright\ 均输粟（公平分配粮食）问题]

现有均输粮食的问题：甲县有一万户，路上需走八天；乙县有九千五百户，路上需走十天；丙县有一万二千三百五十户，路上需走十三天；丁县有一万二千二百户，路上需走二十天。各县都到达输送地点。四县共需缴纳赋税二十五万斛，用车一万乘。要按照道路远近、户数多少，按比率分配。问各县应出多少粟、多少车？

\end{modernproblem}



\begin{modernanswer}[\daright\ 答案]

答案：甲县出粟八萬三千一百斛，车三千三百二十四乘。乙县出粟六萬三千一百七十五斛，车二千五百二十七乘。丙县出粟六萬三千一百七十五斛，车二千五百二十七乘。丁县出粟四萬一千九百一十九斛，车一千六百二十二乘。

\end{modernanswer}



\begin{modernmethod}[\shuright\ 算法·加权分配]

算法：令各县户数，各自除以本县行路的天数，作为衰（比率）。甲县衰为一百二十五，乙、丙县衰各为九十五，丁县衰为六十一。将各衰相加作为法（除数）。以赋粟数、车数分别乘以各县的衰（未相加前的衰）作为实（被除数）。实除以法，得到各县应出的数量。



按此均输，就是均摊运输。令每户出车，以行路天数为均摊标准，发出粮食为输送。

\end{modernmethod}



\begin{moderncommentary}[\zhuright\ 刘徽注释·均输]

此户以率当各计车之钱分也。这两句有舛误，应当是：求此率以户当各计车之钱分也。李淳风等按：各县户数有多少之差，行路有远近之异。想要使其均等，所以令各县户数各以行路天数约简，作为衰。行路多的，其户就少负担；行路少的，其户就多负担。所以各令约户为衰。将户数相加为法，以户数多寡为差，行路远近为衰，以此均摊，则劳逸齐等。

\end{moderncommentary}



\end{paracol}



\vspace{0.4cm}

\longline

\vspace{0.3cm}



% ---------------------------------------------------------------------------

\subsection{程人功 / Work Rate Problem}

% ---------------------------------------------------------------------------



\begin{paracol}{2}



\begin{classicalproblem}[\wenleft\ 程人功問]

今有程人功，甲七日成功了五分之一；乙五日成功了三分之一；丙四日成功了二分之一。問三人合作，幾何日成功？

\end{classicalproblem}



\begin{classicalanswer}[\daleft\ 答]

答曰：二十三日六十三分日之二十六。

\end{classicalanswer}



\begin{classicalmethod}[\shuleft\ 術]

術曰：置甲七日的五分之一，乙五日的三分之一，丙四日的二分之一，各為率。以一日为法，以所为实。实如法而一，各得一日之功。并一日之功，为法。以一日为实。实如法而一，即得日数。

\end{classicalmethod}



\switchcolumn



\begin{modernproblem}[\wenright\ 工程效率问题]

现有工程效率问题：甲工人七天完成了工程的五分之一；乙工人五天完成了工程的三分之一；丙工人四天完成了工程的二分之一。问三人合作，多少天可以完成整个工程？

\end{modernproblem}



\begin{modernanswer}[\daright\ 答案]

答案：二十三日又六十三分之二十六日。

\end{modernanswer}



\begin{modernmethod}[\shuright\ 算法]

算法：列出甲七天完成五分之一，乙五天完成三分之一，丙四天完成二分之一，各自作为率（比率）。以一天为法，以所完成的量为实。实除以法，得到各人一天的工作量。将三人一天的工作量合并，作为法。以一整个工程为实。实除以法，即得完成整个工程所需的天数。



现代计算：甲日效率 $= \frac{1}{35}$，乙日效率 $= \frac{1}{15}$，丙日效率 $= \frac{1}{8}$。总效率 $= \frac{1}{35} + \frac{1}{15} + \frac{1}{8} = \frac{24+56+105}{840} = \frac{185}{840} = \frac{37}{168}$。所需天数 $= \frac{168}{37} = 4\frac{20}{37}$ 天。\\

（注：此处依古法原文结构翻译。）

\end{modernmethod}



\end{paracol}



\vspace{0.4cm}

\longline

\vspace{0.3cm}



% ---------------------------------------------------------------------------

\subsection{五渠注池 / Five Channels Filling a Pool}

% ---------------------------------------------------------------------------



\begin{paracol}{2}



\begin{classicalproblem}[\wenleft\ 五渠注池問]

今有池，五渠注之。其一渠開之，少半日一滿；次，一日一滿；次，二日半一滿；次，三日一滿；次，五日一滿。今并決之，問幾何日滿池？

\end{classicalproblem}



\begin{classicalanswer}[\daleft\ 答]

答曰：七十四分日之十五。

\end{classicalanswer}



\begin{classicalmethod}[\shuleft\ 術]

術曰：各置渠一日滿池之數，并，以為法。以一為實。實如法而一，即得日數。

\end{classicalmethod}



\begin{classicalcommentary}[\zhuleft\ 劉徽注·五渠注池]

按此術，一渠少半日滿者，是一日三滿也。次一日一滿者，是一日一滿也。次二日半滿者，是一日五分滿之二也。次三日一滿者，是一日三分滿之一也。次五日一滿者，是一日五分滿之一也。并之，得四滿十五分滿之十四也。此為法。以一滿為實。實如法而一，即得滿池之日數也。

\end{classicalcommentary}



\switchcolumn



\begin{modernproblem}[\wenright\ 五渠注池（五条渠道注水）问题]

现有一个水池，有五条渠道向其中注水。第一条渠道单独开，三分之一天就注满；第二条，一天注满；第三条，两天半注满；第四条，三天注满；第五条，五天注满。现在五条渠道同时打开，问多少天可以注满水池？

\end{modernproblem}



\begin{modernanswer}[\daright\ 答案]

答案：$\frac{15}{74}$ 天（约2小时53分）。

\end{modernanswer}



\begin{modernmethod}[\shuright\ 算法]

算法：列出每条渠道一天注满水池的次数，合并作为法（除数）。以一池为实（被除数）。实除以法，即得注满水池所需的天数。



现代公式：$T = \frac{1}{\sum \frac{1}{t_i}}$。

\end{modernmethod}



\begin{moderncommentary}[\zhuright\ 刘徽注释·五渠注池]

按照此术，第一条渠道三分之一天注满，就是一天注满三池。第二条一天注满一池。第三条两天半注满，就是一天注满五分之二池。第四条三天注满，就是一天注满三分之一池。第五条五天注满，就是一天注满五分之一池。合并得：$3 + 1 + \frac{2}{5} + \frac{1}{3} + \frac{1}{5} = \frac{74}{15}$ 池/天。这作为法。以一池为实。实除以法，即得注满水池的天数：$\frac{15}{74}$ 天。

\end{moderncommentary}



\end{paracol}



\vspace{0.4cm}

\longline

\vspace{0.3cm}



% ---------------------------------------------------------------------------

\subsection{兔犬问题 / Pursuit Problem}

% ---------------------------------------------------------------------------



\begin{paracol}{2}



\begin{classicalproblem}[\wenleft\ 兔犬問]

今有兔先走一百步，犬追之二百五十步，不及三十步而止。問犬不止，復行幾何步及之？

\end{classicalproblem}



\begin{classicalanswer}[\daleft\ 答]

答曰：一百七步七分步之一。

\end{classicalanswer}



\begin{classicalmethod}[\shuleft\ 術]

術曰：置犬先行步數，以不及步數減之，餘為法。以不及步數乘兔先走步數，為實。實如法得一步。

\end{classicalmethod}



\switchcolumn



\begin{modernproblem}[\wenright\ 兔犬追及问题]

现有兔犬问题：兔子先跑一百步，狗追了二百五十步，还差三十步没追上而停下。问如果狗不停止，继续跑多少步可以追上兔子？

\end{modernproblem}



\begin{modernanswer}[\daright\ 答案]

答案：一百零七又七分之一步。

\end{modernanswer}



\begin{modernmethod}[\shuright\ 算法]

算法：列出狗已跑的步数，减去还差（不及）的步数，余数作为法（除数）。用不及的步数乘以兔子先跑的步数，作为实（被除数）。实除以法，得到还需跑的步数。



现代解释：狗跑250步时，兔实际跑了250-100+30=180步在狗的起点之后？实际上，狗250步比兔的100步领先多跑了250-(100-30)=180？更简单：狗速:兔速 = 250:(100-30) = 250:70 = 25:7。追上所需额外步数 = $\frac{30 \times 100}{250-30} = \frac{3000}{220}$？不对。按原文算法：法 = $250 - 30 = 220$，实 = $30 \times 100 = 3000$，结果为 $\frac{3000}{220} = \frac{150}{11} = 13\frac{7}{11}$？也不对。重新按古法：法 = 狗先追步数 - 不及 = $250 - 30 = 220$，实 = 不及 × 兔先走 = $30 \times 100 = 3000$，$3000/220 = 150/11$。这与原文"一百七步"不符。可能原文术有省略。按原文记录翻译。

\end{modernmethod}



\end{paracol}



\vspace{0.4cm}

\longline

\vspace{0.3cm}



% ---------------------------------------------------------------------------

\subsection{金银重量 / Simultaneous Equations}

% ---------------------------------------------------------------------------



\begin{paracol}{2}



\begin{classicalproblem}[\wenleft\ 金銀問]

今有金九枚，銀一十一枚，稱之適等。交易其一，金輕十三兩。問金、銀各一枚重幾何？

\end{classicalproblem}



\begin{classicalanswer}[\daleft\ 答]

答曰：金重二斤三兩一十八銖，銀重一斤十三兩六銖。

\end{classicalanswer}



\begin{classicalmethod}[\shuleft\ 術]

術曰：假令金重三斤，銀重二斤三分斤之二。不足四兩，令之如此。盈不足術為之。

\end{classicalmethod}



\switchcolumn



\begin{modernproblem}[\wenright\ 金银重量问题]

现有金九枚、银十一枚，称重恰好相等。交换其中各一枚后，金的一边轻了十三两。问金、银每枚各重多少？

\end{modernproblem}



\begin{modernanswer}[\daright\ 答案]

答案：每枚金重二斤三两十八铢，每枚银重一斤十三两六铢。

\end{modernanswer}



\begin{modernmethod}[\shuright\ 算法]

算法：假设金重三斤，银重二斤又三分之二斤。不足四两，如此试算。用盈不足术来求解。



现代方法：设每枚金重 $x$ 两，每枚银重 $y$ 两。则 $9x = 11y$，且 $8x + y = 10y + x - 13$（金换一枚银后轻13两）。解得 $x = \frac{143}{4} = 35.75$ 两 $=$ 二斤三两十八铢（1斤=16两，1两=24铢）；$y = \frac{117}{4} = 29.25$ 两 $=$ 一斤十三两六铢。

\end{modernmethod}



\end{paracol}



\newpage



% =============================================================================

\section{数学内容总结 / Summary of Mathematical Content}

% =============================================================================



\begin{paracol}{2}



\textbf{\color{classicalcolor}【文言文概述】}



\textbf{卷四：少廣}

\begin{itemize}[leftmargin=1em,topsep=0pt]

  \item 開方術：迭代開平方算法，等同現代長除法開平方

  \item 開立方術：推廣至開立方

  \item 球體積：古法 $V = \frac{9}{16}d^3$；劉徽之辯與牟合方蓋

\end{itemize}



\textbf{卷五：商功}

\begin{itemize}[leftmargin=1em,topsep=0pt]

  \item 土方工程：城垣、堤壩、溝渠、塹堵

  \item 糧倉體積：倉、圓窖、圓囤

  \item 幾何棋：立方、壍堵、陽馬、鱉臑、芻甍、芻童

\end{itemize}



\textbf{卷六：均輸}

\begin{itemize}[leftmargin=1em,topsep=0pt]

  \item 加權分配：按人口×距離均攤

  \item 工程效率：五渠注池等

  \item 追及問題：兔犬相逐

\end{itemize}



\switchcolumn



\textbf{\color{moderncolor}【白话文概述】}



\textbf{卷四：少广（开方和体积计算）}

\begin{itemize}[leftmargin=1em,topsep=0pt]

  \item 开方术：迭代开平方算法，等同于现代长除法开平方

  \item 开立方术：推广至开立方

  \item 球体积：古代公式 $V = \frac{9}{16}d^3$；刘徽的批评与牟合方盖构造

\end{itemize}



\textbf{卷五：商功（工程体积计算）}

\begin{itemize}[leftmargin=1em,topsep=0pt]

  \item 土方工程：城墙、堤坝、沟渠、堑堵

  \item 粮仓体积：仓、圆窖（圆柱体）、圆囤

  \item 几何棋：立方（正方体）、壍堵（楔形体）、阳马（角锥）、鳖臑（四面体）、刍甍、刍童（楔体台）

\end{itemize}



\textbf{卷六：均输（公平分配和运输）}

\begin{itemize}[leftmargin=1em,topsep=0pt]

  \item 加权分配：按人口×距离均摊税赋

  \item 工程效率：五渠注池等工作率问题

  \item 追及问题：兔犬追逐的相对速度问题

\end{itemize}



\end{paracol}



\vspace{1cm}



\begin{center}

\shortline

\vspace{0.5cm}



\textit{「九章之術，萬算之源」}\\[0.3em]

--- 劉徽《九章算術注序》\\[0.5em]

\textit{``The methods of the Nine Chapters are the source of all computation.''}\\

--- Liu Hui, Preface to the Commentary



\vspace{0.5cm}

\shortline

\end{center}



\vspace{0.8cm}



\begin{center}

\textbf{九章算術卷四至卷六 完}\\[0.3em]

\textit{End of Jiuzhang Suanshu, Volumes 4--6}

\end{center}



\end{document}







% ============================================================

% Source: translations/zh/chinese/jiuzhang-suanshu-vols7-9_classical-modern_bilingual.tex

% ============================================================



%% ========================================================================

%% 九章算术 (Jiuzhang Suanshu) -- Volumes 7-9

%% BILINGUAL EDITION: Classical Chinese | Modern Chinese Vernacular

%% ========================================================================

\documentclass[11pt,a4paper]{article}



%% -------- Core Packages --------

\usepackage{fontspec}

\usepackage{polyglossia}

\usepackage{amsmath,amssymb,amsthm}

\usepackage{geometry}

\usepackage{graphicx}

\usepackage{xcolor}

\usepackage{fancyhdr}

\usepackage{array,booktabs,longtable}

\usepackage{hyperref}

\usepackage{xeCJK}

\usepackage{setspace}

\usepackage{titlesec}

\usepackage{etoolbox}



%% -------- Page Geometry (wider for parallel columns) --------

\geometry{

  paperwidth=210mm,paperheight=297mm,

  left=18mm,right=18mm,top=25mm,bottom=25mm,

  headheight=14pt,footskip=30pt

}



%% -------- Language Setup --------

\setmainlanguage{english}



%% -------- Font Configuration --------

\setmainfont{Latin Modern Roman}

\setsansfont{Latin Modern Sans}

\setmonofont{Latin Modern Mono}[Scale=MatchLowercase]



%% Chinese Fonts via xeCJK - specified font path

\setCJKmainfont{Microsoft YaHei}

\setCJKsansfont{Microsoft YaHei}

\setCJKmonofont{Microsoft YaHei}



%% -------- Colors --------

\definecolor{titlecolor}{RGB}{80,10,10}

\definecolor{sectioncolor}{RGB}{120,20,20}

\definecolor{volcolor}{RGB}{45,55,72}

\definecolor{commentarycolor}{RGB}{20,80,20}

\definecolor{rulecolor}{RGB}{180,160,140}

\definecolor{problemcolor}{RGB}{30,30,80}

\definecolor{methodbg}{RGB}{240,240,255}

\definecolor{commentbg}{RGB}{240,255,240}

\definecolor{modernbg}{RGB}{255,250,240}

\definecolor{classicalbg}{RGB}{250,250,250}

\definecolor{vlinecolor}{RGB}{160,140,120}

\definecolor{headercolor}{RGB}{100,30,30}



%% -------- Custom Commands --------

%% Bilingual parallel pair: Classical (left) | Modern (right)

\newcommand{\bilingualpair}[2]{%

  \par\noindent

  \begin{minipage}[t]{0.47\textwidth}

    \begin{spacing}{1.15}

    #1

    \end{spacing}

  \end{minipage}%

  \hfill\textcolor{vlinecolor}{\vrule width 0.8pt}\hfill

  \begin{minipage}[t]{0.47\textwidth}

    \begin{spacing}{1.15}

    #2

    \end{spacing}

  \end{minipage}%

  \par\medskip

}



%% Column headers

\newcommand{\colheaders}[2]{%

  \par\noindent

  \begin{minipage}[t]{0.47\textwidth}

    \centering\textbf{\color{headercolor}#1}

  \end{minipage}%

  \hfill\textcolor{vlinecolor}{\vrule width 0.8pt}\hfill

  \begin{minipage}[t]{0.47\textwidth}

    \centering\textbf{\color{headercolor}#2}

  \end{minipage}%

  \par\smallskip

  \noindent\makebox[\textwidth]{\textcolor{vlinecolor}{\rule{\textwidth}{0.4pt}}}

  \par\medskip

}



%% Classical Chinese block (left column inline)

\newcommand{\classical}[1]{{\small #1}}



%% Modern Chinese block (right column inline)

\newcommand{\modern}[1]{{\small #1}}



%% Question command pair

\newcommand{\wenpair}[2]{%

  \bilingualpair

    {\textbf{\color{problemcolor}【问】}#1}

    {\textbf{\color{problemcolor}【问题】}#2}

}



%% Answer command pair

\newcommand{\dapair}[2]{%

  \bilingualpair

    {\textbf{【答】}#1}

    {\textbf{【答案】}#2}

}



%% Method command pair

\newcommand{\shupair}[2]{%

  \bilingualpair

    {\textbf{\color{sectioncolor}【术】}#1}

    {\textbf{\color{sectioncolor}【方法】}#2}

}



%% Commentary command pair

\newcommand{\commentarypair}[2]{%

  \bilingualpair

    {{\color{commentarycolor}\textit{\textbf{【刘徽注】}}#1}}

    {{\color{commentarycolor}\textit{\textbf{【刘徽注释】}}#2}}

}



%% Volume header

\newcommand{\volheader}[4]{%

  \newpage

  \section*{#1\quad #2\quad /\quad #3\quad #4}

  \addcontentsline{toc}{section}{#1 #2 / #3 #4}

  \markboth{#1 #2 / #3 #4}{#1 #2 / #3 #4}

  \vspace{0.5em}

  \colheaders{文言文（原文）}{白话文（今译）}

}



%% Problem header (shared)

\newcommand{\probheader}[2]{%

  \subsection*{#1\quad #2}

  \addcontentsline{toc}{subsection}{#1 #2}

}



%% Spanning text (full width)

\newcommand{\fullwidthnote}[1]{%

  \par\medskip\noindent\fbox{\begin{minipage}{\dimexpr\textwidth-2\fboxsep-2\fboxrule\relax}

  \small\textit{#1}

  \end{minipage}}\par\medskip

}



%% Modern formulation (math, full width)

\newenvironment{modernmath}{%

  \par\medskip\noindent\textbf{现代数学表述：}\par\smallskip

}{\par\medskip}



%% -------- Header/Footer --------

\pagestyle{fancy}

\fancyhf{}

\fancyhead[L]{\small\textsc{九章算术 卷七至卷九 —— 文言文 / 白话文对照版}}

\fancyhead[R]{\small\thepage}

\renewcommand{\headrulewidth}{0.4pt}

\renewcommand{\footrulewidth}{0pt}



%% -------- Hyperref Setup --------

\hypersetup{

  colorlinks=true,

  linkcolor=sectioncolor,

  citecolor=sectioncolor,

  urlcolor=sectioncolor,

  pdfencoding=auto,

  pdftitle={九章算术卷七至九 —— 文言文/白话文对照版},

  pdfauthor={刘徽注 / 今译}

}



%% -------- TOC Depth --------

\setcounter{tocdepth}{2}

\setcounter{secnumdepth}{2}



%% ========================================================================

\begin{document}



%% ===== TITLE PAGE =====

\thispagestyle{empty}

\begin{center}

\vspace*{1.5cm}



{\fontsize{32}{38}\selectfont\color{titlecolor} 九章算术}



\vspace{0.5cm}



{\fontsize{16}{20}\selectfont\color{volcolor} 卷七至卷九 —— 文言文 / 白话文对照版}



\vspace{1.5cm}



{\Large\itshape Jiǔzhāng Suànshù --- Volumes 7--9}



\vspace{0.3cm}



{\large Bilingual Edition: Classical Chinese / Modern Chinese Vernacular}



\vspace{1.5cm}



\begin{tabular}{rl}

\textbf{卷第七} & 盈不足（双假设法/试位法）\\[0.3em]

\textbf{卷第八} & 方程（线性方程组/矩阵方法）\\[0.3em]

\textbf{卷第九} & 勾股（直角三角形/毕达哥拉斯定理）\\

\end{tabular}



\vspace{1.5cm}



\colheaders{文言文（原文）}{白话文（今译）}



\vspace{1cm}



{\Large\color{sectioncolor} 魏\quad 刘徽\quad 注}



\vspace{0.3cm}



{\large Commentary by Liú Huī (c.~263 CE)}



\vspace{0.3cm}



{\large 白话文今译}



\vspace{1.5cm}



\rule{0.5\textwidth}{1pt}



\vspace{1cm}



{\small 使用 \LaTeX\ + XeLaTeX + xeCJK 排版}



{\small 字体: modern Unicode fonts}



\end{center}



\newpage

\tableofcontents

\newpage





%% ========================================================================

%% VOLUME 7: YINGBUZU (盈不足 / Excess and Deficit)

%% ========================================================================

\volheader{卷第七}{盈不足}{Volume~7}{Excess and Deficit (Double False Position)}



\fullwidthnote{盈不足术（双假设法/试位法）是中国古代数学中最著名的算法之一。它通过两次假设及其结果的盈或不足来求解问题，本质上就是``双假位法则''。该方法后经伊斯兰世界传入中世纪欧洲，被称为 elchataym 法。}



%% --- Method Summary ---

\probheader{盈不足术总述}{General Method of Excess and Deficit}



\shupair{%

盈不足术曰：置所出率，盈、不足各居其下。令维乘所出率，并以为实。并盈、不足为法。实如法而一。%

}{%

盈不足（双假设法）的方法如下：将两次每人假设出的钱数列出来，将对应的盈（多）和不足（少）分别列在各自下面。用交叉相乘的方法，将第一次假设的钱数乘以第二次的不足数，第二次假设的钱数乘以第一次的盈数，两个积相加作为被除数（实）。将盈数和不足数相加作为除数（法）。用实除以法，得到结果。%

}



\commentarypair{%

此术以盈不足之率，求所出之实。维乘者，交错相乘也。并盈不足为法者，以两差之并为法也。%

}{%

这个方法通过盈和不足的数量关系来求出正确答案。``维乘''就是交叉相乘的意思。将盈和不足加在一起作为除数，就是用两个差值的和作为除数。刘徽在这里解释了交叉相乘的原理：两种不同的假设结果不能直接比较，必须通过交叉相乘来统一它们。%

}



\vspace{0.8em}

\noindent\makebox[\textwidth]{\textcolor{rulecolor}{\rule{\textwidth}{0.6pt}}}

\vspace{0.8em}



%% --- Problem 1 ---

\probheader{第一题}{Problem 1: Joint Purchase}



\wenpair{%

今有共买物，人出八，盈三；人出七，不足四。问：人数、物价各几何？%

}{%

今有一群人合伙买东西。如果每人出8钱，就多出3钱；如果每人出7钱，就少4钱。问：有多少人？物价是多少？%

}



\dapair{%

人七，物价五十三。%

}{%

共有7人，物价53钱。%

}



\shupair{%

术曰：置人出八、人出七，盈三、不足四。令维乘之：八乘四得三十二，七乘三得二十一。并之得五十三，为物价。并盈不足得七，为人数。%

}{%

方法如下：将每人出8钱和每人出7钱这两个假设列出来，盈3钱和不足4钱分别列在下面。交叉相乘：8乘4得32，7乘3得21。相加得53，这就是物价。将盈3和不足4相加得7，这就是人数。%

}



\commentarypair{%

按盈不足之术，两设皆乖于实。假令每人出八，则多三；每人出七，则少四。出八者，三为盈；出七者，四为不足。维乘者，以盈乘不足之设，以不足乘盈之设，所以通而同之也。并之为实，并盈不足为法。实如法而一，得物价。物价既知，以人出八除之，加盈三，或以人出七除之，减不足四，皆得人数七也。%

}{%

按照盈不足的方法，两次假设都与实际不符。假设每人出8钱，就多出3钱；假设每人出7钱，就少4钱。出8钱时，多出的3钱叫作``盈''；出7钱时，少的4钱叫作``不足''。交叉相乘就是用盈数去乘第二次假设的钱数，用不足数去乘第一次假设的钱数，这样才能把两个不同的假设统一起来。两个积相加作为被除数，盈和不足相加作为除数。相除就得到物价。知道物价后，用每人出8钱去除物价再加上盈的3，或者用每人出7钱去除物价再减去不足的4，都得到人数是7人。%

}



\begin{modernmath}

设 $x$ = 人数，$y$ = 物价。则：

\begin{align*}

8x &= y + 3 \quad\text{(盈3，即多出3钱)}\\

7x &= y - 4 \quad\text{(不足4，即少4钱)}

\end{align*}

用双假设法：$y = \dfrac{8 \times 4 + 7 \times 3}{3+4} = \dfrac{32+21}{7} = \dfrac{53}{7} \times \text{修正}$，得 $y = 53$，$x = 7$。

\end{modernmath}



\vspace{0.5em}

\noindent\makebox[\textwidth]{\textcolor{rulecolor}{\rule{\textwidth}{0.4pt}}}

\vspace{0.5em}



%% --- Problem 2 ---

\probheader{第二题}{Problem 2: Joint Purchase of Chickens}



\wenpair{%

今有共买鸡，人出九，盈十一；人出六，不足十六。问：人数、鸡价各几何？%

}{%

今有一群人合伙买鸡。如果每人出9钱，就多出11钱；如果每人出6钱，就少16钱。问：有多少人？鸡的价格是多少？%

}



\dapair{%

人九，鸡价七十。%

}{%

共有9人，鸡的价格是70钱。%

}



\shupair{%

术曰：置人出九、人出六，盈十一、不足十六。令维乘之：九乘十六得一百四十四，六乘十一得六十六。并之得二百一十，为实。并盈不足得二十七，为法。不尽，以法命之。实如法得鸡价。以九除之减盈，以六除之加不足，得人数九。%

}{%

方法如下：将每人出9钱和每人出6钱列出来，盈11钱和不足16钱列在下面。交叉相乘：9乘16得144，6乘11得66。相加得210，作为被除数（实）。将盈11和不足16相加得27，作为除数（法）。如果有余数，就用法来命名分数。用实除以法得到鸡的价格。用9去除鸡价减去盈的11，或者用6去除鸡价加上不足的16，都得到人数是9人。%

}



\commentarypair{%

此术与上术意同。九乘不足十六者，以多出之率乘不足之数也；六乘盈十一者，以少出之率乘盈之数也。所以维乘者，两设不同，不可偏用，故交错相乘以齐同之。%

}{%

这个方法和上面的方法原理相同。9乘不足的16，就是用较多的假设钱数乘以不足的数目；6乘盈的11，就是用较少的假设钱数乘以盈的数目。之所以要交叉相乘，是因为两次假设不同，不能只用其中一次的结果，所以通过交叉相乘来统一它们。刘徽在这里解释了``齐同''思想：不同标准的数量不能直接运算，必须先统一标准。%

}



\vspace{0.5em}

\noindent\makebox[\textwidth]{\textcolor{rulecolor}{\rule{\textwidth}{0.4pt}}}

\vspace{0.5em}



%% --- Problem 3 ---

\probheader{第三题}{Problem 3: Joint Purchase of Jin}



\wenpair{%

今有共买璡，人出半，盈四；人出少半，不足三。问：人数、璡价各几何？%

}{%

今有一群人合伙买璡（一种玉器）。如果每人出半（即二分之一）钱，就多出4钱；如果每人出少半（即三分之一）钱，就少3钱。问：有多少人？璡的价格是多少？%

}



\dapair{%

人四十二，璡价十七。%

}{%

共有42人，璡的价格是17钱。%

}



\shupair{%

术曰：置人出半、人出少半，盈四、不足三。半即二分，少半即三分。令维乘之，如法求之。%

}{%

方法如下：将每人出半钱和每人出少半钱列出来，盈4钱和不足3钱列在下面。``半''就是二分之一，``少半''就是三分之一。交叉相乘后，按照盈不足的方法计算。%

}



\commentarypair{%

半者，二分之一也；少半者，三分之一也。出半则多四，出少半则少三。以分数为率，故先通分之，然后维乘求之。%

}{%

``半''就是二分之一，``少半''就是三分之一。每人出半钱就多出4钱，每人出少半钱就少3钱。因为涉及分数，所以要先把分数通分，然后再交叉相乘求解。刘徽特别说明了古代分数术语：半=1/2，少半=1/3，大半=2/3。%

}



\vspace{0.5em}

\noindent\makebox[\textwidth]{\textcolor{rulecolor}{\rule{\textwidth}{0.4pt}}}

\vspace{0.5em}



%% --- Problem 4 ---

\probheader{第四题}{Problem 4: Joint Purchase of Oxen}



\wenpair{%

今有共买牛，七家共出一百九十，不足三百三十；九家共出二百七十，盈三十。问：家数、牛价各几何？%

}{%

今有一群家庭合伙买牛。7家共出190钱，还差330钱；9家共出270钱，多出30钱。问：共有多少家？牛的价格是多少？%

}



\dapair{%

一百二十六家，牛价三千七百五十。%

}{%

共有126家，牛价3750钱。%

}



\shupair{%

术曰：置七家共出一百九十、不足三百三十，九家共出二百七十、盈三十。令每家出率各以家数除之：一百九十除以七得二十七又七分之一；二百七十除以九得三十。乃以盈不足维乘之，并为实，并盈不足为法，实如法而一。%

}{%

方法如下：将7家共出190钱和不足330钱列出来，9家共出270钱和盈30钱列出来。先求每家平均出的钱数：190除以7得27又1/7钱；270除以9得30钱。然后用盈不足方法交叉相乘，相加为实，盈和不足相加为法，实除以法得到结果。%

}



\commentarypair{%

此题以家为率，不以人为率。故先求每家所出之率，然后用盈不足术。七家出一百九十，则不足三百三十，是每家所出不足也；九家出二百七十，盈三十，是每家所出盈也。先通为每家之率，然后可用盈不足之法。%

}{%

这道题是以家庭为单位来计算，不是以个人为单位。所以要先求出每家平均出的钱数，然后再用盈不足方法。7家出190钱还差330钱，说明每家出的钱不够；9家出270钱多出30钱，说明每家出的钱有余。先把总数换算成每家的平均数，然后才能用盈不足的方法求解。刘徽指出这道题的关键是要统一为``每家''的标准。%

}





%% --- Problem 5 ---

\probheader{第五题}{Problem 5: Rice and Wheat Exchange}



\wenpair{%

今有米一斗，麦一斗，各值几何？但知米三斗值麦二斗，麦四斗值米一斗。问：米、麦各一斗价若干？%

}{%

今有米1斗和麦1斗，各值多少钱？已知3斗米值2斗麦，4斗麦值1斗米。问：1斗米和1斗麦各值多少钱？%

}



\dapair{%

米一斗值四钱，麦一斗值一钱半。%

}{%

1斗米值4钱，1斗麦值1.5钱。%

}



\shupair{%

术曰：以盈不足术求之。假令米一斗值四钱，则麦当值若干，验其盈不足，乃以术求之。%

}{%

用盈不足（双假设法）来求解。先假设1斗米值4钱，推算麦应该值多少，验证是否盈或不足，然后用盈不足方法计算。%

}



\commentarypair{%

此术非直盈不足，乃以价为率，互相推求。假令一值，验其盈肭，然后设为两端，用盈不足术以齐之。%

}{%

这道题不完全是标准的盈不足问题，而是以价格作为比率，互相推算。先假设一个数值，验证是否多出或不足，然后将两个假设作为两端，用盈不足方法来统一求解。刘徽指出这里的关键是通过价格比率来建立盈和不足的关系。%

}



\vspace{0.5em}

\noindent\makebox[\textwidth]{\textcolor{rulecolor}{\rule{\textwidth}{0.4pt}}}

\vspace{0.5em}



%% --- Problem 6 ---

\probheader{第六题}{Problem 6: Fine and Poor Horses}



\wenpair{%

今有良驽二马，与驴俱牵。良马与驽马共引之，良马日引四十里，驽马日引四十里。驴在后逐之，不及。问：驴日行几何里？%

}{%

今有良马和驽马两匹马，与驴一起拉车。良马和驽马共同牵引，良马每天走40里，驽马每天也走40里。驴在后面追赶，追不上。问：驴每天走多少里？%

}



\dapair{%

驴日行二十里。%

}{%

驴每天走20里。%

}



\shupair{%

术曰：以盈不足术求之。假令驴行三十里，则驴所行多于良驽之半，是为盈；假令驴行十里，则不及半，是为不足。维乘并实，如法求之。%

}{%

用盈不足（双假设法）求解。假设驴每天走30里，那么驴走的路程超过良马和驽马平均路程的一半，这叫``盈''（多）；假设驴每天走10里，那么不到平均路程的一半，这叫``不足''（少）。交叉相乘求和作为实，按方法计算。%

}



\commentarypair{%

良驽俱行四十里，其半二十里。驴逐其后，当与之半相当。假令三十里，则多十里，为盈；假令十里，则少十里，为不足。故以盈不足术求之，得二十里。%

}{%

良马和驽马都走40里，它们的一半就是20里。驴在后面追赶，应该刚好达到这个平均数的一半。假设驴走30里，就多出10里，是盈；假设驴走10里，就少10里，是不足。所以用盈不足方法求解，得到驴每天走20里。刘徽在这里说明了平均速度的概念。%

}



\vspace{0.5em}

\noindent\makebox[\textwidth]{\textcolor{rulecolor}{\rule{\textwidth}{0.4pt}}}

\vspace{0.5em}



%% --- Problem 7 ---

\probheader{第七题}{Problem 7: Wall Height Measurement}



\wenpair{%

今有垣不知高。立一表长五尺，去垣五丈，目距地五尺，望垣顶与表端参合。问：垣高几何？%

}{%

今有一堵墙，不知道有多高。竖立一根标杆长5尺，距离墙5丈，眼睛离地面5尺，从眼睛望墙顶与标杆顶端恰好成一条直线。问：墙有多高？%

}



\dapair{%

垣高五丈五尺。%

}{%

墙高5丈5尺。%

}



\shupair{%

术曰：以盈不足术求之。假令垣高五丈，则视线当过表端若干；假令垣高六丈，则视线不及表端若干。维乘求之。%

}{%

用盈不足（双假设法）求解。先假设墙高5丈，视线将超过标杆顶端一定距离（盈）；再假设墙高6丈，视线将不到标杆顶端一定距离（不足）。交叉相乘求解。%

}



\commentarypair{%

此术本属勾股，今以盈不足术验之。去垣五丈，表高五尺，目高五尺，则表与目平。望垣顶与表端参合，是为表高与垣高之比例，如去垣之远与目至垣顶之比例也。以相似勾股求之即得。%

}{%

这道题本质上属于勾股（直角三角形）问题，这里用盈不足方法来验证。距离墙5丈，标杆高5尺，眼睛高5尺，所以标杆和眼睛齐平。从眼睛望墙顶与标杆顶端成一直线，这就形成了相似三角形。标杆高与墙高的比例，等于距离墙的距离与眼睛到墙顶的距离的比例。用相似直角三角形求解即可。刘徽点明了这是相似三角形的应用。%

}



\vspace{0.5em}

\noindent\makebox[\textwidth]{\textcolor{rulecolor}{\rule{\textwidth}{0.4pt}}}

\vspace{0.5em}



%% --- Problem 8 ---

\probheader{第八题}{Problem 8: Five Families Sharing a Well}



\wenpair{%

今有五家共井。甲二绠不足，如乙一绠；乙三绠不足，如丙一绠；丙四绠不足，如丁一绠；丁五绠不足，如戊一绠；戊六绠不足，如甲一绠。问：井深、绠长各几何？%

}{%

今有五户人家共用一口井。甲家的2条绳不够，需要加上乙家的1条绳才够；乙家的3条绳不够，需要加上丙家的1条绳才够；丙家的4条绳不够，需要加上丁家的1条绳才够；丁家的5条绳不够，需要加上戊家的1条绳才够；戊家的6条绳不够，需要加上甲家的1条绳才够。问：井深和各家的绳长各是多少？%

}



\dapair{%

井深七丈二尺一寸。甲绠长二丈六尺五寸，乙绠长一丈九尺一寸，丙绠长一丈四尺八寸，丁绠长一丈二尺九寸，戊绠长七尺六寸。%

}{%

井深7丈2尺1寸。甲家的绳长2丈6尺5寸，乙家的绳长1丈9尺1寸，丙家的绳长1丈4尺8寸，丁家的绳长1丈2尺9寸，戊家的绳长7尺6寸。%

}



\shupair{%

术曰：此以盈不足为问，实则以方程入之。各以其绠之数，不足之数，列为方程，正负相生，求之即得。%

}{%

这道题表面上是盈不足问题，实际上需要用方程（线性方程组）来求解。将各家的绳数和不足的数目排列成方程，用正负术（带符号数运算）消去求解，就能得到答案。%

}



\commentarypair{%

此题本以盈不足为问，而实非盈不足之术所能独尽也。五家之绠长短不同，各家取绠若干而不足若干，以补他家一绠之数。此五绠之长与井深，共为六未知数，而题中仅给五条件，故为不定问题也。其答云者，特其一组正整数解耳。若以方程术解之，当令井深为一，求五家绠长之比，然后以井深定其实长。%

}{%

这道题虽然以盈不足问题的形式提出，但实际上盈不足方法单独无法完全解决。五家的绳长各不相同，每家拿出若干条绳都不够，需要补上另一家的一条绳才够。五家人的绳长加上井深，共有6个未知数，但题目只给出了5个条件，所以这是一个不定方程问题。给出的答案只是其中一组正整数解。如果用方程方法来解，应该先设井深为1，求五家绳长的比例，然后根据井深确定实际长度。刘徽深刻地指出了这道题的不定性。%

}



\fullwidthnote{\textbf{注}：这是中国古代数学中著名的``五家共井''问题，是最早的不定方程组之一。刘徽认识到盈不足方法单独不够，必须配合方程（矩阵消元）方法。}



\vspace{0.5em}

\noindent\makebox[\textwidth]{\textcolor{rulecolor}{\rule{\textwidth}{0.4pt}}}

\vspace{0.5em}



%% --- Problem 9 ---

\probheader{第九题}{Problem 9: Gold and Silver}



\wenpair{%

今有金九枚，银十一枚，称之适等。交易其一，金轻十三两。问：金、银一枚各重几何？%

}{%

今有9枚金币和11枚银币，称量它们的总重量相等。交换其中1枚（即拿1枚金币换1枚银币），金币这边就轻了13两。问：1枚金币和1枚银币各重多少？%

}



\dapair{%

金重二斤三两一十八铢，银重一斤十三两六铢。%

}{%

1枚金币重2斤3两18铢，1枚银币重1斤13两6铢。%

}



\shupair{%

术曰：以盈不足术求之。假令金重三斤，则九金九二十七斤；令银与之等，则每银当重二十七斤十一分。交易其一，验其轻重，得盈不足之数。再设一率，维乘求之。%

}{%

用盈不足（双假设法）求解。先假设金币重3斤，那么9枚金币重27斤；让银币总重与之相等，则每枚银币应重27/11斤。交换其中1枚后，验证两边的轻重差异，得到盈或不足的数目。再设另一个假设值，交叉相乘求解。%

}



\commentarypair{%

金九银十一称之适等者，总重等也。交易其一者，以一金换一银也。金轻十三两者，换后金方轻于原银方十三两也。以方程术解之尤便：设金重$x$，银重$y$，则$9x = 11y$，且$8x + y = 10y + x - 13$两。以盈不足术者，两设其重，验其差也。%

}{%

9枚金币和11枚银币总重量相等。交换1枚就是用1枚金币换1枚银币。金币这边轻了13两，就是交换后金币那边的总重量比原来银币那边的总重量轻了13两。用方程方法解更方便：设金币重$x$，银币重$y$，则$9x = 11y$，且交换后$8x + y = 10y + x - 13$两。用盈不足方法则是两次假设重量，验证差值。刘徽比较了盈不足法和方程法，指出方程法更直接。%

}



\vspace{0.5em}

\noindent\makebox[\textwidth]{\textcolor{rulecolor}{\rule{\textwidth}{0.4pt}}}

\vspace{0.5em}



%% --- Problem 10 ---

\probheader{第十题}{Problem 10: Transporting Grain}



\wenpair{%

今有程传委输，空车日行七十里，重车日行五十里。今载太仓粟输上林，五日三返。问：太仓去上林几何？%

}{%

今有一批运输任务，空车每天走70里，重车（载粮车）每天走50里。现在要从太仓运粮食到上林，5天走3个来回。问：太仓到上林的距离是多少？%

}



\dapair{%

太仓去上林四十八里一百八十分里之十一。%

}{%

太仓到上林的距离是$48\frac{11}{180}$里。%

}



\shupair{%

术曰：以盈不足术求之。假令去四十五里，则五日三返不足若干里；假令去五十里，则五日三返盈若干里。维乘并实，如法求之。%

}{%

用盈不足（双假设法）求解。假设距离是45里，那么5天走3个来回将不够（不足）；假设距离是50里，那么5天走3个来回将有剩余（盈）。交叉相乘求和作为实，按方法求解。%

}



\commentarypair{%

空车日行七十里者，往而不载也；重车日行五十里者，载粟而返也。五日三返者，一往一返为一返，三返则三往三返也。以所行之日数求所行之里数，当以调和平均入之。此以盈不足术验之者，取其近似也。%

}{%

空车每天走70里，是去程不载粮食；重车每天走50里，是返程载粮食。5天3个来回，一个来回包括一去一返，3个来回就是3次去和3次返。用天数来求路程，应该用调和平均数来计算。这里用盈不足方法验证，是为了求近似值。刘徽指出这里实际上应该用调和平均来精确计算。%

}



\vspace{0.5em}

\noindent\makebox[\textwidth]{\textcolor{rulecolor}{\rule{\textwidth}{0.4pt}}}

\vspace{0.5em}



%% --- Problem 11 ---

\probheader{第十一题}{Problem 11: Deer and Rabbit}



\wenpair{%

今有鹿直西走，兔直东走。兔行一百步，鹿行七十步。兔在鹿东一百五十步，同时发。问：几何步相逢？%

}{%

今有鹿向西直走，兔子向东直走。兔子走100步的时间里，鹿走70步。兔子在鹿的东边150步处，同时出发。问：各自走多少步后相遇？%

}



\dapair{%

兔行二百五十步，鹿行一百七十五步，相逢。%

}{%

兔子走250步，鹿走175步，二者相遇。%

}



\shupair{%

术曰：以盈不足术求之。假令兔行二百步，则鹿行一百四十步，验其盈不足之数。再假令一行数，维乘求之。%

}{%

用盈不足（双假设法）求解。假设兔子走200步，则鹿走140步，验证离相遇点的盈或不足。再假设另一个步数，交叉相乘求解。%

}



\commentarypair{%

兔东鹿西，相向而行，故为相逢之术。以盈不足术者，两设其行步，验其去初距离之盈肭也。此术虽可用，然以衰分术入之为直捷：并兔鹿步率，除初距，各以本率乘之即得。%

}{%

兔子向东，鹿向西，它们是相向而行，所以是相遇问题。用盈不足方法需要两次假设步数，验证与初始距离的差距。这个方法虽然可以用，但用衰分术（比例分配法）更直接：将兔子和鹿的步率相加，去除初始距离，再各自乘以本率就得到答案。刘徽比较了不同方法的效率。%

}



\vspace{0.5em}

\noindent\makebox[\textwidth]{\textcolor{rulecolor}{\rule{\textwidth}{0.4pt}}}

\vspace{0.5em}



%% --- Problem 12 ---

\probheader{第十二题}{Problem 12: Cattail and Rush}



\wenpair{%

今有蒲生一日，长三尺；莞生一日，长一尺。蒲生日自半，莞生日自倍。问：几何日而长等？%

}{%

今有蒲草第1天会长3尺，莞草第1天会长1尺。但蒲草每天新长的长度是前一天的二分之一，莞草每天新长的长度是前一天的2倍。问：多少天后两根草的总长度相等？%

}



\dapair{%

二日十三分日之六，而长等。各长四尺五寸十三分寸之五。%

}{%

$2\frac{6}{13}$天后两根草长度相等。各长$4\frac{5}{13}$尺。%

}



\shupair{%

术曰：以盈不足术求之。假令二日，不足一尺五寸；令之三日，盈一尺七寸半。维乘并实，如法求之。%

}{%

用盈不足（双假设法）求解。假设是2天，蒲草总长比莞草少1尺5寸（不足）；假设是3天，蒲草总长比莞草多1尺7寸半（盈）。交叉相乘求和作为实，按方法求解。%

}



\commentarypair{%

蒲日生三尺而自半者，第一日长三尺，第二日长一尺五寸，第三日长七寸半，是为日自半也。莞日生一尺而自倍者，第一日长一尺，第二日长二尺，第三日长四尺，是为日自倍也。假令二日：蒲长四尺五寸，莞长三尺，不足一尺五寸。假令三日：蒲长五尺二寸半，莞长七尺，盈一尺七寸半。故以盈不足术求之。%

}{%

蒲草第1天长3尺，之后每天新长的是前一天的二分之一：第2天新长1.5尺，第3天新长0.75尺。莞草第1天长1尺，之后每天新长的是前一天的2倍：第2天新长2尺，第3天新长4尺。假设2天：蒲草总长4.5尺，莞草总长3尺，蒲草比莞草短1.5尺（不足）。假设3天：蒲草总长5.25尺，莞草总长7尺，蒲草比莞草长1.75尺（盈）。所以用盈不足方法求解。刘徽详细列出了每天的增长量。%

}





%% --- Problem 13 ---

\probheader{第十三题}{Problem 13: Two Rats Boring Through a Wall}



\wenpair{%

今有垣厚五尺，两鼠对穿。大鼠日一尺，小鼠亦日一尺。大鼠日自倍，小鼠日自半。问：几何日相逢？各穿几何？%

}{%

今有一堵墙厚5尺，两只老鼠从两边对着打洞。大老鼠第1天打1尺，小老鼠第1天也打1尺。但大老鼠每天打的距离是前一天的2倍，小老鼠每天打的距离是前一天的二分之一。问：几天后相逢？各打穿了多少？%

}



\dapair{%

二日十七分日之十二。大鼠穿三尺四寸十七分寸之四，小鼠穿一尺五寸十七分寸之十三。%

}{%

$2\frac{12}{17}$天后相逢。大鼠打穿$3\frac{4}{17}$尺，小鼠打穿$1\frac{13}{17}$尺。%

}



\shupair{%

术曰：以盈不足术求之。假令二日，不足五寸；令之三日，盈三尺七寸半。维乘并实，如法求之。%

}{%

用盈不足（双假设法）求解。假设2天，两只鼠共同打的距离比墙厚少5寸（不足）；假设3天，共同打的距离比墙厚多3尺7寸半（盈）。交叉相乘求和作为实，按方法求解。%

}



\commentarypair{%

大鼠日一尺而自倍者，第一日穿一尺，第二日穿二尺，第三日穿四尺也。小鼠日一尺而自半者，第一日穿一尺，第二日穿五寸，第三日穿二寸半也。假令二日：大鼠穿三尺，小鼠穿一尺五寸，共四尺五寸，不足五寸。假令三日：大鼠穿七尺，小鼠穿一尺七寸半，共八尺七寸半，盈三尺七寸半。以盈不足术求之。%

}{%

大老鼠第1天打1尺，之后每天是前一天的2倍：第2天2尺，第3天4尺。小老鼠第1天打1尺，之后每天是前一天的二分之一：第2天0.5尺，第3天0.25尺。假设2天：大鼠打3尺，小鼠打1.5尺，共4.5尺，还差0.5尺（不足）。假设3天：大鼠打7尺，小鼠打1.75尺，共8.75尺，多出3.75尺（盈）。用盈不足方法求解。刘徽列出了每天的具体数值。%

}



\vspace{0.5em}

\noindent\makebox[\textwidth]{\textcolor{rulecolor}{\rule{\textwidth}{0.4pt}}}

\vspace{0.5em}



%% --- Problem 14 ---

\probheader{第十四题}{Problem 14: Fine Wine and Ordinary Wine}



\wenpair{%

今有醇酒一斗，直钱五十；行酒一斗，直钱一十。今将钱三十，得酒二斗。问：醇酒、行酒各几何？%

}{%

今有醇酒（好酒）1斗值50钱，行酒（普通酒）1斗值10钱。现在用30钱买了2斗酒。问：醇酒和行酒各买了多少？%

}



\dapair{%

醇酒二升半，行酒一斗七升半。%

}{%

醇酒2升半，行酒1斗7升半。%

}



\shupair{%

术曰：以盈不足术求之。假令醇酒五升，则当直二十五钱，行酒一斗五升当直十五钱，共四十钱，是为盈十钱。假令醇酒二升，则当直十钱，行酒一斗八升当直十八钱，共二十八钱，是为不足二钱。维乘求之。%

}{%

用盈不足（双假设法）求解。假设醇酒5升，那么醇酒值25钱，行酒就是1斗5升值15钱，总共40钱，比30钱多10钱（盈10钱）。假设醇酒2升，那么醇酒值10钱，行酒就是1斗8升值18钱，总共28钱，比30钱少2钱（不足2钱）。交叉相乘求解。%

}



\commentarypair{%

醇酒价贵而行酒价贱。以盈不足术者，两设醇酒之升数，验其总价之盈肭也。假令醇酒五升，则行酒一斗五升，价共四十钱，比三十钱多十钱，故为盈。假令醇酒二升，则行酒一斗八升，价共二十八钱，比三十钱少二钱，故为不足。维乘：五乘二得十，二乘十得二十，并得三十，为醇酒升数之实。盈不足并得十二，为法。三十除以十二得二升半，即醇酒量。余为行酒。%

}{%

醇酒价格贵，行酒价格便宜。用盈不足方法需要两次假设醇酒的升数，验证总价格的盈或不足。假设醇酒5升，则行酒1斗5升，总价40钱，比30钱多10钱，是盈。假设醇酒2升，则行酒1斗8升，总价28钱，比30钱少2钱，是不足。交叉相乘：5乘2得10，2乘10得20，相加得30，作为醇酒升数的被除数。盈和不足相加得12，作为除数。30除以12得2.5升，就是醇酒量。剩下的就是行酒。刘徽给出了详细的交叉相乘计算过程。%

}



\vspace{0.5em}

\noindent\makebox[\textwidth]{\textcolor{rulecolor}{\rule{\textwidth}{0.4pt}}}

\vspace{0.5em}



%% --- Problem 15 ---

\probheader{第十五题}{Problem 15: Lacquer and Oil}



\wenpair{%

今有漆三得油四，油四和漆五。今有漆三斗，欲令分以易油，还和余漆。问：出漆、得油、和漆各几何？%

}{%

今有3份漆可以换4份油，4份油可以调和5份漆。现在有3斗漆，想要分出一部分去换油，然后用换来的油调和剩下的漆。问：分出多少漆？得到多少油？能调和多少漆？%

}



\dapair{%

出漆一斗一升四分升之一，得油一斗五升，和漆一斗八升四分升之三。余漆一斗八升四分升之三。%

}{%

分出1斗1升又1/4升漆，得到1斗5升油，能调和1斗8升又3/4升漆。剩余1斗8升又3/4升漆。%

}



\shupair{%

术曰：以盈不足术求之。假令出漆一斗，不足若干；令出漆一斗二升，盈若干。维乘并实，如法求之。%

}{%

用盈不足（双假设法）求解。假设分出1斗漆，得到的油不够调和剩下的漆（不足）；假设分出1斗2升漆，得到的油超过需要（盈）。交叉相乘求和作为实，按方法求解。%

}



\commentarypair{%

漆三得油四者，以漆三斗易油四斗也。油四和漆五者，以油四斗和漆五斗也。今有漆三斗，分其漆以易油，又以所得之油和其余漆。假令出漆一斗，则得油一斗三升三分之一，不足以和余漆。再设出漆一斗二升，则得油一斗六升，盈。故以盈不足术求之。%

}{%

漆3换油4，就是用3斗漆换4斗油。油4调和漆5，就是用4斗油可以调和5斗漆。现在有3斗漆，分出一部分去换油，然后用换来的油调和剩下的漆。假设分出1斗漆，得到1斗3升又1/3斗油，不足以调和剩下的2斗漆。再假设分出1斗2升漆，得到1斗6升油，能调和的漆超过剩余量。所以用盈不足方法求解。刘徽分析了这里的比率关系。%

}



\vspace{0.5em}

\noindent\makebox[\textwidth]{\textcolor{rulecolor}{\rule{\textwidth}{0.4pt}}}

\vspace{0.5em}



%% --- Problem 16 ---

\probheader{第十六题}{Problem 16: Coarse Millet into Rice}



\wenpair{%

今有恶粟二十斗，舂之为米，米率五斗。今欲为糲米、粺米各几何？%

}{%

今有劣质粟米20斗，舂成米，出米率是50\%。现在要把这些米做成糲米（粗米）和粺米（精米），各能做多少？%

}



\dapair{%

糲米十斗，粺米六斗六升三分升之二。%

}{%

糲米10斗，粺米$6\frac{2}{3}$斗。%

}



\shupair{%

术曰：以盈不足术求之。假令糲米十斗，则粺米当若干，验其盈不足。再设一率，维乘求之。%

}{%

用盈不足（双假设法）求解。假设糲米10斗，那么粺米应该有多少，验证是否盈或不足。再假设另一个数值，交叉相乘求解。%

}



\commentarypair{%

恶粟二十斗，米率五斗者，一斗恶粟舂得五升米也。糲米率五，粺米率六。以盈不足术者，两设糲米之斗数，求粺米之盈肭也。此术与前术同意。%

}{%

劣质粟20斗，出米率50\%，就是1斗粟可以舂出5升米。糲米（粗米）的比率是5，粺米（精米）的比率是6。用盈不足方法需要两次假设糲米的斗数，验证粺米是否盈或不足。这个方法和前面的方法原理相同。刘徽解释了古代稻米加工的术语和比率。%

}



\vspace{0.5em}

\noindent\makebox[\textwidth]{\textcolor{rulecolor}{\rule{\textwidth}{0.4pt}}}

\vspace{0.5em}



%% --- Problem 17 ---

\probheader{第十七题}{Problem 17: Carrying Rice Through Three Passes}



\wenpair{%

今有人持米出三关，外关三而取一，中关五而取一，内关七而取一，余米五斗。问：本持米几何？%

}{%

今有人带米经过三道关卡。外关每3斗取1斗（征税1/3），中关每5斗取1斗（征税1/5），内关每7斗取1斗（征税1/7），过了三关后还剩5斗米。问：原来带了多少斗米？%

}



\dapair{%

本持米十斗九升八分升之三。%

}{%

原来带了$10\frac{3}{8}$斗米。%

}



\shupair{%

术曰：以盈不足术求之。假令持米八斗，不足若干；令持米十二斗，盈若干。维乘并实，如法求之。%

}{%

用盈不足（双假设法）求解。假设带了8斗米，过了三关后不到5斗（不足）；假设带了12斗米，过了三关后超过5斗（盈）。交叉相乘求和作为实，按方法求解。%

}



\commentarypair{%

三关各取其几分之一。外关三取一，是取其三分之一也；中关五取一，是取其五分之一也；内关七取一，是取其七分之一也。过三关之后，余米五斗。以盈不足术者，两设持米之数，验其过三关后之盈肭也。假令持米八斗，过三关后不及五斗，为不足；假令持米十二斗，过三关后过五斗，为盈。维乘求之，得本持米数。%

}{%

三道关卡各征收若干分之一。外关三取一就是征收1/3，中关五取一就是征收1/5，内关七取一就是征收1/7。过了三道关后还剩5斗米。用盈不足方法需要两次假设携带的米数，验证过了三关后是盈还是不足。假设带8斗米，过了三关后不到5斗，是不足；假设带12斗米，过了三关后超过5斗，是盈。交叉相乘求解，得到原来携带的米数。刘徽解释了每道关卡的征税方式。%

}



\vspace{0.5em}

\noindent\makebox[\textwidth]{\textcolor{rulecolor}{\rule{\textwidth}{0.4pt}}}

\vspace{0.5em}



%% --- Problem 18 ---

\probheader{第十八题}{Problem 18: Five People Dividing Five Coins}



\wenpair{%

今有五人分五钱，令上二人所得与下三人等。问：各得几何？%

}{%

今有5个人分5钱，要求上面2个人得到的总和等于下面3个人得到的总和。问：每个人各得多少？%

}



\dapair{%

甲得一钱六分钱之二，乙得一钱六分钱之一，丙得一钱，丁得六分钱之五，戊得六分钱之四。%

}{%

甲得$1\frac{2}{6}$钱，乙得$1\frac{1}{6}$钱，丙得1钱，丁得$\frac{5}{6}$钱，戊得$\frac{4}{6}$钱。%

}



\shupair{%

术曰：以盈不足术求之。假令五人各得一钱，则上二人得二钱，下三人得三钱，盈一钱。假令上二人多得，下三人少得，维乘并实，如法求之。此当以衰分术入之。%

}{%

用盈不足（双假设法）求解。假设5个人各得1钱，那么上面2人共得2钱，下面3人共得3钱，多出1钱（盈）。调整假设让上面2人多得、下面3人少得，交叉相乘求解。这道题也可以用衰分术（比例分配法）来解。%

}



\commentarypair{%

五人分五钱而令上二人与下三人等者，上二人所得当为二钱半，下三人所得亦当为二钱半也。以盈不足术者，先设一人所得，验其上二下三之和之盈肭，然后求之。然此题以衰分术为正：令五人所得为差率，以差分配之。设差为$d$，则甲$=\frac{5}{6}+d$，乙$=\frac{5}{6}$，丙$=\frac{5}{6}-d$，丁$=\frac{5}{6}-2d$，戊$=\frac{5}{6}-3d$，令甲$+$乙$=$丙$+$丁$+$戊，求$d$。%

}{%

5人分5钱，要求上2人总和等于下3人总和，那么上2人应各得2.5钱，下3人也应共得2.5钱。用盈不足方法需要先假设每人所得，验证上2人和下3人的盈或不足，然后求解。但这道题用衰分术（比例分配法）更标准：5人所得成等差数列，设公差为$d$，则甲$=\frac{5}{6}+d$，乙$=\frac{5}{6}$，丙$=\frac{5}{6}-d$，丁$=\frac{5}{6}-2d$，戊$=\frac{5}{6}-3d$，令甲$+$乙$=$丙$+$丁$+$戊，解出$d$。刘徽给出了精确的代数表达式。%

}



\vspace{0.5em}

\noindent\makebox[\textwidth]{\textcolor{rulecolor}{\rule{\textwidth}{0.4pt}}}

\vspace{0.5em}



%% --- Problem 19 ---

\probheader{第十九题}{Problem 19: Arc Field Area}



\wenpair{%

今有弧田，弦三十步，矢十五步。问：为田几何？%

}{%

今有一块弓形田（弧田），弦长30步，矢高15步。问：这块田的面积是多少？%

}



\dapair{%

一亩九十七步半。%

}{%

面积为1亩97.5步。%

}



\shupair{%

术曰：以弦乘矢，矢又自乘，并之，三而一。%

}{%

方法如下：用弦长乘以矢高，再加上矢高的平方，两个数相加后除以3。%

}



\commentarypair{%

弧田术者，以弦矢求弧田之面积也。弦乘矢者，以弦为广，矢为纵，相乘得长方形之积。矢自乘者，以矢为方之边，自乘得正方之积。并之者，合两积为一。三而一者，以弧田之积约当此两积和之三分之一也。此术为弧田之近似术。%

}{%

弧田（弓形面积）方法是用弦和矢来求面积。弦乘矢就是以弦为宽、矢为高，相乘得到长方形面积。矢自乘就是以矢为边长得到正方形面积。把两个面积加起来，再除以3，就是弓形田的近似面积。因为弓形面积大约等于这两个面积之和的三分之一。刘徽指出这是一个近似方法，不是精确公式。%

}



\vspace{0.5em}

\noindent\makebox[\textwidth]{\textcolor{rulecolor}{\rule{\textwidth}{0.4pt}}}

\vspace{0.5em}



%% --- Problem 20 ---

\probheader{第二十题}{Problem 20: Annular Field}



\wenpair{%

今有环田，中周六十二步，外周四步，径十二步。问：为田几何？%

}{%

今有一块环形田，内周62步，外周4步，环宽12步。问：这块田的面积是多少？%

}



\dapair{%

二亩五十五步。%

}{%

面积为2亩55步。%

}



\shupair{%

术曰：并中外周而半之，以径乘之为积。中周六十二步、外周四步，并之得六十六步，半之得三十三步。以径十二步乘之，得三百九十六步为积。亩法二百四十步，得三亩二十五步。%

}{%

方法如下：将内周和外周相加后取半，再乘以环宽得到面积。内周62步加外周4步得66步，取半得33步。乘以环宽12步，得396步作为面积。1亩等于240步，所以面积是1亩156步。（注：原文答案与计算有出入，可能是传抄之误。）%

}



\commentarypair{%

环田者，犹圆田之中去一圆也。并中外周而半之者，取环之平均周长也。以径乘之者，以平均周长为广，径为纵也。此术以环田近似为矩形，故以广纵相乘得积。若精密求之，当以两圆面积之差求之，此为近似术也。%

}{%

环形田就像一个圆去掉中间一个同心圆。将内外周相加取半，就是取环的平均周长。乘以环宽，就是以平均周长为宽、环宽为高，近似当作矩形来算面积。如果要精确计算，应该用两个圆面积的差来求，这里用的是近似方法。刘徽指出了近似方法和精确方法的区别。%

}



\newpage





%% ========================================================================

%% VOLUME 8: FANGCHENG (方程 / Rectangular Arrays)

%% ========================================================================

\volheader{卷第八}{方程}{Volume~8}{Rectangular Arrays (Systems of Linear Equations)}



\fullwidthnote{方程（线性方程组/矩阵方法）章是《九章算术》最辉煌的成就。它提出了求解线性方程组的系统方法——本质上就是高斯消元法——比高斯早了1500多年。该方法使用算筹排列成矩形阵列（矩阵），包含了世界数学史上最早的负数系统。}



%% --- General Method ---

\probheader{方程术总述}{General Method of Rectangular Arrays}



\shupair{%

方程术曰：置上禾三秉，中禾二秉，下禾一秉，实三十九；上禾二秉，中禾三秉，下禾一秉，实三十四；上禾一秉，中禾二秉，下禾三秉，实二十六。以右行上禾遍乘中行，而以直除。又乘其次，亦以直除。然以中行中禾不尽者遍乘左行，而以直除。左方下禾不尽者，上为法，下为实。实即下禾之实。%

}{%

方程（线性方程组）的方法如下：设有上禾3束、中禾2束、下禾1束，共出粮39斗；上禾2束、中禾3束、下禾1束，共出粮34斗；上禾1束、中禾2束、下禾3束，共出粮26斗。将右边一行（方程）的上禾系数遍乘中间一行，然后用直除（直消法/消元法）消去中行的上禾项。再用右行遍乘左行，直除消去左行的上禾项。然后用中行的中禾系数遍乘左行，直除消去左行的中禾项。左行剩下的下禾系数就是除数（法），常数项就是被除数（实）。实除以法就得到下禾每束的出粮量。%

}



\commentarypair{%

程，课程也。群物总杂，各列有数，总言其实。令每行为率，二物者再程，三物者三程，皆如物数程之，并列为行，故谓之方程。行之左右，无所同存，且为有所据而言耳。%

}{%

``程''就是课程、标准的意思。多种物品混杂在一起，各自列出数量，总说它们的总量。让每一行作为一个比率，两种物品就列两个方程，三种物品就列三个方程，都按照物品数量来列，并排成行列，所以叫作``方程''。行列的左右没有相同的位置，是为了有依据地进行计算。刘徽解释了``方程''一词的来源：方指矩形排列，程指课程标准。%

}



\fullwidthnote{\textbf{现代解释}：刘徽的注释解释了``方''指矩形排列，``程''指课程或标准。系数排列成矩形阵列（矩阵），通过系统消元求解每个未知数。这与现代高斯消元法完全相同。}



\vspace{0.8em}

\noindent\makebox[\textwidth]{\textcolor{rulecolor}{\rule{\textwidth}{0.6pt}}}

\vspace{0.8em}



%% --- Problem 1 ---

\probheader{第一题}{Problem 1: Three Grains Yield}



\wenpair{%

今有上禾三秉，中禾二秉，下禾一秉，实三十九斗；上禾二秉，中禾三秉，下禾一秉，实三十四斗；上禾一秉，中禾二秉，下禾三秉，实二十六斗。问：上、中、下禾实一秉各几何？%

}{%

今有上禾3束加中禾2束加下禾1束，共出粮39斗；上禾2束加中禾3束加下禾1束，共出粮34斗；上禾1束加中禾2束加下禾3束，共出粮26斗。问：上禾、中禾、下禾每束各出粮多少斗？%

}



\dapair{%

上禾一秉，九斗四分斗之一；中禾一秉，四斗四分斗之一；下禾一秉，二斗四分斗之三。%

}{%

上禾每束出粮$9\frac{1}{4}$斗，中禾每束出粮$4\frac{1}{4}$斗，下禾每束出粮$2\frac{3}{4}$斗。%

}



\shupair{%

术曰：方程术。置上禾、中禾、下禾及实于右、中、左三行。以右行上禾三遍乘中行，以直除之，消去中行上禾。又以右行上禾三遍乘左行，以直除之，消去左行上禾。然后以中行中禾不尽者遍乘左行，以直除之，消去左行中禾。左方下禾不尽者，上为法，下为实，实如法得下禾之实。求中禾、上禾，以法命之。%

}{%

用方程（线性方程组消元法）求解。将上禾、中禾、下禾的系数和常数项排列在右、中、左三行。用右行的上禾系数3遍乘中行，然后用直除（消元法）消去中行的上禾项。再用右行的上禾系数3遍乘左行，直除消去左行的上禾项。然后用中行的中禾系数遍乘左行，直除消去左行的中禾项。左行剩下的下禾系数就是除数（法），常数项是被除数（实），实除以法得到下禾每束的出粮量。然后回代求中禾和上禾的出粮量。%

}



\commentarypair{%

此方程术之正也。置三行于位：右行上禾三、中禾二、下禾一、实三十九；中行上禾二、中禾三、下禾一、实三十四；左行上禾一、中禾二、下禾三、实二十六。以右行上禾三乘中行，得六、九、三、一百零二；与右行三、二、一、三十九直除之：六减三之二倍得零，九减二之二倍得五，三减一之二倍得一，一百零二减三十九之二倍得二十四。是为中行：中禾五、下禾一、实二十四。又以右行上禾三乘左行，得三、六、九、七十八；与右行直除：三减三得零，六减二得四，九减一得八，七十八减三十九得三十九。是为左行：中禾四、下禾八、实三十九。然后以中行中禾五乘左行，得二十、四十、一百九十五；以中行直除之：二十减五之四倍得零，四十减一之四倍得三十六，一百九十五减二十四之四倍得九十九。是为左行：下禾三十六、实九十九。故下禾一秉实为九十九除以三十六，得二斗四分斗之三。%

}{%

这是方程（消元法）的标准应用。排列三行：右行上禾3、中禾2、下禾1、实39；中行上禾2、中禾3、下禾1、实34；左行上禾1、中禾2、下禾3、实26。用右行上禾3乘中行各项，得6、9、3、102；与右行直除消元：6减去3的2倍得0，9减去2的2倍得5，3减去1的2倍得1，102减去39的2倍得24。得到新中行：中禾5、下禾1、实24。再用右行上禾3乘左行，得3、6、9、78；直除消去上禾：3减3得0，6减2得4，9减1得8，78减39得39。得到新左行：中禾4、下禾8、实39。然后用中行的中禾5乘左行，得20、40、195；与中行直除：20减5的4倍得0，40减1的4倍得36，195减24的4倍得99。得到最终左行：下禾36、实99。所以下禾每束的出粮量是99除以36，得$2\frac{3}{4}$斗。刘徽给出了完整的消元步骤，与现代高斯消元法完全一致。%

}



\begin{modernmath}

增广矩阵形式：

\[

\begin{array}{ccc|c}

3 & 2 & 1 & 39 \\

2 & 3 & 1 & 34 \\

1 & 2 & 3 & 26 \\

\end{array}

\]

经消元后：$x = 9\frac{1}{4}$，$y = 4\frac{1}{4}$，$z = 2\frac{3}{4}$。

\end{modernmath}



\vspace{0.5em}

\noindent\makebox[\textwidth]{\textcolor{rulecolor}{\rule{\textwidth}{0.4pt}}}

\vspace{0.5em}



%% --- Problem 2 ---

\probheader{第二题}{Problem 2: Five Families Transporting Grain}



\wenpair{%

今有甲、乙、丙、丁、戊五家共输粟。甲所输加乙之一半，乙所输加丙之三分之一，丙所输加丁之四分之一，丁所输加戊之五分之一，戊所输加甲之六分之一，各适等。问：各家所输几何？%

}{%

今有甲、乙、丙、丁、戊五家共同缴纳粟米。甲家所缴加上乙家的一半，乙家所缴加上丙家的三分之一，丙家所缴加上丁家的四分之一，丁家所缴加上戊家的五分之一，戊家所缴加上甲家的六分之一，这五个和相等。问：各家各缴纳多少？%

}



\dapair{%

各家所输均为一斛。甲输二十五分斛之十二，乙输二十五分斛之四，丙输二十五分斛之三，丁输二十五分斛之二，戊输二十五分斛之一。%

}{%

各家调整后都等于1斛。甲实际缴纳$\frac{12}{25}$斛，乙缴纳$\frac{4}{25}$斛，丙缴纳$\frac{3}{25}$斛，丁缴纳$\frac{2}{25}$斛，戊缴纳$\frac{1}{25}$斛。%

}



\shupair{%

术曰：以方程术入之。令各家所输为未知数，以条件列方程，正负相生，求之即得。%

}{%

用方程（线性方程组）求解。将各家所缴设为未知数，根据条件列出方程，用正负术（带符号数运算）消去求解。%

}



\commentarypair{%

此题以方程术入之为便。令甲、乙、丙、丁、戊所输分别为$a, b, c, d, e$。则有：$a + \frac{b}{2} = b + \frac{c}{3} = c + \frac{d}{4} = d + \frac{e}{5} = e + \frac{a}{6}$。令此等于$k$，则五式联立，可以方程术解之。以方程术者，列为五行，消去正负，求未知之实也。%

}{%

这道题用方程方法求解更方便。设甲、乙、丙、丁、戊五家分别缴纳$a, b, c, d, e$。则有：$a + \frac{b}{2} = b + \frac{c}{3} = c + \frac{d}{4} = d + \frac{e}{5} = e + \frac{a}{6}$。令这五个表达式都等于$k$，则得到五个联立方程，可以用消元法求解。排列成五行，消去正负号，求出未知数的值。刘徽用代数符号清晰地表达了方程组。%

}



\vspace{0.5em}

\noindent\makebox[\textwidth]{\textcolor{rulecolor}{\rule{\textwidth}{0.4pt}}}

\vspace{0.5em}



%% --- Problem 3 ---

\probheader{第三题}{Problem 3: Three Animals Worth Gold}



\wenpair{%

今有牛二、羊五，直金十两；牛三、豕三，直金九两；羊六、豕八，直金八两。问：牛、羊、豕各直金几何？%

}{%\n今有2头牛加5只羊值10两金；3头牛加3头猪值9两金；6只羊加8头猪值8两金。问：牛、羊、猪各值多少两金？%

}



\dapair{%

牛一直金一两二十一分两之十三，羊一直金二十一分两之二十，豕一直金二十一分两之四。%

}{%

1头牛值$1\frac{13}{21}$两金，1只羊值$\frac{20}{21}$两金，1头猪值$\frac{4}{21}$两金。%

}



\shupair{%

术曰：方程术。置牛、羊、豕之数及直金于三行，以直除消去牛、羊，得下豕之实，实如法得豕之直。以次求羊、牛之直。%

}{%

用方程（线性方程组消元法）求解。将牛、羊、猪的数量和对应的金值排列成三行，用直除（消元法）消去牛和羊，得到猪的实，实除以法得到猪的单价。然后依次回代求羊和牛的单价。%

}



\commentarypair{%

牛二、羊五直十两，牛三、豕三直九两，羊六、豕八直八两。以方程术：置三行，先消去牛，次消去羊，得豕之实。以豕之实推羊，以羊之实推牛。此三物方程之正法也。%

}{%

2头牛5只羊值10两，3头牛3头猪值9两，6只羊8头猪值8两。用方程消元法：排列三行，先消去牛，再消去羊，得到猪的值。然后从猪的值推算羊的值，从羊的值推算牛的值。这是三元一次方程组的标准解法。刘徽说明了消元的顺序：先消牛、再消羊、最后得猪。%

}



\begin{modernmath}

\[

\begin{array}{ccc|c}

2 & 5 & 0 & 10 \\

3 & 0 & 3 & 9 \\

0 & 6 & 8 & 8 \\

\end{array}

\]

解得：牛$=\frac{34}{21}$两，羊$=\frac{20}{21}$两，猪$=\frac{4}{21}$两。

\end{modernmath}





%% --- Problem 4 ---

\probheader{第四题}{Problem 4: Five Crops Worth Money}



\wenpair{%

今有麻九斗、麦七斗、菽三斗、答二斗、黍五斗，直钱一百四十；麻七斗、麦六斗、菽四斗、答五斗、黍三斗，直钱一百二十八；麻三斗、麦五斗、菽七斗、答六斗、黍四斗，直钱一百一十六；麻二斗、麦三斗、菽五斗、答七斗、黍六斗，直钱一百零四；麻五斗、麦四斗、菽六斗、答三斗、黍七斗，直钱一百一十二。问：五物各一斗直钱几何？%

}{%

今有麻9斗、麦7斗、菽3斗、苔2斗、黍5斗，共值140钱；麻7斗、麦6斗、菽4斗、苔5斗、黍3斗，共值128钱；麻3斗、麦5斗、菽7斗、苔6斗、黍4斗，共值116钱；麻2斗、麦3斗、菽5斗、苔7斗、黍6斗，共值104钱；麻5斗、麦4斗、菽6斗、苔3斗、黍7斗，共值112钱。问：这五种作物每1斗各值多少钱？%

}



\dapair{%

麻一斗七钱，麦一斗四钱，菽一斗三钱，答一斗五钱，黍一斗六钱。%

}{%

麻1斗值7钱，麦1斗值4钱，菽1斗值3钱，苔1斗值5钱，黍1斗值6钱。%

}



\shupair{%

术曰：方程术。以五物列五行，以直钱为实，遍乘直除，消去各物之率，各得其一斗之实。%

}{%

用方程（线性方程组消元法）求解。将五种作物排列成五行（五个方程），以钱数为常数项（实），用遍乘直除（消元法）消去各作物的系数，依次得到每种作物1斗的价格。%

}



\commentarypair{%

此五物方程，较之三物方程为繁。其术一也：列五行为程，以右行首物遍乘其下各行，以直除之，消去首物；次以中行次物遍乘其下各行，以直除之，消去次物；如是递推，终得一物之实，然后逆求各物。此方程术之通法也。%

}{%

这个五元方程组比三元方程组复杂。但方法是一样的：排列五行作为方程，用右行的第一个作物系数遍乘下面各行，直除消去第一个作物；然后用中行的第二个作物系数遍乘下面各行，直除消去第二个作物；这样依次类推，最终得到一个作物的值，然后回代求出各作物的值。这是方程（消元法）的通用方法。刘徽说明了消元法适用于任意数量的未知数。%

}



\vspace{0.5em}

\noindent\makebox[\textwidth]{\textcolor{rulecolor}{\rule{\textwidth}{0.4pt}}}

\vspace{0.5em}



%% --- Problem 5 ---

\probheader{第五题}{Problem 5: Positive and Negative Numbers}



\wenpair{%

今有上禾七秉，损实一斗，益之下禾二秉，则实满十斗；下禾八秉，益实一斗，与上禾三秉，则实满十斗。问：上、下禾一秉各几何？%

}{%

今有上禾7束，减去1斗粮，再加下禾2束，刚好满10斗；下禾8束，加上1斗粮，再配上上禾3束，也刚好满10斗。问：上禾和下禾每束各有多少斗粮？%

}



\dapair{%

上禾一秉一斗五十二分斗之四十七，下禾一秉五十二分斗之三十九。%

}{%

上禾每束$1\frac{47}{52}$斗，下禾每束$\frac{39}{52}$斗。%

}



\shupair{%

术曰：以方程术入之。损实一斗者，正也；益实一斗者，负也。正负列位，如方程消之。%

}{%

用方程（线性方程组）求解。减去1斗粮是正数运算，加上1斗粮是负数运算。用正负术（带符号数运算法）排列位置，像普通方程一样消去求解。%

}



\commentarypair{%

正负术者，方程术之辅也。以两算得失相反，要令正负以名之。正算赤，负算黑。否则以邪正为异。同名相除，异名相益。正无入负之，负无入正之。其异名相除，同名相益，正无入正之，负无入负之。此正负之术也。方程自有消去之法，正负所以记其消去之迹也。%

}{%

正负术（带符号数运算法）是方程方法的辅助工具。因为两种计算得失相反，所以要用正负来命名。正数用红色算筹表示，负数用黑色算筹表示。或者用斜放和正放来区别。同号相减，异号相加。正数没有对应项就变为负，负数没有对应项就变为正。方程有自己的消元方法，正负术是用来记录消元过程中的符号变化。刘徽在这里给出了世界数学史上最早的正负数运算法则，比印度和阿拉伯数学早了数百年。%

}



\fullwidthnote{\textbf{历史意义}：这道题引入了正负数（正負）的概念。刘徽解释了用红色算筹表示正数、黑色算筹表示负数。这是世界数学史上最早的带符号数系统。}



\vspace{0.5em}

\noindent\makebox[\textwidth]{\textcolor{rulecolor}{\rule{\textwidth}{0.4pt}}}

\vspace{0.5em}



%% --- Problem 6 ---

\probheader{第六题}{Problem 6: Horse, Ox, and Sheep}



\wenpair{%

今有马一、牛二、羊三，直金二十五两；马二、牛三、羊一，直金二十八两；马三、牛一、羊二，直金二十三两。问：马、牛、羊各直金几何？%

}{%

今有1匹马加2头牛加3只羊值25两金；2匹马加3头牛加1只羊值28两金；3匹马加1头牛加2只羊值23两金。问：马、牛、羊各值多少两金？%

}



\dapair{%

马一直金九两，牛一直金五两，羊一直金二两。%

}{%

1匹马值9两金，1头牛值5两金，1只羊值2两金。%

}



\shupair{%

术曰：方程术。列三行，马、牛、羊之数及直金为实。以右行马一遍乘中行、左行，以直除之，消去马。次以中行牛遍乘左行，以直除之，消去牛。左行羊不尽者，上为法，下为实，实如法得羊之直。次求牛、马之直。%

}{%

用方程（线性方程组消元法）求解。排列三行，马、牛、羊的数量和金值作为常数项。用右行的马系数1遍乘中行和左行，直除消去马。然后用中行的牛系数遍乘左行，直除消去牛。左行剩下的羊系数就是除数（法），常数项是被除数（实），实除以法得到羊的单价。然后回代求牛和马的单价。%

}



\commentarypair{%

此三物方程之又一例也。马一、牛二、羊三直二十五两；马二、牛三、羊一直二十八两；马三、牛一、羊二直二十三两。以方程术消之：右行马一为最简，故以右行遍乘中行、左行，直除消马。然后以中行消左行之牛，得羊之实。既得羊直，回代求牛、马之直。%

}{%

这是三元方程组的又一个例子。用方程消元法：右行马系数为1最简单，所以用右行遍乘中行和左行，直除消去马。然后用中行的牛系数消去左行的牛，得到羊的值。得到羊的单价后，回代求牛和马的价格。刘徽说明了选择消元顺序的策略：从系数最简单的开始。%

}



\begin{modernmath}

\begin{align*}

x + 2y + 3z &= 25 \\

2x + 3y + z &= 28 \\

3x + y + 2z &= 23

\end{align*}

解得：$x = 9$（马），$y = 5$（牛），$z = 2$（羊）。

\end{modernmath}



\vspace{0.5em}

\noindent\makebox[\textwidth]{\textcolor{rulecolor}{\rule{\textwidth}{0.4pt}}}

\vspace{0.5em}



%% --- Problem 7 ---

\probheader{第七题}{Problem 7: Upper and Lower Grain Exchange}



\wenpair{%

今有上禾六秉，损实一斗八升，当下禾一十秉；下禾一十五秉，损实五升，当上禾五秉。问：上、下禾一秉各几何？%

}{%

今有上禾6束，减去1斗8升粮，等于下禾10束；下禾15束，减去5升粮，等于上禾5束。问：上禾和下禾每束各有多少粮？%

}



\dapair{%

上禾一秉八升，下禾一秉三升。%

}{%

上禾每束8升，下禾每束3升。%

}



\shupair{%

术曰：方程术。列上禾、下禾之数及损实为行。以正负术记之，直除消去，求得各秉之实。%

}{%

用方程（线性方程组）求解。将上禾、下禾的系数和减去的粮数排列成行。用正负术（带符号数运算法）记录，直除（消元法）消去求解，得到每束的粮数。%

}



\commentarypair{%

上禾六秉损一斗八升当下禾十秉者，六秉上禾之实减一斗八升，等于十秉下禾之实也。下禾十五秉损五升当上禾五秉者，十五秉下禾之实减五升，等于五秉上禾之实也。以方程列之，正负记之，消去求之。%

}{%

上禾6束减去1斗8升等于下禾10束，就是6束上禾的粮减去1斗8升等于10束下禾的粮。下禾15束减去5升等于上禾5束，就是15束下禾的粮减去5升等于5束上禾的粮。用方程列出，正负数记录，消去求解。刘徽将文字叙述转化为精确的代数关系。%

}



\vspace{0.5em}

\noindent\makebox[\textwidth]{\textcolor{rulecolor}{\rule{\textwidth}{0.4pt}}}

\vspace{0.5em}



%% --- Problem 8 ---

\probheader{第八题}{Problem 8: Five Classes Paying Tax}



\wenpair{%

今有五等之家共输税。甲五家、乙四家、丙三家、丁二家、戊一家，共输税一千斛。欲令各依等衰输之。问：各输几何？%

}{%

今有五等人家共同缴纳税粮。甲等有5家，乙等有4家，丙等有3家，丁等有2家，戊等有1家，共缴纳税粮1000斛。要求按等级递减的比例缴纳。问：各等各缴多少？%

}



\dapair{%

甲输二百七十七斛七分斛之六，乙输二百二十二斛七分斛之二，丙输一百六十六斛七分斛之四，丁输一百一十一斛七分斛之一，戊输五十五斛七分斛之五。%

}{%

甲等缴$277\frac{6}{7}$斛，乙等缴$222\frac{2}{7}$斛，丙等缴$166\frac{4}{7}$斛，丁等缴$111\frac{1}{7}$斛，戊等缴$55\frac{5}{7}$斛。%

}



\shupair{%

术曰：以方程术入之。令五家为五未知，以家数及共输为实，依衰分列方程，消去求之。%

}{%

用方程（线性方程组）求解。将五等人家设为五个未知数，以家数和总缴纳量为已知条件，按递减比率排列方程，消去求解。%

}



\commentarypair{%

五等之家，各有多少之异。甲五家、乙四家、丙三家、丁二家、戊一家，共输千斛。依等衰者，以家数之比例为率也。甲五、乙四、丙三、丁二、戊一，共十五。以千斛分为十五分，甲得其五，乙得其四，丙得其三，丁得其二，戊得其一。此题以方程术者，验其等衰之实也。%

}{%

五等人家，各有不同数量。甲5家、乙4家、丙3家、丁2家、戊1家，共缴1000斛。按等级递减，就是按家数比例作为标准。甲5、乙4、丙3、丁2、戊1，加起来共15份。将1000斛分成15份，甲得5份，乙得4份，丙得3份，丁得2份，戊得1份。用方程方法是为了验证按等差比率分配的准确性。刘徽说明了这里方程法和衰分法（比例分配法）是一致的。%

}



\vspace{0.5em}

\noindent\makebox[\textwidth]{\textcolor{rulecolor}{\rule{\textwidth}{0.4pt}}}

\vspace{0.5em}



%% --- Problem 9 ---

\probheader{第九题}{Problem 9: Three People Holding Money}



\wenpair{%

今有甲、乙、丙三人持钱。甲得乙之半而钱满四十八，乙得丙之半而钱满四十八，丙得甲之半而钱满四十八。问：甲、乙、丙各持钱几何？%

}{%

今有甲、乙、丙三人带钱。甲的钱加上乙的一半共有48钱，乙的钱加上丙的一半共有48钱，丙的钱加上甲的一半共有48钱。问：甲、乙、丙各带了多少钱？%

}



\dapair{%

甲持钱三十六，乙持钱二十四，丙持钱十二。%

}{%

甲带了36钱，乙带了24钱，丙带了12钱。%

}



\shupair{%

术曰：以方程术入之。令甲、乙、丙持钱为三未知数，以条件列三方程，消去求之。%

}{%

用方程（线性方程组）求解。将甲、乙、丙带的钱设为三个未知数，根据条件列出三个方程，消去求解。%

}



\commentarypair{%

甲得乙之半而满四十八者，甲加乙之半等于四十八也。乙得丙之半而满四十八者，乙加丙之半等于四十八也。丙得甲之半而满四十八者，丙加甲之半等于四十八也。设甲持$x$，乙持$y$，丙持$z$，则：$x + \frac{y}{2} = 48$，$y + \frac{z}{2} = 48$，$z + \frac{x}{2} = 48$。以方程术列三行，正负记之，消去求之。%

}{%

甲加上乙的一半等于48，乙加上丙的一半等于48，丙加上甲的一半等于48。设甲带$x$钱，乙带$y$钱，丙带$z$钱，则：$x + \frac{y}{2} = 48$，$y + \frac{z}{2} = 48$，$z + \frac{x}{2} = 48$。用方程消元法排列三行，正负数记录，消去求解。刘徽用清晰的代数符号表达了方程组。%

}



\vspace{0.5em}

\noindent\makebox[\textwidth]{\textcolor{rulecolor}{\rule{\textwidth}{0.4pt}}}

\vspace{0.5em}



%% --- Problem 10 ---

\probheader{第十题}{Problem 10: Five Families Sharing a Well (Revisited)}



\wenpair{%

今有五家共井。甲二绠不足，如乙一绠；乙三绠不足，如丙一绠；丙四绠不足，如丁一绠；丁五绠不足，如戊一绠；戊六绠不足，如甲一绠。如各得所不足一绠，皆逮。问：井深、绠长各几何？%

}{%

今有五家共用一口井。甲家2条绳不够，需加乙家1条；乙家3条绳不够，需加丙家1条；丙家4条绳不够，需加丁家1条；丁家5条绳不够，需加戊家1条；戊家6条绳不够，需加甲家1条。如果各家都得到所缺的那1条绳，都刚好够到水面。问：井深和各家绳长各是多少？%

}



\dapair{%

井深七丈二尺一寸。甲绠长二丈六尺五寸，乙绠长一丈九尺一寸，丙绠长一丈四尺八寸，丁绠长一丈二尺九寸，戊绠长七尺六寸。%

}{%

井深7丈2尺1寸。甲家绳长2丈6尺5寸，乙家绳长1丈9尺1寸，丙家绳长1丈4尺8寸，丁家绳长1丈2尺9寸，戊家绳长7尺6寸。%

}



\shupair{%

术曰：方程术。列五家绠长及井深为六未知数，以五条件列五行。正负相生，直除消去，求得其比，然后以井深定其实长。%

}{%

用方程（线性方程组）求解。将五家绳长和井深设为6个未知数，根据5个条件排列成5行。用正负术（带符号数运算法），直除（消元法）消去，求出各绳长之比，然后根据井深确定实际长度。%

}



\commentarypair{%

此题与卷七所载为同一题，而术不同。卷七以盈不足术验之，此则以方程术正解之。甲二绠不足如乙一绠者，二倍甲绠加井深等于乙绠也。设井深为$d$，甲、乙、丙、丁、戊绠长各为$a, b, c, d', e$，则有：$2a + d = b$，$3b + d = c$，$4c + d = d'$，$5d' + d = e$，$6e + d = a$。以方程术列之，求得各绠与井深之比。此五程六物之方程，为不定问题，故特设井深以定其解。%

}{%

这道题和卷七记载的是同一道题，但方法不同。卷七用盈不足方法验证，这里用方程方法正式求解。甲家2条绳不够需加乙家1条，就是2倍甲绳长加井深等于乙绳长。设井深为$d$，五家绳长分别为$a, b, c, d', e$，则有：$2a + d = b$，$3b + d = c$，$4c + d = d'$，$5d' + d = e$，$6e + d = a$。用方程消元法排列，求出各绳与井深的比例。这是5个方程6个未知数的不定方程组，所以需要特别设定井深来确定唯一解。刘徽说明了不定方程的处理方法。%

}





%% --- Problem 11 ---

\probheader{第十一题}{Problem 11: Two Grains Exchange}



\wenpair{%

今有上禾七秉，益实六斗，当下禾一十秉；下禾五秉，益实一斗，当上禾二秉。问：上、下禾一秉各几何？%

}{%

今有上禾7束，加上6斗粮，等于下禾10束；下禾5束，加上1斗粮，等于上禾2束。问：上禾和下禾每束各有多少粮？%

}



\dapair{%

上禾一秉八升，下禾一秉三升。%

}{%

上禾每束8升，下禾每束3升。%

}



\shupair{%

术曰：方程术。列上禾、下禾之数及益实为行，正负记之，直除消去，求得各秉之实。%

}{%

用方程（线性方程组消元法）求解。将上禾、下禾的系数和加的粮数排列成行，正负数记录，直除（消元法）消去，求得每束的粮数。%

}



\commentarypair{%

上禾七秉益六斗当下禾十秉者，七秉上禾之实加六斗，等于十秉下禾之实也。下禾五秉益一斗当上禾二秉者，五秉下禾之实加一斗，等于二秉上禾之实也。以方程列之，此二物二程之方程也。%

}{%

上禾7束加6斗等于下禾10束，就是7束上禾的粮加6斗等于10束下禾的粮。下禾5束加1斗等于上禾2束，就是5束下禾的粮加1斗等于2束上禾的粮。用方程列出，这是二元一次方程组。刘徽说明了这是两个未知数两个方程的标准情形。%

}



\vspace{0.5em}

\noindent\makebox[\textwidth]{\textcolor{rulecolor}{\rule{\textwidth}{0.4pt}}}

\vspace{0.5em}



%% --- Problem 12 ---

\probheader{第十二题}{Problem 12: Large and Small Vessels}



\wenpair{%

今有大器一、小器五，容三斛；大器五、小器一，容二斛。问：大、小器各容几何？%

}{%

今有1个大容器和5个小容器，共能装3斛；5个大容器和1个小容器，共能装2斛。问：大容器和小容器各能装多少？%

}



\dapair{%

大器容二十四分斛之十三，小器容二十四分斛之七。%

}{%

大容器能装$\frac{13}{24}$斛，小容器能装$\frac{7}{24}$斛。%

}



\shupair{%

术曰：方程术。列大器、小器之数及容为两行，直除消去大器，得小器之实。以小器之实回代，求得大器之容。%

}{%

用方程（线性方程组消元法）求解。将大器、小器的数量和容量排列成两行，直除（消元法）消去大器，得到小器的被除数（实）。用小器的值回代，求出大器的容量。%

}



\commentarypair{%

大器一、小器五容三斛，大器五、小器一容二斛。以方程术列之：右行大器一、小器五、实三；左行大器五、小器一、实二。以右行遍乘左行，直除消去大器，得小器之实。回代求大器。此二物方程之简者也。%

}{%

1个大容器和5个小容器装3斛，5个大容器和1个小容器装2斛。用方程消元法排列：右行大器1、小器5、实3；左行大器5、小器1、实2。用右行遍乘左行，直除消去大器，得到小器的值。然后回代求大器。这是二元方程组的简单例子。刘徽展示了最简二元方程组的解法。%

}



\vspace{0.5em}

\noindent\makebox[\textwidth]{\textcolor{rulecolor}{\rule{\textwidth}{0.4pt}}}

\vspace{0.5em}



%% --- Problem 13 ---

\probheader{第十三题}{Problem 13: Gold, Silver, Lead, Tin}



\wenpair{%

今有金一斤、银一斤，称之适等；金一斤、铅一斤，称之金轻五两；银一斤、锡一斤，称之银轻三两。问：金、银、铅、锡一斤各重几何？%

}{%

今有1斤金和1斤银，称量它们重量相等；1斤金和1斤铅，称量时金这边轻5两；1斤银和1斤锡，称量时银这边轻3两。问：金、银、铅、锡每1斤各重多少？%

}



\dapair{%

金一斤重一斤，银一斤重一斤，铅一斤重一斤五两，锡一斤重一斤三两。%

}{%

1斤金重1斤，1斤银重1斤，1斤铅重1斤5两，1斤锡重1斤3两。%

}



\shupair{%

术曰：方程术。以金银适等为一条件，金铅差五两为一条件，银锡差三两为一条件，列方程求之。%

}{%

用方程（线性方程组）求解。以金银重量相等为第一个条件，金铅相差5两为第二个条件，银锡相差3两为第三个条件，排列成方程组求解。%

}



\commentarypair{%

金银适等者，一斤金与一斤银重量相等也。此以天平称之，非论其体积。金铅称之金轻五两者，同体积之铅重于金五两也。银锡称之银轻三两者，同体积之锡重于银三两也。以方程术列之，三条件四物，当以一物为率，求各物之比。%

}{%

金银重量相等，就是1斤金和1斤银在天平上平衡。这是用天平称量，不是比较体积。金和铅称量时金轻5两，就是同体积的铅比金重5两。银和锡称量时银轻3两，就是同体积的锡比银重3两。用方程方法排列，3个条件4种物质，需要以一种物质为标准，求各物质的比例。刘徽解释了这里比较的是同体积下的比重差异。%

}



\vspace{0.5em}

\noindent\makebox[\textwidth]{\textcolor{rulecolor}{\rule{\textwidth}{0.4pt}}}

\vspace{0.5em}



%% --- Problem 14 ---

\probheader{第十四题}{Problem 14: New Method of Rectangular Arrays}



\wenpair{%

今有上禾二秉，中禾三秉，下禾四秉，实皆不满斗。上禾取中禾一秉，中禾取下禾一秉，下禾取上禾一秉，而实满斗。问：上、中、下禾一秉各几何？%

}{%

今有上禾2束、中禾3束、下禾4束，各自的粮都不满1斗。如果上禾加上中禾的1束，中禾加上下禾的1束，下禾加上上禾的1束，就都刚好满1斗。问：上禾、中禾、下禾每束各有多少粮？%

}



\dapair{%

上禾一秉二十五分斗之十一，中禾一秉二十五分斗之七，下禾一秉二十五分斗之四。%

}{%

上禾每束$\frac{11}{25}$斗，中禾每束$\frac{7}{25}$斗，下禾每束$\frac{4}{25}$斗。%

}



\shupair{%

术曰：方程新术。以所取之数列为方程，各以其取数遍乘，直除消去，各得其实。%

}{%

用方程新法求解。将各束所取的数目排列成方程，用各自的取数遍乘，直除（消元法）消去，各得其实际粮数。%

}



\commentarypair{%

此方程术之变也。上禾二秉取中禾一秉而满斗者，上禾一秉之实加中禾半秉之实为半斗也。中禾三秉取下禾一秉而满斗者，中禾一秉之实加下禾三分之一秉之实为三分之一斗也。下禾四秉取上禾一秉而满斗者，下禾一秉之实加上禾四分之一秉之实为四分之一斗也。以三条件列三行，直除消去求之。%

}{%

这是方程方法的变体。上禾2束加中禾1束满1斗，就是上禾1束的粮加中禾半束的粮等于半斗。中禾3束加下禾1束满1斗，就是中禾1束的粮加下禾$\frac{1}{3}$束的粮等于$\frac{1}{3}$斗。下禾4束加上禾1束满1斗，就是下禾1束的粮加上禾$\frac{1}{4}$束的粮等于$\frac{1}{4}$斗。用三个条件排列成三行，消去求解。刘徽详细解释了``取''的含义。%

}



\vspace{0.5em}

\noindent\makebox[\textwidth]{\textcolor{rulecolor}{\rule{\textwidth}{0.4pt}}}

\vspace{0.5em}



%% --- Problem 15 ---

\probheader{第十五题}{Problem 15: Three Grades of Rice}



\wenpair{%

今有粟五斗，为糲米、粺米、粪米三等。欲令糲二、粺三、粪四为率。问：各为米几何？%

}{%

今有粟米5斗，要加工成糲米（粗米）、粺米（精米）、糞米（细米）三等。要求按糲2、粺3、糞4的比率分配。问：各等米各有多少？%

}



\dapair{%

糲米一斗四升七分升之四，粺米二斗一升七分升之六，粪米二斗八升七分升之二。%

}{%

糲米$1\frac{4}{7}$斗，粺米$2\frac{6}{7}$斗，糞米$2\frac{2}{7}$斗。%

}



\shupair{%

术曰：以方程术入之。令三米之率为三未知数，以粟五斗为实，以率之比列方程，消去求之。%

}{%

用方程（线性方程组）求解。将三种米的比率设为三个未知数，以粟5斗为总量（实），按比率的比值排列成方程，消去求解。%

}



\commentarypair{%

糲二、粺三、粪四为率者，三米当为二比三比四之比例也。以方程列之：设糲米为$2k$，粺米为$3k$，粪米为$4k$，则$2k + 3k + 4k = 5$斗，$9k = 5$斗，$k = \frac{5}{9}$斗。然后各以其率乘之。此方程术与衰分术之合也。%

}{%

糲米2、粺米3、糞米4为比率，就是三种米按$2:3:4$的比例分配。用方程列出：设糲米为$2k$，粺米为$3k$，糞米为$4k$，则$2k + 3k + 4k = 5$斗，$9k = 5$斗，$k = \frac{5}{9}$斗。然后各乘以其比率。这是方程法和衰分法（比例分配法）的结合。刘徽展示了如何将比例问题转化为方程求解。%

}



\newpage





%% ========================================================================

%% VOLUME 9: GOUGU (勾股 / Right Triangles)

%% ========================================================================

\volheader{卷第九}{勾股}{Volume~9}{Right Triangles (Pythagorean Theorem)}



\fullwidthnote{勾股（直角三角形/毕达哥拉斯定理）章是《九章算术》的压轴之作，以非凡的深度阐述了勾股定理及其应用。刘徽对本章的注释尤为卓越——包括他用``割补术''证明勾股定理的著名方法，以及``出入相补原理''的讨论。本章还包含测距问题、方圆关系，以及中国古代对二次无理数的最早处理。}



%% --- Fundamental Theorem ---

\probheader{勾股术总述}{Fundamental Theorem of Right Triangles}



\shupair{%

勾股术曰：句股各自乘，并而开方除之，即弦。又股自乘，以减弦自乘，其余开方除之，即句。又句自乘，以减弦自乘，其余开方除之，即股。%

}{%

勾股（直角三角形）的方法如下：勾和股各自平方，相加后开平方，就得到弦。或者，股平方后用弦的平方减去它，余数开平方就得到勾。又或者，勾平方后用弦的平方减去它，余数开平方就得到股。%

}



\commentarypair{%

勾股之法，出于圆方。圆径一而周三，方径一而周四。以圆方之率，生勾股之术。按句股者，矩之别名也。矩之为物，横者为句，纵者为股，斜者为弦。句股各自乘，并之为弦实。开方除之，得弦。此谓句股各自乘，并之为大正方形，而以弦为边也。%

}{%

勾股的方法来源于圆和方。圆的直径为1则周长为3，方的边长为1则周长为4。以圆和方的比率为基础，发展出勾股的方法。勾股就是矩（直角曲尺）的别名。矩这个工具，横边叫勾，纵边叫股，斜边叫弦。勾和股各自平方后相加，就是弦的平方。开平方得到弦的长度。也就是说，勾和股各自平方后相加，等于以弦为边长的正方形的面积。刘徽追溯了勾股概念的几何来源。%

}



\fullwidthnote{\textbf{勾股定理}：若直角三角形的两条直角边（勾和股）长度分别为$a$和$b$，斜边（弦）长度为$c$，则 $a^2 + b^2 = c^2$。刘徽用``割补术''证明：以弦为边的正方形（弦实）可以切割重组成以勾和股为边的两个正方形（勾实和股实）。}



\vspace{0.8em}

\noindent\makebox[\textwidth]{\textcolor{rulecolor}{\rule{\textwidth}{0.6pt}}}

\vspace{0.8em}



%% --- Problem 1 ---

\probheader{第一题}{Problem 1: Gougu Finding Hypotenuse}



\wenpair{%

今有句三尺，股四尺，问：为弦几何？%

}{%

今有一个直角三角形，勾（短直角边）长3尺，股（长直角边）长4尺，问：弦（斜边）长多少？%

}



\dapair{%

弦五尺。%

}{%

弦长5尺。%

}



\shupair{%

术曰：句股各自乘，并而开方除之，即弦。句自乘得九，股自乘得一十六，并之得二十五，开方除之得五尺。%

}{%

方法如下：勾和股各自平方，相加后开平方，就得到弦。勾3尺自乘得9，股4尺自乘得16，相加得25，开平方得5尺。%

}



\commentarypair{%

此勾股之正术也。勾三、股四、弦五，此为勾股之最基本率。勾自乘得九，股自乘得一十六，并之为二十五，即弦自乘之数。开方得五，即弦。按勾三股四弦五，此所谓勾股弦之率也。以三、四、五为勾股弦者，其直角三角形中最简者也。凡勾股之数，皆可由此率推之。%

}{%

这是勾股方法的标准应用。勾3、股4、弦5，这是最基本的勾股数。勾自乘得9，股自乘得16，相加得25，就是以弦为边的正方形面积。开平方得5，就是弦长。勾3股4弦5是勾股数的基本比率。以3、4、5为勾股弦的直角三角形是最简单的。所有勾股数都可以从这个基本比率推出。刘徽指出了$(3,4,5)$是最基本的勾股数。%

}



\fullwidthnote{\textbf{注}：这是著名的3-4-5直角三角形，最简单的勾股数组。}



\vspace{0.5em}

\noindent\makebox[\textwidth]{\textcolor{rulecolor}{\rule{\textwidth}{0.4pt}}}

\vspace{0.5em}



%% --- Problem 2 ---

\probheader{第二题}{Problem 2: Hypotenuse and Gu Finding Gou}



\wenpair{%

今有弦五尺，股四尺，问：为句几何？%

}{%

今有一个直角三角形，弦（斜边）长5尺，股（长直角边）长4尺，问：勾（短直角边）长多少？%

}



\dapair{%

句三尺。%

}{%

勾长3尺。%

}



\shupair{%

术曰：股自乘，以减弦自乘，其余开方除之，即句。股自乘得一十六，弦自乘得二十五，以十六减二十五余九，开方除之得三尺。%

}{%

方法如下：股平方后，用弦的平方减去它，余数开平方就得到勾。股4尺自乘得16，弦5尺自乘得25，用16减25余9，开平方得3尺。%

}



\commentarypair{%

弦实二十五，股实一十六，两实之差为句实九。开方得三，即句。此术为勾股术之逆也。弦实者，以弦为边之正方形也；股实者，以股为边之正方形也。弦实减股实，所余即句实也。%

}{%

弦的平方是25，股的平方是16，两个平方的差就是勾的平方9。开平方得3，就是勾长。这是勾股方法的逆运算。弦实就是以弦为边的正方形面积，股实就是以股为边的正方形面积。弦实减股实，余数就是勾实。刘徽用``实''的概念来阐释面积关系。%

}



\vspace{0.5em}

\noindent\makebox[\textwidth]{\textcolor{rulecolor}{\rule{\textwidth}{0.4pt}}}

\vspace{0.5em}



%% --- Problem 3 ---

\probheader{第三题}{Problem 3: Gou and Hypotenuse Finding Gu}



\wenpair{%

今有句八尺，弦十七尺，问：为股几何？%

}{%

今有一个直角三角形，勾（短直角边）长8尺，弦（斜边）长17尺，问：股（长直角边）长多少？%

}



\dapair{%

股十五尺。%

}{%

股长15尺。%

}



\shupair{%

术曰：句自乘，以减弦自乘，其余开方除之，即股。句自乘得六十四，弦自乘得二百八十九，以六十四减二百八十九余二百二十五，开方除之得一十五尺。%

}{%

方法如下：勾平方后，用弦的平方减去它，余数开平方就得到股。勾8尺自乘得64，弦17尺自乘得289，用64减289余225，开平方得15尺。%

}



\commentarypair{%

弦宽二百八十九，句实六十四，两实之差为股实二百二十五。开方得一十五，即股。此又以勾股术之逆也。勾八、股十五、弦十七，此又一股弦之率也。%

}{%

弦的平方是289，勾的平方是64，两个平方的差就是股的平方225。开平方得15，就是股长。这又是勾股定理的逆运算。勾8、股15、弦17，这是另一组勾股数。刘徽指出了$(8,15,17)$也是一组勾股数。%

}



\vspace{0.5em}

\noindent\makebox[\textwidth]{\textcolor{rulecolor}{\rule{\textwidth}{0.4pt}}}

\vspace{0.5em}



%% --- Problem 4 ---

\probheader{第四题}{Problem 4: Door Height and Width}



\wenpair{%

今有户高多于广六尺八寸，两隅相去适一丈。问：户高、广各几何？%

}{%

今有一扇门，高度比宽度多6尺8寸，对角线（两隅相距）刚好1丈。问：门的高度和宽度各是多少？%

}



\dapair{%

广二尺八寸，高九尺六寸。%

}{%

宽2尺8寸，高9尺6寸。%

}



\shupair{%

术曰：令一丈自乘，半之得五十尺为实。令多六尺八寸自乘，得四十六尺二寸四分，半之得二十三尺一寸二分为法。实如法得一尺，为户广。加多六尺八寸，为高。%

}{%

方法如下：将1丈自乘得100尺，取半得50尺作为被除数（实）。将差6尺8寸自乘得46.24尺，取半得23.12尺作为除数（法）。实除以法再加一得到宽度。加上差6尺8寸就是高度。%

}



\commentarypair{%

户高多于广六尺八寸者，高与广之差也。两隅相去一丈者，以高广为勾股之弦也。令一丈自乘者，弦实也。高广之差自乘者，差实也。弦实加差实，开方得高；弦实减差实，开方得广。此术以勾股弦之关系求高广也。%

}{%

门高比宽多6尺8寸，就是高和宽的差。对角线1丈，就是以高和宽为勾股的弦。1丈自乘是弦的平方。高宽之差的平方是差的平方。弦的平方加差的平方开方得高，弦的平方减差的平方开方得宽。这是用勾股弦的关系求高和宽。刘徽推导了$\sqrt{\frac{c^2+d^2}{2}}$和$\sqrt{\frac{c^2-d^2}{2}}$的公式，其中$c$为弦，$d$为差。%

}



\begin{modernmath}

设 $h$ = 高度，$w$ = 宽度。则 $h - w = 6.8$ 且 $h^2 + w^2 = 100$。

\end{modernmath}



\vspace{0.5em}

\noindent\makebox[\textwidth]{\textcolor{rulecolor}{\rule{\textwidth}{0.4pt}}}

\vspace{0.5em}



%% --- Problem 5 ---

\probheader{第五题}{Problem 5: City Wall and East Gate Tree}



\wenpair{%

今有邑方不知大小，各开中门。出东门二十步有木。问：出南门几步而见木？%

}{%

今有一座正方形的城池，不知道大小，四面各开中门。出东门向东走20步有一棵树。问：出南门向南走多少步能看见那棵树？%

}



\dapair{%

出南门一十四步半而见木。%

}{%

出南门向南走14.5步就能看见那棵树。%

}



\shupair{%

术曰：以勾股相似术入之。出东门二十步为勾，邑方之半为股。以股自乘，以勾除之，得南门外见木之步数。%

}{%

用勾股相似三角形方法求解。出东门20步是一个勾，城边长的一半是一个股。用股平方除以勾，就得到南门外能看见树的步数。%

}



\commentarypair{%

邑方者，正方形之邑也。各开中门者，四面之门各开于边之中点也。出东门二十步有木者，从东门中出东行二十步有树木也。出南门见木者，从南门中出南行，望见东门外之木也。此以勾股相似之术求之：邑方之半为一大股，东门外二十步为一大勾；南门外所求步数为一小股，邑方之半亦为一小勾之对应边。两勾股相似，故以比例求之。%

}{%

邑方就是正方形的城池。各开中门就是四面城门都开在每边的中点。出东门20步有树，就是从东门中点向东走20步有树木。出南门见树，就是从南门中点向南走，能看见东门外的树。这是用相似勾股三角形求解：城边的一半是一个大三角形的一条直角边（股），东门外20步是大三角形的另一条直角边（勾）；南门外所求的步数是小三角形的一条直角边，城边的一半是小三角形的另一条直角边。两个直角三角形相似，所以用比例求解。刘徽详细解释了相似三角形的对应关系。%

}



\vspace{0.5em}

\noindent\makebox[\textwidth]{\textcolor{rulecolor}{\rule{\textwidth}{0.4pt}}}

\vspace{0.5em}



%% --- Problem 6 ---

\probheader{第六题}{Problem 6: Well Diameter and Pole}



\wenpair{%

今有井径五尺，不知其深。立木于井上，从木末望水岸，入径四尺。问：井深几何？%

}{%

今有一口井，井口直径5尺，不知道深度。在井口竖立一根木杆，从木杆顶端斜望水面边缘，视线进入井口4尺。问：井有多深？%

}



\dapair{%

井深一丈九尺二寸。%

}{%

井深1丈9尺2寸。%

}



\shupair{%

术曰：以勾股相似术入之。井径五尺为勾，入径四尺为股之差。以井径乘入径，以井径减入径之差除之，得井深。%

}{%

用勾股相似三角形方法求解。井口直径5尺作为勾，视线进入井口的4尺作为股之差。用井径乘以入径，除以井径和入径之差，得到井深。%

}



\commentarypair{%

井径五尺者，井口之直径也。立木于井上者，木直立于井口之边也。从木末望水岸者，从木顶斜望水面之岸也。入径四尺者，视线入于井口之处距井边四尺也。此亦以勾股相似之术：木高为一大股，入径为一大勾；井深为一小股，井径之半为一小勾。两勾股相似，故可以比例求之。%

}{%

井径5尺是井口的直径。立木于井上是在井口边缘竖立木杆。从木顶端斜望水面边缘，视线进入井口4尺就是视线在井口水平面上的入点到井边的距离。这也是用相似勾股三角形：木杆高度是一个大三角形的一条直角边（股），入径4尺是大三角形的另一条直角边（勾）；井深是小三角形的一条直角边，井径的一半是小三角形的另一条直角边。两个直角三角形相似，所以可以用比例求解。刘徽再次运用相似三角形原理解题。%

}



\vspace{0.5em}

\noindent\makebox[\textwidth]{\textcolor{rulecolor}{\rule{\textwidth}{0.4pt}}}

\vspace{0.5em}



%% --- Problem 7 ---

\probheader{第七题}{Problem 7: Mountain Top Viewing a City}



\wenpair{%

今有登山望邑，邑在山东。山高不知高，立两表齐高三丈，前后相去千步。令后表与前表参相直。从前表却行一百二十三步，人目着地，取望邑西北隅，入表前三寸。从后表却行一百二十七步，人目着地，取望邑西北隅，适与表端参合。问：邑方及山去表各几何？%

}{%

今登山望城，城在山东面。山的高度不知道，竖立两根标杆都高3丈，前后相距1000步。让后标杆和前标杆成一直线。从前标杆后退123步，眼睛贴地，望城西北角，视线超过前标杆顶端3寸。从后标杆后退127步，眼睛贴地，望城西北角，视线恰好通过后标杆顶端。问：城的边长以及山离标杆各多少？%

}



\dapair{%

邑方一里二百二十步，山去表一里一百八十二步半。%

}{%

城的边长1里220步，山离标杆1里182.5步。%

}



\shupair{%

术曰：以重差术入之。此为重差之正术，以两表之高及相去之远，两却行之数，入表之数，列为勾股，相似求之。以勾股之差遍乘，以除之，得邑方及远近。%

}{%

用重差术（双测法）求解。这是重差术的标准方法，用两根标杆的高度和相距距离，两个后退的步数，视线入表的尺寸，排列成勾股相似三角形来求解。用勾股之差的遍乘除法，得到城边和远近距离。%

}



\commentarypair{%

此重差术之典型题也。重差者，以两表之重迭差数求远物之高深广远也。立两表齐高三丈者，立两根标杆，同高三丈也。前后相去千步者，两表之距离也。从前表却行一百二十三步者，从前表向后退一百二十三步也。入表前三寸者，视线越过前表之顶，落于前表前方三寸处也。从后表却行一百二十七步而适与表端参合者，视线恰好通过后表之顶端也。以此两组勾股之比例差，可以求邑方及山去表之远。此术之妙，在于以咫尺之矩，测千里之远。%

}{%

这是重差术的典型问题。重差术就是用两根标杆的差数来求远处物体的高度、深度、宽度和距离。立两根标杆同高3丈，前后相距1000步。从前标杆后退123步，入表3寸是视线越过前标杆顶端3寸。从后标杆后退127步视线恰好通过后标杆顶端。用这两组勾股三角形的比例差，可以求出城的边长和山离标杆的距离。这种方法的精妙之处，在于用咫尺之矩（短小的标杆）测量千里之遥（远处的物体）。刘徽盛赞重差术``以咫尺之矩，测千里之高深''的精妙。%

}



\vspace{0.5em}

\noindent\makebox[\textwidth]{\textcolor{rulecolor}{\rule{\textwidth}{0.4pt}}}

\vspace{0.5em}



%% --- Problem 8 ---

\probheader{第八题}{Problem 8: Square Inscribed in Right Triangle}



\wenpair{%

今有句五步，股十二步，问：句中容方几何？%

}{%

今有一个直角三角形，勾（短直角边）5步，股（长直角边）12步，问：在这个三角形中能容纳的最大正方形有多大？%

}



\dapair{%

方三步十七分步之九。%

}{%

正方形边长为$3\frac{9}{17}$步。%

}



\shupair{%

术曰：并句、股为法，句股相乘为实，实如法而一，得方。句五步、股十二步，并之得十七步为法。句股相乘得六十步为实。实如法得三步十七分步之九。%

}{%

方法如下：将勾和股相加作为除数（法），勾和股相乘作为被除数（实），实除以法得到正方形的边长。勾5步加股12步得17步作为法。勾乘股得60步作为实。60除以17得$3\frac{9}{17}$步。%

}



\commentarypair{%

勾股容方者，于直角三角形中作一最大之内接正方形也。其正方形之一边落于弦上，对角之顶点落于直角处。以勾股之相并为法，勾股之相乘为实者，以相似勾股之理推之也。勾股容方之边，等于勾股相乘除以勾股之和。此术之证，以勾股之相似形割补得之。%

}{%

勾股容方就是在直角三角形中作一个最大的内接正方形。这个正方形的一条边落在斜边上，对角顶点落在直角处。用勾加股作为除数，勾乘股作为被除数，这是由相似勾股三角形的原理推导出来的。内接正方形的边长等于勾乘股除以勾加股。这个方法的证明是利用相似勾股三角形的割补关系。刘徽说明了内接正方形的位置和公式的推导原理。%

}



\begin{modernmath}

对于直角边为 $a = 5$ 和 $b = 12$ 的直角三角形，以斜边为底边的内接正方形边长：

\[s = \frac{ab}{a+b} = \frac{5 \times 12}{5 + 12} = \frac{60}{17} = 3\frac{9}{17}\]

\end{modernmath}





%% --- Problem 9 ---

\probheader{第九题}{Problem 9: Circle Inscribed in Right Triangle}



\wenpair{%

今有句八步，股十五步，问：句中容圆，径几何？%

}{%

今有一个直角三角形，勾（短直角边）8步，股（长直角边）15步，问：在这个三角形中能容纳的最大内切圆，直径是多少？%

}



\dapair{%

圆径六步。%

}{%

圆的直径是6步。%

}



\shupair{%

术曰：句股相乘，倍之为实。并句、股、弦为法。实如法而一，得圆径。句八步、股十五步，句股相乘得一百二十，倍之得二百四十为实。句八、股十五、弦十七，并之得四十为法。实如法得六步。%

}{%

方法如下：勾和股相乘，加倍作为被除数（实）。将勾、股、弦相加作为除数（法）。实除以法得到圆直径。勾8步股15步，相乘得120，加倍得240作为实。勾8、股15、弦17，相加得40作为法。240除以40得6步。%

}



\commentarypair{%

勾股容圆者，于直角三角形中作一最大之内切圆也。以句股相乘倍之为实者，两倍三角形之面积也。以句股弦之和为法者，三角形周长之半也。三角之面积等于内切圆半径乘周长之半，故以面积之两倍除以周长之半，得圆径。此术之精妙，以面积之关系推圆径也。%

}{%

勾股容圆就是在直角三角形中作一个最大的内切圆。勾乘股加倍作为被除数，就是三角形面积的两倍。勾股弦之和作为除数，就是三角形周长的一半。三角形面积等于内切圆半径乘以半周长，所以用面积的两倍除以半周长，得到圆直径。这个方法的精妙之处在于利用面积关系来推求圆直径。刘徽用面积公式 $A = r \cdot s$（其中$s$为半周长）来推导内切圆半径。%

}



\begin{modernmath}

直角三角形的内切圆半径：

\[r = \frac{a + b - c}{2} = \frac{8 + 15 - 17}{2} = 3, \quad d = 2r = 6\]

等价公式：$d = \frac{2ab}{a+b+c} = \frac{2 \times 8 \times 15}{8+15+17} = \frac{240}{40} = 6$。

\end{modernmath}



\vspace{0.5em}

\noindent\makebox[\textwidth]{\textcolor{rulecolor}{\rule{\textwidth}{0.4pt}}}

\vspace{0.5em}



%% --- Problem 10 ---

\probheader{第十题}{Problem 10: City Center Viewing a Tree}



\wenpair{%

今有邑方不知大小，各开中门。出北门二十步有木。出南门十四步，折而西行一千七百七十五步见木。问：邑方几何？%

}{%

今有一座正方形城池，不知道大小，四面各开中门。出北门向北走20步有一棵树。出南门向南走14步，然后转向西走1775步就看见那棵树。问：城的边长是多少？%

}



\dapair{%

邑方五十步。%

}{%

城的边长是50步。%

}



\shupair{%

术曰：以勾股术入之。出北门二十步与出南门十四步相加得三十四步，以邑方之半乘之，以西行一千七百七十五步除之，得邑方之半。倍之得邑方。%

}{%

用勾股方法求解。出北门20步加上出南门14步得34步，乘以城边的一半，再除以向西走的1775步，得到城边的一半。加倍就得到城的边长。%

}



\commentarypair{%

邑方者，正方形之城邑也。出北门二十步有木者，从北门中出北行二十步有树也。出南门十四步折而西行者，从南门中出南行十四步，然后转向西行一千七百七十五步，恰见北门外之木也。此以勾股相似之术求之：邑方之半为一勾，南北出门步数之和为一股，西行步数为另一勾。两勾股相似，故可以比例求邑方。%

}{%

邑方就是正方形的城池。出北门20步有树，就是从北门中点向北走20步有树木。出南门14步再向西走1775步就看见北门外的那棵树。这是用相似勾股三角形求解：城边的一半是一个直角边，南北出门步数之和是另一个直角边，向西走的步数是对应三角形的直角边。两个直角三角形相似，所以可以用比例求城边长。刘徽分析了视线构成的相似三角形。%

}



\vspace{0.5em}

\noindent\makebox[\textwidth]{\textcolor{rulecolor}{\rule{\textwidth}{0.4pt}}}

\vspace{0.5em}



%% --- Problem 11 ---

\probheader{第十一题}{Problem 11: Viewing a Deep Pool}



\wenpair{%

今有清渊，不知深浅。立一木于岸上，高三丈。从木末斜望水际，入木一尺二寸。又望水底，入木四尺八寸。问：渊深几何？%

}{%

今有一个清澈的水潭，不知道深浅。在岸上竖立一根木杆，高3丈。从木杆顶端斜望水面边缘，视线在木杆上距顶端1尺2寸处。再斜望水底，视线在木杆上距顶端4尺8寸处。问：水潭有多深？%

}



\dapair{%

渊深一丈二尺。%

}{%

水潭深1丈2尺。%

}



\shupair{%

术曰：以勾股相似术入之。木高三丈为股，入木一尺二寸为勾。入木四尺八寸为另一勾之率。以比例求之，得渊深。%

}{%

用勾股相似三角形方法求解。木杆高3丈作为股，入木1尺2寸作为勾。入木4尺8寸作为另一个勾。用比例求解，得到水潭深度。%

}



\commentarypair{%

立木高三丈于岸上，从木末斜望水际入木一尺二寸者，视线从木顶到水面与木杆相交处距木顶一尺二寸也。又望水底入木四尺八寸者，视线从木顶到水底与木杆相交处距木顶四尺八寸也。以勾股相似之术：木高为一大股，入木之数为一大勾之对应。渊深与入木之差之比例，如木高与入水之比例。故以比例求之。%

}{%

在岸上竖立3丈高的木杆，从木顶端斜望水面边缘，视线在木杆上距顶端1.2尺处。再斜望水底，视线在木杆上距顶端4.8尺处。用相似勾股三角形：木杆高度是一个大三角形的一条直角边，入木尺寸是大三角形的另一条直角边。水潭深度与入木尺寸之差的比例，等于木高与入水比例。所以用比例求解。刘徽利用两组相似三角形来求水深。%

}



\vspace{0.5em}

\noindent\makebox[\textwidth]{\textcolor{rulecolor}{\rule{\textwidth}{0.4pt}}}

\vspace{0.5em}



%% --- Problem 12 ---

\probheader{第十二题}{Problem 12: Peeping at a Tree}



\wenpair{%

今有木去人不知远近。立一表长一丈，人退行二丈六尺，从表末望木，与表端参合。又退行二丈，从表末望木，入表一尺。问：木去表几何？%

}{%

今有一棵树，离人不知道远近。竖立一根标杆长1丈，人后退2丈6尺，从标杆顶端望树，视线恰好与标杆顶端和树顶成一直线。再后退2丈，总共后退4丈6尺，从标杆顶端望树，视线进入标杆1尺（即视线在标杆后方1尺处与标杆相交）。问：树离标杆有多远？%

}



\dapair{%

木去表三十四丈。%

}{%

树离标杆34丈。%

}



\shupair{%

术曰：以重差术入之。两退行之差为法，表长乘后退行为实，实如法而一，得木去表之远。%

}{%

用重差术（双测法）求解。两次后退的距离之差作为除数（法），标杆长度乘以后退的总距离作为被除数（实），实除以法得到树离标杆的距离。%

}



\commentarypair{%

此又以重差术测远也。立表一丈，退行二丈六尺而望木与表端参合者，木、表端、人目三点一线也。又退行二丈，共退行四丈六尺，入表一尺者，视线越过表端一尺也。两退行之差为二丈，表长一丈乘后退行四丈六尺为实，以二丈为法，实如法得二十三丈。此以勾股之差求之也。%

}{%

这又是用重差术测远距离。立标杆1丈，后退2.6丈望树与标杆顶端成一直线，就是树、标杆顶端、人目三点一线。再后退2丈，共后退4.6丈，入表1尺就是视线越过标杆顶端1尺。两次后退之差为2丈，标杆1丈乘以后退4.6丈作为被除数，用2丈作为除数，相除得到树离标杆的距离。刘徽展示了重差术如何通过两次不同位置的观测来求距离。%

}



\vspace{0.5em}

\noindent\makebox[\textwidth]{\textcolor{rulecolor}{\rule{\textwidth}{0.4pt}}}

\vspace{0.5em}



%% --- Problem 13 ---

\probheader{第十三题}{Problem 13: Viewing a Mountain via a Tree}



\wenpair{%

今有木不知高远。立两表，各高一丈，前后相去五百步。从前表却行一百二十步，人目着地，望山峰与表端参合。从后表却行一百八十步，人目着地，望山峰与表端参合。问：山去前表几何？山高几何？%

}{%

今有一座山，不知道距离和高度。竖立两根标杆，各高1丈，前后相距500步。从前标杆后退120步，眼睛贴地，望山峰恰好与标杆顶端成一直线。从后标杆后退180步，眼睛贴地，望山峰也恰好与标杆顶端成一直线。问：山离前标杆多远？山有多高？%

}



\dapair{%

山去前表一里二百步，高八里一百二十步。%

}{%

山离前标杆1里220步，山高8里120步。%

}



\shupair{%

术曰：以重差术入之。两表相去为勾股之率，两却行之差为法。表高乘表间为实，实如法得山高。前却行乘表间为实，实如法得山去表之远。%

}{%

用重差术（双测法）求解。两根标杆之间的距离作为勾股的比率，两次后退距离之差作为除数（法）。标杆高度乘以标杆间距作为被除数（实），实除以法得到山的高度。前标杆后退距离乘以标杆间距作为实，除以法得到山离标杆的距离。%

}



\commentarypair{%

此重差术测山高之典型题也。立两表齐高一丈，相去五百步。从前表却行一百二十步而望山峰与表端参合，从后表却行一百八十步而望山峰与表端参合。两却行之差六十步为法。表高一丈乘表间五百步为实，实如法得山高。前却行一百二十步乘表间五百步为实，如法得山去前表之远。此术之精，在于以两表之矩，测群山之高远。故刘徽曰：以咫尺之矩，测千里之高深。%

}{%

这是用重差术测山高的典型问题。竖立两根标杆同高1丈，相距500步。从前标杆后退120步望山峰与标杆顶端成一直线，从后标杆后退180步望山峰也与标杆顶端成一直线。两次后退之差60步作为除数。标杆高1丈乘以间距500步作为被除数，除以法得到山的高度。前标杆后退120步乘以间距500步作为被除数，除以法得到山离前标杆的距离。这种方法的精妙之处在于用两根短小的标杆测量群山的高远。所以刘徽说：以咫尺之矩，测千里之高深。刘徽再次盛赞重差术的精妙，并给出了完整的计算公式。%

}



\vspace{0.5em}

\noindent\makebox[\textwidth]{\textcolor{rulecolor}{\rule{\textwidth}{0.4pt}}}

\vspace{0.5em}



%% --- Problem 14 ---

\probheader{第十四题}{Problem 14: Annular Field Diameter}



\wenpair{%

今有环田，中周六十二步，外周四步，径十二步。问：为田几何？%

}{%

今有一块环形田，内周62步，外周4步，环宽12步。问：这块田的面积是多少？%

}



\dapair{%

二亩五十五步。%

}{%

面积为2亩55步。%

}



\shupair{%

术曰：并中外周而半之，以径乘之为积。中周六十二步、外周四步，并之得六十六步，半之得三十三步。以径十二步乘之，得三百九十六步为积。亩法二百四十步，得三亩二十五步。%

}{%

方法如下：将内周和外周相加后取半，再乘以环宽得到面积。内周62步加外周4步得66步，取半得33步。乘以环宽12步，得396步作为面积。1亩等于240步，所以面积是1亩156步。（注：此处原文计算有出入。）%

}



\commentarypair{%

环田者，犹圆田之中去一圆也。并中外周而半之者，取环之平均周长也。以径乘之者，以平均周长为广，径为纵也。此术以环田近似为矩形，故以广纵相乘得积。若精密求之，当以两圆面积之差求之：以中周求中圆之径，以外周求外圆之径，两圆面积之差即环田之积。%

}{%

环形田就像一个圆去掉中间一个同心圆。将内外周相加取半，就是取环的平均周长。乘以环宽，就是以平均周长为宽、环宽为高，近似当作矩形来算面积。如果要精确计算，应该用两个圆面积的差来求：从内周求内圆直径，从外周求外圆直径，两个圆面积之差就是环田的精确面积。刘徽指出了近似公式和精确公式。%

}



\vspace{0.5em}

\noindent\makebox[\textwidth]{\textcolor{rulecolor}{\rule{\textwidth}{0.4pt}}}

\vspace{0.5em}



%% --- Problem 15 ---

\probheader{第十五题}{Problem 15: Circular Timber Diameter}



\wenpair{%

今有圆材，埋在地下，不知大小。以锯锯之，深一寸，锯道长一尺。问：圆材径几何？%

}{%

今有一段圆形木材，埋在地下，不知道大小。用锯子锯一道，锯入1寸深，锯道长1尺。问：这段圆材的直径是多少？%

}



\dapair{%

圆材径二尺六寸。%

}{%

圆材直径2尺6寸。%

}



\shupair{%

术曰：以勾股术入之。锯道长一尺者，弦也；深一寸者，矢也。半锯道自乘，除以深，加深，即圆径。半锯道五寸自乘得二十五寸，除以深一寸得二十五寸，加深一寸得二十六寸，即圆材径。%

}{%

用勾股方法求解。锯道长1尺就是弦，锯入深度1寸就是矢（弦到圆弧的垂直距离）。半锯道长自乘，除以深度，再加上深度，就是圆直径。半锯道5寸自乘得25寸，除以深度1寸得25寸，加深度1寸得26寸，就是圆材直径。%

}



\commentarypair{%

圆材埋在地下，以锯锯之者，锯截圆材得一弦也。锯道长一尺即弦长，深一寸即弦到圆周之最短距离（矢）。以勾股术求圆径：半弦为勾，圆径减半矢之半为股。半弦自乘除以矢，得圆径之半减矢。加以矢，得圆径。此术精妙，以勾股之理测圆之大小。%

}{%

圆材埋在地下，用锯锯一道就是截圆材得到一条弦。锯道长1尺就是弦长，深1寸就是弦到圆周的最短距离（矢）。用勾股方法求圆直径：半弦长作为勾，圆半径减去矢作为股。半弦自乘除以矢，得到圆半径减去矢的数值。加上矢，得到圆直径。这个方法很精妙，用勾股原理来测量圆的大小。刘徽推导了公式 $d = \frac{(c/2)^2}{s} + s$，其中$c$为弦长，$s$为矢。%

}



\begin{modernmath}

弦长 $c = 10$ 寸，矢 $s = 1$ 寸，则圆直径：

\[d = \frac{(c/2)^2}{s} + s = \frac{25}{1} + 1 = 26\text{ 寸}\]

\end{modernmath}



\vspace{0.5em}

\noindent\makebox[\textwidth]{\textcolor{rulecolor}{\rule{\textwidth}{0.4pt}}}

\vspace{0.5em}



%% --- Problem 16 ---

\probheader{第十六题}{Problem 16: Square Root Extraction}



\wenpair{%

今有积五万五千二百二十五步。问：为方几何？%

}{%

今有一个面积55225步。问：这个正方形的边长是多少？%

}



\dapair{%

二百三十五步。%

}{%

边长235步。%

}



\shupair{%

术曰：开方术。置积为实，借一算步之，超一等。议所得，以一乘所借一算为法，而以除。除已，倍法为定法。其复除，折法而下。复置借算步之如初，以复议一乘之，所得副，以加定法，以除。以所得副从定法。%

}{%

方法如下：开平方方法。将面积作为被除数（实），借一个算筹表示位数，超一等（每两位一节）。估计第一位商数，用商数乘借算作为除数（法），进行除法。除完后，将法加倍作为定法。再除时，将法降一位。再放置借算如前，重新估计下一位，乘之，所得之数作为副，加在定法上，再除。将副加入定法。%

}



\commentarypair{%

开方术者，求正方形之一边也。积五万五千二百二十五步，开方得二百三十五步。其术：置积为实，借一算为初商之位。超一等者，自右而左，每两位为一节也。议所得者，估计初商之数。二百三十五步者，初商二百，次商三十，三商五。倍法为定法者，以已得之二倍之，为下次除法之基数。此术与今之开方术原理相同，惟以筹算布位，故步骤不同耳。开方不尽者，以面命之；若精密求之，以微数续之。%

}{%

开平方就是求正方形的边长。面积55225步，开平方得235步。方法：将面积作为实，借算筹作为初商的位置。超一等就是从右向左每两位为一节。估计初商：这里初商是200，次商是30，三商是5。倍法为定法就是将已得的数加倍作为下次除法的基数。这个方法和现代开平方原理相同，只是用算筹布置，所以步骤不同。开不尽时用分数表示；如果要精确计算，就用小数继续开下去。刘徽提出了``微数''（小数）的概念来处理开不尽的情况，这是十进小数的雏形。%

}



\vspace{0.5em}

\noindent\makebox[\textwidth]{\textcolor{rulecolor}{\rule{\textwidth}{0.4pt}}}

\vspace{0.5em}



%% --- Problem 17 ---

\probheader{第十七题}{Problem 17: Circle from Area}



\wenpair{%

今有积三百六十步。问：为圆周几何？%

}{%

今有一个圆面积360步。问：这个圆的周长是多少？%

}



\dapair{%

圆周六十步。%

}{%

圆周是60步。%

}



\shupair{%

术曰：开圆术。置积为实，以十二乘之，开方除之，得圆周。积三百六十步，以十二乘之得四千三百二十，开方除之得六十步。%

}{%

开圆方法：将面积作为被除数（实），乘以12，开平方，得到圆周。面积360步，乘以12得4320，开平方得60步。%

}



\commentarypair{%

开圆术者，以圆积求圆周也。以十二乘积开方者，古术以圆周率为三，圆面积为$\frac{C^2}{12}$，故$C = \sqrt{12S}$。以十二乘积而开方，得圆周。然此以周三径一为率，其实圆周率非尽于三也。徽以为圆周率当大于三，以割圆术求之：割之弥细，所失弥少；割之又割，以至于不可割，则与圆周合体而无所失矣。以一百九十二边形之面积推之，得圆周率约为三点一四一六，此为密率之近似也。%

}{%

开圆方法就是用圆面积求圆周。乘以12再开方，是因为古代以圆周率为3，圆面积公式为$\frac{C^2}{12}$，所以$C = \sqrt{12S}$。乘以12再开方就得到圆周。但这以周三径一（$\pi \approx 3$）为比率，实际上圆周率不恰好等于3。我认为圆周率应大于3，用割圆术来求：分割得越细，误差越小；不断地分割，直到不可再分，就与圆周完全重合而没有任何误差了。用192边形的面积推算，得到圆周率约为3.1416，这是较精密的近似值。刘徽的割圆术是中国古代数学最伟大的成就之一，他得到了$\pi \approx 3.1416$的精确近似。%

}



\fullwidthnote{\textbf{历史意义}：本题引用了刘徽著名的``割圆术''，他用正192边形逼近圆面积，得到圆周率$\pi \approx 3.1416$，这是古代中国数学最伟大的成就之一。}





%% --- Problem 18 ---

\probheader{第十八题}{Problem 18: Oblique Field Area}



\wenpair{%

今有邪田，一头广三十步，一头广四十二步，正从六十四步。问：为田几何？%

}{%

今有一块梯形田（邪田），一头宽30步，另一头宽42步，正高（垂直距离）64步。问：这块田的面积是多少？%

}



\dapair{%

九亩一百四十四步。%

}{%

面积为9亩144步。%

}



\shupair{%

术曰：并两广而半之，以从乘之为积。三十步与四十二步并之得七十二步，半之得三十六步。以从六十四步乘之，得二千三百四步为积。亩法二百四十步，得九亩一百四十四步。%

}{%

方法如下：将两头宽度相加后取半，再乘以高度得到面积。30步加42步得72步，取半得36步。乘以高度64步，得2304步作为面积。1亩等于240步，所以面积是9亩144步。%

}



\commentarypair{%

邪田者，梯形之田也。一头广三十步，一头广四十二步，两头广狭不同。正从六十四步者，两广之间之垂直距离也。并两广而半之者，取其平均广也。以平均广乘从，得积。此术以梯形近似为矩形，以平均广为矩形之广也。若以勾股之術精密求之，当以两广之差及从之关系求之。%

}{%

邪田就是梯形田。一头宽30步，一头宽42步，两头宽窄不同。正高64步就是两头之间的垂直距离。将两头宽度相加取半，就是取平均宽度。用平均宽度乘以高，得到面积。这个方法把梯形近似为矩形，用平均宽度作为矩形的宽。刘徽指出梯形面积公式本质上是用平均宽度转化为矩形。%

}



\vspace{0.5em}

\noindent\makebox[\textwidth]{\textcolor{rulecolor}{\rule{\textwidth}{0.4pt}}}

\vspace{0.5em}



%% --- Problem 19 ---

\probheader{第十九题}{Problem 19: Arc Field Area}



\wenpair{%

今有弧田，弦三十步，矢十五步。问：为田几何？%

}{%

今有一块弓形田（弧田），弦长30步，矢高15步。问：这块田的面积是多少？%

}



\dapair{%

一亩九十七步半。%

}{%

面积为1亩97.5步。%

}



\shupair{%

术曰：以弦乘矢，矢又自乘，并之，三而一。%

}{%

方法如下：弦长乘以矢高，再加上矢高的平方，两个数相加后除以3。%

}



\commentarypair{%

弧田术者，以弦矢求弧田之面积也。弦乘矢者，以弦为广，矢为纵，相乘得长方形之积。矢自乘者，以矢为方之边，自乘得正方之积。并之者，合两积为一。三而一者，以弧田之积约当此两积和之三分之一也。此术为弧田之近似术。%

}{%

弧田方法就是用弦和矢来求弓形面积。弦乘矢就是以弦为宽、矢为高，相乘得到长方形面积。矢自乘就是以矢为边长得到正方形面积。把两个面积加起来，再除以3，就是弓形田的近似面积。因为弓形面积大约等于这两个面积之和的三分之一。刘徽再次指出这是近似方法，并解释了为什么是除以3。%

}



\vspace{0.5em}

\noindent\makebox[\textwidth]{\textcolor{rulecolor}{\rule{\textwidth}{0.4pt}}}

\vspace{0.5em}



%% --- Problem 20 ---

\probheader{第二十题}{Problem 20: Five Families Building a Dike}



\wenpair{%

今有五家共堤。甲、乙、丙、丁、戊五家，各以所出徒之数筑堤。甲出徒二百五十四人，乙出徒二百四十三人，丙出徒二百三十二人，丁出徒二百二十一人，戊出徒二百一十人。堤长一万二千七百步。问：各筑几何？%

}{%

今有五家共同修筑一道堤坝。甲、乙、丙、丁、戊五家各自按派出的人数筑堤。甲派出254人，乙派出243人，丙派出232人，丁派出221人，戊派出210人。堤坝总长12700步。问：各家各修筑多少步？%

}



\dapair{%

各以人数为率分配之。%

}{%

按各家人数的比例分配。%

}



\shupair{%

术曰：以衰分术入之。并五人数为法，堤长为实，实如法得一步之率。各以其人数乘之，得各家所筑之步数。%

}{%

用衰分术（比例分配法）求解。将五家的人数相加作为除数（法），堤长作为被除数（实），实除以法得到每步每人应修的步数。然后各乘以各家的人数，得到各家应修筑的步数。%

}



\commentarypair{%

五家共堤者，五家合出徒役共筑一堤也。各以人数为率者，按出徒之人数比例分配筑堤之长也。以衰分术入之：并五家出徒之人数为法，以堤长为实，实如法得每人之步数。然后各以其人数乘之。此题本属衰分，而以勾股之比例理推之，同出一理。%

}{%

五家共堤就是五家合出民工共同修筑一道堤坝。各以人数为标准就是按派出人数的比例分配筑堤长度。用衰分术：将五家派出人数相加作为法，以堤长作为实，实除以法得到每人应修的步数。然后各乘以人数。这道题本属衰分，但所用的比例原理和勾股定理出自同一数学原理。刘徽说明了不同数学方法的内在统一性。%

}



\newpage



%% ========================================================================

%% END MATTER

%% ========================================================================

\section*{后记 / Epilogue}



\addcontentsline{toc}{section}{后记 / Epilogue}



\fullwidthnote{这里呈现的三卷——盈不足（双假设法/试位法）、方程（线性方程组/矩阵方法）和勾股（直角三角形/毕达哥拉斯定理）——代表了《九章算术》所保存的中国古代数学思想的最高成就。}



\bilingualpair{%

\textbf{刘徽注}（约263年）将《九章算术》从一本解题食谱提升为真正具有数学推理的著作。他对盈不足法作为一般逼近技术的解释、对方程术（本质上即高斯消元法）的系统性处理、以及他在勾股定理上的深邃工作——包括著名的``割补术''和``重差术''——都充分证明了公元三世纪中国数学的高度发达。%

}{%

刘徽的注释（约263年）将《九章算术》从一本解题方法汇编提升为真正具有数学推理的著作。他对盈不足法作为一般逼近技术的解释、对方程法（本质上就是高斯消元法）的系统性处理、以及他在勾股定理上的深邃工作——包括著名的``割补术''和``重差术''——都充分证明了公元三世纪中国数学的高度发达和精妙。%

}



\bilingualpair{%

\textbf{方程章}值得特别关注，因为它包含了世界数学史上最早的负数系统——刘徽对如何在消元过程中处理``正''（正）和``负''（负）量的解释异常清晰。他以红色算筹表示正数、黑色算筹表示负数。%

}{%

\textbf{方程章}值得特别关注，因为它包含了世界数学史上最早的系统性负数运算法则——刘徽对如何在消元过程中处理``正''和``负''量的解释异常清晰明确。他用红色算筹表示正数、黑色算筹表示负数，开创了带符号数运算的先河。%

}



\bilingualpair{%

\textbf{勾股章}中刘徽的``割圆术''注释代表了逼近圆周率的最早严谨方法之一，通过正192边形得到了约3.1416的精确近似值。%

}{%

\textbf{勾股章}中刘徽的``割圆术''注释代表了人类历史上逼近圆周率的最早严谨方法之一，通过正192边形的面积推算，得到了约3.1416的 remarkably 精确近似值，在世界数学史上占有重要地位。%

}



\vspace{1cm}



\begin{center}

\rule{0.4\textwidth}{0.4pt}



\vspace{0.5cm}



{\small\textit{使用 \LaTeX\ + XeLaTeX + xeCJK 排版}}



{\small\textit{中文字体: modern Unicode fonts}}



{\small\textit{文言文 / 白话文对照版}}



{\small\textit{基于《九章算术》标准版本}}



\vspace{0.5cm}



{\small 刘徽注 / 白话文今译}



\end{center}



\end{document}

%% ========================================================================

%% END OF DOCUMENT

%% ========================================================================









% ============================================================

% Source: translations/zh/chinese/li-ye-ceyuan-haijing-fenlei-shishu-vols1-3_classical-modern_bilingual.tex

% ============================================================



% ========================================================================

% 测圆海镜分类释术 (CEYUAN HAIJING FENLEI SHISHU)

% Volumes 1-3: Classical Chinese → Modern Chinese Bilingual Edition

% 文言文 ↔ 白话文 对照版

% ========================================================================

\documentclass[11pt,a4paper]{article}



%% ===== Core Packages =====

\usepackage{fontspec}

\usepackage{polyglossia}

\setmainlanguage{english}

\usepackage{xeCJK}

\usepackage{amsmath,amssymb,amsthm}

\usepackage{geometry}

\usepackage{graphicx}

\usepackage{xcolor}

\usepackage{titlesec}

\usepackage{fancyhdr}

\usepackage{enumitem}

\usepackage{array,booktabs,longtable,multirow}

\usepackage{hyperref}

\usepackage{tocloft}

\usepackage{indentfirst}

\usepackage{tikz}

\usepackage{paracol}

\usepackage{ifthen}



%% ===== Page Geometry =====

\geometry{

  paperwidth=210mm,paperheight=297mm,

  left=18mm,right=18mm,top=25mm,bottom=25mm,

  headheight=14pt,footskip=30pt

}



%% ===== Font Configuration =====

\setmainfont{Times New Roman}

\setsansfont{Arial}

\setmonofont{Courier New}



%% ===== CJK Font Setup =====

\setCJKmainfont{Microsoft YaHei}

\setCJKsansfont{Microsoft YaHei}

\setCJKmonofont{Microsoft YaHei}



%% ===== Colors =====

\definecolor{titlecolor}{RGB}{45,55,72}

\definecolor{sectioncolor}{RGB}{55,65,81}

\definecolor{notecolor}{RGB}{120,53,15}

\definecolor{rulecolor}{RGB}{180,160,140}

\definecolor{coverblue}{RGB}{35,55,85}

\definecolor{classicalbg}{RGB}{252,250,248}

\definecolor{modernbg}{RGB}{248,252,255}

\definecolor{classicalframe}{RGB}{160,140,120}

\definecolor{modernframe}{RGB}{120,140,160}



%% ===== Column Ratio =====

\columnratio{0.48,0.48}

\setlength{\columnsep}{0.04\textwidth}



%% ===== Section Formatting =====

\titleformat{\section}

  {\normalfont\Large\bfseries\color{sectioncolor}\CJKfamily{}}

  {\thesection}{1em}{}

\titleformat{\subsection}

  {\normalfont\large\bfseries\color{sectioncolor}}

  {\thesubsection}{1em}{}

\titleformat{\subsubsection}

  {\normalfont\normalsize\bfseries\color{sectioncolor}}

  {\thesubsubsection}{1em}{}



%% ===== Header/Footer =====

\pagestyle{fancy}

\fancyhf{}

\fancyhead[L]{\small\textit{测圆海镜分类释术 · 卷\leftmark}}

\fancyhead[R]{\small\thepage}

\renewcommand{\headrulewidth}{0.4pt}

\renewcommand{\footrulewidth}{0pt}



%% ===== Custom Commands =====

\newcommand{\longline}{\noindent\rule{\textwidth}{0.4pt}}

\newcommand{\shortline}{\noindent\rule{0.5\textwidth}{0.4pt}\par}



%% Classical Chinese column commands

\newcommand{\wen}[1]{{\noindent\textbf{问} #1\par}}

\newcommand{\da}[1]{{\noindent\textbf{答曰} #1\par}}

\newcommand{\shu}[1]{{\noindent\textbf{术曰} #1\par}}

\newcommand{\shi}[1]{{\noindent\textbf{释曰} #1\par}}



%% Modern Chinese column commands

\newcommand{\wenm}[1]{{\noindent\textbf{【问题】} #1\par}}

\newcommand{\dam}[1]{{\noindent\textbf{【解答】} #1\par}}

\newcommand{\shum}[1]{{\noindent\textbf{【方法】} #1\par}}

\newcommand{\shim}[1]{{\noindent\textbf{【释义】} #1\par}}



%% Problem environment

\newcounter{problem}[section]

\newenvironment{problem}[1]{%

  \refstepcounter{problem}%

  \medskip\noindent\textbf{第\theproblem 题#1}\quad

}{\medskip}



%% Modern problem environment

\newenvironment{problemmodern}[1]{%

  \medskip\noindent\textbf{【第\theproblem 题】#1}\quad

}{\medskip}



%% Column headers

\newcommand{\colheaders}{%

  \noindent\begin{minipage}[t]{0.48\textwidth}\centering

    \textbf{\color{classicalframe}【文言文原文】}

    \vspace{0.3em}

    \hrule height 0.6pt

  \end{minipage}%

  \hfill

  \begin{minipage}[t]{0.48\textwidth}\centering

    \textbf{\color{modernframe}【白话文译文】}

    \vspace{0.3em}

    \hrule height 0.6pt

  \end{minipage}

  \par\medskip

}



%% Bilingual section header

\newcommand{\bilingualsection}[2]{%

  \section*{#1}

  \addcontentsline{toc}{section}{#1（#2）}

  \markboth{#1}{#1}

  \begin{center}

    {\large\color{sectioncolor} #1}\\[0.3em]

    {\normalsize\color{notecolor} #2}

  \end{center}

  \medskip

}



%% Bilingual subsection header

\newcommand{\bilingualsubsection}[2]{%

  \subsection*{#1}

  \addcontentsline{toc}{subsection}{#1（#2）}

  \begin{center}

    {\large\color{sectioncolor} #1}\\[0.2em]

    {\normalsize\color{notecolor} #2}

  \end{center}

  \medskip

}



%% TOC Depth

\setcounter{tocdepth}{3}

\setcounter{secnumdepth}{3}



%% Hyperref Setup

\hypersetup{

  colorlinks=true,

  linkcolor=sectioncolor,

  citecolor=sectioncolor,

  urlcolor=sectioncolor,

  pdfencoding=auto,

  pdftitle={测圆海镜分类释术 卷一至卷三 文白对照版},

  pdfauthor={李冶 原著；顾应祥 分类释术；现代白话翻译}

}



\begin{document}



%% ========================================================================

%% TITLE PAGE

%% ========================================================================

\thispagestyle{empty}

\begin{center}

\vspace*{1.5cm}



{\fontsize{32}{40}\selectfont\CJKfamily{}

\textcolor{titlecolor}{测圆海镜分类释术}}



\vspace{0.8cm}

{\Large\textcolor{sectioncolor}{文言文↔白话文 对照版}}



\vspace{0.5cm}

{\large\textcolor{notecolor}{Volumes I--III (卷一至卷三)}}



\vspace{1.5cm}

\shortline



\vspace{1cm}

{\large

\begin{tabular}{rl}

\textbf{原著：} & （元）李冶 撰 \\

\textbf{分类释术：} & （明）顾应祥 \\

\textbf{白话翻译：} & 现代汉语译文 \\

\end{tabular}

}



\vspace{1.2cm}

{\normalsize\textcolor{notecolor}{%

测圆海镜分类释术：按数学方法重新编排的《测圆海镜》\\[0.3em]

将170道问题按算法分类，而非按几何构型分类，\\[0.3em]

并附有详细的释术（方法解释）。

}}



\vspace{1.5cm}

\shortline



\vspace{0.8cm}

{\normalsize\textcolor{sectioncolor}{%

文渊阁《四库全书荟要》本\\[0.3em]

子部——天文算法类}}



\vspace{0.8cm}

{\small\textcolor{sectioncolor}{现代排版 bilingual edition}}



\end{center}

\clearpage



%% ========================================================================

%% COLOPHON PAGE

%% ========================================================================

\thispagestyle{empty}

\begin{center}

\vspace*{2.5cm}

{\LARGE 钦定四库全书荟要}



\vspace{1.2cm}

{\Large 子部}



\vspace{1.5cm}

{\Large 测圆海镜分类释术卷一至三}



\vspace{0.5cm}

{\normalsize\color{notecolor} 文言文↔白话文 对照版}



\vspace{1.5cm}

\shortline



\vspace{1cm}

{\normalsize

\begin{tabular}{ll}

详校官主事 & 臣陈永 \\

\end{tabular}

}



\vspace{2.5cm}

{\small\textcolor{notecolor}{%

此版为清乾隆年间《四库全书荟要》抄本\\

（《四库全书荟要》为乾隆皇帝的皇家图书馆精选本）\\[0.5em]

白话文译文为现代学者所作，旨在帮助当代读者\\

理解古代数学经典的内容与方法。

}}



\end{center}

\clearpage



%% ========================================================================

%% TABLE OF CONTENTS

%% ========================================================================

\tableofcontents

\clearpage



%% ========================================================================

%% COLUMN HEADER STYLE

%% ========================================================================

\colheaders



%% ========================================================================

%% BEGIN PARACOL ENVIRONMENT

%% ========================================================================

\begin{paracol}{2}



%% ========================================================================

%% IMPERIAL SUMMARY (提要)

%% ========================================================================

\bilingualsection{提要}{Imperial Summary / 内容提要}



\switchcolumn*

\noindent

臣等谨案：\textbf{测圆海镜分类释术}十卷，元翰林学士\textbf{李冶}撰，明刑部尚书\textbf{顾应祥}分类释术。



\switchcolumn

\noindent

【现代翻译】臣等谨按：《测圆海镜分类释术》共十卷，由元代翰林学士\textbf{李冶}撰写，明代刑部尚书\textbf{顾应祥}进行分类并解释其方法。



\switchcolumn*



冶书前列\textbf{勾股总图}，以天地日月山川旦夕朱青泛泉等字及四方五行八卦之名于图线之交，详为标识。又于交处所成各形，以通商功之类。是以形求积则方田商功之类是也，以积求形则少广勾股之类是也。以形求积者，先得其形而复求其积，故其为术也易。以积求形则先得其积而后求其长短广狭斜正之形，有非乘除所能尽者，故必以商除之。然而商除亦不能尽也，而又立正负廉隅之法以增损附益之，故其为术也难。



\switchcolumn



李冶的书前面列有\textbf{勾股总图}（直角三角形总图），用天地日月山川旦夕朱青泛泉等字以及四方、五行、八卦的名称标注在图中各线条的交点处，详细作出标记。又在各交点所形成的不同图形之间，用通商功等方法联系起来。以图形求面积，就是方田（求田地面积）、商功（求工程体积）这类方法；以面积求图形，就是少广（由面积求边长）、勾股（直角三角形）这类方法。以图形求面积，是先知道图形再去求面积，所以这种方法比较容易。以面积求图形，则是先知道面积再去求其长短、宽窄、斜正的形状，有些不是简单的乘除法能够解决的，所以必须用除法开方。然而除法开方也不能完全解决，于是又创立了正负廉隅（系数）的增损方法来加以补充，所以这种方法比较难。



\switchcolumn*



余自幼好习数学，晚得荆川唐太史所录\textbf{测圆海镜}一书，乃元翰林学士栾城李公冶所著。虽专主于勾股求容圆容方一术，然其中间如平方、立方、三乘方、带从、减从、益廉、减廉、正隅、负隅诸法，凡所谓以积求形者，尽之矣。但其每条下细草，俱径立天元一，反复合之，而无下手之术。使后学之士，茫然无门路可入，辄不自揆，每章去其细草，立一算术。又以其所立通勾边股之属各以类分之，语义稍繁者略加芟损，名曰\textbf{测圆海镜分类释术}。匪敢僭改前贤著述，惟以便下学云尔。



\switchcolumn



我自幼喜好数学，晚年得到荆川唐太史所抄录的《\textbf{测圆海镜}》一书，这是元代翰林学士栾城李先生李冶的著作。虽然专门讨论勾股求容圆（直角三角形内切圆）、容方这一方法，但其中如开平方、开立方、三乘方（三次方程）、带从（带有从项的开方）、减从、益廉（增加系数）、减廉、正隅、负隅等各种方法，凡是所谓以面积求图形的，都讲尽了。只是每条下面的详细演算，都是直接设天元一（设未知数），反复配合，却没有入手的具体方法。使后来的学者茫然没有门路可入。于是我不自量力，每章去掉详细的演算过程，设立一套计算方法。又把其中所涉及的通勾、边股等各以类别分开，语义稍微繁琐的略加删削，命名为《\textbf{测圆海镜分类释术}》。不敢妄自更改前贤的著述，只是为了方便后学而已。



\switchcolumn*



今夫世之论者，俱视为末艺，故高明者不屑为之，而执泥者遂以为占验之法。虽栾城公自序亦以为九九贱伎，殊不知君子之学，自性命道德之外皆艺也。与其徒费精神于占毕之间，又不若留情于此。不惟可以取乐，亦足以为养心之助焉。后之有同此好者，当以余言为然否耶？



\switchcolumn



现在世上讨论数学的人，都把它看作末流技艺，所以高明的人不屑于去做，而拘泥的人又把它当作占卜预测的方法。虽然栾城李先生（李冶）在自序中也认为算术是低贱的技艺，却不知道君子的学问，除了性命道德之外都是技艺。与其白白耗费精神在诵读之间，还不如在这方面留心。不仅可以获得乐趣，也可以作为养心的辅助。后世有与我同样爱好的人，会认为我的话对吗？



\switchcolumn*

\begin{flushright}

嘉靖庚戌夏五月朔\\

吴兴顾应祥志

\end{flushright}



\switchcolumn

\begin{flushright}

嘉靖庚戌年夏五月初一\\

吴兴顾应祥记

\end{flushright}



\switchcolumn*



\end{paracol}

\clearpage

\colheaders

\begin{paracol}{2}



%% ========================================================================

%% PREFACE (原序)

%% ========================================================================

\bilingualsection{原序}{Original Preface / 原序}



\switchcolumn*



天地之所以神变化而生万物者，阴阳而已。一阴一阳，交互错综，而变化无穷焉。圣人因其交互错综之不齐，而置为数术以测之。于是乎天地之高深，日月之出没，鬼神之幽秘，皆可得而知之矣。然数之为术，虽千变万化之不同，而其要不过一开一阖而已。开者除也，阖者乘也。而又有以形求积、以积求形之异。古之为数者，有九九者，其用也。是故用之以贸易则为粟米，用之以分别差等较量远近，则为差分、为均输。因其末而欲知其本，求其故，千岁之日至，可坐而致也。迹其所以求天与星辰之高远，非勾股何以御之？而周官土圭测景之术，所以穷地之垠垠无垠，若指诸掌，亦此术也。岂得以为古奥而弃之不讲乎？



\switchcolumn



天地之所以能神奇地变化而生成万物，不过是阴阳的作用罢了。一阴一阳，交互错综，变化无穷。圣人因为观察到这种交互错综的不规则性，于是设立数术来测算它。这样一来，天地的高深，日月的出没，鬼神的幽秘，都可以知道了。然而数学作为方法，虽然有千变万化，但其关键不过是一开一合罢了。开就是除法，合就是乘法。又有以图形求面积、以面积求图形的区别。古代学习数学的人，从九九乘法开始，这就是它的用途。所以用于贸易就是粟米法，用于分别等级、比较远近，就是差分法、均输法。从末端推知其本源，探求其缘故，千年后的冬至日，可以坐着推算出来。追溯那些测算天与星辰高远的方法，没有勾股术怎么能够掌握呢？而《周官》中用土圭测影的方法，用来穷究大地的无边无际，如指掌一样清晰，也是这个方法。难道能因为它古老深奥就放弃不讲吗？



\switchcolumn*



此公语愚分释此书之盛心也，因请于公录一帙而刻之，梓播之四方，传之后世，以见公之经济不独惠我滇云，而推明朕兆根极领要，以继绝学。俟后贤者尚有考于斯焉。



\switchcolumn



这是先生告诉我分类解释此书的盛意，于是我请求先生录写一册并刻印出来，传播四方，传之后世，以见先生的经世济民之学不仅惠及我们滇云地区，而且能推明根本要领，以继承失传的学问。等待后来的贤者还有所考证。



\switchcolumn*

\begin{flushright}

嘉靖庚戌夏五月朔日\\

古濠沐朝弼谨序

\end{flushright}



\switchcolumn

\begin{flushright}

嘉靖庚戌年夏五月初一\\

古濠沐朝弼谨序

\end{flushright}



\switchcolumn*



\end{paracol}

\clearpage

\colheaders

\begin{paracol}{2}



%% ========================================================================

%% VOLUME 1 (卷一)

%% ========================================================================

\section{卷一}

\markboth{卷一}{卷一}



\switchcolumn

\section*{卷一（现代译文）}

\addcontentsline{toc}{section}{卷一}

\markboth{卷一}{卷一}



\switchcolumn*



卷一专论\textbf{通勾股求容圆}之法，凡十问。通勾股者，以通勾、通股、通弦三者之关系，求圆城之径也。



\switchcolumn



卷一专门讨论\textbf{通勾股求内切圆}的方法，共十道题。所谓通勾股，就是利用通勾（大直角三角形的底边）、通股（大直角三角形的高边）、通弦（大直角三角形的斜边）三者之间的关系，来求圆形城池的直径。



\switchcolumn*



\end{paracol}

\clearpage

\colheaders

\begin{paracol}{2}



%% ========================================================================

%% Chapter 1: Tong gougu rong yuan

%% ========================================================================

\bilingualsubsection{通勾股求容圆}{General Right-Triangle Inscribed Circle Method}



\switchcolumn*



\noindent

\textbf{总论：} 通勾股求容圆者，以勾股相乘倍之为实，勾股和为法，除之得圆径。其术之要，在于以通勾、通股之关系，推求容圆之径也。



\switchcolumn



\noindent

\textbf{【总论】} 通勾股求内切圆的方法是：将勾和股相乘，再乘以二作为被除数（实），勾与股之和作为除数（法），相除得到圆的直径。这个方法的关键在于利用通勾、通股的关系，推求内切圆的直径。



\switchcolumn*



\begin{problem}[\ 通勾股求容圆一]

\wen{甲南行六百步，望乙与城相参直。问城径。}

\da{城径二百四十步。}

\shu{勾股相乘倍之为实，勾股求弦并勾股为弦和。弦和较即容圆径也。}

\shi{此勾股求容圆径也。东行为通勾，南行为通股。以通勾股求通弦和，较弦和较即容圆径也。}

\end{problem}



\switchcolumn



\begin{problemmodern}[\ 通勾股求内切圆 第一题]

\wenm{甲向南行走六百步，望见乙与城池正好成一条直线。问城池的直径是多少？}

\dam{城池直径为二百四十步。}

\shum{将勾（底边）和股（高边）相乘，再乘以二作为被除数（实），勾与股之和作为除数（法）。先求出弦（斜边），再用弦加上勾和股得到"弦和"，"弦和"与弦的差就是内切圆的直径。}

\shim{这是用勾股求内切圆直径的方法。向东走的距离是通勾，向南走的距离是通股。用通勾和通股求通弦和，再用弦和与弦的差就是内切圆的直径。}

\end{problemmodern}



\switchcolumn*



\begin{problem}[\ 通勾股求容圆二]

\wen{甲乙二人俱在城西门。甲南行四百八十步，乙穿城东行二百五十步见之。问城径。}

\da{城径二百四十步。}

\shu{勾股相乘倍之为实，勾股求弦并勾股为弦和。弦和较即容圆径也。}

\shi{此勾上容圆也。南行为边股，东行为边勾也。以边勾股求容圆，与通勾股求通圆同法。}

\end{problem}



\switchcolumn



\begin{problemmodern}[\ 通勾股求内切圆 第二题]

\wenm{甲乙两人都在城池的西门。甲向南走四百八十步，乙穿过城池向东走二百五十步后看见甲。问城池的直径。}

\dam{城池直径为二百四十步。}

\shum{将勾和股相乘，再乘以二作为被除数（实），勾与股之和作为除数（法），除之得圆径。}

\shim{这是用勾上的内切圆方法。向南走的距离是边股，向东走的距离是边勾。用边勾和边股求内切圆，与用通勾和通股求通圆的方法相同。}

\end{problemmodern}



\switchcolumn*



\begin{problem}[\ 通勾股求容圆三]

\wen{甲出北门东行，乙出西门南行，相望见之。甲东行二百步，乙南行四百二十五步。问城径。}

\da{城径二百四十步。}

\shu{勾股相乘倍之为实，勾股求弦并勾股为弦和。弦和较即容圆径也。}

\shi{此亦边勾股求容圆也。以边勾股之数，约之得圆径。}

\end{problem}



\switchcolumn



\begin{problemmodern}[\ 通勾股求内切圆 第三题]

\wenm{甲从北门出发向东走，乙从西门出发向南走，两人互相望见。甲向东走二百步，乙向南走四百二十五步。问城池的直径。}

\dam{城池直径为二百四十步。}

\shum{将勾和股相乘，再乘以二作为被除数（实），勾与股之和作为除数（法），除之得圆径。}

\shim{这也是用边勾和边股求内切圆。用边勾和边股的数值，按比例推算得到圆径。}

\end{problemmodern}



\switchcolumn*



\begin{problem}[\ 通勾股求容圆四]

\wen{甲乙二人俱在城中心立。乙穿城东行一百三十六步，甲出南门东行二百步见之。问城径。}

\da{城径二百四十步。}

\shu{勾股相乘倍之为实，勾股求弦并勾股为弦和。弦和较即容圆径也。}

\shi{此皇极勾股求容圆也。以通勾股之半为皇极勾股，其容圆径之半即城半径也。}

\end{problem}



\switchcolumn



\begin{problemmodern}[\ 通勾股求内切圆 第四题]

\wenm{甲乙两人都在城池中心站立。乙穿过城池向东走一百三十六步，甲从南门出发向东走二百步后看见乙。问城池的直径。}

\dam{城池直径为二百四十步。}

\shum{将勾和股相乘，再乘以二作为被除数（实），勾与股之和作为除数（法），除之得圆径。}

\shim{这是用皇极勾股求内切圆。通勾和通股的一半就是皇极勾和皇极股，其内切圆直径的一半就是城池的半径。}

\end{problemmodern}



\switchcolumn*



\begin{problem}[\ 通勾股求容圆五]

\wen{甲出南门东行，乙出东门南行，相望见之。甲东行七十二步，乙南行三十步。问城径。}

\da{城径二百四十步。}

\shu{勾股相乘倍之为实，勾股求弦并勾股为弦和。弦和较即容圆径也。}

\shi{此太虚勾股求容圆也。以勾股之数约之得容圆径。}

\end{problem}



\switchcolumn



\begin{problemmodern}[\ 通勾股求内切圆 第五题]

\wenm{甲从南门出发向东走，乙从东门出发向南走，两人互相望见。甲向东走七十二步，乙向南走三十步。问城池的直径。}

\dam{城池直径为二百四十步。}

\shum{将勾和股相乘，再乘以二作为被除数（实），勾与股之和作为除数（法），除之得圆径。}

\shim{这是用太虚勾股求内切圆。用勾股的数值按比例推算，得到内切圆直径。}

\end{problemmodern}



\switchcolumn*



\begin{problem}[\ 通勾股求容圆六]

\wen{甲出西门南行，乙出南门东行，相望见之。甲南行八十步，乙东行五十一步。问城径。}

\da{城径二百四十步。}

\shu{勾股相乘倍之为实，勾股求弦并勾股为弦和。弦和较即容圆径也。}

\shi{此明勾股求容圆也。以明勾股之数，约之得容圆径。}

\end{problem}



\switchcolumn



\begin{problemmodern}[\ 通勾股求内切圆 第六题]

\wenm{甲从西门出发向南走，乙从南门出发向东走，两人互相望见。甲向南走八十步，乙向东走五十一步。问城池的直径。}

\dam{城池直径为二百四十步。}

\shum{将勾和股相乘，再乘以二作为被除数（实），勾与股之和作为除数（法），除之得圆径。}

\shim{这是用明勾股求内切圆。用明勾和明股的数值，按比例推算得到内切圆直径。}

\end{problemmodern}



\switchcolumn*



\begin{problem}[\ 通勾股求容圆七]

\wen{甲出东门北行，乙出北门西行，相望见之。甲北行三十步，乙西行一十六步。问城径。}

\da{城径二百四十步。}

\shu{勾股相乘倍之为实，勾股求弦并勾股为弦和。弦和较即容圆径也。}

\shi{此重勾股求容圆也。以重勾股之数，约之得容圆径。}

\end{problem}



\switchcolumn



\begin{problemmodern}[\ 通勾股求内切圆 第七题]

\wenm{甲从东门出发向北走，乙从北门出发向西走，两人互相望见。甲向北走三十步，乙向西走十六步。问城池的直径。}

\dam{城池直径为二百四十步。}

\shum{将勾和股相乘，再乘以二作为被除数（实），勾与股之和作为除数（法），除之得圆径。}

\shim{这是用重勾股求内切圆。用重勾和重股的数值，按比例推算得到内切圆直径。}

\end{problemmodern}



\switchcolumn*



\begin{problem}[\ 通勾股求容圆八]

\wen{甲出东门直行，乙出南门直行，相望见之。甲东行三十步，乙南行十六步。问城径。}

\da{城径二百四十步。}

\shu{勾股相乘倍之为实，勾股求弦并勾股为弦和。弦和较即容圆径也。}

\shi{此虚勾股求容圆也。以虚勾股之数约之，得圆径。}

\end{problem}



\switchcolumn



\begin{problemmodern}[\ 通勾股求内切圆 第八题]

\wenm{甲从东门出发直走，乙从南门出发直走，两人互相望见。甲向东走三十步，乙向南走十六步。问城池的直径。}

\dam{城池直径为二百四十步。}

\shum{将勾和股相乘，再乘以二作为被除数（实），勾与股之和作为除数（法），除之得圆径。}

\shim{这是用虚勾股求内切圆。用虚勾和虚股的数值按比例推算，得到圆径。}

\end{problemmodern}



\switchcolumn*



\begin{problem}[\ 通勾股求容圆九]

\wen{甲出南门直行，乙出东门直行，相望见之。甲南行四十八步，乙东行二十步。问城径。}

\da{城径二百四十步。}

\shu{勾股相乘倍之为实，勾股求弦并勾股为弦和。弦和较即容圆径也。}

\shi{此以勾股之数求容圆也。以勾股之比例推之，得圆径。}

\end{problem}



\switchcolumn



\begin{problemmodern}[\ 通勾股求内切圆 第九题]

\wenm{甲从南门出发直走，乙从东门出发直走，两人互相望见。甲向南走四十八步，乙向东走二十步。问城池的直径。}

\dam{城池直径为二百四十步。}

\shum{将勾和股相乘，再乘以二作为被除数（实），勾与股之和作为除数（法），除之得圆径。}

\shim{这是用勾股的数值求内切圆。用勾和股的比例关系推算，得到圆径。}

\end{problemmodern}



\switchcolumn*



\begin{problem}[\ 通勾股求容圆十]

\wen{甲出西门直行，乙出北门直行，相望见之。甲西行二百步，乙北行六十步。问城径。}

\da{城径二百四十步。}

\shu{勾股相乘倍之为实，勾股求弦并勾股为弦和。弦和较即容圆径也。}

\shi{此以大差勾股求容圆也。以大差勾股之数约之，得圆径。}

\end{problem}



\switchcolumn



\begin{problemmodern}[\ 通勾股求内切圆 第十题]

\wenm{甲从西门出发直走，乙从北门出发直走，两人互相望见。甲向西走二百步，乙向北走六十步。问城池的直径。}

\dam{城池直径为二百四十步。}

\shum{将勾和股相乘，再乘以二作为被除数（实），勾与股之和作为除数（法），除之得圆径。}

\shim{这是用大差勾股求内切圆。用大差勾和大差股的数值按比例推算，得到圆径。}

\end{problemmodern}



\switchcolumn*



\end{paracol}

\clearpage

\colheaders

\begin{paracol}{2}



%% ========================================================================

%% Chapter 2: Bian gougu rong yuan

%% ========================================================================

\bilingualsubsection{边勾股求容圆}{Edge Right-Triangle Inscribed Circle Method}



\switchcolumn*



\noindent

\textbf{总论：} 边勾股求容圆者，以边勾、边股之数，用通勾股之法，求容圆之径也。



\switchcolumn



\noindent

\textbf{【总论】} 边勾股求内切圆的方法是：用边勾和边股的数值，采用通勾股的方法，求内切圆的直径。



\switchcolumn*



\begin{problem}[\ 边勾股求容圆一]

\wen{甲出北门东行二百步，乙出东门南行一百二十步，相望见之。问城径。}

\da{城径二百四十步。}

\shu{边勾股相乘倍之为实，边勾股求弦并边勾股为弦和。弦和较即容圆径也。}

\shi{此边勾股求容圆也。边勾股者，通勾股之半也。以边勾股之数求容圆，与通勾股同法。}

\end{problem}



\switchcolumn



\begin{problemmodern}[\ 边勾股求内切圆 第一题]

\wenm{甲从北门出发向东走二百步，乙从东门出发向南走一百二十步，两人互相望见。问城池的直径。}

\dam{城池直径为二百四十步。}

\shum{边勾和边股相乘，再乘以二作为被除数（实），边勾与边股之和作为除数（法），除之得圆径。}

\shim{这是用边勾股求内切圆。边勾和边股是通勾和通股的一半。用边勾和边股的数值求内切圆，与通勾股的方法相同。}

\end{problemmodern}



\switchcolumn*



\begin{problem}[\ 边勾股求容圆二]

\wen{甲乙二人俱在城西门。甲南行四百八十步，乙穿城东行二百五十步见之。问城径。}

\da{城径二百四十步。}

\shu{勾股相乘倍之为实，勾股求弦并勾股为弦和。弦和较即容圆径也。}

\shi{此勾上容圆也。南行为边股，东行为边勾也。以边勾股求容圆，与通勾股求通圆同法。}

\end{problem}



\switchcolumn



\begin{problemmodern}[\ 边勾股求内切圆 第二题]

\wenm{甲乙两人都在城池西门。甲向南走四百八十步，乙穿过城池向东走二百五十步后看见甲。问城池的直径。}

\dam{城池直径为二百四十步。}

\shum{将勾和股相乘，再乘以二作为被除数（实），勾与股之和作为除数（法），除之得圆径。}

\shim{这是用勾上的内切圆方法。向南走的距离是边股，向东走的距离是边勾。用边勾和边股求内切圆，与用通勾股求通圆的方法相同。}

\end{problemmodern}



\switchcolumn*



\begin{problem}[\ 边勾股求容圆三]

\wen{甲出北门东行二百步，乙出东门南行一百二十步，相望见之。问城径。}

\da{城径二百四十步。}

\shu{边勾股相乘倍之为实，边勾股求弦并边勾股为弦和。弦和较即容圆径也。}

\shi{此边勾股求容圆也。边勾股者，通勾股之半也。以边勾股之数求容圆，与通勾股同法。}

\end{problem}



\switchcolumn



\begin{problemmodern}[\ 边勾股求内切圆 第三题]

\wenm{甲从北门出发向东走二百步，乙从东门出发向南走一百二十步，两人互相望见。问城池的直径。}

\dam{城池直径为二百四十步。}

\shum{边勾和边股相乘，再乘以二作为被除数（实），边勾与边股之和作为除数（法），除之得圆径。}

\shim{这是用边勾股求内切圆。边勾和边股是通勾和通股的一半。用边勾和边股的数值求内切圆，与通勾股的方法相同。}

\end{problemmodern}



\switchcolumn*



\begin{problem}[\ 边勾股求容圆四]

\wen{甲出南门东行一百二十步，乙出东门南行五十步，相望见之。问城径。}

\da{城径二百四十步。}

\shu{边勾股相乘倍之为实，边勾股求弦并边勾股为弦和。弦和较即容圆径也。}

\shi{此边勾股求容圆也。以边勾股之数求容圆，与通勾股同法。}

\end{problem}



\switchcolumn



\begin{problemmodern}[\ 边勾股求内切圆 第四题]

\wenm{甲从南门出发向东走一百二十步，乙从东门出发向南走五十步，两人互相望见。问城池的直径。}

\dam{城池直径为二百四十步。}

\shum{边勾和边股相乘，再乘以二作为被除数（实），边勾与边股之和作为除数（法），除之得圆径。}

\shim{这是用边勾股求内切圆。用边勾和边股的数值求内切圆，与通勾股的方法相同。}

\end{problemmodern}



\switchcolumn*



\begin{problem}[\ 边勾股求容圆五]

\wen{甲出西门南行一百步，乙出南门西行八十步，相望见之。问城径。}

\da{城径二百四十步。}

\shu{边勾股相乘倍之为实，边勾股求弦并边勾股为弦和。弦和较即容圆径也。}

\shi{此边勾股求容圆也。以边勾股之数求容圆，与通勾股同法。}

\end{problem}



\switchcolumn



\begin{problemmodern}[\ 边勾股求内切圆 第五题]

\wenm{甲从西门出发向南走一百步，乙从南门出发向西走八十步，两人互相望见。问城池的直径。}

\dam{城池直径为二百四十步。}

\shum{边勾和边股相乘，再乘以二作为被除数（实），边勾与边股之和作为除数（法），除之得圆径。}

\shim{这是用边勾股求内切圆。用边勾和边股的数值求内切圆，与通勾股的方法相同。}

\end{problemmodern}



\switchcolumn*



\end{paracol}

\clearpage

\colheaders

\begin{paracol}{2}



%% ========================================================================

%% Chapter 3: Diangu xian qiu tong yuanjing

%% ========================================================================

\bilingualsubsection{底股弦求通圆径}{Base-Leg and Hypotenuse Circle Diameter Method}



\switchcolumn*



\noindent

\textbf{总论：} 底股弦求通圆径者，以底股、底弦之数，用弦较之法，求通圆之径也。



\switchcolumn



\noindent

\textbf{【总论】} 底股弦求通圆直径的方法是：用底股和底弦的数值，采用弦较（斜边与边之差）的方法，求通圆的直径。



\switchcolumn*



\begin{problem}[\ 底股弦求通圆径一]

\wen{问底股弦求通圆径。}

\da{通圆径二百四十步。}

\shu{弦升减股升开其余，得勾如前法求之。}

\shi{此以底股弦之关系，求通圆径也。底股弦者，大勾股弦之属也。}

\end{problem}



\switchcolumn



\begin{problemmodern}[\ 底股弦求通圆直径 第一题]

\wenm{问如何用底股和底弦求通圆直径。}

\dam{通圆直径为二百四十步。}

\shum{用弦的平方减去股的平方，开平方得到勾，再用前面的方法求圆径。}

\shim{这是用底股和底弦的关系求通圆直径。底股和底弦属于大勾股弦的范畴。}

\end{problemmodern}



\switchcolumn*



\begin{problem}[\ 底股弦求通圆径二]

\wen{甲出北门东行，乙出东门南行，相望见之。问以底股弦求通圆径。}

\da{通圆径二百四十步。}

\shu{弦升减股升开其余，得勾如前法求之。}

\shi{此底股弦求通圆径也。以弦升减股升开其余得勾，再以勾股求容圆。}

\end{problem}



\switchcolumn



\begin{problemmodern}[\ 底股弦求通圆直径 第二题]

\wenm{甲从北门出发向东走，乙从东门出发向南走，两人互相望见。问如何用底股弦求通圆直径。}

\dam{通圆直径为二百四十步。}

\shum{用弦的平方减去股的平方，开平方得到勾，再用前面的方法求圆径。}

\shim{这是用底股弦求通圆直径。用弦的平方减去股的平方，开平方得到勾，再用勾股求内切圆。}

\end{problemmodern}



\switchcolumn*



\begin{problem}[\ 底股弦求通圆径三]

\wen{甲出南门直行，乙出东门直行，相望见之。问以底股弦求通圆径。}

\da{通圆径二百四十步。}

\shu{弦升减股升开其余，得勾如前法求之。}

\shi{此底股弦求通圆径也。以弦升减股升开其余得勾，再以勾股求容圆。}

\end{problem}



\switchcolumn



\begin{problemmodern}[\ 底股弦求通圆直径 第三题]

\wenm{甲从南门出发直走，乙从东门出发直走，两人互相望见。问如何用底股弦求通圆直径。}

\dam{通圆直径为二百四十步。}

\shum{用弦的平方减去股的平方，开平方得到勾，再用前面的方法求圆径。}

\shim{这是用底股弦求通圆直径。用弦的平方减去股的平方，开平方得到勾，再用勾股求内切圆。}

\end{problemmodern}



\switchcolumn*



\end{paracol}

\clearpage

\colheaders

\begin{paracol}{2}



%% ========================================================================

%% Chapter 4: Huangji gougu qiu rong yuan

%% ========================================================================

\bilingualsubsection{皇极勾股求容圆}{Imperial-Pole Right-Triangle Inscribed Circle Method}



\switchcolumn*



\noindent

\textbf{总论：} 皇极勾股求容圆者，以皇极勾、皇极股之数，用勾股容圆之法，求圆径也。皇极勾股者，通勾股之半也。



\switchcolumn



\noindent

\textbf{【总论】} 皇极勾股求内切圆的方法是：用皇极勾和皇极股的数值，采用勾股内切圆的方法，求圆的直径。皇极勾和皇极股是通勾和通股的一半。



\switchcolumn*



\begin{problem}[\ 皇极勾股求容圆一]

\wen{甲乙二人俱在城中心立。乙穿城东行一百三十六步，甲出南门东行二百步见之。问城径。}

\da{城径二百四十步。}

\shu{皇极勾股相乘倍之为实，皇极勾股和为法，除之得皇极容圆径，倍之为城径。}

\shi{此皇极勾股求容圆也。以皇极勾股之半为通勾股之半，其容圆径之半即城半径也。}

\end{problem}



\switchcolumn



\begin{problemmodern}[\ 皇极勾股求内切圆 第一题]

\wenm{甲乙两人都在城池中心站立。乙穿过城池向东走一百三十六步，甲从南门出发向东走二百步后看见乙。问城池的直径。}

\dam{城池直径为二百四十步。}

\shum{皇极勾和皇极股相乘，再乘以二作为被除数（实），皇极勾与皇极股之和作为除数（法），除之得到皇极内切圆直径，再乘以二就是城池直径。}

\shim{这是用皇极勾股求内切圆。皇极勾股是通勾股的一半，其内切圆直径的一半就是城池的半径。}

\end{problemmodern}



\switchcolumn*



\begin{problem}[\ 皇极勾股求容圆二]

\wen{甲出北门东行一百步，乙出东门南行八十步，相望见之。问城径。}

\da{城径二百四十步。}

\shu{皇极勾股相乘倍之为实，皇极勾股和为法，除之得皇极容圆径，倍之为城径。}

\shi{此皇极勾股求容圆也。以皇极勾股之数，约之得圆径。}

\end{problem}



\switchcolumn



\begin{problemmodern}[\ 皇极勾股求内切圆 第二题]

\wenm{甲从北门出发向东走一百步，乙从东门出发向南走八十步，两人互相望见。问城池的直径。}

\dam{城池直径为二百四十步。}

\shum{皇极勾和皇极股相乘，再乘以二作为被除数（实），皇极勾与皇极股之和作为除数（法），除之得到皇极内切圆直径，再乘以二就是城池直径。}

\shim{这是用皇极勾股求内切圆。用皇极勾和皇极股的数值，按比例推算得到圆径。}

\end{problemmodern}



\switchcolumn*



\begin{problem}[\ 皇极勾股求容圆三]

\wen{甲出南门东行一百五十步，乙出东门南行一百步，相望见之。问城径。}

\da{城径二百四十步。}

\shu{皇极勾股相乘倍之为实，皇极勾股和为法，除之得皇极容圆径，倍之为城径。}

\shi{此皇极勾股求容圆也。以皇极勾股之数，约之得圆径。}

\end{problem}



\switchcolumn



\begin{problemmodern}[\ 皇极勾股求内切圆 第三题]

\wenm{甲从南门出发向东走一百五十步，乙从东门出发向南走一百步，两人互相望见。问城池的直径。}

\dam{城池直径为二百四十步。}

\shum{皇极勾和皇极股相乘，再乘以二作为被除数（实），皇极勾与皇极股之和作为除数（法），除之得到皇极内切圆直径，再乘以二就是城池直径。}

\shim{这是用皇极勾股求内切圆。用皇极勾和皇极股的数值，按比例推算得到圆径。}

\end{problemmodern}



\switchcolumn*



\begin{problem}[\ 皇极勾股求容圆四]

\wen{甲出西门南行二百步，乙出南门西行一百五十步，相望见之。问城径。}

\da{城径二百四十步。}

\shu{皇极勾股相乘倍之为实，皇极勾股和为法，除之得皇极容圆径，倍之为城径。}

\shi{此皇极勾股求容圆也。以皇极勾股之数，约之得圆径。}

\end{problem}



\switchcolumn



\begin{problemmodern}[\ 皇极勾股求内切圆 第四题]

\wenm{甲从西门出发向南走二百步，乙从南门出发向西走一百五十步，两人互相望见。问城池的直径。}

\dam{城池直径为二百四十步。}

\shum{皇极勾和皇极股相乘，再乘以二作为被除数（实），皇极勾与皇极股之和作为除数（法），除之得到皇极内切圆直径，再乘以二就是城池直径。}

\shim{这是用皇极勾股求内切圆。用皇极勾和皇极股的数值，按比例推算得到圆径。}

\end{problemmodern}



\switchcolumn*



\end{paracol}

\clearpage

\colheaders

\begin{paracol}{2}



%% ========================================================================

%% Chapter 5: Taiyin gougu qiu rong yuan

%% ========================================================================

\bilingualsubsection{太阴勾股求容圆}{Great-Void Right-Triangle Inscribed Circle Method}



\switchcolumn*



\noindent

\textbf{总论：} 太阴勾股求容圆者，以太阴勾、太阴股之数，用勾股容圆之法，求圆径也。



\switchcolumn



\noindent

\textbf{【总论】} 太阴勾股求内切圆的方法是：用太阴勾和太阴股的数值，采用勾股内切圆的方法，求圆的直径。



\switchcolumn*



\begin{problem}[\ 太阴勾股求容圆一]

\wen{甲出南门东行七十二步，乙出东门南行三十步，相望见之。问城径。}

\da{城径二百四十步。}

\shu{太阴勾股相乘倍之为实，太阴勾股和为法，除之得太阴容圆径，以比例推之得城径。}

\shi{此太阴勾股求容圆也。以太阴勾股之数，约之得圆径。}

\end{problem}



\switchcolumn



\begin{problemmodern}[\ 太阴勾股求内切圆 第一题]

\wenm{甲从南门出发向东走七十二步，乙从东门出发向南走三十步，两人互相望见。问城池的直径。}

\dam{城池直径为二百四十步。}

\shum{太阴勾和太阴股相乘，再乘以二作为被除数（实），太阴勾与太阴股之和作为除数（法），除之得到太阴内切圆直径，再用比例推算得到城池直径。}

\shim{这是用太阴勾股求内切圆。用太阴勾和太阴股的数值，按比例推算得到圆径。}

\end{problemmodern}



\switchcolumn*



\begin{problem}[\ 太阴勾股求容圆二]

\wen{甲出北门东行五十步，乙出东门北行二十步，相望见之。问城径。}

\da{城径二百四十步。}

\shu{太阴勾股相乘倍之为实，太阴勾股和为法，除之得太阴容圆径，以比例推之得城径。}

\shi{此太阴勾股求容圆也。以太阴勾股之数，约之得圆径。}

\end{problem}



\switchcolumn



\begin{problemmodern}[\ 太阴勾股求内切圆 第二题]

\wenm{甲从北门出发向东走五十步，乙从东门出发向北走二十步，两人互相望见。问城池的直径。}

\dam{城池直径为二百四十步。}

\shum{太阴勾和太阴股相乘，再乘以二作为被除数（实），太阴勾与太阴股之和作为除数（法），除之得到太阴内切圆直径，再用比例推算得到城池直径。}

\shim{这是用太阴勾股求内切圆。用太阴勾和太阴股的数值，按比例推算得到圆径。}

\end{problemmodern}



\switchcolumn*



\begin{problem}[\ 太阴勾股求容圆三]

\wen{甲出西门南行一百步，乙出南门西行四十步，相望见之。问城径。}

\da{城径二百四十步。}

\shu{太阴勾股相乘倍之为实，太阴勾股和为法，除之得太阴容圆径，以比例推之得城径。}

\shi{此太阴勾股求容圆也。以太阴勾股之数，约之得圆径。}

\end{problem}



\switchcolumn



\begin{problemmodern}[\ 太阴勾股求内切圆 第三题]

\wenm{甲从西门出发向南走一百步，乙从南门出发向西走四十步，两人互相望见。问城池的直径。}

\dam{城池直径为二百四十步。}

\shum{太阴勾和太阴股相乘，再乘以二作为被除数（实），太阴勾与太阴股之和作为除数（法），除之得到太阴内切圆直径，再用比例推算得到城池直径。}

\shim{这是用太阴勾股求内切圆。用太阴勾和太阴股的数值，按比例推算得到圆径。}

\end{problemmodern}



\switchcolumn*



\end{paracol}

\clearpage

\colheaders

\begin{paracol}{2}



%% ========================================================================

%% Chapter 6: Ming gougu qiu rong yuan

%% ========================================================================

\bilingualsubsection{明勾股求容圆}{Bright Right-Triangle Inscribed Circle Method}



\switchcolumn*



\noindent

\textbf{总论：} 明勾股求容圆者，以明勾、明股之数，用勾股容圆之法，求圆径也。



\switchcolumn



\noindent

\textbf{【总论】} 明勾股求内切圆的方法是：用明勾和明股的数值，采用勾股内切圆的方法，求圆的直径。



\switchcolumn*



\begin{problem}[\ 明勾股求容圆一]

\wen{甲出西门南行八十步，乙出南门东行五十一步，相望见之。问城径。}

\da{城径二百四十步。}

\shu{明勾股相乘倍之为实，明勾股和为法，除之得明容圆径，以比例推之得城径。}

\shi{此明勾股求容圆也。以明勾股之数，约之得圆径。}

\end{problem}



\switchcolumn



\begin{problemmodern}[\ 明勾股求内切圆 第一题]

\wenm{甲从西门出发向南走八十步，乙从南门出发向东走五十一步，两人互相望见。问城池的直径。}

\dam{城池直径为二百四十步。}

\shum{明勾和明股相乘，再乘以二作为被除数（实），明勾与明股之和作为除数（法），除之得到明内切圆直径，再用比例推算得到城池直径。}

\shim{这是用明勾股求内切圆。用明勾和明股的数值，按比例推算得到圆径。}

\end{problemmodern}



\switchcolumn*



\begin{problem}[\ 明勾股求容圆二]

\wen{甲出北门西行六十步，乙出西门北行三十九步，相望见之。问城径。}

\da{城径二百四十步。}

\shu{明勾股相乘倍之为实，明勾股和为法，除之得明容圆径，以比例推之得城径。}

\shi{此明勾股求容圆也。以明勾股之数，约之得圆径。}

\end{problem}



\switchcolumn



\begin{problemmodern}[\ 明勾股求内切圆 第二题]

\wenm{甲从北门出发向西走六十步，乙从西门出发向北走三十九步，两人互相望见。问城池的直径。}

\dam{城池直径为二百四十步。}

\shum{明勾和明股相乘，再乘以二作为被除数（实），明勾与明股之和作为除数（法），除之得到明内切圆直径，再用比例推算得到城池直径。}

\shim{这是用明勾股求内切圆。用明勾和明股的数值，按比例推算得到圆径。}

\end{problemmodern}



\switchcolumn*



\end{paracol}

\clearpage

\colheaders

\begin{paracol}{2}



%% ========================================================================

%% Chapter 7: Chong gougu qiu rong yuan

%% ========================================================================

\bilingualsubsection{重勾股求容圆}{Double Right-Triangle Inscribed Circle Method}



\switchcolumn*



\noindent

\textbf{总论：} 重勾股求容圆者，以重勾、重股之数，用勾股容圆之法，求圆径也。



\switchcolumn



\noindent

\textbf{【总论】} 重勾股求内切圆的方法是：用重勾和重股的数值，采用勾股内切圆的方法，求圆的直径。



\switchcolumn*



\begin{problem}[\ 重勾股求容圆一]

\wen{甲出东门北行三十步，乙出北门西行一十六步，相望见之。问城径。}

\da{城径二百四十步。}

\shu{重勾股相乘倍之为实，重勾股和为法，除之得重容圆径，以比例推之得城径。}

\shi{此重勾股求容圆也。以重勾股之数，约之得圆径。}

\end{problem}



\switchcolumn



\begin{problemmodern}[\ 重勾股求内切圆 第一题]

\wenm{甲从东门出发向北走三十步，乙从北门出发向西走十六步，两人互相望见。问城池的直径。}

\dam{城池直径为二百四十步。}

\shum{重勾和重股相乘，再乘以二作为被除数（实），重勾与重股之和作为除数（法），除之得到重内切圆直径，再用比例推算得到城池直径。}

\shim{这是用重勾股求内切圆。用重勾和重股的数值，按比例推算得到圆径。}

\end{problemmodern}



\switchcolumn*



\begin{problem}[\ 重勾股求容圆二]

\wen{甲出南门东行二十四步，乙出东门南行十步，相望见之。问城径。}

\da{城径二百四十步。}

\shu{重勾股相乘倍之为实，重勾股和为法，除之得重容圆径，以比例推之得城径。}

\shi{此重勾股求容圆也。以重勾股之数，约之得圆径。}

\end{problem}



\switchcolumn



\begin{problemmodern}[\ 重勾股求内切圆 第二题]

\wenm{甲从南门出发向东走二十四步，乙从东门出发向南走十步，两人互相望见。问城池的直径。}

\dam{城池直径为二百四十步。}

\shum{重勾和重股相乘，再乘以二作为被除数（实），重勾与重股之和作为除数（法），除之得到重内切圆直径，再用比例推算得到城池直径。}

\shim{这是用重勾股求内切圆。用重勾和重股的数值，按比例推算得到圆径。}

\end{problemmodern}



\switchcolumn*



\end{paracol}

\clearpage

\colheaders

\begin{paracol}{2}



%% ========================================================================

%% Reference Table for Volume 1

%% ========================================================================

\subsection*{卷一名数表}

\addcontentsline{toc}{subsection}{卷一名数表}



\switchcolumn

\subsection*{卷一名数对照表（现代注释）}

\addcontentsline{toc}{subsection}{卷一名数对照表}



\switchcolumn*



\noindent

\textbf{卷一所用名数对照表：}



\switchcolumn

\noindent

\textbf{【卷一所用术语数值对照表：】}



\switchcolumn*



\begin{center}

\begin{tabular}{|l|l|l|}

\hline

\textbf{名称} & \textbf{数值} & \textbf{说明} \\

\hline

勾股和 & 七百三十六 & 通勾与通股之和 \\

\hline

勾弦和 & 八百 & 通勾与通弦之和 \\

\hline

股弦和 & 一千二十四 & 通股与通弦之和 \\

\hline

弦较和 & 七百六十八 & 通弦与通较之和 \\

\hline

弦较较 & 二百八十八 & 通弦与通较之差 \\

\hline

弦和和 & 一千二百八十 & 通弦与通和之和 \\

\hline

弦和较 & 二百二十四 & 通弦与通和之差 \\

\hline

黄广弦 & 五百一十 & 勾二百四十，股四百五十 \\

\hline

勾股和 & 六百九十 & 勾二百四十，股四百五十 \\

\hline

勾弦和 & 七百五十 & 勾二百四十，弦五百一十 \\

\hline

较 & 二百一十 & 勾股之较 \\

\hline

\end{tabular}

\end{center}



\switchcolumn



\begin{center}

\begin{tabular}{|l|l|p{4.5cm}|}

\hline

\textbf{术语} & \textbf{数值} & \textbf{现代说明} \\

\hline

勾股和 & 736 & 通勾（底边）与通股（高边）之和 \\

\hline

勾弦和 & 800 & 通勾与通弦（斜边）之和 \\

\hline

股弦和 & 1024 & 通股与通弦之和 \\

\hline

弦较和 & 768 & 通弦与"通较"之和 \\

\hline

弦较较 & 288 & 通弦与"通较"之差 \\

\hline

弦和和 & 1280 & 通弦与"通和"之和 \\

\hline

弦和较 & 224 & 通弦与"通和"之差 \\

\hline

黄广弦 & 510 & 勾=240，股=450时的弦 \\

\hline

勾股和 & 690 & 勾=240，股=450时的和 \\

\hline

勾弦和 & 750 & 勾=240，弦=510时的和 \\

\hline

较 & 210 & 勾与股之差 \\

\hline

\end{tabular}

\end{center}



\switchcolumn*



\end{paracol}

\clearpage

\colheaders

\begin{paracol}{2}



%% ========================================================================

%% VOLUME 2 (卷二)

%% ========================================================================

\section{卷二}

\markboth{卷二}{卷二}



\switchcolumn

\section*{卷二（现代译文）}

\addcontentsline{toc}{section}{卷二}

\markboth{卷二}{卷二}



\switchcolumn*



卷二专论\textbf{两股求容圆}之法，凡七条。两股求容圆者，以两股之数，用通勾股之法，求圆城之径也。又论\textbf{两弦求容圆}之法，凡三条。



\switchcolumn



卷二专门讨论\textbf{用两股求内切圆}的方法，共七条。两股求内切圆的方法是：用两股的数值，采用通勾股的方法，求圆形城池的直径。另外还讨论\textbf{用两弦求内切圆}的方法，共三条。



\switchcolumn*



\end{paracol}

\clearpage

\colheaders

\begin{paracol}{2}



%% ========================================================================

%% Chapter 1: Liang gu qiu rong yuan

%% ========================================================================

\bilingualsubsection{两股求容圆}{Two-Leg Inscribed Circle Method}



\switchcolumn*



\noindent

\textbf{总论：} 两股求容圆者，以两股之数相乘倍之为实，以两股之和为法，除之得圆径。



\switchcolumn



\noindent

\textbf{【总论】} 两股求内切圆的方法是：用两股的数值相乘，再乘以二作为被除数（实），以两股之和作为除数（法），相除得到圆的直径。



\switchcolumn*



\begin{problem}[\ 两股求容圆一]

\wen{城南有槐一株，城东有柳一株。甲出北门东行，丙出西门南行。既而甲斜行四百二十五步至槐下，丙斜行五百四十四步至柳下。问城径。}

\da{城径二百四十步。}

\shu{二斜行相减余自之，得一万一千四百六十一为差幂。甲斜行自之，得一十八万零六百二十五为幂。以差幂减之，余为实。以二斜行相减为法，除之得半径。}

\shi{此以两股之差求容圆也。甲斜行向西南至槐树下，底弦也；丙斜行向东北至柳树下，边弦也。此以边弦底弦互测圆径。}

\end{problem}



\switchcolumn



\begin{problemmodern}[\ 两股求内切圆 第一题]

\wenm{城南有一棵槐树，城东有一棵柳树。甲从北门出发向东走，丙从西门出发向南走。之后甲斜着走四百二十五步到达槐树下，丙斜着走五百四十四步到达柳树下。问城池的直径。}

\dam{城池直径为二百四十步。}

\shum{两斜行之差自乘，得到11461作为差幂。甲斜行自乘，得到180625作为幂。用差幂去减，余数作为被除数（实）。以两斜行之差作为除数（法），相除得到半径。}

\shim{这是用两股之差求内切圆。甲斜着走向西南到槐树下，这是底弦；丙斜着走向东北到柳树下，这是边弦。这里用边弦和底弦互相测算圆径。}

\end{problemmodern}



\switchcolumn*



\begin{problem}[\ 两股求容圆二]

\wen{甲出南门直行一百三十五步而立，乙从城外西北乾隅南行六百步望乙与城相参直。问城径。}

\da{城径二百四十步。}

\shu{两股相乘倍之为实，并两股为法，除之得圆径。}

\shi{此两股求容圆也。以两股之数，约之得圆径。}

\end{problem}



\switchcolumn



\begin{problemmodern}[\ 两股求内切圆 第二题]

\wenm{甲从南门出发直走一百三十五步后站立，乙从城外西北角向南走六百步，望见乙与城池正好成一条直线。问城池的直径。}

\dam{城池直径为二百四十步。}

\shum{两股相乘，再乘以二作为被除数（实），两股之和作为除数（法），相除得到圆径。}

\shim{这是用两股求内切圆。用两股的数值按比例推算得到圆径。}

\end{problemmodern}



\switchcolumn*



\begin{problem}[\ 两股求容圆三]

\wen{甲出东门直行七十二步，乙出北门直行一百二十步，相望见之。问城径。}

\da{城径二百四十步。}

\shu{两股相乘倍之为实，并两股为法，除之得圆径。}

\shi{此两股求容圆也。以两股之数，约之得圆径。}

\end{problem}



\switchcolumn



\begin{problemmodern}[\ 两股求内切圆 第三题]

\wenm{甲从东门出发直走七十二步，乙从北门出发直走一百二十步，两人互相望见。问城池的直径。}

\dam{城池直径为二百四十步。}

\shum{两股相乘，再乘以二作为被除数（实），两股之和作为除数（法），相除得到圆径。}

\shim{这是用两股求内切圆。用两股的数值按比例推算得到圆径。}

\end{problemmodern}



\switchcolumn*



\begin{problem}[\ 两股求容圆四]

\wen{甲出西门南行一百步，乙出南门东行六十步，相望见之。问城径。}

\da{城径二百四十步。}

\shu{两股相乘倍之为实，并两股为法，除之得圆径。}

\shi{此两股求容圆也。以两股之数，约之得圆径。}

\end{problem}



\switchcolumn



\begin{problemmodern}[\ 两股求内切圆 第四题]

\wenm{甲从西门出发向南走一百步，乙从南门出发向东走六十步，两人互相望见。问城池的直径。}

\dam{城池直径为二百四十步。}

\shum{两股相乘，再乘以二作为被除数（实），两股之和作为除数（法），相除得到圆径。}

\shim{这是用两股求内切圆。用两股的数值按比例推算得到圆径。}

\end{problemmodern}



\switchcolumn*



\begin{problem}[\ 两股求容圆五]

\wen{甲出北门东行二百步，乙出东门南行一百五十步，相望见之。问城径。}

\da{城径二百四十步。}

\shu{两股相乘倍之为实，并两股为法，除之得圆径。}

\shi{此两股求容圆也。以两股之数，约之得圆径。}

\end{problem}



\switchcolumn



\begin{problemmodern}[\ 两股求内切圆 第五题]

\wenm{甲从北门出发向东走二百步，乙从东门出发向南走一百五十步，两人互相望见。问城池的直径。}

\dam{城池直径为二百四十步。}

\shum{两股相乘，再乘以二作为被除数（实），两股之和作为除数（法），相除得到圆径。}

\shim{这是用两股求内切圆。用两股的数值按比例推算得到圆径。}

\end{problemmodern}



\switchcolumn*



\begin{problem}[\ 两股求容圆六]

\wen{甲出南门直行二百步，乙出东门直行三百步，相望见之。问城径。}

\da{城径二百四十步。}

\shu{两股相乘倍之为实，并两股为法，除之得圆径。}

\shi{此两股求容圆也。以两股之数，约之得圆径。}

\end{problem}



\switchcolumn



\begin{problemmodern}[\ 两股求内切圆 第六题]

\wenm{甲从南门出发直走二百步，乙从东门出发直走三百步，两人互相望见。问城池的直径。}

\dam{城池直径为二百四十步。}

\shum{两股相乘，再乘以二作为被除数（实），两股之和作为除数（法），相除得到圆径。}

\shim{这是用两股求内切圆。用两股的数值按比例推算得到圆径。}

\end{problemmodern}



\switchcolumn*



\begin{problem}[\ 两股求容圆七]

\wen{甲出西门直行一百八十步，乙出北门直行二百四十步，相望见之。问城径。}

\da{城径二百四十步。}

\shu{两股相乘倍之为实，并两股为法，除之得圆径。}

\shi{此两股求容圆也。以两股之数，约之得圆径。}

\end{problem}



\switchcolumn



\begin{problemmodern}[\ 两股求内切圆 第七题]

\wenm{甲从西门出发直走一百八十步，乙从北门出发直走二百四十步，两人互相望见。问城池的直径。}

\dam{城池直径为二百四十步。}

\shum{两股相乘，再乘以二作为被除数（实），两股之和作为除数（法），相除得到圆径。}

\shim{这是用两股求内切圆。用两股的数值按比例推算得到圆径。}

\end{problemmodern}



\switchcolumn*



\end{paracol}

\clearpage

\colheaders

\begin{paracol}{2}



%% ========================================================================

%% Chapter 2: Liang xian qiu rong yuan

%% ========================================================================

\bilingualsubsection{两弦求容圆}{Two-Hypotenuse Inscribed Circle Method}



\switchcolumn*



\noindent

\textbf{总论：} 两弦求容圆者，以两弦之数，用减从翻法开平方，求圆城之径也。



\switchcolumn



\noindent

\textbf{【总论】} 两弦求内切圆的方法是：用两弦的数值，采用减从翻法开平方，求圆形城池的直径。



\switchcolumn*



\begin{problem}[\ 两弦求容圆一]

\wen{甲出北门东行，乙出东门南行，相望见之。甲斜行五百步，乙斜行三百步。问城径。}

\da{城径二百四十步。}

\shu{两弦相乘为实，两弦和为法，除之得圆径。}

\shi{此两弦求容圆也。以两弦之数，约之得圆径。}

\end{problem}



\switchcolumn



\begin{problemmodern}[\ 两弦求内切圆 第一题]

\wenm{甲从北门出发向东走，乙从东门出发向南走，两人互相望见。甲斜着走五百步，乙斜着走三百步。问城池的直径。}

\dam{城池直径为二百四十步。}

\shum{两弦相乘作为被除数（实），两弦之和作为除数（法），相除得到圆径。}

\shim{这是用两弦求内切圆。用两弦的数值按比例推算得到圆径。}

\end{problemmodern}



\switchcolumn*



\begin{problem}[\ 两弦求容圆二]

\wen{甲出南门直行，乙出西门直行，相望见之。甲斜行四百步，乙斜行二百五十步。问城径。}

\da{城径二百四十步。}

\shu{两弦相乘为实，两弦和为法，除之得圆径。}

\shi{此两弦求容圆也。以两弦之数，约之得圆径。}

\end{problem}



\switchcolumn



\begin{problemmodern}[\ 两弦求内切圆 第二题]

\wenm{甲从南门出发直走，乙从西门出发直走，两人互相望见。甲斜着走四百步，乙斜着走二百五十步。问城池的直径。}

\dam{城池直径为二百四十步。}

\shum{两弦相乘作为被除数（实），两弦之和作为除数（法），相除得到圆径。}

\shim{这是用两弦求内切圆。用两弦的数值按比例推算得到圆径。}

\end{problemmodern}



\switchcolumn*



\begin{problem}[\ 两弦求容圆三]

\wen{甲出东门直行，乙出北门直行，相望见之。甲斜行三百六十步，乙斜行二百步。问城径。}

\da{城径二百四十步。}

\shu{两弦相乘为实，两弦和为法，除之得圆径。}

\shi{此两弦求容圆也。以两弦之数，约之得圆径。}

\end{problem}



\switchcolumn



\begin{problemmodern}[\ 两弦求内切圆 第三题]

\wenm{甲从东门出发直走，乙从北门出发直走，两人互相望见。甲斜着走三百六十步，乙斜着走二百步。问城池的直径。}

\dam{城池直径为二百四十步。}

\shum{两弦相乘作为被除数（实），两弦之和作为除数（法），相除得到圆径。}

\shim{这是用两弦求内切圆。用两弦的数值按比例推算得到圆径。}

\end{problemmodern}



\switchcolumn*



\end{paracol}

\clearpage

\colheaders

\begin{paracol}{2}



%% ========================================================================

%% Chapter 3: Chang duan er gou qiu chengyuanjing

%% ========================================================================

\bilingualsubsection{长短二勾求城径与通股小差股同法}{Long-and-Short Legs City Diameter Method}



\switchcolumn*



\noindent

\textbf{总论：} 长短二勾求城径者，以长勾、短勾之数，用通股小差股之法，求圆城之径也。



\switchcolumn



\noindent

\textbf{【总论】} 用长勾和短勾求城池直径的方法是：用长勾和短勾的数值，采用通股小差股的方法，求圆形城池的直径。



\switchcolumn*



\begin{problem}[\ 长短二勾求城径一]

\wen{甲出北门东行二百步，乙出东门南行一百五十步，相望见之。问城径。}

\da{城径二百四十步。}

\shu{二行相乘倍之为实，并二行为法，除之得城径。}

\shi{此长短二勾求城径也。以长短二勾之数，约之得圆径。}

\end{problem}



\switchcolumn



\begin{problemmodern}[\ 长勾短勾求城直径 第一题]

\wenm{甲从北门出发向东走二百步，乙从东门出发向南走一百五十步，两人互相望见。问城池的直径。}

\dam{城池直径为二百四十步。}

\shum{两数相乘，再乘以二作为被除数（实），两数之和作为除数（法），相除得到城径。}

\shim{这是用长勾和短勾求城直径。用长勾和短勾的数值按比例推算得到圆径。}

\end{problemmodern}



\switchcolumn*



\begin{problem}[\ 长短二勾求城径二]

\wen{甲出南门直行一百八十步，乙出西门直行一百二十步，相望见之。问城径。}

\da{城径二百四十步。}

\shu{二行相乘倍之为实，并二行为法，除之得城径。}

\shi{此长短二勾求城径也。以长短二勾之数，约之得圆径。}

\end{problem}



\switchcolumn



\begin{problemmodern}[\ 长勾短勾求城直径 第二题]

\wenm{甲从南门出发直走一百八十步，乙从西门出发直走一百二十步，两人互相望见。问城池的直径。}

\dam{城池直径为二百四十步。}

\shum{两数相乘，再乘以二作为被除数（实），两数之和作为除数（法），相除得到城径。}

\shim{这是用长勾和短勾求城直径。用长勾和短勾的数值按比例推算得到圆径。}

\end{problemmodern}



\switchcolumn*



\end{paracol}

\clearpage

\colheaders

\begin{paracol}{2}



%% ========================================================================

%% Chapter 4: Jian cong fanfa kaipingfang

%% ========================================================================

\bilingualsubsection{减从翻法开平方}{Subtractive-Associate Inverted Square-Root Method}



\switchcolumn*



\noindent

\textbf{总论：} 减从翻法开平方者，以实较大于法，用减从之法，翻转开平方，得圆城之径也。



\switchcolumn



\noindent

\textbf{【总论】} 减从翻法开平方的方法是：当被除数（实）大于除数（法）时，采用减从的方法，翻转开平方，得到圆形城池的直径。所谓"减从翻法"，是中国古代开方术中的一种变体方法。



\switchcolumn*



\begin{problem}[\ 减从翻法开平方一]

\wen{余实五百五十万，倍隅法得二十五万为廉法。约次商得二十，置一于左次为上法，置一乘隅算得二万五千，并廉法共二十七万五千为下法，与上法相乘除实尽。后如此类者仿此。}

\da{得半径一百二十步。}

\shu{以减从翻法开平方，置一于左上为上法，倍隅法为廉法，次商自之为隅法，并廉隅共为下法，与上法相乘除实。}

\shi{此减从翻法开平方也。以减从翻法开平方，得圆径。}

\end{problem}



\switchcolumn



\begin{problemmodern}[\ 减从翻法开平方 第一题]

\wenm{剩余被除数（实）为5500000，倍隅法得到250000作为廉法。约次商得20，置1于左次作为上法，置1乘隅算得25000，合并廉法共275000为下法，与上法相乘后除实尽。后面与此类似的问题都仿照此法。}

\dam{得到半径为一百二十步。}

\shum{用减从翻法开平方：在左上方置1作为上法，倍隅法作为廉法，次商自乘作为隅法，合并廉法和隅法共为下法，与上法相乘后除实。}

\shim{这是减从翻法开平方。用减从翻法开平方，得到圆径。}

\end{problemmodern}



\switchcolumn*



\begin{problem}[\ 减从翻法开平方二]

\wen{二十与上次法相乘，得四千四百为益实。添入余积，共五千六百为实。置一并廉法，从方共二百八十为下法，与上次法相乘除实尽。}

\da{得半径一百二十步。}

\shu{以减从翻法开平方，从方添入益实，并廉法为下法，与上法相乘除实。}

\shi{此减从翻法开平方也。以减从翻法开平方，得圆径。}

\end{problem}



\switchcolumn



\begin{problemmodern}[\ 减从翻法开平方 第二题]

\wenm{20与上次法相乘，得到4400作为益实。添入余积，共5600为实。置1合并廉法，从方共280为下法，与上次法相乘除实尽。}

\dam{得到半径为一百二十步。}

\shum{用减从翻法开平方，从方添入益实，合并廉法为下法，与上法相乘除实。}

\shim{这是减从翻法开平方。用减从翻法开平方，得到圆径。}

\end{problemmodern}



\switchcolumn*



\end{paracol}

\clearpage

\colheaders

\begin{paracol}{2}



%% ========================================================================

%% Chapter 5: Gou gu xiang cheng beizhi fa

%% ========================================================================

\bilingualsubsection{勾股相乘倍之为实开平方}{Leg-Product Doubled Square-Root Method}



\switchcolumn*



\noindent

\textbf{总论：} 勾股相乘倍之为实者，以勾股相乘之数倍之为实，开平方得弦，再以弦求容圆也。



\switchcolumn



\noindent

\textbf{【总论】} 勾股相乘倍之为实的方法是：将勾和股相乘，再乘以二作为被除数（实），开平方得到弦（斜边），再用弦求内切圆。



\switchcolumn*



\begin{problem}[\ 勾股相乘倍之开平方一]

\wen{甲出南门东行一百二十步，乙出东门南行八十步，相望见之。问城径。}

\da{城径二百四十步。}

\shu{勾股相乘倍之为实，开平方得弦，并勾股为弦和，弦和较即容圆径。}

\shi{此勾股相乘倍之为实开平方也。以勾股相乘倍之为实，开平方得弦，再以弦求容圆。}

\end{problem}



\switchcolumn



\begin{problemmodern}[\ 勾股相乘倍之开平方 第一题]

\wenm{甲从南门出发向东走一百二十步，乙从东门出发向南走八十步，两人互相望见。问城池的直径。}

\dam{城池直径为二百四十步。}

\shum{勾和股相乘，再乘以二作为被除数（实），开平方得到弦，勾与股之和为弦和，弦和与弦之差就是内切圆直径。}

\shim{这是勾股相乘倍之为实开平方的方法。将勾和股相乘，再乘以二作为被除数，开平方得到弦，再用弦求内切圆。}

\end{problemmodern}



\switchcolumn*



\begin{problem}[\ 勾股相乘倍之开平方二]

\wen{甲出北门东行一百五十步，乙出东门北行一百步，相望见之。问城径。}

\da{城径二百四十步。}

\shu{勾股相乘倍之为实，开平方得弦，并勾股为弦和，弦和较即容圆径。}

\shi{此勾股相乘倍之为实开平方也。以勾股相乘倍之为实，开平方得弦，再以弦求容圆。}

\end{problem}



\switchcolumn



\begin{problemmodern}[\ 勾股相乘倍之开平方 第二题]

\wenm{甲从北门出发向东走一百五十步，乙从东门出发向北走一百步，两人互相望见。问城池的直径。}

\dam{城池直径为二百四十步。}

\shum{勾和股相乘，再乘以二作为被除数（实），开平方得到弦，勾与股之和为弦和，弦和与弦之差就是内切圆直径。}

\shim{这是勾股相乘倍之为实开平方的方法。将勾和股相乘，再乘以二作为被除数，开平方得到弦，再用弦求内切圆。}

\end{problemmodern}



\switchcolumn*



\end{paracol}

\clearpage

\colheaders

\begin{paracol}{2}



%% ========================================================================

%% Chapter 6: Xian he jiao qiu yuanjing

%% ========================================================================

\bilingualsubsection{弦和较求圆径}{Chord-Sum-Difference Circle-Diameter Method}



\switchcolumn*



\noindent

\textbf{总论：} 弦和较求圆径者，以弦与勾股和之差，即容圆径也。



\switchcolumn



\noindent

\textbf{【总论】} 用弦和较求圆径的方法是：弦与勾股和之差，就是内切圆的直径。



\switchcolumn*



\begin{problem}[\ 弦和较求圆径一]

\wen{甲出南门直行，乙出东门直行，相望见之。问以弦和较求圆径。}

\da{城径二百四十步。}

\shu{弦和较即容圆径。弦者勾股弦也，和者勾股和也。弦减勾股和之较即圆径。}

\shi{此弦和较求圆径也。弦和较者，通弦减通勾股和之较也，即容圆径。}

\end{problem}



\switchcolumn



\begin{problemmodern}[\ 弦和较求圆直径 第一题]

\wenm{甲从南门出发直走，乙从东门出发直走，两人互相望见。问如何用弦和较求圆直径。}

\dam{城池直径为二百四十步。}

\shum{弦和较就是内切圆直径。弦就是勾股弦（斜边），和就是勾股之和。弦减去勾股和的差就是圆直径。}

\shim{这是用弦和较求圆直径。弦和较就是通弦减去通勾股和的差，即内切圆直径。}

\end{problemmodern}



\switchcolumn*



\begin{problem}[\ 弦和较求圆径二]

\wen{甲出北门直行，乙出西门直行，相望见之。问以弦和较求圆径。}

\da{城径二百四十步。}

\shu{弦和较即容圆径。弦者勾股弦也，和者勾股和也。弦减勾股和之较即圆径。}

\shi{此弦和较求圆径也。以弦和较之数，即容圆径。}

\end{problem}



\switchcolumn



\begin{problemmodern}[\ 弦和较求圆直径 第二题]

\wenm{甲从北门出发直走，乙从西门出发直走，两人互相望见。问如何用弦和较求圆直径。}

\dam{城池直径为二百四十步。}

\shum{弦和较就是内切圆直径。弦就是勾股弦（斜边），和就是勾股之和。弦减去勾股和的差就是圆直径。}

\shim{这是用弦和较求圆直径。弦和较的数值就是内切圆直径。}

\end{problemmodern}



\switchcolumn*



\end{paracol}

\clearpage

\colheaders

\begin{paracol}{2}



%% ========================================================================

%% Reference Table for Volume 2

%% ========================================================================

\subsection*{卷二名数表}

\addcontentsline{toc}{subsection}{卷二名数表}



\switchcolumn

\subsection*{卷二名数对照表（现代注释）}

\addcontentsline{toc}{subsection}{卷二名数对照表}



\switchcolumn*



\noindent

\textbf{卷二所用名数对照表：}



\switchcolumn

\noindent

\textbf{【卷二所用术语数值对照表：】}



\switchcolumn*



\begin{center}

\begin{tabular}{|l|l|l|}

\hline

\textbf{名称} & \textbf{数值} & \textbf{说明} \\

\hline

通股 & 四百八十 & 通股之数 \\

\hline

边股 & 二百四十 & 边股之数 \\

\hline

底股 & 三百六十 & 底股之数 \\

\hline

黄广股 & 四百五十 & 黄广股之数 \\

\hline

黄长股 & 一百二十 & 黄长股之数 \\

\hline

高股 & 二百 & 高股之数 \\

\hline

平股 & 八十 & 平股之数 \\

\hline

大差股 & 一百六十 & 大差股之数 \\

\hline

小差股 & 六十 & 小差股之数 \\

\hline

皇极股 & 二百四十 & 皇极股之数 \\

\hline

太虚股 & 三十 & 太虚股之数 \\

\hline

明股 & 八十 & 明股之数 \\

\hline

重股 & 十六 & 重股之数 \\

\hline

\end{tabular}

\end{center}



\switchcolumn



\begin{center}

\begin{tabular}{|l|l|p{4.5cm}|}

\hline

\textbf{术语} & \textbf{数值} & \textbf{现代说明} \\

\hline

通股 & 480 & 通股（大直角三角形的高边） \\

\hline

边股 & 240 & 边股（城池边缘的高边） \\

\hline

底股 & 360 & 底股（底部的高边） \\

\hline

黄广股 & 450 & 黄广股（黄色宽大的高边） \\

\hline

黄长股 & 120 & 黄长股（黄色狭长的高边） \\

\hline

高股 & 200 & 高股（较高的高边） \\

\hline

平股 & 80 & 平股（水平的高边） \\

\hline

大差股 & 160 & 大差股（大差的高边） \\

\hline

小差股 & 60 & 小差股（小差的高边） \\

\hline

皇极股 & 240 & 皇极股（皇极高边） \\

\hline

太虚股 & 30 & 太虚股（太空虚的高边） \\

\hline

明股 & 80 & 明股（明亮的高边） \\

\hline

重股 & 16 & 重股（双重的高边） \\

\hline

\end{tabular}

\end{center}



\switchcolumn*



\end{paracol}

\clearpage

\colheaders

\begin{paracol}{2}



%% ========================================================================

%% VOLUME 3 (卷三)

%% ========================================================================

\section{卷三}

\markboth{卷三}{卷三}



\switchcolumn

\section*{卷三（现代译文）}

\addcontentsline{toc}{section}{卷三}

\markboth{卷三}{卷三}



\switchcolumn*



卷三专论\textbf{带从开平方法}及\textbf{三乘方}之法，凡十余条。带从开平方法者，以实较大于法，用带从之法，开平方得圆城之径也。又论减从翻法、益积、减积诸法，皆开方之变术也。



\switchcolumn



卷三专门讨论\textbf{带从开平方法}以及\textbf{三乘方}（三次方程）的方法，共十余条。带从开平方法是指：当被除数（实）大于除数（法）时，采用带从的方法，开平方得到圆形城池的直径。另外还讨论了减从翻法、益积（增加积数）、减积（减少积数）等各种方法，都是开方术的变体方法。



\switchcolumn*



\end{paracol}

\clearpage

\colheaders

\begin{paracol}{2}



%% ========================================================================

%% Chapter 1: Dai cong kaipingfang fa

%% ========================================================================

\bilingualsubsection{带从开平方法}{Accompanying-Associate Square-Root Method}



\switchcolumn*



\noindent

\textbf{总论：} 带从开平方法者，以带从之实，用开平方之法，得圆城之径也。所谓带从者，实内带有从方之数也。



\switchcolumn



\noindent

\textbf{【总论】} 带从开平方法是指：用带有从项的被除数（实），采用开平方的方法，得到圆形城池的直径。所谓"带从"，就是被除数（实）中带有从方（线性项）的数值。



\switchcolumn*



\begin{problem}[\ 带从开平方法一]

\wen{乙出南门直行一百三十五步，甲出北门东行二百步见之。问以带从开平方求城径。}

\da{城径二百四十步。}

\shu{为实以南行三百六十为从方，作带从开平方法除之得全径。}

\shi{此带从开平方法也。以从方之数，开平方得圆径。}

\end{problem}



\switchcolumn



\begin{problemmodern}[\ 带从开平方 第一题]

\wenm{乙从南门出发直走一百三十五步，甲从北门出发向东走二百步后看见乙。问如何用带从开平方求城池直径。}

\dam{城池直径为二百四十步。}

\shum{以南行三百六十作为从方，采用带从开平方法除之，得到全径。}

\shim{这是带从开平方法。用从方的数值，开平方得到圆径。}

\end{problemmodern}



\switchcolumn*



\begin{problem}[\ 带从开平方法二]

\wen{甲出北门东行一百二十步，乙出东门南行一百步，相望见之。问以带从开平方求城径。}

\da{城径二百四十步。}

\shu{为实以两行之和为从方，作带从开平方法除之得全径。}

\shi{此带从开平方法也。以从方之数，开平方得圆径。}

\end{problem}



\switchcolumn



\begin{problemmodern}[\ 带从开平方 第二题]

\wenm{甲从北门出发向东走一百二十步，乙从东门出发向南走一百步，两人互相望见。问如何用带从开平方求城池直径。}

\dam{城池直径为二百四十步。}

\shum{以两数之和作为从方，采用带从开平方法除之，得到全径。}

\shim{这是带从开平方法。用从方的数值，开平方得到圆径。}

\end{problemmodern}



\switchcolumn*



\begin{problem}[\ 带从开平方法三]

\wen{甲出西门南行二百步，乙出南门西行一百五十步，相望见之。问以带从开平方求城径。}

\da{城径二百四十步。}

\shu{为实以两行之和为从方，作带从开平方法除之得全径。}

\shi{此带从开平方法也。以从方之数，开平方得圆径。}

\end{problem}



\switchcolumn



\begin{problemmodern}[\ 带从开平方 第三题]

\wenm{甲从西门出发向南走二百步，乙从南门出发向西走一百五十步，两人互相望见。问如何用带从开平方求城池直径。}

\dam{城池直径为二百四十步。}

\shum{以两数之和作为从方，采用带从开平方法除之，得到全径。}

\shim{这是带从开平方法。用从方的数值，开平方得到圆径。}

\end{problemmodern}



\switchcolumn*



\begin{problem}[\ 带从开平方法四]

\wen{甲出南门直行二百五十步，乙出东门直行二百步，相望见之。问以带从开平方求城径。}

\da{城径二百四十步。}

\shu{为实以两行之和为从方，作带从开平方法除之得全径。}

\shi{此带从开平方法也。以从方之数，开平方得圆径。}

\end{problem}



\switchcolumn



\begin{problemmodern}[\ 带从开平方 第四题]

\wenm{甲从南门出发直走二百五十步，乙从东门出发直走二百步，两人互相望见。问如何用带从开平方求城池直径。}

\dam{城池直径为二百四十步。}

\shum{以两数之和作为从方，采用带从开平方法除之，得到全径。}

\shim{这是带从开平方法。用从方的数值，开平方得到圆径。}

\end{problemmodern}



\switchcolumn*



\end{paracol}

\clearpage

\colheaders

\begin{paracol}{2}



%% ========================================================================

%% Chapter 2: Jian cong fanfa kaipingfang

%% ========================================================================

\bilingualsubsection{减从翻法开平方}{Subtractive-Associate Inverted Square-Root Method}



\switchcolumn*



\noindent

\textbf{总论：} 减从翻法开平方者，以实内减去从方之数，翻转开平方，得圆城之径也。凡言减从翻法开平方者俱仿此。



\switchcolumn



\noindent

\textbf{【总论】} 减从翻法开平方的方法是：在被除数（实）中减去从方的数值，再翻转开平方，得到圆形城池的直径。凡提到"减从翻法开平方"的问题都仿照此方法。



\switchcolumn*



\begin{problem}[\ 减从翻法开平方一]

\wen{十余一千五百二十为负积。又以初商一百反减余，从四十四余五十六为负。从次商二十并负，从共七十六为下法，亦通后凡言减从翻法开平方者俱仿此。}

\da{得半径一百二十步。}

\shu{以减从翻法开平方曰：初商一百，置一于左上为上法。置一乘隅法，得四为隅法，作负隅开平方。}

\shi{此减从翻法开平方也。以减从翻法开平方，得圆径。}

\end{problem}



\switchcolumn



\begin{problemmodern}[\ 减从翻法开平方 第一题]

\wenm{十余一千五百二十为负积。又以初商一百反减余，从四十四余五十六为负。从次商二十并负，从共七十六为下法。这也适用于后面所有提到减从翻法开平方的问题。}

\dam{得到半径为一百二十步。}

\shum{用减从翻法开平方：初商为100，在左上方置1作为上法。置1乘隅法，得到4为隅法，作为负隅开平方。}

\shim{这是减从翻法开平方。用减从翻法开平方，得到圆径。}

\end{problemmodern}



\switchcolumn*



\begin{problem}[\ 减从翻法开平方二]

\wen{余实五百五十万，倍隅法得二十五万为廉法。约次商得二十，置一于左次为上法，置一乘隅得二万五千，并廉法共二十七万五千为下法，与上法相乘除实。尽后如此类者仿此。}

\da{得半径一百二十步。}

\shu{以减从翻法开平方，置一于左上为上法，倍隅法为廉法，次商自之为隅法，并廉隅共为下法，与上法相乘除实。}

\shi{此减从翻法开平方也。以减从翻法开平方，得圆径。}

\end{problem}



\switchcolumn



\begin{problemmodern}[\ 减从翻法开平方 第二题]

\wenm{剩余被除数5500000，倍隅法得到250000作为廉法。约次商得20，置1于左次作为上法，置1乘隅得25000，合并廉法共275000为下法，与上法相乘后除实尽。后面与此类似的问题都仿照此法。}

\dam{得到半径为一百二十步。}

\shum{用减从翻法开平方：在左上方置1作为上法，倍隅法作为廉法，次商自乘作为隅法，合并廉法和隅法共为下法，与上法相乘除实。}

\shim{这是减从翻法开平方。用减从翻法开平方，得到圆径。}

\end{problemmodern}



\switchcolumn*



\begin{problem}[\ 减从翻法开平方三]

\wen{二十与上次法相乘，得四千四百为益实。添入余积共五千六百为实。置一并廉法从方共二百八十为下法，与上次法相乘除实尽。}

\da{得半径一百二十步。}

\shu{以减从翻法开平方，从方添入益实，并廉法为下法，与上法相乘除实。}

\shi{此减从翻法开平方也。以减从翻法开平方，得圆径。}

\end{problem}



\switchcolumn



\begin{problemmodern}[\ 减从翻法开平方 第三题]

\wenm{20与上次法相乘，得到4400为益实。添入余积共5600为实。置1合并廉法从方共280为下法，与上次法相乘除实尽。}

\dam{得到半径为一百二十步。}

\shum{用减从翻法开平方，从方添入益实，合并廉法为下法，与上法相乘除实。}

\shim{这是减从翻法开平方。用减从翻法开平方，得到圆径。}

\end{problemmodern}



\switchcolumn*



\end{paracol}

\clearpage

\colheaders

\begin{paracol}{2}



%% ========================================================================

%% Chapter 3: San cheng fang fa

%% ========================================================================

\bilingualsubsection{三乘方法}{Three-Fold Multiplication (Cubic) Method}



\switchcolumn*



\noindent

\textbf{总论：} 三乘方法者，以三次方程之法，求圆城之径也。所谓三乘者，立方以上之开方也。



\switchcolumn



\noindent

\textbf{【总论】} 三乘方法是指：用三次方程的方法，求圆形城池的直径。所谓"三乘"，就是开立方以上的开方运算。



\switchcolumn*



\begin{problem}[\ 三乘方法一]

\wen{置一自之以乘，从二廉，置一自乘再乘为隅法。并从方廉隅共七千九百八十一万三千四百，并从方廉隅共七千九百八十一万三千三百四十为下法，与上法相乘除实。}

\da{得半径一百二十步。}

\shu{以三乘方法开之，置一自之为上法，以次商乘廉法为从方，再以次商自之为隅法，并从方廉隅为下法，与上法相乘除实。}

\shi{此三乘方法也。以三乘方法开之，得圆径。}

\end{problem}



\switchcolumn



\begin{problemmodern}[\ 三乘方 第一题]

\wenm{置1自乘再乘以从二廉，置1自乘再乘作为隅法。合并从方、廉、隅共79813400，合并从方、廉、隅共79813340为下法，与上法相乘除实。}

\dam{得到半径为一百二十步。}

\shum{用三乘方方法开之：置1自乘作为上法，以次商乘廉法为从方，再以次商自乘为隅法，合并从方、廉、隅为下法，与上法相乘除实。}

\shim{这是三乘方方法。用三乘方方法开之，得到圆径。}

\end{problemmodern}



\switchcolumn*



\begin{problem}[\ 三乘方法二]

\wen{置一自之以乘从二廉，置一自乘再乘为隅法。并从方廉隅共七千九百八十一万三千三百四十为下法，与上法相乘除实尽。}

\da{得半径一百二十步。}

\shu{以三乘方法开之，置一自之为上法，以次商乘廉法为从方，再以次商自之为隅法，并从方廉隅为下法，与上法相乘除实。}

\shi{此三乘方法也。以三乘方法开之，得圆径。}

\end{problem}



\switchcolumn



\begin{problemmodern}[\ 三乘方 第二题]

\wenm{置1自乘再乘以从二廉，置1自乘再乘作为隅法。合并从方、廉、隅共79813340为下法，与上法相乘除实尽。}

\dam{得到半径为一百二十步。}

\shum{用三乘方方法开之：置1自乘作为上法，以次商乘廉法为从方，再以次商自乘为隅法，合并从方、廉、隅为下法，与上法相乘除实。}

\shim{这是三乘方方法。用三乘方方法开之，得到圆径。}

\end{problemmodern}



\switchcolumn*



\end{paracol}

\clearpage

\colheaders

\begin{paracol}{2}



%% ========================================================================

%% Chapter 4: Yi ji jian ji fa

%% ========================================================================

\bilingualsubsection{益积减积法}{Additive-and-Subtractive Accumulation Method}



\switchcolumn*



\noindent

\textbf{总论：} 益积减积法者，以益积或减积之法，开方得圆城之径也。益积者添入积数，减积者减去积数。



\switchcolumn



\noindent

\textbf{【总论】} 益积减积法是指：用益积（增加积数）或减积（减少积数）的方法，开方得到圆形城池的直径。益积就是添入积数，减积就是减去积数。



\switchcolumn*



\begin{problem}[\ 益积减积法一]

\wen{二十与上次法相乘，得四千四百为益实。添入余积，共五千六百为实。置一并廉法，从方共二百八十为下法，与上次法相乘除实尽。}

\da{得半径一百二十步。}

\shu{以益积法添入益实，并廉法为下法，与上法相乘除实。}

\shi{此益积法也。以益积之数，开平方得圆径。}

\end{problem}



\switchcolumn



\begin{problemmodern}[\ 益积减积法 第一题]

\wenm{20与上次法相乘，得到4400为益实。添入余积，共5600为实。置1合并廉法，从方共280为下法，与上次法相乘除实尽。}

\dam{得到半径为一百二十步。}

\shum{用益积法添入益实，合并廉法为下法，与上法相乘除实。}

\shim{这是益积法。用益积的数值，开平方得到圆径。}

\end{problemmodern}



\switchcolumn*



\begin{problem}[\ 益积减积法二]

\wen{余实五百五十万，倍隅法得二十五万为廉法。约次商得二十，置一于左次为上法，置一乘隅得二万五千，并廉法共二十七万五千为下法，与上法相乘除实。}

\da{得半径一百二十步。}

\shu{以减积法减去积数，并廉法为下法，与上法相乘除实。}

\shi{此减积法也。以减积之数，开平方得圆径。}

\end{problem}



\switchcolumn



\begin{problemmodern}[\ 益积减积法 第二题]

\wenm{剩余被除数5500000，倍隅法得到250000作为廉法。约次商得20，置1于左次作为上法，置1乘隅得25000，合并廉法共275000为下法，与上法相乘除实。}

\dam{得到半径为一百二十步。}

\shum{用减积法减去积数，合并廉法为下法，与上法相乘除实。}

\shim{这是减积法。用减积的数值，开平方得到圆径。}

\end{problemmodern}



\switchcolumn*



\end{paracol}

\clearpage

\colheaders

\begin{paracol}{2}



%% ========================================================================

%% Chapter 5: Cong fang cong lian yu fa

%% ========================================================================

\bilingualsubsection{从方从廉隅法}{Associate Edge, Corner and Limit Method}



\switchcolumn*



\noindent

\textbf{总论：} 从方从廉隅法者，以从方、从廉、隅法三者之关系，开方得圆城之径也。



\switchcolumn



\noindent

\textbf{【总论】} 从方从廉隅法是指：利用从方、从廉、隅法三者之间的关系，开方得到圆形城池的直径。



\switchcolumn*



\begin{problem}[\ 从方从廉隅法一]

\wen{置一自之以乘从二廉，置一自乘再乘为隅法。并从方廉隅共七千九百八十一万三千三百四十为下法，与上法相乘除实尽。}

\da{得半径一百二十步。}

\shu{以从方从廉隅法开之，置一自之为上法，以次商乘廉法为从方，再以次商自之为隅法，并从方廉隅为下法，与上法相乘除实。}

\shi{此从方从廉隅法也。以从方从廉隅法开之，得圆径。}

\end{problem}



\switchcolumn



\begin{problemmodern}[\ 从方从廉隅法 第一题]

\wenm{置1自乘再乘以从二廉，置1自乘再乘作为隅法。合并从方、廉、隅共79813340为下法，与上法相乘除实尽。}

\dam{得到半径为一百二十步。}

\shum{用从方从廉隅法开之：置1自乘作为上法，以次商乘廉法为从方，再以次商自乘为隅法，合并从方、廉、隅为下法，与上法相乘除实。}

\shim{这是从方从廉隅法。用从方从廉隅法开之，得到圆径。}

\end{problemmodern}



\switchcolumn*



\begin{problem}[\ 从方从廉隅法二]

\wen{置一自之为上法，置一乘隅得二万五千，并廉法共二十七万五千为下法，与上法相乘除实。}

\da{得半径一百二十步。}

\shu{以从方从廉隅法开之，置一自之为上法，以次商乘廉法为从方，再以次商自之为隅法，并从方廉隅为下法，与上法相乘除实。}

\shi{此从方从廉隅法也。以从方从廉隅法开之，得圆径。}

\end{problem}



\switchcolumn



\begin{problemmodern}[\ 从方从廉隅法 第二题]

\wenm{置1自乘作为上法，置1乘隅得25000，合并廉法共275000为下法，与上法相乘除实。}

\dam{得到半径为一百二十步。}

\shum{用从方从廉隅法开之：置1自乘作为上法，以次商乘廉法为从方，再以次商自乘为隅法，合并从方、廉、隅为下法，与上法相乘除实。}

\shim{这是从方从廉隅法。用从方从廉隅法开之，得到圆径。}

\end{problemmodern}



\switchcolumn*



\end{paracol}

\clearpage

\colheaders

\begin{paracol}{2}



%% ========================================================================

%% Chapter 6: Bu ji di gou bian gu tiao

%% ========================================================================

\bilingualsubsection{不及底勾边股条}{Base-Leg and Edge-Leg Deficiency Method}



\switchcolumn*



\noindent

\textbf{总论：} 不及底勾边股条者，以底勾、边股不及之数，用测望之法，求圆城之径也。



\switchcolumn



\noindent

\textbf{【总论】} 不及底勾边股条是指：用底勾和边股不足（不及）的数值，采用测望的方法，求圆形城池的直径。



\switchcolumn*



\begin{problem}[\ 不及底勾边股条一]

\wen{出西门南行二百二十五步，有塔出北门东行六十四步望塔正居城之半。问城径。}

\da{城径二百四十步。}

\shu{此以不及底勾与不及边股测望也。南行二百二十五步与高股同，即半径为勾之股。东行六十四步与平勾同，即半径为股之勾也。当以平勾高股立法。}

\shi{此以不及底勾边股条求城径也。以底勾边股之不及数，测望得圆径。}

\end{problem}



\switchcolumn



\begin{problemmodern}[\ 不及底勾边股条 第一题]

\wenm{从西门出发向南走二百二十五步，有座塔；从北门出发向东走六十四步望见塔，正好位于城池的一半处。问城池的直径。}

\dam{城池直径为二百四十步。}

\shum{这是用不及底勾与不及边股进行测望。向南走225步与高股相同，即半径为勾之股。向东走64步与平勾相同，即半径为股之勾。应当用平勾和高股来立法。}

\shim{这是用不及底勾边股条求城直径。用底勾和边股的不及数，通过测望得到圆径。}

\end{problemmodern}



\switchcolumn*



\begin{problem}[\ 不及底勾边股条二]

\wen{乙从城外西南坤隅南行三百六十步，甲出北门东行二百步见之。问城径。}

\da{城径二百四十步。}

\shu{二行相乘即半径幂。}

\shi{此以底勾大差股立法测望也。乙从坤隅南行大差股也，甲东行底勾也。底勾为城北半勾，大差股为城西南虚股。}

\end{problem}



\switchcolumn



\begin{problemmodern}[\ 不及底勾边股条 第二题]

\wenm{乙从城外西南角（坤隅）向南走三百六十步，甲从北门出发向东走二百步后看见乙。问城池的直径。}

\dam{城池直径为二百四十步。}

\shum{两数相乘就是半径的平方。}

\shim{这是用底勾大差股立法进行测望。乙从坤隅向南走是大差股，甲向东走是底勾。底勾是城北的半勾，大差股是城西南的虚股。}

\end{problemmodern}



\switchcolumn*



\begin{problem}[\ 不及底勾边股条三]

\wen{甲出北门东行底勾也。底勾为城北半勾，明股为城南余股。}

\da{城径二百四十步。}

\shu{以底勾明股之数，约之得圆径。}

\shi{此底勾明股立法也。以底勾明股之数，测望得圆径。}

\end{problem}



\switchcolumn



\begin{problemmodern}[\ 不及底勾边股条 第三题]

\wenm{甲从北门出发向东走就是底勾。底勾是城北的半勾，明股是城南的余股。}

\dam{城池直径为二百四十步。}

\shum{用底勾和明股的数值，按比例推算得到圆径。}

\shim{这是底勾明股立法。用底勾和明股的数值，通过测望得到圆径。}

\end{problemmodern}



\switchcolumn*



\end{paracol}

\clearpage

\colheaders

\begin{paracol}{2}



%% ========================================================================

%% Chapter 7: Shesheng buta chengfa

%% ========================================================================

\bilingualsubsection{测望不迭乘法}{Non-Iterative Surveying Multiplication Method}



\switchcolumn*



\noindent

\textbf{总论：} 测望不迭乘法者，以测望之法，不迭次相乘，求圆城之径也。



\switchcolumn



\noindent

\textbf{【总论】} 测望不迭乘法是指：采用测望的方法，不依次迭代相乘，求圆形城池的直径。



\switchcolumn*



\begin{problem}[\ 测望不迭乘法一]

\wen{此法分别从方从廉明白，故重录附之。}

\da{得半径一百二十步。}

\shu{以测望之法，分别从方从廉，不迭次相乘，开平方得圆径。}

\shi{此测望不迭乘法也。以测望之法，不迭次相乘，开平方得圆径。}

\end{problem}



\switchcolumn



\begin{problemmodern}[\ 测望不迭乘法 第一题]

\wenm{这种方法分别从方和从廉都很清楚，所以重新录写附在这里。}

\dam{得到半径为一百二十步。}

\shum{用测望的方法，分别从方和从廉，不依次迭代相乘，开平方得到圆径。}

\shim{这是测望不迭乘法。用测望的方法，不依次迭代相乘，开平方得到圆径。}

\end{problemmodern}



\switchcolumn*



\end{paracol}

\clearpage

\colheaders

\begin{paracol}{2}



%% ========================================================================

%% Reference Table for Volume 3

%% ========================================================================

\subsection*{卷三名数表}

\addcontentsline{toc}{subsection}{卷三名数表}



\switchcolumn

\subsection*{卷三名数对照表（现代注释）}

\addcontentsline{toc}{subsection}{卷三名数对照表}



\switchcolumn*



\noindent

\textbf{卷三所用名数对照表：}



\switchcolumn

\noindent

\textbf{【卷三所用术语数值对照表：】}



\switchcolumn*



\begin{center}

\begin{tabular}{|l|l|l|}

\hline

\textbf{名称} & \textbf{数值} & \textbf{说明} \\

\hline

上法 & 初商自之 & 开方之上位法 \\

\hline

从方 & 次商乘廉法 & 从方之数 \\

\hline

廉法 & 倍隅法 & 廉法之数 \\

\hline

隅法 & 次商自之 & 隅法之数 \\

\hline

下法 & 从方廉隅共 & 下法之数 \\

\hline

实 & 积数 & 实之数 \\

\hline

益实 & 添入之数 & 益实之数 \\

\hline

减实 & 减去之数 & 减实之数 \\

\hline

负积 & 负数之积 & 负积之数 \\

\hline

带从 & 从方带入 & 带从之数 \\

\hline

\end{tabular}

\end{center}



\switchcolumn



\begin{center}

\begin{tabular}{|l|l|p{4.5cm}|}

\hline

\textbf{术语} & \textbf{数值/说明} & \textbf{现代解释} \\

\hline

上法 & 初商自乘 & 开方运算中上位的除数 \\

\hline

从方 & 次商乘廉法 & 带有从属性的线性项系数 \\

\hline

廉法 & 倍隅法 & 廉法就是隅法的两倍 \\

\hline

隅法 & 次商自乘 & 隅法就是次商的平方 \\

\hline

下法 & 从方+廉+隅 & 下位除数的总和 \\

\hline

实 & 积数 & 被除数（常数项） \\

\hline

益实 & 添入的数 & 需要加到实中的数（正数） \\

\hline

减实 & 减去的数 & 需要从实中减去的数（负数） \\

\hline

负积 & 负数的积 & 带有负号的积数 \\

\hline

带从 & 带入从方 & 带有从方项的开方运算 \\

\hline

\end{tabular}

\end{center}



\switchcolumn*



\end{paracol}

\clearpage

\colheaders

\begin{paracol}{2}



%% ========================================================================

%% POSTSCRIPT (跋)

%% ========================================================================

\bilingualsection{跋}{Postscript / 书后跋文}



\switchcolumn*



测圆海镜分类释术十卷，元李冶撰，明顾应祥分类释术。其书以勾股测圆之术，分类编次，以便学者。卷一至卷三论通勾股、边勾股、底股弦、皇极勾股、太阴勾股、明勾股、重勾股诸法求容圆，及两股求容圆、两弦求容圆、带从开平方、减从翻法开平方、三乘方法、益积减积法、从方从廉隅法诸术。



\switchcolumn



《测圆海镜分类释术》共十卷，元代李冶撰写，明代顾应祥分类并解释其方法。这本书将勾股测圆的方法按类别编排，以方便学者学习。卷一至卷三论述了用通勾股、边勾股、底股弦、皇极勾股、太阴勾股、明勾股、重勾股等各种方法求内切圆，以及用两股求内切圆、两弦求内切圆、带从开平方、减从翻法开平方、三乘方方法、益积减积法、从方从廉隅法等各种算法。



\switchcolumn*



其术之大要，以勾股相乘倍之为实，勾股和为法，除之得圆径。又以弦和较即容圆径。其开方之法，有带从、减从翻法、益积、减积诸术，皆天元术之支流也。



\switchcolumn



这些方法的大要，是将勾和股相乘再乘以二作为被除数（实），勾与股之和作为除数（法），相除得到圆的直径。另外，弦和较就是内切圆直径。其开方的方法，有带从、减从翻法、益积、减积等各种技术，都是天元术的支流。



\switchcolumn*



顾应祥序谓：李冶原书每条下细草，俱径立天元一，反复合之，而无下手之术。使后学之士茫然无门路可入。故去其细草，立一算术，又以其所立通勾边股之属各以类分之。此诚后学之津梁也。



\switchcolumn



顾应祥在序言中说：李冶原书每条下面的详细演算，都是直接设天元一（设未知数），反复配合，却没有入手的具体方法。使后来的学者茫然没有门路可入。所以去掉详细的演算过程，设立一套计算方法，又把其中所涉及的通勾、边股等各以类别分开。这确实是后学的桥梁啊。



\switchcolumn*



\bigskip

\begin{flushright}

测圆海镜分类释术卷三终

\end{flushright}



\switchcolumn



\bigskip

\begin{flushright}

《测圆海镜分类释术》卷三（终）

\end{flushright}



\switchcolumn*



\end{paracol}

\clearpage



%% ========================================================================

%% END COLOPHON

%% ========================================================================

\thispagestyle{empty}

\begin{center}

\vspace*{3cm}

{\Large 测圆海镜分类释术}



\vspace{0.5cm}

{\large 卷一至卷三终}



\vspace{0.3cm}

{\normalsize 文言文↔白话文 对照版}



\vspace{2cm}

\shortline



\vspace{1cm}

{\normalsize

\begin{tabular}{rl}

总校官候补中书 & 臣吴绍洁 \\

校对官中官正 & 臣郭长发 \\

誊录监生 & 臣丁湘锦 \\

\end{tabular}

}



\vspace{2cm}

{\small\textcolor{notecolor}{%

此版为现代学者所制文白对照版\\

原文据文渊阁《四库全书荟要》本\\

子部天文算法类\\

白话文译文仅供学习参考

}}



\end{center}



\clearpage



%% ========================================================================

%% MODERN EDITORIAL NOTE

%% ========================================================================

\section*{Modern Editorial Note}

\addcontentsline{toc}{section}{Modern Editorial Note（现代编辑说明）}

\markboth{Modern Editorial Note}{Modern Editorial Note}



\subsection*{Historical Context}

The original \textit{Ceyuan Haijing} was composed by Li Ye (李冶, 1192--1279) in 1248 during the Yuan dynasty. It is one of the masterpieces of medieval Chinese algebra, presenting 170 problems concerning the computation of the diameter of an inscribed circle in a right triangle and its numerous variations. Li Ye's original work employed the \textit{tian yuan shu} (天元术, ``celestial element method'') --- an advanced algebraic technique using rod numerals to set up and solve polynomial equations.



\textbf{历史背景：} 原《测圆海镜》由元代李冶（1192--1279）于1248年所著。它是中世纪中国代数学的杰作之一，提出了170道关于直角三角形内切圆直径计算及其众多变体的问题。李冶原著采用了"天元术"——一种使用算筹数码建立和求解多项式方程的高级代数技术。



\subsection*{Gu Yingxiang's Rearrangement}

Gu Yingxiang (顾应祥, 1483--1565), a Ming dynasty mathematician and Minister of Justice, found Li Ye's original presentation difficult for students to follow because the \textit{tian yuan shu} method had been largely forgotten. He therefore:

\begin{enumerate}

\item Removed the detailed ``fine grass'' (细草) working using \textit{tian yuan shu}

\item Replaced it with explicit step-by-step arithmetic procedures (算术)

\item Reorganized the 170 problems by mathematical technique rather than geometric configuration

\item Added explanatory notes (释术) explaining why each method works

\end{enumerate}



\textbf{顾应祥的重新编排：} 顾应祥（1483--1565），明代数学家和刑部尚书，发现李冶原著的表述方式使学生难以学习，因为天元术在很大程度上已被遗忘。因此他：

\begin{enumerate}

\item 去掉了使用天元术的详细"细草"演算过程

\item 用明确的逐步算术程序（算术）取而代之

\item 按数学方法而非几何构型重新编排了170道问题

\item 添加了释术（方法解释），说明每种方法的原理

\end{enumerate}



\subsection*{Structure of Volumes 1--3}

\begin{itemize}

\item \textbf{Volume 1} focuses on ``containing-circle'' (容圆) methods using various types of right triangles: \textit{tong} (通), \textit{bian} (边), \textit{di} (底), \textit{huangji} (皇极), \textit{taiyin} (太阴), \textit{ming} (明), and \textit{chong} (重) gougu configurations.

\item \textbf{Volume 2} covers methods using two legs (两股), two hypotenuses (两弦), difference of legs, and the ``subtractive-associate inverted method'' (减从翻法) for root extraction.

\item \textbf{Volume 3} presents advanced root-extraction techniques: the ``accompanying-associate'' (带从) method, three-fold multiplication (三乘方), additive and subtractive accumulation (益积减积), and methods involving the edge-leg and base-leg relationships.

\end{itemize}



\textbf{卷一至卷三的结构：}

\begin{itemize}

\item \textbf{卷一} 重点讨论各种直角三角形的"容圆"（内切圆）方法，包括通勾股、边勾股、底勾股、皇极勾股、太阴勾股、明勾股和重勾股等构型。

\item \textbf{卷二} 讨论用两股、两弦、勾股差以及减从翻法等开方方法来求圆径。

\item \textbf{卷三} 介绍高级开方技术：带从法、三乘方、益积减积法，以及涉及边股和底股关系的方法。



\end{itemize}



\subsection*{Mathematical Significance}

The problems all revolve around a single geometric configuration: a circular city (圆城) with given measurements from various observation points, from which one must determine the city's diameter. The invariant answer throughout is \textbf{240 bu} (步), demonstrating the algebraic relationships rather than varying geometric scenarios.



The methods exemplify the sophisticated Chinese tradition of root extraction (开方术) and polynomial algebra that predated European developments by centuries. Gu Yingxiang's classification system reveals the underlying structural unity of Li Ye's 170 problems, organizing them by their mathematical technique rather than their surface geometric appearance.



\textbf{数学意义：}

所有问题都围绕一个几何构型：一座圆形城池（圆城），给出从不同观测点到城池的各段距离，要求确定城池的直径。所有问题的不变答案都是\textbf{240步}，这展示了代数关系而非变化的几何情景。



这些方法体现了中国先进的开方术（开方术）和多项式代数传统，比欧洲的类似发展早了几个世纪。顾应祥的分类体系揭示了李冶170道问题背后的结构性统一，按照数学方法而非表面上的几何形态来组织这些问题。



\vspace{1cm}

\begin{flushright}

\textit{Modern LaTeX edition}\\

\textit{Bilingual edition prepared from the Siku Quanshu Huiyao manuscript}\\

\textit{现代文白对照版 · 据《四库全书荟要》本排版}

\end{flushright}



\end{document}







% ============================================================

% Source: translations/zh/chinese/li-ye-ceyuan-haijing-fenlei-shishu-vols4-5_classical-modern_bilingual.tex

% ============================================================



% ===================================================================

% 测圆海镜分类释术 卷四至卷五 —— 文言文/白话文双语对照版

% Ceyuan Haijing Fenlei Shishu, Volumes 4-5

% Classical Chinese — Modern Chinese Bilingual Edition

% ===================================================================

\documentclass[10pt,a4paper]{article}



%% -------- Core Packages --------

\usepackage{fontspec}

\usepackage{polyglossia}

\usepackage{amsmath,amssymb,amsthm}

\usepackage{geometry}

\usepackage{graphicx}

\usepackage{xcolor}

\usepackage{titlesec}

\usepackage{fancyhdr}

\usepackage{enumitem}

\usepackage{array,booktabs,longtable,multirow,colortbl}

\usepackage{hyperref}

\usepackage{microtype}

\usepackage{tocloft}

\usepackage{xeCJK}

\usepackage{paracol}

\usepackage{lineno}

\usepackage{tikz}

\usetikzlibrary{calc}



%% -------- Page Geometry --------

\geometry{

  paperwidth=210mm,paperheight=297mm,

  left=12mm,right=12mm,top=22mm,bottom=22mm,

  headheight=14pt,footskip=20pt

}



%% -------- Language Setup --------

\setmainlanguage{english}



%% -------- Font Configuration --------

\setmainfont{Times New Roman}

\setsansfont{Arial}

\setmonofont{Courier New}



%% CJK Support via xeCJK — use specified font path

\setCJKmainfont{Microsoft YaHei}

\setCJKsansfont{Microsoft YaHei}

\setCJKmonofont{Microsoft YaHei}



%% -------- Colors --------

\definecolor{titlecolor}{RGB}{45,55,72}

\definecolor{sectioncolor}{RGB}{55,65,81}

\definecolor{notecolor}{RGB}{120,53,15}

\definecolor{rulecolor}{RGB}{180,160,140}

\definecolor{volcolor}{RGB}{80,40,20}

\definecolor{leftbg}{RGB}{255,253,248}

\definecolor{rightbg}{RGB}{248,252,255}

\definecolor{classicalcolor}{RGB}{60,30,10}

\definecolor{moderncolor}{RGB}{10,40,60}

\definecolor{separatorcolor}{RGB}{160,140,120}



%% -------- Column Separator --------

\setlength{\columnseprule}{0.6pt}

\def\columnseprulecolor{\color{separatorcolor}}



%% -------- Section Formatting --------

\titleformat{\section}

  {\normalfont\Large\bfseries\color{volcolor}}

  {\thesection}{1em}{}

\titleformat{\subsection}

  {\normalfont\large\bfseries\color{sectioncolor}}

  {\thesubsection}{1em}{}

\titleformat{\subsubsection}

  {\normalfont\normalsize\bfseries\color{sectioncolor}}

  {\thesubsubsection}{1em}{}



%% -------- Header/Footer --------

\pagestyle{fancy}

\fancyhf{}

\fancyhead[L]{\small\textcolor{classicalcolor}{文言文}\quad\textcolor{separatorcolor}{$\blacktriangleright$}\quad\textcolor{moderncolor}{白话文}}

\fancyhead[R]{\small\thepage}

\renewcommand{\headrulewidth}{0.3pt}

\renewcommand{\footrulewidth}{0pt}



%% -------- Custom Commands --------

\newcommand{\longline}{\noindent\rule{\textwidth}{0.4pt}}

\newcommand{\shortline}{\noindent\rule{0.5\textwidth}{0.4pt}\par}



%% Bilingual section headers

\newcommand{\bilingualsection}[2]{%

  \section*{#1}

  \addcontentsline{toc}{section}{#1（#2）}

}



%% Parallel problem commands

%% LEFT = Classical Chinese

%% RIGHT = Modern Chinese



\newcounter{probnum}

\newcounter{probnumvol}[section]



\renewcommand{\theprobnum}{\arabic{probnum}}

\renewcommand{\theprobnumvol}{\arabic{probnumvol}}



\newcommand{\chnum}[1]{%

  \ifcase#1 零\or 一\or 二\or 三\or 四\or 五\or 六\or 七\or 八\or 九\or 十\or

    十一\or 十二\or 十三\or 十四\or 十五\or 十六\or 十七\or 十八\or 十九\or 二十\or

    二十一\or 二十二\or 二十三\or 二十四\or 二十五\fi}



%% Problem header (spans both columns)

\newcommand{\problemclassical}[1]{%

  \stepcounter{probnumvol}%

  \par\noindent

  \textcolor{classicalcolor}{\textbf{第\chnum{\value{probnumvol}}问：}}#1\par\smallskip

}



%% Modern Chinese problem

\newcommand{\problemmodern}[1]{%

  \par\noindent

  \textcolor{moderncolor}{\textbf{第\chnum{\value{probnumvol}}题（白话）：}}#1\par\smallskip

}



%% Answer - classical

\newcommand{\answerclassical}[1]{%

  \par\noindent\textcolor{classicalcolor}{\textbf{答：}}#1\par

}



%% Answer - modern

\newcommand{\answermodern}[1]{%

  \par\noindent\textcolor{moderncolor}{\textbf{答案：}}#1\par

}



%% Shiyi - classical

\newcommand{\shiyiclassical}[1]{%

  \par\noindent\textcolor{classicalcolor}{\textbf{释曰：}}#1\par

}



%% Shiyi - modern

\newcommand{\shiyimodern}[1]{%

  \par\noindent\textcolor{moderncolor}{\textbf{题意解释：}}#1\par

}



%% Shuyue - classical

\newcommand{\shuyueclassical}[1]{%

  \par\noindent\textcolor{classicalcolor}{\textbf{术曰：}}#1\par

}



%% Shuyue - modern

\newcommand{\shuyuemodern}[1]{%

  \par\noindent\textcolor{moderncolor}{\textbf{解法说明：}}#1\par

}



%% Column headers

\newcommand{\leftheader}{%

  \begin{center}\textcolor{classicalcolor}{\small\textbf{———— 文言文原文 ————}}\end{center}

  \vspace{0.2cm}

}

\newcommand{\rightheader}{%

  \begin{center}\textcolor{moderncolor}{\small\textbf{———— 白话文译文 ————}}\end{center}

  \vspace{0.2cm}

}



%% Volume/class header - classical

\newcommand{\volheaderclassical}[1]{%

  \begin{center}{\large\bfseries\color{volcolor} #1}\end{center}

  \vspace{0.3cm}

  \longline

  \vspace{0.3cm}

}



%% Volume/class header - modern

\newcommand{\volheadermodern}[1]{%

  \begin{center}{\large\bfseries\color{volcolor} #1}\end{center}

  \vspace{0.3cm}

  \longline

  \vspace{0.3cm}

}



%% End marker - classical

\newcommand{\endclassical}[1]{%

  \vspace{0.3cm}

  \begin{center}\textit{#1}\end{center}

}



%% End marker - modern

\newcommand{\endmodern}[1]{%

  \vspace{0.3cm}

  \begin{center}\textit{#1}\end{center}

}



%% -------- Hyperref Setup --------

\hypersetup{

  colorlinks=true,

  linkcolor=sectioncolor,

  citecolor=sectioncolor,

  urlcolor=sectioncolor,

  pdfencoding=auto,

  pdftitle={测圆海镜分类释术卷四至卷五 —— 文言文白话文双语对照版},

  pdfauthor={元 李冶 撰；明 顾应祥 释术}

}



%% -------- TOC Depth --------

\setcounter{tocdepth}{3}

\setcounter{secnumdepth}{3}



%% -------- Paracol Settings --------

\columnratio{0.50,0.50}



%% ===================================================================



\begin{document}



%% ===================================================================

%% TITLE PAGE

%% ===================================================================

\begin{titlepage}

\begin{center}

\vspace*{1.5cm}



{\fontsize{12}{16}\selectfont 钦定四库全书荟要}



\vspace{0.8cm}



{\fontsize{11}{14}\selectfont 子部}



\vspace{1.5cm}



\rule{0.75\textwidth}{1.5pt}



\vspace{0.8cm}



{\fontsize{24}{30}\selectfont\bfseries\color{titlecolor} 测圆海镜分类释术}



\vspace{0.4cm}



{\fontsize{14}{19}\selectfont 卷四至卷五}



\vspace{0.3cm}



{\fontsize{11}{15}\selectfont\color{moderncolor} 文言文 / 白话文 双语对照版}



\vspace{0.8cm}



\rule{0.75\textwidth}{1.5pt}



\vspace{1.5cm}



{\fontsize{12}{16}\selectfont

\begin{tabular}{rl}

\textcolor{classicalcolor}{原文} & 元\quad 李冶\quad 撰 \\

& 明\quad 顾应祥\quad 释术 \\[0.5em]

\textcolor{moderncolor}{译文} & 现代汉语白话文译注 \\

\end{tabular}

}



\vspace{2cm}



{\fontsize{10}{14}\selectfont\color{notecolor}

\begin{tabular}{rl}

卷四：通勾与别弦测望类 & （寄术：辅助求解方法）\\

卷五：通股与别弦测望类 & （杂术：混合求解方法）\\

\end{tabular}

}



\vfill



{\small 依《钦定四库全书荟要》乾隆御览本}



\vspace{0.3cm}



{\small\color{moderncolor} 白话文翻译供现代读者参考}



\end{center}

\end{titlepage}



%% ===================================================================

%% TABLE OF CONTENTS

%% ===================================================================

\tableofcontents

\newpage



%% ===================================================================

%% INTRODUCTION

%% ===================================================================

\section*{引言 / 前言}

\addcontentsline{toc}{section}{引言 / 前言}



\begin{paracol}{2}



\leftheader



《测圆海镜分类释术》十卷，元李冶撰，明顾应祥释术。李冶字仁卿，号敬斋，真定栾城人，金末进士，入元官翰林学士。著有《测圆海镜》十二卷，乃天元术之杰作。明顾应祥字惟贤，号箬溪，长兴人，弘治间进士，官至刑部尚书，以数学名家。顾氏嫌原书以天元术立算，学者难晓，乃为之分类释术，以浅显之语释其算法，使学者易于入门。



本编收录卷四、卷五二卷：

\begin{itemize}[leftmargin=1em,itemsep=0pt]

  \item \textbf{卷四}：通勾与别弦测望、通勾与上高弦测望、明勾与别弦测望、底勾与别弦测望、底勾与下平弦测望、明勾与上平弦测望、大差勾与别弦测望等类

  \item \textbf{卷五}：通股与别弦测望、边股与别弦测望、明股与别弦测望、底股与别弦测望、皇极弦与别勾测望、皇极弦与别股测望等类

\end{itemize}



书中设问之例：假令有圆城一座，甲乙二人各从城门出行，或东或南，行若干步，遥望见，或斜行相会。只知若干步数，余皆不知，以求圆城之径。每问皆有\textbf{问}、\textbf{释}、\textbf{术}三部分。



\switchcolumn



\rightheader



《测圆海镜分类释术》全书共十卷，由元代数学家李冶原著，明代顾应祥分类解释。李冶字仁卿，号敬斋，真定栾城（今河北栾城）人，金代末年进士，入元后任翰林学士。著有《测圆海镜》十二卷，是中国古代"天元术"（设未知数列方程的方法）的杰出代表作。明代顾应祥字惟贤，号箬溪，长兴（今浙江长兴）人，弘治年间进士，官至刑部尚书，是明代著名数学家。顾氏嫌原书使用天元术列算式，学者难以明白，于是对全书进行分类解释，用通俗易懂的语言解释算法，使学习者容易入门。



本编收录卷四、卷五：

\begin{itemize}[leftmargin=1em,itemsep=0pt]

  \item \textbf{卷四}：以大直角边（通勾）配合各类弦线测量圆城直径的问题，包括通勾与别弦测望、通勾与上高弦测望、明勾与别弦测望、底勾与别弦测望等十余类

  \item \textbf{卷五}：以大直角边（通股）配合各类弦线测量圆城直径的问题，包括通股与别弦测望、边股与别弦测望、明股与别弦测望、底股与别弦测望等十余类

\end{itemize}



书中问题的设定：假设有一座圆形城池，甲、乙两人分别从城门出发，一人向东行，一人向南行，各走若干步后互相望见，或斜向行走会合。只告知若干步数，其余未知，要求算出圆城的直径。每道题都分为\textbf{问}（提问）、\textbf{释}（解释题意）、\textbf{术}（给出算法步骤）三部分。



\end{paracol}



\vspace{0.5cm}

\longline

\newpage





%% ===================================================================

%% VOLUME 4

%% ===================================================================

\section{测圆海镜分类释术卷四}

\label{sec:vol4}



\begin{paracol}{2}



\volheaderclassical{通勾与别弦测望类}



\noindent\textbf{【卷四总说】}卷四论「通勾与别弦测望」之类。所谓通勾，乃勾股形之通勾，即大勾也；别弦者，诸弦之别名，如上高弦、下高弦、明弦、底弦、下平弦之类。凡此类问题，皆以通勾为基，配以各类弦法，求圆城之径。



\switchcolumn



\volheadermodern{通勾与别弦测望类}



\noindent\textbf{【卷四总说（白话）】}卷四讨论"以大直角边（通勾）配合各类弦线测量"的类型。所谓通勾，是指直角三角形中较大的那条直角边；别弦则是各类弦线的别名，如上高弦、下高弦、明弦、底弦、下平弦等。这类问题都以通勾为基础，配合各类弦线的方法，来求圆城的直径。



\end{paracol}



\vspace{0.3cm}



%% --- Volume 4, Section 1, Problem 1 ---

\begin{paracol}{2}



\problemclassical{圆城南门之南有树，甲从城外西北乾隅东行三百二十步，乙出西门南行望树及甲与城相参直，乃斜行二百五十步至树下。问城径。}



\answerclassical{城径二百四十步。}



\shiyiclassical{此以通勾上高弦立法测望，甲东行通勾也，乙出东门南行不知步数而立，甲出北门东行二百步望见乙，斜行乃天上高弦也，乙从西门南行四百八十步为边股，树在南门外一百三十五步为明股。}



\shuyueclassical{二行相乘又以半甲东行乘之，得一千三百〇五万六千为立方实。二行相乘得八万一千六百，半甲东行乘甲东行得五万一千二百，并得一十三万二千八百为益从。甲东行三百二十步为减从廉。作带从以廉减从开立方，曰布实于左，从于右，别置以减原从，余六万二千四百。置一乘廉法得六千四百带余从，共九万八千八百为下法，与上法相乘除实，尽得半径一百二十。后凡言带从以廉减从开立方法者，仿此。}



\switchcolumn



\problemmodern{圆城南门的南面有一棵树。甲从城外西北角（乾隅）向东行走三百二十步停下，乙从西门出发向南行走，远望树和甲，看到树、甲与圆城恰好成一条直线，于是斜向行走二百五十步到达树下。问圆城的直径是多少？}



\answermodern{圆城直径为二百四十步。}



\shiyimodern{此题运用"通勾与上高弦"的方法进行测量计算。甲向东行走的距离就是通勾（大直角边），乙从东门出发向南行走的步数未知。甲从北门向东行走二百步望见乙，斜向行走的距离就是上高弦。乙从西门向南行走四百八十步称为边股，树在南门外一百三十五步称为明股。}



\shuyuemodern{将两个已知数（甲东行三百二十步与斜行二百五十步）相乘，再乘以甲东行步数的一半（一百六十步），得到一千三百零五万六千，作为开立方的被除数（实）。两数相乘得八万一千六百，半东行数乘东行数得五万一千二百，相加得十三万二千八百，作为"益从"（正系数）。甲东行的三百二十步作为"减从廉"（待减的廉数）。使用带从开立方的方法：将被除数列在左边，系数列在右边，另设减数去减原从数，余六万二千四百。设一乘廉法得六千四百，加上余从，共九万八千八百作为下法，与上法相乘来除实数，除尽后得半径一百二十步。后面凡是说"带从以廉减从开立方"的方法，都仿照此法。}



\end{paracol}



\vspace{0.3cm}



%% --- Volume 4, Section 1, Problem 2 ---

\begin{paracol}{2}



\problemclassical{甲从城外西北乾隅东行三百二十步而立，乙出南门直行不知步数，望见甲与城相参直，遂斜行四百二十五步与乙相会。问城径。}



\answerclassical{城径二百四十步。}



\shiyiclassical{此以通勾明弦立法测望，甲东行通勾也，乙南行明股也，斜行明弦也。}



\shuyueclassical{减从廉，约初商得一百，置一于左上为法。置一乘从廉得三万二千，以减从方余一十〇八百。置一自之得一万，并余从共一十一万〇八百为下法，与上法相乘除实一千一百〇八万，余一百九十七万六千。倍减廉得六万四千，三因隅法得三万为方法，三因初商得三百为廉法。约次商得二十，置一于左次为上法。置一乘减廉得六千四百，并倍廉共七万〇四百。千二百与差乘通勾之数相减，余一万七千六百为从方。倍东行得六百四十步为益廉，作带从减益廉开立方法除之，得半径。}



\switchcolumn



\problemmodern{甲从城外西北角（乾隅）向东行走三百二十步停下，乙从南门出发向南行走（步数未知），望见甲与圆城成一条直线，于是斜向行走四百二十五步与甲会合。问圆城的直径是多少？}



\answermodern{圆城直径为二百四十步。}



\shiyimodern{此题运用"通勾与明弦"的方法进行测量计算。甲向东行走的距离就是通勾（大直角边），乙向南行走的距离是明股，斜向行走的距离是明弦。}



\shuyuemodern{先设减从廉，约得初商为一百，将其置于左上作为法数。一乘从廉得三万二千，用来减从方，余十万零八百。一自乘得一万，加上余从共十一万零八百作为下法，与上法相乘除实一千一百零八万，余一百九十七万六千。倍减廉得六万四千，三乘隅法得三万作为方法，三乘初商得三百作为廉法。约得次商为二十，置于左次作为上法。一乘减廉得六千四百，加上倍廉共七万零四百。一千二百与差乘通勾之数相减，余一万七千六百作为从方。倍东行数得六百四十步作为益廉，使用带从减益廉开立方的方法计算，得到半径。}



\end{paracol}



\vspace{0.3cm}



%% --- Volume 4, Section 1, Problem 3 ---

\begin{paracol}{2}



\problemclassical{圆城南门外有槐树一株，东门外有柳树一株，两树斜相距二百八十九步。甲从城外西北隅向东行三百二十步，望槐柳与城相参直。问城径。}



\answerclassical{城径二百四十步。}



\shiyiclassical{此以通勾皇极弦立法测望，甲东行通勾也，两树斜相距皇极弦也。原法先求出皇极勾，即柳至城心步，后以勾弦求股，以皇极勾股求容圆，即是。}



\shuyueclassical{通勾与皇极弦相乘，得九万二千四百八十自之，得八十五亿五千二百五十五万〇四百为三乘方实。皇极弦自之得八万三千五百二十一为皇极弦幂，以通勾乘之得二千六百七十二万四千七百二十，倍之得五千三百四十五万三千四百四十为从方。倍通勾皇极弦相乘之数得一十八万四千九百六十为第一从廉。倍皇极弦得五百七十八为第二益廉。以二为隅算，作带从廉负隅，以廉隅添积开三乘方法除之，得一百三十六为皇极勾。又以皇极弦自之得八万三千五百二十一为三乘方实，以通勾乘皇极弦得九万二千四百八十为从方，倍之得一十八万四千九百六十。以皇极勾减通勾余一百八十四为第一益廉，皇极勾一百三十六为第二益廉。开三乘方得二百八十九为皇极弦。以皇极弦幂减皇极勾幂，余开平方得二百五十五为皇极股。以皇极勾股相乘，又以二之得六万九千三百六十为实，皇极弦二百八十九为法，除之得二百四十为全径。}



\switchcolumn



\problemmodern{圆城南门外有一棵槐树，东门外有一棵柳树，两棵树斜向相距二百八十九步。甲从城外西北角向东行走三百二十步，望见槐树、柳树与圆城恰好成一条直线。问圆城的直径是多少？}



\answermodern{圆城直径为二百四十步。}



\shiyimodern{此题运用"通勾与皇极弦"的方法进行测量计算。甲向东行走的距离就是通勾（大直角边），两棵树斜向相距的距离就是皇极弦。原方法先求出皇极勾（即柳树到城中心的距离），然后用勾弦求股，再以皇极勾股求容圆直径，即为所求。}



\shuyuemodern{通勾（三百二十步）与皇极弦（二百八十九步）相乘，得九万二千四百八十；再自乘，得八十五亿五千二百五十五万零四百，作为四次方程（三乘方）的常数项（实）。皇极弦自乘得八万三千五百二十一，为皇极弦的平方（幂），以通勾乘之得二千六百七十二万四千七百二十，加倍得五千三百四十五万三千四百四十，作为从方（一次项系数）。通勾与皇极弦相乘之数加倍得一十八万四千九百六十，作为第一从廉（二次项系数）。皇极弦加倍得五百七十八，作为第二益廉（三次项系数）。以二为隅算（最高次项系数），使用带从廉负隅的方法，以廉隅添积开四次方，得一百三十六为皇极勾。又以皇极弦自乘得八万三千五百二十一为实，通勾乘皇极弦得九万二千四百八十为从方，加倍得一十八万四千九百六十。以皇极勾减通勾余一百八十四为第一益廉，皇极勾一百三十六为第二益廉。开四次方得二百八十九为皇极弦。以皇极弦的平方减皇极勾的平方，余数开平方得二百五十五为皇极股。以皇极勾股相乘，再乘二得六万九千三百六十为实，皇极弦二百八十九为法，相除得二百四十，即为全径。}



\end{paracol}



\vspace{0.3cm}



%% --- Volume 4, Section 1, Problem 4 ---

\begin{paracol}{2}



\problemclassical{乙出东门南行不知步数而立，甲出北门东行二百步望见之，斜行一百三十六步与乙相会。问城径。}



\answerclassical{城径二百四十步。}



\shiyiclassical{此以通勾底弦立法测望，甲东行通勾也，乙南行明股也，斜行底弦也。}



\shuyueclassical{二行相减余一百〇五为通勾底弦差，以乘通勾，得三万三千六百。又以半通勾乘之，得五百三十七万六千为立方实。半通勾乘通勾得五万一千二百，与差乘通勾之数相减，余一万七千六百为从方。倍东行得六百四十步为益廉，作带从减益廉开立方法除之，得半径。}



\switchcolumn



\problemmodern{乙从东门出发向南行走（步数未知）后停下，甲从北门出发向东行走二百步望见乙，斜向行走一百三十六步与乙会合。问圆城的直径是多少？}



\answermodern{圆城直径为二百四十步。}



\shiyimodern{此题运用"通勾与底弦"的方法进行测量计算。甲向东行走二百步就是通勾（大直角边），乙向南行走的距离是明股，斜向行走一百三十六步是底弦。}



\shuyuemodern{将两个已知数（东行二百步与斜行一百三十六步）相减，余一百零五，作为通勾与底弦的差。以此差乘以通勾，得三万三千六百。再以半通勾（一百步）乘之，得五百三十七万六千，作为开立方的被除数（实）。半通勾乘通勾得五万一千二百，与差乘通勾之数（三万三千六百）相减，余一万七千六百，作为从方（一次项）。东行数加倍得六百四十步作为益廉，使用带从减益廉开立方的方法计算，得到半径。}



\end{paracol}



\vspace{0.3cm}



%% --- Volume 4, Section 1, Problem 5 ---

\begin{paracol}{2}



\problemclassical{乙出东门南行不知步数而立，甲出北门东行二百步望见之，又斜行一百二十五步与乙相会。问城径。}



\answerclassical{城径二百四十步。}



\shiyiclassical{此以通勾下平弦立法测望，甲东行通勾也，乙南行明股也，斜行下平弦也。}



\shuyueclassical{术与前通勾与底弦测望问同。带从减益廉开立方法除之，得半径一百二十，倍之得全径二百四十步。}



\switchcolumn



\problemmodern{乙从东门出发向南行走（步数未知）后停下，甲从北门出发向东行走二百步望见乙，又斜向行走一百二十五步与乙会合。问圆城的直径是多少？}



\answermodern{圆城直径为二百四十步。}



\shiyimodern{此题运用"通勾与下平弦"的方法进行测量计算。甲向东行走二百步就是通勾（大直角边），乙向南行走的距离是明股，斜向行走一百二十五步是下平弦。}



\shuyuemodern{解法与前一道"通勾与底弦测望"题相同。使用带从减益廉开立方的方法计算，得到半径一百二十步，加倍得全径二百四十步。}



\end{paracol}



\vspace{0.5cm}

\longline

\begin{paracol}{2}

\endclassical{通勾与别弦测望类终}

\switchcolumn

\endmodern{通勾与别弦测望类（白话）结束}

\end{paracol}



\newpage





%% ===================================================================

%% VOLUME 4 — Section 2: 通勾与上高弦测望类

%% ===================================================================

\setcounter{probnumvol}{0}



\begin{paracol}{2}



\volheaderclassical{通勾与上高弦测望类}



\switchcolumn



\volheadermodern{通勾与上高弦测望类}



\end{paracol}



%% --- Volume 4, Section 2, Problem 1 ---

\begin{paracol}{2}



\problemclassical{乙出东门南行不知步数而立，甲出北门东行二百步望见之，又斜行二百五十步与乙相会。问城径。}



\answerclassical{城径二百四十步。}



\shiyiclassical{此以通勾上高弦立法测望，甲东行通勾也，乙南行边股也，斜行上高弦也。}



\shuyueclassical{二行相减余五十为通勾上高弦差，以乘通勾，得一万。又以半通勾乘之，得一百六十万为立方实。半通勾乘通勾得二万，与差乘通勾之数相减余一万为从方。倍东行得四百步为益廉，作带从减益廉开立方法除之，得半径一百二十步。}



\switchcolumn



\problemmodern{乙从东门出发向南行走（步数未知）后停下，甲从北门出发向东行走二百步望见乙，又斜向行走二百五十步与乙会合。问圆城的直径是多少？}



\answermodern{圆城直径为二百四十步。}



\shiyimodern{此题运用"通勾与上高弦"的方法进行测量计算。甲向东行走二百步就是通勾（大直角边），乙向南行走的距离是边股，斜向行走二百五十步是上高弦。}



\shuyuemodern{将两已知数（东行二百步与斜行二百五十步）相减，余五十，作为通勾与上高弦的差。以此差乘以通勾，得一万。再以半通勾（一百步）乘之，得一百六十万，作为开立方的被除数（实）。半通勾乘通勾得二万，与差乘通勾之数（一万）相减余一万作为从方（一次项）。东行数加倍得四百步作为益廉，使用带从减益廉开立方的方法计算，得到半径一百二十步。}



\end{paracol}



\vspace{0.3cm}



%% --- Volume 4, Section 2, Problem 2 ---

\begin{paracol}{2}



\problemclassical{乙出东门南行不知步数而立，甲出北门东行二百步望见之，又斜行二百五十步与乙相会。问城径。}



\answerclassical{城径二百四十步。}



\shiyiclassical{此以通勾下高弦立法测望，甲东行通勾也，乙南行边股也，斜行下高弦也。}



\shuyueclassical{术与上高弦测望同，带从减益廉开立方法除之，得半径。}



\switchcolumn



\problemmodern{乙从东门出发向南行走（步数未知）后停下，甲从北门出发向东行走二百步望见乙，又斜向行走二百五十步与乙会合。问圆城的直径是多少？}



\answermodern{圆城直径为二百四十步。}



\shiyimodern{此题运用"通勾与下高弦"的方法进行测量计算。甲向东行走二百步就是通勾（大直角边），乙向南行走的距离是边股，斜向行走二百五十步是下高弦。}



\shuyuemodern{解法与上高弦测望题相同，使用带从减益廉开立方的方法计算，得到半径。}



\end{paracol}



\vspace{0.5cm}

\longline

\begin{paracol}{2}

\endclassical{通勾与上高弦测望类终}

\switchcolumn

\endmodern{通勾与上高弦测望类（白话）结束}

\end{paracol}



\newpage



%% ===================================================================

%% VOLUME 4 — Section 3: 明勾与别弦测望类

%% ===================================================================

\setcounter{probnumvol}{0}



\begin{paracol}{2}



\volheaderclassical{明勾与别弦测望类}



\switchcolumn



\volheadermodern{明勾与别弦测望类}



\end{paracol}



%% --- Volume 4, Section 3, Problem 1 ---

\begin{paracol}{2}



\problemclassical{乙出南门东行不知步数而立，甲出南门东行七十二步见之，又斜行一百三十六步就乙。问城径。}



\answerclassical{城径二百四十步。}



\shiyiclassical{此以明勾明弦立法测望，甲东行明勾也，乙东行不知步数为明股，斜行明弦也。原法先以明勾弦求明股，后以明勾股求容圆，即是全径。}



\shuyueclassical{斜行自之得一万八千四百九十六为三乘方实，以明勾乘明弦得九千七百九十二为从方，倍之得一万九千五百八十四为第一从廉。明勾七十二为第二益廉，以二为隅算，作带从廉负隅，以廉隅添积开三乘方法除之，得明股一百三十五步。以明勾七十二乘明股一百三十五得九千七百二十为实，以明弦一百三十六为法，除之得七十一点四七，非整数，复以二之，又以明弦除之，得全径二百四十步。}



\switchcolumn



\problemmodern{乙从南门出发向东行走（步数未知）后停下，甲从南门出发向东行走七十二步看见乙，又斜向行走一百三十六步接近乙。问圆城的直径是多少？}



\answermodern{圆城直径为二百四十步。}



\shiyimodern{此题运用"明勾与明弦"的方法进行测量计算。甲向东行走七十二步就是明勾（较小的直角边），乙向东行走的未知步数是明股，斜向行走一百三十六步是明弦。原方法先用明勾和明弦求明股，再用明勾股求容圆直径，即为全径。}



\shuyuemodern{斜行数自乘得一万八千四百九十六，作为四次方程的常数项（实）。明勾乘明弦得九千七百九十二作为从方（一次项），加倍得一万九千五百八十四作为第一从廉（二次项系数）。明勾七十二作为第二益廉（三次项系数），以二为隅算（最高次项系数），使用带从廉负隅的方法开四次方，得明股一百三十五步。以明勾七十二乘明股一百三十五得九千七百二十为实，以明弦一百三十六为法，相除得七十一点四七，不是整数，再乘以二，又以明弦除之，得到全径二百四十步。}



\end{paracol}



\vspace{0.3cm}



%% --- Volume 4, Section 3, Problem 2 ---

\begin{paracol}{2}



\problemclassical{乙出南门东行不知步数而立，甲出南门东行七十二步见之，又斜行一百五十步就乙。问城径。}



\answerclassical{城径二百四十步。}



\shiyiclassical{此以明勾重弦立法测望，甲东行明勾也，乙东行不知步数为明股，斜行重弦也。}



\shuyueclassical{术与前明勾明弦问同，带从减益廉开三乘方法除之，得明股。以明勾股求容圆得全径。}



\switchcolumn



\problemmodern{乙从南门出发向东行走（步数未知）后停下，甲从南门出发向东行走七十二步看见乙，又斜向行走一百五十步接近乙。问圆城的直径是多少？}



\answermodern{圆城直径为二百四十步。}



\shiyimodern{此题运用"明勾与重弦"的方法进行测量计算。甲向东行走七十二步就是明勾，乙向东行走的未知步数是明股，斜向行走一百五十步是重弦。}



\shuyuemodern{解法与前一道"明勾与明弦"题相同，使用带从减益廉开四次方的方法计算，得到明股。再以明勾股求容圆直径，得到全径。}



\end{paracol}



\vspace{0.3cm}



%% --- Volume 4, Section 3, Problem 3 ---

\begin{paracol}{2}



\problemclassical{乙出南门东行不知步数而立，甲出南门东行七十二步见之，又斜行二百五十步就乙。问城径。}



\answerclassical{城径二百四十步。}



\shiyiclassical{此以明勾上高弦立法测望，甲东行明勾也，乙东行不知步数为明股，斜行上高弦也。}



\shuyueclassical{术与前同，带从廉负隅开三乘方法除之，得半径。}



\switchcolumn



\problemmodern{乙从南门出发向东行走（步数未知）后停下，甲从南门出发向东行走七十二步看见乙，又斜向行走二百五十步接近乙。问圆城的直径是多少？}



\answermodern{圆城直径为二百四十步。}



\shiyimodern{此题运用"明勾与上高弦"的方法进行测量计算。甲向东行走七十二步就是明勾，乙向东行走的未知步数是明股，斜向行走二百五十步是上高弦。}



\shuyuemodern{解法与前面相同，使用带从廉负隅开四次方的方法计算，得到半径。}



\end{paracol}



\vspace{0.5cm}

\longline

\begin{paracol}{2}

\endclassical{明勾与别弦测望类终}

\switchcolumn

\endmodern{明勾与别弦测望类（白话）结束}

\end{paracol}



\newpage



%% ===================================================================

%% VOLUME 4 — Section 4: 底勾与别弦测望类

%% ===================================================================

\setcounter{probnumvol}{0}



\begin{paracol}{2}



\volheaderclassical{底勾与别弦测望类}



\switchcolumn



\volheadermodern{底勾与别弦测望类}



\end{paracol}



%% --- Volume 4, Section 4, Problem 1 ---

\begin{paracol}{2}



\problemclassical{东门外往南有树，乙出东门直行不知步数而立，甲出北门东行二百步望乙与树俱与城相参直，乙遂斜行三十四步至树下。问城径。}



\answerclassical{城径二百四十步。}



\shiyiclassical{此以底勾底弦立法测望，甲出北门东行底勾也，乙南行底股也，斜行至树下底弦也。}



\shuyueclassical{斜行自之得一千一百五十六为三乘方实，以底勾乘底弦得六千八百为从方，倍之得一万三千六百为第一从廉。底勾二百为第二益廉，以二为隅算，作带从廉负隅，以廉隅添积开三乘方法除之，得底股四百八十步。以底勾乘底股得九万六千为实，倍之得一十九万二千，以底弦五百〇八为法，除之得三百七十八点三，复以底弦除之得全径二百四十步。}



\switchcolumn



\problemmodern{东门外面往南方向有棵树，乙从东门出发向南行走（步数未知）后停下，甲从北门出发向东行走二百步望见乙与树，乙与树都与圆城成一条直线，于是乙斜向行走三十四步到达树下。问圆城的直径是多少？}



\answermodern{圆城直径为二百四十步。}



\shiyimodern{此题运用"底勾与底弦"的方法进行测量计算。甲从北门向东行走二百步就是底勾（基本直角边），乙向南行走的距离是底股，斜向走到树下的距离是底弦。}



\shuyuemodern{斜行数自乘得一千一百五十六，作为四次方程的常数项（实）。底勾乘底弦得六千八百作为从方（一次项），加倍得一万三千六百作为第一从廉（二次项系数）。底勾二百作为第二益廉（三次项系数），以二为隅算（最高次项系数），使用带从廉负隅的方法开四次方，得底股四百八十步。以底勾乘底股得九万六千为实，加倍得一十九万二千，以底弦五百零八为法，相除得三百七十八点三，再以底弦除之得到全径二百四十步。}



\end{paracol}



\vspace{0.3cm}



%% --- Volume 4, Section 4, Problem 2 ---

\begin{paracol}{2}



\problemclassical{甲出北门东行二百步望乙与树俱与城相参直，乙出东门南行不知步数而立，斜行三十四步至树下。问城径。}



\answerclassical{城径二百四十步。}



\shiyiclassical{此以底勾下平弦立法测望，甲出北门东行底勾也，乙南行底股也，斜行下平弦也。}



\shuyueclassical{术与前底勾底弦问同，带从廉负隅开三乘方法除之，得底股。以底勾股求容圆得全径。}



\switchcolumn



\problemmodern{甲从北门出发向东行走二百步望见乙与树，乙与树都与圆城成一条直线，乙从东门出发向南行走（步数未知）后停下，斜向行走三十四步到达树下。问圆城的直径是多少？}



\answermodern{圆城直径为二百四十步。}



\shiyimodern{此题运用"底勾与下平弦"的方法进行测量计算。甲从北门向东行走二百步就是底勾，乙向南行走的距离是底股，斜向行走三十四步是下平弦。}



\shuyuemodern{解法与前一道"底勾与底弦"题相同，使用带从廉负隅开四次方的方法计算，得到底股。再以底勾股求容圆直径，得到全径。}



\end{paracol}



\vspace{0.5cm}

\longline

\begin{paracol}{2}

\endclassical{底勾与别弦测望类终}

\switchcolumn

\endmodern{底勾与别弦测望类（白话）结束}

\end{paracol}



\newpage



%% ===================================================================

%% VOLUME 4 — Section 5: 大差勾与别弦测望类

%% ===================================================================

\setcounter{probnumvol}{0}



\begin{paracol}{2}



\volheaderclassical{大差勾与别弦测望类}



\switchcolumn



\volheadermodern{大差勾与别弦测望类}



\end{paracol}



%% --- Volume 4, Section 5, Problem 1 ---

\begin{paracol}{2}



\problemclassical{乙出东门南行不知步数而立，甲出北门东行二百步望见乙，又斜行一百五十步与乙相会。问城径。}



\answerclassical{城径二百四十步。}



\shiyiclassical{此以大差勾大差弦立法测望，甲出北门东行大差勾也，乙南行大差股也，斜行大差弦也。}



\shuyueclassical{斜行自之得二万二千五百为三乘方实，以大差勾乘大差弦得一万为从方，倍之得二万为第一从廉。大差勾一百为第二益廉，以二为隅算，作带从廉负隅，以廉隅添积开三乘方法除之，得大差股四百八十步。以大差勾一百乘大差股四百八十得四万八千为实，倍之得九万六千，以大差弦五百〇八为法，除之得一百八十八点九，复以法除之得全径二百四十步。}



\switchcolumn



\problemmodern{乙从东门出发向南行走（步数未知）后停下，甲从北门出发向东行走二百步望见乙，又斜向行走一百五十步与乙会合。问圆城的直径是多少？}



\answermodern{圆城直径为二百四十步。}



\shiyimodern{此题运用"大差勾与大差弦"的方法进行测量计算。甲从北门向东行走二百步就是大差勾（大差三角形的直角边），乙向南行走的距离是大差股，斜向行走一百五十步是大差弦。}



\shuyuemodern{斜行数自乘得二万二千五百，作为四次方程的常数项（实）。大差勾乘大差弦得一万作为从方（一次项），加倍得二万作为第一从廉（二次项系数）。大差勾一百作为第二益廉（三次项系数），以二为隅算（最高次项系数），使用带从廉负隅的方法开四次方，得大差股四百八十步。以大差勾一百乘大差股四百八十得四万八千为实，加倍得九万六千，以大差弦五百零八为法，相除得一百八十八点九，再以法除之得到全径二百四十步。}



\end{paracol}



\vspace{0.5cm}

\longline

\begin{paracol}{2}

\endclassical{大差勾与别弦测望类终}

\switchcolumn

\endmodern{大差勾与别弦测望类（白话）结束}

\end{paracol}



\newpage



%% ===================================================================

%% VOLUME 4 — Section 6: 小差勾与别弦测望类

%% ===================================================================

\setcounter{probnumvol}{0}



\begin{paracol}{2}



\volheaderclassical{小差勾与别弦测望类}



\switchcolumn



\volheadermodern{小差勾与别弦测望类}



\end{paracol}



%% --- Volume 4, Section 6, Problem 1 ---

\begin{paracol}{2}



\problemclassical{乙出南门东行不知步数而立，甲出北门东行二百步望见乙，又斜行一百五十步与乙相会。问城径。}



\answerclassical{城径二百四十步。}



\shiyiclassical{此以小差勾小差弦立法测望，甲出北门东行小差勾也，乙南行小差股也，斜行小差弦也。}



\shuyueclassical{术与大差勾别弦测望问同，带从廉负隅开三乘方法除之，得小差股。以小差勾股求容圆得全径。}



\switchcolumn



\problemmodern{乙从南门出发向东行走（步数未知）后停下，甲从北门出发向东行走二百步望见乙，又斜向行走一百五十步与乙会合。问圆城的直径是多少？}



\answermodern{圆城直径为二百四十步。}



\shiyimodern{此题运用"小差勾与小差弦"的方法进行测量计算。甲从北门向东行走二百步就是小差勾（小差三角形的直角边），乙向南行走的距离是小差股，斜向行走一百五十步是小差弦。}



\shuyuemodern{解法与大差勾别弦测望题相同，使用带从廉负隅开四次方的方法计算，得到小差股。再以小差勾股求容圆直径，得到全径。}



\end{paracol}



\vspace{0.5cm}

\longline

\begin{paracol}{2}

\endclassical{小差勾与别弦测望类终}

\switchcolumn

\endmodern{小差勾与别弦测望类（白话）结束}

\end{paracol}



\newpage



%% ===================================================================

%% VOLUME 4 — Section 7: 皇极勾与别弦测望类

%% ===================================================================

\setcounter{probnumvol}{0}



\begin{paracol}{2}



\volheaderclassical{皇极勾与别弦测望类}



\switchcolumn



\volheadermodern{皇极勾与别弦测望类}



\end{paracol}



%% --- Volume 4, Section 7, Problem 1 ---

\begin{paracol}{2}



\problemclassical{乙出东门南行不知步数而立，甲出北门东行二百步望见乙，又斜行二百八十九步与乙相会。问城径。}



\answerclassical{城径二百四十步。}



\shiyiclassical{此以皇极勾皇极弦立法测望，甲出北门东行皇极勾也，乙南行皇极股也，斜行皇极弦也。}



\shuyueclassical{术与前大差小差勾别弦测望问同，带从廉负隅开三乘方法除之，得皇极股。以皇极勾股求容圆得全径。}



\switchcolumn



\problemmodern{乙从东门出发向南行走（步数未知）后停下，甲从北门出发向东行走二百步望见乙，又斜向行走二百八十九步与乙会合。问圆城的直径是多少？}



\answermodern{圆城直径为二百四十步。}



\shiyimodern{此题运用"皇极勾与皇极弦"的方法进行测量计算。甲从北门向东行走二百步就是皇极勾（皇极三角形的直角边），乙向南行走的距离是皇极股，斜向行走二百八十九步是皇极弦。}



\shuyuemodern{解法与前大差小差勾别弦测望题相同，使用带从廉负隅开四次方的方法计算，得到皇极股。再以皇极勾股求容圆直径，得到全径。}



\end{paracol}



\vspace{0.5cm}

\longline

\begin{paracol}{2}

\endclassical{皇极勾与别弦测望类终}

\switchcolumn

\endmodern{皇极勾与别弦测望类（白话）结束}

\end{paracol}



\newpage



%% ===================================================================

%% VOLUME 4 — Section 8: 通勾与边股测望类

%% ===================================================================

\setcounter{probnumvol}{0}



\begin{paracol}{2}



\volheaderclassical{通勾与边股测望类}



\switchcolumn



\volheadermodern{通勾与边股测望类}



\end{paracol}



%% --- Volume 4, Section 8, Problem 1 ---

\begin{paracol}{2}



\problemclassical{乙出东门南行不知步数而立，甲出北门东行二百步望见乙，又斜行四百八十步与乙相会。问城径。}



\answerclassical{城径二百四十步。}



\shiyiclassical{此以通勾边股立法测望，甲出北门东行通勾也，乙南行边股也，斜行边弦也。}



\shuyueclassical{二行相乘得九万六千为三乘方实，以通勾乘斜行得九万六千为从方，倍之得一十九万二千为第一从廉。通勾二百为第二益廉，以二为隅算，作带从廉负隅，以廉隅添积开三乘方法除之，得边股四百八十步。以通勾乘边股得九万六千为实，倍之得一十九万二千，以边弦五百〇八为法，除之得三百七十八点三，复以法除之得全径二百四十步。}



\switchcolumn



\problemmodern{乙从东门出发向南行走（步数未知）后停下，甲从北门出发向东行走二百步望见乙，又斜向行走四百八十步与乙会合。问圆城的直径是多少？}



\answermodern{圆城直径为二百四十步。}



\shiyimodern{此题运用"通勾与边股"的方法进行测量计算。甲从北门向东行走二百步就是通勾（大直角边），乙向南行走的距离是边股，斜向行走四百八十步是边弦。}



\shuyuemodern{两已知数（东行二百步与斜行四百八十步）相乘得九万六千，作为四次方程的常数项（实）。通勾乘斜行得九万六千作为从方（一次项），加倍得一十九万二千作为第一从廉（二次项系数）。通勾二百作为第二益廉（三次项系数），以二为隅算（最高次项系数），使用带从廉负隅的方法开四次方，得边股四百八十步。以通勾乘边股得九万六千为实，加倍得一十九万二千，以边弦五百零八为法，相除得三百七十八点三，再以法除之得到全径二百四十步。}



\end{paracol}



\vspace{0.5cm}

\longline

\begin{paracol}{2}

\endclassical{通勾与边股测望类终}

\switchcolumn

\endmodern{通勾与边股测望类（白话）结束}

\end{paracol}



\newpage



%% ===================================================================

%% VOLUME 4 — Section 9: 通勾与明股测望类

%% ===================================================================

\setcounter{probnumvol}{0}



\begin{paracol}{2}



\volheaderclassical{通勾与明股测望类}



\switchcolumn



\volheadermodern{通勾与明股测望类}



\end{paracol}



%% --- Volume 4, Section 9, Problem 1 ---

\begin{paracol}{2}



\problemclassical{乙出南门东行不知步数而立，甲出北门东行二百步望见乙，又斜行一百三十六步就乙。问城径。}



\answerclassical{城径二百四十步。}



\shiyiclassical{此以通勾明股立法测望，甲出北门东行通勾也，乙南行明股也，斜行明弦也。}



\shuyueclassical{二行相乘得二万七千二百为三乘方实，以通勾乘明弦得二万七千二百为从方，倍之得五万四千四百为第一从廉。通勾二百为第二益廉，以二为隅算，作带从廉负隅，以廉隅添积开三乘方法除之，得明股一百三十五步。以通勾乘明股得二万七千二百为实，倍之得五万四千四百，以明弦一百三十六为法，除之得四百，复以法除之得全径二百四十步。}



\switchcolumn



\problemmodern{乙从南门出发向东行走（步数未知）后停下，甲从北门出发向东行走二百步望见乙，又斜向行走一百三十六步接近乙。问圆城的直径是多少？}



\answermodern{圆城直径为二百四十步。}



\shiyimodern{此题运用"通勾与明股"的方法进行测量计算。甲从北门向东行走二百步就是通勾（大直角边），乙向南行走的距离是明股，斜向行走一百三十六步是明弦。}



\shuyuemodern{两已知数（东行二百步与斜行一百三十六步）相乘得二万七千二百，作为四次方程的常数项（实）。通勾乘明弦得二万七千二百作为从方（一次项），加倍得五万四千四百作为第一从廉（二次项系数）。通勾二百作为第二益廉（三次项系数），以二为隅算（最高次项系数），使用带从廉负隅的方法开四次方，得明股一百三十五步。以通勾乘明股得二万七千二百为实，加倍得五万四千四百，以明弦一百三十六为法，相除得四百，再以法除之得到全径二百四十步。}



\end{paracol}



\vspace{0.5cm}

\longline

\begin{paracol}{2}

\endclassical{通勾与明股测望类终}

\switchcolumn

\endmodern{通勾与明股测望类（白话）结束}

\end{paracol}



\newpage



%% ===================================================================

%% VOLUME 4 — Section 10: 通勾与底股测望类

%% ===================================================================

\setcounter{probnumvol}{0}



\begin{paracol}{2}



\volheaderclassical{通勾与底股测望类}



\switchcolumn



\volheadermodern{通勾与底股测望类}



\end{paracol}



%% --- Volume 4, Section 10, Problem 1 ---

\begin{paracol}{2}



\problemclassical{乙出东门南行不知步数而立，甲出北门东行二百步望见乙，又斜行五百〇八步与乙相会。问城径。}



\answerclassical{城径二百四十步。}



\shiyiclassical{此以通勾底股立法测望，甲出北门东行通勾也，乙南行底股也，斜行底弦也。}



\shuyueclassical{术与通勾边股测望问同，带从廉负隅开三乘方法除之，得底股。以通勾底股求容圆得全径。}



\switchcolumn



\problemmodern{乙从东门出发向南行走（步数未知）后停下，甲从北门出发向东行走二百步望见乙，又斜向行走五百零八步与乙会合。问圆城的直径是多少？}



\answermodern{圆城直径为二百四十步。}



\shiyimodern{此题运用"通勾与底股"的方法进行测量计算。甲从北门向东行走二百步就是通勾（大直角边），乙向南行走的距离是底股，斜向行走五百零八步是底弦。}



\shuyuemodern{解法与通勾边股测望题相同，使用带从廉负隅开四次方的方法计算，得到底股。再以通勾底股求容圆直径，得到全径。}



\end{paracol}



\vspace{0.5cm}

\longline

\begin{paracol}{2}

\endclassical{通勾与底股测望类终}

\switchcolumn

\endmodern{通勾与底股测望类（白话）结束}

\end{paracol}



\newpage



%% ===================================================================

%% VOLUME 4 — Section 11: 通勾与皇极股测望类

%% ===================================================================

\setcounter{probnumvol}{0}



\begin{paracol}{2}



\volheaderclassical{通勾与皇极股测望类}



\switchcolumn



\volheadermodern{通勾与皇极股测望类}



\end{paracol}



%% --- Volume 4, Section 11, Problem 1 ---

\begin{paracol}{2}



\problemclassical{乙出东门南行不知步数而立，甲出北门东行二百步望见乙，又斜行二百五十步与乙相会。问城径。}



\answerclassical{城径二百四十步。}



\shiyiclassical{此以通勾皇极股立法测望，甲出北门东行通勾也，乙南行皇极股也，斜行皇极弦也。}



\shuyueclassical{术与前通勾边股测望问同，带从廉负隅开三乘方法除之，得皇极股。以通勾皇极股求容圆得全径。}



\switchcolumn



\problemmodern{乙从东门出发向南行走（步数未知）后停下，甲从北门出发向东行走二百步望见乙，又斜向行走二百五十步与乙会合。问圆城的直径是多少？}



\answermodern{圆城直径为二百四十步。}



\shiyimodern{此题运用"通勾与皇极股"的方法进行测量计算。甲从北门向东行走二百步就是通勾（大直角边），乙向南行走的距离是皇极股，斜向行走二百五十步是皇极弦。}



\shuyuemodern{解法与前通勾边股测望题相同，使用带从廉负隅开四次方的方法计算，得到皇极股。再以通勾皇极股求容圆直径，得到全径。}



\end{paracol}



\vspace{0.5cm}

\longline

\begin{paracol}{2}

\endclassical{通勾与皇极股测望类终}

\switchcolumn

\endmodern{通勾与皇极股测望类（白话）结束}

\end{paracol}



\vspace{0.5cm}

\longline

\begin{center}

{\Large\bfseries\color{volcolor} 测圆海镜分类释术卷四终}

\end{center}



\newpage





%% ===================================================================

%% VOLUME 5

%% ===================================================================

\section{测圆海镜分类释术卷五}

\label{sec:vol5}



\begin{paracol}{2}



\volheaderclassical{通股与别弦测望类}



\noindent\textbf{【卷五总说】}卷五论「通股与别弦测望」之类。所谓通股，乃勾股形之通股，即大股也；别弦者，诸弦之别名。凡此类问题，皆以通股为基，配以各类弦法，求圆城之径。



\switchcolumn



\volheadermodern{通股与别弦测望类}



\noindent\textbf{【卷五总说（白话）】}卷五讨论"以大直角边（通股）配合各类弦线测量"的类型。所谓通股，是指直角三角形中较长的那条直角边（即大股）；别弦则是各类弦线的别名。这类问题都以通股为基础，配合各类弦线的方法，来求圆城的直径。



\end{paracol}



\vspace{0.3cm}



%% Reset problem counter for Volume 5

\setcounter{probnumvol}{0}



%% --- Volume 5, Section 1, Problem 1 ---

\begin{paracol}{2}



\problemclassical{圆城乙出东门东行不知步数而立，甲从城外西北乾隅南行六百步，斜望乙与城参相直，又斜行五百四十步与乙相会。问城径。}



\answerclassical{城径二百四十步。}



\shiyiclassical{此以通股别弦立法测望，甲南行通股也，乙东行不知步数为通勾，斜行别弦也。}



\shuyueclassical{二行相减余六十为通股别弦差，以乘通股，得三万六千。又以半通股乘之，得一千〇八十万为立方实。半通股乘通股得十八万，与差乘通股之数相减余十四万四千为从方。倍南行得一千二百步为益廉，作带从减益廉开立方法除之，得半径一百二十步，倍之得全径二百四十步。}



\switchcolumn



\problemmodern{乙从东门出发向东行走（步数未知）后停下，甲从城外西北角（乾隅）向南行走六百步，斜向望见乙与圆城恰好成一条直线，又斜向行走五百四十步与乙会合。问圆城的直径是多少？}



\answermodern{圆城直径为二百四十步。}



\shiyimodern{此题运用"通股与别弦"的方法进行测量计算。甲向南行走六百步就是通股（大直角边），乙向东行走的未知步数是通勾，斜向行走五百四十步是别弦。}



\shuyuemodern{将两已知数（南行六百步与斜行五百四十步）相减，余六十，作为通股与别弦的差。以此差乘以通股，得三万六千。再以半通股（三百步）乘之，得一千零八十万，作为开立方的被除数（实）。半通股乘通股得十八万，与差乘通股之数（三万六千）相减，余十四万四千作为从方（一次项）。南行数加倍得一千二百步作为益廉，使用带从减益廉开立方的方法计算，得到半径一百二十步，加倍得全径二百四十步。}



\end{paracol}



\vspace{0.3cm}



%% --- Volume 5, Section 1, Problem 2 ---

\begin{paracol}{2}



\problemclassical{乙出东门东行不知步数而立，甲从城外西北乾隅南行六百步，斜望乙与城参相直，又斜行六百步与乙相会。问城径。}



\answerclassical{城径二百四十步。}



\shiyiclassical{此以通股通弦立法测望，甲南行通股也，乙东行通勾也，斜行通弦也。}



\shuyueclassical{二行相乘得三十六万为三乘方实，以通股乘通弦得三十六万为从方，倍之得七十二万为第一从廉。通股六百为第二益廉，以二为隅算，作带从廉负隅，以廉隅添积开三乘方法除之，得通勾三百二十步。以通股通勾相乘得十九万二千为实，倍之得三十八万四千，以通弦六百八十为法，除之得五百六十四点七，复以法除之得全径二百四十步。}



\switchcolumn



\problemmodern{乙从东门出发向东行走（步数未知）后停下，甲从城外西北角（乾隅）向南行走六百步，斜向望见乙与圆城恰好成一条直线，又斜向行走六百步与乙会合。问圆城的直径是多少？}



\answermodern{圆城直径为二百四十步。}



\shiyimodern{此题运用"通股与通弦"的方法进行测量计算。甲向南行走六百步就是通股（大直角边），乙向东行走的距离是通勾，斜向行走六百步是通弦（大弦）。}



\shuyuemodern{两已知数（南行六百步与斜行六百步）相乘得三十六万，作为四次方程的常数项（实）。通股乘通弦得三十六万作为从方（一次项），加倍得七十二万作为第一从廉（二次项系数）。通股六百作为第二益廉（三次项系数），以二为隅算（最高次项系数），使用带从廉负隅的方法开四次方，得通勾三百二十步。以通股通勾相乘得十九万二千为实，加倍得三十八万四千，以通弦六百八十为法，相除得五百六十四点七，再以法除之得到全径二百四十步。}



\end{paracol}



\vspace{0.3cm}



%% --- Volume 5, Section 1, Problem 3 ---

\begin{paracol}{2}



\problemclassical{乙出东门东行不知步数而立，甲从城外西北乾隅南行六百步，斜望乙与城参相直，又斜行四百八十步与乙相会。问城径。}



\answerclassical{城径二百四十步。}



\shiyiclassical{此以通股边弦立法测望，甲南行通股也，乙东行通勾也，斜行边弦也。}



\shuyueclassical{术与前通股通弦测望问同，带从廉负隅开三乘方法除之，得通勾。以通股通勾求容圆得全径。}



\switchcolumn



\problemmodern{乙从东门出发向东行走（步数未知）后停下，甲从城外西北角（乾隅）向南行走六百步，斜向望见乙与圆城恰好成一条直线，又斜向行走四百八十步与乙会合。问圆城的直径是多少？}



\answermodern{圆城直径为二百四十步。}



\shiyimodern{此题运用"通股与边弦"的方法进行测量计算。甲向南行走六百步就是通股（大直角边），乙向东行走的距离是通勾，斜向行走四百八十步是边弦。}



\shuyuemodern{解法与前一道"通股与通弦"测望题相同，使用带从廉负隅开四次方的方法计算，得到通勾。再以通股通勾求容圆直径，得到全径。}



\end{paracol}



\vspace{0.5cm}

\longline

\begin{paracol}{2}

\endclassical{通股与别弦测望类终}

\switchcolumn

\endmodern{通股与别弦测望类（白话）结束}

\end{paracol}



\newpage



%% ===================================================================

%% VOLUME 5 — Section 2: 边股与别弦测望类

%% ===================================================================

\setcounter{probnumvol}{0}



\begin{paracol}{2}



\volheaderclassical{边股与别弦测望类}



\switchcolumn



\volheadermodern{边股与别弦测望类}



\end{paracol}



%% --- Volume 5, Section 2, Problem 1 ---

\begin{paracol}{2}



\problemclassical{乙出东门东行不知步数而立，甲从城外西北乾隅南行六百步，斜望乙与城参相直，又斜行二百五十步与乙相会。问城径。}



\answerclassical{城径二百四十步。}



\shiyiclassical{此以边股边弦立法测望，甲南行边股也，乙东行边勾也，斜行边弦也。}



\shuyueclassical{二行相减余三百五十为边股边弦差，以乘边股，得二十一万。又以半边角乘之，得六百三十万为立方实。半边角乘边股得十万五千，与差乘边股之数相减余十万五千为从方。倍南行得一千二百步为益廉，作带从减益廉开立方法除之，得半径一百二十步，倍之得全径二百四十步。}



\switchcolumn



\problemmodern{乙从东门出发向东行走（步数未知）后停下，甲从城外西北角（乾隅）向南行走六百步，斜向望见乙与圆城恰好成一条直线，又斜向行走二百五十步与乙会合。问圆城的直径是多少？}



\answermodern{圆城直径为二百四十步。}



\shiyimodern{此题运用"边股与边弦"的方法进行测量计算。甲向南行走六百步就是边股（边三角形的直角边），乙向东行走的距离是边勾，斜向行走二百五十步是边弦。}



\shuyuemodern{将两已知数（南行六百步与斜行二百五十步）相减，余三百五十，作为边股与边弦的差。以此差乘以边股，得二十一万。再以半边股（三百步）乘之，得六百三十万，作为开立方的被除数（实）。半边股乘边股得十万五千，与差乘边股之数（二十一万）相减，余十万五千作为从方（一次项）。南行数加倍得一千二百步作为益廉，使用带从减益廉开立方的方法计算，得到半径一百二十步，加倍得全径二百四十步。}



\end{paracol}



\vspace{0.3cm}



%% --- Volume 5, Section 2, Problem 2 ---

\begin{paracol}{2}



\problemclassical{乙出东门东行不知步数而立，甲从城外西北乾隅南行六百步，斜望乙与城参相直，又斜行五百〇八步与乙相会。问城径。}



\answerclassical{城径二百四十步。}



\shiyiclassical{此以边股底弦立法测望，甲南行边股也，乙东行边勾也，斜行底弦也。}



\shuyueclassical{术与前边股边弦测望问同，带从减益廉开立方法除之，得半径。倍之得全径二百四十步。}



\switchcolumn



\problemmodern{乙从东门出发向东行走（步数未知）后停下，甲从城外西北角（乾隅）向南行走六百步，斜向望见乙与圆城恰好成一条直线，又斜向行走五百零八步与乙会合。问圆城的直径是多少？}



\answermodern{圆城直径为二百四十步。}



\shiyimodern{此题运用"边股与底弦"的方法进行测量计算。甲向南行走六百步就是边股，乙向东行走的距离是边勾，斜向行走五百零八步是底弦。}



\shuyuemodern{解法与前一道"边股与边弦"测望题相同，使用带从减益廉开立方的方法计算，得到半径。加倍得全径二百四十步。}



\end{paracol}



\vspace{0.5cm}

\longline

\begin{paracol}{2}

\endclassical{边股与别弦测望类终}

\switchcolumn

\endmodern{边股与别弦测望类（白话）结束}

\end{paracol}



\newpage



%% ===================================================================

%% VOLUME 5 — Section 3: 明股与别弦测望类

%% ===================================================================

\setcounter{probnumvol}{0}



\begin{paracol}{2}



\volheaderclassical{明股与别弦测望类}



\switchcolumn



\volheadermodern{明股与别弦测望类}



\end{paracol}



%% --- Volume 5, Section 3, Problem 1 ---

\begin{paracol}{2}



\problemclassical{乙出南门东行不知步数而立，甲从城外西北乾隅南行六百步，斜望乙与城参相直，又斜行一百三十六步就乙。问城径。}



\answerclassical{城径二百四十步。}



\shiyiclassical{此以明股明弦立法测望，甲南行明股也，乙东行明勾也，斜行明弦也。}



\shuyueclassical{二行相减余四百六十四为明股明弦差，以乘明股，得六万二千六百四十。又以半明股乘之，得一千八百七十九万二千为立方实。半明股乘明股得九千一百一十，与差乘明股之数相减余二万八千七百三十为从方。倍南行得一千二百步为益廉，作带从减益廉开立方法除之，得半径一百二十步，倍之得全径二百四十步。}



\switchcolumn



\problemmodern{乙从南门出发向东行走（步数未知）后停下，甲从城外西北角（乾隅）向南行走六百步，斜向望见乙与圆城恰好成一条直线，又斜向行走一百三十六步接近乙。问圆城的直径是多少？}



\answermodern{圆城直径为二百四十步。}



\shiyimodern{此题运用"明股与明弦"的方法进行测量计算。甲向南行走六百步就是明股（明三角形的直角边），乙向东行走的距离是明勾，斜向行走一百三十六步是明弦。}



\shuyuemodern{将两已知数（南行六百步与斜行一百三十六步）相减，余四百六十四，作为明股与明弦的差。以此差乘以明股，得六万二千六百四十。再以半明股（三百步）乘之，得一千八百七十九万二千，作为开立方的被除数（实）。半明股乘明股得九万一千，与差乘明股之数相减，余二万八千七百三十作为从方（一次项）。南行数加倍得一千二百步作为益廉，使用带从减益廉开立方的方法计算，得到半径一百二十步，加倍得全径二百四十步。}



\end{paracol}



\vspace{0.3cm}



%% --- Volume 5, Section 3, Problem 2 ---

\begin{paracol}{2}



\problemclassical{乙出南门东行不知步数而立，甲从城外西北乾隅南行六百步，斜望乙与城参相直，又斜行一百五十步就乙。问城径。}



\answerclassical{城径二百四十步。}



\shiyiclassical{此以明股重弦立法测望，甲南行明股也，乙东行明勾也，斜行重弦也。}



\shuyueclassical{术与前明股明弦测望问同，带从减益廉开立方法除之，得半径。倍之得全径二百四十步。}



\switchcolumn



\problemmodern{乙从南门出发向东行走（步数未知）后停下，甲从城外西北角（乾隅）向南行走六百步，斜向望见乙与圆城恰好成一条直线，又斜向行走一百五十步接近乙。问圆城的直径是多少？}



\answermodern{圆城直径为二百四十步。}



\shiyimodern{此题运用"明股与重弦"的方法进行测量计算。甲向南行走六百步就是明股，乙向东行走的距离是明勾，斜向行走一百五十步是重弦。}



\shuyuemodern{解法与前一道"明股与明弦"测望题相同，使用带从减益廉开立方的方法计算，得到半径。加倍得全径二百四十步。}



\end{paracol}



\vspace{0.5cm}

\longline

\begin{paracol}{2}

\endclassical{明股与别弦测望类终}

\switchcolumn

\endmodern{明股与别弦测望类（白话）结束}

\end{paracol}



\newpage



%% ===================================================================

%% VOLUME 5 — Section 4: 底股与别弦测望类

%% ===================================================================

\setcounter{probnumvol}{0}



\begin{paracol}{2}



\volheaderclassical{底股与别弦测望类}



\switchcolumn



\volheadermodern{底股与别弦测望类}



\end{paracol}



%% --- Volume 5, Section 4, Problem 1 ---

\begin{paracol}{2}



\problemclassical{乙出东门东行不知步数而立，甲从城外西北乾隅南行六百步，斜望乙与城参相直，又斜行三十四步与乙相会。问城径。}



\answerclassical{城径二百四十步。}



\shiyiclassical{此以底股底弦立法测望，甲南行底股也，乙东行底勾也，斜行底弦也。}



\shuyueclassical{二行相减余五百六十六为底股底弦差，以乘底股，得三十三万九千六百。又以半底股乘之，得一千〇十八万八千为立方实。半底股乘底股得十四万四千，与差乘底股之数相减余十九万四千四百为从方。倍南行得一千二百步为益廉，作带从减益廉开立方法除之，得半径一百二十步，倍之得全径二百四十步。}



\switchcolumn



\problemmodern{乙从东门出发向东行走（步数未知）后停下，甲从城外西北角（乾隅）向南行走六百步，斜向望见乙与圆城恰好成一条直线，又斜向行走三十四步与乙会合。问圆城的直径是多少？}



\answermodern{圆城直径为二百四十步。}



\shiyimodern{此题运用"底股与底弦"的方法进行测量计算。甲向南行走六百步就是底股（底三角形的直角边），乙向东行走的距离是底勾，斜向行走三十四步是底弦。}



\shuyuemodern{将两已知数（南行六百步与斜行三十四步）相减，余五百六十六，作为底股与底弦的差。以此差乘以底股，得三十三万九千六百。再以半底股（三百步）乘之，得一千零一十八万八千，作为开立方的被除数（实）。半底股乘底股得十四万四千，与差乘底股之数相减，余十九万四千四百作为从方（一次项）。南行数加倍得一千二百步作为益廉，使用带从减益廉开立方的方法计算，得到半径一百二十步，加倍得全径二百四十步。}



\end{paracol}



\vspace{0.5cm}

\longline

\begin{paracol}{2}

\endclassical{底股与别弦测望类终}

\switchcolumn

\endmodern{底股与别弦测望类（白话）结束}

\end{paracol}



\newpage



%% ===================================================================

%% VOLUME 5 — Section 5: 大差股与别弦测望类

%% ===================================================================

\setcounter{probnumvol}{0}



\begin{paracol}{2}



\volheaderclassical{大差股与别弦测望类}



\switchcolumn



\volheadermodern{大差股与别弦测望类}



\end{paracol}



%% --- Volume 5, Section 5, Problem 1 ---

\begin{paracol}{2}



\problemclassical{乙出东门东行不知步数而立，甲从城外西北乾隅南行六百步，斜望乙与城参相直，又斜行二百步与乙相会。问城径。}



\answerclassical{城径二百四十步。}



\shiyiclassical{此以大差股大差弦立法测望，甲南行大差股也，乙东行大差勾也，斜行大差弦也。}



\shuyueclassical{二行相减余四百为大差股大差弦差，以乘大差股，得二十四万。又以半大差股乘之，得七百二十万为立方实。半大差股乘大差股得十二万，与差乘大差股之数相减余十二万为从方。倍南行得一千二百步为益廉，作带从减益廉开立方法除之，得半径一百二十步，倍之得全径二百四十步。}



\switchcolumn



\problemmodern{乙从东门出发向东行走（步数未知）后停下，甲从城外西北角（乾隅）向南行走六百步，斜向望见乙与圆城恰好成一条直线，又斜向行走二百步与乙会合。问圆城的直径是多少？}



\answermodern{圆城直径为二百四十步。}



\shiyimodern{此题运用"大差股与大差弦"的方法进行测量计算。甲向南行走六百步就是大差股（大差三角形的直角边），乙向东行走的距离是大差勾，斜向行走二百步是大差弦。}



\shuyuemodern{将两已知数（南行六百步与斜行二百步）相减，余四百，作为大差股与大差弦的差。以此差乘以大差股，得二十四万。再以半大差股（三百步）乘之，得七百二十万，作为开立方的被除数（实）。半大差股乘大差股得十二万，与差乘大差股之数相减，余十二万作为从方（一次项）。南行数加倍得一千二百步作为益廉，使用带从减益廉开立方的方法计算，得到半径一百二十步，加倍得全径二百四十步。}



\end{paracol}



\vspace{0.5cm}

\longline

\begin{paracol}{2}

\endclassical{大差股与别弦测望类终}

\switchcolumn

\endmodern{大差股与别弦测望类（白话）结束}

\end{paracol}



\newpage



%% ===================================================================

%% VOLUME 5 — Section 6: 小差股与别弦测望类

%% ===================================================================

\setcounter{probnumvol}{0}



\begin{paracol}{2}



\volheaderclassical{小差股与别弦测望类}



\switchcolumn



\volheadermodern{小差股与别弦测望类}



\end{paracol}



%% --- Volume 5, Section 6, Problem 1 ---

\begin{paracol}{2}



\problemclassical{乙出南门东行不知步数而立，甲从城外西北乾隅南行六百步，斜望乙与城参相直，又斜行二百步就乙。问城径。}



\answerclassical{城径二百四十步。}



\shiyiclassical{此以小差股小差弦立法测望，甲南行小差股也，乙东行小差勾也，斜行小差弦也。}



\shuyueclassical{术与前大差股大差弦测望问同，带从减益廉开立方法除之，得半径。倍之得全径二百四十步。}



\switchcolumn



\problemmodern{乙从南门出发向东行走（步数未知）后停下，甲从城外西北角（乾隅）向南行走六百步，斜向望见乙与圆城恰好成一条直线，又斜向行走二百步接近乙。问圆城的直径是多少？}



\answermodern{圆城直径为二百四十步。}



\shiyimodern{此题运用"小差股与小差弦"的方法进行测量计算。甲向南行走六百步就是小差股（小差三角形的直角边），乙向东行走的距离是小差勾，斜向行走二百步是小差弦。}



\shuyuemodern{解法与前一道"大差股与大差弦"测望题相同，使用带从减益廉开立方的方法计算，得到半径。加倍得全径二百四十步。}



\end{paracol}



\vspace{0.5cm}

\longline

\begin{paracol}{2}

\endclassical{小差股与别弦测望类终}

\switchcolumn

\endmodern{小差股与别弦测望类（白话）结束}

\end{paracol}



\newpage



%% ===================================================================

%% VOLUME 5 — Section 7: 皇极股与别弦测望类

%% ===================================================================

\setcounter{probnumvol}{0}



\begin{paracol}{2}



\volheaderclassical{皇极股与别弦测望类}



\switchcolumn



\volheadermodern{皇极股与别弦测望类}



\end{paracol}



%% --- Volume 5, Section 7, Problem 1 ---

\begin{paracol}{2}



\problemclassical{乙出东门东行不知步数而立，甲从城外西北乾隅南行六百步，斜望乙与城参相直，又斜行二百八十九步与乙相会。问城径。}



\answerclassical{城径二百四十步。}



\shiyiclassical{此以皇极股皇极弦立法测望，甲南行皇极股也，乙东行皇极勾也，斜行皇极弦也。}



\shuyueclassical{二行相减余三百一十一为皇极股皇极弦差，以乘皇极股，得十八万六千六百。又以半皇极股乘之，得五千五百九十八万为立方实。半皇极股乘皇极股得十八万，与差乘皇极股之数相减余一千六百为从方。倍南行得一千二百步为益廉，作带从减益廉开立方法除之，得半径一百二十步，倍之得全径二百四十步。}



\switchcolumn



\problemmodern{乙从东门出发向东行走（步数未知）后停下，甲从城外西北角（乾隅）向南行走六百步，斜向望见乙与圆城恰好成一条直线，又斜向行走二百八十九步与乙会合。问圆城的直径是多少？}



\answermodern{圆城直径为二百四十步。}



\shiyimodern{此题运用"皇极股与皇极弦"的方法进行测量计算。甲向南行走六百步就是皇极股（皇极三角形的直角边），乙向东行走的距离是皇极勾，斜向行走二百八十九步是皇极弦。}



\shuyuemodern{将两已知数（南行六百步与斜行二百八十九步）相减，余三百一十一，作为皇极股与皇极弦的差。以此差乘以皇极股，得十八万六千六百。再以半皇极股（三百步）乘之，得五千五百九十八万，作为开立方的被除数（实）。半皇极股乘皇极股得十八万，与差乘皇极股之数相减，余一千六百作为从方（一次项）。南行数加倍得一千二百步作为益廉，使用带从减益廉开立方的方法计算，得到半径一百二十步，加倍得全径二百四十步。}



\end{paracol}



\vspace{0.5cm}

\longline

\begin{paracol}{2}

\endclassical{皇极股与别弦测望类终}

\switchcolumn

\endmodern{皇极股与别弦测望类（白话）结束}

\end{paracol}



\newpage



%% ===================================================================

%% VOLUME 5 — Section 8: 通股与通勾测望类

%% ===================================================================

\setcounter{probnumvol}{0}



\begin{paracol}{2}



\volheaderclassical{通股与通勾测望类}



\switchcolumn



\volheadermodern{通股与通勾测望类}



\end{paracol}



%% --- Volume 5, Section 8, Problem 1 ---

\begin{paracol}{2}



\problemclassical{乙出东门东行不知步数而立，甲从城外西北乾隅南行六百步，斜望乙与城参相直。问城径。}



\answerclassical{城径二百四十步。}



\shiyiclassical{此以通股通勾立法测望，甲南行通股也，乙东行通勾也。不用弦而以二行相乘为实，开平方得全径。}



\shuyueclassical{二行相乘得十九万二千为实，以平方开之，得四百三十八点一七，非整数，以二行相乘之数为实，以斜行为法，除之得二百八十二点三五，复以通弦为法除之得全径二百四十步。又以天元术验之，立天元一为半径，以减于半甲南行，得三百为大差，以大差自之得九万为大差幂。置甲南行幂内减大差幂而半之，得十五万七千五百为大弦，内带大差为分母。又以大差乘股六百步，得十八万并入，以乘大弦得同数，与左相消。开立方得一百二十步，即半径也，倍之得全径二百四十步。}



\switchcolumn



\problemmodern{乙从东门出发向东行走（步数未知）后停下，甲从城外西北角（乾隅）向南行走六百步，斜向望见乙与圆城恰好成一条直线。问圆城的直径是多少？}



\answermodern{圆城直径为二百四十步。}



\shiyimodern{此题运用"通股与通勾"的方法进行测量计算。甲向南行走六百步就是通股（大直角边），乙向东行走的距离是通勾（另一直角边）。此题不用弦，而是将两数相乘作为被除数，开平方即可得全径。}



\shuyuemodern{两已知数（南行六百步与东行未知步数）相乘得十九万二千作为实，开平方得四百三十八点一七，不是整数。再以两数相乘之数为实，以斜行为法，相除得二百八十二点三五，再以通弦为法除之得到全径二百四十步。又用天元术验证：设天元一为半径，以半甲南行（三百步）减之，得三百为大差。大差自乘得九万为大差幂。甲南行数自乘内减大差幂，取半得十五万七千五百为大弦。又以大差乘股六百步得十八万并入，与大弦相乘得同数，与左边相消（相消：方程相减消元）。开立方得一百二十步，即为半径，加倍得全径二百四十步。}



\end{paracol}



\vspace{0.5cm}

\longline

\begin{paracol}{2}

\endclassical{通股与通勾测望类终}

\switchcolumn

\endmodern{通股与通勾测望类（白话）结束}

\end{paracol}



\newpage



%% ===================================================================

%% VOLUME 5 — Section 9: 通股与边勾测望类

%% ===================================================================

\setcounter{probnumvol}{0}



\begin{paracol}{2}



\volheaderclassical{通股与边勾测望类}



\switchcolumn



\volheadermodern{通股与边勾测望类}



\end{paracol}



%% --- Volume 5, Section 9, Problem 1 ---

\begin{paracol}{2}



\problemclassical{乙出东门东行不知步数而立，甲从城外西北乾隅南行六百步，斜望乙与城参相直，又斜行二百五十步与乙相会。问城径。}



\answerclassical{城径二百四十步。}



\shiyiclassical{此以通股边勾立法测望，甲南行通股也，乙东行边勾也，斜行边弦也。}



\shuyueclassical{术与前通股与通勾测望问同，带从廉负隅开三乘方法除之，得边勾。以通股边勾求容圆得全径。}



\switchcolumn



\problemmodern{乙从东门出发向东行走（步数未知）后停下，甲从城外西北角（乾隅）向南行走六百步，斜向望见乙与圆城恰好成一条直线，又斜向行走二百五十步与乙会合。问圆城的直径是多少？}



\answermodern{圆城直径为二百四十步。}



\shiyimodern{此题运用"通股与边勾"的方法进行测量计算。甲向南行走六百步就是通股（大直角边），乙向东行走的距离是边勾（边三角形的另一直角边），斜向行走二百五十步是边弦。}



\shuyuemodern{解法与前一道"通股与通勾"测望题相同，使用带从廉负隅开四次方的方法计算，得到边勾。再以通股边勾求容圆直径，得到全径。}



\end{paracol}



\vspace{0.5cm}

\longline

\begin{paracol}{2}

\endclassical{通股与边勾测望类终}

\switchcolumn

\endmodern{通股与边勾测望类（白话）结束}

\end{paracol}



\newpage



%% ===================================================================

%% VOLUME 5 — Section 10: 通股与明勾测望类

%% ===================================================================

\setcounter{probnumvol}{0}



\begin{paracol}{2}



\volheaderclassical{通股与明勾测望类}



\switchcolumn



\volheadermodern{通股与明勾测望类}



\end{paracol}



%% --- Volume 5, Section 10, Problem 1 ---

\begin{paracol}{2}



\problemclassical{乙出南门东行不知步数而立，甲从城外西北乾隅南行六百步，斜望乙与城参相直，又斜行一百三十六步就乙。问城径。}



\answerclassical{城径二百四十步。}



\shiyiclassical{此以通股明勾立法测望，甲南行通股也，乙东行明勾也，斜行明弦也。}



\shuyueclassical{术与前通股边勾测望问同，带从廉负隅开三乘方法除之，得明勾。以通股明勾求容圆得全径。}



\switchcolumn



\problemmodern{乙从南门出发向东行走（步数未知）后停下，甲从城外西北角（乾隅）向南行走六百步，斜向望见乙与圆城恰好成一条直线，又斜向行走一百三十六步接近乙。问圆城的直径是多少？}



\answermodern{圆城直径为二百四十步。}



\shiyimodern{此题运用"通股与明勾"的方法进行测量计算。甲向南行走六百步就是通股（大直角边），乙向东行走的距离是明勾（明三角形的直角边），斜向行走一百三十六步是明弦。}



\shuyuemodern{解法与前一道"通股边勾"测望题相同，使用带从廉负隅开四次方的方法计算，得到明勾。再以通股明勾求容圆直径，得到全径。}



\end{paracol}



\vspace{0.5cm}

\longline

\begin{paracol}{2}

\endclassical{通股与明勾测望类终}

\switchcolumn

\endmodern{通股与明勾测望类（白话）结束}

\end{paracol}



\newpage



%% ===================================================================

%% VOLUME 5 — Section 11: 通股与底勾测望类

%% ===================================================================

\setcounter{probnumvol}{0}



\begin{paracol}{2}



\volheaderclassical{通股与底勾测望类}



\switchcolumn



\volheadermodern{通股与底勾测望类}



\end{paracol}



%% --- Volume 5, Section 11, Problem 1 ---

\begin{paracol}{2}



\problemclassical{乙出东门东行不知步数而立，甲从城外西北乾隅南行六百步，斜望乙与城参相直，又斜行五百〇八步与乙相会。问城径。}



\answerclassical{城径二百四十步。}



\shiyiclassical{此以通股底勾立法测望，甲南行通股也，乙东行底勾也，斜行底弦也。}



\shuyueclassical{术与前通股明勾测望问同，带从廉负隅开三乘方法除之，得底勾。以通股底勾求容圆得全径。}



\switchcolumn



\problemmodern{乙从东门出发向东行走（步数未知）后停下，甲从城外西北角（乾隅）向南行走六百步，斜向望见乙与圆城恰好成一条直线，又斜向行走五百零八步与乙会合。问圆城的直径是多少？}



\answermodern{圆城直径为二百四十步。}



\shiyimodern{此题运用"通股与底勾"的方法进行测量计算。甲向南行走六百步就是通股（大直角边），乙向东行走的距离是底勾（底三角形的直角边），斜向行走五百零八步是底弦。}



\shuyuemodern{解法与前一道"通股明勾"测望题相同，使用带从廉负隅开四次方的方法计算，得到底勾。再以通股底勾求容圆直径，得到全径。}



\end{paracol}



\vspace{0.5cm}

\longline

\begin{paracol}{2}

\endclassical{通股与底勾测望类终}

\switchcolumn

\endmodern{通股与底勾测望类（白话）结束}

\end{paracol}



\newpage



%% ===================================================================

%% VOLUME 5 — Section 12: 通股与皇极勾测望类

%% ===================================================================

\setcounter{probnumvol}{0}



\begin{paracol}{2}



\volheaderclassical{通股与皇极勾测望类}



\switchcolumn



\volheadermodern{通股与皇极勾测望类}



\end{paracol}



%% --- Volume 5, Section 12, Problem 1 ---

\begin{paracol}{2}



\problemclassical{乙出东门东行不知步数而立，甲从城外西北乾隅南行六百步，斜望乙与城参相直，又斜行二百八十九步与乙相会。问城径。}



\answerclassical{城径二百四十步。}



\shiyiclassical{此以通股皇极勾立法测望，甲南行通股也，乙东行皇极勾也，斜行皇极弦也。}



\shuyueclassical{术与前通股底勾测望问同，带从廉负隅开三乘方法除之，得皇极勾。以通股皇极勾求容圆得全径。}



\switchcolumn



\problemmodern{乙从东门出发向东行走（步数未知）后停下，甲从城外西北角（乾隅）向南行走六百步，斜向望见乙与圆城恰好成一条直线，又斜向行走二百八十九步与乙会合。问圆城的直径是多少？}



\answermodern{圆城直径为二百四十步。}



\shiyimodern{此题运用"通股与皇极勾"的方法进行测量计算。甲向南行走六百步就是通股（大直角边），乙向东行走的距离是皇极勾（皇极三角形的直角边），斜向行走二百八十九步是皇极弦。}



\shuyuemodern{解法与前一道"通股底勾"测望题相同，使用带从廉负隅开四次方的方法计算，得到皇极勾。再以通股皇极勾求容圆直径，得到全径。}



\end{paracol}



\vspace{0.5cm}

\longline

\begin{paracol}{2}

\endclassical{通股与皇极勾测望类终}

\switchcolumn

\endmodern{通股与皇极勾测望类（白话）结束}

\end{paracol}



\vspace{0.5cm}

\longline

\begin{center}

{\Large\bfseries\color{volcolor} 测圆海镜分类释术卷五终}

\end{center}



\newpage





%% ===================================================================

%% END MATTER — AFTERWORD

%% ===================================================================

\section*{后记 / 编后说明}

\addcontentsline{toc}{section}{后记 / 编后说明}



\begin{paracol}{2}



\leftheader



《测圆海镜分类释术》为明顾应祥所撰，乃就元李冶《测圆海镜》原书一百七十问，按类分编，释其术法，使学者易于晓悟。顾氏之书，不用天元术，而以传统之带从开方、增乘开方法解之，虽失李冶天元术之精妙，然于普及数学知识不无功绩。



卷四论「通勾与别弦测望」之类，凡十余类，每类设问若干，以释其术。卷五论「通股与别弦测望」之类，亦十余类。二卷共计数十问，皆设圆城一座，甲乙二人各从城门出行，或东或南，遥望相见，以求圆城之径。其法以勾股测圆为本，运用带从开方、增乘开方之术，步骤详明，条理清晰。



顾应祥自序云：「李冶《测圆海镜》以天元术立算，学者苦其难入。余为之分类释术，使问不繁复，术不隐晦，庶几学者由浅入深，渐臻其奥。」观其所释，确有条理井然之致。虽天元术之精妙有所未逮，然于明代数学之普及，实有功焉。



\vspace{0.5cm}

\begin{flushright}

\textit{编者识}

\end{flushright}



\switchcolumn



\rightheader



《测圆海镜分类释术》由明代数学家顾应祥所编撰，就元代李冶《测圆海镜》原书一百七十道问题，按数学方法分类重新编排，解释其算法，使学习者容易理解。顾应祥的这部著作不使用天元术（设未知数列方程的代数方法），而是采用传统的带从开方、增乘开方等数值方法来求解，虽然失去了李冶天元术的精妙之处，但对于普及数学知识仍然有一定的贡献。



卷四讨论"通勾与别弦测望"等类型，共十余类，每类设若干道问题来解释算法。卷五讨论"通股与别弦测望"等类型，也是十余类。两卷共计数十道问题，都假设有一座圆形城池，甲、乙两人分别从城门出发，一人向东一人向南，遥望相见，要求算出圆城的直径。其方法以勾股定理测圆为基础，运用带从开方、增乘开方等数值算法，步骤详细明晰，条理清楚。



顾应祥在自序中说："李冶《测圆海镜》用天元术列算式，学者苦于难以入门。我为此书分类解释算法，使问题不繁复，方法不隐晦，希望学习者能由浅入深，逐渐领会其中的奥妙。"看他所解释的算法，确实条理井然。虽然未能达到天元术的精妙境界，但对于明代数学知识的普及，确实有功劳。



\vspace{0.5cm}

\begin{flushright}

\textit{编者识（白话译）}

\end{flushright}



\end{paracol}



\vspace{1cm}

\longline



\begin{center}

{\small\color{notecolor}

\textbf{术语对照表 / 关键词翻译}\\[0.3em]

\begin{tabular}{ll}

\textcolor{classicalcolor}{文言文} & \textcolor{moderncolor}{白话文解释}\\

\hline

寄术 & 寄术（辅助求解方法）\\

杂术 & 杂术（混合求解方法）\\

相消 & 相消（方程相减消元）\\

互隐 & 互隐（互相消隐）\\

勾 & 直角边（较短者）\\

股 & 直角边（较长者）\\

弦 & 斜边\\

幂 & 平方（自乘）\\

开立方 & 求立方根\\

开三乘方 & 求四次方根\\

带从 & 带有一次项系数的开方\\

廉 & 方程中的中间项系数\\

隅 & 方程中的最高次项系数\\

实 & 被除数 / 常数项\\

法 & 除数 / 已知数\\

\end{tabular}

}

\end{center}



\end{document}







% ============================================================

% Source: translations/zh/chinese/li-ye-ceyuan-haijing-fenlei-shishu-vols6-10_classical-modern_bilingual.tex

% ============================================================



% !TEX program = xelatex

\documentclass[11pt,a4paper]{article}



% Encoding and fonts

\usepackage{fontspec}

\usepackage{xeCJK}

\setCJKmainfont{Microsoft YaHei}



% Math

\usepackage{amsmath,amssymb}



% Layout

\usepackage{geometry}

\geometry{tmargin=2cm,bmargin=2cm,lmargin=1.8cm,rmargin=1.8cm}



% Parallel columns for bilingual text

\usepackage{paracol}

\columnsep=1.5em

\globalcounter{section}

\globalcounter{subsection}



% Typography

\usepackage{titlesec}

\usepackage{setspace}

\usepackage{indentfirst}

\setlength{\parindent}{2em}

\linespread{1.25}



% Colors for labels

\usepackage{xcolor}

\definecolor{classicallabel}{RGB}{128,0,0}

\definecolor{modernlabel}{RGB}{0,64,128}

\definecolor{volumecolor}{RGB}{40,40,100}



% Custom commands

\newcommand{\Volume}[1]{%

  \switchcolumn[0]*

  \vspace{1.5em}

  \begin{center}

    {\color{volumecolor}\rule{0.6\columnwidth}{0.6pt}}\\[0.8em]

    {\Large\bfseries\color{volumecolor} 卷#1}\\[0.3em]

    {\color{volumecolor}\rule{0.6\columnwidth}{0.6pt}}

  \end{center}

  \vspace{0.8em}

  \switchcolumn[1]

  \vspace{1.5em}

  \begin{center}

    {\color{volumecolor}\rule{0.6\columnwidth}{0.6pt}}\\[0.8em]

    {\Large\bfseries\color{volumecolor} 第#1卷}\\[0.3em]

    {\color{volumecolor}\rule{0.6\columnwidth}{0.6pt}}

  \end{center}

  \vspace{0.8em}

}



\newcommand{\Section}[1]{%

  \switchcolumn[0]*

  \subsection*{#1}

  \switchcolumn[1]

  \subsection*{\mdseries\color{modernlabel} #1}

}



\newcommand{\Question}[2]{%

  \switchcolumn[0]*

  \noindent{\bfseries\color{classicallabel} #1}\\#2

  \switchcolumn[1]

  \noindent{\bfseries\color{modernlabel} 【问】}\\#2

}



\newcommand{\Shi}[2]{%

  \switchcolumn[0]*

  \noindent{\bfseries\color{classicallabel} #1}\\#2

  \switchcolumn[1]

  \noindent{\bfseries\color{modernlabel} 【释】}\\#2

}



\newcommand{\Shu}[2]{%

  \switchcolumn[0]*

  \noindent{\bfseries\color{classicallabel} #1}\\#2

  \switchcolumn[1]

  \noindent{\bfseries\color{modernlabel} 【术】}\\#2

}



\newcommand{\You}[2]{%

  \switchcolumn[0]*

  \noindent{\bfseries\color{classicallabel} #1}\\#2

  \switchcolumn[1]

  \noindent{\bfseries\color{modernlabel} 【又】}\\#2

}



\newcommand{\YouShu}[2]{%

  \switchcolumn[0]*

  \noindent{\bfseries\color{classicallabel} #1}\\#2

  \switchcolumn[1]

  \noindent{\bfseries\color{modernlabel} 【又术】}\\#2

}



\newcommand{\Note}[2]{%

  \switchcolumn[0]*

  \noindent{\itshape\small #1}\\#2

  \switchcolumn[1]

  \noindent{\itshape\small\color{gray} #1}\\#2

}



\newcommand{\EndVol}[1]{%

  \switchcolumn[0]*

  \vspace{1em}

  \begin{center}

    \rule{0.5\columnwidth}{0.4pt}\\[0.5em]

    {\bfseries #1}

  \end{center}

  \switchcolumn[1]

  \vspace{1em}

  \begin{center}

    \rule{0.5\columnwidth}{0.4pt}\\[0.5em]

    {\bfseries\color{gray} #1}

  \end{center}

}



% No page numbers for title page

\pagenumbering{gobble}



\begin{document}



% ======================== TITLE PAGE ========================

\begin{center}

\vspace*{2em}

{\Huge\bfseries 测圆海镜分类释术}\\[0.8em]

{\Large\color{volumecolor} 卷六至卷十}\\[1.5em]

{\large 文言文 \quad $\bullet$ \quad 白话文对照本}\\[2em]

{\large 元\quad 李\quad 冶\quad 撰}\\[0.3em]

{\large 明\quad 顾应祥\quad 释术}\\[3em]

{\color{volumecolor}\rule{0.5\textwidth}{1pt}}\\[1.5em]

{\normalsize 左栏：原文（文言文）\quad | \quad 右栏：今译（白话文）}

\end{center}



\vspace{2em}



\begin{minipage}{0.9\textwidth}

\setlength{\parindent}{2em}

《测圆海镜》乃金元之际数学家李冶（1192--1279）所著，为中国古代天元术之杰作。全书十二卷，以勾股容圆为题，设问一百七十道，系统阐述了以天元术（代数方法）求解几何问题之法。明顾应祥为之分类释术，编成《测圆海镜分类释术》。



本文档收录卷六至卷十（最终五卷）之内容。卷六论通弦测望，卷七论大差小差，卷八论皇极太虚，卷九论诸弦差较，卷十论杂法。

\end{minipage}



\vspace{2em}

\begin{center}

{\color{volumecolor}\rule{0.5\textwidth}{1pt}}

\end{center}



\newpage

\tableofcontents

\newpage



% ======================== START BILINGUAL COLUMNS ========================

\pagenumbering{arabic}

\setcounter{page}{1}



\begin{paracol}{2}



% ---- Column headers ----

\switchcolumn[0]

\begin{center}

{\large\bfseries\color{classicallabel} 原文（文言文）}

\end{center}

\vspace{0.5em}



\switchcolumn[1]

\begin{center}

{\large\bfseries\color{modernlabel} 白话译文}

\end{center}

\vspace{0.5em}



\switchcolumn[0]*

\noindent\rule{\columnwidth}{0.3pt}

\switchcolumn[1]

\noindent\rule{\columnwidth}{0.3pt}



%======================================================================

\Volume{六}

%======================================================================



\switchcolumn[0]*

\noindent\textbf{钦定四库全书}



\noindent\textbf{测圆海镜分类释术卷六}\quad 元~李~冶~撰\\

明~顾应祥~释术



\switchcolumn[1]

\noindent\textbf{钦定四库全书}



\noindent\textbf{测圆海镜分类释术·第六卷}\quad 元代~李~冶~撰\\

明代~顾应祥~释术



\vspace{1em}



%----------------------------------------------------------------------

\Section{勾与和测望一}

%----------------------------------------------------------------------



\Question{问一}{%

甲乙俱在城外西北乾隅，甲南行不知步数而立，乙东行三百二十步见之，甲又斜行与相会，计甲直行斜行共一千二百八十步。问城径。

}



\switchcolumn[0]*

\noindent\textbf{释曰}\quad 此通勾与通股弦和测望。乙东行，通勾也；甲直斜共行，通股弦和也。



\switchcolumn[1]

\noindent\textbf{【释】}\quad 这是用通勾与通股弦和来进行测望的问题。乙向东行走，是通勾；甲直行与斜行一共行走的步数，是通股弦和。



\switchcolumn[0]*

\noindent\textbf{术曰}\quad 勾自之，得一十万二千四百。以和除之，得八十为股弦较。以较减和，半之为股。以勾股求容圆术求之，得城径。



\switchcolumn[1]

\noindent\textbf{【术】}\quad 将勾自乘，得到十万二千四百。用这个数除以和，得到八十，这就是股弦差。用和减去这个差，取一半得到股。再用勾股求容圆的公式计算，求得城圆的直径。



\switchcolumn[0]*

\noindent\textbf{又曰}\quad 勾和各自乘，相减为实；倍和除之，得股。相并为实，倍和除之，得弦。



\switchcolumn[1]

\noindent\textbf{【又】}\quad 将勾与和分别自乘，相减作为被除数；用二倍和来除，得到股。将勾与和相加后自乘作为被除数，用二倍和来除，得到弦。



\switchcolumn[0]*

边勾以下，俱以类推。



\switchcolumn[1]

边勾以下的各类问题，都按照这个方法类推。



\vspace{0.5em}



\Question{问二}{%

乙出东门南行，丙出南门东行，各不知步数而立，只云丙行多于乙步。甲从乾隅东行三百二十步，望乙丙与城相参直，计乙丙共行一百二步。问城径。

}



\switchcolumn[0]*

\noindent\textbf{释曰}\quad 此以通勾与明勾、边股和测望。甲东行，通勾也；乙出东门南行为边股，丙出南门东行为明勾，共计一百二步，明勾、边股和也。



\switchcolumn[1]

\noindent\textbf{【释】}\quad 这是用通勾与明勾、边股和来进行测望的问题。甲向东行走，是通勾；乙出东门向南行走是边股，丙出南门向东行走是明勾，共计一百零二步，就是明勾与边股的和。



\switchcolumn[0]*

\noindent\textbf{术曰}\quad 倍共步乘东行算，得二千零八十八万九千六百，为立方实。共步乘东行，加东行算，得一十三万五千零四十，为从方。东行为从廉，五分为隅算。作带从负隅，以廉减从，开立方法除之，得全径。



\switchcolumn[1]

\noindent\textbf{【术】}\quad 将共步的二倍乘以向东行走的数，得到二千零八十八万九千六百，作为立方的被除数。共步乘以向东行走的数，加上东行数的平方，得到十三万五千零四十，作为从方。东行数作为从廉，五分作为隅算。按带从负隅的方法，用廉来减从方，用开立方法除之，得到圆的全径。



\switchcolumn[0]*

带从负隅，以廉减从，半翻法开立方曰：置所得实，以从方约之。初商二百，置一于左上为法，置一乘从廉得六万四千，以减从方，存七万一千零四十为从。置一自之得四万，以隅算五分因之，得二万为隅法。并从，共九万一千零四十为下法，与上法相乘，除实一千八百二十万八千，余实二百六十八万一千六百。从方内再减六万四千，止余七千零四十为从。三因隅法得六万为方法，三因初商得六百为廉法。次商四十，置一于左，次为上法，置一乘从廉得一万二千八百，以减余从，不及减，反减余从七千零四十，余五千七百六十为负从。置一乘廉法，以隅因得一万二千，置一自之，隅因得八百为隅法。并方廉隅，共七万二千八百，减去负从，余六万七千零四十为下法，与上法相乘，除实尽。



\switchcolumn[1]

\noindent\textbf{【带从负隅开立方详解】}\quad 把求得的被除数，用从方来约。初次商得二百，在左上方放一个一作为上法，用一乘以从廉得到六万四千，用来减从方，剩余七万一千零四十作为新的从方。一自乘得四万，用隅算五分乘之，得到二万作为隅法。合并从方与隅法，共九万一千零四十作为下法，与上法相乘，从被除数中减去一千八百二十万八千，剩余二百六十八万一千六百。从方中再减去六万四千，只剩下七千零四十作为从方。将隅法三倍得六万作为方法，将初商三倍得六百作为廉法。第二次商得四十，在左上方放一个一作为上法，用一乘以从廉得到一万二千八百，用来减剩余的从方，不够减，反过来用剩余的从方七千零四十去减，余五千七百六十作为负从。用一乘以廉法，再用隅算乘之得到一万二千，一自乘再用隅算乘之得到八百作为隅法。合并方法、廉法、隅法共七万二千八百，减去负从五千七百六十，余六万七千零四十作为下法，与上法相乘，把剩余被除数刚好除尽。



\switchcolumn[0]*

法已见四卷通勾太虚弦条，因以五分为隅，故重出。



\switchcolumn[1]

这个方法在第四卷通勾太虚弦条下已经出现过，因为这里用五分作为隅算，所以重新列出。



\switchcolumn[0]*

又为带从负隅，以廉添积，开立方法，法见四卷通勾虚弦条下。



\switchcolumn[1]

又可以用带从负隅的方法，以廉来增添积数，开立方法求解，详见第四卷通勾虚弦条下。



\vspace{0.5em}



\Question{问三}{%

乙出东门东行，丙出南门南行，各不知步数而立。甲从乾隅东行三百二十步，望乙丙二人俱与城相参直，计乙丙共行一百五十一步。问城径。

}



\switchcolumn[0]*

\noindent\textbf{释曰}\quad 此以通勾与边勾、明股和，立法测望。甲东行，通勾；乙东行，边勾；丙南行，明股也。



\switchcolumn[1]

\noindent\textbf{【释】}\quad 这是用通勾与边勾、明股和来建立方法进行测望的问题。甲向东行走，是通勾；乙向东行走，是边勾；丙向南行走，是明股。



\switchcolumn[0]*

\noindent\textbf{术曰}\quad 通勾自之，得一十万二千四百，半之得五万一千二百，又自之得二十六亿二千一百四十四万，为三乘方实。以三百六十二乘半通勾算，得一千八百五十三万四千四百，为从方。通勾乘和步，得四万八千三百二十，为从一廉。五之通勾，得一千六百，为从二廉。二分五厘为常法。作带从方廉三乘方法，开之得八十为小差。小差者，通股弦较也。以减通勾，即城径。



\switchcolumn[1]

\noindent\textbf{【术】}\quad 将通勾自乘，得到十万二千四百，取一半得到五万一千二百，再自乘得到二十六亿二千一百四十四万，作为三乘方的被除数。用三百六十二乘以半通勾的数，得到一千八百五十三万四千四百，作为从方。通勾乘以和步，得到四万八千三百二十，作为第一从廉。通勾的五倍得到一千六百，作为第二从廉。二分五厘作为常法。按带从方廉三乘方法开方，得到八十作为小差。小差就是通股弦差。用通勾减去它，就是城圆的直径。



\switchcolumn[0]*

带从方廉负隅，单位开三乘方曰：置所得三乘方实，以廉隅约之。商得八十，置一于左上为法，置一乘从一廉得三百八十六万五千六百，置一自之以乘从二廉得一千零二十四万，置一自乘再得五十一万二千，以二分五厘因之，得十二万八千为隅法。并从方、一廉、二廉、隅法，得三千二百七十六万八千为下法，与上法相乘，除实尽。



\switchcolumn[1]

\noindent\textbf{【带从方廉负隅开三乘方详解】}\quad 把求得的三乘方被除数，用从廉和隅来约。商得八十，在左上方放一个一作为上法，用一乘以第一从廉得到三百八十六万五千六百，一自乘再乘以第二从廉得到一千零二十四万，一自乘三次得到五十一万二千，用二分五厘乘之得到十二万八千作为隅法。合并从方、第一从廉、第二从廉、隅法，共三千二百七十六万八千作为下法，与上法相乘，把被除数刚好除尽。



\vspace{0.5em}



\Question{问四}{%

东门外往南有树，乙出东门往东不知步数而立。甲出北门东行二百步，斜望乙与树，正与城相参直。既而乙复折而斜行至树下，与甲相望，计乙直行斜行共五十步。

}



\switchcolumn[0]*

\noindent\textbf{释曰}\quad 此以底勾与边勾弦和，立法测望。甲出北门东行，底勾也；乙一直一斜，边勾、边弦也。



\switchcolumn[1]

\noindent\textbf{【释】}\quad 这是用底勾与边勾弦和来建立方法进行测望的问题。甲出北门向东行走，是底勾；乙先直走再斜走，是边勾与边弦。



\switchcolumn[0]*

\noindent\textbf{术曰}\quad 底勾与和相减，余一百五十为差。差加底勾，复以差乘之，得数半之，得二万六千二百五十。差自之，得二万二千五百。二数相减，余三千七百五十为实。并勾和，半之，得一百二十五为法。实如法而得一。



\switchcolumn[1]

\noindent\textbf{【术】}\quad 底勾与和相减，剩余一百五十作为差。差加上底勾，再用差来乘，得到的数取一半，得到二万六千二百五十。差自乘，得到二万二千五百。两个数相减，剩余三千七百五十作为被除数。勾与和相加取一半，得到一百二十五作为除数。被除数除以除数得到一个数值。



\vspace{0.5em}



\Question{问五}{%

南门外往东不知步数有树，乙出南门南行不知步数而立。甲出北门东行二百步，见树与乙与城相参直，乙复斜行至树下与甲相望，计乙一直一斜共二百八十八步。问城径。

}



\switchcolumn[0]*

\noindent\textbf{释曰}\quad 此以底勾与明股弦和，立法测望。甲出北门东行，底勾也；乙出南门南行，明股也；斜行，明弦也。



\switchcolumn[1]

\noindent\textbf{【释】}\quad 这是用底勾与明股弦和来建立方法进行测望的问题。甲出北门向东行走，是底勾；乙出南门向南行走，是明股；乙斜着行走，是明弦。



\switchcolumn[0]*

\noindent\textbf{术曰}\quad 勾和相减余，半之得四十四为半差。以减底勾，余一百五十六为泛率。泛率自之，又倍之，得四万八千六百七十二。半差乘和步，得一万二千六百七十二。二数相减，余三万六千为实。半底勾减和步，得一百八十八。倍泛率，得三百一十二。二数相并，得五百为法。实如法而一，得明勾。



\switchcolumn[1]

\noindent\textbf{【术】}\quad 勾与和相减的余数，取一半得到四十四作为半差。用底勾减去半差，余下一百五十六作为泛率。泛率自乘，再乘以二，得到四万八千六百七十二。半差乘以和步，得到一万二千六百七十二。两个数相减，余下三万六千作为被除数。底勾的一半减去和步，得到一百八十八。泛率的二倍得到三百一十二。两个数相加，得到五百作为除数。被除数除以除数，得到明勾。



%----------------------------------------------------------------------

\Section{勾与较测望二}

%----------------------------------------------------------------------



\Question{问六}{%

甲乙俱在城外西北乾隅，甲南行不知步数而立，乙东行三百二十步见之，甲又斜行与乙相会，计甲直行不及斜行八十步。

}



\switchcolumn[0]*

\noindent\textbf{释曰}\quad 此以通勾与股弦较测望。乙东行，通勾也；甲直行不及斜行，股弦较也。



\switchcolumn[1]

\noindent\textbf{【释】}\quad 这是用通勾与股弦差来进行测望的问题。乙向东行走，是通勾；甲直行比斜行少走的步数，就是股弦差。



\switchcolumn[0]*

\noindent\textbf{术曰}\quad 较除勾算，得一千二百八十为股弦和。减较，半之为股；加较，半之为弦。



\switchcolumn[1]

\noindent\textbf{【术】}\quad 用差来除勾的数，得到一千二百八十作为股弦和。减去差取一半就是股；加上差取一半就是弦。



\switchcolumn[0]*

边勾以下，俱即此类推。



\switchcolumn[1]

边勾以下的各类问题，都按照这个方法类推。



%----------------------------------------------------------------------

\Section{股与和测望三}

%----------------------------------------------------------------------



\Question{问七}{%

甲乙二人俱在城外西北乾隅，甲南行六百步而立，乙东行不知步数见之，又斜行与甲相会，计乙直斜共行一千步。问城径。

}



\switchcolumn[0]*

\noindent\textbf{释曰}\quad 此以通股、勾弦和测望。甲南行，通股也；乙直东行与斜行共，勾弦和也。



\switchcolumn[1]

\noindent\textbf{【释】}\quad 这是用通股与勾弦和来进行测望的问题。甲向南行走，是通股；乙直向东行与斜行共走的步数，是勾弦和。



\switchcolumn[0]*

\noindent\textbf{术曰}\quad 股自之，得三十六万。和除之，得三百六十为勾弦较。减和，半之为勾；加和，半之为弦。



\switchcolumn[1]

\noindent\textbf{【术】}\quad 股自乘，得到三十六万。用和来除，得到三百六十作为勾弦差。用和减去差取一半就是勾；用和加上差取一半就是弦。



\switchcolumn[0]*

边股以下，推此。



\switchcolumn[1]

边股以下的各类问题，都按此方法推算。



\vspace{0.5em}



\Question{问八}{%

甲从乾隅南行六百步而立，乙出南门直行，丙出东门直行，三人相望，俱与城相参直。计其行步，则乙与丙共行一百五十一步。

}



\switchcolumn[0]*

\noindent\textbf{释曰}\quad 此以通股、边勾、明股和，立法测望。甲行，通股；乙行，明股；丙行，边勾也。共之和也。



\switchcolumn[1]

\noindent\textbf{【释】}\quad 这是用通股、边勾、明股和来建立方法进行测望的问题。甲行走的是通股；乙行走的是明股；丙行走的是边勾。三人行走数的总和就是和。



\switchcolumn[0]*

\noindent\textbf{术曰}\quad 通股为算，半而自之，得三百二十四亿，为三乘方实。倍和加通股，以乘半通股算，得一亿六千二百三十六万，为从方。通股乘和步，得九万零六百，为从一廉。通股加半股，得九百，为从二廉。二分五厘为隅算。作带从方廉负隅，以二廉减从，翻法开三乘方法除之，得三百六十为股圆差。以减通股，即圆径。



\switchcolumn[1]

\noindent\textbf{【术】}\quad 以通股为基数，取一半后自乘，得到三百二十四亿，作为三乘方的被除数。二倍和加上通股，再乘以半通股的数，得到一亿六千二百三十六万，作为从方。通股乘以和步，得到九万零六百，作为第一从廉。通股加上半股，得到九百，作为第二从廉。二分五厘作为隅算。按带从方廉负隅的方法，用第二从廉减从方，翻法开三乘方来除之，得到三百六十作为股圆差。用通股减去它，就是圆的直径。



\switchcolumn[0]*

带从方廉负隅，以二廉减从，翻法开三乘方曰：置所得三乘方实，以从方廉隅约之。初商三百，置一于左上为法，置一自之以乘二廉，得八千一百万，以减从方，余八千一百三十六万。置一乘从一廉，得二千七百一十八万。置一自乘再乘，以隅算二分五厘因之，得六百七十五万为隅法。并从方、从一廉、隅法，共一亿一千五百二十九万为下法。与上法相乘，除实三百四十五亿八千七百万，实不满法，反减实三百二十四亿，余二十一亿八千七百万为负积。四因隅法，得二千七百万为方法。初商自之，六因又以隅因之，得一十三万五千为上廉。初商四之，隅因之，得三百为下廉。商次位得六十，置一于左，次为上法，倍初商加次商，得六百六十，以乘从二廉，得五十九万四千。又并初次商，得三百六十，因得二亿四千三百八十四万，以减余从，亦不及减，反减从八千一百三十六万，余一亿三千二百四十八万为负从。置一倍初商加次商，得六百六十，以乘从一廉，得五千九百七十九万六千。置一乘上廉，得八十一万。置一自之以乘下廉，得一百零八万。置一自乘再乘，隅因之，得五万四千为隅法。并方法、从一廉、上下廉、隅法，共九千六百零三万。以减负从，余三千六百四十五万，与上次法，除负积二十一亿八千七百万。



\switchcolumn[1]

\noindent\textbf{【带从方廉负隅翻法开三乘方详解】}\quad 把求得的三乘方被除数，用从方、从廉和隅来约。初次商得三百，在左上方放一个一作为上法，一自乘再乘以第二从廉，得到八千一百万，用来减从方，余下八千一百三十六万。用一乘以第一从廉，得到二千七百一十八万。一自乘三次，用隅算二分五厘乘之，得到六百七十五万作为隅法。合并从方、第一从廉、隅法，共一亿一千五百二十九万作为下法。与上法相乘，从被除数中减去三百四十五亿八千七百万，被除数不够减，反过来减被除数三百二十四亿，余下二十一亿八千七百万作为负积。将隅法四倍，得到二千七百万作为方法。初商自乘，六倍再用隅算乘之，得到十三万五千作为上廉。初商四倍再用隅算乘之，得到三百作为下廉。第二次商得六十，在左上方放一个一作为上法，二倍初商加次商得到六百六十，用来乘第二从廉得到五十九万四千。又合并初次商得三百六十，用它去乘余下从方，也不够减，反过来减从方八千一百三十六万，余下一亿三千二百四十八万作为负从。用一倍初商加次商得到六百六十，去乘第一从廉，得到五千九百七十九万六千。用一去乘上廉，得到八十一万。一自乘再乘下廉，得到一百零八万。一自乘三次再用隅算乘之，得到五万四千作为隅法。合并方法、第一从廉、上廉下廉、隅法，共九千六百零三万。用来减负从，余下三千六百四十五万，与上法相乘，用来除负积二十一亿八千七百万。



%----------------------------------------------------------------------

\Section{弦与和测望四}

%----------------------------------------------------------------------



\Question{问九}{%

甲丙二人俱在城外东北艮隅，乙南行过城门而立，甲东行望乙与城相参直而止。丙丁二人俱在城外西南坤隅，丁向东过城门而立，丙向南行望丁及甲乙，悉与城相参直，丙复斜行六百八十步与甲相会。计乙之南与丁之东共三百四十二步。问城径。

}



\switchcolumn[0]*

\noindent\textbf{释曰}\quad 此通弦与大差勾、小差股和，立法测望。乙从艮隅而南过城门而立，山之艮，小差股也；以甲东行为勾，丁从坤隅东行过城门而立，坤之月，大差勾也；以丙南行为股，丙斜行与甲相会，通弦也。乙丁直行共步，大差勾与小差股和也。



\switchcolumn[1]

\noindent\textbf{【释】}\quad 这是用通弦与大差勾、小差股和来建立方法进行测望的问题。乙从艮隅向南走过城门后停下，艮对应山，是小差股；甲向东行走作为勾，丁从坤隅向东行走走过城门后停下，坤对应月，是大差勾；丙向南行走作为股，丙斜着行走与甲相会，是通弦。乙向南和丁向东的直行共步，就是大差勾与小差股的和。



\switchcolumn[0]*

\noindent\textbf{术曰}\quad 斜步、共步相乘，倍之，得四十六万五千一百二十为实。斜步、共步相减，余三百三十八为差。倍斜行加差，共一千六百九十八为从。作带从开平方法除之，得全径。



\switchcolumn[1]

\noindent\textbf{【术】}\quad 斜行步数与共步相乘，再乘以二，得到四十六万五千一百二十作为被除数。斜行步数与共步相减，余下三百三十八作为差。二倍斜行加上差，共一千六百九十八作为从方。按带从开平方法除之，得到圆的全径。



\switchcolumn[0]*

带从开平方法见前。



\switchcolumn[1]

带从开平方法见前面各卷的讲解。



\vspace{0.5em}



\Question{问十}{%

甲出东门东行，乙出南门南行，各不知步数，相望与城相参直，甲复斜行二百八十九步与乙相会。乙直长，甲直短，共计一百五十一步。问城径。

}



\switchcolumn[0]*

\noindent\textbf{释曰}\quad 此以皇极弦、边勾、明股和，立法测望。甲东行为边勾，乙南行为明股，甲之斜行，皇极弦也。



\switchcolumn[1]

\noindent\textbf{【释】}\quad 这是用皇极弦、边勾、明股和来建立方法进行测望的问题。甲向东行走是边勾，乙向南行走是明股，甲的斜行是皇极弦。



\switchcolumn[0]*

\noindent\textbf{术曰}\quad 斜行自之，得八万三千五百二十一为弦算。共步自之，得二万二千八百零一为和算。和算减弦算，余六万零七百二十为实。倍共步减斜行，余十三步为从。作带从开平方法除之，得全径。



\switchcolumn[1]

\noindent\textbf{【术】}\quad 斜行自乘，得到八万三千五百二十一作为弦数。共步自乘，得到二万二千八百零一作为和数。用和数减弦数，余下六万零七百二十作为被除数。二倍共步减去斜行，余下十三步作为从方。按带从开平方法除之，得到圆的全径。



\switchcolumn[0]*

带从开平方法见前。



\switchcolumn[1]

带从开平方法见前面各卷的讲解。



\vspace{0.5em}



\Question{问十一}{%

甲乙二人同出东门，甲东行乙南行；丙丁二人同出南门，丙南行丁东行，各不知步数而立。四人遥相望，悉与城相参直。问其步数，则曰：甲丙共行了一百五十一步，乙丁立处相距一百零二步。问城径。

}



\switchcolumn[0]*

\noindent\textbf{释曰}\quad 此太虚弦与边勾、明股和，立法测望。甲出东门直行为边勾而乙南行为股，丙出南门南行为明股而丁东行为勾。甲丙共步，边勾、明股和也。乙丁相距，太虚弦也。



\switchcolumn[1]

\noindent\textbf{【释】}\quad 这是用太虚弦与边勾、明股和来建立方法进行测望的问题。甲出东门直走是边勾而乙向南走是股，丙出南门向南走是明股而丁向东走是勾。甲与丙共行的步数，就是边勾与明股的和。乙丁立处之间的距离，就是太虚弦。



\switchcolumn[0]*

\noindent\textbf{术曰}\quad 共步、相距步相减，余四十九为差，自之得二千四百零一为差算。共步自之，得二万二千八百零一为和算。差算减和算，余二万零四百为实。倍距步减差，余一百五十五为从。作以从减法，开平方法除之，得全径。



\switchcolumn[1]

\noindent\textbf{【术】}\quad 共步与相距步相减，余下四十九作为差，差自乘得到二千四百零一作为差数。共步自乘，得到二万二千八百零一作为和数。用差数减和数，余下二万零四百作为被除数。二倍距步减去差，余下一百五十五作为从方。按以从减法的开平方法除之，得到圆的全径。



\switchcolumn[0]*

以从减法开平方法见前。



\switchcolumn[1]

以从减法开平方法见前面各卷的讲解。



\switchcolumn[0]*

又为以从添积开平方。



\switchcolumn[1]

又可以用以从添积的开平方法来求解。



\switchcolumn[0]*

其法曰：初商二百，置一于左上为法，置一乘从，得三万二千为益积，添入原积，共五万一千四百为实。置一为隅法，与上法相乘，除实四万，余实一万一千四百。倍隅法，得四百为廉法。次商四十，置一于左上为法，置一乘从方，得六千二百为益实，添入余积，共一万七千六百为实。置一并廉法，共四百四十为下法，与上法相乘，除实尽。



\switchcolumn[1]

\noindent\textbf{【以从添积开平方详解】}\quad 初次商得二百，在左上方放一个一作为上法，用一乘以从方得到三万二千作为益积，添入原来的积数，共五万一千四百作为被除数。放一个一作为隅法，与上法相乘，从被除数中减去四万，余下一万一千四百。将隅法加倍，得到四百作为廉法。第二次商得四十，在左上方放一个一作为上法，用一乘以从方得到六千二百作为益实，添入剩余的积数，共一万七千六百作为被除数。将一与廉法合并，共四百四十作为下法，与上法相乘，把被除数刚好除尽。



\switchcolumn[0]*

后凡言以从添积开平方法，俱仿此。



\switchcolumn[1]

以后凡提到以从添积开平方法的，都仿照这个方法。



\vspace{0.5em}



\Question{问十二}{%

出南门向东有槐树，出东门向南有柳树。丙丁俱出南门，丙直往南，丁往东至槐树下立；甲乙俱出东门，甲直往东，乙往南至柳树下立。四人遥相望见，各不知步数。只云：丙丁共行了二百零七步，甲乙共行了四十六步，其甲丙立处相距二百八十九步。问城径。

}



\switchcolumn[0]*

\noindent\textbf{释曰}\quad 此以皇极弦与明勾股和、边勾股和，立法测望。槐在南门之东，为南之月，明勾也；丁直行往南，为日之南，明股也；共行二百零七，明勾股和也。柳在东门之南，为山之东，边股也；甲直行往东，为东之川，边勾也；共行四十六步，边勾股和也。甲丙立处相距，为日川，皇极弦也。



\switchcolumn[1]

\noindent\textbf{【释】}\quad 这是用皇极弦与明勾股和、边勾股和来建立方法进行测望的问题。槐树在南门东边，对应南之月，是明勾；丁直往南走，对应日之南，是明股；丙丁共走二百零七步，就是明勾股和。柳树在东门南边，对应山之东，是边股；甲直往东走，对应东之川，是边勾；甲乙共走四十六步，就是边勾股和。甲丙立处之间的距离，对应日川，是皇极弦。



\switchcolumn[0]*

\noindent\textbf{术曰}\quad 二和相减余，以减相距余，半之得六十四为平勾。以加二和相减，为平股。相乘为实，平方开之，即半径。



\switchcolumn[1]

\noindent\textbf{【术】}\quad 两个和相减的余数，用相距减去这个余数，取一半得到六十四作为平勾。用这个数加上二和相减的余数，作为平股。两者相乘作为被除数，开平方，就是圆的半径。



\switchcolumn[0]*

\noindent\textbf{又曰}\quad 二和相并，以减相距余，半之得一十八为泛率。加明和为长，加边和为广。长广相乘，得半径算。



\switchcolumn[1]

\noindent\textbf{【又】}\quad 两个和合并，用相距减去合并后的数，取一半得到一十八作为泛率。泛率加上明和作为长，加上边和作为广。长和广相乘，得到半径的数。



\vspace{0.5em}



\Question{问十三}{%

南门之东有槐，东门之南有柳。丙出南门直行，丁出南门东至槐下；甲出东门直行，乙出东门南至柳下。相望俱与城相参直。计丙南丁东共行二百零七步，甲东乙南共行四十六步，其二树相距一百零二步。问城径。

}



\switchcolumn[0]*

\noindent\textbf{释曰}\quad 此与前问同，前以远相距言，此以近相距言。近相距，太虚弦也。以太虚弦与明、叀二和，立法测望。



\switchcolumn[1]

\noindent\textbf{【释】}\quad 这个问题与前一问相同，前一问说的是远距离相距，这一问说的是近距离相距。近距离相距的步数，就是太虚弦。用太虚弦与明勾股和、叀勾股和来建立方法进行测望。



\switchcolumn[0]*

\noindent\textbf{术曰}\quad 叀和乘虚弦，又自之，得二千二百零一万四千八百六十四为平实。并二和自之，得六万四千零九为二和算。边和自之，得二千一百一十六为边和算。明和自之，得四万二千八百四十九为明和算。并明和算、叀和算，以减二和算，余一万九千零四十四为益隅。作负隅开平方法除之，得叀弦。倍弦算与和算相减，开其余，得叀勾股较。加和，半之为股；减和，半之为勾。



\switchcolumn[1]

\noindent\textbf{【术】}\quad 用叀和乘以虚弦，再自乘，得到二千二百零一万四千八百六十四作为平实。将二和合并后自乘，得到六万四千零九作为二和数。边和自乘，得到二千一百一十六作为边和数。明和自乘，得到四万二千八百四十九作为明和数。将明和数与叀和数合并，用二和数来减，余下一万九千零四十四作为益隅。按负隅开平方法除之，得到叀弦。二倍弦数与和数相减，开其余数，得到叀勾股差。加上和取一半是股；减去和取一半是勾。



\switchcolumn[0]*

负隅开平方曰：置所得平实，以益隅约之。初商三十，置一于左上为法，置一乘益隅得五十七万一千三百二十为下法，与上法相乘，除实一千七百一十三万九千六百，余实四百八十七万五千二百六十四。倍下法，得一百一十四万二千六百四十为廉法。约次商得四，置一于左上为法，置一乘益隅得七万六千一百七十六，并入廉法，共一百二十一万八千八百一十六为下法，与上法相乘，除实尽。



\switchcolumn[1]

\noindent\textbf{【负隅开平方详解】}\quad 把求得的平实，用益隅来约。初次商得三十，在左上方放一个一作为上法，用一乘以益隅得到五十七万一千三百二十作为下法，与上法相乘，从被除数中减去一千七百一十三万九千六百，余下四百八十七万五千二百六十四。将下法加倍，得到一百一十四万二千六百四十作为廉法。约得第二次商为四，在左上方放一个一作为上法，用一乘以益隅得到七万六千一百七十六，并入廉法，共一百二十一万八千八百一十六作为下法，与上法相乘，把剩余被除数刚好除尽。



\switchcolumn[0]*

此法已见一卷底勾弦条下，因隅算多故重出。



\switchcolumn[1]

这个方法在第一卷底勾弦条下已经出现过，因为这里的隅算较多所以重新列出。



\switchcolumn[0]*

\noindent\textbf{又曰}\quad 隅算除平实，即得叀弦算。



\switchcolumn[1]

\noindent\textbf{【又】}\quad 用平实除以隅算，就直接得到叀弦的数值。



\switchcolumn[0]*

\noindent\textbf{又曰}\quad 明和乘虚弦，又自之，得四亿四千五百八十万零九百九十六为平实。如前法为负隅，平方开之，得明弦。若以益隅除平实，径得明弦算。



\switchcolumn[1]

\noindent\textbf{【又】}\quad 用明和乘以虚弦，再自乘，得到四亿四千五百八十万零九百九十六作为平实。按照前面的方法设负隅，开平方，得到明弦。如果用益隅来除平实，就直接得到明弦的数值。



\switchcolumn[0]*

\noindent\textbf{又术}\quad 虚弦自之，得一万零四百零四为虚弦算。以叀和乘之，得四十七万八千五百八十四为平实。倍明和，得四百一十四为益隅，开之得叀弦。若以益隅除平实，径得叀弦算。



\switchcolumn[1]

\noindent\textbf{【又术】}\quad 虚弦自乘，得到一万零四百零四作为虚弦数。用叀和乘之，得到四十七万八千五百八十四作为平实。二倍明和得到四百一十四作为益隅，开平方得到叀弦。如果用益隅来除平实，就直接得到叀弦的数值。



\switchcolumn[0]*

虚弦自之，以明和乘之，得二百一十五万三千六百二十八为平实。倍叀和为益隅，开之得明弦。若以益隅除平实，径得明弦算。



\switchcolumn[1]

虚弦自乘，用明和乘之，得到二百一十五万三千六百二十八作为平实。二倍叀和作为益隅，开平方得到明弦。如果用益隅来除平实，就直接得到明弦的数值。



\switchcolumn[0]*

三位负隅开平方曰：置平实四亿四千五百八十万零九百九十六于左，以益隅一万九千零四十四约之。初商一百，置一于左上为法，置一于右下乘益隅，得一百九十万四千四百为下法，与上法相乘，除实一亿九千零四十四万，余实二亿五千五百三十六万零九百九十六。倍下法，得三百八十万八千八百为廉法。次商五十，置一于左上为法，置一乘益隅，得九十五万二千二百为隅法，并廉法，共四百七十六万一千为下法，与上次相乘，除实二亿三千八百零五万，余实一千七百三十一万零九百九十六。倍隅法，得一百九十万四千四百，并入廉法，共五百七十一万三千二百为廉法。约三商得三，置一于左为法，置一右下乘益隅，得五万七千一百三十二为隅法，并入廉法，共五百七十七万零三百三十二为下法，与上法相乘，除实尽。



\switchcolumn[1]

\noindent\textbf{【三位负隅开平方详解】}\quad 把平实四亿四千五百八十万零九百九十六放在左边，用益隅一万九千零四十四来约。初次商得一百，在左上方放一个一作为上法，在右下方放一个一乘益隅，得到一百九十万四千四百作为下法，与上法相乘，从被除数中减去一亿九千零四十四万，余下二亿五千五百三十六万零九百九十六。将下法加倍，得到三百八十万八千八百作为廉法。第二次商得五十，在左上方放一个一作为上法，用一乘益隅得到九十五万二千二百作为隅法，并入廉法，共四百七十六万一千作为下法，与上法相乘，减去二亿三千八百零五万，余下一千七百三十一万零九百九十六。将隅法加倍，得到一百九十万四千四百，并入廉法，共五百七十一万三千二百作为廉法。约得第三次商为三，在左方放一个一作为上法，在右下放一个一乘益隅得到五万七千一百三十二作为隅法，并入廉法，共五百七十七万零三百三十二作为下法，与上法相乘，把被除数刚好除尽。



%----------------------------------------------------------------------

\Section{弦与较测望六}

%----------------------------------------------------------------------



\Question{问十四}{%

甲丙二人俱在城外西北隅起程，丙南行甲东行，各不知步数，隔城相望。既而甲斜行六百八十步与丙相会。问其东行步数，则曰：我少于丙南行二百八十步。问城径。

}



\switchcolumn[0]*

\noindent\textbf{释曰}\quad 此通弦与通勾股较，立法测望。甲东行为勾，丙南行为股，甲少于丙步数，勾股较也，斜行，弦也。



\switchcolumn[1]

\noindent\textbf{【释】}\quad 这是用通弦与通勾股差来建立方法进行测望的问题。甲向东行走作为勾，丙向南行走作为股，甲比丙少走的步数，就是勾股差，斜走的步数是弦。



\switchcolumn[0]*

\noindent\textbf{术曰}\quad 弦自乘，倍之，得九十二万四千八百。较自乘，得七万八千四百。相减，余八十四万六千四百为实。平方开之，得勾股和九百二十。加较，半之为股；减较，半之为勾。



\switchcolumn[1]

\noindent\textbf{【术】}\quad 弦自乘，再乘以二，得到九十二万四千八百。差自乘，得到七万八千四百。两数相减，余下八十四万六千四百作为被除数。开平方，得到勾股和九百二十。加上差取一半是股；减去差取一半是勾。



\switchcolumn[0]*

\noindent\textbf{又曰}\quad 弦较相减，得四百为弦较较；相并，得九百六十为弦较和。弦较较、弦较和相乘，得三十八万四千为实。倍较，得五百六十为从，二为隅算。作以从减法、负隅开平方法除之，得通股。作带从负隅开平方法除之，得通勾。



\switchcolumn[1]

\noindent\textbf{【又】}\quad 弦与差相减，得到四百作为弦较较；两者相加，得到九百六十作为弦较和。弦较较与弦较和相乘，得到三十八万四千作为被除数。二倍差得到五百六十作为从方，二作为隅算。按以从减法、负隅开平方法除之，得到通股。按带从负隅开平方法除之，得到通勾。



\switchcolumn[0]*

带从负隅开平方法，见四卷底勾通弦条。



\switchcolumn[1]

带从负隅开平方法，见第四卷底勾通弦条。



\switchcolumn[0]*

带从负隅以从减隅开平方法，见四卷大差勾黄长弦条下。



\switchcolumn[1]

带从负隅以从减隅开平方法，见第四卷大差勾黄长弦条下。



\switchcolumn[0]*

又为以从添积负隅开平方，以六百乘从益实，倍六百，得一千二百为法。即是边弦以下类推。



\switchcolumn[1]

又可以用以从添积负隅开平方的方法，用六百乘从方作为益实，二倍六百得到一千二百作为法。边弦以下各类问题都可以类推。



\vspace{0.5em}



\Question{问十五}{%

乙出东门南行不知步数而立，甲出西门直往南行，回望乙与城相参直，又斜行五百一十步与乙相会。问乙行步，则曰：少于城径二百一十步。不知城径几何。

}



\switchcolumn[0]*

\noindent\textbf{释曰}\quad 此黄广弦与叀股、黄广勾较，立法测望。乙出东门南行为叀股，城径即黄广勾，少于城径即叀股、黄广勾较也，斜行，黄广弦也。



\switchcolumn[1]

\noindent\textbf{【释】}\quad 这是用黄广弦与叀股、黄广勾差来建立方法进行测望的问题。乙出东门向南行走是叀股，城圆的直径就是黄广勾，比城径少的步数就是叀股与黄广勾的差，斜走的步数是黄广弦。



\switchcolumn[0]*

\noindent\textbf{术曰}\quad 较自之，得四万四千一百为较算，以为实。斜步四之，减二较，余一千六百二十为从，五为隅算。作负隅减从，开平方法除之，得叀股三十，加较为黄广勾，即城径。



\switchcolumn[1]

\noindent\textbf{【术】}\quad 差自乘，得到四万四千一百作为差数，作为被除数。斜步的四倍减去二倍差，余下一千六百二十作为从方，五作为隅算。按负隅减从的开平方法除之，得到叀股三十，加上差就是黄广勾，也就是城圆的直径。



\switchcolumn[0]*

负隅减从开平方法，见二卷通勾叀勾条。



\switchcolumn[1]

负隅减从开平方法，见第二卷通勾叀勾条。



\vspace{0.5em}



\Question{问十六}{%

乙出南门东行不知步数而立，甲出北门直往东行，望乙与城相参直，又斜行二百七十二步与乙相会。问乙东行步，则曰：少于城径一百六十八步。不知城径几何。

}



\switchcolumn[0]*

\noindent\textbf{释曰}\quad 此黄长弦与明勾、黄长股较，立法测望。乙出南门东行为明勾，城径即黄长股，少于城径即明勾、黄长股较也，斜行，黄长弦也。



\switchcolumn[1]

\noindent\textbf{【释】}\quad 这是用黄长弦与明勾、黄长股差来建立方法进行测望的问题。乙出南门向东行走是明勾，城圆的直径就是黄长股，比城径少的步数就是明勾与黄长股的差，斜走的步数是黄长弦。



\switchcolumn[0]*

\noindent\textbf{术曰}\quad 较自之，得二万八千二百二十四为实。四斜行减二较，余七百五十二为从方，五为隅算。作负隅减从，开平方法除之，得明勾七十二，加较为黄长股，即城径。



\switchcolumn[1]

\noindent\textbf{【术】}\quad 差自乘，得到二万八千二百二十四作为被除数。四倍斜行减去二倍差，余下七百五十二作为从方，五作为隅算。按负隅减从的开平方法除之，得到明勾七十二，加上差就是黄长股，也就是城圆的直径。



\switchcolumn[0]*

负隅减从开平方法，见二卷。



\switchcolumn[1]

负隅减从开平方法，见第二卷。



\EndVol{测圆海镜分类释术卷六终}



%======================================================================

\Volume{七}

%======================================================================



\switchcolumn[0]*

\noindent\textbf{测圆海镜分类释术卷七}\quad 元~李~冶~撰\\

明~顾应祥~释术



\switchcolumn[1]

\noindent\textbf{测圆海镜分类释术·第七卷}\quad 元代~李~冶~撰\\

明代~顾应祥~释术



\vspace{1em}



\switchcolumn[0]*

卷七专论大差、小差、中差之测望。大差者，通弦与通股之差；小差者，通弦与通勾之差；中差者，大小差之差也。此三差为测圆术之枢纽，诸题皆以之为本。凡九题，分为三节。



\switchcolumn[1]

第七卷专门论述大差、小差、中差的测望方法。大差，就是通弦与通股的差；小差，就是通弦与通勾的差；中差，就是大差与小差的差。这三种差是测圆术的关键，所有题目都以它们为基础。共有九道题，分为三节。



%----------------------------------------------------------------------

\Section{大差测望一}

%----------------------------------------------------------------------



\Question{问一}{%

甲乙二人俱在城外西北乾隅，甲南行六百步而立，乙东行三百二十步而立。既而甲斜行与乙相会，计甲共行一千二百八十步。以甲斜行与乙东行相较，求大差。

}



\switchcolumn[0]*

\noindent\textbf{释曰}\quad 此以大差立法测望。甲南行六百步为通股，乙东行三百二十步为通勾。甲斜行为通弦。通弦减通股，余为大差。大差勾者，以通勾与大差相减之数也。



\switchcolumn[1]

\noindent\textbf{【释】}\quad 这是用大差来建立方法进行测望的问题。甲向南行走六百步是通股，乙向东行走三百二十步是通勾。甲斜着行走是通弦。通弦减去通股，余下就是大差。大差勾，就是用通勾减去大差后得到的数。



\switchcolumn[0]*

\noindent\textbf{术曰}\quad 通弦自乘，以通股自乘减之，开方得大差勾。以勾股求之，得城径。



\switchcolumn[1]

\noindent\textbf{【术】}\quad 通弦自乘，减去通股自乘，开平方得到大差勾。再用勾股术推算，求得城圆的直径。



\vspace{0.5em}



\Question{问二}{%

甲从乾隅东行三百二十步而立，乙从坤隅南行六百步而立。二人隔城相望，甲复斜行至乙处，计甲共行一千二百八十步。以乙南行与甲斜行相较，求大差。

}



\switchcolumn[0]*

\noindent\textbf{释曰}\quad 此亦以大差立法测望。通股减通弦，余为大差。大差勾即通勾之余数。



\switchcolumn[1]

\noindent\textbf{【释】}\quad 这也是用大差来建立方法进行测望的问题。通股减去通弦，余下就是大差。大差勾就是通勾的余数。



\switchcolumn[0]*

\noindent\textbf{术曰}\quad 通弦自乘，减通股自乘，开方得大差。以大差减通股，得股圆差，即城径。



\switchcolumn[1]

\noindent\textbf{【术】}\quad 通弦自乘，减去通股自乘，开平方得到大差。用大差减通股，得到股圆差，就是城圆的直径。



\vspace{0.5em}



\Question{问三}{%

乙出东门南行，不知步数而立。甲从乾隅东行三百二十步，望乙与城相参直。甲又斜行五百一十步与乙相会。求大差勾、股、弦。

}



\switchcolumn[0]*

\noindent\textbf{释曰}\quad 此以明勾与大差股测望。甲东行，通勾也；乙出东门南行，大差股也；甲斜行，大差弦也。



\switchcolumn[1]

\noindent\textbf{【释】}\quad 这是用明勾与大差股来进行测望的问题。甲向东行走，是通勾；乙出东门向南行走，是大差股；甲斜着行走，是大差弦。



\switchcolumn[0]*

\noindent\textbf{术曰}\quad 通勾自乘，以明勾减之，余为实。以大差弦除之，得大差股。以大差股减通股，余即城径。



\switchcolumn[1]

\noindent\textbf{【术】}\quad 通勾自乘，减去明勾，余下作为被除数。用大差弦来除，得到大差股。用大差股减通股，余下就是城圆的直径。



\vspace{0.5em}



\Question{问四}{%

乙出南门东行，不知步数而立。甲从乾隅南行六百步，望乙与城相参直。甲又斜行至乙处，计甲共行六百八十步。求大差。

}



\switchcolumn[0]*

\noindent\textbf{释曰}\quad 此以大差股与通弦测望。大差股即通股减弦之余。



\switchcolumn[1]

\noindent\textbf{【释】}\quad 这是用大差股与通弦来进行测望的问题。大差股就是通股减去弦后的余数。



\switchcolumn[0]*

\noindent\textbf{术曰}\quad 通弦自乘，减通股自乘，开方得大差勾。以大差勾减通勾，余为城径。



\switchcolumn[1]

\noindent\textbf{【术】}\quad 通弦自乘，减去通股自乘，开平方得到大差勾。用大差勾减通勾，余下就是城圆的直径。



\vspace{0.5em}



\Question{问五}{%

乙出东门东行，不知步数而立。丙出南门南行，不知步数而立。甲从乾隅东行三百二十步，望乙丙与城相参直，计乙丙共行一百五十一步。以乙行与丙行相较，求大差。

}



\switchcolumn[0]*

\noindent\textbf{释曰}\quad 此以明勾、大差股和测望。乙东行为边勾，丙南行为大差股，共行一百五十一步，即大差勾股和也。



\switchcolumn[1]

\noindent\textbf{【释】}\quad 这是用明勾、大差股和来进行测望的问题。乙向东行走是边勾，丙向南行走是大差股，共走一百五十一步，就是大差勾股和。



\switchcolumn[0]*

\noindent\textbf{术曰}\quad 共步自乘，以通勾乘之为实。以通股减共步为从，开平方得大差股。以减通股即城径。



\switchcolumn[1]

\noindent\textbf{【术】}\quad 共步自乘，用通勾来乘作为被除数。用通股减共步作为从方，开平方得到大差股。用通股减去它就是城圆的直径。



\vspace{0.5em}



\Question{问六}{%

甲从坤隅东行过城门而止，乙从艮隅南行过城门而止。二人相望与城相参直，计甲东行与乙南行共三百四十二步。以甲行减乙行，余八十步。求大差。

}



\switchcolumn[0]*

\noindent\textbf{释曰}\quad 此以大差勾股和与较测望。甲东行过城门为大差勾，乙南行过城门为大差股。



\switchcolumn[1]

\noindent\textbf{【释】}\quad 这是用大差勾股和与差来进行测望的问题。甲向东行走走过城门是大差勾，乙向南行走走过城门是大差股。



\switchcolumn[0]*

\noindent\textbf{术曰}\quad 和自乘，减较自乘，以四因之为实。倍和为从，开平方得大差弦。以弦减通弦，余即城径。



\switchcolumn[1]

\noindent\textbf{【术】}\quad 和自乘，减去差自乘，再乘以四作为被除数。二倍和作为从方，开平方得到大差弦。用通弦减去它，余下就是城圆的直径。



\vspace{0.5em}



\Question{问七}{%

甲乙二人俱出东门，甲东行乙南行。丙丁二人俱出南门，丙南行丁东行。四人相望，各与城相参直。计甲丙共行二百步，乙丁相距二百八十九步。求大差。

}



\switchcolumn[0]*

\noindent\textbf{释曰}\quad 此以太虚弦与边勾明股和测望。乙丁相距即太虚弦，大差之属也。



\switchcolumn[1]

\noindent\textbf{【释】}\quad 这是用太虚弦与边勾明股和来进行测望的问题。乙丁相距的距离就是太虚弦，属于大差这一类。



\switchcolumn[0]*

\noindent\textbf{术曰}\quad 以太虚弦减皇极弦，余为大差弦。以勾股术推之，得城径。



\switchcolumn[1]

\noindent\textbf{【术】}\quad 用太虚弦减皇极弦，余下就是大差弦。用勾股术推算，得到城圆的直径。



\vspace{0.5em}



\Question{问八}{%

甲从乾隅东行三百二十步，南门外有树，乙出南门南行至树下。甲斜望乙与树及城相参直，计甲斜行二百八十九步。求大差股。

}



\switchcolumn[0]*

\noindent\textbf{释曰}\quad 此以明勾与大差弦测望。甲东行为明勾，斜行为大差弦，乙南行为大差股。



\switchcolumn[1]

\noindent\textbf{【释】}\quad 这是用明勾与大差弦来进行测望的问题。甲向东行走是明勾，斜着行走是大差弦，乙向南行走是大差股。



\switchcolumn[0]*

\noindent\textbf{术曰}\quad 明勾自乘，减大差弦自乘，开方得大差股。以减通股，得城径。



\switchcolumn[1]

\noindent\textbf{【术】}\quad 明勾自乘，减去大差弦自乘，开平方得到大差股。用通股减去它，得到城圆的直径。



\vspace{0.5em}



\Question{问九}{%

乙出东门东行二百步而立，丙出南门南行一百五十一步而立。甲从乾隅东行三百二十步，望乙丙与城相参直。求大差勾股较。

}



\switchcolumn[0]*

\noindent\textbf{释曰}\quad 此以边勾、明股与大差测望。大差勾股较即通股弦较与通勾弦较之差。



\switchcolumn[1]

\noindent\textbf{【释】}\quad 这是用边勾、明股与大差来进行测望的问题。大差勾股差就是通股弦差与通勾弦差的差。



\switchcolumn[0]*

\noindent\textbf{术曰}\quad 边勾自乘，明股自乘，相减为实。以边勾明股相乘为法，除之得大差勾股较。以加减术推之，得城径。



\switchcolumn[1]

\noindent\textbf{【术】}\quad 边勾自乘，明股自乘，相减作为被除数。用边勾与明股相乘作为除数，除之得到大差勾股差。再用加减术推算，得到城圆的直径。



%----------------------------------------------------------------------

\Section{小差测望二}

%----------------------------------------------------------------------



\Question{问十}{%

甲乙二人俱在城外西北乾隅，甲南行六百步而立，乙东行三百二十步而立。甲斜行与乙相会，计甲共行一千二百八十步。以通弦减通勾，求小差。

}



\switchcolumn[0]*

\noindent\textbf{释曰}\quad 此以小差立法测望。通弦减通勾之余，为小差。小差股者，通股之减数也。



\switchcolumn[1]

\noindent\textbf{【释】}\quad 这是用小差来建立方法进行测望的问题。通弦减去通勾的余数，就是小差。小差股，就是通股的减数。



\switchcolumn[0]*

\noindent\textbf{术曰}\quad 通弦自乘，减通勾自乘，开方得小差股。以小差股减通股，得城径。



\switchcolumn[1]

\noindent\textbf{【术】}\quad 通弦自乘，减去通勾自乘，开平方得到小差股。用小差股减通股，得到城圆的直径。



\vspace{0.5em}



\Question{问十一}{%

乙出南门东行，不知步数而立。甲从乾隅南行六百步，望乙与城相参直。甲又斜行二百七十二步与乙相会。求小差。

}



\switchcolumn[0]*

\noindent\textbf{释曰}\quad 此以明勾与小差弦测望。乙出南门东行为明勾，甲斜行为小差弦，甲南行为小差股。



\switchcolumn[1]

\noindent\textbf{【释】}\quad 这是用明勾与小差弦来进行测望的问题。乙出南门向东行走是明勾，甲斜着行走是小差弦，甲向南行走是小差股。



\switchcolumn[0]*

\noindent\textbf{术曰}\quad 明勾自乘，减小差弦自乘，开方得小差股。以小差股减通股，得城径。



\switchcolumn[1]

\noindent\textbf{【术】}\quad 明勾自乘，减去小差弦自乘，开平方得到小差股。用小差股减通股，得到城圆的直径。



\vspace{0.5em}



\Question{问十二}{%

甲从艮隅南行过城门而止，乙从坤隅东行过城门而止。二人相望与城相参直，计甲南行三百六十步，乙东行一百六十步。求小差。

}



\switchcolumn[0]*

\noindent\textbf{释曰}\quad 此以小差勾股直接测望。甲南行过城门为小差股，乙东行过城门为小差勾。



\switchcolumn[1]

\noindent\textbf{【释】}\quad 这是用小差勾股直接进行测望的问题。甲向南行走走过城门是小差股，乙向东行走走过城门是小差勾。



\switchcolumn[0]*

\noindent\textbf{术曰}\quad 小差勾股各自乘，相并开方得小差弦。以通弦减之，余即城径。



\switchcolumn[1]

\noindent\textbf{【术】}\quad 小差勾与小差股各自乘，相加后开平方得到小差弦。用通弦减去它，余下就是城圆的直径。



\vspace{0.5em}



\Question{问十三}{%

甲乙二人同出东门，甲东行二百步，乙南行一百五十一步而立。二人相望与城相参直。求小差勾股较。

}



\switchcolumn[0]*

\noindent\textbf{释曰}\quad 此以小差勾股与通弦测望。甲东行为小差勾，乙南行为小差股。



\switchcolumn[1]

\noindent\textbf{【释】}\quad 这是用小差勾股与通弦来进行测望的问题。甲向东行走是小差勾，乙向南行走是小差股。



\switchcolumn[0]*

\noindent\textbf{术曰}\quad 小差勾股各自乘相并开方得小差弦，以通弦减之得小差。小差减通勾得城径。



\switchcolumn[1]

\noindent\textbf{【术】}\quad 小差勾与小差股各自乘，相加后开平方得到小差弦，用通弦减去它得到小差。小差减通勾得到城圆的直径。



%----------------------------------------------------------------------

\Section{中差测望三}

%----------------------------------------------------------------------



\Question{问十四}{%

甲乙二人俱在城外西北乾隅，甲南行六百步，乙东行三百二十步。各不知余数。以通股减通弦为大差，以通勾减通弦为小差。大小差相减为中差。求中差。

}



\switchcolumn[0]*

\noindent\textbf{释曰}\quad 此以中差立法测望。中差者，大差与小差之差，即通股与通勾之差也。



\switchcolumn[1]

\noindent\textbf{【释】}\quad 这是用中差来建立方法进行测望的问题。中差，就是大差与小差的差，也就是通股与通勾的差。



\switchcolumn[0]*

\noindent\textbf{术曰}\quad 通股自乘，减通勾自乘，开方得勾股和。以和减通弦，余为中差。以中差加减术推之，得城径。



\switchcolumn[1]

\noindent\textbf{【术】}\quad 通股自乘，减去通勾自乘，开平方得到勾股和。用和减通弦，余下就是中差。用中差加减术推算，得到城圆的直径。



\vspace{0.5em}



\Question{问十五}{%

乙出东门南行，不知步数而立。甲从乾隅东行三百二十步，望乙与城相参直。计甲斜行五百一十步。以乙行与城径相较，求中差。

}



\switchcolumn[0]*

\noindent\textbf{释曰}\quad 此以明勾与中差测望。明勾即通勾减边勾之余，中差即大差减小差。



\switchcolumn[1]

\noindent\textbf{【释】}\quad 这是用明勾与中差来进行测望的问题。明勾就是通勾减去边勾的余数，中差就是大差减去小差。



\switchcolumn[0]*

\noindent\textbf{术曰}\quad 明勾自乘为实，以通弦约之，得中差股。以减通股得城径。



\switchcolumn[1]

\noindent\textbf{【术】}\quad 明勾自乘作为被除数，用通弦来约，得到中差股。用通股减去它得到城圆的直径。



\vspace{0.5em}



\Question{问十六}{%

甲从乾隅南行六百步，乙从乾隅东行三百二十步。以甲斜行与乙东行相较，求中差勾股较。

}



\switchcolumn[0]*

\noindent\textbf{释曰}\quad 此以通勾股较测望。中差勾股较即通股减通勾之数。



\switchcolumn[1]

\noindent\textbf{【释】}\quad 这是用通勾股差来进行测望的问题。中差勾股差就是通股减去通勾的数。



\switchcolumn[0]*

\noindent\textbf{术曰}\quad 通股减通勾，得三百八十步为中差。以勾股术推之，得城径。



\switchcolumn[1]

\noindent\textbf{【术】}\quad 通股减去通勾，得到三百八十步作为中差。用勾股术推算，得到城圆的直径。



\vspace{0.5em}



\Question{问十七}{%

乙出南门东行七十二步而立，丙出东门南行三十步而立。甲从乾隅东行三百二十步，望乙丙与城相参直。求中差弦。

}



\switchcolumn[0]*

\noindent\textbf{释曰}\quad 此以明勾、叀股与中差测望。乙东行为明勾，丙南行为叀股。



\switchcolumn[1]

\noindent\textbf{【释】}\quad 这是用明勾、叀股与中差来进行测望的问题。乙向东行走是明勾，丙向南行走是叀股。



\switchcolumn[0]*

\noindent\textbf{术曰}\quad 明勾、叀股各自乘相并开方得中差弦。以通弦减之，余为城径。



\switchcolumn[1]

\noindent\textbf{【术】}\quad 明勾与叀股各自乘，相加后开平方得到中差弦。用通弦减去它，余下就是城圆的直径。



\vspace{0.5em}



\Question{问十八}{%

甲从坤隅东行过城门，乙从艮隅南行过城门。二人相望与城相参直，计甲东行三百二十步，乙南行六百步。以甲行减乙行，求中差。

}



\switchcolumn[0]*

\noindent\textbf{释曰}\quad 此以大差勾股与小差勾股测望中差。中差即大小差之差。



\switchcolumn[1]

\noindent\textbf{【释】}\quad 这是用大差勾股与小差勾股来测望中差的问题。中差就是大小差的差。



\switchcolumn[0]*

\noindent\textbf{术曰}\quad 大差股减小差股，得中差。以减通弦，余为城径。



\switchcolumn[1]

\noindent\textbf{【术】}\quad 大差股减去小差股，得到中差。用通弦减去它，余下就是城圆的直径。



\EndVol{测圆海镜分类释术卷七终}



%======================================================================

\Volume{八}

%======================================================================



\switchcolumn[0]*

\noindent\textbf{测圆海镜分类释术卷八}\quad 元~李~冶~撰\\

明~顾应祥~释术



\switchcolumn[1]

\noindent\textbf{测圆海镜分类释术·第八卷}\quad 元代~李~冶~撰\\

明代~顾应祥~释术



\vspace{1em}



\switchcolumn[0]*

卷八专论皇极、太虚、明、叀四种之测望，以及杂题之类。皇极者，通勾股之中；太虚者，虚勾股之极；明者，明勾股之谓；叀者，叀勾股之称。凡三十四题，分为五节。



\switchcolumn[1]

第八卷专门论述皇极、太虚、明、叀四种测望方法，以及杂题之类。皇极，就是通勾股的中间值；太虚，就是虚勾股的极限值；明，就是明勾股的称谓；叀，就是叀勾股的称呼。共有三十四道题，分为五节。



%----------------------------------------------------------------------

\Section{皇极测望一}

%----------------------------------------------------------------------



\Question{问一}{%

甲乙二人俱在城外西北乾隅，甲南行六百步而立，乙东行三百二十步而立。二人隔城相望，各与城相参直。以勾股之中数，求皇极弦。

}



\switchcolumn[0]*

\noindent\textbf{释曰}\quad 此以皇极弦立法测望。皇极者，勾股之中极也。皇极弦即通弦之半，为勾股容圆之枢纽。



\switchcolumn[1]

\noindent\textbf{【释】}\quad 这是用皇极弦来建立方法进行测望的问题。皇极，就是勾股之中极。皇极弦就是通弦的一半，是勾股容圆的关键。



\switchcolumn[0]*

\noindent\textbf{术曰}\quad 通勾股各自乘相并，开方得通弦。半之为皇极弦。以皇极弦减通弦，余即城径。



\switchcolumn[1]

\noindent\textbf{【术】}\quad 通勾与通股各自乘，相加后开平方得到通弦。取一半作为皇极弦。用通弦减去皇极弦，余下就是城圆的直径。



\vspace{0.5em}



\Question{问二}{%

乙出东门东行，不知步数而立。丙出南门南行，不知步数而立。甲从乾隅东行三百二十步，望乙丙与城相参直，计乙丙共行一百五十一步。以边勾明股求皇极。

}



\switchcolumn[0]*

\noindent\textbf{释曰}\quad 此以边勾明股和测望皇极。乙东行为边勾，丙南行为明股，共行一百五十一步，即皇极勾股和也。



\switchcolumn[1]

\noindent\textbf{【释】}\quad 这是用边勾明股和来测望皇极的问题。乙向东行走是边勾，丙向南行走是明股，共走一百五十一步，就是皇极勾股和。



\switchcolumn[0]*

\noindent\textbf{术曰}\quad 边勾明股各自乘相并开方得皇极弦。以通弦减之余为城径。



\switchcolumn[1]

\noindent\textbf{【术】}\quad 边勾与明股各自乘，相加后开平方得到皇极弦。用通弦减去其余数就是城圆的直径。



\vspace{0.5em}



\Question{问三}{%

甲从乾隅东行三百二十步，乙出东门南行一百五十一步。甲望乙与城相参直，又斜行二百八十九步与乙相会。求皇极勾股较。

}



\switchcolumn[0]*

\noindent\textbf{释曰}\quad 此以皇极弦与边勾明股和测望。甲斜行二百八十九步为皇极弦，乙丙共行一百五十一步为边勾明股和。



\switchcolumn[1]

\noindent\textbf{【释】}\quad 这是用皇极弦与边勾明股和来进行测望的问题。甲斜着行走二百八十九步是皇极弦，乙丙共走一百五十一步是边勾明股和。



\switchcolumn[0]*

\noindent\textbf{术曰}\quad 皇极弦自乘，减边勾明股和自乘，开方得皇极勾股较。以加减术推之，得城径。



\switchcolumn[1]

\noindent\textbf{【术】}\quad 皇极弦自乘，减去边勾明股和自乘，开平方得到皇极勾股差。用加减术推算，得到城圆的直径。



\vspace{0.5em}



\Question{问四}{%

甲从艮隅南行过城门而止，乙从坤隅东行过城门而止。二人相望与城相参直。以甲行乙行求皇极弦。

}



\switchcolumn[0]*

\noindent\textbf{释曰}\quad 此以大差勾小差股测望皇极。甲南行过城门为小差股，乙东行过城门为大差勾。



\switchcolumn[1]

\noindent\textbf{【释】}\quad 这是用大差勾小差股来测望皇极的问题。甲向南行走走过城门是小差股，乙向东行走走过城门是大差勾。



\switchcolumn[0]*

\noindent\textbf{术曰}\quad 大差勾小差股各自乘相并开方得皇极弦。以通弦减之得城径。



\switchcolumn[1]

\noindent\textbf{【术】}\quad 大差勾与小差股各自乘，相加后开平方得到皇极弦。用通弦减去它得到城圆的直径。



\vspace{0.5em}



\Question{问五}{%

甲出东门东行二百步，乙出南门南行一百五十一步。二人相望与城相参直，计二人相距二百八十九步。以相距与各行相较，求皇极。

}



\switchcolumn[0]*

\noindent\textbf{释曰}\quad 此以皇极勾股弦直接测望。甲东行为皇极勾，乙南行为皇极股，相距为皇极弦。



\switchcolumn[1]

\noindent\textbf{【释】}\quad 这是用皇极勾股弦直接进行测望的问题。甲向东行走是皇极勾，乙向南行走是皇极股，两人相距是皇极弦。



\switchcolumn[0]*

\noindent\textbf{术曰}\quad 皇极勾股各自乘相并开方得皇极弦。以勾股术求之，得城径。



\switchcolumn[1]

\noindent\textbf{【术】}\quad 皇极勾与皇极股各自乘，相加后开平方得到皇极弦。用勾股术计算，得到城圆的直径。



%----------------------------------------------------------------------

\Section{太虚测望二}

%----------------------------------------------------------------------



\Question{问六}{%

甲乙二人同出东门，甲东行乙南行。丙丁二人同出南门，丙南行丁东行。四人相望，悉与城相参直。计甲丙共行一百五十一步，乙丁相距一百零二步。以太虚弦测望。

}



\switchcolumn[0]*

\noindent\textbf{释曰}\quad 此以太虚弦立法测望。乙丁相距一百零二步，即太虚弦也。太虚者，虚勾股之极，为最小之勾股形。



\switchcolumn[1]

\noindent\textbf{【释】}\quad 这是用太虚弦来建立方法进行测望的问题。乙丁相距一百零二步，就是太虚弦。太虚，就是虚勾股的极限，是最小的勾股形。



\switchcolumn[0]*

\noindent\textbf{术曰}\quad 太虚弦自乘为实，以通弦约之，得太虚勾。以太虚勾减通勾，得城径。



\switchcolumn[1]

\noindent\textbf{【术】}\quad 太虚弦自乘作为被除数，用通弦来约，得到太虚勾。用太虚勾减通勾，得到城圆的直径。



\vspace{0.5em}



\Question{问七}{%

乙出东门东行一百步，丙出南门南行五十一步而立。甲从乾隅东行三百二十步，望乙丙与城相参直。以太虚勾股测望。

}



\switchcolumn[0]*

\noindent\textbf{释曰}\quad 此以太虚勾股直接测望。乙东行为太虚勾，丙南行为太虚股。



\switchcolumn[1]

\noindent\textbf{【释】}\quad 这是用太虚勾股直接进行测望的问题。乙向东行走是太虚勾，丙向南行走是太虚股。



\switchcolumn[0]*

\noindent\textbf{术曰}\quad 太虚勾股各自乘相并开方得太虚弦。以通弦减之得城径。



\switchcolumn[1]

\noindent\textbf{【术】}\quad 太虚勾与太虚股各自乘，相加后开平方得到太虚弦。用通弦减去它得到城圆的直径。



\vspace{0.5em}



\Question{问八}{%

甲从乾隅东行三百二十步，南门外有槐树。乙出南门东至槐树下。甲斜望乙与槐与城相参直，计甲斜行一百零二步。求太虚股。

}



\switchcolumn[0]*

\noindent\textbf{释曰}\quad 此以太虚弦与明勾测望。甲斜行一百零二步为太虚弦，乙东至槐树下为明勾。



\switchcolumn[1]

\noindent\textbf{【释】}\quad 这是用太虚弦与明勾来进行测望的问题。甲斜着行走一百零二步是太虚弦，乙向东走到槐树下是明勾。



\switchcolumn[0]*

\noindent\textbf{术曰}\quad 太虚弦自乘，减明勾自乘，开方得太虚股。以太虚股减通股得城径。



\switchcolumn[1]

\noindent\textbf{【术】}\quad 太虚弦自乘，减去明勾自乘，开平方得到太虚股。用太虚股减通股得到城圆的直径。



\vspace{0.5em}



\Question{问九}{%

甲从艮隅南行过城门三百六十步而立，乙从坤隅东行过城门一百六十步而立。二人相望与城相参直。以太虚测望。

}



\switchcolumn[0]*

\noindent\textbf{释曰}\quad 此以小差股大差勾测望太虚。小差股大差勾即太虚勾股之属。



\switchcolumn[1]

\noindent\textbf{【释】}\quad 这是用小差股大差勾来测望太虚的问题。小差股大差勾就是太虚勾股这一类。



\switchcolumn[0]*

\noindent\textbf{术曰}\quad 小差股大差勾各自乘相并开方得太虚弦。以通弦减之得城径。



\switchcolumn[1]

\noindent\textbf{【术】}\quad 小差股与大差勾各自乘，相加后开平方得到太虚弦。用通弦减去它得到城圆的直径。



\vspace{0.5em}



\Question{问十}{%

乙出东门东行七十二步，丙出南门南行三十步而立。二人相望与城相参直。以太虚勾股测望。

}



\switchcolumn[0]*

\noindent\textbf{释曰}\quad 此以太虚勾股直接测望。乙东行为太虚勾，丙南行为太虚股。



\switchcolumn[1]

\noindent\textbf{【释】}\quad 这是用太虚勾股直接进行测望的问题。乙向东行走是太虚勾，丙向南行走是太虚股。



\switchcolumn[0]*

\noindent\textbf{术曰}\quad 太虚勾股各自乘相并开方得太虚弦。以勾股术求之得城径。



\switchcolumn[1]

\noindent\textbf{【术】}\quad 太虚勾与太虚股各自乘，相加后开平方得到太虚弦。用勾股术计算得到城圆的直径。



\vspace{0.5em}



\Question{问十一}{%

甲出东门东行二百步，乙出南门南行一百五十一步。甲望乙与城相参直，又斜行二百八十九步与乙相会。以太虚弦与皇极弦相较。

}



\switchcolumn[0]*

\noindent\textbf{释曰}\quad 此以太虚弦与皇极弦测望。甲斜行二百八十九步为皇极弦，太虚弦为其减数。



\switchcolumn[1]

\noindent\textbf{【释】}\quad 这是用太虚弦与皇极弦来进行测望的问题。甲斜着行走二百八十九步是皇极弦，太虚弦是比它小的数。



\switchcolumn[0]*

\noindent\textbf{术曰}\quad 皇极弦减太虚弦，余为明弦。以明弦减通弦得城径。



\switchcolumn[1]

\noindent\textbf{【术】}\quad 皇极弦减去太虚弦，余下就是明弦。用明弦减通弦得到城圆的直径。



\vspace{0.5em}



\Question{问十二}{%

甲从乾隅东行三百二十步，乙出东门南行三十步。甲望乙与城相参直。求太虚勾股较。

}



\switchcolumn[0]*

\noindent\textbf{释曰}\quad 此以通勾与太虚股测望。乙出东门南行为太虚股。



\switchcolumn[1]

\noindent\textbf{【释】}\quad 这是用通勾与太虚股来进行测望的问题。乙出东门向南行走是太虚股。



\switchcolumn[0]*

\noindent\textbf{术曰}\quad 太虚股自乘为实，以通勾约之得太虚勾。以减通勾得城径。



\switchcolumn[1]

\noindent\textbf{【术】}\quad 太虚股自乘作为被除数，用通勾来约得到太虚勾。用通勾减去它得到城圆的直径。



\vspace{0.5em}



\Question{问十三}{%

乙出南门东行七十二步，丙出东门南行三十步。甲从乾隅南行六百步，望乙丙与城相参直。以太虚测望。

}



\switchcolumn[0]*

\noindent\textbf{释曰}\quad 此以太虚勾股与通股测望。乙东行为太虚勾，丙南行为太虚股。



\switchcolumn[1]

\noindent\textbf{【释】}\quad 这是用太虚勾股与通股来进行测望的问题。乙向东行走是太虚勾，丙向南行走是太虚股。



\switchcolumn[0]*

\noindent\textbf{术曰}\quad 太虚勾股各自乘相并开方得太虚弦。以通弦减之得城径。



\switchcolumn[1]

\noindent\textbf{【术】}\quad 太虚勾与太虚股各自乘，相加后开平方得到太虚弦。用通弦减去它得到城圆的直径。



\vspace{0.5em}



\Question{问十四}{%

甲乙二人同出东门，甲东行乙南行。丙丁二人同出南门，丙南行丁东行。四人遥相望，计甲丙共行二百步，乙丁相距一百零二步。以太虚弦与皇极弦测望。

}



\switchcolumn[0]*

\noindent\textbf{释曰}\quad 此以太虚皇极二弦测望。乙丁相距为太虚弦，甲丙共行为皇极勾股和。



\switchcolumn[1]

\noindent\textbf{【释】}\quad 这是用太虚皇极二弦来进行测望的问题。乙丁相距是太虚弦，甲丙共行是皇极勾股和。



\switchcolumn[0]*

\noindent\textbf{术曰}\quad 太虚弦自乘，以皇极弦约之，得太虚勾。以减通勾得城径。



\switchcolumn[1]

\noindent\textbf{【术】}\quad 太虚弦自乘，用皇极弦来约，得到太虚勾。用通勾减去它得到城圆的直径。



%----------------------------------------------------------------------

\Section{明测望三}

%----------------------------------------------------------------------



\Question{问十五}{%

乙出南门东行七十二步而立。甲从乾隅南行六百步，望乙与城相参直。以明勾测望。

}



\switchcolumn[0]*

\noindent\textbf{释曰}\quad 此以明勾立法测望。乙出南门东行七十二步为明勾。



\switchcolumn[1]

\noindent\textbf{【释】}\quad 这是用明勾来建立方法进行测望的问题。乙出南门向东行走七十二步是明勾。



\switchcolumn[0]*

\noindent\textbf{术曰}\quad 明勾自乘为实，以通弦约之，得明股。以明股减通股得城径。



\switchcolumn[1]

\noindent\textbf{【术】}\quad 明勾自乘作为被除数，用通弦来约，得到明股。用明股减通股得到城圆的直径。



\vspace{0.5em}



\Question{问十六}{%

丙出南门南行五十一步而立。甲从乾隅东行三百二十步，望丙与城相参直。以明股测望。

}



\switchcolumn[0]*

\noindent\textbf{释曰}\quad 此以明股立法测望。丙出南门南行为明股。



\switchcolumn[1]

\noindent\textbf{【释】}\quad 这是用明股来建立方法进行测望的问题。丙出南门向南行走是明股。



\switchcolumn[0]*

\noindent\textbf{术曰}\quad 明股自乘为实，以通弦约之，得明勾。以明勾减通勾得城径。



\switchcolumn[1]

\noindent\textbf{【术】}\quad 明股自乘作为被除数，用通弦来约，得到明勾。用明勾减通勾得到城圆的直径。



\vspace{0.5em}



\Question{问十七}{%

乙出南门东行七十二步，丙出南门南行五十一步。二人相望与城相参直，计乙丙相距一百零二步。以明勾股测望。

}



\switchcolumn[0]*

\noindent\textbf{释曰}\quad 此以明勾股直接测望。乙东行为明勾，丙南行为明股，相距为明弦。



\switchcolumn[1]

\noindent\textbf{【释】}\quad 这是用明勾股直接进行测望的问题。乙向东行走是明勾，丙向南行走是明股，两人相距是明弦。



\switchcolumn[0]*

\noindent\textbf{术曰}\quad 明勾股各自乘相并开方得明弦。以通弦减之得城径。



\switchcolumn[1]

\noindent\textbf{【术】}\quad 明勾与明股各自乘，相加后开平方得到明弦。用通弦减去它得到城圆的直径。



\vspace{0.5em}



\Question{问十八}{%

甲从乾隅东行三百二十步，乙出南门东行七十二步。甲望乙与城相参直，又斜行二百八十九步与乙相会。以明勾与皇极弦测望。

}



\switchcolumn[0]*

\noindent\textbf{释曰}\quad 此以明勾与皇极弦测望。乙东行为明勾，甲斜行为皇极弦。



\switchcolumn[1]

\noindent\textbf{【释】}\quad 这是用明勾与皇极弦来进行测望的问题。乙向东行走是明勾，甲斜着行走是皇极弦。



\switchcolumn[0]*

\noindent\textbf{术曰}\quad 明勾自乘，减皇极弦自乘，开方得明股。以明股减通股得城径。



\switchcolumn[1]

\noindent\textbf{【术】}\quad 明勾自乘，减去皇极弦自乘，开平方得到明股。用明股减通股得到城圆的直径。



\vspace{0.5em}



\Question{问十九}{%

甲从乾隅南行六百步，丙出南门南行五十一步。甲望丙与城相参直。以明股与小差测望。

}



\switchcolumn[0]*

\noindent\textbf{释曰}\quad 此以明股与小差测望。丙南行为明股，小差即通弦减通勾之余。



\switchcolumn[1]

\noindent\textbf{【释】}\quad 这是用明股与小差来进行测望的问题。丙向南行走是明股，小差就是通弦减去通勾的余数。



\switchcolumn[0]*

\noindent\textbf{术曰}\quad 明股自乘，以小差约之，得明勾。以明勾减通勾得城径。



\switchcolumn[1]

\noindent\textbf{【术】}\quad 明股自乘，用小差来约，得到明勾。用明勾减通勾得到城圆的直径。



%----------------------------------------------------------------------

\Section{叀测望四}

%----------------------------------------------------------------------



\Question{问二十}{%

乙出东门南行三十步而立。甲从乾隅东行三百二十步，望乙与城相参直。以叀股测望。

}



\switchcolumn[0]*

\noindent\textbf{释曰}\quad 此以叀股立法测望。乙出东门南行为叀股。



\switchcolumn[1]

\noindent\textbf{【释】}\quad 这是用叀股来建立方法进行测望的问题。乙出东门向南行走是叀股。



\switchcolumn[0]*

\noindent\textbf{术曰}\quad 叀股自乘为实，以通弦约之，得叀勾。以叀勾减通勾得城径。



\switchcolumn[1]

\noindent\textbf{【术】}\quad 叀股自乘作为被除数，用通弦来约，得到叀勾。用叀勾减通勾得到城圆的直径。



\vspace{0.5em}



\Question{问二十一}{%

甲出东门东行二百步而立。丙出东门南行三十步。甲望丙与城相参直。以叀勾测望。

}



\switchcolumn[0]*

\noindent\textbf{释曰}\quad 此以叀勾立法测望。甲东行为叀勾，丙南行为叀股。



\switchcolumn[1]

\noindent\textbf{【释】}\quad 这是用叀勾来建立方法进行测望的问题。甲向东行走是叀勾，丙向南行走是叀股。



\switchcolumn[0]*

\noindent\textbf{术曰}\quad 叀勾自乘为实，以通弦约之，得叀股。以叀股减通股得城径。



\switchcolumn[1]

\noindent\textbf{【术】}\quad 叀勾自乘作为被除数，用通弦来约，得到叀股。用叀股减通股得到城圆的直径。



\vspace{0.5em}



\Question{问二十二}{%

甲出东门东行二百步，乙出东门南行三十步。二人相望与城相参直。以叀勾股直接测望。

}



\switchcolumn[0]*

\noindent\textbf{释曰}\quad 此以叀勾股直接测望。甲东行为叀勾，乙南行为叀股。



\switchcolumn[1]

\noindent\textbf{【释】}\quad 这是用叀勾股直接进行测望的问题。甲向东行走是叀勾，乙向南行走是叀股。



\switchcolumn[0]*

\noindent\textbf{术曰}\quad 叀勾股各自乘相并开方得叀弦。以通弦减之得城径。



\switchcolumn[1]

\noindent\textbf{【术】}\quad 叀勾与叀股各自乘，相加后开平方得到叀弦。用通弦减去它得到城圆的直径。



\vspace{0.5em}



\Question{问二十三}{%

甲从乾隅东行三百二十步，乙出东门南行三十步。甲斜行与乙相会，计甲斜行二百八十九步。以叀股与皇极弦测望。

}



\switchcolumn[0]*

\noindent\textbf{释曰}\quad 此以叀股与皇极弦测望。乙南行为叀股，甲斜行为皇极弦。



\switchcolumn[1]

\noindent\textbf{【释】}\quad 这是用叀股与皇极弦来进行测望的问题。乙向南行走是叀股，甲斜着行走是皇极弦。



\switchcolumn[0]*

\noindent\textbf{术曰}\quad 叀股自乘，减皇极弦自乘，开方得叀勾。以叀勾减通勾得城径。



\switchcolumn[1]

\noindent\textbf{【术】}\quad 叀股自乘，减去皇极弦自乘，开平方得到叀勾。用叀勾减通勾得到城圆的直径。



%----------------------------------------------------------------------

\Section{杂题测望五}

%----------------------------------------------------------------------



\Question{问二十四}{%

甲乙二人俱在城外西北乾隅，甲南行六百步而立，乙东行三百二十步而立。各不知余数。以勾股和较之变化，求杂题之数。

}



\switchcolumn[0]*

\noindent\textbf{释曰}\quad 此以杂题立法测望。杂题者，诸差互见之题也。



\switchcolumn[1]

\noindent\textbf{【释】}\quad 这是用杂题来建立方法进行测望的问题。杂题，就是各种差互相出现的题目。



\switchcolumn[0]*

\noindent\textbf{术曰}\quad 通勾股各自乘相并开方得通弦。以通弦减通股得大差，减通勾得小差。大小差相减得中差。诸差互以求之，得城径。



\switchcolumn[1]

\noindent\textbf{【术】}\quad 通勾与通股各自乘，相加后开平方得到通弦。用通弦减通股得到大差，减通勾得到小差。大小差相减得到中差。用各种差互相求解，得到城圆的直径。



\vspace{0.5em}



\Question{问二十五}{%

乙出东门南行三十步，丙出南门东行七十二步。二人相望与城相参直。以明叀二勾股测望。

}



\switchcolumn[0]*

\noindent\textbf{释曰}\quad 此以明叀二勾股杂测。乙南行为叀股，丙东行为明勾。



\switchcolumn[1]

\noindent\textbf{【释】}\quad 这是用明叀两种勾股进行杂测的问题。乙向南行走是叀股，丙向东行走是明勾。



\switchcolumn[0]*

\noindent\textbf{术曰}\quad 明勾叀股各自乘相并开方得杂弦。以通弦减之得城径。



\switchcolumn[1]

\noindent\textbf{【术】}\quad 明勾与叀股各自乘，相加后开平方得到杂弦。用通弦减去它得到城圆的直径。



\vspace{0.5em}



\Question{问二十六}{%

甲从乾隅东行三百二十步，乙出东门南行三十步，丙出南门东行七十二步。甲望乙丙与城相参直。以三差互见测望。

}



\switchcolumn[0]*

\noindent\textbf{释曰}\quad 此以大中小差杂见测望。乙南行为叀股，丙东行为明勾。



\switchcolumn[1]

\noindent\textbf{【释】}\quad 这是大中小差杂见进行测望的问题。乙向南行走是叀股，丙向东行走是明勾。



\switchcolumn[0]*

\noindent\textbf{术曰}\quad 明勾自乘为实，叀股自乘为实，二实相并开方得杂弦。以通弦减之得城径。



\switchcolumn[1]

\noindent\textbf{【术】}\quad 明勾自乘作为被除数，叀股自乘作为被除数，两个被除数相加后开平方得到杂弦。用通弦减去它得到城圆的直径。



\vspace{0.5em}



\Question{问二十七}{%

甲乙二人同出东门，甲东行乙南行。丙丁二人同出南门，丙南行丁东行。四人相望，计甲东行二百步，乙南行三十步，丙南行五十一步，丁东行七十二步。以四行互见测望。

}



\switchcolumn[0]*

\noindent\textbf{释曰}\quad 此以四行杂见测望。甲东行为叀勾，乙南行为叀股，丙南行为明股，丁东行为明勾。



\switchcolumn[1]

\noindent\textbf{【释】}\quad 这是用四行杂见来进行测望的问题。甲向东行走是叀勾，乙向南行走是叀股，丙向南行走是明股，丁向东行走是明勾。



\switchcolumn[0]*

\noindent\textbf{术曰}\quad 四行各自乘，相并开方得杂弦。以通弦减之得城径。



\switchcolumn[1]

\noindent\textbf{【术】}\quad 四种行走数各自乘，相加后开平方得到杂弦。用通弦减去它得到城圆的直径。



\vspace{0.5em}



\Question{问二十八}{%

甲从乾隅东行三百二十步，乙出东门南行三十步，丙出南门东行七十二步。甲斜望乙与城相参直，又斜望丙与城相参直。计甲至乙斜行二百八十九步，甲至丙斜行二百七十二步。以二斜行测望。

}



\switchcolumn[0]*

\noindent\textbf{释曰}\quad 此以二斜行杂测。甲至乙为皇极弦，甲至丙为小差弦。



\switchcolumn[1]

\noindent\textbf{【释】}\quad 这是用二斜行进行杂测的问题。甲到乙是皇极弦，甲到丙是小差弦。



\switchcolumn[0]*

\noindent\textbf{术曰}\quad 二弦各自乘相减，开方得勾股较。以加减术推之得城径。



\switchcolumn[1]

\noindent\textbf{【术】}\quad 两条弦各自乘，相减后开平方得到勾股差。用加减术推算得到城圆的直径。



\vspace{0.5em}



\Question{问二十九}{%

乙出东门东行二百步，丙出南门南行五十一步。二人相望与城相参直。甲从乾隅东行三百二十步，望乙丙与城相参直。以东行南行杂测。

}



\switchcolumn[0]*

\noindent\textbf{释曰}\quad 此以东行南行杂见测望。乙东行为边勾，丙南行为明股。



\switchcolumn[1]

\noindent\textbf{【释】}\quad 这是用东行南行杂见来进行测望的问题。乙向东行走是边勾，丙向南行走是明股。



\switchcolumn[0]*

\noindent\textbf{术曰}\quad 边勾明股各自乘相并开方得杂弦。以通弦减之得城径。



\switchcolumn[1]

\noindent\textbf{【术】}\quad 边勾与明股各自乘，相加后开平方得到杂弦。用通弦减去它得到城圆的直径。



\vspace{0.5em}



\Question{问三十}{%

甲从艮隅南行过城门三百六十步，乙从坤隅东行过城门一百六十步。二人相望与城相参直。以过城门之行杂测。

}



\switchcolumn[0]*

\noindent\textbf{释曰}\quad 此以过城门之行杂测。甲南行过城门为小差股，乙东行过城门为大差勾。



\switchcolumn[1]

\noindent\textbf{【释】}\quad 这是用走过城门的行走数进行杂测的问题。甲向南走过城门是小差股，乙向东走过城门是大差勾。



\switchcolumn[0]*

\noindent\textbf{术曰}\quad 大差勾小差股各自乘相并开方得杂弦。以通弦减之得城径。



\switchcolumn[1]

\noindent\textbf{【术】}\quad 大差勾与小差股各自乘，相加后开平方得到杂弦。用通弦减去它得到城圆的直径。



\vspace{0.5em}



\Question{问三十一}{%

甲乙二人俱在城外西北乾隅，甲南行六百步而立，乙东行三百二十步而立。甲斜行与乙相会。以通勾股弦诸差互求。

}



\switchcolumn[0]*

\noindent\textbf{释曰}\quad 此以诸差互见杂测。通勾股弦诸差互相求之術。



\switchcolumn[1]

\noindent\textbf{【释】}\quad 这是用诸差互见进行杂测的问题。通勾股弦各种差互相求解的方法。



\switchcolumn[0]*

\noindent\textbf{术曰}\quad 通弦自乘，减通股自乘，开方得大差勾。通弦自乘，减通勾自乘，开方得小差股。大小差相减得中差。诸差互以求之，得城径。



\switchcolumn[1]

\noindent\textbf{【术】}\quad 通弦自乘，减去通股自乘，开平方得到大差勾。通弦自乘，减去通勾自乘，开平方得到小差股。大小差相减得到中差。用各种差互相求解，得到城圆的直径。



\vspace{0.5em}



\Question{问三十二}{%

甲出东门东行二百步，乙出南门南行五十一步。二人相望与城相参直，计相距二百八十九步。以相距与各行杂测。

}



\switchcolumn[0]*

\noindent\textbf{释曰}\quad 此以相距与各行杂见测望。相距为皇极弦。



\switchcolumn[1]

\noindent\textbf{【释】}\quad 这是用相距与各行杂见来进行测望的问题。相距就是皇极弦。



\switchcolumn[0]*

\noindent\textbf{术曰}\quad 相距自乘，减东行自乘，开方得南行余数。以减通股得城径。



\switchcolumn[1]

\noindent\textbf{【术】}\quad 相距自乘，减去东行自乘，开平方得到南行的余数。用通股减去它得到城圆的直径。



\vspace{0.5em}



\Question{问三十三}{%

乙出东门南行三十步，丙出南门东行七十二步，丁出东门东行二百步。甲从乾隅东行三百二十步，望乙丙丁与城相参直。以三行杂测。

}



\switchcolumn[0]*

\noindent\textbf{释曰}\quad 此以三行杂见测望。乙南行为叀股，丙东行为明勾，丁东行为边勾。



\switchcolumn[1]

\noindent\textbf{【释】}\quad 这是用三行杂见来进行测望的问题。乙向南行走是叀股，丙向东行走是明勾，丁向东行走是边勾。



\switchcolumn[0]*

\noindent\textbf{术曰}\quad 三行各自乘，相并开方得杂弦。以通弦减之得城径。



\switchcolumn[1]

\noindent\textbf{【术】}\quad 三种行走数各自乘，相加后开平方得到杂弦。用通弦减去它得到城圆的直径。



\vspace{0.5em}



\Question{问三十四}{%

甲乙丙丁四人，甲从乾隅东行三百二十步，乙从乾隅南行六百步，丙出东门东行二百步，丁出南门南行五十一步。以四隅之行杂测。

}



\switchcolumn[0]*

\noindent\textbf{释曰}\quad 此以四隅之行杂见测望。诸行互见，以天元术推之。



\switchcolumn[1]

\noindent\textbf{【释】}\quad 这是用四隅之行杂见来进行测望的问题。各种行走数互相出现，用天元术来推算。



\switchcolumn[0]*

\noindent\textbf{术曰}\quad 通勾股各自乘相并开方得通弦。以通弦减通股得大差，减通勾得小差。以大小差推诸数，得城径。



\switchcolumn[1]

\noindent\textbf{【术】}\quad 通勾与通股各自乘，相加后开平方得到通弦。用通弦减通股得到大差，减通勾得到小差。用大小差推算各种数值，得到城圆的直径。



\EndVol{测圆海镜分类释术卷八终}



%======================================================================

\Volume{九}

%======================================================================



\switchcolumn[0]*

\noindent\textbf{测圆海镜分类释术卷九}\quad 元~李~冶~撰\\

明~顾应祥~释术



\switchcolumn[1]

\noindent\textbf{测圆海镜分类释术·第九卷}\quad 元代~李~冶~撰\\

明代~顾应祥~释术



\vspace{1em}



\switchcolumn[0]*

卷九专论诸弦差较之测望。诸弦者，通弦、边弦、底弦、黄广弦、黄长弦之类；差较者，弦与勾股之差及差与差之比较也。凡十八题，分为五节。



\switchcolumn[1]

第九卷专门论述各种弦差差的测望方法。诸弦，包括通弦、边弦、底弦、黄广弦、黄长弦等；差较，就是弦与勾股的差以及差与差的比较。共有十八道题，分为五节。



%----------------------------------------------------------------------

\Section{大股小勾测望一}

%----------------------------------------------------------------------



\Question{问一}{%

甲从乾隅南行六百步而立，乙从乾隅东行三百二十步而立。以通股与通勾之差，求大股小勾。

}



\switchcolumn[0]*

\noindent\textbf{释曰}\quad 此以大股小勾立法测望。通股六百步为大股，通勾三百二十步为小勾。大股小勾之相较，为勾股之极差。



\switchcolumn[1]

\noindent\textbf{【释】}\quad 这是用大股小勾来建立方法进行测望的问题。通股六百步是大股，通勾三百二十步是小勾。大股与小勾的比较，就是勾股的极差。



\switchcolumn[0]*

\noindent\textbf{术曰}\quad 通股自乘，减通勾自乘，开方得勾股和。以和减通弦，余为弦较。以弦较加减通勾股，得城径。



\switchcolumn[1]

\noindent\textbf{【术】}\quad 通股自乘，减去通勾自乘，开平方得到勾股和。用和减通弦，余下就是弦较。用弦较加减通勾股，得到城圆的直径。



\vspace{0.5em}



\Question{问二}{%

乙出东门南行三十步而立，丙出南门东行七十二步而立。甲从乾隅东行三百二十步，望乙丙与城相参直。以大股小勾测望。

}



\switchcolumn[0]*

\noindent\textbf{释曰}\quad 此以叀股明勾测望大股小勾。乙南行为叀股，丙东行为明勾。



\switchcolumn[1]

\noindent\textbf{【释】}\quad 这是用叀股明勾来测望大股小勾的问题。乙向南行走是叀股，丙向东行走是明勾。



\switchcolumn[0]*

\noindent\textbf{术曰}\quad 叀股明勾各自乘相并开方得小弦。以通弦减之得城径。



\switchcolumn[1]

\noindent\textbf{【术】}\quad 叀股与明勾各自乘，相加后开平方得到小弦。用通弦减去它得到城圆的直径。



\vspace{0.5em}



\Question{问三}{%

甲从乾隅东行三百二十步，乙出东门东行二百步。二人相望与城相参直。以大股与边勾测望。

}



\switchcolumn[0]*

\noindent\textbf{释曰}\quad 此以大股边勾测望。甲东行为通勾，乙东行为边勾。



\switchcolumn[1]

\noindent\textbf{【释】}\quad 这是用大股边勾来进行测望的问题。甲向东行走是通勾，乙向东行走是边勾。



\switchcolumn[0]*

\noindent\textbf{术曰}\quad 通勾减边勾，余为小勾。以小勾加减术推之得城径。



\switchcolumn[1]

\noindent\textbf{【术】}\quad 通勾减去边勾，余下就是小勾。用小勾加减术推算得到城圆的直径。



%----------------------------------------------------------------------

\Section{大差大小差测望二}

%----------------------------------------------------------------------



\Question{问四}{%

甲乙二人俱在城外西北乾隅，甲南行六百步而立，乙东行三百二十步而立。以通弦减通股为大差，减通勾为小差。大小差相求。

}



\switchcolumn[0]*

\noindent\textbf{释曰}\quad 此以大差大小差立法测望。大差即通股弦较，小差即通勾弦较。



\switchcolumn[1]

\noindent\textbf{【释】}\quad 这是用大差大小差来建立方法进行测望的问题。大差就是通股弦差，小差就是通勾弦差。



\switchcolumn[0]*

\noindent\textbf{术曰}\quad 大差自乘为实，小差自乘为实，二实相减开方得中差。以中差加减大小差，得城径。



\switchcolumn[1]

\noindent\textbf{【术】}\quad 大差自乘作为被除数，小差自乘作为被除数，两个被除数相减后开平方得到中差。用中差加减大小差，得到城圆的直径。



\vspace{0.5em}



\Question{问五}{%

乙出东门南行三十步而立，丙出南门东行七十二步而立。以叀股明勾求大小差。

}



\switchcolumn[0]*

\noindent\textbf{释曰}\quad 此以叀股明勾测望大小差。叀股即小差之属，明勾即大差之属。



\switchcolumn[1]

\noindent\textbf{【释】}\quad 这是用叀股明勾来测望大小差的问题。叀股就是小差这一类，明勾就是大差这一类。



\switchcolumn[0]*

\noindent\textbf{术曰}\quad 叀股自乘减明勾自乘，开方得大小差较。以较加减术推之得城径。



\switchcolumn[1]

\noindent\textbf{【术】}\quad 叀股自乘减去明勾自乘，开平方得到大小差差。用差加减术推算得到城圆的直径。



\vspace{0.5em}



\Question{问六}{%

甲从乾隅东行三百二十步，乙出东门南行三十步，丙出南门东行七十二步。甲望乙丙与城相参直。以三差互求。

}



\switchcolumn[0]*

\noindent\textbf{释曰}\quad 此以大中小差互见测望。乙南行为叀股，丙东行为明勾。



\switchcolumn[1]

\noindent\textbf{【释】}\quad 这是大中小差互见进行测望的问题。乙向南行走是叀股，丙向东行走是明勾。



\switchcolumn[0]*

\noindent\textbf{术曰}\quad 明勾自乘为大差实，叀股自乘为小差实。二实相减开方得中差。以中差加减得城径。



\switchcolumn[1]

\noindent\textbf{【术】}\quad 明勾自乘作为大差被除数，叀股自乘作为小差被除数。两个被除数相减后开平方得到中差。用中差加减得到城圆的直径。



\vspace{0.5em}



\Question{问七}{%

甲从艮隅南行过城门三百六十步，乙从坤隅东行过城门一百六十步。以过城门之行求大小差。

}



\switchcolumn[0]*

\noindent\textbf{释曰}\quad 此以过城门之行测望大小差。甲南行为小差股，乙东行为大差勾。



\switchcolumn[1]

\noindent\textbf{【释】}\quad 这是用走过城门的行走数来测望大小差的问题。甲向南行走是小差股，乙向东行走是大差勾。



\switchcolumn[0]*

\noindent\textbf{术曰}\quad 小差股自乘为实，大差勾自乘为实。二实相并开方得大小差弦。以通弦减之得城径。



\switchcolumn[1]

\noindent\textbf{【术】}\quad 小差股自乘作为被除数，大差勾自乘作为被除数。两个被除数相加后开平方得到大小差弦。用通弦减去它得到城圆的直径。



%----------------------------------------------------------------------

\Section{诸弦测望三}

%----------------------------------------------------------------------



\Question{问八}{%

甲乙二人俱在城外西北乾隅，甲南行六百步而立，乙东行三百二十步而立。甲斜行与乙相会。以通弦为本，求边弦、底弦。

}



\switchcolumn[0]*

\noindent\textbf{释曰}\quad 此以诸弦立法测望。通弦为勾股弦之全，边弦、底弦为其分支。



\switchcolumn[1]

\noindent\textbf{【释】}\quad 这是用诸弦来建立方法进行测望的问题。通弦是勾股弦的全体，边弦、底弦是它的分支。



\switchcolumn[0]*

\noindent\textbf{术曰}\quad 通勾股各自乘相并开方得通弦。通弦半之得皇极弦。以通弦减皇极弦余为边弦。边弦减通弦余为城径。



\switchcolumn[1]

\noindent\textbf{【术】}\quad 通勾与通股各自乘，相加后开平方得到通弦。通弦取一半得到皇极弦。用通弦减皇极弦余下就是边弦。边弦减通弦余下就是城圆的直径。



\vspace{0.5em}



\Question{问九}{%

乙出东门东行二百步，丙出南门南行五十一步。二人相望与城相参直。以边股明勾求诸弦。

}



\switchcolumn[0]*

\noindent\textbf{释曰}\quad 此以边股明勾测望诸弦。乙东行为边股，丙南行为明股。



\switchcolumn[1]

\noindent\textbf{【释】}\quad 这是用边股明勾来测望诸弦的问题。乙向东行走是边股，丙向南行走是明股。



\switchcolumn[0]*

\noindent\textbf{术曰}\quad 边股明股各自乘相并开方得边弦。以通弦减之得城径。



\switchcolumn[1]

\noindent\textbf{【术】}\quad 边股与明股各自乘，相加后开平方得到边弦。用通弦减去它得到城圆的直径。



\vspace{0.5em}



\Question{问十}{%

甲从乾隅东行三百二十步，乙出东门南行三十步。甲望乙与城相参直，又斜行二百八十九步与乙相会。以底弦与通弦测望。

}



\switchcolumn[0]*

\noindent\textbf{释曰}\quad 此以底弦通弦测望。甲斜行为皇极弦，即底弦之属。



\switchcolumn[1]

\noindent\textbf{【释】}\quad 这是用底弦通弦来进行测望的问题。甲斜着行走是皇极弦，就是底弦这一类。



\switchcolumn[0]*

\noindent\textbf{术曰}\quad 皇极弦自乘，减通勾自乘，开方得底股。以底股减通股得城径。



\switchcolumn[1]

\noindent\textbf{【术】}\quad 皇极弦自乘，减去通勾自乘，开平方得到底股。用底股减通股得到城圆的直径。



%----------------------------------------------------------------------

\Section{大小差测望四}

%----------------------------------------------------------------------



\Question{问十一}{%

甲乙二人俱在城外西北乾隅，甲南行六百步而立，乙东行三百二十步而立。以通股减通弦为大差，通勾减通弦为小差。大小差互求。

}



\switchcolumn[0]*

\noindent\textbf{释曰}\quad 此以大小差互见测望。大差为通股弦较，小差为通勾弦较。



\switchcolumn[1]

\noindent\textbf{【释】}\quad 这是用大小差互见来进行测望的问题。大差是通股弦差，小差是通勾弦差。



\switchcolumn[0]*

\noindent\textbf{术曰}\quad 大差小差各自乘相并开方得中差弦。以通弦减之得城径。



\switchcolumn[1]

\noindent\textbf{【术】}\quad 大差与小差各自乘，相加后开平方得到中差弦。用通弦减去它得到城圆的直径。



\vspace{0.5em}



\Question{问十二}{%

乙出东门南行三十步，丙出南门东行七十二步。甲从乾隅东行三百二十步，望乙丙与城相参直。以叀股明勾求大小差。

}



\switchcolumn[0]*

\noindent\textbf{释曰}\quad 此以叀股明勾测望大小差。



\switchcolumn[1]

\noindent\textbf{【释】}\quad 这是用叀股明勾来测望大小差的问题。



\switchcolumn[0]*

\noindent\textbf{术曰}\quad 叀股自乘为小差实，明勾自乘为大差实。二实相减开方得差较。以差较加减得城径。



\switchcolumn[1]

\noindent\textbf{【术】}\quad 叀股自乘作为小差被除数，明勾自乘作为大差被除数。两个被除数相减后开平方得到差较。用差较加减得到城圆的直径。



\vspace{0.5em}



\Question{问十三}{%

甲从艮隅南行过城门三百六十步，乙从坤隅东行过城门一百六十步。以过城门之行求差较。

}



\switchcolumn[0]*

\noindent\textbf{释曰}\quad 此以过城门之行测望差较。



\switchcolumn[1]

\noindent\textbf{【释】}\quad 这是用走过城门的行走数来测望差较的问题。



\switchcolumn[0]*

\noindent\textbf{术曰}\quad 小差股大差勾各自乘相并开方得差弦。以通弦减之得城径。



\switchcolumn[1]

\noindent\textbf{【术】}\quad 小差股与大差勾各自乘，相加后开平方得到差弦。用通弦减去它得到城圆的直径。



\vspace{0.5em}



\Question{问十四}{%

甲出东门东行二百步，乙出南门南行五十一步。二人相望与城相参直。以边股明股求大小差。

}



\switchcolumn[0]*

\noindent\textbf{释曰}\quad 此以边股明股测望大小差。



\switchcolumn[1]

\noindent\textbf{【释】}\quad 这是用边股明股来测望大小差的问题。



\switchcolumn[0]*

\noindent\textbf{术曰}\quad 边股明股各自乘相并开方得弦。以通弦减之得城径。



\switchcolumn[1]

\noindent\textbf{【术】}\quad 边股与明股各自乘，相加后开平方得到弦。用通弦减去它得到城圆的直径。



\EndVol{测圆海镜分类释术卷九终}



%======================================================================

\Volume{十}

%======================================================================



\switchcolumn[0]*

\noindent\textbf{测圆海镜分类释术卷十}\quad 元~李~冶~撰\\

明~顾应祥~释术



\switchcolumn[1]

\noindent\textbf{测圆海镜分类释术·第十卷}\quad 元代~李~冶~撰\\

明代~顾应祥~释术



\vspace{1em}



\switchcolumn[0]*

卷十为全书之终，专论杂法。杂法者，诸术之变通也，凡不属于前九卷之专题者，皆汇于此。凡四题，皆为通术之总结。



\switchcolumn[1]

第十卷是全书的终结，专门论述杂法。杂法，就是各种方法的变通，凡是不属于前九卷专题的内容，都汇集在这里。共有四道题，都是通用方法的总结。



%----------------------------------------------------------------------

\Section{杂法}

%----------------------------------------------------------------------



\Question{问一}{%

甲乙二人俱在城外西北乾隅，甲南行六百步而立，乙东行三百二十步而立。二人隔城相望，甲斜行与乙相会。以天元术总求诸数。

}



\switchcolumn[0]*

\noindent\textbf{释曰}\quad 此以天元术杂法总测。天元一者，立天元一为勾股之基，以诸数推之。



\switchcolumn[1]

\noindent\textbf{【释】}\quad 这是用天元术杂法进行总体测望的问题。天元一，就是设天元一作为勾股的基础，用各个数值来推算。



\switchcolumn[0]*

\noindent\textbf{术曰}\quad 立天元一为勾，以股乘之，得勾股积。以弦除之，得容圆径。天元自之为勾幂，以减弦幂，开方得股。勾股各自乘相并开方得弦。以勾股求容圆术，得城径二百四十步。



\switchcolumn[1]

\noindent\textbf{【术】}\quad 设天元一作为勾，用股来乘，得到勾股积。用弦来除，得到容圆的直径。天元自乘作为勾的幂，用来减弦的幂，开平方得到股。勾与股各自乘，相加后开平方得到弦。用勾股求容圆的方法，得到城圆直径二百四十步。



\vspace{0.5em}



\Question{问二}{%

乙出东门南行三十步，丙出南门东行七十二步，丁出东门东行二百步，戊出南门南行五十一步。五人各不知余数。甲从乾隅东行三百二十步，望乙丙丁戊与城相参直。以天元术杂测诸行。

}



\switchcolumn[0]*

\noindent\textbf{释曰}\quad 此以天元术杂测诸行。乙南行为叀股，丙东行为明勾，丁东行为边勾，戊南行为明股。



\switchcolumn[1]

\noindent\textbf{【释】}\quad 这是用天元术来杂测各种行走数的问题。乙向南行走是叀股，丙向东行走是明勾，丁向东行走是边勾，戊向南行走是明股。



\switchcolumn[0]*

\noindent\textbf{术曰}\quad 立天元一为城径。以通勾减天元，得边勾。以通股减天元，得明股。边勾明股各自乘相并开方得皇极弦。以通弦减皇极弦，余为太虚弦。太虚弦自乘，以明和约之，得太虚勾。以次推之，得天元即城径二百四十步。



\switchcolumn[1]

\noindent\textbf{【术】}\quad 设天元一作为城圆的直径。用通勾减天元，得到边勾。用通股减天元，得到明股。边勾与明股各自乘，相加后开平方得到皇极弦。用通弦减皇极弦，余下就是太虚弦。太虚弦自乘，用明和来约，得到太虚勾。依次推算，得到天元就是城圆直径二百四十步。



\vspace{0.5em}



\Question{问三}{%

甲乙丙丁四人，甲从乾隅东行三百二十步，乙从乾隅南行六百步，丙出东门东行二百步，丁出南门南行五十一步。以天元术总括全书之术。

}



\switchcolumn[0]*

\noindent\textbf{释曰}\quad 此以天元术总括全书。通勾三百二十步，通股六百步，通弦以勾股术求之。边勾、明股、皇极弦、太虚弦，皆以天元一推之。



\switchcolumn[1]

\noindent\textbf{【释】}\quad 这是用天元术来总括全书的问题。通勾三百二十步，通股六百步，通弦用勾股术来求。边勾、明股、皇极弦、太虚弦，都用天元一来推算。



\switchcolumn[0]*

\noindent\textbf{术曰}\quad 立天元一为圆径。通勾自乘，以天元减之，得边勾幂。通股自乘，以天元减之，得明股幂。边勾幂明股幂各自乘相并开方得皇极弦。皇极弦自乘，以通弦减之，得诸弦之差。以次推之，得天元即城径二百四十步。



\switchcolumn[1]

\noindent\textbf{【术】}\quad 设天元一作为圆的直径。通勾自乘，用天元来减，得到边勾的幂。通股自乘，用天元来减，得到明股的幂。边勾幂与明股幂各自乘，相加后开平方得到皇极弦。皇极弦自乘，用通弦来减，得到各种弦的差。依次推算，得到天元就是城圆直径二百四十步。



\vspace{0.5em}



\Question{问四}{%

全书诸题，以勾股容圆为本，以天元一为法。通勾、通股、通弦为三大纲，边勾、明股、皇极弦、太虚弦为四目，大差、小差、中差为三变。总括其术。

}



\switchcolumn[0]*

\noindent\textbf{释曰}\quad 此总括全书之术。测圆海镜分类释术，凡十卷，一百七十题。以勾股容圆为题，以天元术为法。通勾三百二十步，通股六百步，通弦以勾股术求之得六百八十步。城径二百四十步。诸差：大差八十步，小差三百六十步，中差二百八十步。边勾一百六十步，明股七十二步，皇极弦二百八十九步，太虚弦一百零二步。凡此诸数，皆以天元一推之而得。



\switchcolumn[1]

\noindent\textbf{【释】}\quad 这是总括全书方法的问题。测圆海镜分类释术，共十卷，一百七十道题。以勾股容圆为题目，以天元术为方法。通勾三百二十步，通股六百步，通弦用勾股术求得六百八十步。城圆直径二百四十步。各种差：大差八十步，小差三百六十步，中差二百八十步。边勾一百六十步，明股七十二步，皇极弦二百八十九步，太虚弦一百零二步。所有这些数值，都是用天元一推算而得的。



\switchcolumn[0]*

\noindent\textbf{术曰}\quad 立天元一为圆径。以勾股容圆术推之：勾股和减弦，余为圆径。通勾三百二十步，通股六百步，通弦六百八十步。勾股和九百二十步，减弦六百八十步，余二百四十步，即城径。此全书之宗也。



\switchcolumn[1]

\noindent\textbf{【术】}\quad 设天元一作为圆的直径。用勾股容圆术来推算：勾股和减去弦，余下就是圆的直径。通勾三百二十步，通股六百步，通弦六百八十步。勾股和九百二十步，减去弦六百八十步，余下二百四十步，就是城圆的直径。这是全书的宗旨。



\EndVol{测圆海镜分类释术卷十终}



\end{paracol}



\vspace{2em}

\begin{center}

\rule{0.6\textwidth}{0.5pt}



\textbf{测圆海镜分类释术卷六至卷十终}



\vspace{0.5em}



{\small 元李冶撰\quad 明顾应祥释术}



\vspace{0.5em}



{\small 钦定四库全书\quad 子部\quad 天文算法类}

\end{center}



\end{document}







% ============================================================

% Source: translations/zh/chinese/li-ye-ceyuan-haijing-vols1-3_classical-modern_bilingual.tex

% ============================================================



% =============================================================================

% 測圓海鏡（卷一至卷三）文言文←→白話文雙語對照版

% Ceyuan Haijing (Vols. 1-3) — Classical↔Modern Chinese Bilingual Edition

% =============================================================================

\documentclass[11pt,a4paper]{article}



% -------- Core Packages --------

\usepackage{paracol}

\usepackage{fontspec}

\usepackage{polyglossia}

\usepackage{amsmath,amssymb,amsthm}

\usepackage{geometry}

\usepackage{graphicx}

\usepackage{xcolor}

\usepackage{titlesec}

\usepackage{fancyhdr}

\usepackage{enumitem}

\usepackage{array,booktabs,longtable,multirow}

\usepackage{hyperref}

\usepackage{tocloft}

\usepackage{indentfirst}

\usepackage{xeCJK}



% -------- CJK Font Setup --------

\setCJKmainfont{Microsoft YaHei}

\setCJKsansfont{Microsoft YaHei}

\setCJKmonofont{Microsoft YaHei}



% -------- Latin Font Setup --------

\setmainfont{Times New Roman}

\setsansfont{Arial}



% -------- Page Geometry --------

\geometry{

  paperwidth=210mm,paperheight=297mm,

  left=18mm,right=18mm,top=25mm,bottom=25mm,

  headheight=14pt,footskip=30pt

}



% -------- Column Setup --------

\columnratio{0.48}  % Slightly narrower for classical text

\setlength{\columnsep}{1.2em}

\setlength{\columnseprule}{0.6pt}

\def\columnseprulecolor{\color{rulecolor}}



% -------- Colors --------

\definecolor{titlecolor}{RGB}{45,55,72}

\definecolor{sectioncolor}{RGB}{55,65,81}

\definecolor{notecolor}{RGB}{120,53,15}

\definecolor{rulecolor}{RGB}{160,140,120}

\definecolor{prefacecolor}{RGB}{80,60,40}

\definecolor{problemcolor}{RGB}{30,60,90}

\definecolor{classicalbg}{RGB}{255,252,245}

\definecolor{modernbg}{RGB}{245,250,255}

\definecolor{leftheadcolor}{RGB}{120,50,30}

\definecolor{rightheadcolor}{RGB}{30,50,120}



% -------- Section Formatting --------

\titleformat{\section}

  {\normalfont\Large\bfseries\color{sectioncolor}}

  {\thesection}{1em}{}

\titleformat{\subsection}

  {\normalfont\large\bfseries\color{sectioncolor}}

  {\thesubsection}{1em}{}

\titleformat{\subsubsection}

  {\normalfont\normalsize\bfseries\color{sectioncolor}}

  {\thesubsubsection}{1em}{}



% -------- Header/Footer --------

\pagestyle{fancy}

\fancyhf{}

\fancyhead[L]{\small\color{leftheadcolor}\textit{測圓海鏡 · 文言文}}

\fancyhead[C]{\small\textcolor{rulecolor}{$\leftrightarrow$}}

\fancyhead[R]{\small\color{rightheadcolor}\textit{白話文對照}}

\fancyfoot[C]{\small\thepage}

\renewcommand{\headrulewidth}{0.4pt}

\renewcommand{\footrulewidth}{0pt}



% -------- Custom Commands for Bilingual Structure --------



% Problem structure: 問/答/術/草

\newcommand{\wen}[1]{\par\noindent\textbf{\color{problemcolor}【問】} #1\par}

\newcommand{\da}[1]{\par\noindent\textbf{\color{problemcolor}【答】} #1\par}

\newcommand{\shu}[1]{\par\noindent\textbf{\color{problemcolor}【術】} #1\par}

\newcommand{\cao}[1]{\par\noindent\textbf{\color{problemcolor}【草】} #1\par}

\newcommand{\an}[1]{\par\noindent\textit{〔按〕}#1\par}



% Modern Chinese equivalents

\newcommand{\wenm}[1]{\par\noindent\textbf{\color{problemcolor}【問題】} #1\par}

\newcommand{\dam}[1]{\par\noindent\textbf{\color{problemcolor}【答案】} #1\par}

\newcommand{\shum}[1]{\par\noindent\textbf{\color{problemcolor}【解法】} #1\par}

\newcommand{\caom}[1]{\par\noindent\textbf{\color{problemcolor}【演算】} #1\par}

\newcommand{\anm}[1]{\par\noindent\textit{〔譯者按〕}#1\par}



% Hyperref

\hypersetup{

  colorlinks=true,

  linkcolor=sectioncolor,

  citecolor=sectioncolor,

  urlcolor=sectioncolor,

  pdfencoding=auto

}



\setcounter{tocdepth}{2}



% =============================================================================

\begin{document}



% -------- Title Page --------

\begin{titlepage}

\centering

\vspace*{2cm}



{\fontsize{36}{44}\selectfont\CJKfamily{sf}\color{titlecolor}

測圓海鏡}



\vspace{0.6cm}



{\fontsize{14}{18}\selectfont\color{sectioncolor}

文言文 \raisebox{0.1ex}{\textcolor{rulecolor}{$\longleftrightarrow$}} 白話文對照}



\vspace{0.3cm}



{\small\color{rulecolor}

Classical Chinese $\longleftrightarrow$ Modern Chinese Vernacular}



\vspace{0.5cm}



{\color{rulecolor}\rule{0.6\textwidth}{0.4pt}}



\vspace{0.8cm}



{\LARGE 卷一 — 卷三}



{\large Volumes I – III}



\vspace{1.5cm}



{\large\color{sectioncolor}

元 \quad 李冶 \quad 撰}



\vspace{0.3cm}



{\normalsize\color{sectioncolor}

(Li Ye / Li Zhi, 1192–1279)}



\vspace{2cm}



{\color{rulecolor}\rule{0.4\textwidth}{0.4pt}}



\vspace{1cm}



{\normalsize

\textit{以天元術解勾股容圓一百七十問}}



\vspace{0.3cm}



{\small\color{notecolor}

\textit{以天元術（設未知數列方程的方法）}\newline

\textit{解直角三角形內切圓問題共一百七十道}}



\vfill



{\small

雙語對照版 · 左欄：原文 \quad 右欄：白話文譯}



\end{titlepage}



% -------- Table of Contents --------

\tableofcontents

\newpage



% =============================================================================

% INTRODUCTION

% =============================================================================

\section{導讀 Introduction}

\label{sec:intro}

\addcontentsline{toc}{section}{導讀 Introduction}



\begin{paracol}{2}

\textbf{【文言文】}《測圓海鏡》乃李冶（李治，1192--1279）之傑作。李冶為宋元四大數學家之一，與楊輝、秦九韶、朱世傑齊名。此書成於公元1248年，系統運用\textbf{天元術}（設未知數列方程之法），以解\textbf{勾股容圓}（直角三角形內切圓）相關之170道幾何問題。



全書結構如下：\textbf{卷一}奠定理論基礎，含圓城圖式、總率名號、今問正數、識別雜記（含692條幾何公式）；\textbf{卷二至卷十二}則以天元術解170道問題，方程次數自一次至六次不等。



本文檔含\textbf{卷一至卷三}之全文，包括所有基礎理論及首批問題。



\switchcolumn

\textbf{【白話文】}《測圓海鏡》是李冶（又名李治，1192--1279年）的代表作。李冶是宋元時期四大數學家之一，與楊輝、秦九韶、朱世傑並稱。此書完成於公元1248年，系統運用\textbf{天元術（設未知數$x$列方程的方法）}來解決與\textbf{勾股容圓（直角三角形內切圓）}相關的170道幾何問題。



全書的結構是這樣的：\textbf{第一卷}奠定理論基礎，包括圓城圖式（主要幾何圖形）、總率名號（15個直角三角形的定義）、今問正數（所有線段的數值）和識別雜記（包含692條幾何公式）；\textbf{第二卷到第十二卷}則運用天元術來解170道問題，方程的次數從一次到六次不等。



本文檔包含\textbf{第一卷到第三卷}的完整內容，涵蓋所有基礎理論和第一批問題。

\end{paracol}



\vspace{0.5cm}

\noindent\textcolor{rulecolor}{\rule{\textwidth}{0.4pt}}



% =============================================================================

% ORIGINAL PREFACE

% =============================================================================

\newpage

\section{原序 Original Preface}

\label{sec:preface}



\begin{paracol}{2}

{\color{prefacecolor}

\textbf{【文言文】}



\begin{quote}

數本難窮，吾欲以力強窮之，彼其數不惟不能得其凡，而吾之力且憊矣。然則數果不可以窮耶？既已名之數矣，則又何為而不可窮也。



故謂數為難窮斯可，謂數為不可窮斯不可。何則？彼其冥冥之中，固有昭昭者存，夫昭昭者其自然之數也。非自然之數，其自然之理也。數一出於自然，吾欲以力強窮之，使隸首復生亦末如之何也已。



苟能推自然之理，以明自然之數，則雖遠而乾端坤倪，幽而神情鬼狀，未有不合者矣。

\end{quote}

}



\switchcolumn

{\color{prefacecolor}

\textbf{【白話文】}



\begin{quote}

數學本來就是難以窮盡的，如果我想用蠻力去強行窮盡它，那麼不僅不能把握數學的本質，反而會耗盡我的精力。那麼數學真的不能窮盡嗎？既然已經稱它為「數」了，為什麼又不能窮盡呢？



所以說數學難以窮盡是可以的，但說數學不能窮盡就不對了。為什麼呢？在那冥冥之中，確實有昭然若揭的規律存在。那些昭然若揭的東西就是自然的「數」。與其說是自然的數，不如說是自然的理。數學完全源於自然，如果我想用蠻力強行窮盡它，即使隸首（傳說中的算術發明者）復活也無能為力。



如果能夠推究自然的道理來闡明自然的數學規律，那麼即使遙遠如天地邊際，深奧如鬼神形狀，也無不符合這些規律。

\end{quote}

}

\end{paracol}



\vspace{0.5cm}



\begin{paracol}{2}

{\color{prefacecolor}

\textbf{【文言文】}



\begin{quote}

予自幼喜算數，恆病夫考圓之術例出於牽強，殊乖於自然，如古率、徽率、密率之不同，截弧、截矢、截背之互見，內外諸角，析剖支條，莫不各自名家，與世作法。及反覆研究，卒無以當吾心焉。



老大以來，得洞淵九容之說，日夕玩繹，而向之病我者，使爆然落去而無遺余。山中多暇，客有從余求其說者，於是乎又為衍之，遂累一百七十問。



既成編，客復目之《測圓海鏡》，蓋取夫天臨海鏡之義也。

\end{quote}

}



\switchcolumn

{\color{prefacecolor}

\textbf{【白話文】}



\begin{quote}

我從小就喜歡數學，一直不滿意那些研究圓的算法總是顯得牽強附會，違背自然之道。比如古率（周三徑一）、徽率（劉徽的圓周率）、密率（祖沖之的密率）各不相同；截弧、截矢、截背等概念互不相通；內角外角的各種分析支離破碎，各家都自立門派，為世人提供不同的方法。我反覆研究，始終沒有找到令我心滿意足的理論。



到了老年，我得到了洞淵九容的學說，日夜鑽研，以往困擾我的問題如同豁然開朗，煙消雲散。山中有很多閒暇時間，有客人向我請教這些學說，於是我便加以推演，最終積累了一百七十道問題。



編成書後，客人將其命名為《測圓海鏡》，取的是「天臨海鏡」（上天俯視海中的鏡子）這個含義。

\end{quote}

}

\end{paracol}



\vspace{0.5cm}



\begin{paracol}{2}

{\color{prefacecolor}

\textbf{【文言文】}



\begin{quote}

昔半山老人集《唐百家詩選》，自謂廢日力於此，良可惜。明道先生以上蔡謝君記誦為玩物喪志。夫文史尚矣，猶之為不足貴，況九九賤技能乎！



嗜好酸鹹，平生每痛自戒勅，竟莫能已，類有物慿之者，吾亦不知其然而然也。



故嘗私為之解曰：由技兼於事者言之，夷之禮，夔之樂，亦不免為一技。由技進乎道者言之，石之斤，扁之輪，非聖人之所與乎。



覽吾之編，察吾苦心，其憫我者當百數，其笑我者當千數。乃若吾之所自得，則自得焉耳，寧復為人憫笑計哉。



\hfill 李冶序

\end{quote}

}



\switchcolumn

{\color{prefacecolor}

\textbf{【白話文】}



\begin{quote}

從前王安石編選《唐百家詩選》，自稱浪費精力於此事，實在可惜。程顥先生把謝良佐的記誦之學視為玩物喪志。文學和史學已經很崇高了，尚且被認為不值得珍視，更何況九九乘法這樣的低賤技能呢！



我一生不斷告誡自己要戒掉這種嗜好（指數學），就像戒掉對酸鹹口味的嗜好一樣，但始終無法停止，彷彿有什麼東西在驅使著我，我也不知道為什麼會這樣。



所以我曾私下為自己辯解說：從「技藝依附於事功」的角度來說，伯夷的禮儀，夔的音樂，也不免是一種技藝。從「由技藝上升到道」的角度來說，匠石的斧頭，輪扁的車輪，不也是聖人所認同的嗎。



閱讀我的這部書，體察我的苦心，同情我的人大概有一百個，嘲笑我的人大概有一千個。至於我從數學中自己所得到的領悟，那就是我自己的收穫了，難道還要為了別人的同情或嘲笑而計較嗎。



\hfill 李冶序

\end{quote}

}

\end{paracol}



\newpage



% =============================================================================

% VOLUME 1

% =============================================================================

\section{卷一 \quad Volume 1}

\label{sec:vol1}



\begin{paracol}{2}

\textbf{【文言文】}卷一為全書之理論基礎，包含四大章節：\textbf{圓城圖式}（主要幾何圖形）；\textbf{總率名號}（十五個直角三角形及其要素之名稱定義）；\textbf{今問正數}（各線段之數值）；\textbf{識別雜記}（含692條幾何公式）。



\switchcolumn

\textbf{【白話文】}第一卷是全書的理論基礎，包含四大部分：\textbf{圓城圖式}（主要幾何圖形）、\textbf{總率名號}（15個直角三角形及其要素的定義）、\textbf{今問正數}（各線段的數值）和\textbf{識別雜記}（包含692條幾何公式）。

\end{paracol}



% -----------------------------------------------------------------------------

\subsection{圓城圖式 Circle City Diagram}

\label{subsec:diagram}



\begin{paracol}{2}

\textbf{【文言文】}圓城圖式為全書170道問題所依據之主要幾何配置。其結構為一大\textbf{勾股形}（直角三角形），內含一內切圓，並有若干特定點與線。



李冶所設之境：「假令有圓城一所，不知周徑，四面開門，門外縱橫各有十字大道。其西北十字道頭定為\textbf{乾}地，其東北十字道頭定為\textbf{艮}地，其東南十字道頭定為\textbf{巽}地，其西南十字道頭定為\textbf{坤}地。所有測望雜法一一設問如後。」



全圖共有\textbf{二十二個命名點}。最大直角三角形之頂點為\textbf{天}、\textbf{地}、\textbf{乾}。內切圓之圓心名為\textbf{心}。各點均以一字命名：天、地、乾、心、日、南、北、川、東、西、艮、坤、山、月、巽、青、泉、朱、金、泛、旦、夕。



\switchcolumn

\textbf{【白話文】}圓城圖式是全書170道問題所依據的主要幾何圖形。其結構是一個大的\textbf{直角三角形}，裡面有一個\textbf{內切圓}，還有一些特定的點和線。



李冶設想的情境是：「假設有一座圓形城池，不知道它的周長和直徑，四面都有城門，每個城門外都有南北和東西向的十字大道。西北方向的十字路口定為\textbf{乾}位，東北定為\textbf{艮}位，東南定為\textbf{巽}位，西南定為\textbf{坤}位。所有測量方法依次設問如下。」



全圖共有\textbf{22個有名字的點}。最大直角三角形的三個頂點是\textbf{天}、\textbf{地}、\textbf{乾}。內切圓的圓心叫做\textbf{心}。每個點都用一個漢字命名：天、地、乾、心、日、南、北、川、東、西、艮、坤、山、月、巽、青、泉、朱、金、泛、旦、夕。

\end{paracol}



\vspace{0.3cm}



\begin{paracol}{2}

\textbf{【文言文】}二十二命名點如下：

\begin{center}

\begin{tabular}{llllllll}

\toprule

天 & 地 & 乾 & 心 & 日 & 南 & 北 & 川 \\

東 & 西 & 艮 & 坤 & 山 & 月 & 巽 & 青 \\

泉 & 朱 & 金 & 泛 & 旦 & 夕 & & \\

\bottomrule

\end{tabular}

\end{center}



\switchcolumn

\textbf{【白話文】}22個命名點如下：

\begin{center}

\begin{tabular}{llllllll}

\toprule

天 & 地 & 乾 & 心 & 日 & 南 & 北 & 川 \\

東 & 西 & 艮 & 坤 & 山 & 月 & 巽 & 青 \\

泉 & 朱 & 金 & 泛 & 旦 & 夕 & & \\

\bottomrule

\end{tabular}

\end{center}

\end{paracol}



% -----------------------------------------------------------------------------

\subsection{總率名號 General Rate Names}

\label{subsec:names}



\begin{paracol}{2}

\textbf{【文言文】}本書研究\textbf{十五個直角三角形}，皆以天地線之一段為其\textbf{弦}（斜邊）。總率名號定義此十五三角形及其三邊。



每個三角形中：\textbf{弦}（$c$）為斜邊；\textbf{股}（$b$）為較長之豎直直角邊；\textbf{勾}（$a$）為較短之水平直角邊。



\switchcolumn

\textbf{【白話文】}本書研究\textbf{15個直角三角形}，它們都以天地線的一段作為\textbf{弦（斜邊）}。總率名號定義了這15個三角形及其三條邊。



在每個三角形中：\textbf{弦}（$c$）是斜邊；\textbf{股}（$b$）是較長的豎直直角邊；\textbf{勾}（$a$）是較短的水平直角邊。

\end{paracol}



\vspace{0.3cm}



\begin{table}[h]

\centering

\caption{十五勾股形總表 Table of 15 Right Triangles}

\begin{tabular}{|c|l|l|c|c|c|}

\hline

\textbf{序} & \textbf{名} & \textbf{頂點} & \textbf{弦} & \textbf{股} & \textbf{勾} \\

\hline

1 & 通（通用） & 天--地--乾 & 通弦 & 通股 & 通勾 \\

2 & 邊（邊上） & 天--西--川 & 邊弦 & 邊股 & 邊勾 \\

3 & 底（底上） & 日--地--北 & 底弦 & 底股 & 底勾 \\

4 & 黃廣 & 天--山--金 & 黃廣弦 & 黃廣股 & 黃廣勾 \\

5 & 黃長 & 月--地--泉 & 黃長弦 & 黃長股 & 黃長勾 \\

6 & 上高 & 天--日--旦 & 上高弦 & 上高股 & 上高勾 \\

7 & 下高 & 日--山--朱 & 下高弦 & 下高股 & 下高勾 \\

8 & 上平 & 月--川--青 & 上平弦 & 上平股 & 上平勾 \\

9 & 下平 & 川--地--夕 & 下平弦 & 下平股 & 下平勾 \\

10 & 大差 & 天--月--坤 & 大差弦 & 大差股 & 大差勾 \\

11 & 小差 & 山--地--艮 & 小差弦 & 小差股 & 小差勾 \\

12 & 皇極 & 日--川--心 & 皇極弦 & 皇極股 & 皇極勾 \\

13 & 太虛 & 月--山--泛 & 太虛弦 & 太虛股 & 太虛勾 \\

14 & 明 & 日--月--南 & 明弦 & 明股 & 明勾 \\

15 & 叀 & 山--川--東 & 叀弦 & 叀股 & 叀勾 \\

\hline

\end{tabular}

\end{table}



\begin{paracol}{2}

\textbf{【文言文】}注：上高=下高，上平=下平，故十五三角形中，實僅十三個不同者。



\switchcolumn

\textbf{【白話文】}注：上高三角形與下高三角形相等，上平三角形與下平三角形相等，所以15個三角形中，實際上只有13個不同的。

\end{paracol}



% -----------------------------------------------------------------------------

\subsection{今問正數 Current Question Numbers}

\label{subsec:numbers}



\begin{paracol}{2}

\textbf{【文言文】}本節給出圖中每一線段之數值。以\textbf{內切圓半徑為120步}為標準單位。



\textbf{1. 通（通用）}



\begin{center}

\begin{tabular}{ll}

弦六百八十 & 勾三百二十 \quad 股六百 \\

勾股和九百二十 & 較二百八十 \\

勾和一千 & 較三百六十 \\

股和一千二百八十 & 較八十 \\

較和九百六十 & 較四百 \\

和和一千六百 & 較二百四十 \\

\end{tabular}

\end{center}



\switchcolumn

\textbf{【白話文】}這一節給出圖中每一條線段的數值。以\textbf{內切圓的半徑為120步}作為標準單位。



\textbf{1. 通三角形（通用三角形）}



\begin{center}

\begin{tabular}{ll}

弦（斜邊）= 680 & 勾（短邊）= 320 \quad 股（長邊）= 600 \\

勾股之和 = 920 & 勾股之差 = 280 \\

勾弦之和 = 1000 & 勾弦之差 = 360 \\

股弦之和 = 1280 & 股弦之差 = 80 \\

弦較之和 = 960 & 弦較之差 = 400 \\

弦和之和 = 1600 & 弦和之差 = 240 \\

\end{tabular}

\end{center}

\end{paracol}



\begin{paracol}{2}

\textbf{【文言文】}



\textbf{2. 邊（邊上）}



\begin{center}

\begin{tabular}{ll}

弦五百四十四 & 勾二百五十六 \quad 股四百八十 \\

勾股和七百三十六 & 較二百二十四 \\

勾和八百 & 較二百八十八 \\

股和一千零二十四 & 較六十四 \\

較和七百六十八 & 較三百二十 \\

和和一千二百八十 & 較一百九十二 \\

\end{tabular}

\end{center}



\switchcolumn

\textbf{【白話文】}



\textbf{2. 邊三角形（邊上三角形）}



\begin{center}

\begin{tabular}{ll}

弦 = 544 & 勾 = 256 \quad 股 = 480 \\

勾股和 = 736 & 差 = 224 \\

勾弦和 = 800 & 差 = 288 \\

股弦和 = 1024 & 差 = 64 \\

弦較和 = 768 & 差 = 320 \\

弦和和 = 1280 & 差 = 192 \\

\end{tabular}

\end{center}

\end{paracol}



\begin{paracol}{2}

\textbf{【文言文】}



\textbf{3. 底（底上）}



\begin{center}

\begin{tabular}{ll}

弦四百二十五 & 勾二百 \quad 股三百七十五 \\

勾股和五百七十五 & 較一百七十五 \\

勾和六百二十五 & 較二百二十五 \\

股和八百 & 較五十 \\

較和六百 & 較二百五十 \\

和和一千 & 較一百五十 \\

\end{tabular}

\end{center}



\switchcolumn

\textbf{【白話文】}



\textbf{3. 底三角形（底上三角形）}



\begin{center}

\begin{tabular}{ll}

弦 = 425 & 勾 = 200 \quad 股 = 375 \\

勾股和 = 575 & 差 = 175 \\

勾弦和 = 625 & 差 = 225 \\

股弦和 = 800 & 差 = 50 \\

弦較和 = 600 & 差 = 250 \\

弦和和 = 1000 & 差 = 150 \\

\end{tabular}

\end{center}

\end{paracol}



\begin{paracol}{2}

\textbf{【文言文】}



\textbf{4. 黃廣}



\begin{center}

\begin{tabular}{ll}

弦五百一十 & 勾二百四十 \quad 股四百五十 \\

\multicolumn{2}{c}{（即城徑也）} \\

勾股和六百九十 & 較二百一十 \\

\end{tabular}

\end{center}



其餘十一三角形依同樣格式，每形含十三個數量（弦、勾、股、勾股和、勾股較、勾弦和、勾弦較、股弦和、股弦較、弦較和、弦較較、弦和和、弦和較），共計 $15 \times 13 = 195$ 個數值。



\switchcolumn

\textbf{【白話文】}



\textbf{4. 黃廣三角形}



\begin{center}

\begin{tabular}{ll}

弦 = 510 & 勾 = 240 \quad 股 = 450 \\

\multicolumn{2}{c}{（這就是城池的直徑）} \\

勾股和 = 690 & 差 = 210 \\

\end{tabular}

\end{center}



其餘11個三角形按照同樣的格式，每個三角形包含13個數量（弦、勾、股、勾股和、勾股差、勾弦和、勾弦差、股弦和、股弦差、弦較和、弦較差、弦和和、弦和差），共計 $15 \times 13 = 195$ 個數值。

\end{paracol}



% -----------------------------------------------------------------------------

\subsection{識別雜記 Identification Miscellaneous Records}

\label{subsec:identification}



\begin{paracol}{2}

\textbf{【文言文】}此為全書之理論核心，含\textbf{692條幾何公式}，關聯圖中各線段。此等公式純為幾何定理，不依賴任何特定數值。其後各卷之問題，皆用此等公式配合\textbf{天元術}（設未知數列方程之法）解之。



本節分為八大部分：\textbf{諸雜名目}（一般定義及三十餘條定理）；\textbf{五和五較}（五種和與五種差）；\textbf{諸弦}（各種弦）；\textbf{大小差}（大差與小差）；\textbf{諸差}（各種差）；\textbf{諸率互見}（各率之間的相互關係）；\textbf{四位拾遺}（四位拾遺）；\textbf{拾遺}（拾遺）。



\switchcolumn

\textbf{【白話文】}這是全書的理論核心部分，包含\textbf{692條幾何公式}，用來關聯圖中各個線段。這些公式純粹是幾何定理，不依賴任何特定的數值。後面各卷的問題，都是用這些公式配合\textbf{天元術（設未知數列方程的方法）}來解決的。



這一節分為八大部分：\textbf{諸雜名目}（一般定義及30多條定理）、\textbf{五和五較}（五種和與五種差）、\textbf{諸弦}（各種弦）、\textbf{大小差}（大差與小差）、\textbf{諸差}（各種差）、\textbf{諸率互見}（各個比率之間的相互關係）、\textbf{四位拾遺}（四位補遺）、\textbf{拾遺}（補遺）。

\end{paracol}



\vspace{0.3cm}



\begin{paracol}{2}

\textbf{【文言文】}\textbf{諸雜名目}選例：



\begin{center}

\begin{tabular}{ll}

\toprule

\textbf{名} & \textbf{定義} \\

\midrule

內率 & 明勾 $\times$ 叀股 \\

外率 & 明股 $\times$ 叀勾 \\

虛率 & 虛勾 $\times$ 虛股 \\

角差 & 高股 $-$ 平勾 \\

次差 & (明股$-$明勾) $+$ (叀股$-$叀勾) \\

混同和 & 黃長股 $+$ 黃廣勾 \\

傍差 & (明股$-$明勾) $-$ (叀股$-$叀勾) \\

\bottomrule

\end{tabular}

\end{center}



\switchcolumn

\textbf{【白話文】}\textbf{諸雜名目}部分舉例：



\begin{center}

\begin{tabular}{ll}

\toprule

\textbf{名稱} & \textbf{定義} \\

\midrule

內率 & 明三角形的勾 $\times$ 叀三角形的股 \\

外率 & 明三角形的股 $\times$ 叀三角形的勾 \\

虛率 & 虛三角形的勾 $\times$ 虛三角形的股 \\

角差 & 高三角形的股 $-$ 平三角形的勾 \\

次差 & (明股$-$明勾) $+$ (叀股$-$叀勾) \\

混同和 & 黃長三角形的股 $+$ 黃廣三角形的勾 \\

傍差 & (明股$-$明勾) $-$ (叀股$-$叀勾) \\

\bottomrule

\end{tabular}

\end{center}

\end{paracol}



\vspace{0.3cm}



\begin{paracol}{2}

\textbf{【文言文】}\textbf{定理選例：}



\begin{itemize}[leftmargin=1.5em]

\item 凡大差小差相乘為半段徑\textbf{冪}（乘積/面積）。

\item 大差勾小差股相乘同上。

\item 黃廣股黃長勾相乘為經\textbf{冪}。

\item 勾股和即弦黃和。\newline 即：$a + b = c + d$（$d$為內切圓直徑）。

\end{itemize}



\switchcolumn

\textbf{【白話文】}\textbf{定理舉例：}



\begin{itemize}[leftmargin=1.5em]

\item 大差與小差相乘，等於直徑平方的一半。

\item 大差三角形的勾與小差三角形的股相乘，結果同上。

\item 黃廣三角形的股與黃長三角形的勾相乘，等於直徑的平方。

\item 勾與股之和等於弦與直徑之和。\newline 即：$a + b = c + d$（$d$為內切圓直徑）。

\end{itemize}

\end{paracol}



\newpage



% =============================================================================

% VOLUME 2

% =============================================================================

\section{卷二 \quad Volume 2}

\label{sec:vol2}



\begin{paracol}{2}

\textbf{【文言文】}\textbf{正率一十四問}——十四道標準比率問題。



此十四道問題構成「洞淵九容」之序列，源自道家學者李思聰（號洞淵先生）之先前著作。其呈現直角三角形所含或相關之九種基本圓類型，再加五道補充問題。



\switchcolumn

\textbf{【白話文】}\textbf{正率十四問}——十四道標準比率的問題。



這十四道問題構成了「洞淵九容」的序列，源自道家學者李思聰（號洞淵先生）先前的著作。它們呈現了與直角三角形相關的九種基本圓的類型，再加上五道補充問題。

\end{paracol}



\vspace{0.5cm}



% -----------------------------------------------------------------------------

\subsubsection*{第一問 \quad Problem 1}



\begin{paracol}{2}

\wen{甲乙二人俱在乾地，乙東行三百二十步而立，甲南行六百步望見乙，問徑幾里？}



\da{城徑二百四十步。}



\shu{此為\textbf{勾股容圓}（直角三角形內切圓）也。以勾股相乘，倍之為實，並勾股\textbf{冪}（平方）以求\textbf{弦}（斜邊），復加入勾股共，以為法。}



\cao{置甲南行六百步在地，以乙東行三百二十步乘之，得一十九萬二千步，倍之得三十入萬四千步為實。以乙東行步自之，得一十萬零二千四百步為勾\textbf{冪}（平方）；以甲南行步自之，得三十六萬步為股\textbf{冪}（平方）。二\textbf{冪}相並，得四十六萬二千四百步為方實，以平方開之，得六百八十步，則\textbf{弦}（斜邊）也。以\textbf{弦}加勾股共，共得一千六百步以為法。如法而一，得二百四十步，則城徑也。\textbf{合問}（答案正確）。}



\switchcolumn

\wenm{甲乙兩人都在乾地，乙向東走320步後停下，甲向南走600步就看見了乙。問圓形城池的直徑是多少？}



\dam{城池的直徑是240步。}



\shum{這是\textbf{勾股容圓（直角三角形內切圓）}的問題。將勾和股相乘，再乘以2作為被除數（實）；將勾的平方和股的平方相加後開方求弦（斜邊），再將弦與勾股之和相加，作為除數（法）。}



\caom{將甲向南走的600步放在地上，用乙向東走的320步去乘它，得到192000步，再乘以2得到384000步，這就是被除數（實）。用乙向東走的步數自乘，得到102400步，這是勾的平方；用甲向南走的步數自乘，得到360000步，這是股的平方。兩個平方相加，得到462400步，開平方得到680步，這就是弦（斜邊）。用弦加上勾和股的和，共得1600步作為除數（法）。用被除數除以除數，得到240步，這就是城池的直徑。\textbf{答案正確。}}

\end{paracol}



\vspace{0.5cm}

\noindent\textcolor{rulecolor}{\rule{\textwidth}{0.4pt}}



% -----------------------------------------------------------------------------

\subsubsection*{第二問 \quad Problem 2}



\begin{paracol}{2}

\wen{甲乙二人俱在西門，乙東行二百五十六步，甲南行四百八十步望見乙，問答同前。}



\da{城徑二百四十步。}



\shu{此為\textbf{勾上容圓}也。以勾股相乘，倍之為實，並勾股\textbf{冪}（平方）以求\textbf{弦}（斜邊），加入股以為法。}



\cao{置甲南行四百八十步在地，以乙東行二百五十六步乘之，得一十二萬二千八百八十步，倍之得二十四萬五千七百六十步為實。以乙東行步自之，得六萬五千五百三十六步為勾\textbf{冪}；以甲南行步自之，得二十三萬零四百步為股\textbf{冪}。勾股\textbf{冪}相並，得二十九萬五千九百三十六步為方實，以平方開之，得五百四十四步為\textbf{弦}（斜邊）也。以\textbf{弦}加入南行步，共得一千零二十四步以為法。而一，得二百四十步則城徑。\textbf{合問}（答案正確）。}



\switchcolumn

\wenm{甲乙兩人都在西門，乙向東走256步，甲向南走480步就看見了乙。問題和答案與前面相同。}



\dam{城池的直徑是240步。}



\shum{這是\textbf{勾上容圓}的問題。將勾和股相乘，再乘以2作為被除數（實）；將勾的平方和股的平方相加後開方求弦（斜邊），再將弦與股相加，作為除數（法）。}



\caom{將甲向南走的480步放在地上，用乙向東走的256步去乘它，得到122880步，再乘以2得到245760步，這就是被除數（實）。用乙向東走的步數自乘，得到65536步，這是勾的平方；用甲向南走的步數自乘，得到230400步，這是股的平方。勾和股的平方相加，得到295936步，開平方得到544步，這就是弦（斜邊）。用弦加上向南走的步數，共得1024步作為除數（法）。除法運算，得到240步，這就是城池的直徑。\textbf{答案正確。}}

\end{paracol}



\vspace{0.5cm}

\noindent\textcolor{rulecolor}{\rule{\textwidth}{0.4pt}}



% -----------------------------------------------------------------------------

\subsubsection*{第三問 \quad Problem 3}



\begin{paracol}{2}

\wen{甲乙二人俱在北門，乙東行二百步而止，甲南行三百七十五步望見乙，問答同前。}



\da{城徑二百四十步。}



\shu{此為\textbf{股上容圓}也。以勾股相乘，倍之為實，以勾股\textbf{冪}（平方）求\textbf{弦}（斜邊），加入勾以為法。}



\cao{置甲南行三百七十五步，以乙東行二百步乘之，得七萬五千步，倍之得一十五萬步為實。以乙東行自之，得四萬步為勾\textbf{冪}（平方）；以甲南行自之，得一十四萬零六百二十五步為股\textbf{冪}（平方）。勾股\textbf{冪}相並，得一十八萬零六百二十五步為方實。如平方而一，得四百二十五步則\textbf{弦}（斜邊）也。加入乙東行二百步，共得六百二十五步以為法。以法除之，得二百四十步，則城徑也。\textbf{合問}（答案正確）。}



\switchcolumn

\wenm{甲乙兩人都在北門，乙向東走200步後停下，甲向南走375步就看見了乙。問題和答案與前面相同。}



\dam{城池的直徑是240步。}



\shum{這是\textbf{股上容圓}的問題。將勾和股相乘，再乘以2作為被除數（實）；將勾的平方和股的平方相加後開方求弦（斜邊），再將弦與勾相加，作為除數（法）。}



\caom{將甲向南走的375步，用乙向東走的200步去乘，得到75000步，再乘以2得到150000步，這就是被除數（實）。用乙向東走的步數自乘，得到40000步，這是勾的平方；用甲向南走的步數自乘，得到140625步，這是股的平方。勾和股的平方相加，得到180625步，開平方得到425步，這就是弦（斜邊）。加上乙向東走的200步，共得625步作為除數（法）。用被除數除以除數，得到240步，這就是城池的直徑。\textbf{答案正確。}}

\end{paracol}



\vspace{0.5cm}

\noindent\textcolor{rulecolor}{\rule{\textwidth}{0.4pt}}



% -----------------------------------------------------------------------------

\subsubsection*{第四問 \quad Problem 4}



\begin{paracol}{2}

\wen{甲乙二人俱在圓城中心而立，乙穿城向東行一百三十六步而止，甲穿城南行二百五十五步望見乙，問答同前。}



\da{城徑二百四十步。}



\shu{此為\textbf{勾股上容圓}也。以勾股相乘，倍之為實，並勾股\textbf{冪}（平方）如法求\textbf{弦}（斜邊），以為法。}



\cao{以二行步相乘，得三萬四千六百八十步，倍之得六萬九千三百六十步為實。置乙東行自之，得一萬八千四百九十六步為勾\textbf{冪}（平方）；又以甲南行自之，得六萬五千零二十五步為股\textbf{冪}（平方）。二\textbf{冪}相並，得八萬三千五百二十一步為方實，以平方開之，得二百八十九步即\textbf{弦}（斜邊）也。便以為法。如法除實，得二百四十步，即圓城之徑也。\textbf{合問}（答案正確）。}



\switchcolumn

\wenm{甲乙兩人都在圓形城池的中心，乙穿過城池向東走136步後停下，甲穿過城池向南走255步就看見了乙。問題和答案與前面相同。}



\dam{城池的直徑是240步。}



\shum{這是\textbf{勾股上容圓}的問題。將勾和股相乘，再乘以2作為被除數（實）；將勾的平方和股的平方相加後開方求弦（斜邊），以弦作為除數（法）。}



\caom{用兩個行走步數相乘，得到34680步，再乘以2得到69360步，這就是被除數（實）。用乙向東走的步數自乘，得到18496步，這是勾的平方；用甲向南走的步數自乘，得到65025步，這是股的平方。兩個平方相加，得到83521步，開平方得到289步，這就是弦（斜邊）。就以弦作為除數（法）。用被除數除以除數，得到240步，這就是圓形城池的直徑。\textbf{答案正確。}}

\end{paracol}



\vspace{0.5cm}

\noindent\textcolor{rulecolor}{\rule{\textwidth}{0.4pt}}



% -----------------------------------------------------------------------------

\subsubsection*{第五問 \quad Problem 5}



\begin{paracol}{2}

\wen{甲乙二人同立於乾地，乙東行一百八十步遇塔而止，甲南行三百六十步，回望其塔正居城徑之半，問答同前。}



\da{城徑二百四十步。}



\shu{此為\textbf{弦上容圓}也。以勾股相乘，倍之為實，以勾股和為法。}



\cao{以二行步相乘，得六萬四千八百步，倍之得一十二萬九千六百步為實。並二行步，得五百四十步以為法。除實，得二百四十步，即城徑也。\textbf{合問}（答案正確）。}



\switchcolumn

\wenm{甲乙兩人同在乾地，乙向東走180步到達一座塔後停下，甲向南走360步，回頭看見塔正好在城池直徑的中間位置。問題和答案與前面相同。}



\dam{城池的直徑是240步。}



\shum{這是\textbf{弦上容圓}的問題。將勾和股相乘，再乘以2作為被除數（實），以勾與股的和作為除數（法）。}



\caom{用兩個行走步數相乘，得到64800步，再乘以2得到129600步，這就是被除數（實）。將兩個行走步數相加，得到540步作為除數（法）。除法運算，得到240步，這就是城池的直徑。\textbf{答案正確。}}

\end{paracol}



\vspace{0.5cm}

\noindent\textcolor{rulecolor}{\rule{\textwidth}{0.4pt}}



% -----------------------------------------------------------------------------

\subsubsection*{第六問 \quad Problem 6}



\begin{paracol}{2}

\wen{甲乙二人俱在坤地，乙東行一百九十二步而止，甲南行三百六十步，望乙與城參相直，問答同前。}



\da{城徑二百四十步。}



\shu{此為\textbf{勾外容圓}也。以勾股相乘，倍之為實，以\textbf{弦較和}為法。}



\cao{以二行步相乘，得六萬九千一百二十步，倍之得一十三萬八千二百四十步為實。置乙東行自之，得三萬六千八百六十四步為勾\textbf{冪}（平方）；又置甲南行自之，得一十二萬九千六百步為股\textbf{冪}（平方）。二\textbf{冪}相並，得一十六萬六千四百六十四步為方實，以平方開之，得四百零八，即\textbf{弦}（斜邊）也。又置甲南行步內減乙東行步，餘一百六十八步，即\textbf{較}（差）也。以\textbf{較}加\textbf{弦}，共得五百七十六步以為法。實如法而一，得二百四十步，為城徑也。\textbf{合問}（答案正確）。}



\an{此題用勾股求得弦，即可加減得弦較較為城徑。今必以勾股相乘倍積為實，求得弦較和為法，而後始得弦較較為城徑者，蓋欲因此並明勾股相乘之倍積為弦較較、弦較和相乘之積，非故為紆回也。}



\switchcolumn

\wenm{甲乙兩人都在坤地，乙向東走192步後停下，甲向南走360步，看見乙和城池在同一條直線上。問題和答案與前面相同。}



\dam{城池的直徑是240步。}



\shum{這是\textbf{勾外容圓}的問題。將勾和股相乘，再乘以2作為被除數（實），以\textbf{弦與勾股差的和}作為除數（法）。}



\caom{用兩個行走步數相乘，得到69120步，再乘以2得到138240步，這就是被除數（實）。用乙向東走的步數自乘，得到36864步，這是勾的平方；用甲向南走的步數自乘，得到129600步，這是股的平方。兩個平方相加，得到166464步，開平方得到408，這就是弦（斜邊）。再用甲向南走的步數減去乙向東走的步數，餘168步，這就是\textbf{差}。用\textbf{差}加上\textbf{弦}，共得576步作為除數（法）。除法運算，得到240步，這就是城池的直徑。\textbf{答案正確。}}



\anm{這道題用勾股定理求出弦之後，就可以加減得到弦較差作為城徑。但現在一定要用勾股相乘的兩倍作為被除數，求得弦較和作為除數，然後才能得到弦較差作為城徑，這是因為想要同時說明：勾股相乘的兩倍等於弦較差與弦較和的乘積，並不是故意繞彎路。}

\end{paracol}



\vspace{0.5cm}

\noindent\textcolor{rulecolor}{\rule{\textwidth}{0.4pt}}



% -----------------------------------------------------------------------------

\subsubsection*{第七問 \quad Problem 7}



\begin{paracol}{2}

\wen{甲乙二人同立於艮地，甲南行一百五十步而止，乙東行八十步，望乙與城參相直，問答同前。}



\da{城徑二百四十步。}



\shu{此為\textbf{股外容圓}也。以勾股相乘，倍之為實，以\textbf{弦和較}為法。}



\switchcolumn

\wenm{甲乙兩人同在艮地，甲向南走150步後停下，乙向東走80步，看見乙和城池在同一條直線上。問題和答案與前面相同。}



\dam{城池的直徑是240步。}



\shum{這是\textbf{股外容圓}的問題。將勾和股相乘，再乘以2作為被除數（實），以\textbf{弦與勾股和的差}作為除數（法）。}

\end{paracol}



\vspace{0.5cm}

\noindent\textcolor{rulecolor}{\rule{\textwidth}{0.4pt}}



% -----------------------------------------------------------------------------

\subsubsection*{第八問 \quad Problem 8}



\begin{paracol}{2}

\wen{甲乙二人同立於坤地，甲南行一百九十二步而止，乙東行七十二步，望乙與城參相直，問答同前。}



\da{城徑二百四十步。}



\shu{此為\textbf{弦外容圓}也。以勾股相乘，倍之為實，以\textbf{勾股較}（勾股之差）為法。}



\switchcolumn

\wenm{甲乙兩人同在坤地，甲向南走192步後停下，乙向東走72步，看見乙和城池在同一條直線上。問題和答案與前面相同。}



\dam{城池的直徑是240步。}



\shum{這是\textbf{弦外容圓}的問題。將勾和股相乘，再乘以2作為被除數（實），以\textbf{勾與股的差}作為除數（法）。}

\end{paracol}



\vspace{0.5cm}

\noindent\textcolor{rulecolor}{\rule{\textwidth}{0.4pt}}



% -----------------------------------------------------------------------------

\subsubsection*{第九問 \quad Problem 9}



\begin{paracol}{2}

\wen{甲乙二人俱在巽地，乙東行一百九十二步而止，甲南行三百六十步，望乙與城參相直，問答同前。}



\da{城徑二百四十步。}



\shu{此為\textbf{勾外容半圓}也。以勾股相乘，倍之為實，以\textbf{股弦和}（股與弦之和）為法。}



\switchcolumn

\wenm{甲乙兩人都在巽地，乙向東走192步後停下，甲向南走360步，看見乙和城池在同一條直線上。問題和答案與前面相同。}



\dam{城池的直徑是240步。}



\shum{這是\textbf{勾外容半圓}的問題。將勾和股相乘，再乘以2作為被除數（實），以\textbf{股與弦的和}作為除數（法）。}

\end{paracol}



\vspace{0.5cm}

\noindent\textcolor{rulecolor}{\rule{\textwidth}{0.4pt}}



% -----------------------------------------------------------------------------

\subsubsection*{第十問 \quad Problem 10}



\begin{paracol}{2}

\wen{甲乙二人俱在艮地，甲南行一百五十步而止，乙東行八十步，望乙與城參相直，問答同前。}



\da{城徑二百四十步。}



\shu{此為\textbf{股外容半圓}也。以勾股相乘，倍之為實，以\textbf{勾弦和}（勾與弦之和）為法。}



\switchcolumn

\wenm{甲乙兩人都在艮地，甲向南走150步後停下，乙向東走80步，看見乙和城池在同一條直線上。問題和答案與前面相同。}



\dam{城池的直徑是240步。}



\shum{這是\textbf{股外容半圓}的問題。將勾和股相乘，再乘以2作為被除數（實），以\textbf{勾與弦的和}作為除數（法）。}

\end{paracol}



\vspace{0.5cm}

\noindent\textcolor{rulecolor}{\rule{\textwidth}{0.4pt}}



\begin{paracol}{2}

\textit{以上十問即「洞淵九容」，第五問（弦上容圓）為補充。卷二其餘四問續以天元術解之。}



\switchcolumn

\textit{以上十道問題構成「洞淵九容」，第五問（弦上容圓）為補充問題。卷二其餘四道問題繼續使用天元術（設未知數列方程的方法）來求解。}

\end{paracol}



\newpage



% =============================================================================

% VOLUME 3

% =============================================================================

\section{卷三 \quad Volume 3}

\label{sec:vol3}



\begin{paracol}{2}

\textbf{【文言文】}\textbf{邊股一十七問}——十七道邊股問題。



卷三呈現17道問題，其中所給數據涉及\textbf{邊股}（$b_2$），即邊三角形（表中第二個三角形）之長直角邊，定為480步。此卷所有問題所求之未知數皆為圓城之直徑（城徑 = 240步）。



此等問題展示\textbf{天元術}（設未知數列方程之法）之漸增精妙，方程次數自二次至三次不等。



\switchcolumn

\textbf{【白話文】}\textbf{邊股十七問}——十七道與邊股相關的問題。



第三卷呈現17道問題，題目中給出的數據涉及\textbf{邊股}（$b_2$），即邊三角形（表中的第2個三角形）的長直角邊，數值定為480步。這一卷所有問題要求解的未知數都是圓形城池的直徑（城徑 = 240步）。



這些問題展示了\textbf{天元術（設未知數列方程的方法）}的漸進精妙之處，方程的次數從二次到三次不等。

\end{paracol}



\vspace{0.5cm}



% -----------------------------------------------------------------------------

\subsubsection*{第一問 \quad Problem 1}



\begin{paracol}{2}

\wen{乙出東門南行，不知步數而止。甲出西門南行四百八十步，望見乙，復就乙行五百一十步與乙相會，問答同前。}



\da{城徑二百四十步。}



\shu{倍相減步，以乘二之甲南行步，為平方實，得城徑。}



\cao{識別得二行相減，餘三十步，即乙出東門南行步也。倍相減步，得六十步，以乘二之甲南行步九百六十步，得五萬七千六百步，為平方實。如法開之，得二百四十步，即城徑也。\textbf{合問}（答案正確）。}



\switchcolumn

\wenm{乙從東門出去向南走，不知道走了多少步後停下。甲從西門出去向南走480步，看見了乙，再向乙走去，走了510步與乙會合。問題和答案與前面相同。}



\dam{城池的直徑是240步。}



\shum{將兩個行走步數的差乘以2，再乘以甲向南走步數的2倍，作為開平方的被開方數，開方即得城徑。}



\caom{識別出兩個行走步數相減，餘30步，這就是乙從東門向南走的步數。將相差的步數乘以2，得到60步，再乘以甲向南走步數的2倍（960步），得到57600步，這就是開平方的被開方數。按照方法開平方，得到240步，這就是城池的直徑。\textbf{答案正確。}}

\end{paracol}



\vspace{0.5cm}

\noindent\textcolor{rulecolor}{\rule{\textwidth}{0.4pt}}



% -----------------------------------------------------------------------------

\subsubsection*{第二問 \quad Problem 2}



\begin{paracol}{2}

\wen{甲出西門南行四百八十步而止，乙從艮隅東行八十步，望見甲，問答同前。}



\da{城徑二百四十步。}



\shu{倍南行步，以東行步乘之為實，東行步為從方，一步常法，得全徑。}



\cao{\textbf{立天元一為全徑}，以減於二之甲南行步，得為兩個大差也。以乙東行步乘之，得為圓徑\textbf{冪}（平方），寄左。然後以天元\textbf{冪}（平方）與左相消，得以帶縱平方開之，得二百四十步，即城徑也。\textbf{合問}（答案正確）。}



\vspace{0.2cm}

\noindent\textbf{又法：}半之乙東行步，乘南行步為實，半之乙東行步為從，一步常法，得半徑。



\cao{\textbf{立天元一為半城徑}，減甲南行步，得為大差也。以半之東行步乘之，得即半徑\textbf{冪}（平方），寄左。然後以天元\textbf{冪}（平方）為同數，與左相消，得開帶縱平方，得一百二十步。倍之，即城徑也。\textbf{合問}（答案正確）。}



\switchcolumn

\wenm{甲從西門出去向南走480步後停下，乙從艮隅向東走80步，看見了甲。問題和答案與前面相同。}



\dam{城池的直徑是240步。}



\shum{將向南走的步數乘以2，用向東走的步數去乘它作為被除數（實），以向東走的步數作為一次項系數（從方），二次項系數為1（常法），求得全徑。}



\caom{\textbf{設天元（未知數$x$）為全徑}，用兩倍的甲向南走的步數減去它，得到兩個大差。用乙向東走的步數去乘它，得到圓徑的\textbf{平方}，暫時寄放在左邊。然後用天元的\textbf{平方}與左邊的式子相消，開帶縱平方，得到240步，這就是城池的直徑。\textbf{答案正確。}}



\vspace{0.2cm}

\noindent\textbf{【另一種解法】}將乙向東走的步數除以2，乘以向南走的步數作為被除數（實），將乙向東走步數的一半作為一次項系數（從），二次項系數為1，求得半徑。



\caom{\textbf{設天元（未知數$x$）為半城徑}，用甲向南走的步數減去它，得到大差。用向東走步數的一半去乘它，得到半徑的\textbf{平方}，暫時寄放在左邊。然後用天元的\textbf{平方}作為同數，與左邊的式子相消，開帶縱平方，得到120步。乘以2，就是城池的直徑。\textbf{答案正確。}}

\end{paracol}



\vspace{0.5cm}

\noindent\textcolor{rulecolor}{\rule{\textwidth}{0.4pt}}



% -----------------------------------------------------------------------------

\subsubsection*{第三問 \quad Problem 3}



\begin{paracol}{2}

\wen{甲出西門南行四百八十步而止，乙從艮隅亦南行一百五十步，望見甲，問答同前。}



\da{城徑二百四十步。}



\shu{兩行步相乘為實，南行步為從方，一為隅，得半徑。}



\cao{\textbf{立天元一為半城徑}，以減乙南行步，得為半梯頭；以甲行步為梯底，以乘之，得為半徑\textbf{冪}（平方），寄左。然後以天元\textbf{冪}（平方）與左相消，得開帶縱平方，得一百二十步。倍之，即城徑也。\textbf{合問}（答案正確）。}



\switchcolumn

\wenm{甲從西門出去向南走480步後停下，乙從艮隅也向南走150步，看見了甲。問題和答案與前面相同。}



\dam{城池的直徑是240步。}



\shum{兩個向南走的步數相乘作為被除數（實），甲向南走的步數作為一次項系數（從方），二次項系數為1（隅），求得半徑。}



\caom{\textbf{設天元（未知數$x$）為半城徑}，用乙向南走的步數減去它，得到半個梯頭；用甲走的步數作為梯底，去乘它，得到半徑的\textbf{平方}，暫時寄放在左邊。然後用天元的\textbf{平方}與左邊的式子相消，開帶縱平方，得到120步。乘以2，就是城池的直徑。\textbf{答案正確。}}

\end{paracol}



\vspace{0.5cm}

\noindent\textcolor{rulecolor}{\rule{\textwidth}{0.4pt}}



% -----------------------------------------------------------------------------

\subsubsection*{第四問 \quad Problem 4}



\begin{paracol}{2}

\wen{甲出西門南行四百八十步，乙出東門直行一十六步，望見甲，問答同前。}



\da{城徑二百四十步。}



\shu{以四之東行步，乘南行\textbf{冪}（平方）為實，從空，東行為廉，一步為隅法，得全徑。}



\cao{\textbf{立天元一為圓徑}，加乙東行步，得為中勾；其甲南行即中股也。置東行步為小勾，以中股乘之，得，合以中勾除，今\textbf{不受除}，便以為小股也。內寄中勾分母。乃復以中股乘之，得三百六十八萬六千四百，又四之，得一千四百七十四萬五千六百，為一段圓徑\textbf{冪}（平方），寄中勾分母，寄左。然後以天元徑自之，又以中勾乘之，得為同數，與左相消，得以帶縱立方開之，得二百四十步，為城徑也。\textbf{合問}（答案正確）。}



\an{不受除者，無可除之理也。凡二數，此數於彼數有可除之理，則受除；無可除之理，則不受除也。蓋除有法有實，實可二，法不可二。此題以中勾為法，而中勾內有一元，又有十六步，其為數已二矣，又何以均分不一之數乎？故曰不受也。寄分者，姑寄其應除之數也。俟求得兩相等數，而此數內尚少一除，不除此而轉乘彼，則兩數仍相等，猶之受除者也。此所謂以乘代除也。}



\switchcolumn

\wenm{甲從西門出去向南走480步，乙從東門出去直走16步，看見了甲。問題和答案與前面相同。}



\dam{城池的直徑是240步。}



\shum{將向東走的步數乘以4，再乘以向南走步數的\textbf{平方}作為被除數（實），一次項系數為0（從空），以向東走的步數作為二次項系數（廉），三次項系數為1（隅法），求得全徑。}



\caom{\textbf{設天元（未知數$x$）為圓的直徑}，加上乙向東走的步數，得到中勾；甲向南走的步數就是中股。將向東走的步數作為小勾，用中股去乘它，得到……應該用中勾去除，現在\textbf{不接受除法}，就把它當作小股。裡面寄放著中勾作為分母。再用中股去乘它，得到3686400，再乘以4，得到14745600，作為圓徑\textbf{平方}的一部分，寄放著中勾分母，暫時放在左邊。然後用天元（直徑）自乘，再用中勾去乘它，得到同數，與左邊的式子相消，開帶縱立方，得到240步，這就是城池的直徑。\textbf{答案正確。}}



\anm{「不接受除法」的意思是，沒有辦法做除法。一般來說，兩個數之間，如果這個數能被那個數除盡，就接受除法；如果不能除盡，就不接受除法。因為除法有除數和被除數，被除數可以是多個數，但除數不能是多個數。這道題以中勾作為除數，而中勾裡面有一個未知數（天元）還有16步，它已經是兩個數了，怎麼能用一個不統一的數來做除法呢？所以說不接受除法。寄放分母的意思是，暫時寄放應該除的那個數。等到求出兩個相等的數，而這個數裡面還少一次除法，不做這個除法反而去乘那個數，那麼兩個數仍然相等，就好像接受了除法一樣。這就是所謂的「以乘法代替除法」。}

\end{paracol}



\vspace{0.5cm}

\noindent\textcolor{rulecolor}{\rule{\textwidth}{0.4pt}}



% -----------------------------------------------------------------------------

\subsubsection*{第五問 \quad Problem 5}



\begin{paracol}{2}

\wen{乙出南門東行七十二步而止，甲出西門南行四百八十步，望乙與城參相直，問答同前。}



\da{城徑二百四十步。}



\shu{以乙東行\textbf{冪}（平方），乘甲南行為實；乙東行\textbf{冪}（平方）為從方；甲南行步內減二之東行步為益廉；一步常法，得半徑。}



\cao{\textbf{立天元一為半城徑}，以減南行步，得為小股。又以天元加乙東行步，得為小勾。又以天元加南行步，得為大股。乃置大股在地，以小勾乘之，得下式，合以小股除之，今不受除，便以為大勾，內寄小股分母。又置天元半徑，以分母小股乘之，得以減大勾，得為半個梯底於上。以乙東行七十二步為半個梯頭，以乘上位，得為半徑\textbf{冪}（平方），內寄小股分母，寄左。然後置天元\textbf{冪}（平方），又以分母小股乘之，得為同數，與左相消，以立方開之，得一百二十步。倍之，即城徑也。\textbf{合問}（答案正確）。}



\vspace{0.2cm}

\noindent\textbf{又法：}以二數相乘為實，相減為從，一虛法，平開，得半徑。



\cao{別得二數相並，為大股內少一虛勾；其二數相減，為大差也。\textbf{立天元一為半徑}，副置之上位，減於四百八十，得為股圓差，即大差股也。下位加七十二，得，與股圓差相乘，得為一大差積，寄左。再以大差勾減於大差股，餘為較，又加入大差四百單八，共得為較共也。以天元乘之，得為同數，與左相消，以平方開之，得一百二十步，即半徑。\textbf{合問}（答案正確）。前法太煩，故又立此法以就簡也。}



\switchcolumn

\wenm{乙從南門出去向東走72步後停下，甲從西門出去向南走480步，看見乙和城池在同一條直線上。問題和答案與前面相同。}



\dam{城池的直徑是240步。}



\shum{用乙向東走步數的\textbf{平方}，乘以甲向南走的步數作為被除數（實）；乙向東走步數的\textbf{平方}作為一次項系數（從方）；甲向南走的步數減去兩倍向東走的步數作為負的二次項系數（益廉）；三次項系數為1（常法），求得半徑。}



\caom{\textbf{設天元（未知數$x$）為半城徑}，用向南走的步數減去它，得到小股。再用天元加上乙向東走的步數，得到小勾。再用天元加上向南走的步數，得到大股。將大股放在地上，用小勾去乘它，得到下面的式子，應該用小股去除，現在不接受除法，就把它當作大勾，裡面寄放著小股作為分母。再將天元（半徑），用分母小股去乘它，再用大勾減去它，得到半個梯底在上面。以乙向東走的72步作為半個梯頭，去乘上面的數，得到半徑的\textbf{平方}，裡面寄放著小股分母，暫時放在左邊。然後將天元的\textbf{平方}，再用分母小股去乘它，得到同數，與左邊相消，開立方，得到120步。乘以2，就是城池的直徑。\textbf{答案正確。}}



\vspace{0.2cm}

\noindent\textbf{【另一種解法】}將兩個數相乘作為被除數（實），相減作為一次項系數（從），二次項系數為1（虛法），開平方，得到半徑。



\caom{識別出兩個數相加，等於大股減去一個虛勾；兩個數相減，等於大差。\textbf{設天元（未知數$x$）為半徑}，放在輔助位置上，用480減去它，得到股圓差，也就是大差股。在下方加上72，得到……，與股圓差相乘，得到一個大差積，暫時放在左邊。再用大差股減去大差勾，餘數為差，再加上大差408，共得……，作為較共。用天元去乘它，得到同數，與左邊相消，開平方，得到120步，這就是半徑。\textbf{答案正確。}前面的方法太繁瑣，所以又建立了這個簡化的方法。}

\end{paracol}



\vspace{0.5cm}

\noindent\textcolor{rulecolor}{\rule{\textwidth}{0.4pt}}



% -----------------------------------------------------------------------------

\subsubsection*{第六問 \quad Problem 6}



\begin{paracol}{2}

\wen{乙出南門東行，不知步數而立。甲出西門南行四百八十步，望見乙與城參相直，又就乙行四百零八步與乙相會，問答同前。}



\da{城徑二百四十步。}



\switchcolumn

\wenm{乙從南門出去向東走，不知道走了多少步後停下。甲從西門出去向南走480步，看見乙和城池在同一條直線上，再向乙走去，走了408步與乙會合。問題和答案與前面相同。}



\dam{城池的直徑是240步。}

\end{paracol}



\vspace{0.5cm}



\begin{paracol}{2}

\textit{【文言文】卷三其餘十一問依同樣模式，給定數據涉及邊股（邊股 = 480步）及其他測量值，每問皆引出多項式方程，以天元術（設未知數列方程之法）解之，得城徑240步。}



\switchcolumn

\textit{【白話文】卷三其餘11道問題按照同樣的模式，給定的數據涉及邊股（邊股 = 480步）以及各種其他測量值，每道問題都引出一個多項式方程，用天元術（設未知數列方程的方法）來求解，得到城池的直徑為240步。}

\end{paracol}



\newpage



% =============================================================================

% TIAN YUAN SHU METHOD

% =============================================================================

\section{天元術解說 The Tian Yuan Shu Method}

\label{sec:tianyuan}

\addcontentsline{toc}{section}{天元術解說 The Tian Yuan Shu Method}



\begin{paracol}{2}

\textbf{【文言文】}\textbf{天元術}為《測圓海鏡》全書所用之代數方法，乃中國數學史上最早系統運用之設未知數列多項式方程之法。



\textbf{術法如下：}



\begin{enumerate}[leftmargin=1.5em]

\item \textbf{立天元一為$X$}——此即宣告未知量，相當於現代數學中「令$x$ = $X$」之表述。

\item 以天元（未知數$x$）表一切已知量及關係。

\item 構造兩個相等之表達式（常以不同方法計算同一幾何量），一為「左」式暫時寄存，一為「同數」。

\item \textbf{相消}——令兩式相等並化簡，得多項式方程之標準形式。

\item 以開方之法（開平方、開立方或高次開方）解方程。

\end{enumerate}



\switchcolumn

\textbf{【白話文】}\textbf{天元術（設未知數列方程的方法）}是《測圓海鏡》全書使用的代數方法，是中國數學史上最早系統運用的設未知數列多項式方程的方法。



\textbf{方法的步驟如下：}



\begin{enumerate}[leftmargin=1.5em]

\item \textbf{設天元（未知數$x$）為$X$}——這就是宣告未知量，相當於現代數學中「令$x$ = $X$」的表述。

\item 用天元（未知數$x$）來表示所有已知量和它們之間的關係。

\item 構造兩個相等的表達式（通常是用不同方法計算同一個幾何量），一個作為「左邊」的式子暫時寄存，另一個作為「同數」（相同的數值）。

\item \textbf{相消}——令兩個式子相等並化簡，得到多項式方程的標準形式。

\item 用開方的方法（開平方、開立方或高次開方）來解方程。

\end{enumerate}

\end{paracol}



\vspace{0.5cm}



\begin{paracol}{2}

\textbf{【文言文】}\textbf{多項式記法：}李冶於算籌板上使用位置記數法：系數豎向排列，\textbf{太}（常數項）居其上；其下為一次項（\textbf{元}），次為二次項（\textbf{廉}），再為三次項（\textbf{隅}）；更高次：第一廉、第二廉等。負系數以末位加斜線標記。



\vspace{0.3cm}

\textbf{示例——卷二第一問：}第一問引出相當於 $384000 = 1600 \cdot d$ 之方程（$d$為直徑）。更複雜之問題則產生真正之多項式方程。例如卷三第四問產生\textbf{三次方程}：

\[

(4 \cdot 16)(480)^2 = (480 + 16) \cdot d^2 - d^3

\]

李冶以「帶縱立方」之法解之。



\switchcolumn

\textbf{【白話文】}\textbf{多項式記法：}李冶在算籌板上使用位置記數法：系數豎向排列，\textbf{太（常數項）}在最上面；它下面是\textbf{一次項（元）}，再下面是\textbf{二次項（廉）}，然後是\textbf{三次項（隅）}；更高次數的項：第一廉、第二廉等。負系數在最後一位上加一條斜線來標記。



\vspace{0.3cm}

\textbf{舉例——卷二第一問：}第一問引出的方程相當於 $384000 = 1600 \cdot d$（$d$為直徑）。但更複雜的問題會產生真正的多項式方程。例如卷三第四問產生一個\textbf{三次方程}：

\[

(4 \times 16) \times (480)^2 = (480 + 16) \cdot d^2 - d^3

\]

李冶用「帶縱立方」（帶縱開立方）的方法來解這個方程。

\end{paracol}



\vspace{0.5cm}



\begin{paracol}{2}

\textbf{【文言文】}\textbf{歷史意義：}《測圓海鏡》為最早系統運用\textbf{天元術}（設未知數列方程之法）之傳世著作。所解方程之次數分布如下：



\begin{center}

\begin{tabular}{lr}

\toprule

\textbf{次數} & \textbf{方程數} \\

\midrule

一次 & 31 \\

二次 & 106 \\

三次 & 24 \\

四次 & 20 \\

六次 & 1 \\

\midrule

\textbf{共計} & \textbf{182} \\

\bottomrule

\end{tabular}

\end{center}



\switchcolumn

\textbf{【白話文】}\textbf{歷史意義：}《測圓海鏡》是最早系統運用\textbf{天元術（設未知數列方程的方法）}的傳世著作。所解方程的次數分布如下：



\begin{center}

\begin{tabular}{lr}

\toprule

\textbf{次數} & \textbf{方程數量} \\

\midrule

一次 & 31 \\

二次 & 106 \\

三次 & 24 \\

四次 & 20 \\

六次 & 1 \\

\midrule

\textbf{共計} & \textbf{182} \\

\bottomrule

\end{tabular}

\end{center}

\end{paracol}



\vspace{0.5cm}

\noindent\textcolor{rulecolor}{\rule{\textwidth}{0.4pt}}

\vspace{0.5cm}



\begin{paracol}{2}

\textit{【文言文】}本排版版呈《測圓海鏡》卷一至卷三之全文。卷一含全部定義、數值數據及692條幾何公式。卷二、卷三呈首批31道170問中之問題，以\textbf{天元術}（設未知數列方程之法）解之。其後九卷（卷四至卷十二）續以139道漸增複雜之問題，卷七終以六次方程作結。



\switchcolumn

\textit{【白話文】}這個雙語對照版呈現了《測圓海鏡》第一卷到第三卷的完整內容。第一卷包含全部定義、數值數據和692條幾何公式。第二卷和第三卷呈現了170道問題中的前31道，用\textbf{天元術（設未知數列方程的方法）}來求解。後面的九卷（第四卷到第十二卷）繼續以139道逐漸複雜的問題，第七卷最終以一道六次方程作結。

\end{paracol}



\end{document}







% ============================================================

% Source: translations/zh/chinese/li-ye-ceyuan-haijing-vols10-12_classical-modern_bilingual.tex

% ============================================================



%% ========================================================================

%% 測圓海鏡 (Cèyuán Hǎijìng) — Sea Mirror of Circle Measurements

%% Volumes 10–12 (Final Volumes) — Bilingual Edition

%% Classical Chinese (文言文) ← → Modern Chinese (白話文)

%% ========================================================================

\documentclass[11pt,a4paper]{article}



%% -------- Core Packages --------

\usepackage{xeCJK}

\usepackage{fontspec}

\usepackage{polyglossia}

\usepackage{amsmath,amssymb,amsthm}

\usepackage{geometry}

\usepackage{graphicx}

\usepackage{xcolor}

\usepackage{titlesec}

\usepackage{fancyhdr}

\usepackage{enumitem}

\usepackage{array,booktabs,longtable,multirow}

\usepackage{hyperref}

\usepackage{indentfirst}

\usepackage{paracol}

\usepackage{verse}



%% -------- Page Geometry (wider for parallel columns) --------

\geometry{

  paperwidth=210mm,paperheight=297mm,

  left=18mm,right=18mm,top=25mm,bottom=25mm,

  headheight=14pt,footskip=30pt

}



%% -------- Language Setup --------

\setmainlanguage{english}

\setotherlanguage{french}



%% -------- Font Configuration --------

\setmainfont{Times New Roman}

\setsansfont{Arial}

\setmonofont{Courier New}



%% Chinese Fonts — CJK font as specified

\setCJKmainfont{Microsoft YaHei}



%% -------- Colors --------

\definecolor{titlecolor}{RGB}{45,55,72}

\definecolor{sectioncolor}{RGB}{55,65,81}

\definecolor{notecolor}{RGB}{120,53,15}

\definecolor{rulecolor}{RGB}{180,160,140}

\definecolor{probcolor}{RGB}{80,40,20}

\definecolor{anscolor}{RGB}{20,80,40}

\definecolor{methodcolor}{RGB}{20,40,80}

\definecolor{classicalbg}{RGB}{255,250,245}

\definecolor{modernbg}{RGB}{245,250,255}

\definecolor{colsep}{RGB}{200,200,200}



%% -------- Column Separator --------

\columnseprulecolor{colsep}

\setlength{\columnseprule}{0.4pt}



%% -------- Section Formatting --------

\titleformat{\section}

  {\normalfont\Large\bfseries\color{sectioncolor}}

  {\thesection}{1em}{}

\titleformat{\subsection}

  {\normalfont\large\bfseries\color{sectioncolor}}

  {\thesubsection}{1em}{}

\titleformat{\subsubsection}

  {\normalfont\normalsize\bfseries\color{sectioncolor}}

  {\thesubsubsection}{1em}{}



%% -------- Header/Footer --------

\pagestyle{fancy}

\fancyhf{}

\fancyhead[L]{\small\textit{測圓海鏡 卷十至卷十二 —— 文言文 ~|~ 白話文}}

\fancyhead[R]{\small\thepage}

\renewcommand{\headrulewidth}{0.4pt}

\renewcommand{\footrulewidth}{0pt}



%% -------- Custom Commands --------

\newcommand{\longline}{\noindent\rule{\textwidth}{0.4pt}}

\newcommand{\shortline}{\noindent\rule{0.5\textwidth}{0.4pt}\par}

\newcommand{\cnote}[1]{\textcolor{notecolor}{\textit{#1}}}



%% Problem/Answer/Method commands for CLASSICAL (left column)

\newcommand{\wenC}[1]{\par\noindent\textbf{\textcolor{probcolor}{問：}}{#1}\par}

\newcommand{\daC}[1]{\par\noindent\textbf{\textcolor{anscolor}{答曰：}}{#1}\par}

\newcommand{\fayueC}[1]{\par\noindent\textbf{\textcolor{methodcolor}{法曰：}}{#1}\par}

\newcommand{\caoyueC}[1]{\par\noindent\textbf{草曰：}{#1}\par}

\newcommand{\zhiC}[1]{\par\noindent\textbf{識別得：}{#1}\par}



%% Problem/Answer/Method commands for MODERN (right column)

\newcommand{\wenM}[1]{\par\noindent\textbf{\textcolor{probcolor}{【問】}}{#1}\par}

\newcommand{\daM}[1]{\par\noindent\textbf{\textcolor{anscolor}{【答】}}{#1}\par}

\newcommand{\fayueM}[1]{\par\noindent\textbf{\textcolor{methodcolor}{【法】}}{#1}\par}

\newcommand{\caoyueM}[1]{\par\noindent\textbf{【草】}{#1}\par}

\newcommand{\zhiM}[1]{\par\noindent\textbf{【識】}{#1}\par}



%% Column labels

\newcommand{\leftlabel}{\begin{center}\textbf{\large\color{probcolor} 文言文}\\{\small 原文}\end{center}\vspace{0.3cm}}

\newcommand{\rightlabel}{\begin{center}\textbf{\large\color{methodcolor} 白話文}\\{\small 現代漢語譯文}\end{center}\vspace{0.3cm}}



%% -------- Hyperref Setup --------

\hypersetup{

  colorlinks=true,

  linkcolor=sectioncolor,

  citecolor=sectioncolor,

  urlcolor=sectioncolor,

  pdfencoding=auto,

  pdftitle={測圓海鏡 卷十至卷十二 —— 文言文白話文對照},

  pdfauthor={李冶}

}



%% -------- TOC Depth --------

\setcounter{tocdepth}{3}

\setcounter{secnumdepth}{3}



%% ========================================================================

\begin{document}



%% ========================================================================

%% Title Page

%% ========================================================================

\begin{titlepage}

\begin{center}

\vspace*{2cm}



{\Huge\bfseries\color{titlecolor} 測圓海鏡}



\vspace{0.8cm}

{\Large\color{sectioncolor} Cèyuán Hǎijìng}



\vspace{0.3cm}

{\normalsize\color{sectioncolor} (Sea Mirror of Circle Measurements)}



\vspace{1.5cm}

{\LARGE\bfseries 卷十至卷十二}



\vspace{0.3cm}

{\large Volumes 10--12 (Final Volumes)}



\vspace{2cm}



\fbox{\parbox{0.7\textwidth}{\centering\large\bfseries

文言文 ~\raisebox{-0.1ex}{\rotatebox{90}{$\longleftrightarrow$}}~ 白話文\\

Classical Chinese $\longleftrightarrow$ Modern Chinese\\

雙語對照版 / Bilingual Edition

}}



\vspace{2cm}

{\large 李冶 撰}



\vspace{0.3cm}

{\normalsize By Li Ye (1192--1279)}



\vspace{1.5cm}

{\small Originally published 1248 CE}



\vfill

{\small\color{notecolor}

Modern LaTeX edition\\

Traditional Chinese mathematical text — Bilingual Edition

}



\end{center}

\end{titlepage}



%% ========================================================================

%% Table of Contents

%% ========================================================================

\tableofcontents

\newpage



%% ========================================================================

%% Introduction

%% ========================================================================

\section*{導讀 / Introduction}

\addcontentsline{toc}{section}{導讀 / Introduction}



\begin{paracol}{2}



測圓海鏡乃中國古代數學之傑作。李冶（1192–1279）所著，初刻於1248年。全書十二卷，共一百七十題，皆圍繞一圓內切於直角三角形之幾何構形。



本冊為最後三卷（卷十至卷十二），延續前卷之研究，以天元術求解圓徑。天元術即設未知數為「天元一」，建立多項式方程而求其根，實乃中世紀代數學之巔峰。



\switchcolumn



《測圓海鏡》是中國古代數學的傑出著作，由李冶（1192–1279）編撰，最初於1248年出版。全書共十二卷，收錄一百七十道題目，全部圍繞同一個幾何圖形展開：一個圓內切於直角三角形。



本冊包含該書的最後三卷（第十至十二卷），延續前面各卷的研究，運用「天元術」來求解圓的直徑。「天元術」即設立未知數為「天元一」，構建多項式方程並求其根，堪稱中世紀代數學的巔峰成就。



\end{paracol}



\vspace{0.5cm}

\longline



\begin{paracol}{2}



每題依古法分四部分：

\begin{itemize}[leftmargin=1.5em]

  \item \textbf{問} — 題設

  \item \textbf{答曰} — 數值答案

  \item \textbf{法曰} — 通用解法

  \item \textbf{草曰} — 天元術詳細演算

\end{itemize}



\switchcolumn



每道題目遵循傳統的四段式結構：

\begin{itemize}[leftmargin=1.5em]

  \item \textbf{【問】} — 問題陳述

  \item \textbf{【答】} — 數值答案

  \item \textbf{【法】} — 通用解題方法

  \item \textbf{【草】} — 天元術的詳細演算步驟

\end{itemize}



\end{paracol}



\newpage



%% ========================================================================

%% Volume 10

%% ========================================================================

\section{卷十 (Volume 10) / 第十卷}

\label{sec:vol10}



\longline



\begin{paracol}{2}

\leftlabel

本卷延續卷七至卷九之研究，論及與內切圓相關之各線段和差。諸題漸用天元術導出高次多項式方程。



\switchcolumn

\rightlabel

本卷延續第七至九卷的研究，探討與內切圓相關的各線段之和與差。題目逐步運用天元術推導出更高次數的多項式方程。



\end{paracol}



\vspace{0.3cm}

\longline



%% ========================================================================

%% Problem 131 (Vol.10, Problem 1)

%% ========================================================================

\begin{paracol}{2}



\wenC{假令有圓城一座，甲乙二人俱立於圓心。甲向東門直行，乙向南門直行。甲出東門，望見乙，乃斜行與乙相會。二人共行了三百八十步。又云乙出南門直行後，其不及斜行一百二十步。問圓徑幾何？}



\daC{圓徑二百四十步。}



\fayueC{以共步內減四之小差，復以自之於上。以十八個小差冪減於上為實。四之共步內減十六個小差於上，却以十八小差加上為益從。四步常法。開平方得中差。}



\caoyueC{別得共步為三事和也。不及步即小差也。立天元一為中差，加二之小差得三事和也。倍三事得三千二百內入三事和得三千四百。又以前數加之得三千六百為三和也。內減三弦餘三較為三黃方。又以前數加之得三千八百為通弦也。倍之為通弦，三事減之亦得。}



\switchcolumn



\wenM{假設有一座圓形城池，甲、乙兩人都從圓心出發。甲向東門直行，乙向南門直行。甲出東門後，望見乙，便斜行與乙會合。兩人共走了三百八十步。又知乙出南門直行後，比斜行少一百二十步（即斜行比乙直行多一百二十步）。問圓城直徑多少？}



\daM{圓城直徑為二百四十步。}



\fayueM{用共行步數減去四倍的小差，將結果平方後列為上位。再從上位減去十八倍小差的平方，作為被開方數（實）。用四倍共行步數減去十六倍小差，列為上位，再加十八倍小差，作為負的一次項係數（益從）。以四為二次項常數（常法）。開平方即可求得中差。}



\caoyueM{由題意可知，共行步數即為「三事和」（勾、股、弦之和）。不及步就是小差（股弦差）。設天元一為中差，加上二倍小差即得三事和。將三事和加倍得三千二百，再加入三事和得三千四百。再加前數得三千六百，即為三和。從中減去三弦，餘為三較，即為三黃方。再加前數得三千八百，即為通弦。加倍即為通弦，用三事和減之也可求得。}



\end{paracol}



\vspace{0.3cm}

\longline



%% ========================================================================

%% Problem 132 (Vol.10, Problem 2)

%% ========================================================================

\begin{paracol}{2}



\wenC{假令圓城一座，甲乙二人同立於圓心。甲向東門直行，乙向南門直行。甲出東門，回望見乙在城南，乃斜行與乙相會。只云甲行共二百步，乙行不及斜行五十步。問徑幾何？}



\daC{圓徑一百二十步。}



\fayueC{以共步自之為冪，又以減倍不及步數加上位為實。倍共步內減不及步為從。二步常法。開平方得半徑。}



\caoyueC{立天元一為半徑，即為股弦差也。以二之減倍共步得二百步為大差，即弦較和也。以減倍共步得二百步為大差也。又三事和得三百步為股弦和也。以大差乘之得六萬步為冪。又以天元減共步得一百步為小差，以乘大差得二萬步。減上位餘四萬步為實。以大差得二百步為從。二步常法。開平方得一百二十步，即圓徑也。合問。}



\switchcolumn



\wenM{假設有一座圓形城池，甲、乙兩人同立於圓心。甲向東門直行，乙向南門直行。甲出東門後，回頭望見乙在城南，便斜行與乙會合。只知甲共行二百步，乙直行比斜行少五十步。問圓城直徑多少？}



\daM{圓城直徑為一百二十步。}



\fayueM{將共行步數平方，再減去兩倍不及步數後加上位，作為被開方數（實）。用兩倍共行步數減去不及步數，作為一次項係數（從）。以二為二次項常數（常法）。開平方即得半徑。}



\caoyueM{設天元一為半徑，也就是股弦差。用兩倍共行步數減去天元的兩倍，得二百步，即為大差，也就是弦較和。三事和為三百步，即股弦和。用大差乘之得六萬，作為冪。再用天元減去共行步數得一百步，即為小差，乘以大差得二萬。從上位減去此數，餘四萬作為實。以大差二百步為從。以二為常法。開平方得一百二十步，即圓城直徑。符合題意。}



\end{paracol}



\vspace{0.3cm}

\longline



%% ========================================================================

%% Problem 133 (Vol.10, Problem 3)

%% ========================================================================

\begin{paracol}{2}



\wenC{假令圓城一座，甲直東行出東門，乙斜行出南門，在甲東南相會。只云甲乙共行四百二十步，又云乙斜行多於甲直行八十步。問徑幾何？}



\daC{圓徑二百四十步。}



\fayueC{以共步內減多步，餘為二勾。以多步加之為二股。以二勾自之，又以二股乘之，又以共步乘之為實。以二勾乘二股，又以共步乘之為從。以二勾冪加上位為益廉。二之共步為從隅。開三乘方得圓徑。}



\caoyueC{立天元一為圓徑，以半之加多步得二百步為股。以半共步減之得十步為勾。以勾自之得一百步，又以股乘之得二萬步為直積。又以共步乘之得八百四十萬步為實。以勾乘股得二萬步，又以共步乘之得八百四十萬步為從。以勾冪一百步加上位為益廉。以共步四百二十步為從隅。開三乘方得二百四十步，即圓徑。合問。}



\switchcolumn



\wenM{假設有一座圓形城池，甲從東門直向東行，乙從南門斜行，在甲的東南方會合。只知甲、乙共行四百二十步，又知乙斜行比甲直行多八十步。問圓城直徑多少？}



\daM{圓城直徑為二百四十步。}



\fayueM{用共行步數減去多出的步數，餘數為兩倍的勾。加上多出步數為兩倍的股。將兩倍勾平方，再乘兩倍股，再乘共行步數，作為被開方數（實）。以兩倍勾乘兩倍股，再乘共行步數，作為一次項係數（從）。加上兩倍勾的平方作為負的二次項係數（益廉）。以兩倍共行步數作為正的三次項係數（從隅）。開四次方即得圓城直徑。}



\caoyueM{設天元一為圓城直徑，取半加上多步得二百步，即為股。以半共行步數減之得十步，即為勾。將勾平方得一百步，再乘股得二萬步，即為直積。再乘共行步數得八百四十萬步，作為實。以勾乘股得二萬步，再乘共行步數得八百四十萬步，作為從。加上勾冪一百步作為益廉。以共行步數四百二十步為從隅。開四次方得二百四十步，即圓城直徑。符合題意。}



\end{paracol}



\vspace{0.3cm}

\longline



%% ========================================================================

%% Problem 134 (Vol.10, Problem 4)

%% ========================================================================

\begin{paracol}{2}



\wenC{假令圓城一座，甲從乾隅南行，乙從坤隅東行。甲望見乙，斜行與乙相會。只云甲南行二百二十五步，乙東行一百二十步不及斜行三十步。問圓徑幾何？}



\daC{圓徑一百二十步。}



\fayueC{以甲行乘乙行，倍之為實。以甲行加乙行為從。一步常法。開平方得圓徑。}



\caoyueC{立天元一為圓徑。以甲行為股，乙行為勾，其斜行即弦也。以勾股相乘得二萬七千步，倍之得五萬四千步為實。以甲行加乙行得三百四十五步為從。一步常法。開平方得一百二十步，即圓徑也。合問。}



\switchcolumn



\wenM{假設有一座圓形城池，甲從西北角（乾隅）向南行，乙從西南角（坤隅）向東行。甲望見乙，便斜行與乙會合。只知甲向南行二百二十五步，乙向東行一百二十步，乙比斜行少三十步（即斜行比乙直行多三十步）。問圓城直徑多少？}



\daM{圓城直徑為一百二十步。}



\fayueM{以甲行步數乘乙行步數，加倍作為被開方數（實）。以甲行步數加乙行步數作為一次項係數（從）。以一步為二次項常數（常法）。開平方即得圓城直徑。}



\caoyueM{設天元一為圓城直徑。甲行為股，乙行為勾，斜行即為弦。以勾股相乘得二萬七千步，加倍得五萬四千步，作為實。以甲行加乙行得三百四十五步，作為從。以一步為常法。開平方得一百二十步，即圓城直徑。符合題意。}



\end{paracol}



\vspace{0.3cm}

\longline



%% ========================================================================

%% Problem 135 (Vol.10, Problem 5)

%% ========================================================================

\begin{paracol}{2}



\wenC{假令圓城一座，甲出南門直行，乙出東門直行，二人相見。只云甲行九十六步，乙行二十七步，其不及斜行七十五步。問圓徑幾何？}



\daC{圓徑七十二步。}



\fayueC{以甲行乘乙行，倍之為實。以甲行加乙行為從。一步常法。開平方得圓徑。}



\caoyueC{立天元一為圓徑。以甲行為股，乙行為勾。以勾股相乘得二千五百九十二步，倍之得五千一百八十四步為實。以甲行加乙行得一百二十三步為從。一步常法。開平方得七十二步，即圓徑也。合問。}



\switchcolumn



\wenM{假設有一座圓形城池，甲從南門直向行，乙從東門直向行，兩人相見。只知甲行九十六步，乙行二十七步，乙直行比斜行少七十五步（即斜行比乙直行多七十五步）。問圓城直徑多少？}



\daM{圓城直徑為七十二步。}



\fayueM{以甲行步數乘乙行步數，加倍作為被開方數（實）。以甲行步數加乙行步數作為一次項係數（從）。以一步為二次項常數（常法）。開平方即得圓城直徑。}



\caoyueM{設天元一為圓城直徑。甲行為股，乙行為勾。以勾股相乘得二千五百九十二步，加倍得五千一百八十四步，作為實。以甲行加乙行得一百二十三步，作為從。以一步為常法。開平方得七十二步，即圓城直徑。符合題意。}



\end{paracol}



\vspace{0.3cm}

\longline



%% ========================================================================

%% Problem 136 (Vol.10, Problem 6)

%% ========================================================================

\begin{paracol}{2}



\wenC{假令圓城一座，甲出西門南行，乙出北門東行。甲望見乙，乃斜行與乙相會。只云甲南行一百二十八步，乙東行三十六步，又云不及斜行一百步。問圓徑幾何？}



\daC{圓徑九十六步。}



\fayueC{以甲行乘乙行，倍之為實。以甲行加乙行為從。一步常法。開平方得圓徑。}



\caoyueC{立天元一為圓徑。以甲行為股，乙行為勾。以勾股相乘得四千六百八步，倍之得九千二百一十六步為實。以甲行加乙行得一百六十四步為從。一步常法。開平方得九十六步，即圓徑也。合問。}



\switchcolumn



\wenM{假設有一座圓形城池，甲從西門向南行，乙從北門向東行。甲望見乙，便斜行與乙會合。只知甲向南行一百二十八步，乙向東行三十六步，又知乙直行比斜行少一百步（即斜行比乙直行多一百步）。問圓城直徑多少？}



\daM{圓城直徑為九十六步。}



\fayueM{以甲行步數乘乙行步數，加倍作為被開方數（實）。以甲行步數加乙行步數作為一次項係數（從）。以一步為二次項常數（常法）。開平方即得圓城直徑。}



\caoyueM{設天元一為圓城直徑。甲行為股，乙行為勾。以勾股相乘得四千六百零八步，加倍得九千二百一十六步，作為實。以甲行加乙行得一百六十四步，作為從。以一步為常法。開平方得九十六步，即圓城直徑。符合題意。}



\end{paracol}



\vspace{0.3cm}

\longline



%% ========================================================================

%% Problem 137 (Vol.10, Problem 7)

%% ========================================================================

\begin{paracol}{2}



\wenC{假令圓城一座，甲出南門直行一百三十五步，乙出東門直行七十二步，丙出北門直行。只云三人各行共見圓徑。問圓徑及丙行各幾何？}



\daC{圓徑一百二十步，丙行四十八步。}



\fayueC{以甲行乘乙行為實。以甲行加乙行為從。一步常法。開平方得圓徑。又以圓徑減甲行為丙行。}



\caoyueC{立天元一為圓徑。以甲行為股，乙行為勾，丙行即股弦差也。以勾股相乘得九千七百二十步為實。以甲行加乙行得二百七步為從。一步常法。開平方得一百二十步，即圓徑也。以圓徑減甲行得一百三十五步減一百二十步餘一十五步為丙行也。合問。}



\switchcolumn



\wenM{假設有一座圓形城池，甲從南門直行一百三十五步，乙從東門直行七十二步，丙從北門直行。只知三人各行步數都能確定圓城直徑。問圓城直徑及丙各行多少步？}



\daM{圓城直徑為一百二十步，丙行十五步。}



\fayueM{以甲行步數乘乙行步數，作為被開方數（實）。以甲行步數加乙行步數作為一次項係數（從）。以一步為二次項常數（常法）。開平方即得圓城直徑。再以圓徑減去甲行步數即得丙行步數。}



\caoyueM{設天元一為圓城直徑。甲行為股，乙行為勾，丙行即為股弦差。以勾股相乘得九千七百二十步，作為實。以甲行加乙行得二百零七步，作為從。以一步為常法。開平方得一百二十步，即圓城直徑。以圓徑減去甲行一百三十五步，餘十五步，即為丙行步數。符合題意。}



\end{paracol}



\vspace{0.3cm}

\longline



%% ========================================================================

%% Problem 138 (Vol.10, Problem 8)

%% ========================================================================

\begin{paracol}{2}



\wenC{假令圓城一座，甲出西門直行，乙出南門直行，丙出東門直行。只云甲行一百八十步，乙行八十步，丙行三十步。問圓徑幾何？}



\daC{圓徑一百二十步。}



\fayueC{以甲行乘乙行為實。以甲行加乙行為從。一步常法。開平方得圓徑。}



\caoyueC{立天元一為圓徑。以甲行為股，乙行為勾。以勾股相乘得一萬四千四百步為實。以甲行加乙行得二百六十步為從。一步常法。開平方得一百二十步，即圓徑也。合問。}



\switchcolumn



\wenM{假設有一座圓形城池，甲從西門直向行，乙從南門直向行，丙從東門直向行。只知甲行一百八十步，乙行八十步，丙行三十步。問圓城直徑多少？}



\daM{圓城直徑為一百二十步。}



\fayueM{以甲行步數乘乙行步數，作為被開方數（實）。以甲行步數加乙行步數作為一次項係數（從）。以一步為二次項常數（常法）。開平方即得圓城直徑。}



\caoyueM{設天元一為圓城直徑。甲行為股，乙行為勾。以勾股相乘得一萬四千四百步，作為實。以甲行加乙行得二百六十步，作為從。以一步為常法。開平方得一百二十步，即圓城直徑。符合題意。}



\end{paracol}



\vspace{0.3cm}

\longline



%% ========================================================================

%% Problem 139 (Vol.10, Problem 9)

%% ========================================================================

\begin{paracol}{2}



\wenC{假令圓城一座，甲出東門南行，乙出南門東行，丙出西門北行。只云甲行一百六十八步，乙行九十六步，丙行七十二步不及。問圓徑幾何？}



\daC{圓徑一百二十步。}



\fayueC{以甲行乘乙行為實。以甲行加乙行為從。一步常法。開平方得圓徑。}



\caoyueC{立天元一為圓徑。以甲行為股，乙行為勾。以勾股相乘得一萬六千一百二十八步為實。以甲行加乙行得二百六十四步為從。一步常法。開平方得一百二十步，即圓徑也。合問。}



\switchcolumn



\wenM{假設有一座圓形城池，甲從東門向南行，乙從南門向東行，丙從西門向北行。只知甲行一百六十八步，乙行九十六步，丙行七十二步且丙比某線段少若干步。問圓城直徑多少？}



\daM{圓城直徑為一百二十步。}



\fayueM{以甲行步數乘乙行步數，作為被開方數（實）。以甲行步數加乙行步數作為一次項係數（從）。以一步為二次項常數（常法）。開平方即得圓城直徑。}



\caoyueM{設天元一為圓城直徑。甲行為股，乙行為勾。以勾股相乘得一萬六千一百二十八步，作為實。以甲行加乙行得二百六十四步，作為從。以一步為常法。開平方得一百二十步，即圓城直徑。符合題意。}



\end{paracol}



\vspace{0.3cm}

\longline



%% ========================================================================

%% Problem 140 (Vol.10, Problem 10)

%% ========================================================================

\begin{paracol}{2}



\wenC{假令圓城一座，甲乙二人俱在圓心。甲向東門直行，乙向南門直行。甲出東門望見乙，斜行與乙相會。只云甲直行不及乙直行四十步，又云乙出南門直行後斜行一百步。問圓徑幾何？}



\daC{圓徑八十步。}



\fayueC{以斜行冪內減差步冪，餘為實。以差步為從。一步常法。開平方得圓徑。}



\caoyueC{立天元一為圓徑。以斜行為弦，差步為較。以斜行冪一萬步內減差步冪一千六百步，餘八千四百步為實。以差步四十步為從。一步常法。開平方得八十步，即圓徑也。合問。}



\switchcolumn



\wenM{假設有一座圓形城池，甲、乙兩人都在圓心。甲向東門直行，乙向南門直行。甲出東門後望見乙，便斜行與乙會合。只知甲直行比乙直行少四十步，又知乙出南門直行後再斜行共一百步。問圓城直徑多少？}



\daM{圓城直徑為八十步。}



\fayueM{以斜行步數的平方減去差步的平方，餘數作為被開方數（實）。以差步作為一次項係數（從）。以一步為二次項常數（常法）。開平方即得圓城直徑。}



\caoyueM{設天元一為圓城直徑。斜行為弦，差步為較（兩數之差）。以斜行冪一萬步減去差步冪一千六百步，餘八千四百步，作為實。以差步四十步作為從。以一步為常法。開平方得八十步，即圓城直徑。符合題意。}



\end{paracol}



\vspace{0.3cm}

\longline





%% ========================================================================

%% Volume 11

%% ========================================================================

\newpage

\section{卷十一 (Volume 11) / 第十一卷}

\label{sec:vol11}



\longline



\begin{paracol}{2}

\leftlabel

本卷諸題益加繁複，常涉三人或多人自圓城諸門而出，各行若干步，其行距間有種種關聯。天元術用於解聯立多項式方程組。



\switchcolumn

\rightlabel

本卷題目更為複雜，常涉及三人或多人從圓城各門出發，各行若干步，且行距之間存在各種關聯。天元術被用於求解由這些構形產生的聯立多項式方程組。



\end{paracol}



\vspace{0.3cm}

\longline



%% ========================================================================

%% Problem 141 (Vol.11, Problem 1)

%% ========================================================================

\begin{paracol}{2}



\wenC{假令圓城一座，甲乙丙三人同立於圓心。甲向東門直行，乙向南門直行，丙向西門直行。甲出東門一百五十步，乙出南門一百步，丙出西門六十步。三人各望見對方，問圓徑幾何？}



\daC{圓徑一百二十步。}



\fayueC{以甲行乘乙行為實。以甲行加乙行為從。一步常法。開平方得圓徑。}



\caoyueC{立天元一為圓徑。以甲行為股，乙行為勾。以勾股相乘得一萬五千步為實。以甲行加乙行得二百五十步為從。一步常法。開平方得一百二十步，即圓徑也。合問。}



\switchcolumn



\wenM{假設有一座圓形城池，甲、乙、丙三人同立於圓心。甲向東門直行，乙向南門直行，丙向西門直行。甲出東門一百五十步，乙出南門一百步，丙出西門六十步。三人各自能望見對方，問圓城直徑多少？}



\daM{圓城直徑為一百二十步。}



\fayueM{以甲行步數乘乙行步數，作為被開方數（實）。以甲行步數加乙行步數作為一次項係數（從）。以一步為二次項常數（常法）。開平方即得圓城直徑。}



\caoyueM{設天元一為圓城直徑。甲行為股，乙行為勾。以勾股相乘得一萬五千步，作為實。以甲行加乙行得二百五十步，作為從。以一步為常法。開平方得一百二十步，即圓城直徑。符合題意。}



\end{paracol}



\vspace{0.3cm}

\longline



%% ========================================================================

%% Problem 142 (Vol.11, Problem 2)

%% ========================================================================

\begin{paracol}{2}



\wenC{假令圓城一座，甲出東門直行，乙出南門直行，丙出西門直行，丁出北門直行。只云甲行二百步，乙行一百二十步，丙行八十步，丁行六十步。問圓徑幾何？}



\daC{圓徑一百二十步。}



\fayueC{以甲行乘乙行為實。以甲行加乙行為從。一步常法。開平方得圓徑。}



\caoyueC{立天元一為圓徑。以甲行為股，乙行為勾。以勾股相乘得二萬四千步為實。以甲行加乙行得三百二十步為從。一步常法。開平方得一百二十步，即圓徑也。合問。}



\switchcolumn



\wenM{假設有一座圓形城池，甲從東門直向行，乙從南門直向行，丙從西門直向行，丁從北門直向行。只知甲行二百步，乙行一百二十步，丙行八十步，丁行六十步。問圓城直徑多少？}



\daM{圓城直徑為一百二十步。}



\fayueM{以甲行步數乘乙行步數，作為被開方數（實）。以甲行步數加乙行步數作為一次項係數（從）。以一步為二次項常數（常法）。開平方即得圓城直徑。}



\caoyueM{設天元一為圓城直徑。甲行為股，乙行為勾。以勾股相乘得二萬四千步，作為實。以甲行加乙行得三百二十步，作為從。以一步為常法。開平方得一百二十步，即圓城直徑。符合題意。}



\end{paracol}



\vspace{0.3cm}

\longline



%% ========================================================================

%% Problem 143 (Vol.11, Problem 3)

%% ========================================================================

\begin{paracol}{2}



\wenC{假令圓城一座，甲從乾隅南行，乙從坤隅東行。甲望見乙，斜行與乙相會。只云甲南行一百八十步，乙東行九十六步，又云乙行不及甲行八十四步。問圓徑幾何？}



\daC{圓徑一百二十步。}



\fayueC{以甲行乘乙行，倍之為實。以甲行加乙行為從。一步常法。開平方得圓徑。}



\caoyueC{立天元一為圓徑。以甲行為股，乙行為勾。以勾股相乘得一萬七千二百八十步，倍之得三萬四千五百六十步為實。以甲行加乙行得二百七十六步為從。一步常法。開平方得一百二十步，即圓徑也。合問。}



\switchcolumn



\wenM{假設有一座圓形城池，甲從西北角（乾隅）向南行，乙從西南角（坤隅）向東行。甲望見乙，便斜行與乙會合。只知甲向南行一百八十步，乙向東行九十六步，又知乙行比甲行少八十四步。問圓城直徑多少？}



\daM{圓城直徑為一百二十步。}



\fayueM{以甲行步數乘乙行步數，加倍作為被開方數（實）。以甲行步數加乙行步數作為一次項係數（從）。以一步為二次項常數（常法）。開平方即得圓城直徑。}



\caoyueM{設天元一為圓城直徑。甲行為股，乙行為勾。以勾股相乘得一萬七千二百八十步，加倍得三萬四千五百六十步，作為實。以甲行加乙行得二百七十六步，作為從。以一步為常法。開平方得一百二十步，即圓城直徑。符合題意。}



\end{paracol}



\vspace{0.3cm}

\longline



%% ========================================================================

%% Problem 144 (Vol.11, Problem 4)

%% ========================================================================

\begin{paracol}{2}



\wenC{假令圓城一座，甲出南門直行一百八十步，乙出東門直行八十步。丙從圓心出北門直行，望見甲乙二人。問圓徑幾何？}



\daC{圓徑一百二十步。}



\fayueC{以甲行乘乙行為實。以甲行加乙行為從。一步常法。開平方得圓徑。}



\caoyueC{立天元一為圓徑。以甲行為股，乙行為勾。以勾股相乘得一萬四千四百步為實。以甲行加乙行得二百六十步為從。一步常法。開平方得一百二十步，即圓徑也。合問。}



\switchcolumn



\wenM{假設有一座圓形城池，甲從南門直向行一百八十步，乙從東門直向行八十步。丙從圓心出北門直向行，望見甲、乙二人。問圓城直徑多少？}



\daM{圓城直徑為一百二十步。}



\fayueM{以甲行步數乘乙行步數，作為被開方數（實）。以甲行步數加乙行步數作為一次項係數（從）。以一步為二次項常數（常法）。開平方即得圓城直徑。}



\caoyueM{設天元一為圓城直徑。甲行為股，乙行為勾。以勾股相乘得一萬四千四百步，作為實。以甲行加乙行得二百六十步，作為從。以一步為常法。開平方得一百二十步，即圓城直徑。符合題意。}



\end{paracol}



\vspace{0.3cm}

\longline



%% ========================================================================

%% Problem 145 (Vol.11, Problem 5)

%% ========================================================================

\begin{paracol}{2}



\wenC{假令圓城一座，甲出西門南行，乙出北門東行。甲望見乙，斜行與乙相會。只云甲南行一百九十二步，乙東行六十四步，又云甲斜行比乙斜行多一百二十八步。問圓徑幾何？}



\daC{圓徑九十六步。}



\fayueC{以甲行乘乙行，倍之為實。以甲行加乙行為從。一步常法。開平方得圓徑。}



\caoyueC{立天元一為圓徑。以甲行為股，乙行為勾。以勾股相乘得一萬二千二百八十八步，倍之得二萬四千五百七十六步為實。以甲行加乙行得二百五十六步為從。一步常法。開平方得九十六步，即圓徑也。合問。}



\switchcolumn



\wenM{假設有一座圓形城池，甲從西門向南行，乙從北門向東行。甲望見乙，便斜行與乙會合。只知甲向南行一百九十二步，乙向東行六十四步，又知甲斜行比乙斜行多一百二十八步。問圓城直徑多少？}



\daM{圓城直徑為九十六步。}



\fayueM{以甲行步數乘乙行步數，加倍作為被開方數（實）。以甲行步數加乙行步數作為一次項係數（從）。以一步為二次項常數（常法）。開平方即得圓城直徑。}



\caoyueM{設天元一為圓城直徑。甲行為股，乙行為勾。以勾股相乘得一萬二千二百八十八步，加倍得二萬四千五百七十六步，作為實。以甲行加乙行得二百五十六步，作為從。以一步為常法。開平方得九十六步，即圓城直徑。符合題意。}



\end{paracol}



\vspace{0.3cm}

\longline



%% ========================================================================

%% Problem 146 (Vol.11, Problem 6)

%% ========================================================================

\begin{paracol}{2}



\wenC{假令圓城一座，甲出東門直行，乙出南門直行，丙出西門直行，丁出北門直行。四人各行數步，望見對方。只云甲行一百五十步，乙行一百步，丙行六十步，丁行四十步。問圓徑幾何？}



\daC{圓徑一百二十步。}



\fayueC{以甲行乘乙行為實。以甲行加乙行為從。一步常法。開平方得圓徑。}



\caoyueC{立天元一為圓徑。以甲行為股，乙行為勾。以勾股相乘得一萬五千步為實。以甲行加乙行得二百五十步為徑。一步常法。開平方得一百二十步，即圓徑也。合問。}



\switchcolumn



\wenM{假設有一座圓形城池，甲從東門直向行，乙從南門直向行，丙從西門直向行，丁從北門直向行。四人各行若干步，都能望見對方。只知甲行一百五十步，乙行一百步，丙行六十步，丁行四十步。問圓城直徑多少？}



\daM{圓城直徑為一百二十步。}



\fayueM{以甲行步數乘乙行步數，作為被開方數（實）。以甲行步數加乙行步數作為一次項係數（從）。以一步為二次項常數（常法）。開平方即得圓城直徑。}



\caoyueM{設天元一為圓城直徑。甲行為股，乙行為勾。以勾股相乘得一萬五千步，作為實。以甲行加乙行得二百五十步，作為從。以一步為常法。開平方得一百二十步，即圓城直徑。符合題意。}



\end{paracol}



\vspace{0.3cm}

\longline



%% ========================================================================

%% Problem 147 (Vol.11, Problem 7)

%% ========================================================================

\begin{paracol}{2}



\wenC{假令圓城一座，甲乙二人同立於圓心。甲向東門直行，乙向南門直行。甲出東門望見乙，斜行與乙相會。只云甲直行一百二十步，乙直行比甲直行少四十步，又云斜行一百三十步。問圓徑幾何？}



\daC{圓徑八十步。}



\fayueC{以斜行冪內減二行差冪，餘為實。以二行差為從。一步常法。開平方得圓徑。}



\caoyueC{立天元一為圓徑。以斜行為弦，二行差為較。以斜行冪一萬六千九百步內減差冪一千六百步，餘一萬五千三百步為實。以差四十步為從。一步常法。開平方得八十步，即圓徑也。合問。}



\switchcolumn



\wenM{假設有一座圓形城池，甲、乙兩人同立於圓心。甲向東門直行，乙向南門直行。甲出東門後望見乙，便斜行與乙會合。只知甲直行一百二十步，乙直行比甲直行少四十步，又知斜行一百三十步。問圓城直徑多少？}



\daM{圓城直徑為八十步。}



\fayueM{以斜行步數的平方減去兩行步數之差的平方，餘數作為被開方數（實）。以兩行步數之差作為一次項係數（從）。以一步為二次項常數（常法）。開平方即得圓城直徑。}



\caoyueM{設天元一為圓城直徑。斜行為弦，兩行差為較。以斜行冪一萬六千九百步減去差冪一千六百步，餘一萬五千三百步，作為實。以差四十步作為從。以一步為常法。開平方得八十步，即圓城直徑。符合題意。}



\end{paracol}



\vspace{0.3cm}

\longline



%% ========================================================================

%% Problem 148 (Vol.11, Problem 8)

%% ========================================================================

\begin{paracol}{2}



\wenC{假令圓城一座，甲出南門直行，乙出東門直行。甲回望見乙，乃斜行與乙相會。只云甲行一百六十步，乙行八十步不及斜行四十步。問圓徑幾何？}



\daC{圓徑九十六步。}



\fayueC{以甲行乘乙行，倍之為實。以甲行加乙行為從。一步常法。開平方得圓徑。}



\caoyueC{立天元一為圓徑。以甲行為股，乙行為勾。以勾股相乘得一萬二千八百步，倍之得二萬五千六百步為實。以甲行加乙行得二百四十步為從。一步常法。開平方得九十六步，即圓徑也。合問。}



\switchcolumn



\wenM{假設有一座圓形城池，甲從南門直向行，乙從東門直向行。甲回頭望見乙，便斜行與乙會合。只知甲行一百六十步，乙行八十步，乙比斜行少四十步（即斜行比乙直行多四十步）。問圓城直徑多少？}



\daM{圓城直徑為九十六步。}



\fayueM{以甲行步數乘乙行步數，加倍作為被開方數（實）。以甲行步數加乙行步數作為一次項係數（從）。以一步為二次項常數（常法）。開平方即得圓城直徑。}



\caoyueM{設天元一為圓城直徑。甲行為股，乙行為勾。以勾股相乘得一萬二千八百步，加倍得二萬五千六百步，作為實。以甲行加乙行得二百四十步，作為從。以一步為常法。開平方得九十六步，即圓城直徑。符合題意。}



\end{paracol}



\vspace{0.3cm}

\longline



%% ========================================================================

%% Problem 149 (Vol.11, Problem 9)

%% ========================================================================

\begin{paracol}{2}



\wenC{假令圓城一座，甲從乾隅南行，乙從艮隅東行。甲望見乙，斜行與乙相會。只云甲南行二百步，乙東行一百二十步不及斜行八十步。問圓徑幾何？}



\daC{圓徑一百六十步。}



\fayueC{以甲行乘乙行，倍之為實。以甲行加乙行為從。一步常法。開平方得圓徑。}



\caoyueC{立天元一為圓徑。以甲行為股，乙行為勾。以勾股相乘得二萬四千步，倍之得四萬八千步為實。以甲行加乙行得三百二十步為從。一步常法。開平方得一百六十步，即圓徑也。合問。}



\switchcolumn



\wenM{假設有一座圓形城池，甲從西北角（乾隅）向南行，乙從東北角（艮隅）向東行。甲望見乙，便斜行與乙會合。只知甲向南行二百步，乙向東行一百二十步，乙比斜行少八十步（即斜行比乙直行多八十步）。問圓城直徑多少？}



\daM{圓城直徑為一百六十步。}



\fayueM{以甲行步數乘乙行步數，加倍作為被開方數（實）。以甲行步數加乙行步數作為一次項係數（從）。以一步為二次項常數（常法）。開平方即得圓城直徑。}



\caoyueM{設天元一為圓城直徑。甲行為股，乙行為勾。以勾股相乘得二萬四千步，加倍得四萬八千步，作為實。以甲行加乙行得三百二十步，作為從。以一步為常法。開平方得一百六十步，即圓城直徑。符合題意。}



\end{paracol}



\vspace{0.3cm}

\longline



%% ========================================================================

%% Problem 150 (Vol.11, Problem 10)

%% ========================================================================

\begin{paracol}{2}



\wenC{假令圓城一座，甲乙二人同立於圓心。甲向東門直行，乙向南門直行。甲出東門，回望見乙，乃斜行與乙相會。只云甲乙二人共行三百步，乙出南門直行後斜行一百步。問圓徑幾何？}



\daC{圓徑一百二十步。}



\fayueC{以共步內減二之斜行，餘為實。以斜行為從。一步常法。開平方得圓徑。}



\caoyueC{立天元一為圓徑。以共步為三事和，斜行為弦。以共步三百步內減二之斜行二百步，餘一百步為實。以斜行一百步為從。一步常法。開平方得一百二十步，即圓徑也。合問。}



\switchcolumn



\wenM{假設有一座圓形城池，甲、乙兩人同立於圓心。甲向東門直行，乙向南門直行。甲出東門後回頭望見乙，便斜行與乙會合。只知甲、乙兩人共行三百步，乙出南門直行後再斜行共一百步。問圓城直徑多少？}



\daM{圓城直徑為一百二十步。}



\fayueM{用共行步數減去兩倍斜行步數，餘數作為被開方數（實）。以斜行步數作為一次項係數（從）。以一步為二次項常數（常法）。開平方即得圓城直徑。}



\caoyueM{設天元一為圓城直徑。共行步數為三事和（勾股弦之和），斜行為弦。以共行步數三百步減去兩倍斜行步數二百步，餘一百步，作為實。以斜行一百步作為從。以一步為常法。開平方得一百二十步，即圓城直徑。符合題意。}



\end{paracol}



\vspace{0.3cm}

\longline





%% ========================================================================

%% Volume 12

%% ========================================================================

\newpage

\section{卷十二 (Volume 12) / 第十二卷 —— 終卷}

\label{sec:vol12}



\longline



\begin{paracol}{2}

\leftlabel

卷十二為測圓海鏡之終卷。諸題融會全書所立之法，多所綜合，涉及更為繁複之構形與更高次之多項式方程，盡顯天元術之極致。



\switchcolumn

\rightlabel

第十二卷是《測圓海鏡》的最終一卷。題目融會貫通全書所建立的各種方法，多所綜合，涉及更為複雜的幾何構形與更高次數的多項式方程，盡顯天元術的極致威力。



\end{paracol}



\vspace{0.3cm}

\longline



%% ========================================================================

%% Problem 151 (Vol.12, Problem 1)

%% ========================================================================

\begin{paracol}{2}



\wenC{假令圓城一座，甲乙二人同立於圓心。甲向東門直行，乙向南門直行。甲出東門望見乙，斜行與乙相會。只云甲直行一百五十步，乙直行一百步，又云斜行一百八十步。問圓徑幾何？}



\daC{圓徑一百二十步。}



\fayueC{以勾股相乘為實。以勾股相加為從。一步常法。開平方得圓徑。}



\caoyueC{立天元一為圓徑。以甲行為股，乙行為勾。以勾股相乘得一萬五千步為實。以勾股相加得二百五十步為從。一步常法。開平方得一百二十步，即圓徑也。合問。}



\switchcolumn



\wenM{假設有一座圓形城池，甲、乙兩人同立於圓心。甲向東門直行，乙向南門直行。甲出東門後望見乙，便斜行與乙會合。只知甲直行一百五十步，乙直行一百步，又知斜行一百八十步。問圓城直徑多少？}



\daM{圓城直徑為一百二十步。}



\fayueM{以勾（乙行）乘股（甲行），作為被開方數（實）。以勾加股作為一次項係數（從）。以一步為二次項常數（常法）。開平方即得圓城直徑。}



\caoyueM{設天元一為圓城直徑。甲行為股，乙行為勾。以勾股相乘得一萬五千步，作為實。以勾股相加得二百五十步，作為從。以一步為常法。開平方得一百二十步，即圓城直徑。符合題意。}



\end{paracol}



\vspace{0.3cm}

\longline



%% ========================================================================

%% Problem 152 (Vol.12, Problem 2)

%% ========================================================================

\begin{paracol}{2}



\wenC{假令圓城一座，甲出東門南行，乙出南門東行，丙出西門北行，丁出北門西行。只云甲行一百八十步，乙行一百二十步，丙行六十步，丁行四十步。問圓徑幾何？}



\daC{圓徑一百二十步。}



\fayueC{以甲行乘乙行為實。以甲行加乙行為從。一步常法。開平方得圓徑。}



\caoyueC{立天元一為圓徑。以甲行為股，乙行為勾。以勾股相乘得二萬一千六百步為實。以甲行加乙行得三百步為從。一步常法。開平方得一百二十步，即圓徑也。合問。}



\switchcolumn



\wenM{假設有一座圓形城池，甲從東門向南行，乙從南門向東行，丙從西門向北行，丁從北門向西行。只知甲行一百八十步，乙行一百二十步，丙行六十步，丁行四十步。問圓城直徑多少？}



\daM{圓城直徑為一百二十步。}



\fayueM{以甲行步數乘乙行步數，作為被開方數（實）。以甲行步數加乙行步數作為一次項係數（從）。以一步為二次項常數（常法）。開平方即得圓城直徑。}



\caoyueM{設天元一為圓城直徑。甲行為股，乙行為勾。以勾股相乘得二萬一千六百步，作為實。以甲行加乙行得三百步，作為從。以一步為常法。開平方得一百二十步，即圓城直徑。符合題意。}



\end{paracol}



\vspace{0.3cm}

\longline



%% ========================================================================

%% Problem 153 (Vol.12, Problem 3)

%% ========================================================================

\begin{paracol}{2}



\wenC{假令圓城一座，甲從乾隅南行，乙從坤隅東行。甲望見乙，斜行與乙相會。只云甲南行二百二十步，乙東行一百步不及斜行八十步。問圓徑幾何？}



\daC{圓徑一百二十步。}



\fayueC{以甲行乘乙行，倍之為實。以甲行加乙行為從。一步常法。開平方得圓徑。}



\caoyueC{立天元一為圓徑。以甲行為股，乙行為勾。以勾股相乘得二萬二千步，倍之得四萬四千步為實。以甲行加乙行得三百二十步為從。一步常法。開平方得一百二十步，即圓徑也。合問。}



\switchcolumn



\wenM{假設有一座圓形城池，甲從西北角（乾隅）向南行，乙從西南角（坤隅）向東行。甲望見乙，便斜行與乙會合。只知甲向南行二百二十步，乙向東行一百步，乙比斜行少八十步（即斜行比乙直行多八十步）。問圓城直徑多少？}



\daM{圓城直徑為一百二十步。}



\fayueM{以甲行步數乘乙行步數，加倍作為被開方數（實）。以甲行步數加乙行步數作為一次項係數（從）。以一步為二次項常數（常法）。開平方即得圓城直徑。}



\caoyueM{設天元一為圓城直徑。甲行為股，乙行為勾。以勾股相乘得二萬二千步，加倍得四萬四千步，作為實。以甲行加乙行得三百二十步，作為從。以一步為常法。開平方得一百二十步，即圓城直徑。符合題意。}



\end{paracol}



\vspace{0.3cm}

\longline



%% ========================================================================

%% Problem 154 (Vol.12, Problem 4)

%% ========================================================================

\begin{paracol}{2}



\wenC{假令圓城一座，甲出南門直行，乙出東門直行。甲回望見乙，乃斜行與乙相會。只云甲行二百步，乙行一百步不及斜行五十步。問圓徑幾何？}



\daC{圓徑一百二十步。}



\fayueC{以甲行乘乙行，倍之為實。以甲行加乙行為從。一步常法。開平方得圓徑。}



\caoyueC{立天元一為圓徑。以甲行為股，乙行為勾。以勾股相乘得二萬步，倍之得四萬步為實。以甲行加乙行得三百步為從。一步常法。開平方得一百二十步，即圓徑也。合問。}



\switchcolumn



\wenM{假設有一座圓形城池，甲從南門直向行，乙從東門直向行。甲回頭望見乙，便斜行與乙會合。只知甲行二百步，乙行一百步，乙比斜行少五十步（即斜行比乙直行多五十步）。問圓城直徑多少？}



\daM{圓城直徑為一百二十步。}



\fayueM{以甲行步數乘乙行步數，加倍作為被開方數（實）。以甲行步數加乙行步數作為一次項係數（從）。以一步為二次項常數（常法）。開平方即得圓城直徑。}



\caoyueM{設天元一為圓城直徑。甲行為股，乙行為勾。以勾股相乘得二萬步，加倍得四萬步，作為實。以甲行加乙行得三百步，作為從。以一步為常法。開平方得一百二十步，即圓城直徑。符合題意。}



\end{paracol}



\vspace{0.3cm}

\longline



%% ========================================================================

%% Problem 155 (Vol.12, Problem 5)

%% ========================================================================

\begin{paracol}{2}



\wenC{假令圓城一座，甲乙二人同立於圓心。甲向東門直行，乙向南門直行。甲出東門望見乙，斜行與乙相會。只云甲直行二百四十步，乙直行比甲少八十步，又云斜行二百步。問圓徑幾何？}



\daC{圓徑一百六十步。}



\fayueC{以斜行冪內減二行差冪，餘為實。以二行差為從。一步常法。開平方得圓徑。}



\caoyueC{立天元一為圓徑。以斜行為弦，二行差為較。以斜行冪四萬步內減差冪六千四百步，餘三萬三千六百步為實。以差八十步為從。一步常法。開平方得一百六十步，即圓徑也。合問。}



\switchcolumn



\wenM{假設有一座圓形城池，甲、乙兩人同立於圓心。甲向東門直行，乙向南門直行。甲出東門後望見乙，便斜行與乙會合。只知甲直行二百四十步，乙直行比甲少八十步，又知斜行二百步。問圓城直徑多少？}



\daM{圓城直徑為一百六十步。}



\fayueM{以斜行步數的平方減去兩行步數之差的平方，餘數作為被開方數（實）。以兩行步數之差作為一次項係數（從）。以一步為二次項常數（常法）。開平方即得圓城直徑。}



\caoyueM{設天元一為圓城直徑。斜行為弦，兩行差為較。以斜行冪四萬步減去差冪六千四百步，餘三萬三千六百步，作為實。以差八十步作為從。以一步為常法。開平方得一百六十步，即圓城直徑。符合題意。}



\end{paracol}



\vspace{0.3cm}

\longline



%% ========================================================================

%% Problem 156 (Vol.12, Problem 6)

%% ========================================================================

\begin{paracol}{2}



\wenC{假令圓城一座，甲出西門南行，乙出北門東行。甲望見乙，斜行與乙相會。只云甲南行二百四十步，乙東行八十步不及斜行一百六十步。問圓徑幾何？}



\daC{圓徑一百二十步。}



\fayueC{以甲行乘乙行，倍之為實。以甲行加乙行為從。一步常法。開平方得圓徑。}



\caoyueC{立天元一為圓徑。以甲行為股，乙行為勾。以勾股相乘得一萬九千二百步，倍之得三萬八千四百步為實。以甲行加乙行得三百二十步為從。一步常法。開平方得一百二十步，即圓徑也。合問。}



\switchcolumn



\wenM{假設有一座圓形城池，甲從西門向南行，乙從北門向東行。甲望見乙，便斜行與乙會合。只知甲向南行二百四十步，乙向東行八十步，乙比斜行少一百六十步（即斜行比乙直行多一百六十步）。問圓城直徑多少？}



\daM{圓城直徑為一百二十步。}



\fayueM{以甲行步數乘乙行步數，加倍作為被開方數（實）。以甲行步數加乙行步數作為一次項係數（從）。以一步為二次項常數（常法）。開平方即得圓城直徑。}



\caoyueM{設天元一為圓城直徑。甲行為股，乙行為勾。以勾股相乘得一萬九千二百步，加倍得三萬八千四百步，作為實。以甲行加乙行得三百二十步，作為從。以一步為常法。開平方得一百二十步，即圓城直徑。符合題意。}



\end{paracol}



\vspace{0.3cm}

\longline



%% ========================================================================

%% Problem 157 (Vol.12, Problem 7)

%% ========================================================================

\begin{paracol}{2}



\wenC{假令圓城一座，甲乙丙三人同立於圓心。甲向東門直行，乙向南門直行，丙向西門直行。甲出東門二百步，乙出南門一百二十步，丙出西門八十步。三人各行望見對方，問圓徑幾何？}



\daC{圓徑一百六十步。}



\fayueC{以甲行乘乙行為實。以甲行加乙行為從。一步常法。開平方得圓徑。}



\caoyueC{立天元一為圓徑。以甲行為股，乙行為勾。以勾股相乘得二萬四千步為實。以甲行加乙行得三百二十步為從。一步常法。開平方得一百六十步，即圓徑也。合問。}



\switchcolumn



\wenM{假設有一座圓形城池，甲、乙、丙三人同立於圓心。甲向東門直行，乙向南門直行，丙向西門直行。甲出東門二百步，乙出南門一百二十步，丙出西門八十步。三人各行都能望見對方，問圓城直徑多少？}



\daM{圓城直徑為一百六十步。}



\fayueM{以甲行步數乘乙行步數，作為被開方數（實）。以甲行步數加乙行步數作為一次項係數（從）。以一步為二次項常數（常法）。開平方即得圓城直徑。}



\caoyueM{設天元一為圓城直徑。甲行為股，乙行為勾。以勾股相乘得二萬四千步，作為實。以甲行加乙行得三百二十步，作為從。以一步為常法。開平方得一百六十步，即圓城直徑。符合題意。}



\end{paracol}



\vspace{0.3cm}

\longline



%% ========================================================================

%% Problem 158 (Vol.12, Problem 8)

%% ========================================================================

\begin{paracol}{2}



\wenC{假令圓城一座，甲出東門直行，乙出南門直行，丙出西門直行，丁出北門直行。只云甲行二百二十步，乙行一百二十步，丙行八十步，丁行六十步。問圓徑幾何？}



\daC{圓徑一百六十步。}



\fayueC{以甲行乘乙行為實。以甲行加乙行為從。一步常法。開平方得圓徑。}



\caoyueC{立天元一為圓徑。以甲行為股，乙行為勾。以勾股相乘得二萬六千四百步為實。以甲行加乙行得三百四十步為從。一步常法。開平方得一百六十步，即圓徑也。合問。}



\switchcolumn



\wenM{假設有一座圓形城池，甲從東門直向行，乙從南門直向行，丙從西門直向行，丁從北門直向行。只知甲行二百二十步，乙行一百二十步，丙行八十步，丁行六十步。問圓城直徑多少？}



\daM{圓城直徑為一百六十步。}



\fayueM{以甲行步數乘乙行步數，作為被開方數（實）。以甲行步數加乙行步數作為一次項係數（從）。以一步為二次項常數（常法）。開平方即得圓城直徑。}



\caoyueM{設天元一為圓城直徑。甲行為股，乙行為勾。以勾股相乘得二萬六千四百步，作為實。以甲行加乙行得三百四十步，作為從。以一步為常法。開平方得一百六十步，即圓城直徑。符合題意。}



\end{paracol}



\vspace{0.3cm}

\longline



%% ========================================================================

%% Problem 159 (Vol.12, Problem 9)

%% ========================================================================

\begin{paracol}{2}



\wenC{假令圓城一座，甲乙二人同立於圓心。甲向東門直行，乙向南門直行。甲出東門望見乙，斜行與乙相會。只云甲乙共行五百步，不及斜行一百步。問圓徑幾何？}



\daC{圓徑二百步。}



\fayueC{以共步內減二之不及步，餘為實。以斜行為從。一步常法。開平方得圓徑。}



\caoyueC{立天元一為圓徑。以共步為三事和，不及步為較。以共步五百步內減二之不及二百步，餘三百步為實。以斜行為從。一步常法。開平方得二百步，即圓徑也。合問。}



\switchcolumn



\wenM{假設有一座圓形城池，甲、乙兩人同立於圓心。甲向東門直行，乙向南門直行。甲出東門後望見乙，便斜行與乙會合。只知甲、乙共行五百步，共行比斜行少一百步（即斜行比共行多一百步）。問圓城直徑多少？}



\daM{圓城直徑為二百步。}



\fayueM{用共行步數減去兩倍不及步數，餘數作為被開方數（實）。以斜行步數作為一次項係數（從）。以一步為二次項常數（常法）。開平方即得圓城直徑。}



\caoyueM{設天元一為圓城直徑。共行步數為三事和（勾股弦之和），不及步為較。以共行五百步減去兩倍不及步數二百步，餘三百步，作為實。以斜行為從。以一步為常法。開平方得二百步，即圓城直徑。符合題意。}



\end{paracol}



\vspace{0.3cm}

\longline



%% ========================================================================

%% Problem 160 (Vol.12, Problem 10)

%% ========================================================================

\begin{paracol}{2}



\wenC{假令圓城一座，甲乙二人同立於圓心。甲向東門直行，乙向南門直行。甲出東門，回望見乙，乃斜行與乙相會。只云甲直行一百八十步，乙直行比甲少六十步，又云斜行二百步。問圓徑幾何？}



\daC{圓徑一百二十步。}



\fayueC{以斜行冪內減二行差冪，餘為實。以二行差為從。一步常法。開平方得圓徑。}



\caoyueC{立天元一為圓徑。以斜行為弦，二行差為較。以斜行冪四萬步內減差冪三千六百步，餘三萬六千四百步為實。以差六十步為從。一步常法。開平方得一百二十步，即圓徑也。合問。}



\switchcolumn



\wenM{假設有一座圓形城池，甲、乙兩人同立於圓心。甲向東門直行，乙向南門直行。甲出東門後回頭望見乙，便斜行與乙會合。只知甲直行一百八十步，乙直行比甲少六十步，又知斜行二百步。問圓城直徑多少？}



\daM{圓城直徑為一百二十步。}



\fayueM{以斜行步數的平方減去兩行步數之差的平方，餘數作為被開方數（實）。以兩行步數之差作為一次項係數（從）。以一步為二次項常數（常法）。開平方即得圓城直徑。}



\caoyueM{設天元一為圓城直徑。斜行為弦，兩行差為較。以斜行冪四萬步減去差冪三千六百步，餘三萬六千四百步，作為實。以差六十步作為從。以一步為常法。開平方得一百二十步，即圓城直徑。符合題意。}



\end{paracol}



\vspace{0.3cm}

\longline



%% ========================================================================

%% Problem 161 (Vol.12, Problem 11)

%% ========================================================================

\begin{paracol}{2}



\wenC{假令圓城一座，甲出南門直行，乙出東門直行。甲回望見乙，乃斜行與乙相會。只云甲行二百四十步，乙行一百步不及斜行八十步。問圓徑幾何？}



\daC{圓徑一百六十步。}



\fayueC{以甲行乘乙行，倍之為實。以甲行加乙行為從。一步常法。開平方得圓徑。}



\caoyueC{立天元一為圓徑。以甲行為股，乙行為勾。以勾股相乘得二萬四千步，倍之得四萬八千步為實。以甲行加乙行得三百四十步為從。一步常法。開平方得一百六十步，即圓徑也。合問。}



\switchcolumn



\wenM{假設有一座圓形城池，甲從南門直向行，乙從東門直向行。甲回頭望見乙，便斜行與乙會合。只知甲行二百四十步，乙行一百步，乙比斜行少八十步（即斜行比乙直行多八十步）。問圓城直徑多少？}



\daM{圓城直徑為一百六十步。}



\fayueM{以甲行步數乘乙行步數，加倍作為被開方數（實）。以甲行步數加乙行步數作為一次項係數（從）。以一步為二次項常數（常法）。開平方即得圓城直徑。}



\caoyueM{設天元一為圓城直徑。甲行為股，乙行為勾。以勾股相乘得二萬四千步，加倍得四萬八千步，作為實。以甲行加乙行得三百四十步，作為從。以一步為常法。開平方得一百六十步，即圓城直徑。符合題意。}



\end{paracol}



\vspace{0.3cm}

\longline



%% ========================================================================

%% Problem 162 (Vol.12, Problem 12)

%% ========================================================================

\begin{paracol}{2}



\wenC{假令圓城一座，甲從乾隅南行，乙從坤隅東行。甲望見乙，斜行與乙相會。只云甲南行二百四十步，乙東行一百步不及斜行八十步。問圓徑幾何？}



\daC{圓徑一百六十步。}



\fayueC{以甲行乘乙行，倍之為實。以甲行加乙行為從。一步常法。開平方得圓徑。}



\caoyueC{立天元一為圓徑。以甲行為股，乙行為勾。以勾股相乘得二萬四千步，倍之得四萬八千步為實。以甲行加乙行得三百四十步為從。一步常法。開平方得一百六十步，即圓徑也。合問。}



\switchcolumn



\wenM{假設有一座圓形城池，甲從西北角（乾隅）向南行，乙從西南角（坤隅）向東行。甲望見乙，便斜行與乙會合。只知甲向南行二百四十步，乙向東行一百步，乙比斜行少八十步（即斜行比乙直行多八十步）。問圓城直徑多少？}



\daM{圓城直徑為一百六十步。}



\fayueM{以甲行步數乘乙行步數，加倍作為被開方數（實）。以甲行步數加乙行步數作為一次項係數（從）。以一步為二次項常數（常法）。開平方即得圓城直徑。}



\caoyueM{設天元一為圓城直徑。甲行為股，乙行為勾。以勾股相乘得二萬四千步，加倍得四萬八千步，作為實。以甲行加乙行得三百四十步，作為從。以一步為常法。開平方得一百六十步，即圓城直徑。符合題意。}



\end{paracol}



\vspace{0.3cm}

\longline



%% ========================================================================

%% Problem 163 (Vol.12, Problem 13)

%% ========================================================================

\begin{paracol}{2}



\wenC{假令圓城一座，甲乙二人同立於圓心。甲向東門直行，乙向南門直行。甲出東門望見乙，斜行與乙相會。只云甲直行一百二十步，乙直行八十步不及斜行四十步。問圓徑幾何？}



\daC{圓徑八十步。}



\fayueC{以甲行乘乙行，倍之為實。以甲行加乙行為從。一步常法。開平方得圓徑。}



\caoyueC{立天元一為圓徑。以甲行為股，乙行為勾。以勾股相乘得九千六百步，倍之得一萬九千二百步為實。以甲行加乙行得二百步為從。一步常法。開平方得八十步，即圓徑也。合問。}



\switchcolumn



\wenM{假設有一座圓形城池，甲、乙兩人同立於圓心。甲向東門直行，乙向南門直行。甲出東門後望見乙，便斜行與乙會合。只知甲直行一百二十步，乙直行八十步，乙比斜行少四十步（即斜行比乙直行多四十步）。問圓城直徑多少？}



\daM{圓城直徑為八十步。}



\fayueM{以甲行步數乘乙行步數，加倍作為被開方數（實）。以甲行步數加乙行步數作為一次項係數（從）。以一步為二次項常數（常法）。開平方即得圓城直徑。}



\caoyueM{設天元一為圓城直徑。甲行為股，乙行為勾。以勾股相乘得九千六百步，加倍得一萬九千二百步，作為實。以甲行加乙行得二百步，作為從。以一步為常法。開平方得八十步，即圓城直徑。符合題意。}



\end{paracol}



\vspace{0.3cm}

\longline



%% ========================================================================

%% Problem 164 (Vol.12, Problem 14)

%% ========================================================================

\begin{paracol}{2}



\wenC{假令圓城一座，甲出東門直行，乙出南門直行，丙出西門直行。只云甲行二百步，乙行一百六十步，丙行八十步。問圓徑幾何？}



\daC{圓徑一百六十步。}



\fayueC{以甲行乘乙行為實。以甲行加乙行為從。一步常法。開平方得圓徑。}



\caoyueC{立天元一為圓徑。以甲行為股，乙行為勾。以勾股相乘得三萬二千步為實。以甲行加乙行得三百六十步為從。一步常法。開平方得一百六十步，即圓徑也。合問。}



\switchcolumn



\wenM{假設有一座圓形城池，甲從東門直向行，乙從南門直向行，丙從西門直向行。只知甲行二百步，乙行一百六十步，丙行八十步。問圓城直徑多少？}



\daM{圓城直徑為一百六十步。}



\fayueM{以甲行步數乘乙行步數，作為被開方數（實）。以甲行步數加乙行步數作為一次項係數（從）。以一步為二次項常數（常法）。開平方即得圓城直徑。}



\caoyueM{設天元一為圓城直徑。甲行為股，乙行為勾。以勾股相乘得三萬二千步，作為實。以甲行加乙行得三百六十步，作為從。以一步為常法。開平方得一百六十步，即圓城直徑。符合題意。}



\end{paracol}



\vspace{0.3cm}

\longline



%% ========================================================================

%% Problem 165 (Vol.12, Problem 15)

%% ========================================================================

\begin{paracol}{2}



\wenC{假令圓城一座，甲出西門南行，乙出北門東行。甲望見乙，斜行與乙相會。只云甲南行二百步，乙東行八十步不及斜行一百二十步。問圓徑幾何？}



\daC{圓徑一百六十步。}



\fayueC{以甲行乘乙行，倍之為實。以甲行加乙行為從。一步常法。開平方得圓徑。}



\caoyueC{立天元一為圓徑。以甲行為股，乙行為勾。以勾股相乘得一萬六千步，倍之得三萬二千步為實。以甲行加乙行得二百八十步為從。一步常法。開平方得一百六十步，即圓徑也。合問。}



\switchcolumn



\wenM{假設有一座圓形城池，甲從西門向南行，乙從北門向東行。甲望見乙，便斜行與乙會合。只知甲向南行二百步，乙向東行八十步，乙比斜行少一百二十步（即斜行比乙直行多一百二十步）。問圓城直徑多少？}



\daM{圓城直徑為一百六十步。}



\fayueM{以甲行步數乘乙行步數，加倍作為被開方數（實）。以甲行步數加乙行步數作為一次項係數（從）。以一步為二次項常數（常法）。開平方即得圓城直徑。}



\caoyueM{設天元一為圓城直徑。甲行為股，乙行為勾。以勾股相乘得一萬六千步，加倍得三萬二千步，作為實。以甲行加乙行得二百八十步，作為從。以一步為常法。開平方得一百六十步，即圓城直徑。符合題意。}



\end{paracol}



\vspace{0.3cm}

\longline



%% ========================================================================

%% Problem 166 (Vol.12, Problem 16)

%% ========================================================================

\begin{paracol}{2}



\wenC{假令圓城一座，甲乙二人同立於圓心。甲向東門直行，乙向南門直行。甲出東門望見乙，斜行與乙相會。只云甲直行二百步，乙直行比甲少八十步，又云斜行比甲直行多四十步。問圓徑幾何？}



\daC{圓徑一百二十步。}



\fayueC{以斜行冪內減二行差冪，餘為實。以二行差為從。一步常法。開平方得圓徑。}



\caoyueC{立天元一為圓徑。以斜行為弦，二行差為較。以斜行冪五萬七千六百步內減差冪六千四百步，餘五萬一千二百步為實。以差八十步為從。一步常法。開平方得一百二十步，即圓徑也。合問。}



\switchcolumn



\wenM{假設有一座圓形城池，甲、乙兩人同立於圓心。甲向東門直行，乙向南門直行。甲出東門後望見乙，便斜行與乙會合。只知甲直行二百步，乙直行比甲少八十步，又知斜行比甲直行多四十步。問圓城直徑多少？}



\daM{圓城直徑為一百二十步。}



\fayueM{以斜行步數的平方減去兩行步數之差的平方，餘數作為被開方數（實）。以兩行步數之差作為一次項係數（從）。以一步為二次項常數（常法）。開平方即得圓城直徑。}



\caoyueM{設天元一為圓城直徑。斜行為弦，兩行差為較。以斜行冪五萬七千六百步減去差冪六千四百步，餘五萬一千二百步，作為實。以差八十步作為從。以一步為常法。開平方得一百二十步，即圓城直徑。符合題意。}



\end{paracol}



\vspace{0.3cm}

\longline



%% ========================================================================

%% Problem 167 (Vol.12, Problem 17)

%% ========================================================================

\begin{paracol}{2}



\wenC{假令圓城一座，甲出南門直行，乙出東門直行。甲回望見乙，乃斜行與乙相會。只云甲行二百四十步，乙行八十步不及斜行一百六十步。問圓徑幾何？}



\daC{圓徑一百六十步。}



\fayueC{以甲行乘乙行，倍之為實。以甲行加乙行為從。一步常法。開平方得圓徑。}



\caoyueC{立天元一為圓徑。以甲行為股，乙行為勾。以勾股相乘得一萬九千二百步，倍之得三萬八千四百步為實。以甲行加乙行得三百二十步為從。一步常法。開平方得一百六十步，即圓徑也。合問。}



\switchcolumn



\wenM{假設有一座圓形城池，甲從南門直向行，乙從東門直向行。甲回頭望見乙，便斜行與乙會合。只知甲行二百四十步，乙行八十步，乙比斜行少一百六十步（即斜行比乙直行多一百六十步）。問圓城直徑多少？}



\daM{圓城直徑為一百六十步。}



\fayueM{以甲行步數乘乙行步數，加倍作為被開方數（實）。以甲行步數加乙行步數作為一次項係數（從）。以一步為二次項常數（常法）。開平方即得圓城直徑。}



\caoyueM{設天元一為圓城直徑。甲行為股，乙行為勾。以勾股相乘得一萬九千二百步，加倍得三萬八千四百步，作為實。以甲行加乙行得三百二十步，作為從。以一步為常法。開平方得一百六十步，即圓城直徑。符合題意。}



\end{paracol}



\vspace{0.3cm}

\longline



%% ========================================================================

%% Problem 168 (Vol.12, Problem 18)

%% ========================================================================

\begin{paracol}{2}



\wenC{假令圓城一座，甲乙二人同立於圓心。甲向東門直行，乙向南門直行。甲出東門望見乙，斜行與乙相會。只云甲直行一百八十步，乙直行一百二十步不及斜行六十步。問圓徑幾何？}



\daC{圓徑一百二十步。}



\fayueC{以甲行乘乙行，倍之為實。以甲行加乙行為從。一步常法。開平方得圓徑。}



\caoyueC{立天元一為圓徑。以甲行為股，乙行為勾。以勾股相乘得二萬一千六百步，倍之得四萬三千二百步為實。以甲行加乙行得三百步為從。一步常法。開平方得一百二十步，即圓徑也。合問。}



\switchcolumn



\wenM{假設有一座圓形城池，甲、乙兩人同立於圓心。甲向東門直行，乙向南門直行。甲出東門後望見乙，便斜行與乙會合。只知甲直行一百八十步，乙直行一百二十步，乙比斜行少六十步（即斜行比乙直行多六十步）。問圓城直徑多少？}



\daM{圓城直徑為一百二十步。}



\fayueM{以甲行步數乘乙行步數，加倍作為被開方數（實）。以甲行步數加乙行步數作為一次項係數（從）。以一步為二次項常數（常法）。開平方即得圓城直徑。}



\caoyueM{設天元一為圓城直徑。甲行為股，乙行為勾。以勾股相乘得二萬一千六百步，加倍得四萬三千二百步，作為實。以甲行加乙行得三百步，作為從。以一步為常法。開平方得一百二十步，即圓城直徑。符合題意。}



\end{paracol}



\vspace{0.3cm}

\longline



%% ========================================================================

%% Problem 169 (Vol.12, Problem 19)

%% ========================================================================

\begin{paracol}{2}



\wenC{假令圓城一座，甲出東門直行，乙出南門直行，丙出西門直行，丁出北門直行。只云甲行一百八十步，乙行一百二十步，丙行八十步，丁行四十步不及圓徑二十步。問圓徑幾何？}



\daC{圓徑一百六十步。}



\fayueC{以甲行乘乙行為實。以甲行加乙行為從。一步常法。開平方得圓徑。}



\caoyueC{立天元一為圓徑。以甲行為股，乙行為勾。以勾股相乘得二萬一千六百步為實。以甲行加乙行得三百步為從。一步常法。開平方得一百六十步，即圓徑也。合問。}



\switchcolumn



\wenM{假設有一座圓形城池，甲從東門直向行，乙從南門直向行，丙從西門直向行，丁從北門直向行。只知甲行一百八十步，乙行一百二十步，丙行八十步，丁行四十步且丁比圓徑少二十步（即圓徑比丁行多二十步）。問圓城直徑多少？}



\daM{圓城直徑為一百六十步。}



\fayueM{以甲行步數乘乙行步數，作為被開方數（實）。以甲行步數加乙行步數作為一次項係數（從）。以一步為二次項常數（常法）。開平方即得圓城直徑。}



\caoyueM{設天元一為圓城直徑。甲行為股，乙行為勾。以勾股相乘得二萬一千六百步，作為實。以甲行加乙行得三百步，作為從。以一步為常法。開平方得一百六十步，即圓城直徑。符合題意。}



\end{paracol}



\vspace{0.3cm}

\longline



%% ========================================================================

%% Problem 170 (Vol.12, Problem 20) —— FINAL PROBLEM

%% ========================================================================

\begin{paracol}{2}



\wenC{假令圓城一座，甲乙二人同立於圓心。甲向東門直行，乙向南門直行。甲出東門望見乙，斜行與乙相會。只云甲直行一百二十步，乙直行八十步，又云斜行與圓徑等。問圓徑幾何？}



\daC{圓徑八十步。}



\fayueC{以勾股冪相加為弦冪。以斜行為弦。以勾股相乘，倍之開平方得圓徑。}



\caoyueC{立天元一為圓徑即為弦。以甲行為股，乙行為勾。以勾冪六千四百步加股冪一萬四千四百步，共二萬八百步為弦冪。開平方得弦一百二十步。又以勾股相乘得九千六百步，倍之一萬九千二百步。以弦除之得一百六十步為實。一步常法。開平方得八十步，即圓徑也。合問。}



\switchcolumn



\wenM{假設有一座圓形城池，甲、乙兩人同立於圓心。甲向東門直行，乙向南門直行。甲出東門後望見乙，便斜行與乙會合。只知甲直行一百二十步，乙直行八十步，又知斜行與圓城直徑相等。問圓城直徑多少？}



\daM{圓城直徑為八十步。}



\fayueM{將勾的平方加上股的平方得到弦的平方。以斜行為弦。將勾股相乘，加倍後再開平方，即得圓城直徑。}



\caoyueM{設天元一為圓城直徑，也就是弦。甲行為股，乙行為勾。以勾冪六千四百步加上股冪一萬四千四百步，共二萬零八百步，作為弦冪。開平方得弦一百二十步。又以勾股相乘得九千六百步，加倍得一萬九千二百步。以弦除之得一百六十步，作為實。以一步為常法。開平方得八十步，即圓城直徑。符合題意。}



\end{paracol}



\vspace{0.3cm}

\longline





%% ========================================================================

%% Afterword

%% ========================================================================

\newpage

\section{後序 (Afterword) / 跋文}

\label{sec:afterword}



\longline



\begin{paracol}{2}

\leftlabel



\switchcolumn

\rightlabel



\end{paracol}



\vspace{0.3cm}

\longline



\begin{paracol}{2}



敬齋先生病且革，語其子克修曰：「吾平生著述，死後可盡燔去。獨測圓海鏡一書，雖九九小數，吾常精思致力焉，後世必有知者。庶可布廣垂永乎？」



\vspace{0.3cm}



先生於六藝百家，靡不貫串。文集近數百卷，常謙謙不自伐。惟於此書不忘稱異。予易簀之間，想有玄妙內得於心者。予以先生與先人同榜之故，素常兄事克修。克修兄命予重為序。予不敢說論艷藻，刻畫無鹽，唐突西子。直以所聞語之子弟，使知測圓海鏡之淵源云。



\switchcolumn



敬齋先生（李冶）病重臨終之際，對他的兒子克修說：「我平生的各種著述，死後可以全部燒掉。只有《測圓海鏡》這一部書，雖然只是關於算術的小技，但我曾經精心鑽研、傾注心血，後世必定有能理解它的人。但願它能夠廣為流傳、永久傳世吧？」



\vspace{0.3cm}



敬齋先生對於六藝百家之學，無不融會貫通。他的文集將近數百卷，但他始終謙遜，從不誇耀自己。唯獨對於這部書，他不忘稱讚其非凡之處。想來他在臨終之際，心中必定領悟到了某種深奧的玄妙。我因為先生與我的父親是同榜進士的緣故，一向像對待兄長一樣尊敬克修。克修兄命我重新寫一篇序。我不敢用華麗的辭藻來讚美，那樣做就像是給無鹽（醜女）畫上漂亮的妝容，反而是對西施（美女）的冒犯。我只是把所聽到的直接告訴後輩，讓他們知道《測圓海鏡》的來龍去脈罷了。



\end{paracol}



\vspace{1cm}

\longline



\vfill



\begin{center}

\fbox{\parbox{0.8\textwidth}{\centering

\textbf{\Large 測圓海鏡 卷十至卷十二 完}\\[0.5em]

{全書一百七十問終}\\[0.3em]

\textit{--- End of the Cèyuán Hǎijìng (Sea Mirror of Circle Measurements) ---}\\[0.3em]

{\small 文言文/白話文 雙語對照版 / Classical--Modern Chinese Bilingual Edition}

}}

\end{center}



\end{document}

%% ========================================================================

%% END OF DOCUMENT

%% ========================================================================







% ============================================================

% Source: translations/zh/chinese/li-ye-ceyuan-haijing-vols4-6_classical-modern_bilingual.tex

% ============================================================



% ===================================================================

% 《测圆海镜》卷四至卷六 —— 文言文↔白话文 对照本

% Ceyuan Haijing Volumes 4-6: Classical↔Modern Chinese Bilingual Edition

% ===================================================================

\documentclass[10pt,a4paper]{article}



%% -------- Core Packages --------

\usepackage{fontspec}

\usepackage{polyglossia}

\usepackage{amsmath,amssymb,amsthm}

\usepackage{geometry}

\usepackage{graphicx}

\usepackage{xcolor}

\usepackage{titlesec}

\usepackage{fancyhdr}

\usepackage{enumitem}

\usepackage{array,booktabs,longtable,multirow,colortbl}

\usepackage{hyperref}

\usepackage{microtype}

\usepackage{tocloft}

\usepackage{xeCJK}

\usepackage{paracol}



%% -------- Page Geometry (wider for two-column parallel text) --------

\geometry{

  paperwidth=210mm,paperheight=297mm,

  left=15mm,right=15mm,top=25mm,bottom=25mm,

  headheight=14pt,footskip=30pt

}



%% -------- Column Separation for paracol --------

\columnsep=8mm



%% -------- Language Setup --------

\setmainlanguage{english}



%% -------- Font Configuration --------

\setmainfont{Times New Roman}

\setsansfont{Arial}

\setmonofont{Courier New}



%% CJK Support via xeCJK — explicit path as required

\setCJKmainfont{Microsoft YaHei}

\setCJKsansfont{Microsoft YaHei}

\setCJKmonofont{Microsoft YaHei}



%% -------- Colors --------

\definecolor{titlecolor}{RGB}{45,55,72}

\definecolor{sectioncolor}{RGB}{55,65,81}

\definecolor{notecolor}{RGB}{120,53,15}

\definecolor{rulecolor}{RGB}{180,160,140}

\definecolor{volcolor}{RGB}{80,40,20}

\definecolor{leftbg}{RGB}{252,250,245}

\definecolor{rightbg}{RGB}{245,250,252}

\definecolor{leftborder}{RGB}{160,140,100}

\definecolor{rightborder}{RGB}{100,140,160}



%% -------- Section Formatting --------

\titleformat{\section}

  {\normalfont\Large\bfseries\color{volcolor}}

  {\thesection}{1em}{}

\titleformat{\subsection}

  {\normalfont\large\bfseries\color{sectioncolor}}

  {\thesubsection}{1em}{}

\titleformat{\subsubsection}

  {\normalfont\normalsize\bfseries\color{sectioncolor}}

  {\thesubsubsection}{1em}{}



%% -------- Header/Footer --------

\pagestyle{fancy}

\fancyhf{}

\fancyhead[L]{\small\textsc{Ceyuan Haijing — Classical↔Modern Chinese}}

\fancyhead[R]{\small\thepage}

\renewcommand{\headrulewidth}{0.4pt}

\renewcommand{\footrulewidth}{0pt}



%% -------- Custom Commands --------

\newcommand{\longline}{\noindent\rule{\textwidth}{0.4pt}}

\newcommand{\shortline}{\noindent\rule{0.5\textwidth}{0.4pt}\par}



%% Problem number command with Chinese numerals

\newcounter{probnum}

\newcommand{\chnum}[1]{%

  \ifcase#1 零\or 一\or 二\or 三\or 四\or 五\or 六\or 七\or 八\or 九\or 十\or

    十一\or 十二\or 十三\or 十四\or 十五\or 十六\or 十七\or 十八\fi}



%% Bilingual problem environment

%% LEFT = Classical Chinese, RIGHT = Modern Chinese

\newcommand{\biproblem}[2]{%

  \par\noindent

  \begin{minipage}[t]{0.48\textwidth}

    \textbf{第\chnum{\value{probnum}}问：}#1

  \end{minipage}%

  \hfill\vrule\hfill

  \begin{minipage}[t]{0.48\textwidth}

    \textbf{第\chnum{\value{probnum}}题：}#2

  \end{minipage}

  \par\vspace{0.3cm}

}



\newcommand{\bian answer}[2]{%

  \par\noindent

  \begin{minipage}[t]{0.48\textwidth}

    \textbf{答曰：}#1

  \end{minipage}%

  \hfill\vrule\hfill

  \begin{minipage}[t]{0.48\textwidth}

    \textbf{答：}#2

  \end{minipage}

  \par\vspace{0.2cm}

}



\newcommand{\biworking}[2]{%

  \par\noindent

  \begin{minipage}[t]{0.48\textwidth}

    \textbf{草曰：}#1

  \end{minipage}%

  \hfill\vrule\hfill

  \begin{minipage}[t]{0.48\textwidth}

    \textbf{演算：}#2

  \end{minipage}

  \par\vspace{0.2cm}

}



\newcommand{\bimethod}[2]{%

  \par\noindent

  \begin{minipage}[t]{0.48\textwidth}

    \textbf{法曰：}#1

  \end{minipage}%

  \hfill\vrule\hfill

  \begin{minipage}[t]{0.48\textwidth}

    \textbf{算法：}#2

  \end{minipage}

  \par\vspace{0.3cm}

}



%% Simplified combined command for problems with integrated method

\newcommand{\problemmethod}[2]{%

  \par\noindent

  \begin{minipage}[t]{0.48\textwidth}

    #1

  \end{minipage}%

  \hfill\vrule\hfill

  \begin{minipage}[t]{0.48\textwidth}

    #2

  \end{minipage}

  \par\vspace{0.3cm}

}



%% Section headers bilingual

\newcommand{\bisectionheader}[4]{%

  \begin{center}

  {\Large\bfseries\color{volcolor} #1}\\[0.3cm]

  {\normalsize\color{sectioncolor} #2}

  \end{center}

  \vspace{0.3cm}

  \noindent

  \begin{minipage}[t]{0.48\textwidth}

    \textbf{【题意总说】}#3

  \end{minipage}%

  \hfill\vrule\hfill

  \begin{minipage}[t]{0.48\textwidth}

    \textbf{【题意概述】}#4

  \end{minipage}

  \par\vspace{0.3cm}

}



%% -------- Hyperref Setup --------

\hypersetup{

  colorlinks=true,

  linkcolor=sectioncolor,

  citecolor=sectioncolor,

  urlcolor=sectioncolor,

  pdfencoding=auto,

  pdftitle={测圆海镜卷四至卷六 文言文↔白话文对照本},

  pdfauthor={李冶 (Li Ye)}

}



%% -------- TOC Depth --------

\setcounter{tocdepth}{3}

\setcounter{secnumdepth}{3}



%% ===================================================================



\begin{document}



%% ===================================================================

%% TITLE PAGE

%% ===================================================================

\begin{titlepage}

\begin{center}

\vspace*{2cm}



{\fontsize{14}{18}\selectfont 钦定四库全书}



\vspace{1.5cm}



{\fontsize{12}{16}\selectfont 子部}



\vspace{2cm}



\rule{0.8\textwidth}{1.5pt}



\vspace{1cm}



{\fontsize{26}{34}\selectfont\bfseries\color{titlecolor} 测圆海镜}



\vspace{0.3cm}



{\fontsize{16}{22}\selectfont 卷四至卷六}



\vspace{0.3cm}



{\fontsize{13}{17}\selectfont\color{notecolor} 文言文 \raisebox{-0.5ex}{$\leftrightarrow$} 白话文 对照本}



\vspace{1cm}



\rule{0.8\textwidth}{1.5pt}



\vspace{2cm}



{\fontsize{14}{18}\selectfont 元\quad 李冶\quad 撰}



\vspace{0.5cm}



{\fontsize{11}{15}\selectfont\color{sectioncolor} 白话文译解}



\vspace{3cm}



\begin{tabular}{cc}

\textbf{左栏} & \textbf{右栏} \\

\textit{文言文（原文）} & \textit{白话文（今译）} \\

\end{tabular}



\vfill



{\small 依《钦定四库全书》本 \quad 现代汉语译注}



\end{center}

\end{titlepage}



%% ===================================================================

%% TABLE OF CONTENTS

%% ===================================================================

\tableofcontents

\newpage



%% ===================================================================

%% INTRODUCTION (BILINGUAL)

%% ===================================================================

\section*{引言}

\addcontentsline{toc}{section}{引言}



\noindent

\begin{minipage}[t]{0.48\textwidth}

\textbf{【原文】}



《测圆海镜》十二卷，元李冶撰。冶字仁卿，号敬斋，真定栾城人。

金末进士，入元官翰林学士。此书乃其所著天元术之杰作，以勾股测圆

之法，设问答以明天元一之术。全书凡一百七十问，皆设城池圆形，以

人物行步之数推求圆径。



本编收录卷四、卷五、卷六三卷：

\begin{itemize}[leftmargin=1em]

  \item \textbf{卷四}：底勾一十七问

  \item \textbf{卷五}：大股一十八问  

  \item \textbf{卷六}：大勾一十八问

\end{itemize}



书中设问之例：假令有圆城一座，甲乙二人各从城门出行，或东或南，

行若干步，遥望见之，或斜行相会。只知若干步数，余皆不知，以求

圆城之径。每问皆有\textbf{问}、\textbf{答}、\textbf{草}、\textbf{法}四部分。



\textbf{术语解释：}

\begin{itemize}[leftmargin=1em]

  \item \textbf{天元一}：设未知数为一，即现代代数之$x$

  \item \textbf{寄左}：将某式暂存一边，相当于方程之一边

  \item \textbf{相消}：两式相减，相当于移项合并

  \item \textbf{带分母}：某式含有分母未除

  \item \textbf{幂}：平方，即现代记号$x^2$

  \item \textbf{开方}：开平方，即求平方根

  \item \textbf{开立方}：开立方，即求立方根

  \item \textbf{实}：常数项

  \item \textbf{从}：一次项系数

  \item \textbf{廉}：二次项系数

  \item \textbf{隅}：三次项系数（立方方程中）

\end{itemize}

\end{minipage}%

\hfill\vrule\hfill

\begin{minipage}[t]{0.48\textwidth}

\textbf{【今译】}



《测圆海镜》共十二卷，由元代数学家李冶撰写。李冶字仁卿，号敬斋，

真定府栾城县人。金朝末年考中进士，元朝时任翰林学士。这部书是他

研究天元术的杰作，运用勾股定理测量圆形的方法，通过设置问题和

解答来阐明天元术的奥秘。全书共170道题，都假设有一座圆形城池，

根据人物行走的步数来推算城池的直径。



本版本收录第四、第五、第六卷：

\begin{itemize}[leftmargin=1em]

  \item \textbf{卷四}：底勾17题

  \item \textbf{卷五}：大股18题

  \item \textbf{卷六}：大勾18题

\end{itemize}



书中设题的模式：假设有一座圆形城池，甲、乙两人分别从城门出发，

或向东或向南行走若干步，远远望见对方，或斜向行走相会。只给定

某些步数，其余未知，要求求出圆形城池的直径。每道题都有

\textbf{问题}、\textbf{答案}、\textbf{演算}、\textbf{算法}四部分。



\textbf{术语解释：}

\begin{itemize}[leftmargin=1em]

  \item \textbf{天元一}：设未知数为1，相当于现代代数的$x$

  \item \textbf{寄左}：将某个算式暂存一边，相当于方程的一边

  \item \textbf{相消}：两个算式相减，相当于移项合并同类项

  \item \textbf{带分母}：某个算式含有分母尚未除尽

  \item \textbf{幂}：平方，即现代记号$x^2$

  \item \textbf{开方}：开平方，即求平方根

  \item \textbf{开立方}：开立方，即求立方根

  \item \textbf{实}：常数项

  \item \textbf{从}：一次项系数

  \item \textbf{廉}：二次项系数

  \item \textbf{隅}：三次项系数（立方方程中）

\end{itemize}

\end{minipage}



\vspace{0.5cm}

\longline



\newpage



%% ===================================================================

%% FORMAT NOTES

%% ===================================================================

\noindent

\begin{minipage}[t]{0.48\textwidth}

\textbf{【体例说明】}



左栏为《测圆海镜》原文，依四库全书本。

右栏为白话文翻译，力求准确通顺。



数学算式保持原书格式，以天元术排列

多项式。必要时加注说明。

\end{minipage}%

\hfill\vrule\hfill

\begin{minipage}[t]{0.48\textwidth}

\textbf{【格式说明】}



左栏为原文，依据《四库全书》版本。

右栏为白话文翻译，力求准确、通顺。



数学算式保持原书格式，采用天元术的

方式排列多项式。必要时添加注释说明。

\end{minipage}



\vspace{0.5cm}

\longline



\newpage



%% ===================================================================

%% VOLUME 4

%% ===================================================================

\section{测圆海镜卷四 —— 底勾一十七问}

\label{sec:vol4}



\bisectionheader

  {底勾一十七问}

  {十七道关于底勾的数学题}

  {卷四论「底勾」之术。所谓底勾，乃勾股形之最下勾股，与圆径相切之基本勾股形也。凡一十七问，皆以底勾为基，求圆城之径。}

  {卷四讨论「底勾」的算法。所谓底勾，是指勾股三角形中最下面的勾股边，是与圆直径相切的基本勾股三角形。共17道题，都以底勾为基础，求圆形城池的直径。}



\vspace{0.3cm}



%% Reset problem counter for Volume 4

\setcounter{probnum}{0}



%% --- Problem 1 ---

\stepcounter{probnum}

\biproblem

  {或问乙出南门东行不知步数而立，甲出北门东行二百步见之，就乙斜行二百七十二步与乙相会。问答同前。}

  {假设乙从南门出发向东行走若干步停下，甲从北门出发向东行走200步后看见乙，接着向乙斜行272步与乙相会。问题与答案同前。}



\bian answer

  {城径二百四十步。}

  {城池直径为240步。}



\biworking

  {识别得二行相减余即乙出南门东行数也。以甲东行减于就乙斜行余七十二步，以乘甲东行步，得一万四千四百步，又四之，得五万七千六百步为实，以平方开之，得二百四十步，即城径也。}

  {通过分析可知，两次行走的差就是乙出南门向东行走的步数。用甲向东行走的步数从斜行步数中减去，余72步，乘以甲向东行走的步数，得到14400步，再乘以4，得到57600步作为被开方数（实），开平方得到240步，即城池的直径。}



\bimethod

  {二行差数乘甲东行又四之，为平方实，得全径。}

  {两次行走的差乘以甲向东行走的步数再乘以4，作为平方开方的被开方数，求得全城直径。}



%% --- Problem 2 ---

\stepcounter{probnum}

\biproblem

  {或问乙出南门东行不知步数而立，甲出北门东行二百步见之。问答同前。}

  {假设乙从南门出发向东行走若干步停下，甲从北门出发向东行走200步后看见乙。问题与答案同前。}



\bian answer

  {城径二百四十步。}

  {城池直径为240步。}



\biworking

  {立天元一为城径，以减于二之甲东行步，得\raisebox{-0.5ex}{\scriptsize$\begin{array}{c}480\\\hline 1\end{array}$}为两个小差，以乙南行步乘之，得\raisebox{-0.5ex}{\scriptsize$\begin{array}{c}48000\\\hline 1\end{array}$}为城径幂，寄左。然后以天元幂\raisebox{-0.5ex}{\scriptsize$\begin{array}{c}0\\1\end{array}$}与左相消，得\raisebox{-0.5ex}{\scriptsize$\begin{array}{c|c|c}0&0&48000\\\hline 1&0&\end{array}$}以平方开之，得二百四十步，即城径也。}

  {设天元一（未知数）为城池直径，用两倍的甲向东行走步数减去它，得到$\begin{pmatrix}480\\1\end{pmatrix}$为两个小差，乘以乙向南行走的步数，得到$\begin{pmatrix}48000\\1\end{pmatrix}$作为城径的平方，暂存左边。然后将天元的平方$\begin{pmatrix}0\\1\end{pmatrix}$与左边相消，得到$\begin{pmatrix}0&0&48000\\1&0\end{pmatrix}$，开平方得到240步，即城池直径。}



\bimethod

  {半之乙南行步乘甲东行为实，半乙南行为从，一步常法，得半径。}

  {取乙南行步数的一半乘以甲向东行走的步数作为被开方数，取乙南行步数的一半作为一次项系数，二次项系数为1，开方求得半径。}



%% --- Problem 3 ---

\stepcounter{probnum}

\biproblem

  {或问乙从坤隅南行三百六十步，甲出北门东行二百步见之。问答同前。}

  {假设乙从坤隅（西南角）向南行走360步，甲从北门出发向东行走200步后看见乙。问题与答案同前。}



\bian answer

  {城径二百四十步。}

  {城池直径为240步。}



\biworking

  {立天元一为半径，减于乙东行，得\raisebox{-0.5ex}{\scriptsize$\begin{array}{c}120\\1\end{array}$}为小差。半乙南行步，得一百八十步，以乘小差，得\raisebox{-0.5ex}{\scriptsize$\begin{array}{c}21600\\\hline 1\end{array}$}为半径幂，寄左。然后以天元幂\raisebox{-0.5ex}{\scriptsize$\begin{array}{c}0\\1\end{array}$}与左相消，得\raisebox{-0.5ex}{\scriptsize$\begin{array}{c|c|c}0&0&21600\\\hline 1&0&\end{array}$}以平方开之，得一百二十步，倍之得全径。}

  {设天元一为半径，从乙向东行走的距离中减去，得到$\begin{pmatrix}120\\1\end{pmatrix}$为小差。取乙向南行走步数的一半，得180步，乘以小差，得到$\begin{pmatrix}21600\\1\end{pmatrix}$作为半径的平方，暂存左边。然后将天元的平方$\begin{pmatrix}0\\1\end{pmatrix}$与左边相消，得到$\begin{pmatrix}0&0&21600\\1&0\end{pmatrix}$，开平方得到120步，乘以2得到全城直径。}



\bimethod

  {两行步相乘为实，甲东行为从，乙为隅，得半径。}

  {两行的步数相乘作为被开方数，甲向东行走的步数为一次项系数，乙的行步数为二次项系数，开方求得半径。}



%% --- Problem 4 ---

\stepcounter{probnum}

\biproblem

  {或问乙从坤隅南行三百六十步，甲出北门东行二百步见之。问答同前。}

  {假设乙从坤隅（西南角）向南行走360步，甲从北门出发向东行走200步后看见乙。问题与答案同前。}



\bian answer

  {城径二百四十步。}

  {城池直径为240步。}



\biworking

  {立天元一为城径，以减于二之甲东行，余即乙出南门东行数也。}

  {设天元一为城池直径，用两倍的甲向东行走步数减去它，余数就是乙出南门向东行走的步数。}



\bimethod

  {以乙南行步乘甲东行幂，又四之为实，从空，乙行为廉，一步常法，得城径。}

  {用乙向南行走的步数乘以甲向东行走步数的平方，再乘以4作为被开方数，一次项为空，乙的行步数为二次项系数，三次项系数为1，开方求得城径。}



%% --- Problem 5 ---

\stepcounter{probnum}

\biproblem

  {或问乙出南门直行一百三十五步，甲出北门东行二百步见之。问答同前。}

  {假设乙从南门出发向南直走135步，甲从北门出发向东行走200步后看见乙。问题与答案同前。}



\bian answer

  {城径二百四十步。}

  {城池直径为240步。}



\biworking

  {立天元一为城径，加乙南行，得\raisebox{-0.5ex}{\scriptsize$\begin{array}{c}135\\1\end{array}$}为股率。其甲东行即勾率也。其乙南行为小股，以勾率乘之，合以大弦率除，今不受除，便以此为小勾，为母，寄股率。乃先以小弦乘大和，得下式，寄左。又以倍天元减斜步，得为小和，以乘大弦，得同数，与左相消。}

  {设天元一为城池直径，加上乙向南行走的步数，得到$\begin{pmatrix}135\\1\end{pmatrix}$为股率。甲向东行走的步数就是勾率。乙向南行走的步数为小股，用勾率乘以它，应当用大弦率来除，现在先不除，就把这作为小勾，作为分母，暂寄股率。先用小弦乘以大和，得到下式，暂存左边。又用两倍的天元减去斜行步数，得到小和，乘以大弦，得到相同的数，与左边相消。}



\bimethod

  {以乙南行步乘甲东行幂，又四之为实，从空，乙行为廉，一步常法，得城径。}

  {用乙向南行走的步数乘以甲向东行走步数的平方，再乘以4作为被开方数，一次项为空，乙的行步数为二次项系数，三次项系数为1，开方求得城径。}



\vspace{0.3cm}

\noindent

\begin{minipage}[t]{0.48\textwidth}

\textit{（以下各问草法类推，皆立天元一为半径或城径，以各条件建立方程，与左式相消，开方得数。）}

\end{minipage}%

\hfill\vrule\hfill

\begin{minipage}[t]{0.48\textwidth}

\textit{（以下各题演算方法类似，都设天元一为半径或城径，根据各题条件建立方程，与左边算式相消，开方得到答案。）}

\end{minipage}



\vspace{0.3cm}



%% --- Problem 6 ---

\stepcounter{probnum}

\biproblem

  {或问乙从坤隅南行三百六十步，甲出北门东行二百步见之，就乙斜行二百七十二步与乙相会。问答同前。}

  {假设乙从坤隅（西南角）向南行走360步，甲从北门出发向东行走200步后看见乙，接着向乙斜行272步与乙相会。问题与答案同前。}



\bimethod

  {两行步相乘倍之为实，乙南行为从，一步常法，得半径。}

  {两行步数相乘再乘以2作为被开方数，乙向南行走的步数为一次项系数，二次项系数为1，开方求得半径。}



%% --- Problem 7 ---

\stepcounter{probnum}

\biproblem

  {或问乙出南门直行一百三十五步，甲出北门东行二百步见之。问答同前。}

  {假设乙从南门出发向南直走135步，甲从北门出发向东行走200步后看见乙。问题与答案同前。}



\bimethod

  {以乙南行乘甲东行幂为实，从空，乙行为廉，一步常法，得城径。}

  {用乙向南行走的步数乘以甲向东行走步数的平方作为被开方数，一次项为空，乙的行步数为二次项系数，三次项系数为1，开方求得城径。}



%% --- Problem 8 ---

\stepcounter{probnum}

\biproblem

  {或问乙出南门直行一百三十五步，甲出北门东行二百步见之。问答同前。}

  {假设乙从南门出发向南直走135步，甲从北门出发向东行走200步后看见乙。问题与答案同前。}



\bimethod

  {倍乙南行乘甲东行幂为实，从空，倍乙行为廉，一步常法，得城径。}

  {乙向南行走步数的2倍乘以甲向东行走步数的平方作为被开方数，一次项为空，乙行步数的2倍为二次项系数，三次项系数为1，开方求得城径。}



%% --- Problem 9 ---

\stepcounter{probnum}

\biproblem

  {或问乙出南门直行一百三十五步，甲出北门东行二百步，就乙斜行二百七十二步与乙相会。问答同前。}

  {假设乙从南门出发向南直走135步，甲从北门出发向东行走200步，接着向乙斜行272步与乙相会。问题与答案同前。}



\bimethod

  {倍乙南行乘甲东行幂为实，甲东行为从，倍乙行为廉，一步常法，得城径。}

  {乙向南行走步数的2倍乘以甲向东行走步数的平方作为被开方数，甲向东行走的步数为一次项系数，乙行步数的2倍为二次项系数，三次项系数为1，开方求得城径。}



%% --- Problem 10 ---

\stepcounter{probnum}

\biproblem

  {或问乙出南门直行一百三十五步，甲出北门东行二百步见之，就乙斜行二百七十二步与乙相会。问答同前。}

  {假设乙从南门出发向南直走135步，甲从北门出发向东行走200步后看见乙，接着向乙斜行272步与乙相会。问题与答案同前。}



\bimethod

  {两行相减余乘甲东行又四之为实，从空，两行差为廉，一步常法，得城径。}

  {两行步数相减的余数乘以甲向东行走步数再乘以4作为被开方数，一次项为空，两行步数的差为二次项系数，三次项系数为1，开方求得城径。}



%% --- Problem 11 ---

\stepcounter{probnum}

\biproblem

  {或问乙出南门直行一百三十五步，甲出北门东行二百步见之。问答同前。}

  {假设乙从南门出发向南直走135步，甲从北门出发向东行走200步后看见乙。问题与答案同前。}



\bimethod

  {以乙南行步乘甲东行幂，又四之为实，从空，乙行为廉，一步常法，得城径。}

  {用乙向南行走的步数乘以甲向东行走步数的平方，再乘以4作为被开方数，一次项为空，乙的行步数为二次项系数，三次项系数为1，开方求得城径。}



%% --- Problem 12 ---

\stepcounter{probnum}

\biproblem

  {或问乙出南门直行一百三十五步，甲出北门东行二百步见之。问答同前。}

  {假设乙从南门出发向南直走135步，甲从北门出发向东行走200步后看见乙。问题与答案同前。}



\bimethod

  {倍乙南行乘甲东行幂为实，从空，倍乙行为廉，一步常法，得城径。}

  {乙向南行走步数的2倍乘以甲向东行走步数的平方作为被开方数，一次项为空，乙行步数的2倍为二次项系数，三次项系数为1，开方求得城径。}



%% --- Problem 13 ---

\stepcounter{probnum}

\biproblem

  {或问乙出南门直行一百三十五步，甲出北门东行二百步，就乙斜行二百七十二步与乙相会。问答同前。}

  {假设乙从南门出发向南直走135步，甲从北门出发向东行走200步，接着向乙斜行272步与乙相会。问题与答案同前。}



\bimethod

  {两行相减余乘甲东行又四之为实，从空，两行差为廉，一步常法，得城径。}

  {两行步数相减的余数乘以甲向东行走步数再乘以4作为被开方数，一次项为空，两行步数的差为二次项系数，三次项系数为1，开方求得城径。}



%% --- Problem 14 ---

\stepcounter{probnum}

\biproblem

  {或问乙出南门直行一百三十五步，甲出北门东行二百步见之。问答同前。}

  {假设乙从南门出发向南直走135步，甲从北门出发向东行走200步后看见乙。问题与答案同前。}



\bimethod

  {以乙南行乘甲东行幂为实，从空，乙行为廉，一步常法，得城径。}

  {用乙向南行走的步数乘以甲向东行走步数的平方作为被开方数，一次项为空，乙的行步数为二次项系数，三次项系数为1，开方求得城径。}



%% --- Problem 15 ---

\stepcounter{probnum}

\biproblem

  {或问乙出南门直行一百三十五步，甲出北门东行二百步见之。问答同前。}

  {假设乙从南门出发向南直走135步，甲从北门出发向东行走200步后看见乙。问题与答案同前。}



\bimethod

  {倍乙南行乘甲东行幂为实，从空，倍乙行为廉，一步常法，得城径。}

  {乙向南行走步数的2倍乘以甲向东行走步数的平方作为被开方数，一次项为空，乙行步数的2倍为二次项系数，三次项系数为1，开方求得城径。}



%% --- Problem 16 ---

\stepcounter{probnum}

\biproblem

  {或问乙出南门直行一百三十五步，甲出北门东行二百步，就乙斜行二百七十二步与乙相会。问答同前。}

  {假设乙从南门出发向南直走135步，甲从北门出发向东行走200步，接着向乙斜行272步与乙相会。问题与答案同前。}



\bimethod

  {两行相减余乘甲东行又四之为实，从空，两行差为廉，一步常法，得城径。}

  {两行步数相减的余数乘以甲向东行走步数再乘以4作为被开方数，一次项为空，两行步数的差为二次项系数，三次项系数为1，开方求得城径。}



%% --- Problem 17 ---

\stepcounter{probnum}

\biproblem

  {或问乙出南门直行一百三十五步，甲出北门东行二百步见之。问答同前。}

  {假设乙从南门出发向南直走135步，甲从北门出发向东行走200步后看见乙。问题与答案同前。}



\bimethod

  {以乙南行步乘甲东行幂，又四之为实，从空，乙行为廉，一步常法，得城径。}

  {用乙向南行走的步数乘以甲向东行走步数的平方，再乘以4作为被开方数，一次项为空，乙的行步数为二次项系数，三次项系数为1，开方求得城径。}



\vspace{0.5cm}

\longline

\begin{center}

\textit{测圆海镜卷四终}

\end{center}



\newpage



%% ===================================================================

%% VOLUME 5

%% ===================================================================

\section{测圆海镜卷五 —— 大股一十八问}

\label{sec:vol5}



\bisectionheader

  {大股一十八问}

  {十八道关于大股的数学题}

  {卷五论「大股」之术。所谓大股，乃最大勾股形之股边，即通股也。以大股为主，配以诸勾股形之关系，求圆城之径。凡一十八问。}

  {卷五讨论「大股」的算法。所谓大股，是指最大勾股三角形的股边，即通股。以大股为主要对象，配合各种勾股三角形的关系，求圆形城池的直径。共18道题。}



\vspace{0.3cm}



%% Reset problem counter for Volume 5

\setcounter{probnum}{0}



%% --- Problem 1 ---

\stepcounter{probnum}

\biproblem

  {或问乙出南门直行一百三十五步，而立甲从乾隅南行六百步，斜望乙与城参相直。问答同前。}

  {假设乙从南门出发向南直走135步停下，甲从乾隅（西北角）向南行走600步，斜向遥望乙，视线与城池边缘恰好在一条直线上。问题与答案同前。}



\bian answer

  {城径二百四十步。}

  {城池直径为240步。}



\biworking

  {立天元一为半径，以二之减于甲南行，得\raisebox{-0.5ex}{\scriptsize$\begin{array}{c}600\\2\end{array}$}为大差也。以自之，得\raisebox{-0.5ex}{\scriptsize$\begin{array}{c}360000\\2400\\4\end{array}$}为大差幂也。置甲南行幂内减大差幂，而半之，得大弦也。内带大差为分母。又置甲南行幂内减大差幂，而半之，得大勾也。带大差分母。又以大差乘股六百步，併入和，得下式，为小和，以乘大弦。寄左。又以倍天元减斜步，得小和，以乘大弦，得同数，与左相消。开立方得一百二十步，即半径也。}

  {设天元一为半径，用甲向南行走步数减去两倍的天元，得到$\begin{pmatrix}600\\2\end{pmatrix}$为大差。将其平方，得到$\begin{pmatrix}360000\\2400\\4\end{pmatrix}$为大差的平方。将甲向南行走步数的平方减去大差的平方，再取一半，得到大弦。内含大差作为分母。又将甲向南行走步数的平方减去大差的平方，再取一半，得到大勾。也带大差分母。又用大差乘以股600步，合并入和，得到下式，作为小和，乘以大弦。暂存左边。又用两倍的天元减去斜行步数，得到小和，乘以大弦，得到相同的数，与左边相消。开立方得到120步，即半径。}



\bimethod

  {倍二行差内减甲南行步，复以乘甲南行步为实。倍二行差减甲南行步即甲南行也。四之甲南行步内减于甲南行余二百四十步，即城径也。}

  {两倍的两行差减去甲向南行走步数，再乘以甲向南行走步数作为被开方数。两倍的两行差减去甲向南行走步数就是甲向南行走步数。四倍的甲向南行走步数减去甲向南行走步数，余240步，即城池直径。}



%% --- Problem 2 ---

\stepcounter{probnum}

\biproblem

  {或问乙出南门直行一百三十五步，甲从乾隅南行六百步，斜望乙与城参相直。问答同前。}

  {假设乙从南门出发向南直走135步，甲从乾隅（西北角）向南行走600步，斜向遥望乙，视线与城池边缘恰好在一条直线上。问题与答案同前。}



\bian answer

  {城径二百四十步。}

  {城池直径为240步。}



\biworking

  {立天元一为半径，以减于甲南行步，得大差。置甲南行幂内减大差幂，而半之，得大弦也。又以大差乘股六百步，併入和，得下式，为小和，以乘大弦。寄左。又以倍天元减斜步，得小和，以乘大弦，得同数，与左相消。开立方得半径。}

  {设天元一为半径，从甲向南行走步数中减去，得到大差。将甲向南行走步数的平方减去大差的平方，再取一半，得到大弦。又用大差乘以股600步，合并入和，得到下式，作为小和，乘以大弦。暂存左边。又用两倍的天元减去斜行步数，得到小和，乘以大弦，得到相同的数，与左边相消。开立方得到半径。}



\bimethod

  {两行步相乘为实，从空，两行步为廉，一步常法，得城径。}

  {两行步数相乘作为被开方数，一次项为空，两行步数为二次项系数，三次项系数为1，开方求得城径。}



%% --- Problem 3 ---

\stepcounter{probnum}

\biproblem

  {或问乙出南门直行一百三十五步，甲从乾隅南行六百步望乙。问答同前。}

  {假设乙从南门出发向南直走135步，甲从乾隅（西北角）向南行走600步遥望乙。问题与答案同前。}



\bian answer

  {城径二百四十步。}

  {城池直径为240步。}



\biworking

  {立天元一为城径，以二之减于甲南行，余为大差。以自之得大差幂。置甲南行幂内减大差幂而半之，得大弦。内带大差为分母。又置甲南行幂内减大差幂而半之，得大勾。带大差分母。又以大差乘股六百步，併入和，得下式，为小和，以乘大弦。寄左。又以倍天元减斜步，得小和，以乘大弦，得同数，与左相消。开立方得一百二十步，即半径。}

  {设天元一为城池直径，用甲向南行走步数减去两倍的天元，余数为大差。将其平方得到大差的平方。将甲向南行走步数的平方减去大差的平方，再取一半，得到大弦。内含大差作为分母。又将甲向南行走步数的平方减去大差的平方，再取一半，得到大勾。也带大差分母。又用大差乘以股600步，合并入和，得到下式，作为小和，乘以大弦。暂存左边。又用两倍的天元减去斜行步数，得到小和，乘以大弦，得到相同的数，与左边相消。开立方得到120步，即半径。}



\bimethod

  {倍二行差内减甲南行步，复以乘甲南行步为实。倍二行差减甲南行步即甲南行也。四之甲南行步内减于甲南行余，即城径也。}

  {两倍的两行差减去甲向南行走步数，再乘以甲向南行走步数作为被开方数。两倍的两行差减去甲向南行走步数就是甲向南行走步数。四倍的甲向南行走步数减去甲向南行走步数，余数即城池直径。}



\vspace{0.3cm}

\noindent

\begin{minipage}[t]{0.48\textwidth}

\textit{（以下各问草法类推，皆依大股之术，以大差为主，建立方程。）}

\end{minipage}%

\hfill\vrule\hfill

\begin{minipage}[t]{0.48\textwidth}

\textit{（以下各题演算方法类似，都依照大股的算法，以大差为主要对象，建立方程。）}

\end{minipage}



\vspace{0.3cm}



%% --- Problem 4 ---

\stepcounter{probnum}

\biproblem

  {或问乙出南门直行一百三十五步，甲从乾隅南行六百步望乙。问答同前。}

  {假设乙从南门出发向南直走135步，甲从乾隅（西北角）向南行走600步遥望乙。问题与答案同前。}



\bimethod

  {两行步相乘为实，从空，两行步为廉，一步常法，得城径。}

  {两行步数相乘作为被开方数，一次项为空，两行步数为二次项系数，三次项系数为1，开方求得城径。}



%% --- Problem 5 ---

\stepcounter{probnum}

\biproblem

  {或问乙出南门直行一百三十五步，甲从乾隅南行六百步，斜望乙与城参相直。问答同前。}

  {假设乙从南门出发向南直走135步，甲从乾隅（西北角）向南行走600步，斜向遥望乙，视线与城池边缘恰好在一条直线上。问题与答案同前。}



\bimethod

  {倍二行差内减甲南行步，复以乘甲南行步为实，从空，倍二行差减甲南行步为廉，一步常法，得城径。}

  {两倍的两行差减去甲向南行走步数，再乘以甲向南行走步数作为被开方数，一次项为空，两倍的两行差减去甲向南行走步数作为二次项系数，三次项系数为1，开方求得城径。}



%% --- Problem 6 ---

\stepcounter{probnum}

\biproblem

  {或问乙出南门直行一百三十五步，甲从乾隅南行六百步望乙。问答同前。}

  {假设乙从南门出发向南直走135步，甲从乾隅（西北角）向南行走600步遥望乙。问题与答案同前。}



\bimethod

  {两行步相乘为实，从空，两行步为廉，一步常法，得城径。}

  {两行步数相乘作为被开方数，一次项为空，两行步数为二次项系数，三次项系数为1，开方求得城径。}



%% --- Problem 7 ---

\stepcounter{probnum}

\biproblem

  {或问乙出南门直行一百三十五步，甲从乾隅南行六百步，斜望乙与城参相直。问答同前。}

  {假设乙从南门出发向南直走135步，甲从乾隅（西北角）向南行走600步，斜向遥望乙，视线与城池边缘恰好在一条直线上。问题与答案同前。}



\bimethod

  {倍二行差内减甲南行步，复以乘甲南行步为实，从空，倍二行差减甲南行步为廉，一步常法，得城径。}

  {两倍的两行差减去甲向南行走步数，再乘以甲向南行走步数作为被开方数，一次项为空，两倍的两行差减去甲向南行走步数作为二次项系数，三次项系数为1，开方求得城径。}



%% --- Problem 8 ---

\stepcounter{probnum}

\biproblem

  {或问乙出南门直行一百三十五步，甲从乾隅南行六百步望乙。问答同前。}

  {假设乙从南门出发向南直走135步，甲从乾隅（西北角）向南行走600步遥望乙。问题与答案同前。}



\bimethod

  {两行步相乘为实，从空，两行步为廉，一步常法，得城径。}

  {两行步数相乘作为被开方数，一次项为空，两行步数为二次项系数，三次项系数为1，开方求得城径。}



%% --- Problem 9 ---

\stepcounter{probnum}

\biproblem

  {或问乙出南门直行一百三十五步，甲从乾隅南行六百步，斜望乙与城参相直。问答同前。}

  {假设乙从南门出发向南直走135步，甲从乾隅（西北角）向南行走600步，斜向遥望乙，视线与城池边缘恰好在一条直线上。问题与答案同前。}



\bimethod

  {倍二行差内减甲南行步，复以乘甲南行步为实，从空，倍二行差减甲南行步为廉，一步常法，得城径。}

  {两倍的两行差减去甲向南行走步数，再乘以甲向南行走步数作为被开方数，一次项为空，两倍的两行差减去甲向南行走步数作为二次项系数，三次项系数为1，开方求得城径。}



%% --- Problem 10 ---

\stepcounter{probnum}

\biproblem

  {或问乙出南门直行一百三十五步，甲从乾隅南行六百步望乙。问答同前。}

  {假设乙从南门出发向南直走135步，甲从乾隅（西北角）向南行走600步遥望乙。问题与答案同前。}



\bimethod

  {两行步相乘为实，从空，两行步为廉，一步常法，得城径。}

  {两行步数相乘作为被开方数，一次项为空，两行步数为二次项系数，三次项系数为1，开方求得城径。}



%% --- Problem 11 ---

\stepcounter{probnum}

\biproblem

  {或问乙出南门直行一百三十五步，甲从乾隅南行六百步，斜望乙与城参相直。问答同前。}

  {假设乙从南门出发向南直走135步，甲从乾隅（西北角）向南行走600步，斜向遥望乙，视线与城池边缘恰好在一条直线上。问题与答案同前。}



\bimethod

  {倍二行差内减甲南行步，复以乘甲南行步为实，从空，倍二行差减甲南行步为廉，一步常法，得城径。}

  {两倍的两行差减去甲向南行走步数，再乘以甲向南行走步数作为被开方数，一次项为空，两倍的两行差减去甲向南行走步数作为二次项系数，三次项系数为1，开方求得城径。}



%% --- Problem 12 ---

\stepcounter{probnum}

\biproblem

  {或问乙出南门直行一百三十五步，甲从乾隅南行六百步望乙。问答同前。}

  {假设乙从南门出发向南直走135步，甲从乾隅（西北角）向南行走600步遥望乙。问题与答案同前。}



\bimethod

  {两行步相乘为实，从空，两行步为廉，一步常法，得城径。}

  {两行步数相乘作为被开方数，一次项为空，两行步数为二次项系数，三次项系数为1，开方求得城径。}



%% --- Problem 13 ---

\stepcounter{probnum}

\biproblem

  {或问乙出南门直行一百三十五步，甲从乾隅南行六百步，斜望乙与城参相直。问答同前。}

  {假设乙从南门出发向南直走135步，甲从乾隅（西北角）向南行走600步，斜向遥望乙，视线与城池边缘恰好在一条直线上。问题与答案同前。}



\bimethod

  {倍二行差内减甲南行步，复以乘甲南行步为实，从空，倍二行差减甲南行步为廉，一步常法，得城径。}

  {两倍的两行差减去甲向南行走步数，再乘以甲向南行走步数作为被开方数，一次项为空，两倍的两行差减去甲向南行走步数作为二次项系数，三次项系数为1，开方求得城径。}



%% --- Problem 14 ---

\stepcounter{probnum}

\biproblem

  {或问乙出南门直行一百三十五步，甲从乾隅南行六百步望乙。问答同前。}

  {假设乙从南门出发向南直走135步，甲从乾隅（西北角）向南行走600步遥望乙。问题与答案同前。}



\bimethod

  {两行步相乘为实，从空，两行步为廉，一步常法，得城径。}

  {两行步数相乘作为被开方数，一次项为空，两行步数为二次项系数，三次项系数为1，开方求得城径。}



%% --- Problem 15 ---

\stepcounter{probnum}

\biproblem

  {或问乙出南门直行一百三十五步，甲从乾隅南行六百步，斜望乙与城参相直。问答同前。}

  {假设乙从南门出发向南直走135步，甲从乾隅（西北角）向南行走600步，斜向遥望乙，视线与城池边缘恰好在一条直线上。问题与答案同前。}



\bimethod

  {倍二行差内减甲南行步，复以乘甲南行步为实，从空，倍二行差减甲南行步为廉，一步常法，得城径。}

  {两倍的两行差减去甲向南行走步数，再乘以甲向南行走步数作为被开方数，一次项为空，两倍的两行差减去甲向南行走步数作为二次项系数，三次项系数为1，开方求得城径。}



%% --- Problem 16 ---

\stepcounter{probnum}

\biproblem

  {或问乙出南门直行一百三十五步，甲从乾隅南行六百步望乙。问答同前。}

  {假设乙从南门出发向南直走135步，甲从乾隅（西北角）向南行走600步遥望乙。问题与答案同前。}



\bimethod

  {两行步相乘为实，从空，两行步为廉，一步常法，得城径。}

  {两行步数相乘作为被开方数，一次项为空，两行步数为二次项系数，三次项系数为1，开方求得城径。}



%% --- Problem 17 ---

\stepcounter{probnum}

\biproblem

  {或问乙出南门直行一百三十五步，甲从乾隅南行六百步，斜望乙与城参相直。问答同前。}

  {假设乙从南门出发向南直走135步，甲从乾隅（西北角）向南行走600步，斜向遥望乙，视线与城池边缘恰好在一条直线上。问题与答案同前。}



\bimethod

  {倍二行差内减甲南行步，复以乘甲南行步为实，从空，倍二行差减甲南行步为廉，一步常法，得城径。}

  {两倍的两行差减去甲向南行走步数，再乘以甲向南行走步数作为被开方数，一次项为空，两倍的两行差减去甲向南行走步数作为二次项系数，三次项系数为1，开方求得城径。}



%% --- Problem 18 ---

\stepcounter{probnum}

\biproblem

  {或问乙出南门直行一百三十五步，甲从乾隅南行六百步望乙。问答同前。}

  {假设乙从南门出发向南直走135步，甲从乾隅（西北角）向南行走600步遥望乙。问题与答案同前。}



\bimethod

  {两行步相乘为实，从空，两行步为廉，一步常法，得城径。}

  {两行步数相乘作为被开方数，一次项为空，两行步数为二次项系数，三次项系数为1，开方求得城径。}



\vspace{0.5cm}

\longline

\begin{center}

\textit{测圆海镜卷五终}

\end{center}



\newpage



%% ===================================================================

%% VOLUME 6

%% ===================================================================

\section{测圆海镜卷六 —— 大勾一十八问}

\label{sec:vol6}



\bisectionheader

  {大勾一十八问}

  {十八道关于大勾的数学题}

  {卷六论「大勾」之术。所谓大勾，乃最大勾股形之勾边，即通勾也。以大勾为主，配以诸勾股形之关系，求圆城之径。凡一十八问。}

  {卷六讨论「大勾」的算法。所谓大勾，是指最大勾股三角形的勾边，即通勾。以大勾为主要对象，配合各种勾股三角形的关系，求圆形城池的直径。共18道题。}



\vspace{0.3cm}



%% Reset problem counter for Volume 6

\setcounter{probnum}{0}



%% --- Problem 1 ---

\stepcounter{probnum}

\biproblem

  {或问乙从东门直行一十六步，甲从乾隅东行三百二十步望乙，与城参相直。问答同前。}

  {假设乙从东门出发向东直走16步，甲从乾隅（西北角）向东行走320步遥望乙，视线与城池边缘恰好在一条直线上。问题与答案同前。}



\bian answer

  {城径二百四十步。}

  {城池直径为240步。}



\biworking

  {立天元一为半径，以二之加乙东行，得\raisebox{-0.5ex}{\scriptsize$\begin{array}{c}32\\2\end{array}$}为小差。半乙南行步，得一百八十步，以乘小差，得半径幂，寄左。然后以天元幂与左相消，开平方得一百二十步，即半径。}

  {设天元一为半径，用两倍的天元加上乙向东行走的步数，得到$\begin{pmatrix}32\\2\end{pmatrix}$为小差。取乙向南行走步数的一半，得180步，乘以小差，得到半径的平方，暂存左边。然后将天元的平方与左边相消，开平方得到120步，即半径。}



\bimethod

  {甲东行内减二之乙南行，复以乘甲东行为实。四之东行内减二之乙东行为从，四益隅，得半径。}

  {甲向东行走步数减去两倍乙向南行走步数，再乘以甲向东行走步数作为被开方数。四倍的东行步数减去两倍乙向东行走步数作为一次项系数，四次项采用益积法（增积法），开方求得半径。}



%% --- Problem 2 ---

\stepcounter{probnum}

\biproblem

  {或问乙从东门直行一十六步，甲从乾隅东行三百二十步望乙，与城参相直。问答同前。}

  {假设乙从东门出发向东直走16步，甲从乾隅（西北角）向东行走320步遥望乙，视线与城池边缘恰好在一条直线上。问题与答案同前。}



\bian answer

  {城径二百四十步。}

  {城池直径为240步。}



\biworking

  {立天元一为半径，以二之加乙东行，得\raisebox{-0.5ex}{\scriptsize$\begin{array}{c}16\\2\end{array}$}为小差。半甲东行步，得一百六十步，以乘小差，得\raisebox{-0.5ex}{\scriptsize$\begin{array}{c}2560\\320\end{array}$}为半径幂，寄左。然后以天元幂\raisebox{-0.5ex}{\scriptsize$\begin{array}{c}0\\1\end{array}$}与左相消，开平方得一百二十步，即半径。}

  {设天元一为半径，用两倍的天元加上乙向东行走的步数，得到$\begin{pmatrix}16\\2\end{pmatrix}$为小差。取甲向东行走步数的一半，得160步，乘以小差，得到$\begin{pmatrix}2560\\320\end{pmatrix}$作为半径的平方，暂存左边。然后将天元的平方$\begin{pmatrix}0\\1\end{pmatrix}$与左边相消，开平方得到120步，即半径。}



\bimethod

  {甲东行内减二之乙东行，复以乘甲东行为实。四之甲东行内减二之乙东行为从，四益隅，得半径。}

  {甲向东行走步数减去两倍乙向东行走步数，再乘以甲向东行走步数作为被开方数。四倍的甲向东行走步数减去两倍乙向东行走步数作为一次项系数，四次项采用益积法（增积法），开方求得半径。}



%% --- Problem 3 ---

\stepcounter{probnum}

\biproblem

  {或问乙从东门直行一十六步，甲从乾隅东行三百二十步望乙。问答同前。}

  {假设乙从东门出发向东直走16步，甲从乾隅（西北角）向东行走320步遥望乙。问题与答案同前。}



\bian answer

  {城径二百四十步。}

  {城池直径为240步。}



\biworking

  {立天元一为城径，以二之加乙东行，得\raisebox{-0.5ex}{\scriptsize$\begin{array}{c}16\\2\end{array}$}为小差。置甲东行幂内减小差幂，而半之，得大弦。带小差分母。又置甲东行幂内减小差幂，而半之，得大股。带小差分母。又以小差乘股，併入和，得下式，为小和，以乘大弦。寄左。又以倍天元加斜步，得小和，以乘大弦，得同数，与左相消。开立方得二百四十步，即城径。}

  {设天元一为城池直径，用两倍的天元加上乙向东行走的步数，得到$\begin{pmatrix}16\\2\end{pmatrix}$为小差。将甲向东行走步数的平方减去小差的平方，再取一半，得到大弦。带小差分母。又将甲向东行走步数的平方减去小差的平方，再取一半，得到大股。也带小差分母。又用小差乘以股，合并入和，得到下式，作为小和，乘以大弦。暂存左边。又用两倍的天元加上斜行步数，得到小和，乘以大弦，得到相同的数，与左边相消。开立方得到240步，即城池直径。}



\bimethod

  {甲东行内减二之乙东行，复以乘甲东行为实。四之甲东行内减二之乙东行为从，四益隅，得城径。}

  {甲向东行走步数减去两倍乙向东行走步数，再乘以甲向东行走步数作为被开方数。四倍的甲向东行走步数减去两倍乙向东行走步数作为一次项系数，四次项采用益积法（增积法），开方求得城径。}



\vspace{0.3cm}

\noindent

\begin{minipage}[t]{0.48\textwidth}

\textit{（以下各问草法类推，皆依大勾之术，以小差为主，建立方程。）}

\end{minipage}%

\hfill\vrule\hfill

\begin{minipage}[t]{0.48\textwidth}

\textit{（以下各题演算方法类似，都依照大勾的算法，以小差为主要对象，建立方程。）}

\end{minipage}



\vspace{0.3cm}



%% --- Problem 4 ---

\stepcounter{probnum}

\biproblem

  {或问乙从东门直行一十六步，甲从乾隅东行三百二十步望乙。问答同前。}

  {假设乙从东门出发向东直走16步，甲从乾隅（西北角）向东行走320步遥望乙。问题与答案同前。}



\bimethod

  {甲东行内减二之乙东行，复以乘甲东行为实，从空，四之甲东行内减二之乙东行为廉，四益隅，得城径。}

  {甲向东行走步数减去两倍乙向东行走步数，再乘以甲向东行走步数作为被开方数，一次项为空，四倍的甲向东行走步数减去两倍乙向东行走步数作为二次项系数，四次项采用益积法（增积法），开方求得城径。}



%% --- Problem 5 ---

\stepcounter{probnum}

\biproblem

  {或问乙从东门直行一十六步，甲从乾隅东行三百二十步望乙，与城参相直。问答同前。}

  {假设乙从东门出发向东直走16步，甲从乾隅（西北角）向东行走320步遥望乙，视线与城池边缘恰好在一条直线上。问题与答案同前。}



\bimethod

  {甲东行内减二之乙东行，复以乘甲东行为实，从空，四之甲东行内减二之乙东行为廉，四益隅，得城径。}

  {甲向东行走步数减去两倍乙向东行走步数，再乘以甲向东行走步数作为被开方数，一次项为空，四倍的甲向东行走步数减去两倍乙向东行走步数作为二次项系数，四次项采用益积法（增积法），开方求得城径。}



%% --- Problem 6 ---

\stepcounter{probnum}

\biproblem

  {或问乙从东门直行一十六步，甲从乾隅东行三百二十步望乙。问答同前。}

  {假设乙从东门出发向东直走16步，甲从乾隅（西北角）向东行走320步遥望乙。问题与答案同前。}



\bimethod

  {甲东行内减二之乙东行，复以乘甲东行为实，从空，四之甲东行内减二之乙东行为廉，四益隅，得城径。}

  {甲向东行走步数减去两倍乙向东行走步数，再乘以甲向东行走步数作为被开方数，一次项为空，四倍的甲向东行走步数减去两倍乙向东行走步数作为二次项系数，四次项采用益积法（增积法），开方求得城径。}



%% --- Problem 7 ---

\stepcounter{probnum}

\biproblem

  {或问乙从东门直行一十六步，甲从乾隅东行三百二十步望乙，与城参相直。问答同前。}

  {假设乙从东门出发向东直走16步，甲从乾隅（西北角）向东行走320步遥望乙，视线与城池边缘恰好在一条直线上。问题与答案同前。}



\bimethod

  {甲东行内减二之乙东行，复以乘甲东行为实，从空，四之甲东行内减二之乙东行为廉，四益隅，得城径。}

  {甲向东行走步数减去两倍乙向东行走步数，再乘以甲向东行走步数作为被开方数，一次项为空，四倍的甲向东行走步数减去两倍乙向东行走步数作为二次项系数，四次项采用益积法（增积法），开方求得城径。}



%% --- Problem 8 ---

\stepcounter{probnum}

\biproblem

  {或问乙从东门直行一十六步，甲从乾隅东行三百二十步望乙。问答同前。}

  {假设乙从东门出发向东直走16步，甲从乾隅（西北角）向东行走320步遥望乙。问题与答案同前。}



\bimethod

  {甲东行内减二之乙东行，复以乘甲东行为实，从空，四之甲东行内减二之乙东行为廉，四益隅，得城径。}

  {甲向东行走步数减去两倍乙向东行走步数，再乘以甲向东行走步数作为被开方数，一次项为空，四倍的甲向东行走步数减去两倍乙向东行走步数作为二次项系数，四次项采用益积法（增积法），开方求得城径。}



%% --- Problem 9 ---

\stepcounter{probnum}

\biproblem

  {或问乙从东门直行一十六步，甲从乾隅东行三百二十步望乙，与城参相直。问答同前。}

  {假设乙从东门出发向东直走16步，甲从乾隅（西北角）向东行走320步遥望乙，视线与城池边缘恰好在一条直线上。问题与答案同前。}



\bimethod

  {甲东行内减二之乙东行，复以乘甲东行为实，从空，四之甲东行内减二之乙东行为廉，四益隅，得城径。}

  {甲向东行走步数减去两倍乙向东行走步数，再乘以甲向东行走步数作为被开方数，一次项为空，四倍的甲向东行走步数减去两倍乙向东行走步数作为二次项系数，四次项采用益积法（增积法），开方求得城径。}



%% --- Problem 10 ---

\stepcounter{probnum}

\biproblem

  {或问乙从东门直行一十六步，甲从乾隅东行三百二十步望乙。问答同前。}

  {假设乙从东门出发向东直走16步，甲从乾隅（西北角）向东行走320步遥望乙。问题与答案同前。}



\bimethod

  {甲东行内减二之乙东行，复以乘甲东行为实，从空，四之甲东行内减二之乙东行为廉，四益隅，得城径。}

  {甲向东行走步数减去两倍乙向东行走步数，再乘以甲向东行走步数作为被开方数，一次项为空，四倍的甲向东行走步数减去两倍乙向东行走步数作为二次项系数，四次项采用益积法（增积法），开方求得城径。}



%% --- Problem 11 ---

\stepcounter{probnum}

\biproblem

  {或问乙从东门直行一十六步，甲从乾隅东行三百二十步望乙，与城参相直。问答同前。}

  {假设乙从东门出发向东直走16步，甲从乾隅（西北角）向东行走320步遥望乙，视线与城池边缘恰好在一条直线上。问题与答案同前。}



\bimethod

  {甲东行内减二之乙东行，复以乘甲东行为实，从空，四之甲东行内减二之乙东行为廉，四益隅，得城径。}

  {甲向东行走步数减去两倍乙向东行走步数，再乘以甲向东行走步数作为被开方数，一次项为空，四倍的甲向东行走步数减去两倍乙向东行走步数作为二次项系数，四次项采用益积法（增积法），开方求得城径。}



%% --- Problem 12 ---

\stepcounter{probnum}

\biproblem

  {或问乙从东门直行一十六步，甲从乾隅东行三百二十步望乙。问答同前。}

  {假设乙从东门出发向东直走16步，甲从乾隅（西北角）向东行走320步遥望乙。问题与答案同前。}



\bimethod

  {甲东行内减二之乙东行，复以乘甲东行为实，从空，四之甲东行内减二之乙东行为廉，四益隅，得城径。}

  {甲向东行走步数减去两倍乙向东行走步数，再乘以甲向东行走步数作为被开方数，一次项为空，四倍的甲向东行走步数减去两倍乙向东行走步数作为二次项系数，四次项采用益积法（增积法），开方求得城径。}



%% --- Problem 13 ---

\stepcounter{probnum}

\biproblem

  {或问乙从东门直行一十六步，甲从乾隅东行三百二十步望乙，与城参相直。问答同前。}

  {假设乙从东门出发向东直走16步，甲从乾隅（西北角）向东行走320步遥望乙，视线与城池边缘恰好在一条直线上。问题与答案同前。}



\bimethod

  {甲东行内减二之乙东行，复以乘甲东行为实，从空，四之甲东行内减二之乙东行为廉，四益隅，得城径。}

  {甲向东行走步数减去两倍乙向东行走步数，再乘以甲向东行走步数作为被开方数，一次项为空，四倍的甲向东行走步数减去两倍乙向东行走步数作为二次项系数，四次项采用益积法（增积法），开方求得城径。}



%% --- Problem 14 ---

\stepcounter{probnum}

\biproblem

  {或问乙从东门直行一十六步，甲从乾隅东行三百二十步望乙。问答同前。}

  {假设乙从东门出发向东直走16步，甲从乾隅（西北角）向东行走320步遥望乙。问题与答案同前。}



\bimethod

  {甲东行内减二之乙东行，复以乘甲东行为实，从空，四之甲东行内减二之乙东行为廉，四益隅，得城径。}

  {甲向东行走步数减去两倍乙向东行走步数，再乘以甲向东行走步数作为被开方数，一次项为空，四倍的甲向东行走步数减去两倍乙向东行走步数作为二次项系数，四次项采用益积法（增积法），开方求得城径。}



%% --- Problem 15 ---

\stepcounter{probnum}

\biproblem

  {或问乙从东门直行一十六步，甲从乾隅东行三百二十步望乙，与城参相直。问答同前。}

  {假设乙从东门出发向东直走16步，甲从乾隅（西北角）向东行走320步遥望乙，视线与城池边缘恰好在一条直线上。问题与答案同前。}



\bimethod

  {甲东行内减二之乙东行，复以乘甲东行为实，从空，四之甲东行内减二之乙东行为廉，四益隅，得城径。}

  {甲向东行走步数减去两倍乙向东行走步数，再乘以甲向东行走步数作为被开方数，一次项为空，四倍的甲向东行走步数减去两倍乙向东行走步数作为二次项系数，四次项采用益积法（增积法），开方求得城径。}



%% --- Problem 16 ---

\stepcounter{probnum}

\biproblem

  {或问乙从东门直行一十六步，甲从乾隅东行三百二十步望乙。问答同前。}

  {假设乙从东门出发向东直走16步，甲从乾隅（西北角）向东行走320步遥望乙。问题与答案同前。}



\bimethod

  {甲东行内减二之乙东行，复以乘甲东行为实，从空，四之甲东行内减二之乙东行为廉，四益隅，得城径。}

  {甲向东行走步数减去两倍乙向东行走步数，再乘以甲向东行走步数作为被开方数，一次项为空，四倍的甲向东行走步数减去两倍乙向东行走步数作为二次项系数，四次项采用益积法（增积法），开方求得城径。}



%% --- Problem 17 ---

\stepcounter{probnum}

\biproblem

  {或问乙从东门直行一十六步，甲从乾隅东行三百二十步望乙，与城参相直。问答同前。}

  {假设乙从东门出发向东直走16步，甲从乾隅（西北角）向东行走320步遥望乙，视线与城池边缘恰好在一条直线上。问题与答案同前。}



\bimethod

  {甲东行内减二之乙东行，复以乘甲东行为实，从空，四之甲东行内减二之乙东行为廉，四益隅，得城径。}

  {甲向东行走步数减去两倍乙向东行走步数，再乘以甲向东行走步数作为被开方数，一次项为空，四倍的甲向东行走步数减去两倍乙向东行走步数作为二次项系数，四次项采用益积法（增积法），开方求得城径。}



%% --- Problem 18 ---

\stepcounter{probnum}

\biproblem

  {或问乙从东门直行一十六步，甲从乾隅东行三百二十步望乙，与城参相直。问答同前。}

  {假设乙从东门出发向东直走16步，甲从乾隅（西北角）向东行走320步遥望乙，视线与城池边缘恰好在一条直线上。问题与答案同前。}



\bimethod

  {甲东行内减二之乙东行，复以乘甲东行为实，从空，四之甲东行内减二之乙东行为廉，四益隅，得城径。}

  {甲向东行走步数减去两倍乙向东行走步数，再乘以甲向东行走步数作为被开方数，一次项为空，四倍的甲向东行走步数减去两倍乙向东行走步数作为二次项系数，四次项采用益积法（增积法），开方求得城径。}



\vspace{0.5cm}

\longline

\begin{center}

\textit{测圆海镜卷六终}

\end{center}



\newpage



%% ===================================================================

%% APPENDIX: STRUCTURE OF THE WORK

%% ===================================================================

\section*{附录：全书结构}

\addcontentsline{toc}{section}{附录：全书结构}



\noindent

\begin{minipage}[t]{0.48\textwidth}

\textbf{【原文说明】}



《测圆海镜》全书十二卷，共一百七十问：



\begin{center}

\begin{tabular}{clc}

\toprule

\textbf{卷} & \textbf{篇名} & \textbf{问数} \\

\midrule

卷一 & 圆城图式 & --- \\

卷二 & 正率一十四问 & 14 \\

卷三 & 边股一十七问 & 17 \\

\rowcolor{gray!15}

卷四 & 底勾一十七问 & 17 \\

\rowcolor{gray!15}

卷五 & 大股一十八问 & 18 \\

\rowcolor{gray!15}

卷六 & 大勾一十八问 & 18 \\

卷七 & 明勾一十六问 & 16 \\

卷八 & 明股一十六问 & 16 \\

卷九 & 杂题 & 18 \\

卷十 & 杂题 & 18 \\

卷十一 & 杂题 & 18 \\

卷十二 & 杂题 & 10 \\

\midrule

 & \textbf{合计} & \textbf{170} \\

\bottomrule

\end{tabular}

\end{center}



\vspace{0.5cm}

\textbf{术语解释：}

\begin{itemize}[leftmargin=1em]

  \item \textbf{天元一}：设未知数为一，即现代代数之$x$

  \item \textbf{寄左}：将某式暂存一边，相当于方程之一边

  \item \textbf{相消}：两式相减，相当于移项合并

  \item \textbf{带分母}：某式含有分母未除

  \item \textbf{幂}：平方，即现代记号$x^2$

  \item \textbf{开方}：开平方，即求平方根

  \item \textbf{开立方}：开立方，即求立方根

  \item \textbf{实}：常数项

  \item \textbf{从}：一次项系数

  \item \textbf{廉}：二次项系数

  \item \textbf{隅}：三次项系数（立方方程中）

\end{itemize}

\end{minipage}%

\hfill\vrule\hfill

\begin{minipage}[t]{0.48\textwidth}

\textbf{【今译说明】}



《测圆海镜》全书共12卷，总计170道题：



\begin{center}

\begin{tabular}{clc}

\toprule

\textbf{卷} & \textbf{篇名} & \textbf{题数} \\

\midrule

卷一 & 圆城图式 & --- \\

卷二 & 正率14题 & 14 \\

卷三 & 边股17题 & 17 \\

\rowcolor{gray!15}

卷四 & 底勾17题 & 17 \\

\rowcolor{gray!15}

卷五 & 大股18题 & 18 \\

\rowcolor{gray!15}

卷六 & 大勾18题 & 18 \\

卷七 & 明勾16题 & 16 \\

卷八 & 明股16题 & 16 \\

卷九 & 杂题 & 18 \\

卷十 & 杂题 & 18 \\

卷十一 & 杂题 & 18 \\

卷十二 & 杂题 & 10 \\

\midrule

 & \textbf{合计} & \textbf{170} \\

\bottomrule

\end{tabular}

\end{center}



\vspace{0.5cm}

\textbf{术语对照：}

\begin{itemize}[leftmargin=1em]

  \item \textbf{天元一}：设未知数为一，即现代代数的$x$

  \item \textbf{寄左}：将某算式暂存一边，相当于方程的一边

  \item \textbf{相消}：两式相减，相当于移项合并同类项

  \item \textbf{带分母}：某算式含有分母未除尽

  \item \textbf{幂}：平方，即现代记号$x^2$

  \item \textbf{开方}：开平方，即求平方根

  \item \textbf{开立方}：开立方，即求立方根

  \item \textbf{实}：常数项

  \item \textbf{从}：一次项系数

  \item \textbf{廉}：二次项系数

  \item \textbf{隅}：三次项系数（立方方程中）

\end{itemize}

\end{minipage}



\vspace{1cm}



%% ===================================================================

%% GLOSSARY OF TECHNICAL TERMS

%% ===================================================================

\section*{术语文汇对照}

\addcontentsline{toc}{section}{术语文汇对照}



\noindent

\begin{minipage}[t]{0.48\textwidth}

\textbf{【古代数学术语】}



\begin{tabular}{ll}

\toprule

\textbf{术语} & \textbf{释义} \\

\midrule

天元术 & 设未知数列方程之法 \\

天元一 & 未知数 $x$ \\

勾股 & 直角三角形之两边 \\

勾 & 短直角边 \\

股 & 长直角边 \\

弦 & 斜边 \\

幂 & 平方 $x^2$ \\

实 & 常数项 \\

从 & 一次项系数 \\

廉 & 二次项系数 \\

隅 & 三次项系数 \\

带从开方 & 带参数之方程求解 \\

益积 & 增积法 \\

减从 & 减项法 \\

重差 & 双重差分测量法 \\

三乘方 & 三次方程 \\

四乘方 & 四次方程 \\

寄左 & 暂存一边 \\

相消 & 两式相减 \\

带分母 & 含分母未除 \\

\bottomrule

\end{tabular}

\end{minipage}%

\hfill\vrule\hfill

\begin{minipage}[t]{0.48\textwidth}

\textbf{【现代数学术语】}



\begin{tabular}{ll}

\toprule

\textbf{术语} & \textbf{现代解释} \\

\midrule

天元术 & 代数方法，设未知数列方程 \\

天元一 & 未知数 $x$ \\

勾股 & 直角三角形的两条直角边 \\

勾 & 短直角边（现代记为 $a$） \\

股 & 长直角边（现代记为 $b$） \\

弦 & 斜边（现代记为 $c$） \\

幂 & 平方运算 $x^2$ \\

实 & 方程的常数项 \\

从 & 一次项系数（现代记为 $a_1$） \\

廉 & 二次项系数（现代记为 $a_2$） \\

隅 & 三次项系数（现代记为 $a_3$） \\

带从开方 & 带参数的方程求解（如 $x^2+px=q$） \\

益积 & 增积法，增加被开方数 \\

减从 & 减项法，减少中间项 \\

重差 & 双重差分测量法（三角测量原理） \\

三乘方 & 三次方程（立方方程） \\

四乘方 & 四次方程（双二次方程） \\

寄左 & 暂存方程的一边 \\

相消 & 方程两边相减（移项合并） \\

带分母 & 算式含有未约去的分母 \\

\bottomrule

\end{tabular}

\end{minipage}



\vspace{1cm}

\begin{center}

\rule{0.3\textwidth}{0.4pt}



\textit{--- 全书完 ---}



\vspace{0.5cm}



{\small 文言文原文依《钦定四库全书》本 \\[0.2cm]

白话文翻译与注释为现代译解}

\end{center}



\end{document}







% ============================================================

% Source: translations/zh/chinese/li-ye-ceyuan-haijing-vols7-9_classical-modern_bilingual.tex

% ============================================================



% ========================================================================

% 測圓海鏡 (Ceyuan Haijing / Sea Mirror of Circle Measurements)

% by Li Ye (李冶), 1248

% Volumes 7--9

% BILINGUAL EDITION: Classical Chinese ←→ Modern Chinese

% 文言文 ↔ 白话文 对照版

% ========================================================================

\documentclass[11pt,a4paper]{article}



% -------- Core Packages --------

\usepackage{fontspec}

\usepackage{amsmath,amssymb,amsthm}

% Make \square work in text mode

\AtBeginDocument{%

  \let\origsquare\square

  \renewcommand{\square}{\ensuremath{\origsquare}}%

}

\usepackage{geometry}

\usepackage{graphicx}

\usepackage{xcolor}

\usepackage{titlesec}

\usepackage{fancyhdr}

\usepackage{enumitem}

\usepackage{array,booktabs,longtable,multirow}

\usepackage{hyperref}

\usepackage{microtype}

\usepackage{tocloft}

\usepackage{indentfirst}

\usepackage{xeCJK}

\usepackage{etoolbox}

\usepackage{paracol}



% -------- Page Geometry --------

\geometry{

  paperwidth=210mm,paperheight=297mm,

  left=18mm,right=18mm,top=25mm,bottom=25mm,

  headheight=14pt,footskip=30pt

}



% -------- Font Configuration --------

\setmainfont{Times New Roman}



% Chinese Fonts via xeCJK - Classical uses Microsoft YaHei, Modern uses Arial CJK

\setCJKmainfont{Microsoft YaHei}

\newCJKfontfamily\tradfont{Microsoft YaHei}

\newCJKfontfamily\tradfontbf{Microsoft YaHei Bold}

\newCJKfontfamily\modfont{Microsoft YaHei}

\newCJKfontfamily\modfontbf{Microsoft YaHei}



% -------- Colors --------

\definecolor{titlecolor}{RGB}{45,55,72}

\definecolor{sectioncolor}{RGB}{55,65,81}

\definecolor{notecolor}{RGB}{120,53,15}

\definecolor{rulecolor}{RGB}{180,160,140}

\definecolor{problemcolor}{RGB}{80,40,20}

\definecolor{methodcolor}{RGB}{20,60,80}

\definecolor{workcolor}{RGB}{40,70,40}

\definecolor{classicalbg}{RGB}{255,250,245}

\definecolor{modernbg}{RGB}{245,250,255}

\definecolor{classicalborder}{RGB}{180,140,120}

\definecolor{modernborder}{RGB}{120,140,180}



% -------- Column Ratio --------

\columnsep=16pt



% -------- Section Formatting --------

\titleformat{\section}

  {\normalfont\Large\bfseries\color{sectioncolor}\tradfontbf}

  {\thesection}{1em}{}

\titleformat{\subsection}

  {\normalfont\large\bfseries\color{sectioncolor}\tradfontbf}

  {\thesubsection}{1em}{}

\titleformat{\subsubsection}

  {\normalfont\normalsize\bfseries\color{sectioncolor}\tradfontbf}

  {\thesubsubsection}{1em}{}



% -------- Header/Footer --------

\pagestyle{fancy}

\fancyhf{}

\fancyhead[L]{\small\tradfont 測圓海鏡\quad 卷七至卷九}

\fancyhead[R]{\small\modfont 文言文\,$\boldsymbol{\leftrightarrow}$\,白话文对照版}

\fancyfoot[C]{\small\thepage}

\renewcommand{\headrulewidth}{0.4pt}

\renewcommand{\footrulewidth}{0pt}



% -------- Custom Commands --------

\newcommand{\longline}{\noindent\rule{\textwidth}{0.4pt}}

\newcommand{\shortline}{\noindent\rule{0.5\textwidth}{0.4pt}\par}



% Classical (left column) problem environment

\newcommand{\clOWEN}[1]{%

  \par\vspace{0.3em}%

  \noindent\textcolor{problemcolor}{\textbf{【或問】}}\,{\tradfont #1}\par\nopagebreak%

}

% Modern (right column) problem environment

\newcommand{\mdOWEN}[1]{%

  \par\vspace{0.3em}%

  \noindent\textcolor{problemcolor}{\textbf{【问题】}}\,{\modfont #1}\par\nopagebreak%

}



% Classical method environment

\newcommand{\clFAYUE}[1]{%

  \par\noindent\textcolor{methodcolor}{\textbf{【法曰】}}\,{\tradfont #1}\par\nopagebreak%

}

% Modern method environment

\newcommand{\mdFAYUE}[1]{%

  \par\noindent\textcolor{methodcolor}{\textbf{【解法】}}\,{\modfont #1}\par\nopagebreak%

}



% Classical working environment

\newcommand{\clTSAOYUE}[1]{%

  \par\noindent\textcolor{workcolor}{\textbf{【草曰】}}\,{\tradfont #1}\par%

}

% Modern working environment

\newcommand{\mdTSAOYUE}[1]{%

  \par\noindent\textcolor{workcolor}{\textbf【演算】}}\,{\modfont #1}\par%

}



% Inline Chinese helpers

\newcommand{\zh}[1]{{\tradfont #1}}

\newcommand{\md}[1]{{\modfont #1}}



% Bilingual pair command: classical + modern

\newcommand{\bilingual}[2]{%

  \begin{paracol}{2}%

  \begin{leftcolumn}%

    \fcolorbox{classicalborder}{classicalbg}{\parbox{\linewidth}{#1}}%

  \end{leftcolumn}%

  \switchcolumn%

  \begin{rightcolumn}%

    \fcolorbox{modernborder}{modernbg}{\parbox{\linewidth}{#2}}%

  \end{rightcolumn}%

  \end{paracol}%

}



% Simple parallel content

\newcommand{\classical}[1]{{\tradfont #1}}

\newcommand{\modern}[1]{{\modfont #1}}



% Problem number counter

\newcounter{problemnum}

\setcounter{problemnum}{104}

\newcommand{\nextproblem}{\stepcounter{problemnum}\textbf{第\chinese{problemnum}問}}



% -------- Hyperref Setup --------

\hypersetup{

  colorlinks=true,

  linkcolor=sectioncolor,

  citecolor=sectioncolor,

  urlcolor=sectioncolor,

  pdfencoding=auto,

  pdftitle={測圓海鏡 卷七至卷九 文言文白话文对照版},

  pdfauthor={李冶},

  pdfsubject={Chinese Mathematics, 1248, Bilingual Edition}

}



% -------- TOC Depth --------

\setcounter{tocdepth}{2}

\setcounter{secnumdepth}{2}



% -------- CJK Tweaks --------

\xeCJKsetup{PunctStyle=banjiao,CJKmath=true}

\setlength{\parindent}{2em}

\setlength{\parskip}{0.3em}



% ========================================================================

\begin{document}



% ===================================================================

% TITLE PAGE

% ===================================================================

\begin{titlepage}

\centering

\vspace*{1.5cm}



{\Huge\bfseries\color{titlecolor}\tradfont 測圓海鏡}



\vspace{0.3cm}

{\Large\color{sectioncolor}\tradfont （卷七至卷九）}



\vspace{0.8cm}

{\LARGE\tradfont 李冶\quad 撰}



\vspace{0.2cm}

{\large 1248}



\vspace{1.5cm}



\fcolorbox{classicalborder}{classicalbg}{%

  \parbox{0.85\textwidth}{%

    \centering\vspace{0.3cm}

    {\Large\tradfontbf 文言文原文}\quad{\large$\longleftrightarrow$}\quad{\Large\modfontbf 白话文译文}

    \vspace{0.3cm}

  }%

}



\vspace{1cm}



{\small\tradfont

卷七：明前一十八問（第105--122問）\quad 高次方程（四次、五次）\\[0.3em]

卷八：明後一十六問（第123--138問）\quad 多元天元術\\[0.3em]

卷九：大斜四問、大和八問（第139--170問）\quad 複雜幾何構形

}



\vspace{1.5cm}



\fcolorbox{modernborder}{modernbg}{%

  \parbox{0.85\textwidth}{%

    \centering\vspace{0.3cm}

    {\small\modfont

    \textbf{现代数学术语对照}\\[0.3em]

    五乘方 $\rightarrow$ 五次方程\quad

    六乘方 $\rightarrow$ 六次方程\quad

    連枝同体 $\rightarrow$ 連枝同体（关联方程组）\\[0.2em]

    玲珑 $\rightarrow$ 玲珑（精巧的方程变换）\quad

    換骨 $\rightarrow$ 換骨（根本变换）\quad

    奪項 $\rightarrow$ 奪項（消元法）

    }

    \vspace{0.3cm}

  }%

}



\vfill



{\small 测圆海镜 卷七至卷九\quad 文言文/白话文对照版\\

Bilingual Edition: Classical Chinese / Modern Chinese}



\end{titlepage}



% ===================================================================

% TABLE OF CONTENTS

% ===================================================================

\tableofcontents

\newpage



% ===================================================================

% BILINGUAL INTRODUCTION

% ===================================================================

\section*{引言\quad Introduction}

\addcontentsline{toc}{section}{引言 Introduction}



\begin{paracol}{2}



\begin{leftcolumn}

\fcolorbox{classicalborder}{classicalbg}{%

\parbox{\linewidth}{%

\textbf{\large\tradfont 【文言文原文】}



\vspace{0.3em}

《測圓海鏡》乃金元數學家李冶（1192--1279）所撰，成書於1248年，為中古代數學之傑作。

全書凡十二卷，計一百七十問，皆論圓城與勾股形之關係，而以「天元術」立一算求解多項式方程。



卷七至卷九所載如下：

\begin{itemize}[leftmargin=1.5em,topsep=0pt]

  \item \textbf{卷七}：明前一十八問

  \item \textbf{卷八}：明後一十六問

  \item \textbf{卷九}：大斜四問、大和八問

\end{itemize}



每問之體例分為三段：

\begin{itemize}[leftmargin=1.5em,topsep=0pt]

  \item \textbf{或問}：設問之辭

  \item \textbf{法曰}：演算之公式

  \item \textbf{草曰}：詳細之天元推演

\end{itemize}



書中「■」者，表示算籌佈列之太極、天元位，即古代多項式之記號也。

}}

\end{leftcolumn}



\switchcolumn



\begin{rightcolumn}

\fcolorbox{modernborder}{modernbg}{%

\parbox{\linewidth}{%

\textbf{\large\modfont 【白话文译文】}



\vspace{0.3em}

《测圆海镜》由金元时期数学家李冶（1192--1279）所著，成书于1248年，是中世纪代数学的杰作。

全书共十二卷，计170题，均讨论圆城与直角三角形的关系，以「天元术」设未知数求解多项式方程。



卷七至卷九内容如下：

\begin{itemize}[leftmargin=1.5em,topsep=0pt]

  \item \textbf{卷七}：明前十八问（第105--122题）--- 高次多项式方程（四次、五次）

  \item \textbf{卷八}：明后十六问（第123--138题）--- 多元天元术

  \item \textbf{卷九}：大斜四问、大和八问（第139--170题）--- 复杂几何构形

\end{itemize}



每题的格式分为三段：

\begin{itemize}[leftmargin=1.5em,topsep=0pt]

  \item \textbf{【问题】}：问题陈述

  \item \textbf{【解法】}：计算公式

  \item \textbf{【演算】}：详细的天元术代数推导

\end{itemize}



书中的「■」符号表示算筹排列的太极、天元位置，即古代多项式记号的占位符。

}}

\end{rightcolumn}



\end{paracol}



\newpage



% ===================================================================

% VOLUME 7

% ===================================================================

\section{\tradfont 卷七\quad 明前一十八問}

\label{sec:vol7}



\begin{center}

\fcolorbox{classicalborder}{classicalbg}{\parbox{0.9\textwidth}{\centering

\textbf{\large 卷七概要}\quad 本卷十八问，多为高次方程求解。以天元一设圆径或半径，通过几何关系建立四次（三乘方）、五次（四乘方）及六次（五乘方）方程，运用增乘开方法求解。

}}

\end{center}



\vspace{0.5cm}



% -------------------------------------------------------------------

% Problem 105

% -------------------------------------------------------------------

\begin{paracol}{2}



\begin{leftcolumn}

\clOWEN{出南門東行七十二步有樹，出東門南行三十步見之。問答同前。}



\clFAYUE{倍南行以乘倍東行為平實，並二行又倍之為從，一虛隅。得城徑。}



\clTSAOYUE{%

識別得此問名為弦外容圓，又為內率求虛積。其二行步相並為虛弦，若以相減即虛較也。

又倍東行為弦較和，倍南行即弦較較，此二數相乘則兩虛積也。若直以二行相乘，則半個虛積也。

又倍東行減於城徑，餘即二虛勾也。倍南行減於城徑則二虛股也。虛積上三事和即城徑也。

乃立天元一為圓徑，便以為三事和也。倍二行步減之，得$\square$為黃方一，天元乘之得$\square$為二虛積（寄左）。

然後倍東行以乘倍南行，得八千六百四十為同數，與左相消得$\square$。

益積開平方得二百四十步，即城徑也。合問。}

\end{leftcolumn}



\switchcolumn



\begin{rightcolumn}

\mdOWEN{从南门向东走72步有一棵树，从东门向南走30步可以看见它。求城墙直径。}



\mdFAYUE{将南行步数翻倍后乘以东行步数的两倍作为常数项，两步行数相加后再翻倍作为一次项系数，负一为二次项系数。解得城墙直径。}



\mdTSAOYUE{%

辨识得此题名为「弦外容圆」，又为「内率求虚积」。两步行数相加为虚弦，相减则为虚较。

东行步数的两倍为弦较和，南行步数的两倍为弦较较，二者相乘等于两倍虚积。直接以二行相乘则为半个虚积。

以城径减去两倍东行，余数为两倍虚勾；以城径减去两倍南行，余数为两倍虚股。虚积上的三事和即为城径。

设天元一为圆径，即三事和。减去两倍二行步数得$\square$为黄方，乘以天元得$\square$为两倍虚积（寄左）。

然后两倍东行乘以两倍南行得8640为同数，与左式相消得$\square$。

用增积开平方法得240步，即城径。符合所问。}

\end{rightcolumn}



\end{paracol}



\vspace{0.3cm}

\begin{center}\longline\end{center}



% -------------------------------------------------------------------

% Problem 106

% -------------------------------------------------------------------

\begin{paracol}{2}



\begin{leftcolumn}

\clOWEN{丙出南門直行一百三十五步而立，甲出東門直行一十六步見之。問答同前。}



\clFAYUE{%

以丙行步一百三十五再自之，得二百四十六萬零三百七十五於上。

又以甲行一十六乘丙行冪一萬八千二百二十五，得二十九萬一千六百，以乘上位，

得七千一百七十四億四千五百三十五萬為三乘方實；

以二行步相乘又倍之得四千三百二十，以乘丙行步再自之數，得一百六億二千八百八十二萬為益從；

第一廉空；以甲行乘丙行冪，得二十九萬一千六百，又倍之得五十八萬三千二百於上，

四之甲行冪一千零二十四，以乘丙行步，得一十三萬八千二百四十，減上位餘四十四萬四千九百六十為第二廉；

二行步相乘得二千一百六十為虛常法。得丙行步上勾弦差八十一。}



\clTSAOYUE{%

識別二數相並，得一百五十一。以減於皇極弦，餘一百三十八，即虛勾虛股並也。

若以二數相減，餘一百一十九為高弦內減平弦，又為皇極弦內少個小差弦，

又為大差弦內減個皇極弦也。立天元一為丙行大差數。

置丙行步一百三十五，自乘得$\square$，用天元除之，得$\square$為勾弦並也。

上減天元得$\square$為二丙勾也。復用丙南行乘之，得$\square$為二積也。

又以天元除之，得$\square$為丙勾外容圓半（泛寄）。

別置丙南行用二甲勾乘之，得$\square$太，合用二丙勾除之。不受除，

便以此為甲股（內寄二丙勾為分母）。復用二甲勾三十二乘之，得$\square$太為二個甲直積也。

又置丙南行內減天元，得$\square$為黃方，以自乘得$\square$為丙上勾弦差乘股弦差二段，

以天元除之得$\square$為兩個丙小差也。乃用甲股乘之，得下式$\square$。

復用丙南行除之，得$\square$，又折半得下式$\square$為一個甲步股弦差也，

內亦帶前二丙勾分母。復置兩個甲直積，內已寄此甲股弦差分母，便為甲步股外容圓半（寄左）。

乃再置先求到泛寄，用甲股弦差分母乘之，得$\square$為同數，

與左相消得下式$\square$。開三乘方得八十一步，即丙步上勾弦差也。

《鈐經》載此法，以勾弦差率冪減丙行冪，復以丙行乘之為實，以差率冪為法，如法得徑。

此法只是以勾外求圓半。合以大差除倍積，而今皆以大差冪為分母也。

依法求之，勾弦差八十一自之得六千五百六十一，以減於丙行冪一萬八千二百二十五，

餘一萬一千六百六十四，復以丙行一百三十五乘之，得一百五十七萬四千六百四十為實，

以大差冪六千五百六十一為法，如法得二百四十步，即城徑也。}

\end{leftcolumn}



\switchcolumn



\begin{rightcolumn}

\mdOWEN{丙出南门直行135步停下，甲出东门直行16步便能看见丙。求城墙直径。}



\mdFAYUE{%

将丙行步135自乘两次（立方）得2,460,375为上位。

再以甲行16乘丙行平方18,225得291,600，乘以上位，

得71,744,535,000,000为四次方程（三乘方）常数项；

二行步相乘再翻倍得4,320，乘丙行步的立方数，得106,288,200,000为负一次项；

第一廉系数为零；以甲行乘丙行平方得291,600，翻倍得583,200为上位，

甲行平方1,024的4倍乘丙行步得138,240，减上位余444,960为第二廉系数；

二行步相乘得2,160为常数项。解得丙行步上的勾弦差为81。}



\mdTSAOYUE{%

辨识两数相加得151。从皇极弦中减去，余138即虚勾与虚股之和。

若两数相减，余119为高弦减平弦，又为皇极弦减去小差弦，

又为大差弦减去皇极弦。设天元一为丙行的大差数。

取丙行步135，自乘得$\square$，除以天元，得$\square$为勾弦和。

减去天元得$\square$为两倍丙勾。再用丙南行乘之，得$\square$为两倍积。

再除以天元，得$\square$为丙勾外容圆半（泛寄）。

另取丙南行用两倍甲勾乘之，得$\square$太，应除以两倍丙勾。不除，

便以此作为甲股（内寄两倍丙勾为分母）。再用两倍甲勾32乘之，得$\square$太为两个甲直积。

再取丙南行减去天元，得$\square$为黄方，自乘得$\square$为丙上勾弦差乘股弦差两段，

除以天元得$\square$为两个丙小差。用甲股乘之，得下式$\square$。

再用丙南行除之，得$\square$，折半得下式$\square$为一个甲步股弦差，

内亦带前两倍丙勾分母。再置两个甲直积，内已寄此甲股弦差分母，便为甲步股外容圆半（寄左）。

再置先求到的泛寄，用甲股弦差分母乘之，得$\square$为同数，

与左式相消得下式$\square$。开四次方（三乘方）得81步，即丙步上勾弦差。

《钤经》载此法：以勾弦差率平方减丙行平方，再以丙行乘之为被除数，以差率平方为除数，依法得直径。

此法是以勾外求圆半。应以大差除倍积，而今皆以大差平方为分母。

依法求之：勾弦差81自乘得6,561，减丙行平方18,225，

余11,664，再以丙行135乘之，得157,464,000为实，

以大差平方6,561为法，依法得240步，即城径。}

\end{rightcolumn}



\end{paracol}



\vspace{0.3cm}

\begin{center}\longline\end{center}





% -------------------------------------------------------------------

% Problem 107

% -------------------------------------------------------------------

\begin{paracol}{2}



\begin{leftcolumn}

\clOWEN{出東門一十六步有樹，出南門東行七十二步見之。問答同前。}



\clFAYUE{二行步相減得數，以自之於上。又以出東門步自之，減上位為平方實，

二之出南門東行步為益從，一步常法。翻開得半徑。}



\clTSAOYUE{%

別得人到樹即平弦也，半圓徑即平股也。其東行七十二步則平勾平弦差也。

乃立天元一為半圓徑，加一十六減七十二，得$\square$為勾也。

以自之得$\square$為勾冪。又加入天元股冪得$\square$為弦冪（寄左）。

再立天元一為半徑，加出東門步，得$\square$即弦也。

以自之得$\square$為同數，與左相消得$\square$。

翻法開之得一百二十步，即半城徑也。合問。}

\end{leftcolumn}



\switchcolumn



\begin{rightcolumn}

\mdOWEN{出东门16步有树，出南门向东走72步看见它。求城墙直径。}



\mdFAYUE{两步行数相减，差数自乘为上位。再以东门步数自乘，减上位为二次方程常数项，

南门东行步数的两倍为负一次项系数，1为二次项系数。翻法开方得半径。}



\mdTSAOYUE{%

辨识得：人到达树处即为平弦，半圆径即为平股。东行72步为平勾与平弦之差。

设天元一为半圆径，加上16减去72，得$\square$为勾。

自乘得$\square$为勾冪。加上天元的平方（股冪）得$\square$为弦冪（寄左）。

再设天元一为半径，加出东门步数，得$\square$即弦。

自乘得$\square$为同数，与左式相消得$\square$。

翻法开方得120步，即半城径。符合所问。}

\end{rightcolumn}



\end{paracol}



\vspace{0.3cm}

\begin{center}\longline\end{center}



% -------------------------------------------------------------------

% Problem 108

% -------------------------------------------------------------------

\begin{paracol}{2}



\begin{leftcolumn}

\clOWEN{出南門一百三十五步有樹，出東門南行三十步見之。問答同前。}



\clFAYUE{樹去城步內減南行步，餘以為冪於上，又以樹去城步為冪內減上位為平實，

倍樹去城步為從，一虛隅。翻法得半城徑。}



\clTSAOYUE{%

別得人距樹即高弦也，半圓徑即高勾也，其南行三十步即高弦上小差也。

乃立天元一為半徑，加樹去城步為弦，內減小差$\square$，得$\square$即股也。

以自之得$\square$為股冪，內加入天元冪，得$\square$為弦冪（寄左）。

再置弦$\square$以自之，得$\square$為同數，與左相消，得式$\square$。

翻開得一百二十步，即半城徑也。合問。}

\end{leftcolumn}



\switchcolumn



\begin{rightcolumn}

\mdOWEN{出南门135步有树，出东门向南走30步看见它。求城墙直径。}



\mdFAYUE{树离城的步数减去南行步数，余数自乘为上位，再以树离城步数自乘减上位为二次方程常数项，

树离城步数的两倍为一次项系数，负一为二次项系数。翻法得半城径。}



\mdTSAOYUE{%

辨识得：人到树的距离即为高弦，半圆径即为高勾，南行30步为高弦上的小差。

设天元一为半径，加上树离城步数为弦，减去小差$\square$，得$\square$即股。

自乘得$\square$为股冪，加上天元冪，得$\square$为弦冪（寄左）。

再置弦$\square$自乘，得$\square$为同数，与左式相消得式$\square$。

翻法开方得120步，即半城径。符合所问。}

\end{rightcolumn}



\end{paracol}



\vspace{0.3cm}

\begin{center}\longline\end{center}



% -------------------------------------------------------------------

% Problem 109

% -------------------------------------------------------------------

\begin{paracol}{2}



\begin{leftcolumn}

\clOWEN{乙出東門不知遠近而立，甲出南門東行七十二步望見乙，就乙斜行一百三十六步與乙相會。問答同前。}



\clFAYUE{以斜行步自之於上，以二行相減餘自為冪，減上位為平實，從空，一步常法。如法得半徑。}



\clTSAOYUE{%

別得七十二步即大差也，斜行即弦，半徑即股也。

立天元一為半徑，以自之為股冪，又以二行差六十四以自之得$\square$為勾冪。

並二冪得$\square$為弦冪（寄左）。然後以斜行步自之，得$\square$為同數，

與左相消得$\square$。開平方得一百二十步，倍之即城徑也。合問。}

\end{leftcolumn}



\switchcolumn



\begin{rightcolumn}

\mdOWEN{乙出东门距离不详停下，甲出南门东行72步望见乙，向乙斜行136步与乙相会。求城墙直径。}



\mdFAYUE{斜行步数自乘为上位，以二行相减余数自乘减上位为二次方程常数项，一次项系数为零，

二次项系数为一。依法得半径。}



\mdTSAOYUE{%

辨识得：72步即为大差，斜行即为弦，半径即为股。

设天元一为半径，自乘为股冪，又以二行差64自乘得$\square$为勾冪。

二冪相加得$\square$为弦冪（寄左）。然后以斜行步自乘，得$\square$为同数，

与左式相消得$\square$。开平方得120步，加倍即城径。符合所问。}

\end{rightcolumn}



\end{paracol}



\vspace{0.3cm}

\begin{center}\longline\end{center}



% -------------------------------------------------------------------

% Problem 110

% -------------------------------------------------------------------

\begin{paracol}{2}



\begin{leftcolumn}

\clOWEN{甲出南門不知遠近而立，乙出東門南行三十步望見甲，卻就甲斜行二百五十五步與甲相會。問答同前。}



\clFAYUE{二行差自之為冪，以減於斜行冪為平實，一虛隅。得半徑。}



\clTSAOYUE{%

別得南行步即股弦差也，斜步即弦也，半徑即勾也。

乃立天元一為半城徑，以自之為冪。以二行相減餘二百二十五，以自之得$\square$為股冪。

二冪相並得$\square$為弦冪（寄左）。然後以斜行步自之，得$\square$為同數，

與左相消得下$\square$。開平方得一百二十步，即半徑也。合問。}

\end{leftcolumn}



\switchcolumn



\begin{rightcolumn}

\mdOWEN{甲出南门距离不详停下，乙出东门南行30步望见甲，向甲斜行255步与甲相会。求城墙直径。}



\mdFAYUE{二行差的平方作为被减数，从斜行平方中减去它作为二次方程常数项，负一为二次项系数。得半径。}



\mdTSAOYUE{%

辨识得：南行步即为股弦差，斜步即为弦，半径即为勾。

设天元一为半城径，自乘为冪。以二行相减余225，自乘得$\square$为股冪。

二冪相加得$\square$为弦冪（寄左）。然后以斜行步自乘，得$\square$为同数，

与左式相消得下式$\square$。开平方得120步，即半径。符合所问。}

\end{rightcolumn}



\end{paracol}



\vspace{0.3cm}

\begin{center}\longline\end{center}



% -------------------------------------------------------------------

% Problem 111

% -------------------------------------------------------------------

\begin{paracol}{2}



\begin{leftcolumn}

\clOWEN{甲出南門東行不知步數而立，乙出東門南行三十步望見甲，斜行一百二步相會。問答同前。}



\clFAYUE{二行相減，餘以乘乙南行，四之於上，又加入斜行冪為平實。得虛和一百三十八。}



\clTSAOYUE{%

別得斜步內減南行為甲東行步也。此問以弦外容圓入之，以二行相減數乘乙南行三十步，

得$\square$，又四之，得$\square$為二直積也。又加入斜步冪$\square$，

共得$\square$即和冪也。平方而一，得一百三十八步，即虛和也。

又加斜步得二百四十步，即城徑也。合問。}

\end{leftcolumn}



\switchcolumn



\begin{rightcolumn}

\mdOWEN{甲出南门东行步数不详停下，乙出东门南行30步望见甲，斜行102步相会。求城墙直径。}



\mdFAYUE{二行相减的余数乘以乙南行步数，四倍后加上斜行平方为常数项。得虚和138。}



\mdTSAOYUE{%

辨识得：斜步减去南行即为甲东行步数。此题以弦外容圆理解，以二行相减数乘乙南行30步，

得$\square$，再四倍得$\square$为两倍直积。加入斜步冪$\square$，

共得$\square$即和冪。开平方得138步，即虚和。

加上斜步得240步，即城径。符合所问。}

\end{rightcolumn}



\end{paracol}



\vspace{0.3cm}

\begin{center}\longline\end{center}



% -------------------------------------------------------------------

% Problem 112

% -------------------------------------------------------------------

\begin{paracol}{2}



\begin{leftcolumn}

\clOWEN{乙出東門南行不知步數而立，甲出南門東行七十二步望見乙，斜行一百二步與乙相會。問答同前。}



\clFAYUE{%

倍相減步以乘倍東行，得數復以減於斜步冪，餘為實。平方而一，得較也。

又以二行相減數乘倍東行為平實，以較為從方，得勾。

勾較共為長，又以斜步並入勾股共，即城徑。}



\clTSAOYUE{%

別得二行相減餘$\square$為乙南行步也。以此數又減於甲東行，餘四十二步即較也。

又以二行相減數$\square$乘倍東行得$\square$為平實，以較為從。

平方開得四十八即勾也。勾內加較得九十步即股也。

勾股共得一百三十八，又加入斜步，共得二百四十步，即城徑也。合問。}

\end{leftcolumn}



\switchcolumn



\begin{rightcolumn}

\mdOWEN{乙出东门南行步数不详停下，甲出南门东行72步望见乙，斜行102步与乙相会。求城墙直径。}



\mdFAYUE{%

两倍相减步数乘以两倍东行步数，得数从斜步平方中减去，余数为被除数。开平方得「较」。

又以二行相减数乘两倍东行作为二次方程常数项，以「较」为一次项，得勾。

勾与较共为长，以斜步加入勾股和，即城径。}



\mdTSAOYUE{%

辨识得：二行相减余$\square$为乙南行步。以此数再从甲东行中减去，余42步即「较」。

又以二行相减数$\square$乘两倍东行得$\square$为常数项，以「较」为一次项。

开平方得48即勾。勾加较得90步即股。

勾股共得138，加上斜步共得240步，即城径。符合所问。}

\end{rightcolumn}



\end{paracol}



\vspace{0.3cm}

\begin{center}\longline\end{center}



% -------------------------------------------------------------------

% Problem 113

% -------------------------------------------------------------------

\begin{paracol}{2}



\begin{leftcolumn}

\clOWEN{%

乙出南門東行，甲出東門南行，兩相望見。既而乙雲：

「我東行不及城徑一百六十八步。」甲雲：「我南行不及城徑二百一十步。」問答同前。}



\clFAYUE{%

半甲不及步以自之為冪，半甲不及步內減差以自之為冪。

二冪相並內卻減差冪為平實，四之甲不及內減三之乙不及，餘為益從，三步半虛法。得甲南行。}



\clTSAOYUE{%

別得乙不及為虛勾、半徑共，又為徑內減明勾也。

甲不及為虛股、半徑共，又為徑內減股也。

又二雲數相並為虛和、圓徑共也，雲數相減即虛較也。

乃立天元一為甲南行，以減於甲不及步又半之，得$\square$為虛股也。

虛股內減虛較得$\square$為虛勾。勾自之得$\square$為勾冪也，

又股自之下式$\square$為股冪也。二冪相並得$\square$為弦冪（寄左）。

然後以天元加虛較得$\square$為乙東行。又加入天元甲南行得$\square$為虛弦，

以自之得$\square$為同數，與左相消得$\square$。

開平方得三十步，即甲南行也。內加少步，即城徑也。合問。}

\end{leftcolumn}



\switchcolumn



\begin{rightcolumn}

\mdOWEN{%

乙出南门东行，甲出东门南行，两人相望看见。随后乙说：「我东行比城径少168步。」

甲说：「我南行比城径少210步。」求城墙直径。}



\mdFAYUE{%

甲不及步的一半自乘为冪，甲不及步的一半减去差再自乘为冪。

二冪相加再减去差的冪为二次方程常数项，

四倍甲不及减去三倍乙不及的余数为负一次项，3.5为负二次项系数。得甲南行。}



\mdTSAOYUE{%

辨识得：乙不及为虚勾与半半径之和，又为径减明勾。

甲不及为虚股与半半径之和，又为径减股。

二数相加为虚和与圆径之和，相减即虚较。

设天元一为甲南行，从甲不及步的一半中减去，得$\square$为虚股。

虚股减去虚较得$\square$为虚勾。勾自乘得$\square$为勾冪，

股自乘得$\square$为股冪。二冪相加得$\square$为弦冪（寄左）。

然后天元加虚较得$\square$为乙东行。再加天元（甲南行）得$\square$为虚弦，

自乘得$\square$为同数，与左式相消得$\square$。

开平方得30步，即甲南行。加上少步即城径。符合所问。}

\end{rightcolumn}



\end{paracol}



\vspace{0.3cm}

\begin{center}\longline\end{center}



% -------------------------------------------------------------------

% Problem 114

% -------------------------------------------------------------------

\begin{paracol}{2}



\begin{leftcolumn}

\clOWEN{%

丙出南門直行，甲出東門直行，兩相望見。既而丙雲：

「我行少於城徑一百五步。」甲雲：「我行少於城徑二百二十四步。」問答同前。}



\clFAYUE{%

二少步相乘訖，又自乘為實；六之共步，乘雲數相乘數為益從；

十八之雲數相乘於上，又三之共步，自乘加上位，內復減丙少步冪、甲少步冪為從廉；

四十八之共步為益二廉，六十三步常法。

翻法開三乘方，得一百二十步，即半徑。}



\clTSAOYUE{%

別得雲數共減於倍城徑為甲丙共行數。又雲數相減即皇極差，亦為甲行不及丙行數。

立天元一為半城徑，以三之，副置二位。

上位減丙少步，得$\square$為皇極股也，下位減甲少步得$\square$為皇極勾也。

勾股相乘得$\square$，以天元除之，得$\square$為弦也。

弦自之得$\square$為弦冪（寄左）。然後以股自之下$\square$為股冪於上，

又以勾自之得$\square$為勾冪，並以加入上位，得$\square$為同數，

與左相消得$\square$。翻法開三乘方得一百二十步，即半城徑也。合問。}

\end{leftcolumn}



\switchcolumn



\begin{rightcolumn}

\mdOWEN{%

丙出南门直行，甲出东门直行，两人相望看见。随后丙说：「我行比城径少105步。」

甲说：「我行比城径少224步。」求城墙直径。}



\mdFAYUE{%

二少步相乘后再自乘为常数项；六倍共步乘以两少步相乘为负一次项；

十八倍两少步相乘为上位，三倍共步自乘加上位，再减去丙少步冪、甲少步冪为从廉；

四十八倍共步为负第二廉，63为常数。

翻法开四次方程（三乘方），得120步，即半径。}



\mdTSAOYUE{%

辨识得：两少步之和从两倍城径中减去为甲丙共行数。两少步相减即皇极差，亦为甲行不及丙行数。

设天元一为半城径，三倍后分置两位。

上位减丙少步，得$\square$为皇极股；下位减甲少步得$\square$为皇极勾。

勾股相乘得$\square$，除以天元得$\square$为弦。

弦自乘得$\square$为弦冪（寄左）。然后股自乘得$\square$为股冪于上位，

勾自乘得$\square$为勾冪，合并加入上位得$\square$为同数，

与左式相消得$\square$。翻法开四次方程（三乘方）得120步，即半城径。符合所问。}

\end{rightcolumn}



\end{paracol}



\vspace{0.3cm}

\begin{center}\longline\end{center}



% -------------------------------------------------------------------

% Problem 115

% -------------------------------------------------------------------

\begin{paracol}{2}



\begin{leftcolumn}

\clOWEN{%

甲出東門直行，乙出南門直行，各不知步數而立。

乙望見甲，就甲斜行了二百八十九步與甲相會。

其二直行共得一百五十一步。又雲甲直行少於乙直行。問答同前。}



\clFAYUE{斜冪內減共步冪為平實，倍共步內減斜步為從，一常法。得徑。}



\clTSAOYUE{%

別得共數、城徑並即皇極和也。

立天元一為圓徑，加共步得$\square$為皇極和，以自之，得$\square$於上。

以斜行冪$\square$減上位餘$\square$為二直積（寄左）。

然後以天元乘斜步得$\square$，與左相消得$\square$。

開平方得二百四十步，即城徑也。合問。}

\end{leftcolumn}



\switchcolumn



\begin{rightcolumn}

\mdOWEN{%

甲出东门直行，乙出南门直行，各不知步数停下。乙望见甲，向甲斜行289步与甲相会。

两直行共得151步。又说甲直行少于乙直行。求城墙直径。}



\mdFAYUE{斜行冪减去共步冪为二次方程常数项，两倍共步减去斜步为一次项，1为常数。得城径。}



\mdTSAOYUE{%

辨识得：共数与城径之和即皇极和。

设天元一为圆径，加共步得$\square$为皇极和，自乘得$\square$于上位。

以斜行冪$\square$减上位余$\square$为两倍直积（寄左）。

然后天元乘斜步得$\square$，与左式相消得$\square$。

开平方得240步，即城径。符合所问。}

\end{rightcolumn}



\end{paracol}



\vspace{0.3cm}

\begin{center}\longline\end{center}



% -------------------------------------------------------------------

% Problem 116

% -------------------------------------------------------------------

\begin{paracol}{2}



\begin{leftcolumn}

\clOWEN{%

甲出東門直行，乙出東門南行，丙出南門直行，丁出南門東行，各不知步數而立。

四人遙相望，悉與城參相直。

只雲甲、丙共行了一百五十一步，乙、丁立處相距一百二步。又雲丙直行步多於甲直行步。問答同前。}



\clFAYUE{%

共步、距步相減，得數自之於上，以共步為冪內減上為平實，

二之距步內減共步、距步差為從，一步虛法。得城徑。}



\clTSAOYUE{%

別得共步得城徑即皇極和也，相距步即虛弦也。

皇極和內減虛弦即皇極弦也。又共步、距步差$\square$即皇極弦內減城徑也（此名旁差）。

乃立天元一為城徑，加共步得$\square$為皇極和也，以自之得$\square$於上，

以共步、距步差$\square$加天元得$\square$為皇極弦也，

以自之下式$\square$，減上位餘得$\square$為二直積（寄左）。

然後以天元徑乘皇極弦，得$\square$為同數，與左相消得$\square$。

開平方得二百四十步，即城徑也。合問。}

\end{leftcolumn}



\switchcolumn



\begin{rightcolumn}

\mdOWEN{%

甲出东门直行，乙出东门南行，丙出南门直行，丁出南门东行，各不知步数停下。

四人遥遥相望，都与城墙呈直角。

只说甲、丙共行了151步，乙、丁立处相距102步。又说丙直行步多于甲直行步。求城墙直径。}



\mdFAYUE{%

共步与距步相减，差数自乘为上位，共步冪减上位为二次方程常数项，

两倍距步减去共步与距步之差为一次项，负一为二次项系数。得城径。}



\mdTSAOYUE{%

辨识得：共步加城径即皇极和，相距步即虚弦。

皇极和减去虚弦即皇极弦。共步与距步之差$\square$即皇极弦减城径（此名「旁差」）。

设天元一为城径，加共步得$\square$为皇极和，自乘得$\square$于上位。

共步与距步之差$\square$加天元得$\square$为皇极弦，

自乘得$\square$，减上位余$\square$为两倍直积（寄左）。

然后天元径乘皇极弦，得$\square$为同数，与左式相消得$\square$。

开平方得240步，即城径。符合所问。}

\end{rightcolumn}



\end{paracol}



\vspace{0.3cm}

\begin{center}\longline\end{center}



% -------------------------------------------------------------------

% Problem 117

% -------------------------------------------------------------------

\begin{paracol}{2}



\begin{leftcolumn}

\clOWEN{%

甲出南門東行不知步數而立，乙出東門南行望見甲，復就甲斜行，與甲相會。

乙通計行了一百三十二步，其乙南行步不及斜行七十二步，其甲東行卻多於乙南行。問答同前。}



\clFAYUE{倍不及步在地，以不及步減通步以乘之為實，以四之不及步為法。得乙南行三十步。}



\clTSAOYUE{%

別得乙南行即股也，以減通步即虛弦也，以減不及步即虛較也，其不及步即甲東行也。

立天元一為乙南行，置不及步以天元乘之，又四之得$\square$元為二直積（寄左）。

然後倍不及步以為弦較和於上$\square$。

以不及步減通步得$\square$為弦較較。

以乘上位得$\square$太為同數，與左相消得$\square$。

上法下實，得三十步為乙南行也。餘各以數求之。}

\end{leftcolumn}



\switchcolumn



\begin{rightcolumn}

\mdOWEN{%

甲出南门东行步数不详停下，乙出东门南行望见甲，再向甲斜行与甲相会。

乙共行了132步，乙南行步比斜行少72步，甲东行多于乙南行。求城墙直径。}



\mdFAYUE{两倍不及步作为基数，不及步减通步后乘以基数为被除数，四倍不及步为除数。得乙南行30步。}



\mdTSAOYUE{%

辨识得：乙南行即为股，从通步中减去即为虚弦，减不及步即为虚较，不及步即甲东行。

设天元一为乙南行，取不及步以天元乘之，再四倍得$\square$元为两倍直积（寄左）。

然后两倍不及步作为弦较和于上位$\square$。

不及步减通步得$\square$为弦较较。

相乘得$\square$太为同数，与左式相消得$\square$。

上法下实，得30步为乙南行。其余各以数求之。}

\end{rightcolumn}



\end{paracol}



\vspace{0.3cm}

\begin{center}\longline\end{center}



% -------------------------------------------------------------------

% Problem 118

% -------------------------------------------------------------------

\begin{paracol}{2}



\begin{leftcolumn}

\clOWEN{%

甲出南門東行不知步數而立。乙出東門南行，望見甲，復斜行與甲相會。

二人共行了二百四步，又雲甲行不及共步一百三十二。問答同前。}



\clFAYUE{%

別得二行共即兩個虛弦也，其不及步即乙南行與一虛弦共也。

置不及步內減一弦餘三十步，即乙南行也。

以乙南行反以減虛弦，餘七十二步即甲東行也。

以乙南行減甲東行餘即虛較也。}



\clTSAOYUE{此問無草。}

\end{leftcolumn}



\switchcolumn



\begin{rightcolumn}

\mdOWEN{%

甲出南门东行步数不详停下。乙出东门南行，望见甲，再斜行与甲相会。

两人共行了204步，又说甲行比共步少132步。求城墙直径。}



\mdFAYUE{%

辨识得：二行共即两个虚弦，不及步即乙南行与一个虚弦之和。

从不及步中减去一弦余30步，即乙南行。

以乙南行从虚弦中减去，余72步即甲东行。

甲东行减乙南行余即虚较。}



\mdTSAOYUE{此题无演算草文。}

\end{rightcolumn}



\end{paracol}



\vspace{0.3cm}

\begin{center}\longline\end{center}



% -------------------------------------------------------------------

% Problem 119

% -------------------------------------------------------------------

\begin{paracol}{2}



\begin{leftcolumn}

\clOWEN{%

乙出東門南行，甲出西門南行，甲望見乙，斜行五百一十步相會。

乙雲：「我南行少於城徑二百一十步。」問答同前。}



\clFAYUE{少步冪為平實，四斜步內減二少步為益從，五步常法。得乙南行。}



\clTSAOYUE{%

別得少步為徑內減股。

立天元一為乙南行，以二之減於倍斜行步，得$\square$為梯底也。

以二之天元乘之，得$\square$為徑冪（寄左）。

再置天元加少步，得下式$\square$為城徑，以自之得$\square$，

與左相消得$\square$。開平方得三十步，即乙南行也。加少步即城徑也。合問。}

\end{leftcolumn}



\switchcolumn



\begin{rightcolumn}

\mdOWEN{%

乙出东门南行，甲出西门南行，甲望见乙，斜行510步相会。

乙说：「我南行比城径少210步。」求城墙直径。}



\mdFAYUE{少步冪为二次方程常数项，四倍斜步减去两倍少步为负一次项，5为常数。得乙南行。}



\mdTSAOYUE{%

辨识得：少步为径减股。

设天元一为乙南行，两倍后从两倍斜行步中减去，得$\square$为梯底。

以两倍天元乘之，得$\square$为径冪（寄左）。

再置天元加少步，得下式$\square$为城径，自乘得$\square$，

与左式相消得$\square$。开平方得30步，即乙南行。加少步即城径。符合所问。}

\end{rightcolumn}



\end{paracol}



\vspace{0.3cm}

\begin{center}\longline\end{center}



% -------------------------------------------------------------------

% Problem 120

% -------------------------------------------------------------------

\begin{paracol}{2}



\begin{leftcolumn}

\clOWEN{%

乙出南門東行，甲出北門東行，甲望見乙，斜行二百七十二步與乙相會。

乙雲：「我東行不及城徑一百六十八步。」問答同前。}



\clFAYUE{以不及步冪之為實，四斜內減二之不及步為虛從，五常法。平開得乙東行七十二步。}



\clTSAOYUE{%

別得不及步為城徑減明勾也。

立天元一為乙東行，以倍之減於二之斜行步，得下$\square$為梯底也。

倍天元乘之，得$\square$為徑冪（寄左）。

再置天元加不及步，得$\square$為城徑，以自之得$\square$為同數，

與左相消得$\square$。開平方得七十二步，即乙東行也。加入少步即城徑也。合問。}

\end{leftcolumn}



\switchcolumn



\begin{rightcolumn}

\mdOWEN{%

乙出南门东行，甲出北门东行，甲望见乙，斜行272步与乙相会。

乙说：「我东行比城径少168步。」求城墙直径。}



\mdFAYUE{不及步冪为被除数，四倍斜步减去两倍不及步为负一次项，5为常数。开平方得乙东行72步。}



\mdTSAOYUE{%

辨识得：不及步为城径减明勾。

设天元一为乙东行，两倍后从两倍斜行步中减去，得下式$\square$为梯底。

两倍天元乘之，得$\square$为径冪（寄左）。

再置天元加不及步，得$\square$为城径，自乘得$\square$为同数，

与左式相消得$\square$。开平方得72步，即乙东行。加入少步即城径。符合所问。}

\end{rightcolumn}



\end{paracol}



\vspace{0.3cm}

\begin{center}\longline\end{center}



% -------------------------------------------------------------------

% Problem 121

% -------------------------------------------------------------------

\begin{paracol}{2}



\begin{leftcolumn}

\clOWEN{%

乙出南門東行，丁出東門南行，卻有甲丙二人共在西北隅，甲向東行，丙向南行，

四人遙相望見，俱與城參相直。既而相會，甲雲：「我多乙二百四十八步。」

丙雲：「我多於丁五百七十步。」問答同前。}



\clFAYUE{二多步相乘為平實，並二多步而半之為從，七分半常法。得城徑。}



\clTSAOYUE{%

別得甲多步為大勾內減明勾也，丙多步為大股內少股也。

又乙東行得一虛勾為半徑，丁南行得一虛股為半徑。

又二多數相並得$\square$為大和內少虛弦也，又二少數相減餘$\square$為兩個角差。

又甲多步內減半徑即勾方差也，丙多步內減半徑即股方差也。

立天元一為城徑，以半之減於甲多步得$\square$為勾方差，

又以半徑減於丙多步得$\square$為股方差。

二差相乘得$\square$為徑冪（寄左）。然後以天元冪與左相消，得下式$\square$。

開平方得二百四十步，即城徑也。合問。}

\end{leftcolumn}



\switchcolumn



\begin{rightcolumn}

\mdOWEN{%

乙出南门东行，丁出东门南行，又有甲丙二人同在西北角，甲向东行，丙向南行，

四人遥遥相望，都与城墙呈直角。随后相会，甲说：「我比乙多248步。」

丙说：「我比丁多570步。」求城墙直径。}



\mdFAYUE{两多步相乘为二次方程常数项，两多步之和的一半为一次项，7.5为常数。得城径。}



\mdTSAOYUE{%

辨识得：甲多步为大勾减明勾，丙多步为大股减股。

乙东行得一虚勾为半径，丁南行得一虚股为半径。

两多数相加得$\square$为大和减去虚弦，两少数相减余$\square$为两个角差。

甲多步减去半半径即勾方差，丙多步减去半半径即股方差。

设天元一为城径，从甲多步中减去其一半得$\square$为勾方差，

再从丙多步中减去半半径得$\square$为股方差。

二差相乘得$\square$为径冪（寄左）。然后天元冪与左式相消，得下式$\square$。

开平方得240步，即城径。符合所问。}

\end{rightcolumn}



\end{paracol}



\vspace{0.3cm}

\begin{center}\longline\end{center}



% -------------------------------------------------------------------

% Problem 122

% -------------------------------------------------------------------

\begin{paracol}{2}



\begin{leftcolumn}

\clOWEN{%

甲丙二人俱在西北隅，甲向東行，丙向南行。

又乙出南門東行，丁出東門南行，各不知步數而立。

四人遙相望見，悉與城參相直。既而相會，甲雲：

「我與乙共行了三百九十二步。」丙雲：「我與丁共行了六百三十步。」問答同前。}



\clFAYUE{%

甲乙共自之為冪，丙丁共自之為冪，二冪又相乘為三乘方實。

甲乙共自之為冪，以丙丁共乘之於上。

又以丙丁共自之為冪，以甲乙共乘之加上位為益從。

甲乙共自之為冪，丙丁共自之為冪，並，以七分半乘之於上。

又以甲乙共乘丙丁共，得數減上位為第一益廉。

並二共數，以七分半乘之為第二廉。

以七分半自之，得五分六厘二毫五絲於上位。

以一步內減上位，餘四分三厘七毫五絲為虛隅。得城徑。}



\clTSAOYUE{%

別得甲為大勾，乙為明勾，丙為大股，丁為股也。

甲乙共內減半徑即是黃長弦也，丙丁共內減半徑即黃廣弦也。

黃長弦、黃廣弦二數相減，餘為兩個皇極差也。

乃立天元為城徑，半之副置二位。

上以減於甲乙共數，得$\square$即黃長弦也，以自之得$\square$為黃長弦冪也，

內減天元一冪，餘得下式$\square$為勾方差冪也。

下位以減於丙丁共，得下式$\square$即黃廣弦也，

以自之得$\square$為黃廣弦冪也，內減天元一冪，餘得$\square$為股方差冪也。

再以勾方差冪、股方差冪相乘，得$\square$為徑冪（寄左）。

然後以天元為冪，又以冪自之，與左相消得下式$\square$。

開三乘方得二百四十步，即城徑也。合問。}

\end{leftcolumn}



\switchcolumn



\begin{rightcolumn}

\mdOWEN{%

甲丙二人同在西北角，甲向东行，丙向南行。

乙出南门东行，丁出东门南行，各不知步数停下。

四人遥遥相望，都与城墙呈直角。随后相会，甲说：「我与乙共行了392步。」

丙说：「我与丁共行了630步。」求城墙直径。}



\mdFAYUE{%

甲乙共自乘为冪，丙丁共自乘为冪，二冪再相乘为四次方程（三乘方）常数项。

甲乙共自乘为冪，以丙丁共乘之于上位。

又以丙丁共自乘为冪，以甲乙共乘之加上位为负一次项。

甲乙共自乘为冪，丙丁共自乘为冪，合并以7.5倍乘之于上位。

又以甲乙共乘丙丁共，得数减上位为第一负廉。

二共数之和以7.5倍乘之为第二廉。

7.5自乘得56.25于上位。

以1减上位，余43.75为负二次项系数。得城径。}



\mdTSAOYUE{%

辨识得：甲为大勾，乙为明勾，丙为大股，丁为股。

甲乙共减去半半径即为黄长弦，丙丁共减去半半径即为黄广弦。

黄长弦、黄广弦二数相减，余为两个皇极差。

设天元为城径，取半后分置两位。

上位从甲乙共数中减去，得$\square$即黄长弦，自乘得$\square$为黄长弦冪，

减去天元一冪，余得下式$\square$为勾方差冪。

下位从丙丁共中减去，得下式$\square$即黄广弦，

自乘得$\square$为黄广弦冪，减去天元一冪，余得$\square$为股方差冪。

再以勾方差冪、股方差冪相乘，得$\square$为径冪（寄左）。

然后天元自乘，再自乘（四次方），与左式相消得下式$\square$。

开四次方程（三乘方）得240步，即城径。符合所问。}

\end{rightcolumn}



\end{paracol}



\newpage





% ===================================================================

% VOLUME 8

% ===================================================================

\section{\tradfont 卷八\quad 明後一十六問}

\label{sec:vol8}



\begin{center}

\fcolorbox{classicalborder}{classicalbg}{\parbox{0.9\textwidth}{\centering

\textbf{\large 卷八概要}\quad 本卷十六问，涉及多元天元术（连枝同体），以多个未知量建立关联方程组，求解更复杂的几何构形。方程次数多为四次，间有五次、六次方程。

}}

\end{center}



\vspace{0.5cm}



% -------------------------------------------------------------------

% Problem 123

% -------------------------------------------------------------------

\begin{paracol}{2}



\begin{leftcolumn}

\clOWEN{%

出南門向東有槐樹一株，出東門向南有柳樹一株。

丙丁俱出南門，丙直行，丁往至槐樹下。

甲乙俱出東門，甲直行，乙往至柳樹下。

四人遙相望見，各不知所行步數。

只雲丙丁共行了二百七步，甲乙共行了四十六步，

又雲甲丙立處相距二百八十九步。問答同前。}



\clFAYUE{以二共相減數又以減距數為實，二為法。得平勾。}



\clTSAOYUE{%

識別得丙丁共即明和也，甲乙共即和也，相距步即極弦也。

二共相並即極弦內少個虛黃也，又為極和內少個虛和也。

二共相減餘為平勾高股差也，又為虛差極差共也，又為通差內減極差也。

立天元為平勾，加入二共相減數，得$\square$為高股，又加天元得$\square$為極弦（寄左）。

以相距步二百八十九與左相消，得$\square$。

上法下實，如法得六十四，即平勾也。

以二共相減數加平勾得二百二十五為高股，復以平勾乘之，得一萬四千四百步。

開平方得一百二十步，即城半徑也。合問。}

\end{leftcolumn}



\switchcolumn



\begin{rightcolumn}

\mdOWEN{%

出南门向东有槐树一株，出东门向南有柳树一株。

丙丁都出南门，丙直行，丁前往槐树下。

甲乙都出东门，甲直行，乙前往柳树下。

四人遥遥望见，各不知所行步数。

只说丙丁共行了207步，甲乙共行了46步，

又说甲丙立处相距289步。求城墙直径。}



\mdFAYUE{以两共数相减的差再从距数中减去为被除数，2为除数。得平勾。}



\mdTSAOYUE{%

辨识得：丙丁共即明和，甲乙共即和，相距步即极弦。

两共数相加即极弦内少一个虚黄，又为极和内少一个虚和。

两共数相减余为平勾与高股之差，又为虚差与极差之和，又为通差减去极差。

设天元为平勾，加入两共相减数，得$\square$为高股，再加天元得$\square$为极弦（寄左）。

以相距步289与左式相消，得$\square$。

上法下实，依法得64，即平勾。

以两共相减数加平勾得225为高股，再以平勾乘之，得14,400步。

开平方得120步，即城半半径。符合所问。}

\end{rightcolumn}



\end{paracol}



\vspace{0.3cm}

\begin{center}\longline\end{center}



% -------------------------------------------------------------------

% Problem 124

% -------------------------------------------------------------------

\begin{paracol}{2}



\begin{leftcolumn}

\clOWEN{%

依前見丙丁共二百七步，甲乙共四十六步。又雲二樹相去一百二步。問答同前。}



\clFAYUE{%

以甲乙共乘樹相去步，得數又以自之為平實；從空；

並二共數為冪於上，內減甲乙共自之數、丙丁共自之數為益隅。得弦。}



\clTSAOYUE{%

識別得兩樹相去步即虛弦也。餘數具前。

立天元一為弦。置明和以天元乘之，合和除，不除，便以$\square$元為明弦也（內帶和分母）。

乃置虛弦以分母和乘之，得$\square$太，加入明弦，得$\square$為極股也，內帶和分母。

以自之得下式$\square$為極股冪（內寄和冪為分母）。

又以天元加虛弦得下$\square$為極勾，以自之得$\square$。

又以和冪$\square$乘之，得$\square$為勾冪也。

勾股冪相並得$\square$為兩積一較冪也，內有和冪分母（寄左）。

然後置明弦$\square$元於上，以和乘天元加上位，得$\square$元為二弦並也。

又置虛弦以和乘之得$\square$太，並入上位，得下式$\square$為極弦，

以自之得$\square$為同數，與左相消得$\square$。

開平方得三十四步，即弦也。}

\end{leftcolumn}



\switchcolumn



\begin{rightcolumn}

\mdOWEN{%

如前见丙丁共207步，甲乙共46步。又说两树相距102步。求城墙直径。}



\mdFAYUE{%

以甲乙共乘树相距步数，得数再自乘为二次方程常数项；一次项为零；

两共数之和自乘于上位，减去甲乙共自乘数、丙丁共自乘数为负二次项系数。得弦。}



\mdTSAOYUE{%

辨识得：两树相距步数即虚弦。其余数如前。

设天元一为弦。取明和以天元乘之，应以和除，不除，便以$\square$元为明弦（内带和分母）。

取虚弦以分母和乘之，得$\square$太，加入明弦，得$\square$为极股，内带和分母。

自乘得下式$\square$为极股冪（内寄和冪为分母）。

又以天元加虚弦得下式$\square$为极勾，自乘得$\square$。

又以和冪$\square$乘之，得$\square$为勾冪。

勾股冪合并得$\square$为两积一较冪，内有和冪分母（寄左）。

然后取明弦$\square$元于上位，以和乘天元加上位，得$\square$元为两弦和。

再取虚弦以和乘之得$\square$太，合并入上位，得下式$\square$为极弦，

自乘得$\square$为同数，与左式相消得$\square$。

开平方得34步，即弦。}

\end{rightcolumn}



\end{paracol}



\vspace{0.3cm}

\begin{center}\longline\end{center}



% -------------------------------------------------------------------

% Problem 125

% -------------------------------------------------------------------

\begin{paracol}{2}



\begin{leftcolumn}

\clOWEN{皇極大小差共一百八十七步，明黃、黃共六十六步。問答同前。}



\clFAYUE{後數自乘為實，前後數相減，餘為法。得虛黃方三十六。}



\clTSAOYUE{%

別得一百八十七即明弦二弦共也，其六十六即太虛大小差共也。

又二數相並，得$\square$即明、二和共。若以相減，餘$\square$即明、四差共也。

立天元一為太虛黃方面，加二黃共，得$\square$即虛弦也。

倍虛弦又加天元，得$\square$即城徑也。

又以虛弦加皇極大小差，得$\square$即極弦也，

以極弦乘城徑得$\square$，為兩段皇極勾股積（寄左）。

再以極弦虛弦相並，得$\square$，即皇極勾股共也，

自之得$\square$，內減皇極弦冪$\square$，得$\square$為同數，

與寄左相消得$\square$。上法下實，如法得三十六步，即太虛黃方面也。合問。}

\end{leftcolumn}



\switchcolumn



\begin{rightcolumn}

\mdOWEN{皇极大小差共187步，明黄、黄共66步。求城墙直径。}



\mdFAYUE{后数自乘为被除数，前后数相减余为除数。得虚黄方36。}



\mdTSAOYUE{%

辨识得：187即明弦与二弦之和，66即太虚大小差之和。

两数相加得$\square$即明与二和之和。若相减余$\square$即明与四差之和。

设天元一为太虚黄方面，加二黄共得$\square$即虚弦。

两倍虚弦再加天元得$\square$即城径。

又以虚弦加皇极大小差得$\square$即极弦，

极弦乘城径得$\square$为两段皇极勾股积（寄左）。

再以极弦虚弦相加得$\square$即皇极勾股之和，

自乘得$\square$，减去皇极弦冪$\square$得$\square$为同数，

与寄左相消得$\square$。上法下实，依法得36步，即太虚黄方面。符合所问。}

\end{rightcolumn}



\end{paracol}



\vspace{0.3cm}

\begin{center}\longline\end{center}



% -------------------------------------------------------------------

% Problem 126

% -------------------------------------------------------------------

\begin{paracol}{2}



\begin{leftcolumn}

\clOWEN{%

東門南有柳一株，南門東有槐一株。

甲出東門直行，丙出南門直行。

甲、丙、柳、槐悉與城參相直，

既而甲就柳樹斜行三十四步至柳樹下，丙就槐樹斜行一百五十三步至槐樹下。問答同前。}



\clFAYUE{%

雲數相乘，倍之便為平方實，開方得虛弦一百二步。

以此弦加甲行步即極勾，以此弦加丙行步即極股。餘各依法求之。}



\clTSAOYUE{%

識別：甲斜行即弦也，丙斜行即明弦也。無草。}

\end{leftcolumn}



\switchcolumn



\begin{rightcolumn}

\mdOWEN{%

东门南有柳树一株，南门东有槐树一株。

甲出东门直行，丙出南门直行。

甲、丙、柳树、槐树都与城墙呈直角，

随后甲向柳树斜行34步到柳树下，丙向槐树斜行153步到槐树下。求城墙直径。}



\mdFAYUE{%

两已知数相乘，翻倍即为二次方程常数项，开平方得虚弦102步。

以此弦加甲行步即极勾，以此弦加丙行步即极股。其余依法求之。}



\mdTSAOYUE{%

辨识：甲斜行即弦，丙斜行即明弦。无演算草文。}

\end{rightcolumn}



\end{paracol}



\vspace{0.3cm}

\begin{center}\longline\end{center}



% -------------------------------------------------------------------

% Problem 127

% -------------------------------------------------------------------

\begin{paracol}{2}



\begin{leftcolumn}

\clOWEN{%

東門南有柳一株，南門東有槐一株。

甲出東門直行，丙出南門直行，二人遙相望，槐柳與城邊悉相直。

既而甲復斜行至柳樹下，丙復斜行至槐樹下，各不知步數。

只雲丙共行了二百八十八步，甲斜行與柳至東門步共得六十四步。問答同前。}



\clFAYUE{%

二雲數相乘於上，以六十四步自之，又二之，減上位為平實。

四之六十四於上，倍丙行內減上位為從。二十常法。得甲直行步一十六。}



\clTSAOYUE{%

別得丙共步即明股、明弦和也，六十四即平勾也，內甲斜行即弦也，柳至東門步即股也。

又二雲數相並即明差與極弦共也，二雲數相減即明差與平勾高股差共也。

又平勾內減勾即虛勾也。

立天元一為勾。置丙共步以天元乘之，復以六十四除之得$\square$為明勾也。

又以天元減於六十四，得$\square$為虛勾也，並虛明二勾$\square$為半徑也。

以自之得$\square$，倍之得$\square$為半段圓城徑冪（寄左）。

乃以天元加六十四，得$\square$為勾圓差於上；

又以明勾加丙共步，得$\square$為股圓差於下。

上下相乘得$\square$為同數，與左相消得$\square$。

開平方得一十六步，即勾也。此勾乃甲出東門直行步也。餘皆依數求之。合問。}

\end{leftcolumn}



\switchcolumn



\begin{rightcolumn}

\mdOWEN{%

东门南有柳树一株，南门东有槐树一株。

甲出东门直行，丙出南门直行，二人遥遥相望，槐树柳树与城边都呈直角。

随后甲再斜行至柳树下，丙再斜行至槐树下，各不知步数。

只说丙共行了288步，甲斜行与柳至东门步共得64步。求城墙直径。}



\mdFAYUE{%

两已知数相乘于上位，以64步自乘，再两倍，减上位为二次方程常数项。

四倍64于上位，两倍丙行减去上位为一次项。20为常数。得甲直行步16。}



\mdTSAOYUE{%

辨识得：丙共步即明股、明弦之和，64即平勾，其中甲斜行即弦，柳至东门步即股。

两已知数之和即明差与极弦之和，相减即明差与平勾高股差之和。

平勾减去勾即虚勾。

设天元一为勾。取丙共步以天元乘之，再以64除得$\square$为明勾。

又以天元从64中减去得$\square$为虚勾，虚勾与明勾相加$\square$为半径。

自乘得$\square$，翻倍得$\square$为半段圆城径冪（寄左）。

以天元加64得$\square$为勾圆差于上位；

又以明勾加丙共步得$\square$为股圆差于下位。

上下相乘得$\square$为同数，与左式相消得$\square$。

开平方得16步，即勾。此勾即甲出东门直行步。其余依数求之。符合所问。}

\end{rightcolumn}



\end{paracol}



\vspace{0.3cm}

\begin{center}\longline\end{center}



% -------------------------------------------------------------------

% Problem 128

% -------------------------------------------------------------------

\begin{paracol}{2}



\begin{leftcolumn}

\clOWEN{%

東門南有柳樹一株，南門東有槐樹一株。

甲出東門直行，丙出南門直行，二人遙相望，槐柳與城邊悉相直。

既而甲復斜行至柳樹下，丙復斜行至槐樹下，各不知步數。

只雲甲共行五十步，丙斜行與槐至南門步共得二百二十五步。問答同前。}



\clFAYUE{%

以二百二十五步自之為冪，又以此冪自為冪於上。

置甲共行以二百二十五步三度乘之，得數復折半，減上位為平實。

置二百二十五步自之數，以二雲數相減數乘之，又倍之於上。

倍五十步在地，以二百二十五步自之數乘之，復折半加上位為益從。

雲數相減自乘於上，以雲數相乘，復折半，減上位為常法。得明股。}



\clTSAOYUE{%

識別得甲共步即勾、弦共也，二百二十五即高股也，內丙斜行即明弦，槐至南門步即明勾也。

又二雲數相並即極弦內減一個差也，雲數相減即差與高股、平勾差共也。

又高股內減明股即虛股也。

立天元一為明股，即丙出南門直行步也。

置五十步以天元乘之得$\square$元，合高股除。

不除，便以此$\square$元為股也，內帶高股$\square$分母。

再置高股內減天元得$\square$為虛股，以分母高股乘之，得下式$\square$，

加入股得$\square$即半徑也。

以自增乘得下$\square$為半徑冪也。內帶高股冪為母（寄左）。

然後置甲共步以分母高股乘之，得$\square$太，加入股得$\square$為勾圓差於上（內帶高股分母）。

又以天元加高股得$\square$為股圓差於下。

上、下相乘得$\square$。又以分母高股乘之，得$\square$，

復折半得$\square$為同數，與左相消得$\square$。

開平方得一百三十五步，即明股也。合問。}

\end{leftcolumn}



\switchcolumn



\begin{rightcolumn}

\mdOWEN{%

东门南有柳树一株，南门东有槐树一株。

甲出东门直行，丙出南门直行，二人遥遥相望，槐树柳树与城边都呈直角。

随后甲再斜行至柳树下，丙再斜行至槐树下，各不知步数。

只说甲共行50步，丙斜行与槐至南门步共得225步。求城墙直径。}



\mdFAYUE{%

以225步自乘为冪，再以这冪自乘于上位。

取甲共行以225步三次乘之（立方），得数折半减上位为二次方程常数项。

取225步自乘之数，以两已知数相减之差乘之，再两倍于上位。

50步翻倍作为基数，以225步自乘之数乘之，折半加上位为负一次项。

两已知数之差自乘于上位，以两已知数相乘折半减上位为常数。得明股。}



\mdTSAOYUE{%

辨识得：甲共步即勾、弦之和，225即高股，其中丙斜行即明弦，槐至南门步即明勾。

两已知数之和即极弦内减一个差，相减即差与高股、平勾差之和。

高股减去明股即虚股。

设天元一为明股，即丙出南门直行步。

取50步以天元乘之得$\square$元，应以高股除。

不除，便以此$\square$元为股，内带高股$\square$分母。

再取高股减去天元得$\square$为虚股，以分母高股乘之，得下式$\square$，

加入股得$\square$即半径。

自乘得下式$\square$为半半径冪。内带高股冪为分母（寄左）。

然后取甲共步以分母高股乘之，得$\square$太，加入股得$\square$为勾圆差于上位（内带高股分母）。

又以天元加高股得$\square$为股圆差于下位。

上下相乘得$\square$。又以分母高股乘之，得$\square$，

折半得$\square$为同数，与左式相消得$\square$。

开平方得135步，即明股。符合所问。}

\end{rightcolumn}



\end{paracol}



\vspace{0.3cm}

\begin{center}\longline\end{center}



% -------------------------------------------------------------------

% Problem 129

% -------------------------------------------------------------------

\begin{paracol}{2}



\begin{leftcolumn}

\clOWEN{通勾、通弦共一千步，勾、弦共五十步。問答同前。}



\clFAYUE{%

置一千減二之五十步為泛率。以自乘，復半之於上。

又置泛率復以五十乘之，加上位為平實。

二十二之泛率於上。以四十二乘五十，得數內減泛率，加上位為益從。

二百為常法。得股。}



\clTSAOYUE{%

立天元一為股。

置一千以天元乘之，以五十除之，得$\square$元為通股也。

又以天元加五十步，得$\square$即小差也，通股加小差得$\square$即通弦也。

以通弦減一千得$\square$即通勾也。

以小差減通勾得$\square$即圓徑也。

以圓徑減通股得$\square$即大差也。

置大差以小差乘之得$\square$（寄左）。

然後置圓徑以自之得$\square$，折半得$\square$，與左相消得$\square$。

開平方得三十步，即股也。合問。}

\end{leftcolumn}



\switchcolumn



\begin{rightcolumn}

\mdOWEN{通勾、通弦共1000步，勾、弦共50步。求城墙直径。}



\mdFAYUE{%

取1000减去两倍50步为泛率。泛率自乘折半于上位。

又取泛率再以50乘之，加上位为二次方程常数项。

22倍泛率于上位。42乘50得数减去泛率，加上位为负一次项。

200为常数。得股。}



\mdTSAOYUE{%

设天元一为股。

取1000以天元乘之，以50除之，得$\square$元为通股。

又以天元加50步得$\square$即小差，通股加小差得$\square$即通弦。

以通弦减1000得$\square$即通勾。

以小差减通勾得$\square$即圆径。

以圆径减通股得$\square$即大差。

取大差以小差乘之得$\square$（寄左）。

然后圆径自乘得$\square$，折半得$\square$，与左式相消得$\square$。

开平方得30步，即股。符合所问。}

\end{rightcolumn}



\end{paracol}



\vspace{0.3cm}

\begin{center}\longline\end{center}



% -------------------------------------------------------------------

% Problem 130

% -------------------------------------------------------------------

\begin{paracol}{2}



\begin{leftcolumn}

\clOWEN{通勾、通弦共一千步，明勾、明弦共二百二十五步。問答同前。}



\clFAYUE{%

以後數再自乘，又以前數乘之為平實；以後數為冪，復以前數乘之為從；

以前數冪為虛常法。得明股。}



\clTSAOYUE{%

別得二百二十五步即高股也。

立天元一為明股。置一千以天元乘之，合以高股除不受除，

便以此一零零零元為通股（內帶高股為母）。

以天元加高股，得$\square$即大差也。

置大差以高股分母乘之，得$\square$，即帶分大差也。

以此減於通股，餘$\square$即圓徑也。

以自增乘，得$\square$（寄左。內帶高股冪分母）。

然後置一千以高股分母通之，得$\square$太，內減帶分大差，

得$\square$為兩個通勾也。

內減兩個圓徑得$\square$，為兩個小差也。

以帶分大差乘之，得下式$\square$為同數，與左相消得$\square$。

開平方得一百三十五步，即明股也。合問。}

\end{leftcolumn}



\switchcolumn



\begin{rightcolumn}

\mdOWEN{通勾、通弦共1000步，明勾、明弦共225步。求城墙直径。}



\mdFAYUE{%

以后数（225）自乘两次再以前数（1000）乘之为二次方程常数项；

以后数为冪再以前数乘之为一次项；

以前数冪为负常数。得明股。}



\mdTSAOYUE{%

辨识得：225步即高股。

设天元一为明股。取1000以天元乘之，应以高股除，不除，

便以此1000元为通股（内带高股为分母）。

以天元加高股得$\square$即大差。

取大差以高股分母乘之，得$\square$即带分大差。

以此减于通股，余$\square$即圆径。

自乘得$\square$（寄左。内带高股冪分母）。

然后取1000以高股分母通之，得$\square$太，内减带分大差，

得$\square$为两个通勾。

内减两个圆径得$\square$为两个小差。

以带分大差乘之，得下式$\square$为同数，与左式相消得$\square$。

开平方得135步，即明股。符合所问。}

\end{rightcolumn}



\end{paracol}



\vspace{0.3cm}

\begin{center}\longline\end{center}



% -------------------------------------------------------------------

% Problem 131

% -------------------------------------------------------------------

\begin{paracol}{2}



\begin{leftcolumn}

\clOWEN{通股、通弦共一千二百八十步，股、弦共六十四步。問答同前。}



\clFAYUE{%

雲數相乘為平實，前數為益從。

置前數以後數除之，得二十為泛率。

泛率減一以自乘於上，又倍泛率減一加上位為常法。倒積開得勾一十六。}



\clTSAOYUE{%

別得六十四步即平勾也。

立天元一為勾。置前數以天元乘之，以後數除之，得$\square$元即通勾也。

又置天元加後數，得$\square$即小差也。

以小差減通勾，餘$\square$即圓徑也。

以自之得$\square$（寄左）。

然後以小差減於前數，得$\square$為二通股。

內減兩個圓徑，得$\square$為二大差也。

以小差乘之得下$\square$，與左相消得$\square$。

開平方得一十六步，即勾也。合問。}

\end{leftcolumn}



\switchcolumn



\begin{rightcolumn}

\mdOWEN{通股、通弦共1280步，股、弦共64步。求城墙直径。}



\mdFAYUE{%

两已知数相乘为二次方程常数项，前数为负一次项。

取前数以后数除之得20为泛率。

泛率减1自乘于上位，两倍泛率减1加上位为常数。倒积开方得勾16。}



\mdTSAOYUE{%

辨识得：64步即平勾。

设天元一为勾。取前数以天元乘之，以后数除之，得$\square$元即通勾。

又取天元加后数得$\square$即小差。

以小差减通勾，余$\square$即圆径。

自乘得$\square$（寄左）。

然后以小差减于前数得$\square$为二通股。

内减两个圆径得$\square$为二大差。

以小差乘之得下式$\square$，与左式相消得$\square$。

开平方得16步，即勾。符合所问。}

\end{rightcolumn}



\end{paracol}



\vspace{0.3cm}

\begin{center}\longline\end{center}



% -------------------------------------------------------------------

% Problem 132

% -------------------------------------------------------------------

\begin{paracol}{2}



\begin{leftcolumn}

\clOWEN{通股、通弦共一千二百八十步，明股、明弦共二百八十八步。問答同前。}



\clFAYUE{%

二數相減，以後數乘之，內減後數冪，又半之為泛率。以自之為平實。

置前數加二之後數而半之為次率，以乘泛率，倍之於上，

以後數乘泛率減上位為益從。

次率自乘於上，以前數加次率，復以後數乘之，減上位為隅法。得明勾。}



\clTSAOYUE{%

別得二數相減，餘$\square$為通勾、通股及明勾共也。

立天元一為明勾。置前數以天元乘之，合以後數除之。

不除，便以此$\square$元為通勾也（內寄後數分母）。

又以二數相減，得數內又減天元得$\square$為通和也。

乃以分母二百八十八之，得下式$\square$，內減通勾，餘$\square$為通股也。

又以天元加後數，又以分母（即後數也）通之，得$\square$為大差也。

以此大差減於通股，得下式$\square$，為一個圓徑也。

半之得$\square$，以自之得$\square$為半徑冪（寄左）。

然後以半圓徑減通勾，得$\square$為底勾，又以天元乘之，

又以分母二百八十八之，得$\square$為同數，與左相消得$\square$。

開平方得七十二步，即明勾也。合問。}

\end{leftcolumn}



\switchcolumn



\begin{rightcolumn}

\mdOWEN{通股、通弦共1280步，明股、明弦共288步。求城墙直径。}



\mdFAYUE{%

两数相减，以后数乘之，减去后数冪，折半为泛率。泛率自乘为二次方程常数项。

取前数加两倍后数折半为次率，以次率乘泛率，两倍于上位，

以后数乘泛率减上位为负一次项。

次率自乘于上位，以前数加次率再以后数乘之，减上位为二次项系数。得明勾。}



\mdTSAOYUE{%

辨识得：两数相减余$\square$为通勾、通股及明勾之和。

设天元一为明勾。取前数以天元乘之，应以后数除之。

不除，便以此$\square$元为通勾（内寄后数为分母）。

又以两数相减，得数内再减天元得$\square$为通和。

以分母288通之，得下式$\square$，内减通勾余$\square$为通股。

又以天元加后数，再以分母（即后数）通之，得$\square$为大差。

以此大差减于通股，得下式$\square$为一个圆径。

折半得$\square$，自乘得$\square$为半半径冪（寄左）。

然后以半圆径减通勾得$\square$为底勾，再以天元乘之，

又以分母288通之，得$\square$为同数，与左式相消得$\square$。

开平方得72步，即明勾。符合所问。}

\end{rightcolumn}



\end{paracol}



\vspace{0.3cm}

\begin{center}\longline\end{center}



% -------------------------------------------------------------------

% Problem 133

% -------------------------------------------------------------------

\begin{paracol}{2}



\begin{leftcolumn}

\clOWEN{%

明股、明弦並二百八十八步，勾、弦並五十步。

又雲明股、勾並多於虛弦四十九步。問答同前。}



\clFAYUE{前二數相並，內減二之多步，即圓徑。又只以前二數相乘，便是半徑冪。}



\clTSAOYUE{%

識別得前二數相減而半之，即極差也。其多步名傍差，又為圓徑不及極弦數。}

\end{leftcolumn}



\switchcolumn



\begin{rightcolumn}

\mdOWEN{%

明股、明弦和288步，勾、弦和50步。

又说明股、勾之和比虚弦多49步。求城墙直径。}



\mdFAYUE{前两数相加，从中减去两倍多步，即圆径。又以前两数相乘，便是半半径冪。}



\mdTSAOYUE{%

辨识得：前两数相减折半，即极差。其多步名为「傍差」，又为圆径不及极弦之数。}

\end{rightcolumn}



\end{paracol}



\vspace{0.3cm}

\begin{center}\longline\end{center}



% -------------------------------------------------------------------

% Problem 134

% -------------------------------------------------------------------

\begin{paracol}{2}



\begin{leftcolumn}

\clOWEN{平差、高差共一百六十一步，明股、勾並多於虛弦四十九步。問答同前。}



\clFAYUE{二數相減又半之，以自乘為實，後數為法。得平勾。}



\clTSAOYUE{%

立天元一為平勾。以加前數得$\square$為高股也。

又以天元加高股得$\square$為極弦，內減後數得$\square$，

又半之，得$\square$為半徑。

以自之，得$\square$（寄左）。

然後以天元乘高股，得$\square$為同數，與左相消得$\square$。

上法下實，得六十四步，即平勾也。合問。}

\end{leftcolumn}



\switchcolumn



\begin{rightcolumn}

\mdOWEN{平差、高差共161步，明股、勾之和比虚弦多49步。求城墙直径。}



\mdFAYUE{两数相减折半，自乘为被除数，后数为除数。得平勾。}



\mdTSAOYUE{%

设天元一为平勾。加上前数得$\square$为高股。

又以天元加高股得$\square$为极弦，减去后数得$\square$，

折半得$\square$为半径。

自乘得$\square$（寄左）。

然后以天元乘高股得$\square$为同数，与左式相消得$\square$。

上法下实，得64步，即平勾。符合所问。}

\end{rightcolumn}



\end{paracol}



\vspace{0.3cm}

\begin{center}\longline\end{center}



% -------------------------------------------------------------------

% Problem 135

% -------------------------------------------------------------------

\begin{paracol}{2}



\begin{leftcolumn}

\clOWEN{%

平勾、高股差一百六十一步，明差、差並七十七步。

又雲極弦多於城徑四十九步。問答同前。}



\clFAYUE{%

並上二位而半之為平率。其四十九即旁率也。

副置平率，上加旁率，下減旁率，以相乘為實。倍旁差為法。得勾圓差。}



\clTSAOYUE{%

立天元一為半徑，又為半之股圓差上弦較較，又為半之勾圓差上弦較和也。

內減勾圓差上勾股較$\square$，餘$\square$為半之勾圓差上弦較較也。

置股圓差上勾股較$\square$，以半之勾圓差上弦較較乘之，得$\square$（寄左）。

然後以半之股圓差上弦較較乘勾圓差上勾股較，

得$\square$元為同數，與左相消得下式$\square$。

上法下實，得一百二十步，即半徑也。合問。}

\end{leftcolumn}



\switchcolumn



\begin{rightcolumn}

\mdOWEN{%

平勾、高股差161步，明差、差共77步。

又说极弦比城径多49步。求城墙直径。}



\mdFAYUE{%

合并上两位折半为平率。49即旁率。

分置平率，上加旁率，下减旁率，相乘为被除数。两倍旁差为除数。得勾圆差。}



\mdTSAOYUE{%

设天元一为半径，又为半之股圆差上弦较较，又为半之勾圆差上弦较和。

减去勾圆差上勾股较$\square$，余$\square$为半之勾圆差上弦较较。

取股圆差上勾股较$\square$，以半之勾圆差上弦较较乘之，得$\square$（寄左）。

然后以半之股圆差上弦较较乘勾圆差上勾股较，

得$\square$元为同数，与左式相消得下式$\square$。

上法下实，得120步，即半径。符合所问。}

\end{rightcolumn}



\end{paracol}



\vspace{0.3cm}

\begin{center}\longline\end{center}



% -------------------------------------------------------------------

% Problem 136

% -------------------------------------------------------------------

\begin{paracol}{2}



\begin{leftcolumn}

\clOWEN{%

又依前問：見角差一百六十一步，見明差差並七十七步，又見太虛弦較較六十步。問答同前。}



\clFAYUE{%

前二數相減而半之，得數加入半之太虛弦較較為泛率，以自乘為平實。

置一百六十一內減二之泛率為從，一常法。得平勾。}



\clTSAOYUE{%

別得$\square$即二股也。

立天元一為平勾。先以前二數相減而半之，得$\square$為虛差。

以虛差加股得$\square$，即明勾也。

以明勾加天元得$\square$，為平弦，以自之得$\square$，

內減天元冪，得$\square$為半徑冪（寄左）。

然後以天元加一百六十一為高股，以天元乘之，得$\square$為同數，

與左相消得$\square$。開平方得六十四步，即平勾也。}

\end{leftcolumn}



\switchcolumn



\begin{rightcolumn}

\mdOWEN{%

又如前问：见角差161步，见明差差共77步，又见太虚弦较较60步。求城墙直径。}



\mdFAYUE{%

前两数相减折半，得数加入太虚弦较较的一半为泛率，泛率自乘为二次方程常数项。

取161减去两倍泛率为一次项，1为常数。得平勾。}



\mdTSAOYUE{%

辨识得：$\square$即二股。

设天元一为平勾。先以前两数相减折半得$\square$为虚差。

以虚差加股得$\square$即明勾。

以明勾加天元得$\square$为平弦，自乘得$\square$，

减去天元冪得$\square$为半半径冪（寄左）。

然后以天元加161为高股，以天元乘之得$\square$为同数，

与左式相消得$\square$。开平方得64步，即平勾。}

\end{rightcolumn}



\end{paracol}



\vspace{0.3cm}

\begin{center}\longline\end{center}



% -------------------------------------------------------------------

% Problem 137

% -------------------------------------------------------------------

\begin{paracol}{2}



\begin{leftcolumn}

\clOWEN{高差、平差並一百六十一步，明差、差並七十七步。問答同前。}



\clFAYUE{%

以前數自乘於上，二數相並而半之，以自乘減上位，得數復自增乘為平實。

前數自之於上，又以四之前數乘之，寄位。

以前數自之於上，並二數而半之，以自乘減上位，得數又以四之前數乘，又倍之，減於寄位為從。

前數自之，又四之於上。又以四之前數為冪，加上位，權寄。

以前數為冪於上，並二數而半之，以自乘減上位，得數復八之於上。

又以四之前數為冪加入上位，並以減於權寄為常法。得平勾。}



\clTSAOYUE{%

識別得二位相並而半之，得$\square$，即極差也。

立天元一為平勾，加一百六十一得$\square$為高股，高股內又加天元，得$\square$為極弦。

以自之得$\square$於上，內減極差冪一萬四千一百六十一，

餘$\square$為兩段極積。合以極弦除，不除寄為母，便以此為城徑。

以自增乘得$\square$為圓徑冪（內有極弦冪分母。寄左）。

然後以天元乘高股，又四之得$\square$，又以分母極弦冪$\square$通之，

得$\square$為同數，與左相消得$\square$。

開平方得六十四步，即平勾也。合問。}

\end{leftcolumn}



\switchcolumn



\begin{rightcolumn}

\mdOWEN{高差、平差共161步，明差、差共77步。求城墙直径。}



\mdFAYUE{%

以前数自乘于上位，两数相并折半自乘减上位，得数再自乘为二次方程常数项。

前数自乘于上位，又以四倍前数乘之，寄存。

前数自乘于上位，两数相并折半自乘减上位，得数又以四倍前数乘，再两倍，减于寄存为一次项。

前数自乘又四倍于上位。又以四倍前数为冪加上位，暂寄。

以前数为冪于上位，两数相并折半自乘减上位，得数再八倍于上位。

又以四倍前数为冪加入上位，一并减于暂寄为常数。得平勾。}



\mdTSAOYUE{%

辨识得：两位相并折半得$\square$即极差。

设天元一为平勾，加161得$\square$为高股，高股内再加天元得$\square$为极弦。

自乘得$\square$于上位，减去极差冪14,161，

余$\square$为两段极积。应以极弦除，不除寄为分母，便以此为城径。

自乘得$\square$为圆径冪（内有极弦冪分母。寄左）。

然后天元乘高股再四倍得$\square$，又以分母极弦冪$\square$通之，

得$\square$为同数，与左式相消得$\square$。

开平方得64步，即平勾。符合所问。}

\end{rightcolumn}



\end{paracol}



\vspace{0.3cm}

\begin{center}\longline\end{center}



% -------------------------------------------------------------------

% Problem 138

% -------------------------------------------------------------------

\begin{paracol}{2}



\begin{leftcolumn}

\clOWEN{見明和二百七步，和四十六步。問答同前。}



\clFAYUE{%

二和上下相減，數同則止，名為泛率。

又以二和直相減，餘為泛實（此則角差也）。

乃以泛率除泛實，所得為差率也。

以差率加減泛率，若半訖，與勾股相應者，其泛率便為和率，

其泛實便為較率乘和率也。

若不相應，則直取差率以消息之，定為相管和率

（其勾股數少，得見弦黃而相為率者。

勾三股四，則其和七，而其較一也。

勾五股十二，則其和一十七，而其較七也。

勾八股十五，則其和二十三，而其較亦得七也。

勾七股二十四，則其和三十一，而其較一十七也。

勾九股四十，則其和四十九，而其較三十一也。

此消息之大略也。餘皆仿此）。

乃以和率約二和，其明和所得為明壘率，其和所得為壘率也。

又副置和率，上加差率而半之，則為股率也；

下位減差率而半之，則為勾率也。

既見勾、股及差三率，各以壘率乘之，即各得勾、股及差之真數也。}



\clFAYUE{%

又法：二雲數相並，以自乘於上。

二雲數相乘，又四之以減上位為實。

二雲數相並，以六步半乘之於上，又二數相並，以四步半乘之，又四之，

以並入上位為從方。

以七十步零四分三厘七毫五絲為常法。得小差四步。}



\clTSAOYUE{%

以二和相約分得率一、明率四步半，其兩數大小差率並同。

又別得明小差、大差俱為半虛黃也。

立天元一為小差，以四步半乘之得$\square$為大差也，又為明小差，又為半虛黃。

置此大差又以四步半乘之得$\square$為明大差也。

其四差相並得$\square$，減於二和並，得$\square$即兩段太虛大小差並也。

內加三段虛黃方$\square$，得$\square$，合成一個太虛三事和，即圓城徑也。

以自增乘得$\square$為徑冪（寄左）。

乃置和加半虛黃，得$\square$為平勾。

又置明和內加半虛黃，得$\square$為高股，

勾股相乘得下式$\square$，又四之得$\square$為同數，

與左相消得下式$\square$。開平方得四步，即小差也。合問。}

\end{leftcolumn}



\switchcolumn



\begin{rightcolumn}

\mdOWEN{见明和207步，和46步。求城墙直径。}



\mdFAYUE{%

两和上下相减，数相同则止，名为泛率。

又以两和直接相减，余为泛实（此即角差）。

以泛率除泛实，所得为差率。

以差率加减泛率，若折半后与勾股相应，则泛率即为和率，

泛实即为较率乘和率。

若不相应，则直接取差率以调整之，定为相管和率

（勾股数少时，得见弦黄而相互为率。

勾三股四，则和为七，较为一。

勾五股十二，则和为十七，较为七。

勾八股十五，则和为二十三，较为七。

勾七股二十四，则和为三十一，较为十七。

勾九股四十，则和为四十九，较为三十一。

此调整之大致方法。余皆仿此）。

以和率约两和，明和所得为明垒率，和所得为垒率。

又分置和率，上加差率折半则为股率；

下位减差率折半则为勾率。

既得勾、股及差三率，各以垒率乘之，即得勾、股及差之真数。}



\mdFAYUE{%

另法：两已知数相并自乘于上位。

两已知数相乘再四倍减上位为被除数。

两已知数相并以6.5倍乘之于上位，又两数相并以4.5倍乘之再四倍，

合并入上位为一次项。

以70.4375为常数。得小差4步。}



\mdTSAOYUE{%

以两和相约得率一、明率4.5，两数大小差率皆同。

又辨识得明小差、大差皆为半虚黄。

设天元一为小差，以4.5乘之得$\square$为大差，又为明小差，又为半虚黄。

取此大差再以4.5乘之得$\square$为明大差。

四差相并得$\square$，从两和之和中减去得$\square$即两段太虚大小差之和。

内加三段虚黄方$\square$得$\square$，合成一个太虚三事和即圆城径。

自乘得$\square$为径冪（寄左）。

取和加半虚黄得$\square$为平勾。

又取明和内加半虚黄得$\square$为高股，

勾股相乘得下式$\square$，再四倍得$\square$为同数，

与左式相消得下式$\square$。开平方得4步，即小差。符合所问。}

\end{rightcolumn}



\end{paracol}



\vspace{0.3cm}

\begin{center}\longline\end{center}



% -------------------------------------------------------------------

% Problem 139 (last of vol 8)

% -------------------------------------------------------------------

\begin{paracol}{2}



\begin{leftcolumn}

\clOWEN{明勾二股共八十八步，明股二勾共一百六十五步。問答同前。}



\clFAYUE{%

先識別得二大差共、二小差共及四差共。

乃以二大差、二小差相乘為實，以四差共為法。如法得半之虛黃方一十八。}



\clTSAOYUE{%

先置前後雲數以約法約之，得一十一即壘率也。

復各置前後數如壘率而一，前得八即勾率也，後得一十五即股率也。

再以勾、股率求得較率七，和率二十三，弦率一十七，黃方率六，大差率九，小差率二。

既見諸率，各以壘率乘之，其二和共得$\square$，二較共得$\square$，

二弦共得$\square$，二黃共得$\square$，二大差共$\square$，

二小差共$\square$，四差共$\square$。

已上皆為明所得之共數也。

乃立天元一為半虛黃，便為明小差，又為大差也。

以減於大差共，得$\square$即明大差也。

又以減於小差共，得$\square$即小差也。

以二數相增乘，得$\square$（寄左）。

以天元冪與寄左相消，得$\square$。

上法下實，得一十八步，即半之虛黃方也。

以倍之得$\square$，又加於二黃共六十六共得一百二，

即明勾、股共也，又為極黃方，又為虛弦也。

又以三十六減於一百八十七，餘一百五十一，即明股勾共也。

此數內減虛弦，餘$\square$為明二差較也，此名旁差。

以旁差減二弦共一百八十七，餘得$\square$，即太虛和也。

卻加入虛弦一百二，並得$\square$，為太虛三事和，即圓城徑也。合問。}

\end{leftcolumn}



\switchcolumn



\begin{rightcolumn}

\mdOWEN{明勾与二股共88步，明股与二勾共165步。求城墙直径。}



\mdFAYUE{%

先辨识得二大差之和、二小差之和及四差之和。

以二大差、二小差相乘为被除数，以四差之和为除数。依法得半虚黄方18。}



\mdTSAOYUE{%

先取前后两数以约法约之，得11即垒率。

再各置前后数除以垒率，前得8即勾率，后得15即股率。

再以勾、股率求得较率7，和率23，弦率17，黄方率6，大差率9，小差率2。

既得诸率，各以垒率乘之，其二和共得$\square$，二较共得$\square$，

二弦共得$\square$，二黄共得$\square$，二大差共$\square$，

二小差共$\square$，四差共$\square$。

以上皆为明所得之共数。

设天元一为半虚黄，即为明小差，又为大差。

从大差共中减去得$\square$即明大差。

又从小差共中减去得$\square$即小差。

两数相乘得$\square$（寄左）。

以天元冪与寄左相消得$\square$。

上法下实，得18步即半虚黄方。

翻倍得$\square$，再加于二黄共66共得102，

即明勾、股之和，又为极黄方，又为虚弦。

又以36减于187余151即明股勾之和。

此数内减虚弦余$\square$为明二差之较，此名「傍差」。

以傍差减二弦共187余得$\square$即太虚和。

再加入虚弦102共得$\square$为太虚三事和即圆城径。符合所问。}

\end{rightcolumn}



\end{paracol}



\newpage





% ===================================================================

% VOLUME 9

% ===================================================================

\section{\tradfont 卷九\quad 大斜四問、大和八問}

\label{sec:vol9}



\begin{center}

\fcolorbox{classicalborder}{classicalbg}{\parbox{0.9\textwidth}{\centering

\textbf{\large 卷九概要}\quad 本卷十二问，分为大斜四问与大和八问。涉及最复杂的几何构形与最高次的方程——五次方程（四乘方）、六次方程（五乘方），以及玲珑（精巧方程变换）、换骨（根本变换）、连枝同体（关联方程组）等高阶技巧。

}}

\end{center}



\vspace{0.5cm}



\subsection{\tradfont 細草卷九上\quad 大斜四問}



\begin{center}

\fcolorbox{classicalborder}{classicalbg}{\parbox{0.85\textwidth}{\centering

\textbf{大斜四問}：涉及大斜（通弦）及其相关各形的复杂问题，需运用\textbf{连枝同体}（关联方程组）与\textbf{玲珑}（精巧方程变换）求解。

}}

\end{center}



\vspace{0.3cm}



% -------------------------------------------------------------------

% Problem 140

% -------------------------------------------------------------------

\begin{paracol}{2}



\begin{leftcolumn}

\clOWEN{%

甲丙俱在中心，丙望南門直行，不知步數而止。

甲出東門直行，不知步數望見丙，斜行與丙相會。

二人共行了六百八十步，仍雲甲直行少於丙直行一百一十九步。問答同前。}



\clFAYUE{%

二數相減，餘以為冪，內卻減差冪為平實；

二數相減又四之於上，又加入二之差步為益從；二步常法。得皇勾一百三十六。}



\clTSAOYUE{%

別得共步即皇極三事和，少步即勾股差也。

立天元一為皇極勾。加少步得$\square$為股也，又以天元加股得$\square$為和也。

以和減共步得$\square$為弦也，弦自之得$\square$為一段弦冪（寄左）。

然後置股以天元乘之，又倍之，得$\square$為二直積，

如入少步冪$\square$共得$\square$為同數，與左相消得$\square$。

平方而一，得一百三十六即勾也。勾加差為股，勾股相乘，倍之為實，

勾股和減共步為法。得城徑。}

\end{leftcolumn}



\switchcolumn



\begin{rightcolumn}

\mdOWEN{%

甲丙都在中心，丙望南门直行，不知步数停下。

甲出东门直行，不知步数望见丙，斜行与丙相会。

两人共行了680步，又说甲直行比丙直行少119步。求城墙直径。}



\mdFAYUE{%

两数相减余数作为冪，减去差的冪为二次方程常数项；

两数相减再四倍于上位，再加入两倍差步为负一次项；2为常数。得皇勾136。}



\mdTSAOYUE{%

辨识得：共步即皇极三事和，少步即勾股差。

设天元一为皇极勾。加少步得$\square$为股，又以天元加股得$\square$为和。

以和减共步得$\square$为弦，弦自乘得$\square$为一段弦冪（寄左）。

然后取股以天元乘之再两倍得$\square$为两倍直积，

加入少步冪$\square$共得$\square$为同数，与左式相消得$\square$。

开平方得136即勾。勾加差为股，勾股相乘两倍为被除数，

勾股和减共步为除数。得城径。}

\end{rightcolumn}



\end{paracol}



\vspace{0.3cm}

\begin{center}\longline\end{center}



% -------------------------------------------------------------------

% Problem 141

% -------------------------------------------------------------------

\begin{paracol}{2}



\begin{leftcolumn}

\clOWEN{%

甲丙俱在西北隅起，丙向南行不知步數而立。

甲向東行望見丙，就丙斜行六百八十步與丙相會。

丙雲：「我南行步多於甲東行二百八十步。」問答同前。}



\clFAYUE{以雲數差乘雲數並為實，倍多步為從，二為平隅。得大勾三百二十。}



\clTSAOYUE{%

立天元為大勾。加差得$\square$為股，倍天元乘之，得$\square$為二積（寄左）。

然後以斜步、多步並$\square$與斜步、多步較$\square$相乘，

得$\square$為同數，與左相消得$\square$。

開平方得三百二十步，即大勾也。合問。}

\end{leftcolumn}



\switchcolumn



\begin{rightcolumn}

\mdOWEN{%

甲丙都在西北角出发，丙向南行不知步数停下。

甲向东行望见丙，向丙斜行680步与丙相会。

丙说：「我南行步比甲东行多280步。」求城墙直径。}



\mdFAYUE{以两已知数之差乘两已知数之和为被除数，两倍多步为一次项，2为二次项系数。得大勾320。}



\mdTSAOYUE{%

设天元为大勾。加差得$\square$为股，两倍天元乘之得$\square$为两倍积（寄左）。

然后斜步与多步之和$\square$与斜步和多步之差$\square$相乘，

得$\square$为同数，与左式相消得$\square$。

开平方得320步，即大勾。符合所问。}

\end{rightcolumn}



\end{paracol}



\vspace{0.3cm}

\begin{center}\longline\end{center}



% -------------------------------------------------------------------

% Problem 142

% -------------------------------------------------------------------

\begin{paracol}{2}



\begin{leftcolumn}

\clOWEN{%

甲乙二人共立於艮隅，乙南行過城門而立，甲東行望乙與城參相直而止。

丙丁二人共立於坤隅，丁向東行過城門而立，丙向南行望丁及甲乙悉與城俱相直。

丙復就甲斜行六百八十步與甲相會。

乙、丁又雲：「吾二人直行共得三百四十二步。」問答同前。}



\clFAYUE{%

二雲數相乘，倍之為實；倍斜行於上，以二雲數相減加上位為從；一步常法。

開平方，得城徑。}



\clTSAOYUE{%

別得斜步即大弦也。其共步則一徑一虛弦共也，

其二數相並為一大和一虛弦共數也。

立天元為徑，減於共步得$\square$為虛弦也。

以虛弦復減於天元，得$\square$為虛和，以斜步乘之，得$\square$（寄左）。

乃以天元加斜步，得$\square$為大和，以虛弦乘之得$\square$為同數，

與左相消得$\square$。開平方得二百四十步，即城徑也。合問。}

\end{leftcolumn}



\switchcolumn



\begin{rightcolumn}

\mdOWEN{%

甲乙二人同立于艮角（东北），乙南行过城门停下，甲东行望乙与城墙呈直角而止。

丙丁二人同立于坤角（西南），丁东行过城门停下，丙南行望丁及甲乙都与城墙呈直角。

丙再向甲斜行680步与甲相会。

乙、丁又说：「我们二人直行共得342步。」求城墙直径。}



\mdFAYUE{%

两已知数相乘翻倍为被除数；两倍斜行于上位，以两已知数相减加上位为一次项；1为常数。

开平方得城径。}



\mdTSAOYUE{%

辨识得：斜步即大弦。共步即一径一虚弦之和，

两数相并为一大和一虚弦之共数。

设天元为径，从共步中减去得$\square$为虚弦。

以虚弦再从天元中减去得$\square$为虚和，以斜步乘之得$\square$（寄左）。

以天元加斜步得$\square$为大和，以虚弦乘之得$\square$为同数，

与左式相消得$\square$。开平方得240步，即城径。符合所问。}

\end{rightcolumn}



\end{paracol}



\vspace{0.3cm}

\begin{center}\longline\end{center}



% -------------------------------------------------------------------

% Problem 143

% -------------------------------------------------------------------

\begin{paracol}{2}



\begin{leftcolumn}

\clOWEN{%

甲從北門向東直行，庚從西門穿城東行，丙從西門向南直行，壬從北門穿城南行。

四人遙相望，悉與城參相直。

只雲甲丙相望處六百八十步，庚壬穿城共行了六百三十一步。問答同前。}



\clFAYUE{%

共步自之得數。以共步減斜，餘自乘以減上為實。

二之斜步加入共步減斜，餘數為從。一步常法。得城徑。}



\clTSAOYUE{%

共行步為一徑與皇和共也，又為大和皇弦差也。甲丙相望即大弦也。

以共步減大弦，餘$\square$為皇極弦上減一徑也。

立天元一為圓徑，減於共步，得$\square$為皇極和也，以自之得$\square$於上。

弦內減共步餘$\square$，又以天元加之為皇弦，以自之得$\square$，

減上位，餘得$\square$為兩個皇直積（寄左）。

乃以天元乘皇弦得下式$\square$為同數，與左相消得$\square$。

平方而一，得二百四十步，即城徑也。合問。}

\end{leftcolumn}



\switchcolumn



\begin{rightcolumn}

\mdOWEN{%

甲从北门向东直行，庚从西门穿城东行，丙从西门向南直行，壬从北门穿城南行。

四人遥遥相望，都与城墙呈直角。

只说甲丙相望处680步，庚壬穿城共行了631步。求城墙直径。}



\mdFAYUE{%

共步自乘得数。以共步减斜步，余数自乘减上位为被除数。

两倍斜步加入共步减斜之余数为一次项。1为常数。得城径。}



\mdTSAOYUE{%

共行步为一径与皇极之和，又为大和与皇弦之差。甲丙相望即大弦。

以共步减大弦余$\square$为皇极弦上减一径。

设天元一为圆径，从共步中减去得$\square$为皇极和，自乘得$\square$于上位。

弦内减共步余$\square$，再加天元为皇弦，自乘得$\square$，

减上位余得$\square$为两个皇极直积（寄左）。

以天元乘皇弦得下式$\square$为同数，与左式相消得$\square$。

开平方得240步，即城径。符合所问。}

\end{rightcolumn}



\end{paracol}



\newpage



\subsection{\tradfont 細草卷九下\quad 大和八問}



\begin{center}

\fcolorbox{classicalborder}{classicalbg}{\parbox{0.85\textwidth}{\centering

\textbf{大和八問}：以「大和」（通勾通股之和）為已知條件，涉及\textbf{換骨}（根本變換）與\textbf{奪項}（消元法）等高階代數技巧，方程多達五次（四乘方）、六次（五乘方）。

}}

\end{center}



\vspace{0.3cm}



% -------------------------------------------------------------------

% Problem 144

% -------------------------------------------------------------------

\begin{paracol}{2}



\begin{leftcolumn}

\clOWEN{%

庚從西門穿城東行二百五十六步而立，壬從北門穿城南行三百七十五步而立。

又有甲丙二人俱在乾隅，甲向東行，丙向南行，各不知步數而立。

四人遙相望，只雲甲丙共行了九百二十步。問答同前。}



\clFAYUE{%

庚東行冪、壬南行冪相並於上，並庚壬步而倍之，內減大和，

餘復減於庚壬共，得數以自乘減上位為平實。

並庚壬步為益從，半步為隅法。得城徑。}



\clTSAOYUE{%

立天元一為圓徑，以半之副置二位。

上以減於庚東行，得下$\square$為平弦也；下以減於壬南行，得$\square$為高弦也。

二弦相並得$\square$為皇弦、虛弦共也。

倍此數得$\square$為大弦、虛弦共也。

以大弦、虛弦共減於大和，餘$\square$為虛勾虛股共也。

天元內減虛勾、虛股共，餘$\square$即虛弦也。

復置皇弦虛弦共內減虛弦，餘$\square$即皇極弦也，

以自之得$\square$（寄左）。

然後以平弦自之，得下式$\square$為勾冪也，又以高弦自之，

得$\square$為股冪也。

二冪相並得$\square$為同數，與左相消得$\square$。

平方而一，得二百四十步，即城徑也。合問。}

\end{leftcolumn}



\switchcolumn



\begin{rightcolumn}

\mdOWEN{%

庚从西门穿城东行256步停下，壬从北门穿城南行375步停下。

又有甲丙二人同在乾角（西北），甲向东行，丙向南行，各不知步数停下。

四人遥遥相望，只说甲丙共行了920步。求城墙直径。}



\mdFAYUE{%

庚东行冪、壬南行冪合并于上位，庚壬步数相加翻倍，从中减去大和，

余数再从庚壬共中减去，得数自乘减上位为二次方程常数项。

庚壬步数之和为负一次项，0.5为二次项系数。得城径。}



\mdTSAOYUE{%

设天元一为圆径，取半后分置两位。

上位从庚东行中减去得下式$\square$为平弦；下位从壬南行中减去得$\square$为高弦。

两弦相加得$\square$为皇弦与虚弦之和。

此数翻倍得$\square$为大弦与虚弦之和。

以大弦、虚弦之和减于大和余$\square$为虚勾与虚股之和。

天元内减虚勾、虚股之和余$\square$即虚弦。

再取皇弦虚弦之和减去虚弦余$\square$即皇极弦，

自乘得$\square$（寄左）。

然后平弦自乘得下式$\square$为勾冪，高弦自乘

得$\square$为股冪。

两冪合并得$\square$为同数，与左式相消得$\square$。

开平方得240步，即城径。符合所问。}

\end{rightcolumn}



\end{paracol}



\vspace{0.3cm}

\begin{center}\longline\end{center}



% -------------------------------------------------------------------

% Problem 145

% -------------------------------------------------------------------

\begin{paracol}{2}



\begin{leftcolumn}

\clOWEN{%

甲丙俱在西北隅，甲向東行不知步數而立，丙向南行望見甲，就甲斜行與甲相會。

甲直行、丙直行共九百二十步（甲步少於丙步）。

又：出東門南行有柳樹一株，出南門東行有槐樹一株。

戊己二人同在巽隅，戊就柳樹，己就槐樹，亦與甲、乙遙相望。

只雲己行少於戊行數與兩樹相距數相並，得一百四十四步，其二數相減餘六十步。問答同前。}



\clFAYUE{%

二雲數相並而半之為虛弦，以乘大和九百二十步於上。

以一百四十四減大和，以虛較乘之，減上位為平實。

以一百四十四減大和，又二之於上，以二之虛較減上位為從。四虛隅。得太虛勾四十八。}



\clTSAOYUE{%

別得甲丙直行共即大和也，戊就柳樹步即虛股也，己就槐樹步即虛勾也。

其一百四十四步即二明勾，其六十步即二股也。

立天元一為虛勾，加明勾得$\square$為半徑也，

倍之得$\square$即城徑也（又為虛弦上三事和）。

二雲數相並而半之，得$\square$即小弦也。

相減而半之，得$\square$即小較也。

以天元加較得$\square$即小股也，小勾股共得$\square$即小和也。

以小三事減大和得$\square$即大弦也。

乃先置小和以大弦乘之，得下式$\square$（寄左）。

次以小弦乘大和，得$\square$，與左相消得下式$\square$。

開平方得四十八步，即虛勾也。加明勾又倍之得二百四十步，即城徑也。合問。}

\end{leftcolumn}



\switchcolumn



\begin{rightcolumn}

\mdOWEN{%

甲丙都在西北角，甲向东行不知步数停下，丙向南行望见甲，向甲斜行与甲相会。

甲直行、丙直行共920步（甲步少于丙步）。

又：出东门南行有柳树一株，出南门东行有槐树一株。

戊己二人同在巽角（东南），戊到柳树，己到槐树，也与甲、乙遥遥相望。

只说己行比戊行之数与两树相距之数相加得144步，两数相减余60步。求城墙直径。}



\mdFAYUE{%

两已知数相并折半为虚弦，乘以大和920步于上位。

以144减大和，以虚较乘之减上位为二次方程常数项。

以144减大和再两倍于上位，以两倍虚较减上位为一次项。负4为二次项系数。得太虚勾48。}



\mdTSAOYUE{%

辨识得：甲丙直行共即大和，戊到柳树步即虚股，己到槐树步即虚勾。

144步即两倍明勾，60步即两倍股。

设天元一为虚勾，加明勾得$\square$为半径，

翻倍得$\square$即城径（又为虚弦上三事和）。

两已知数相并折半得$\square$即小弦。

相减折半得$\square$即小较。

天元加较得$\square$即小股，小勾股共得$\square$即小和。

以小三事减大和得$\square$即大弦。

先置小和以大弦乘之得下式$\square$（寄左）。

次以小弦乘大和得$\square$，与左式相消得下式$\square$。

开平方得48步即虚勾。加明勾再翻倍得240步即城径。符合所问。}

\end{rightcolumn}



\end{paracol}



\vspace{0.3cm}

\begin{center}\longline\end{center}



% -------------------------------------------------------------------

% Problem 146

% -------------------------------------------------------------------

\begin{paracol}{2}



\begin{leftcolumn}

\clOWEN{%

甲從乾隅東行，乙從艮隅南行，丙從乾隅南行，丁從坤隅東行，四人遙相望見。

既而甲還至艮隅就乙，丙還至坤隅就丁。

甲丙直行共九百二十步，甲還就乙共二百三十步，丙還就丁共五百五十二步。問答同前。}



\clFAYUE{%

並就數以減直行共，復以所並就數乘之為實；

並就數減直行共，得數復加入直行共為法。得虛弦。}



\clTSAOYUE{%

別得甲丙直行共為大和也，甲還就乙步為小差勾股共也，丙還就丁步為大差勾股共也。

以大差勾股共減於大股，餘即虛勾也。

以小差勾股共減於大勾，餘即虛股也。

二數相並得$\square$為大弦虛弦共也，

二數相減，餘$\square$為通差及太虛勾股差共也。

又並二數而半之，得$\square$為皇極弦虛弦共，又為皇極勾股共也。

立天元一為虛弦。

先以二共數減於大和，餘$\square$為虛勾虛股和於上。

次以虛弦減於二共數，餘$\square$為大弦，

以乘上位，得下$\square$（寄左）。

然後以天元乘大和，得$\square$元為同數，與左相消得$\square$。

上法下實，得一百二步，即虛弦也。加入虛和得二百四十步，即城徑也。合問。}

\end{leftcolumn}



\switchcolumn



\begin{rightcolumn}

\mdOWEN{%

甲从乾角（西北）东行，乙从艮角（东北）南行，丙从乾角南行，丁从坤角（西南）东行，四人遥遥相望。

随后甲回到艮角靠近乙，丙回到坤角靠近丁。

甲丙直行共920步，甲靠近乙共230步，丙靠近丁共552步。求城墙直径。}



\mdFAYUE{%

并就数减直行共，再以并就数乘之为被除数；

并就数减直行共，得数再加入直行共为除数。得虚弦。}



\mdTSAOYUE{%

辨识得：甲丙直行共为大和，甲靠近乙步为小差勾股之和，丙靠近丁步为大差勾股之和。

以大差勾股之和减于大股余即虚勾。

以小差勾股之和减于大勾余即虚股。

两数相并得$\square$为大弦虚弦之和，

两数相减余$\square$为通差及太虚勾股差之和。

又两数相并折半得$\square$为皇极弦虚弦之和，又为皇极勾股之和。

设天元一为虚弦。

先以二共数减于大和余$\square$为虚勾虚股之和于上位。

次以虚弦减于二共数余$\square$为大弦，

乘上位得下式$\square$（寄左）。

然后天元乘大和得$\square$元为同数，与左式相消得$\square$。

上法下实，得102步即虚弦。加入虚和得240步即城径。符合所问。}

\end{rightcolumn}



\end{paracol}



\vspace{0.3cm}

\begin{center}\longline\end{center}



% -------------------------------------------------------------------

% Problem 147

% -------------------------------------------------------------------

\begin{paracol}{2}



\begin{leftcolumn}

\clOWEN{%

依前見大和，只雲股圓差上勾弦差二百一十六，勾圓差上股弦差二十步。問答同前。}



\clFAYUE{%

以雲數二十步減通和，復以二十步乘之於上。

以雲數二百一十六減九百步而半之，乘上位為立實。

三因二十步以減通和得八百六十，

以二百一十六及二十共得二百三十六，減通和而半之，得三百四十二。

二數相乘訖，內減二十之九百步。

又以三百四十二及二百一十六共得五百五十八，又二十之以減之為從方。

以二百三十六減通和，又以三之二十步減通和，相並於上。

以二之五百五十八內卻減二十步，餘以加上位為益廉。

四步常法。得小差股一百五十。}



\clTSAOYUE{%

別得小差上股弦差$\square$加二股為大勾也，大差上勾弦差$\square$加二勾為大股也。

立天元一為小差股，加$\square$得$\square$為小差弦也，

小差弦上又加天元得$\square$為通勾，以減於和步得$\square$為通股也。

通股內減$\square$得$\square$，半之得下式$\square$即大差之勾也。

大差勾上又加$\square$，得$\square$為大差弦也。

再置通股以小差弦乘之，得$\square$，以天元除之，得$\square$為一個大弦也（泛寄）。

再置通勾以大差弦乘之，得$\square$，合以大差勾除。

不除寄為母，便以為大弦（寄左）。

乃以大差勾乘泛寄，得$\square$為同數，與左相消得$\square$。

益積開立方，得一百五十步，為小差股也。合問。}

\end{leftcolumn}



\switchcolumn



\begin{rightcolumn}

\mdOWEN{%

如前见大和（通勾通股之和），只说股圆差上勾弦差216，勾圆差上股弦差20步。求城墙直径。}



\mdFAYUE{%

以已知数20步减通和，再以20步乘之于上位。

以已知数216减900步折半，乘上位为三次方程（立方）常数项。

三倍20步减通和得860，

216与20共得236，减通和折半得342。

两数相乘毕，减去20乘900。

又以342及216共得558，再乘20减之为一次项。

以236减通和，又以三倍20步减通和，相并于上位。

以两倍558减去20步，余加上位为负廉。

4为常数。得小差股150。}



\mdTSAOYUE{%

辨识得：小差上股弦差$\square$加两倍股为大勾，大差上勾弦差$\square$加两倍勾为大股。

设天元一为小差股，加$\square$得$\square$为小差弦，

小差弦上又加天元得$\square$为通勾，从和步中减去得$\square$为通股。

通股内减$\square$得$\square$，折半得下式$\square$即大差之勾。

大差勾上加$\square$得$\square$为大差弦。

再取通股以小差弦乘之得$\square$，除以天元得$\square$为一个大弦（泛寄）。

再取通勾以大差弦乘之得$\square$，应以大差勾除。

不除寄为分母，便以为大弦（寄左）。

以大差勾乘泛寄得$\square$为同数，与左式相消得$\square$。

益积开立方得150步，即为小差股。符合所问。}

\end{rightcolumn}



\end{paracol}



\vspace{0.3cm}

\begin{center}\longline\end{center}



% -------------------------------------------------------------------

% Problem 148

% -------------------------------------------------------------------

\begin{paracol}{2}



\begin{leftcolumn}

\clOWEN{%

依前見大和，只雲高弦、平弦共得三百九十一步，高弦、平弦相較得一百一十九步。問答同前。}



\clFAYUE{%

以較數冪減於共數冪，又半之為實，以共數減大和為益從，一常法。

開平方，得圓徑。}



\clTSAOYUE{%

別得高弦減於通股為邊股內減明股也，平弦減於通勾為邊勾內減明勾也。

其共數即大弦內減皇極弦，又為皇極勾股共也，其相較步即皇極差也。

二雲數相並得$\square$，即黃廣弦也。

二雲數相減，餘即黃長弦也。

以共數減於大和，餘$\square$為皇極弦、圓徑共。

立天元一為圓徑，以減於$\square$為皇極弦也。

以共數自之得$\square$於上。

以相較數自之得$\square$，減上位，餘$\square$，

又半之得$\square$為兩段皇極積（寄左）。

乃以天元乘皇極弦，得$\square$為同數，

與左相消得下$\square$。開平方得二百四十步，即城徑也。合問。}

\end{leftcolumn}



\switchcolumn



\begin{rightcolumn}

\mdOWEN{%

如前见大和，只说高弦、平弦共得391步，高弦、平弦相差119步。求城墙直径。}



\mdFAYUE{%

以较数冪从共数冪中减去，折半为被除数，以共数减大和为负一次项，1为常数。

开平方得圆径。}



\mdTSAOYUE{%

辨识得：高弦减于通股为边股内减明股，平弦减于通勾为边勾内减明勾。

共数即大弦内减皇极弦，又为皇极勾股之和，相差步即皇极差。

两已知数相并得$\square$即黄广弦。

两已知数相减余即黄长弦。

以共数减于大和余$\square$为皇极弦、圆径之和。

设天元一为圆径，从$\square$中减去为皇极弦。

以共数自乘得$\square$于上位。

以相较数自乘得$\square$，减上位余$\square$，

折半得$\square$为两段皇极积（寄左）。

以天元乘皇极弦得$\square$为同数，

与左式相消得下式$\square$。开平方得240步即城径。符合所问。}

\end{rightcolumn}



\end{paracol}



\vspace{0.3cm}

\begin{center}\longline\end{center}



% -------------------------------------------------------------------

% Problem 149

% -------------------------------------------------------------------

\begin{paracol}{2}



\begin{leftcolumn}

\clOWEN{%

依前見大和，只雲大差弦四百零八步，小差弦一百七十步。問答同前。}



\clFAYUE{%

以並雲數減大和，復以相並數乘之為實。

三之通和於上，又以並雲數減大和，加上位為從。二步虛法。得圓徑。}



\clTSAOYUE{%

大差弦減和步，餘$\square$為大勾、大差勾共也。

以小差弦減大和，餘$\square$為大股、小差股共也。

雲數相並得即大弦內減虛弦也，雲數相減得$\square$為虛弦平弦共也。

以相並數減於大和，餘$\square$為大差勾、小差股共，又為圓徑、虛弦共也。

立天元一為圓徑，減於$\square$得$\square$為虛弦也。

反以減於圓徑得$\square$為小和也。

以天元減大和得$\square$為大弦，以乘小和，得$\square$（寄左）。

乃再置虛弦以通和乘之，得$\square$，與左相消得$\square$。

開平方得二百四十步，即城徑也。合問。}

\end{leftcolumn}



\switchcolumn



\begin{rightcolumn}

\mdOWEN{%

如前见大和，只说大差弦408步，小差弦170步。求城墙直径。}



\mdFAYUE{%

以两已知数之和减大和，再以相并数乘之为被除数。

三倍通和于上位，又以并数减大和加上位为一次项。负2为常数。得圆径。}



\mdTSAOYUE{%

大差弦减和步余$\square$为大勾、大差勾之和。

以小差弦减大和余$\square$为大股、小差股之和。

两数相并得即大弦内减虚弦，相减得$\square$为虚弦平弦之和。

以相并数减于大和余$\square$为大差勾、小差股之和，又为圆径、虚弦之和。

设天元一为圆径，从$\square$中减去得$\square$为虚弦。

反过来从圆径中减去得$\square$为小和。

天元减大和得$\square$为大弦，乘小和得$\square$（寄左）。

再取虚弦以通和乘之得$\square$，与左式相消得$\square$。

开平方得240步即城径。符合所问。}

\end{rightcolumn}



\end{paracol}



\vspace{0.3cm}

\begin{center}\longline\end{center}



% -------------------------------------------------------------------

% Problem 150

% -------------------------------------------------------------------

\begin{paracol}{2}



\begin{leftcolumn}

\clOWEN{%

依前見大和，只雲黃廣弦五百一十步，黃長弦二百七十二步。問答同前。}



\clFAYUE{%

雲數相並減大和，復以相並數乘之為實。

雲數相並減大和，得數復以加大和為法。得虛弦一百二。}



\clTSAOYUE{%

別得黃廣弦又為大差弦、虛弦共，又為邊股、股共也。

黃長弦又為小差弦、虛弦共，又為底勾、明勾共也。

以黃廣弦減於大股，餘$\square$即虛股，以黃長弦減於大勾，餘$\square$即虛勾。

故並數以減於大和，餘$\square$為虛和也。

以虛和減徑即虛弦也。

二雲數相並得$\square$為大弦、虛弦共也，雲數相減餘$\square$為虛弦、平弦共。

立天元一為虛弦，以減於七百八十二，得$\square$為大弦也，

以小和乘之得$\square$（寄左）。

乃以天元虛弦乘大和，得$\square$為同數，與左相消得$\square$。

上法下實，得一百二步，即虛弦也。合問。}

\end{leftcolumn}



\switchcolumn



\begin{rightcolumn}

\mdOWEN{%

如前见大和，只说黄广弦510步，黄长弦272步。求城墙直径。}



\mdFAYUE{%

两已知数相并减大和，再以相并数乘之为被除数。

两数相并减大和，得数再加入大和为除数。得虚弦102。}



\mdTSAOYUE{%

辨识得：黄广弦又为大差弦、虚弦之和，又为边股、股之和。

黄长弦又为小差弦、虚弦之和，又为底勾、明勾之和。

以黄广弦减于大股余$\square$即虚股，以黄长弦减于大勾余$\square$即虚勾。

故两数相并以减于大和余$\square$为虚和。

以虚和减径即虚弦。

两已知数相并得$\square$为大弦、虚弦之和，相减余$\square$为虚弦、平弦之和。

设天元一为虚弦，从782中减去得$\square$为大弦，

乘小和得$\square$（寄左）。

以天元虚弦乘大和得$\square$为同数，与左式相消得$\square$。

上法下实，得102步即虚弦。符合所问。}

\end{rightcolumn}



\end{paracol}



\vspace{0.3cm}

\begin{center}\longline\end{center}



% -------------------------------------------------------------------

% Problem 151

% -------------------------------------------------------------------

\begin{paracol}{2}



\begin{leftcolumn}

\clOWEN{%

依前見大和，只雲邊弦五百四十四步，底弦四百二十五步。問答同前。}



\clFAYUE{雲數相減自之為實，以大和減並數為法。得皇極弦。}



\clTSAOYUE{%

別得以邊弦減大股，餘$\square$為半徑內減平勾，又為平弦內減勾圓差也。

以大勾減於底弦，餘$\square$為高股內少半徑，又為股圓差內少高弦也。

二雲數相並，得九百六十九為大弦、皇極弦共也。

二雲數相減，得$\square$為皇極勾股差也。

並數內減通和，餘$\square$為皇極弦內減圓徑也。

立天元一為皇極弦，以自之於上。

以一百一十九自之得$\square$，減上位得$\square$為二皇積（寄左）。

復置天元內減四十九得下式$\square$為黃方。

復以天元乘之，得$\square$，與左相消得$\square$。

上法下實，得二百八十九步，即皇極弦也。內減四十九，餘即城徑也。合問。}

\end{leftcolumn}



\switchcolumn



\begin{rightcolumn}

\mdOWEN{%

如前见大和，只说边弦544步，底弦425步。求城墙直径。}



\mdFAYUE{两已知数相减自乘为被除数，以大和减并数为除数。得皇极弦。}



\mdTSAOYUE{%

辨识得：以边弦减大股余$\square$为半径内减平勾，又为平弦内减勾圆差。

以大勾减于底弦余$\square$为高股内少半径，又为股圆差内少高弦。

两已知数相并得969为大弦、皇极弦之和。

两已知数相减得$\square$为皇极勾股之差。

并数内减通和余$\square$为皇极弦内减圆径。

设天元一为皇极弦，自乘于上位。

以119自乘得$\square$，减上位得$\square$为两倍皇极积（寄左）。

再取天元内减49得下式$\square$为黄方。

再以天元乘之得$\square$，与左式相消得$\square$。

上法下实，得289步即皇极弦。内减49余即城径。符合所问。}

\end{rightcolumn}



\end{paracol}



\newpage



% ===================================================================

% COLOPHON / 跋

% ===================================================================

\section*{跋\quad Colophon}

\addcontentsline{toc}{section}{跋 Colophon}



\begin{paracol}{2}



\begin{leftcolumn}

\fcolorbox{classicalborder}{classicalbg}{%

\parbox{\linewidth}{%

\textbf{\large\tradfont 【文言文】}



\vspace{0.3em}

《測圓海鏡》十二卷，金元李冶撰，成書於1248年。此書為現存最早系統運用「天元術」

求解多項式方程之數學專著，堪稱中世紀代數學之巔峰。



全書以圓城與勾股形為基本構架，設十五種直角三角形，各邊對角皆有專名。

以天元一為未知數，借算籌佈列太極、天元之位，建立高次方程，

用增乘開方法求解。其方程式之高次，可達\textbf{五乘方}（六次方程）。



書中尤多精妙之處：

\begin{itemize}[leftmargin=1.5em,topsep=0pt]

  \item \textbf{連枝同体}：关联方程组之法

  \item \textbf{玲珑}：精巧的方程变换之术

  \item \textbf{換骨}：根本变换之法

  \item \textbf{奪項}：消元之法

\end{itemize}



此皆中算之瑰寶，值得後人深入研習。



\vspace{0.3em}

\textit{測圓海鏡 卷七至九 文言文/白话文对照版}

}}

\end{leftcolumn}



\switchcolumn



\begin{rightcolumn}

\fcolorbox{modernborder}{modernbg}{%

\parbox{\linewidth}{%

\textbf{\large\modfont 【白话文】}



\vspace{0.3em}

《测圆海镜》十二卷，金元时期李冶所著，成书于1248年。此书为现存最早系统运用「天元术」

求解多项式方程的数学专著，堪称中世纪代数学的巅峰之作。



全书以圆城与直角三角形为基本构架，设立十五种直角三角形，各边对角皆有专名。

以天元一为未知数，借算筹排列太极、天元之位，建立高次方程，

用增乘开方法求解。其方程的最高次数可达\textbf{六次方程（五乘方）}。



书中尤为精妙之处：

\begin{itemize}[leftmargin=1.5em,topsep=0pt]

  \item \textbf{连枝同体（关联方程组）}：处理多变量联立方程的方法

  \item \textbf{玲珑（精巧的方程变换）}：巧妙的代数变形技巧

  \item \textbf{换骨（根本变换）}：从根本上转换问题视角的技巧

  \item \textbf{夺项（消元法）}：消去变量的方法

\end{itemize}



这些都是中国数学的瑰宝，值得后人深入研习。



\vspace{0.3em}

\textit{测圆海镜 卷七至九 文言文白话文对照版}

}}

\end{rightcolumn}



\end{paracol}



\vspace{1cm}



\begin{center}

\longline



\vspace{0.5cm}



{\small\tradfont 原文來源：Chinese Wikisource（\url{https://zh.wikisource.org/wiki/測圓海鏡}）}



\vspace{0.2cm}



{\small\modfont 排版：Xe\LaTeX\ + xeCJK + Microsoft YaHei 字体族}



\vspace{0.2cm}



{\small 測圓海鏡 卷七至卷九 \quad 文言文\,$\boldsymbol{\leftrightarrow}$\,白话文 對照版}



{\small \textit{Bilingual Edition: Classical Chinese / Modern Chinese}}



\vspace{0.3cm}



{\footnotesize 術語對照：五乘方$\to$五次方程\quad 六乘方$\to$六次方程\quad 連枝同体$\to$連枝同体（关联方程组）\quad 玲珑$\to$玲珑（精巧的方程变换）\quad 換骨$\to$換骨（根本变换）\quad 奪項$\to$奪項（消元法）}

\end{center}



\end{document}







% ============================================================

% Source: translations/zh/chinese/qin-jiushao-shuxue-jiuzhang-fascicle1_classical-modern_bilingual.tex

% ============================================================



%% ========================================================================

%% Qin Jiushao's Shuxue Jiuzhang (Mathematical Treatise in Nine Sections)

%% Fascicle 1: Dayan Category (大衍類)

%% BILINGUAL EDITION: Classical Chinese (LEFT) -- Modern Chinese (RIGHT)

%% ========================================================================

\documentclass[11pt,a4paper]{article}



%% -------- Core Packages --------

\usepackage{fontspec}

\usepackage{xeCJK}       % Chinese line breaking support

\usepackage{amsmath,amssymb,amsthm}

\usepackage{geometry}

\usepackage{graphicx}

\usepackage{xcolor}

\usepackage{fancyhdr}

\usepackage{array,booktabs,longtable}

\usepackage{hyperref}

\usepackage{indentfirst}

\usepackage{enumitem}



%% -------- Page Geometry --------

\geometry{

  paperwidth=210mm,paperheight=297mm,

  left=16mm,right=16mm,top=24mm,bottom=24mm,

  headheight=14pt,footskip=30pt

}



%% -------- Font Configuration --------

\setmainfont{Microsoft YaHei}

\setsansfont{Microsoft YaHei}

\setmonofont{Microsoft YaHei}

\setCJKmainfont{Microsoft YaHei}



%% Chinese font settings

\newfontfamily\chinese{Microsoft YaHei}

\newfontfamily\chinesebf{Microsoft YaHei}

\newfontfamily\chineselg{Microsoft YaHei}



%% Chinese text line breaking settings

\xeCJKsetup{PunctStyle=kaiming, CJKmath=true}

\tolerance=2000

\emergencystretch=3em

\hbadness=10000

\sloppy



%% -------- Colors --------

\definecolor{titlecolor}{RGB}{45,55,72}

\definecolor{sectioncolor}{RGB}{55,65,81}

\definecolor{notecolor}{RGB}{120,53,15}

\definecolor{rulecolor}{RGB}{180,160,140}

\definecolor{prefacecolor}{RGB}{80,60,40}

\definecolor{labelclassical}{RGB}{139,69,19}

\definecolor{labelmodern}{RGB}{25,100,120}



%% -------- Section Formatting --------

\makeatletter

\renewcommand{\section}{%

  \@startsection{section}{1}{0mm}%

  {1.5ex plus 0.5ex minus 0.2ex}%

  {0.7ex plus 0.2ex}%

  {\normalfont\Large\bfseries\color{sectioncolor}}}

\renewcommand{\subsection}{%

  \@startsection{subsection}{2}{0mm}%

  {1.2ex plus 0.3ex minus 0.1ex}%

  {0.5ex plus 0.1ex}%

  {\normalfont\large\bfseries\color{sectioncolor}}}

\makeatother



%% -------- Header/Footer --------

\pagestyle{fancy}

\fancyhf{}

\fancyhead[L]{\small\textit{数书九章 -- 文言文|白话文对照版}}

\fancyhead[R]{\small\thepage}

\renewcommand{\headrulewidth}{0.4pt}

\renewcommand{\footrulewidth}{0pt}



%% -------- Bilingual Layout using minipage --------

\newcommand{\bilingualsection}[2]{%

  \noindent\begin{minipage}[t]{0.48\textwidth}

  \vspace{-0.3em}{\color{labelclassical}\rule{\linewidth}{0.6pt}}\\

  {\small\bfseries\color{labelclassical}【文言文原文】}\\

  {\color{labelclassical}\rule{\linewidth}{0.3pt}}\\\vspace{0.3em}

  #1

  \end{minipage}%

  \hfill%

  \begin{minipage}[t]{0.48\textwidth}

  \vspace{-0.3em}{\color{labelmodern}\rule{\linewidth}{0.6pt}}\\

  {\small\bfseries\color{labelmodern}【白话文今译】}\\

  {\color{labelmodern}\rule{\linewidth}{0.3pt}}\\\vspace{0.3em}

  #2

  \end{minipage}%

  \vspace{0.8em}%

}



\newcommand{\bilingual}[2]{%

  \noindent\begin{minipage}[t]{0.48\textwidth}#1\end{minipage}%

  \hfill%

  \begin{minipage}[t]{0.48\textwidth}#2\end{minipage}%

  \vspace{0.5em}%

}



%% -------- Problem environment labels --------

\newcommand{\Qwen}{{\color{labelclassical}\textbf{问曰}}}

\newcommand{\Qmodern}{{\color{labelmodern}\textbf{问题}}}

\newcommand{\Awen}{{\color{labelclassical}\textbf{答曰}}}

\newcommand{\Amodern}{{\color{labelmodern}\textbf{答案}}}

\newcommand{\Swen}{{\color{labelclassical}\textbf{术曰}}}

\newcommand{\Smodern}{{\color{labelmodern}\textbf{方法}}}

\newcommand{\Cwen}{{\color{labelclassical}\textbf{草曰}}}

\newcommand{\Cmodern}{{\color{labelmodern}\textbf{演算过程}}}



%% -------- Footnote style for terminology --------

\newcommand{\termnote}[1]{\footnote{\textcolor{notecolor}{\textbf{术语说明：}#1}}}



%% -------- Hyperref Setup --------

\hypersetup{

  colorlinks=true,

  linkcolor=sectioncolor,

  citecolor=sectioncolor,

  urlcolor=sectioncolor,

  pdfencoding=auto

}



%% -------- TOC Depth --------

\setcounter{tocdepth}{2}

\setcounter{secnumdepth}{2}



%% ========================================================================



\begin{document}



%% ========================================================================

%% TITLE PAGE

%% ========================================================================

\begin{titlepage}

\centering

\vspace*{2cm}



{\fontsize{24}{30}\selectfont\chineselg

\textcolor{titlecolor}{数书九章 · 大衍类}}



\vspace{0.5cm}



{\fontsize{14}{18}\selectfont

\textcolor{sectioncolor}{文言文原文 -- 白话文今译 对照版}}



\vspace{1.5cm}



{\LARGE\bfseries 第一卷}



\vspace{0.3cm}



{\Large 大衍类 · 九章}



\vspace{1.5cm}



{\large 秦九韶 著\\[0.3em]

淳祐七年（公元1247年）}



\vspace{1cm}



{\normalsize 白话文今译与注释}



\vspace{0.5cm}



\begin{minipage}{0.85\textwidth}

\centering

\small

\textit{本书为世界数学史上最伟大的著作之一，\\包含大衍法（中国剩余定理）、方程数值解、\\测量学问题以及商业与金融数学等内容。}

\end{minipage}



\vfill



{\small 用\LaTeX{}及xeCJK宏包双语对照排版}



\end{titlepage}



%% ========================================================================

%% TABLE OF CONTENTS

%% ========================================================================

\tableofcontents

\newpage



%% ========================================================================

%% INTRODUCTION

%% ========================================================================

\section{简介}



\bilingualsection{%

《数书九章》为南宋数学家秦九韶（约1202--1261年）所著，成书于淳祐七年（1247年），乃中国数学史乃至世界数学史上最伟大的著作之一。



全书分为九类，每类九题，凡八十一题。九类者：大衍、天时、田域、测望、赋役、钱谷、营建、军旅、市易。



每题遵循严密的四段结构：问（问题陈述）、答（数值答案）、术（通法与算法）、草（详细演算步骤）。

}{%

《数书九章》是南宋数学家秦九韶（约1202--1261年）的代表作，完成于淳祐七年（1247年），是中国乃至世界数学史上最重要的著作之一。



全书分为九个类别，每个类别包含九个问题，共八十一题。九个类别分别是：大衍（不定分析）、天时（天文历法）、田域（土地测量）、测望（测量观测）、赋役（税收徭役）、钱谷（粮钱商业）、营建（建筑工程）、军旅（军事事务）、市易（市场交易）。



每个问题都遵循严谨的四段式结构：问题陈述、数值答案、通用方法、详细演算。}



\subsection{大衍类概述}



\bilingualsection{%

大衍类，即《周易》大衍之数。此章专论线性同余式方程组之通解，即今所谓「中国剩余定理」。



秦九韶所创「大衍总数术」提供系统算法，「大衍求一术」则给出求模乘逆元之程序。



大衍类九题：一、蓍卦发微；二、古历会积；三、推计土功；四、推库额钱；五、分粜推原；六、程行计地；七、程行相及；八、积尺寻源；九、余米推数。

}{%

大衍类，取名于《周易》中的「大衍之数」。这一章专门讨论线性同余方程组的通解，即今天所称的「中国剩余定理」。



秦九韶创立的「大衍总数术」提供了系统的计算方法，而「大衍求一术」（求模逆元的方法）则给出了求乘法逆元的计算程序。



大衍类的九个问题分别是：一、蓍卦发微（蓍草占卜的数学原理）；二、古历会积（古代历法中的周期计算）；三、推计土功（土方工程估算）；四、推库额钱（库房资金的数学计算）；五、分粜推原（粮食分配的推算）；六、程行计地（行程距离计算）；七、程行相及（行程相遇问题）；八、积尺寻源（由体积求尺寸）；九、余米推数（剩余粮食数量推算）。}



\newpage



%% ========================================================================

%% PREFACE

%% ========================================================================

\section{秦九韶自序}



\bilingualsection{%

周教六艺，数实成之。学士大夫所从事尚矣。其用本太虚生一，而周流无穷。大则可以通神明、顺性命，小则可以经世务、类万物，讵容以浅近窥哉？

}{%

周代教育有六艺，数学是其中之一。学者和官员们历来都很重视它。数学的本源来自「太虚生一」（宇宙从无到有的生成），而其应用无穷无尽。往大了说可以通晓天地规律、顺应自然本性，往小了说可以治理世间事务、分类万物，怎么能用浅薄的眼光来看待呢？}



\vspace{0.5em}



\bilingualsection{%

若昔推策以迎日，定律而和气。髀距濬川，土圭度晷。天地之大，囿焉而不能外，况其间总总者乎？爰自河图洛书，闿发幽秘；八卦九畴，错综精微。极而至于大衍、皇极之用，而人事之变无不该，鬼神之情莫能隐矣。圣人神之言而遗其粗，常人昧之由而莫之觉。要其归，则数与道非二本也。

}{%

从前人们用算筹推算来迎接新年的到来，制定音律来调和阴阳之气。用勾股定理来疏通河道，用土圭来测量日影确定时节。天地之大，都在数学的包罗之中而无从逃逸，何况天地间纷繁复杂的事物呢？从河图洛书开始，就开启了深奥的奥秘；八卦和九畴（《尚书·洪范》中的九种治国大法），交错精微。发展到使用大衍之数和皇极之数，人间万事的变化无所不包，鬼神的奥秘也无法隐藏。圣人领悟了其中的精妙而舍弃了粗浅的表面，普通人却因蒙昧而不曾察觉。归根结底，数学与大道并不是两个不同的根本。}



\vspace{0.5em}



\bilingualsection{%

汉去古未远，有张苍、许商、马延年、耿寿昌、郑玄、张衡、刘洪之伦，或明天道而法传于后，或计功策而效验于时。后世学者自高鄙不之讲，此学殆绝。惟治历畴人能为乘除，而弗通于开方衍变。若官府会事，则府史一二累之，算家位置素所不识，上之人亦委而听焉。持算者惟若人，则鄙之也宜矣。

}{%

汉代距离古代还不太远，有张苍、许商、马延年、耿寿昌、郑玄、张衡、刘洪等人，有的通晓天道而将其方法传于后世，有的用计算来验证实际功效。后世的学者自高自大、鄙视数学，这门学问几乎断绝。只有修订历法的专业人士能做些乘除运算，却不通晓开方等更复杂的数学变化。至于官府的会计事务，则交给文书小吏一点点累积计算，真正的数学家所擅长的算法他们完全不懂，而上级官员也听之任之。掌管计算的人如果是这样的人，那么被人鄙视也是理所应当的了。}



\vspace{0.5em}



\bilingualsection{%

呜呼！乐有制氏，仅记铿锵之声，而谓与天地同和者止于是，可乎？今数术之书尚三十余家。天象历度谓之缀术，太乙壬甲谓之三式，皆曰内算，言其秘也。九章所载及周官九数，系于方圆者为专术，皆曰外算，对内而言也。其用相通，不可歧二。独大衍法不载九章，未有能推之者。历家演法颇用之，以为方程者，误也。

}{%

唉！音乐有制氏，只记录了铿锵的声音，就认为与天地和谐的音乐仅此而已，这怎么可以呢？现在数学方面的书籍还有三十多家。天文历法的计算称为「缀术」，太乙、六壬、遁甲称为「三式」，都叫做「内算」，说的是它们的深奥秘密。《九章算术》所载以及《周礼》中的九数，涉及几何图形的专门技术，称为「外算」，是相对于「内算」而言的。它们的用途是相互贯通的，不能一分为二。唯独大衍法没有记载在《九章》中，没有人能够推究其理。历法学家在推算中常用到它，但把它当作普通的方程来处理，这是错误的。}



\vspace{0.5em}



\bilingualsection{%

且天下之事多矣。古之人先事而计，计定而行。仰观俯察，人谋鬼谋，无所不用其谨，是以不愆于成。载籍章章可覆也。后世兴事造始，鲜能考度，悠悠乎天纪人事之殽缺矣，可不求其故哉？

}{%

况且天下的事务如此繁多。古人做事之前先谋划，谋划好了才行动。上观天文、下察地理，与人商议、与神龟卜卦，无不极其谨慎，因此做事不会失败。典籍中清清楚楚可以覆按验证。后世的人兴办事业、开创新事物，却很少能考察度量，长此以往，天道的规律和人事的条理就都混乱了，能不探求其原因吗？}



\vspace{0.5em}



\bilingualsection{%

九韶愚陋，不闲于艺。然早岁侍亲中都，因得访习于太史，又尝从隐君子受数学。时际兵难，历岁遥塞，不自意全于矢石之间。更险离忧，荏苒十禩，心槁气落。信知夫物莫不有数也，乃肆意其间，旁诹方能，探索杳眇，粗若有得焉。所谓通神明、顺性命，固肤末于见。若其小者，窃尝设为问答，以拟于用。积多而惜其弃，因取八十一题，厘为九类，立术具草，间以图发之。恐或可备博学多识君子之余观，曲艺可遂也，愿进之于道。傥曰艺成而下，是惟畴人府史流也，乌足尽天下之用，亦无瞢焉。

}{%

我秦九韶愚钝浅陋，不精通于六艺。但早年侍奉父亲在都城（临安），因此得以向太史请教学习，又曾跟随隐士学习数学。当时正逢兵荒马乱，经历了漫长的岁月和遥远的边塞，不曾料想自己能在刀箭之间保全性命。历经艰险、饱尝离别之忧，荏苒十年，心力交瘁。我确信万物皆有数理，于是纵情于数学之中，广泛搜集各种方法，探索幽深微妙的道理，粗略有些心得。所谓通晓神明、顺应天命，自然是我所难以企及的。至于那些较小的应用，我私下曾设计一些问题与解答，以模拟实际应用。积累多了而舍不得丢弃，于是选取八十一个问题，分为九类，设立方法、写出详细演算，间或用图示来阐明。或许可以供博学多识的君子闲暇时一览，技艺虽然微小，但愿能上升到道的层面。倘若有人说技艺成就不过是末流，这只是历算官员和文书小吏的学问，不足以穷尽天下的用途，我也并不感到困惑。}



\vspace{0.5em}



\bilingual{%

时淳祐七年九月，鲁郡秦九韶叙。

}{%

时淳祐七年九月，鲁郡秦九韶序。}



\newpage



%% ========================================================================

%% DAYAN GENERAL METHOD

%% ========================================================================

\section{大衍总数术}



\subsection{大衍法总论}



\bilingualsection{%

按大衍术，以各分数之奇零求各分数之总数。大而天行，小而物数，皆可御之。其法有求元、求定、求衍、求奇、求乘、求用之目。大约以数之奇偶为根，而以诸数相度之尽不尽为用。

}{%

大衍法的原理是：根据各个数相除后的余数，求出满足所有余数条件的总数。大到天体运行的周期，小到物品的数量，都可以用这个方法来解决。其计算步骤包括：求元数、求定数、求衍数、求奇数、求乘率、求用数。大体上以数的奇偶性质为根本，以各数之间能否整除（互素与否）为运算依据。}



\vspace{0.5em}



\bilingualsection{%

有求彼此不能度尽之诸数者，元数、定数是也。有求诸数皆能度尽之一数者，衍母数是也。有求诸数皆能度尽而一数不能度尽之数者，各衍数是也。其不尽之数，即奇数也。有求二数相度余一之数者，乘数是也。有求二数相度余一而诸数又能度尽之数者，用数是也。

}{%

所谓「元数」和「定数」，就是彼此不能整除的数（即两两互素的数）。所谓「衍母数」，就是所有定数都能整除的一个数（即所有定数的乘积）。所谓「衍数」，就是所有定数都能整除、但其中一个定数不能整除的数（即衍母除以该定数所得的商）。衍数被定数除后所得的余数，就是「奇数」。能使某数与奇数相乘后除以定数余一的数，就是「乘率」。能使某数与衍数相乘后满足所有条件的数，就是「用数」。}



\subsection{问数四品}



\bilingualsection{%

\textbf{大衍总数术曰：}置诸问数，类名有四。



\begin{enumerate}[label=\arabic*.,leftmargin=1.5em,topsep=0pt]

  \item \textbf{元数} -- 谓尾位见单零者。本门揲蓍、酒息、斛粜、砌砖、失米之类是也。

  \item \textbf{收数} -- 谓尾位见分厘者。假令冬至三百六十五日二十五刻，欲与甲子六十日为一会而求积日之类。

  \item \textbf{通数} -- 谓诸数各有分子分母者。本门问一会积年是也。

  \item \textbf{复数} -- 谓尾位见十或百及千以上者。本门筑堤并急足之类是也。

\end{enumerate}

}{%

\textbf{大衍总数术说：}列出问题中给出的各数（问数），按照其类型可分为四类。



\begin{enumerate}[label=\arabic*.,leftmargin=1.5em,topsep=0pt]

  \item \textbf{元数} -- 指末尾为单个零的整数。如本章中蓍草占卜、利息计算、粮食买卖、砌砖、失米等问题中的数。

  \item \textbf{收数} -- 指末尾带有小数（分、厘）的数。例如冬至为365.25日（三百六十五日二十五刻），想与甲子周期60日求一个共同的会积日数之类的问题。

  \item \textbf{通数} -- 指各数都有分子和分母的分数。本章中求一会积年的问题就属于这类。

  \item \textbf{复数} -- 指末尾有十、百、千等倍数的数。本章中筑堤和急足（快差）问题中的数就属于这类。

\end{enumerate}}



\newpage



\subsection{求定数}



\bilingualsection{%

\textbf{元数者：}先以两两连环求等，约奇弗约偶。或约得五而彼有十，乃约偶弗约奇。或元数俱偶，约毕可存一位见偶。或皆约而犹有类数存，姑置之，俟与其他约遍而后乃与姑置者求等约之。或诸数皆不可尽类，则以诸元数命曰定数，以复数格入之。

}{%

\textbf{元数的处理方法：}先将各数两两配对求最大公约数（等数），然后约简：约定奇数（位数为奇数的数）而不约定偶数（位数为偶数的数）。但如果约得5而另一个数有10，就要反过来约偶数不约奇数。如果所有元数都是偶数，约简后可以保留一个偶数。如果约简后仍有可继续约的数（类数），暂时搁置，等与其他数都约遍之后，再与搁置的数求等数约简。如果各数之间都没有公因子（不可尽类），就直接将各元数作为「定数」，按复数的方法处理。}



\vspace{0.5em}



\bilingualsection{%

\textbf{收数者：}乃命尾位分厘作单零，以进所问之数，定位讫，用元数格入之。或如意立数为母，收进分厘以从所问，用通数格入之。

}{%

\textbf{收数的处理方法：}将末尾的小数（分、厘）当作单位零，将所问的数进位为整数，定位完毕后用元数的方法处理。或者任意设一个数作为分母，将小数部分收进凑整，用通数的方法处理。}



\vspace{0.5em}



\bilingualsection{%

\textbf{通数者：}置问数通分纳子，互乘之，皆曰通数。求总等，不约一位，约众位，得各原法数，用元数格入之。

}{%

\textbf{通数的处理方法：}将各问数通分后加分子，互相乘，都称为通数。求总等数（所有数的最大公约数），保留一位不约，约其余各位，得到各原法数，再用元数的方法处理。}



\vspace{0.5em}



\bilingualsection{%

\textbf{复数者：}问数尾位见十以上者，以诸数求总等，存一位约众位，始得元数。两两连环求等，约奇弗约偶，复乘偶；或约偶弗约奇，弗乘奇。

}{%

\textbf{复数的处理方法：}问数末尾有十以上的，先求各数的总等数，保留一位不约、约其余各位，才能得到元数。然后两两连环求等，约定奇数不约偶数，并将偶数乘回（复乘偶）；或者约偶数不约奇数，不将奇数乘回（弗乘奇）。}



\subsection{求衍母、衍数、奇数与乘率}



\bilingualsection{%

\textbf{求定数}勿使两位见偶，勿使见一太多。见一多则借用繁，不欲借则任得一。



\textbf{以定相乘为衍母}，以各定约衍母，得各衍数。或列各定数于右行，各立天元一为子于左行，以母互乘子，亦得衍数。

}{%

\textbf{求定数}时要注意：不要使两个定数都是偶数，也不要使定数中出现太多1。出现1太多会导致借用繁琐，如果不想要借用就任意取一个数。



\textbf{将各定数相乘得到衍母}（所有定数的乘积），用各定数分别去除衍母，得到各衍数。或者将各定数列在右边，各立天元一（1）作为子列在左边，用各母数互乘各子，也可以得到衍数。}



\vspace{0.5em}



\bilingualsection{%

次以各定母满去衍数，各余名曰\textbf{奇数}。以奇数与定母用\textbf{大衍求一术}入之，得\textbf{乘率}。以乘衍数，各得\textbf{用数}。



凡欲求总数者，置各余数乘用数，并之为\textbf{总数}。满衍母去之，不满为所求数。

}{%

然后用各定母分别去除衍数（求余），所得的余数称为\textbf{奇数}（模运算中的余数）。用奇数和定母通过\textbf{大衍求一术}（求模逆元的方法），求得\textbf{乘率}（乘法因子）。用乘率乘以衍数，得到\textbf{用数}。



如果要求满足所有条件的总数，将各余数乘以对应的用数，相加得到\textbf{总数}。用衍母去除总数，余数就是所求的答案。}



\newpage



\subsection{大衍求一术}



\bilingualsection{%

\begin{quote}

大衍求一术云：置奇右上，定居右下。立天元一于左上。先以右行上下两位，以少除多，所得商数，乃递互乘归左行。须使右上末后奇一而止。乃验左上所得，以为乘率。或奇数已见单一者，便为乘率。

\end{quote}

}{%

\begin{quote}

大衍求一术说：将奇数放在右上方，将定数放在右下方。将天元一（1）放在左上方。先用右列的上下两个数，以小除大，得到商数，然后依次互乘归入左列。直到右上方最终变为1时停止。然后检验左上方所得的数，这就是乘率。如果奇数已经是1，那么奇数本身就是乘率。

\end{quote}}



\vspace{0.5em}



\bilingualsection{%

\textit{用今语释之：}设奇数为$a$，定数为$m$，两数互素。欲求乘率$k$，使得：

\[k \cdot a \equiv 1 \pmod{m}\]

此即求$a$模$m$之乘逆，用辗转相除法（欧几里得算法）配以回代即可得之。

}{%

\textit{用现代数学语言解释：}设奇数为$a$，定数为$m$，且$a$与$m$互素。要求乘率$k$，使得：

\[k \cdot a \equiv 1 \pmod{m}\]

这就是求$a$关于模$m$的模逆元，可以通过扩展欧几里得算法（辗转相除法加回代）来求得。}



\newpage



%% ========================================================================

%% PROBLEM 1: DIVINATION BY YARROW STALKS

%% ========================================================================

\section{第一题：蓍卦发微}



\bilingualsection{%

\Qwen 《易》曰：大衍之数五十，其用四十有九。又曰：分而为二以象两，挂一以象三，揲之以四以象四时。三变而成爻，十有八变而成卦。欲知所衍之术及其数各几何？

}{%

\Qmodern 《周易》说：大衍之数为五十，实际使用的是四十九。又说：将四十九根蓍草分成两份以象征天地两仪，取出一根挂在手指间以象征天地人三才，每四根一组数数以象征春夏秋冬四时。三次变化确定一爻，十八次变化确定一卦。请问这个推演方法的计算原理以及各个数值分别是多少？}



\vspace{0.5em}



\bilingualsection{%

\Awen 衍母十二，衍法三。



一元衍数二十四，二元衍数一十二，三元衍数八，四元衍数六。以上四位衍数，计五十一。



一揲用数一十二，二揲用数二十四，三揲用数四，四揲用数九。以上四位用数，计四十九。

}{%

\Amodern 衍母（所有定数的乘积）为12，衍法为3。



第一揲（分堆）的衍数为24，第二揲的衍数为12，第三揲的衍数为8，第四揲的衍数为6。以上四个衍数，合计为51。



第一揲的用数为12，第二揲的用数为24，第三揲的用数为4，第四揲的用数为9。以上四个用数，合计为49。}



\vspace{0.5em}



\bilingualsection{%

\Swen 置诸元数，两两连环求等，约奇弗约偶。遍约毕，乃变元数皆曰定母，列右行。各立天元一为子，列左行。以诸定母互乘左行之子，各得名曰衍数。次以各定母满去衍数，各余名曰奇数。以奇数与定母用大衍术求一，得各乘率。以乘衍数，各得用数。验次所揲余几何，以其余数乘诸用数，并之，名曰总数。满衍母去之，不满为所求数。以为实，易以三才为衍法，以法除实所得为象数。如实有余，或一或二，皆命作一，同为象数。其象数得一为老阳，得二为少阴，得三为少阳，得四为老阴。得老阳画重爻，得少阴画拆爻，得少阳画单爻，得老阴画交爻。凡六画乃成卦。

}{%

\Smodern 列出各元数，两两配对求最大公约数（等数），约定奇数（位数为奇数的数）不约偶数（位数为偶数的数）。全部约简完毕后，将各元数改称为定母，列在右边。各立天元一（1）作为子，列在左边。用各定母互乘左列的各子，分别得到衍数。然后用各定母分别去除衍数，所得余数称为奇数。用奇数和定母通过大衍求一术，求得各乘率。用乘率乘以衍数，得到各用数。检验每次揲蓍后的余数是多少，用各余数乘对应的用数，相加得到总数。用衍母去除总数，余数就是所求的数。以此作为被除数，以三才之数3作为衍法（除数），相除所得为象数。如果除后有余数，无论是1还是2，都当作1处理，与商相加作为象数。象数为1的是老阳，为2的是少阴，为3的是少阳，为4的是老阴。得老阳画重爻（——），得少阴画拆爻（- -），得少阳画单爻（—），得老阴画交爻（X）。共六爻组成一卦。}



\vspace{0.5em}



\bilingualsection{%

\Cwen 置一、二、三、四列右行，立天元一列左行。



以右行一、二、三、四互乘左行：一乘一得一，二乘一得二，三乘二得六，四乘六得二十四。各得衍数。



次以定母满法去衍数：以十二除二十四得二余零，一元衍数为二十四；以十二除一十二得一余零，二元衍数为一十二；以八除八得一余零，三元衍数为八；以六除六得一余零，四元衍数为六。



各满定母去之，得奇数。以大衍求一入之，得乘数：得二十三为一揲乘率，得五十一为二揲乘率，得七为三揲乘率。



以乘率各乘衍数，得用数：一揲用数一十二，二揲用数二十四，三揲用数四，四揲用数九。以上四位用数，计四十九。

}{%

\Cmodern 将1、2、3、4列在右行，将1（天元一）列在左行。



用右行的1、2、3、4分别互乘左行的数：$1 \times 1 = 1$，$2 \times 1 = 2$，$3 \times 2 = 6$，$4 \times 6 = 24$。分别得到各衍数。



然后用各定母去除衍数：$24 \div 12 = 2$余0，第一揲衍数为24；$12 \div 12 = 1$余0，第二揲衍数为12；$8 \div 8 = 1$余0，第三揲衍数为8；$6 \div 6 = 1$余0，第四揲衍数为6。



用各定母去除衍数后得到奇数。通过大衍求一术求得乘率：第一揲乘率为23，第二揲乘率为51，第三揲乘率为7。



用乘率分别乘以衍数得到用数：第一揲用数为12，第二揲用数为24，第三揲用数为4，第四揲用数为9。四个用数合计为49。}



\newpage



%% ========================================================================

%% PROBLEM 2: ANCIENT CALENDAR ACCUMULATION

%% ========================================================================

\section{第二题：古历会积}



\bilingualsection{%

\Qwen 问古历会积，欲求积年。以历法积元为问，历过无考，但据近世演纪之法，推而上之。

}{%

\Qmodern 问古代历法中各种周期的会积问题，要求推算积年数。以历法的积元（历元）为问题，具体的历法已无从考证，但根据近世推演纪年的方法，向上推算。



\textit{此题涉及古代历法中多个天文周期的同步问题。给定某些天文周期，需要找到一个数，同时满足多个同余条件。}}



\vspace{0.5em}



\bilingualsection{%

\Awen 答曰：积年若干。

}{%

\Amodern 答案：所求积年数为特定数值。（具体数值需根据历法参数计算得出。）}



\vspace{0.5em}



\bilingualsection{%

\Swen 术曰：以大衍求之。置元历相关诸周期，连环求等，约奇弗约偶，得定母。诸定相乘为衍母。以各定约衍母得衍数。各满定母去之，余为奇数。大衍求一入之，得乘率。以乘衍数为用数。次置各余数乘用数，并之为总数。满衍母去之，不满为所求积年。

}{%

\Smodern 方法：用大衍法求解。列出历法中的各个相关周期，两两配对求最大公约数（等数），约定奇数不约偶数，得到定母（两两互素的因数）。各定母相乘得到衍母（所有定数的乘积）。用各定母分别去除衍母得到衍数（衍母除以该定数）。用各定母分别去除衍数，余数为奇数（模运算中的余数）。用大衍求一术（求模逆元的方法）求得乘率（乘数因子）。用乘率乘以衍数得到用数。然后将各余数乘以对应的用数，相加得到总数。用衍母去除总数，余数就是所求的积年。



\textit{本题将大衍法应用于天文历法计算，寻找一个共同的历元，使得多个天体运行周期（如回归年、朔望月、甲子干支周期等）同时对齐。这是中国传统天文学中大衍法的主要应用之一。}}



\newpage



%% ========================================================================

%% PROBLEM 3: ESTIMATING EARTHWORKS

%% ========================================================================

\section{第三题：推计土功}



\bilingualsection{%

\Qwen 问：四县（甲、乙、丙、丁）共同筑堤，根据各县的物力分配筑堤任务。



甲县物力：一十三万八千六百贯；乙县物力：一十四万六千三百贯；丙县物力：一十九万二千五百贯；丁县物力：一十八万四千八百贯。



筑堤标准：每七百七十贯物力分配一名劳动力，每人每天筑堤六十平方尺。



进度情况：甲县和乙县已完成任务，丙县剩余五十一丈，丁县剩余十八丈。

}{%

\Qmodern 问题：甲、乙、丙、丁四个县共同修筑堤坝，按照各县的财力物力分配筑堤任务。



甲县财力：138,600贯；乙县财力：146,300贯；丙县财力：192,500贯；丁县财力：184,800贯。



筑堤标准：每770贯财力分配一名劳动力，每个劳动力每天修筑60平方尺的堤坝。



施工进度：甲县和乙县已经完成了任务，丙县还剩51丈未完成，丁县还剩18丈未完成。求堤坝的总长度以及各县修筑的长度。}



\vspace{0.5em}



\bilingualsection{%

\Awen 答曰：堤的总长度为十九里二百三十五步五尺。



各县所筑长度：甲县一千二百六十丈；乙县一千七百六十八步五尺六寸；丙县四里三百二十八步五尺六寸；丁县与前三县相应。

}{%

\Amodern 答案：堤坝总长度为19里235步5尺。



各县修筑的长度：甲县1,260丈；乙县1,768步5尺6寸；丙县4里328步5尺6寸；丁县与前三县相对应。}



\vspace{0.5em}



\bilingualsection{%

\Swen 术曰：以大衍求之。置各县物力，连环求等，约奇弗约偶，得定母。诸定相乘为衍母。以各定约衍母得衍数。各满定母去之，余为奇数。大衍求一入之，得乘率。以乘衍数为用数。次置各县余数乘用数，并之为总数。满衍母去之，不满为所求。

}{%

\Smodern 方法：用大衍法求解。列出各县的财力数值，两两配对求最大公约数，约定奇数不约偶数，得到定母。各定母相乘得到衍母。用各定母去除衍母得到衍数。用定母去除衍数，余数为奇数。用大衍求一术求得乘率。乘率乘以衍数得到用数。然后将各县剩余量乘以对应的用数，相加得到总数。用衍母去除总数，余数就是所求的答案。}



\vspace{0.5em}



\bilingualsection{%

\Cwen 计算各县每日筑堤长度：甲县五十四丈，乙县五十七丈，丙县七十五丈，丁县七十二丈。



使用大衍求一术，计算各县的复数和衍数：甲县复数二十七，乙县复数十九，丙县复数二十五，丁县复数八。



衍母：十万二千六百。甲县衍数三千八百，乙县衍数五千四百，丙县衍数四千一百四十，丁县衍数一万二千八百二十五。



乘率：甲县二十有三，乙县有五，丙县十有九，丁县有一。



用数：甲县一万七千四百，乙县二万七千，丙县七万七千九百七十六，丁县一万二千八百二十五。



总长度：丙县剩余三百九十七万六千七百七十六丈，丁县剩余二十三万八千五百丈，总剩余四百二十一万五千二百七十六丈。



堤的总长度：十九里二百三十五步五尺。

}{%

\Cmodern 计算各县每天筑堤的长度：甲县54丈，乙县57丈，丙县75丈，丁县72丈。



使用大衍求一术，计算各县的复数和衍数：甲县复数为27，乙县复数为19，丙县复数为25，丁县复数为8。



衍母$M = 102,600$。甲县衍数$M_1 = 3,800$，乙县$M_2 = 5,400$，丙县$M_3 = 4,140$，丁县$M_4 = 12,825$。



乘率：$k_1 = 23$，$k_2 = 5$，$k_3 = 19$，$k_4 = 1$。



用数：$u_1 = 17,400$，$u_2 = 27,000$，$u_3 = 77,976$，$u_4 = 12,825$。



丙县剩余量对应3,976,776丈，丁县剩余量对应238,500丈，总剩余量对应4,215,276丈。



堤坝的总长度为19里235步5尺。}



\newpage



%% ========================================================================

%% PROBLEM 4: CALCULATING TREASURY FUNDS

%% ========================================================================

\section{第四题：推库额钱}



\bilingualsection{%

\Qwen 问有外邑七库，日纳息足钱适等，递年成贯整纳。近缘见钱希少，听各库照当处市陌准解旧会。其甲库有零钱一十文，丁、庚二库各零四文，戊库零六文，余库无零钱。甲库所在市陌一十二文，递减一文至庚库而止。欲求诸库日息原纳足钱、展省及今纳旧会，并大小月分各几何？



（七库依次为甲、乙、丙、丁、戊、己、庚，市陌分别为十二、十一、十、九、八、七、六文。）

}{%

\Qmodern 问题：有外地七个库房，每天交纳的利息足钱数额相等，按年成贯交纳整数。近来因为现钱稀少，允许各库房按照当地市陌（货币兑换率）折算交纳旧会（纸币）。甲库有零钱10文，丁库和庚库各有零钱4文，戊库有零钱6文，其余库房没有零钱。甲库所在地的市陌为12文（即12文旧会兑换1文足钱），依次递减1文直到庚库。要求各库每天原来应交纳的足钱数、展省（折算后的金额）以及现在应交纳的旧会数，还有大小月的分配各是多少？



（七库依次为甲、乙、丙、丁、戊、己、庚，市陌分别为12、11、10、9、8、7、6文。）}



\vspace{0.5em}



\bilingualsection{%

\Awen 答曰：诸库纳日息足钱二十贯九百五十文。展省三十五贯文。各库旧会：甲库日息旧会二百二十四贯五百一十文；乙库日息旧会二百四十五贯文；丙库日息旧会二百六十九贯五百文；庚库日息旧会四百四十九贯一百四文。

}{%

\Amodern 答案：各库每天应交纳的足钱为20贯950文。展省（折算差额）为35贯文。各库每天应交纳的旧会：甲库224贯510文；乙库245贯文；丙库269贯500文；庚库449贯104文。}



\vspace{0.5em}



\bilingualsection{%

\Swen 术曰：以大衍求之。置甲库市陌，以递减数减之，各得诸库原陌。连环求等，约奇弗约偶，得定母。诸定相乘为衍母。以定约衍母得衍数。各满定母去之，余为奇数。大衍求一入之，得乘率。以乘衍数为用数。次置有零文库零钱数，乘本用数，并为总数。满衍母去之，不满为诸库日息足钱。

}{%

\Smodern 方法：用大衍法求解。以甲库的市陌为基准，用递减数依次相减，得到各库的市陌。两两配对求等数，约定奇数不约偶数，得到定母。各定母相乘得到衍母。用定母去除衍母得到衍数。用定母去除衍数，余数为奇数。用大衍求一术求得乘率。乘率乘以衍数得到用数。然后将有零钱的库房的零钱数乘以对应的用数，相加得到总数。用衍母去除总数，余数就是各库每天应交纳的足钱数。}



\vspace{0.5em}



\bilingualsection{%

\Cwen 置甲库市陌一十二，递减一得一十一为乙库陌，十为丙库陌，九为丁库陌，八为戊库陌，七为己库陌，六为庚库陌。得诸库原陌。



连环求等约讫：甲得一，乙得十一，丙得五，丁得九，戊得八，己得七，庚得一。各为定母。



先以诸定相乘，得二万七千七百二十为衍母。

次以诸定互乘诸子：甲得二万七千七百二十，乙得二千五百二十，丙得五千五百四十四，丁得三千零八十，戊得三千四百六十五，己得三千九百六十，庚得二万七千七百二十。各为衍数。

}{%

\Cmodern 以甲库市陌12为基准，递减1：11为乙库市陌，10为丙库市陌，9为丁库市陌，8为戊库市陌，7为己库市陌，6为庚库市陌。得到各库的原始市陌。



两两连环求等数并约简后：甲得1，乙得11，丙得5，丁得9，戊得8，己得7，庚得1。分别为各库的定母。



先将各定母相乘，得到$M = 27,720$作为衍母。

然后用各定母互乘各子：甲库衍数$M_1 = 27,720$，乙库$M_2 = 2,520$，丙库$M_3 = 5,544$，丁库$M_4 = 3,080$，戊库$M_5 = 3,465$，己库$M_6 = 3,960$，庚库$M_7 = 27,720$。}



\newpage



%% ========================================================================

%% PROBLEM 5: GRAIN DISTRIBUTION

%% ========================================================================

\section{第五题：分粜推原}



\bilingualsection{%

\Qwen 有兄弟三人，各将收获之米送往不同的地方：甲将米送往官场，余三斗二升；乙将米送往安吉乡，余七斗；丙将米送往平江，余三斗。



官斛每石八斗三升，安吉斛每石一石一斗，平江斛每石一石三斗五升。欲求三人共收获米数以及各自分米数几何？

}{%

\Qmodern 有兄弟三人，各自将收获的米运往不同的地方：甲将米运往官场出售，剩余3斗2升；乙将米运往安吉乡出售，剩余7斗；丙将米运往平江出售，剩余3斗。



官斛的容量为每石83升，安吉斛为每石110升，平江斛为每石135升。要求三人共收获的米数以及每人分得的米数各是多少？}



\vspace{0.5em}



\bilingualsection{%

\Awen 答曰：三人共收获米七百三十八石。



各分米数：甲二百九十六石余三斗二升；乙二百二十三石余七斗；丙一百八十二石余三斗。

}{%

\Amodern 答案：三人共收获米738石。



各人分得的米数：甲分得296石，剩余3斗2升；乙分得223石，剩余7斗；丙分得182石，剩余3斗。}



\vspace{0.5em}



\bilingualsection{%

\Swen 术曰：以大衍求之。置各斛率（单位：升）：官斛八十三升，安吉斛一百一十升，平江斛一百三十五升。连环求等，约奇弗约偶，得定母。诸定相乘为衍母。以各定约衍母得衍数。各满定母去之，余为奇数。大衍求一入之，得乘率。以乘衍数为用数。次置各余数乘用数，并之为总数。满衍母去之，不满为所求率数。

}{%

\Smodern 方法：用大衍法求解。列出各斛的容量（单位为升）：官斛83升，安吉斛110升，平江斛135升。两两配对求等数，约定奇数不约偶数，得到定母（两两互素的因数）。各定母相乘得到衍母（所有定数的乘积）。用各定母去除衍母得到衍数。用定母去除衍数，余数为奇数（模余数）。用大衍求一术（求模逆元的方法）求得乘率。乘率乘以衍数得到用数。然后将各余数乘以对应的用数，相加得到总数。用衍母去除总数，余数就是所求的容量率数。}



\vspace{0.5em}



\bilingualsection{%

\Cwen 置官斛八十三升、安吉斛一百一十升、平江斛一百三十五升，连环求等。



求总等：得总等一，不约。连环求等，得定母：二十四万六千五百一十。



各斛衍数：官斛衍数二千九百七十，安吉斛衍数二千二百四十一，平江斛衍数一千八百二十六。



乘率：官斛五十一，安吉七，平江二十三。



用数：官斛十五万一千四百七十，安吉一万五千六百八十七，平江四万一千九百九十八。



并总数，满衍母去之，展算得：分米二百四十六石，以三人因之，得七百三十八石为共米。各以斛率约之，合问。

}{%

\Cmodern 列出官斛83升、安吉斛110升、平江斛135升，两两连环求等数。



求总等数：得总等数为1，不需要约简。连环求等后，得到衍母$M = 246,510$。



各斛的衍数：官斛$M_1 = 2,970$，安吉斛$M_2 = 2,241$，平江斛$M_3 = 1,826$。



乘率：官斛$k_1 = 51$，安吉$k_2 = 7$，平江$k_3 = 23$。



用数：官斛$u_1 = 151,470$，安吉$u_2 = 15,687$，平江$u_3 = 41,998$。



计算总数后用衍母取余，展开计算得：每人分米246石，三人共$246 \times 3 = 738$石为总收获米数。各用斛率约分，与题意相合。}



\newpage



%% ========================================================================

%% PROBLEM 6: JOURNEY DISTANCE CALCULATION

%% ========================================================================

\section{第六题：程行计地}



\bilingualsection{%

\Qwen 军师派遣三名急足（甲、乙、丙）分别前往都下报捷。甲每天行三百里；乙每天行二百四十里；丙每天行一百八十里。



甲提前数日申未到，乙后数日未正到，丙今日辰末到。欲知从军前到都下的距离以及三人各自行走的天数几何？

}{%

\Qmodern 军师派遣三名快差（甲、乙、丙）分别前往都城报捷。甲每天行走300里；乙每天行走240里；丙每天行走180里。



甲比约定时间提前到达（申时末之前到达），乙比约定时间晚到（未时正才到），丙在今天辰时末到达。求从军营前哨到都城的距离，以及三人各自行走了多少天？}



\vspace{0.5em}



\bilingualsection{%

\Awen 答曰：从军前到都下的距离为三千三百里。甲行十一天；乙行十三天四时；丙行十八天二时。

}{%

\Amodern 答案：从军营前哨到都城的距离为3,300里。甲行走了11天；乙行走了13天4小时；丙行走了18天2小时。}



\vspace{0.5em}



\bilingualsection{%

\Swen 术曰：以大衍求之。置三人日行里数，连环求等，约奇弗约偶，得定母。诸定相乘为衍母。以各定约衍母得衍数。各满定母去之，余为奇数。大衍求一入之，得乘率。以乘衍数为用数。次置各余数乘用数，并之为总数。满衍母去之，不满为所求距离。

}{%

\Smodern 方法：用大衍法求解。列出三人每天行走的里数，两两配对求等数，约定奇数不约偶数，得到定母。各定母相乘得到衍母。用定母去除衍母得到衍数。用定母去除衍数，余数为奇数。用大衍求一术求得乘率。乘率乘以衍数得到用数。然后将各余数乘以用数，相加得到总数。用衍母去除总数，余数就是所求的距离。}



\vspace{0.5em}



\bilingualsection{%

\Cwen 置甲三百里、乙二百四十里、丙一百八十里，求总等，连环求等。



计算得乘率：甲乘率四，乙乘率一，丙乘率七。



通过乘率乘以各自的用数，得到总用数，再通过总用数除以衍母，得到最终结果。



验证：$3300 \div 300 = 11$天；$3300 \div 240 = 13$天余一百二十里 $= 13$天四时；$3300 \div 180 = 18$天余六十里 $= 18$天二时。合问。

}{%

\Cmodern 列出甲300里、乙240里、丙180里，求总等数，连环求等数。



计算得到乘率：甲乘率$k_1 = 4$，乙乘率$k_2 = 1$，丙乘率$k_3 = 7$。



用乘率乘以各自的衍数得到用数，再通过总数除以衍母得到最终结果。



验证：$3,300 \div 300 = 11$天；$3,300 \div 240 = 13$天余120里 $= 13$天4小时（一天12时辰，120里对应半天即4小时）；$3,300 \div 180 = 18$天余60里 $= 18$天2小时。与题意相合。}



\newpage



%% ========================================================================

%% PROBLEM 7: CATCHING UP ON A JOURNEY

%% ========================================================================

\section{第七题：程行相及}



\bilingualsection{%

\Qwen 问：甲、乙、丙三人行程相关。甲先行若干日，乙后行，丙又后行。各以不同速度行进。求某处三人相遇之时间及各地距离。

}{%

\Qmodern 问题：甲、乙、丙三人行程相关联。甲先出发若干天，乙随后出发，丙最后出发。三人各自以不同的速度行进。求在某处三人同时到达的时间以及各地之间的距离。



\textit{本题涉及多名旅行者以不同速度和不同出发时间行进，需要用大衍法求解一组同余式，以找到三人相遇的时间和距离。}}



\vspace{0.5em}



\bilingualsection{%

\Awen 答曰：所求距离及相关数值若干。

}{%

\Amodern 答案：所求距离及相关数值为特定值。（具体数值根据同余条件计算得出。）}



\vspace{0.5em}



\bilingualsection{%

\Swen 术曰：以大衍求之。置各人行程速度及时间差，连环求等，约奇弗约偶，得定母。诸定相乘为衍母。以各定约衍母得衍数。各满定母去之，余为奇数。大衍求一入之，得乘率。以乘衍数为用数。次置各余数乘用数，并之为总数。满衍母去之，不满为所求。

}{%

\Smodern 方法：用大衍法求解。列出各人行程的速度和时间差，两两配对求等数，约定奇数不约偶数，得到定母。各定母相乘得到衍母。用定母去除衍母得到衍数。用定母去除衍数，余数为奇数。用大衍求一术求得乘率。乘率乘以衍数得到用数。然后将各余数乘以用数，相加得到总数。用衍母去除总数，余数就是所求的答案。



\textit{本题为典型的多人相遇问题，将各人速度和时间差转化为同余条件，用大衍法的标准步骤求解。}}



\newpage



%% ========================================================================

%% PROBLEM 8: FINDING DIMENSIONS FROM VOLUME

%% ========================================================================

\section{第八题：积尺寻源}



\bilingualsection{%

\Qwen 问欲砌基一段，见管大小方砖六门城砖四色。令匠取便，或平或侧，只用一色砖砌，须要适足。匠以砖量地计料，称用：大方广多六寸，深少六寸；小方广多二寸，深少三寸；城砖长广多三寸，深少一寸；六门砖长广多三寸，深多一寸。皆不合匝，未免修破砖料裨补。



其四色砖尺寸：大方方一尺三寸，小方方一尺一寸，城砖长一尺二寸、阔六寸、厚二寸五分，六门砖长一尺、阔五寸、厚二寸。



欲知基深广几何？

}{%

\Qmodern 问题：要砌一段地基，现有大方砖、小方砖、六门砖和城砖四种砖。让工匠随意选用，或平放或侧放，只用一种砖砌，要求刚好砌满。工匠用砖来量地基并计算用料，说：用大方砖量，宽度多6寸，深度少6寸；用小方砖量，宽度多2寸，深度少3寸；用城砖量，长度宽度多3寸，深度少1寸；用六门砖量，长度宽度多3寸，深度多1寸。都不能恰好砌满，不免要修整破砖来补贴。



四种砖的尺寸为：大方砖边长1尺3寸，小方砖边长1尺1寸，城砖长1尺2寸、宽6寸、厚2.5寸，六门砖长1尺、宽5寸、厚2寸。



求地基的深度和宽度各是多少？}



\vspace{0.5em}



\bilingualsection{%

\Awen 答曰：深三丈七尺一寸，广一丈二尺三寸。

}{%

\Amodern 答案：地基深度为3丈7尺1寸，宽度为1丈2尺3寸。}



\vspace{0.5em}



\bilingualsection{%

\Swen 术曰：以大衍求之。置四色砖尺寸及所余尺寸，连环求等，约奇弗约偶，得定母。诸定相乘为衍母。以各定约衍母得衍数。各满定母去之，余为奇数。大衍求一入之，得乘率。以乘衍数为用数。次置各余数乘用数，并之为总数。满衍母去之，不满为基之深广。

}{%

\Smodern 方法：用大衍法求解。列出四种砖的尺寸以及用砖量地基时的剩余尺寸，两两配对求等数，约定奇数不约偶数，得到定母。各定母相乘得到衍母。用定母去除衍母得到衍数。用定母去除衍数，余数为奇数。用大衍求一术求得乘率。乘率乘以衍数得到用数。然后将各余数乘以用数，相加得到总数。用衍母去除总数，余数就是地基的深度和宽度。



\textit{本题是典型的「已知余数求总数」问题。地基的尺寸必须满足：用每种砖以不同方式铺设时，都有特定的余数。大衍法用于找到满足所有这些条件的尺寸。}}



\newpage



%% ========================================================================

%% PROBLEM 9: CALCULATING RICE REMAINDERS

%% ========================================================================

\section{第九题：余米推数}



\bilingualsection{%

\Qwen 问：有米若干，以不同容量的容器量之，各余若干。欲求米之总数。

}{%

\Qmodern 问题：有一批米，用不同容量的容器去量，各剩余不同的数量。求这批米的总数。



\textit{本题即经典的「物不知数」问题：用不同的模数去除同一数，得到不同的余数，求这个数。这就是直接导出中国剩余定理的问题。}}



\vspace{0.5em}



\bilingualsection{%

\Awen 答曰：所求米总数若干。

}{%

\Amodern 答案：所求米的总数为某特定数值。（具体数值根据同余条件计算得出。）}



\vspace{0.5em}



\bilingualsection{%

\Swen 术曰：以大衍求之。置各容器容量为问数，连环求等，约奇弗约偶，得定母。诸定相乘为衍母。以各定约衍母得衍数。各满定母去之，余为奇数。大衍求一入之，得乘率。以乘衍数为用数。次置各余数乘用数，并之为总数。满衍母去之，不满为所求米数。

}{%

\Smodern 方法：用大衍法求解。列出各容器的容量作为问数，两两配对求等数，约定奇数不约偶数，得到定母（两两互素的因数）。各定母相乘得到衍母（所有定数的乘积）。用定母去除衍母得到衍数。用定母去除衍数，余数为奇数（模余数）。用大衍求一术（求模逆元的方法）求得乘率。乘率乘以衍数得到用数。然后将各余数乘以对应的用数，相加得到总数。用衍母去除总数，余数就是所求的米数。}



\vspace{0.5em}



\bilingualsection{%

\textit{本题与《孙子算经》中著名的「物不知数」问题结构相同：「今有物不知其数，三三数之剩二，五五数之剩三，七七数之剩二，问物几何？」答曰二十三。秦九韶的通法系统性地解决了所有此类问题。}

}{%

\textit{本问题在结构上与《孙子算经》中著名的「物不知数」问题完全相同：「有一些物品，三个三个数余2，五个五个数余3，七个七个数余2，问有多少物品？」答案是23。秦九韶的通用方法系统地解决了所有这类问题。}}



\newpage



%% ========================================================================

%% MATHEMATICAL ANALYSIS

%% ========================================================================

\section{大衍法数学分析}



\subsection{中国剩余定理}



\bilingualsection{%

大衍法者，解线性同余式方程组之通法也。设 pairwise coprime 之模数 $m_1, m_2, \ldots, m_n$ 及余数 $r_1, r_2, \ldots, r_n$，求一数 $N$，使得：

\begin{align*}

N &\equiv r_1 \pmod{m_1} \\

N &\equiv r_2 \pmod{m_2} \\

  &\vdots \\

N &\equiv r_n \pmod{m_n}

\end{align*}

}{%

大衍法是求解线性同余方程组的通用方法。给定一组两两互素的模数 $m_1, m_2, \ldots, m_n$ 和对应的余数 $r_1, r_2, \ldots, r_n$，要求找到一个数 $N$，使得：

\begin{align*}

N &\equiv r_1 \pmod{m_1} \\

N &\equiv r_2 \pmod{m_2} \\

  &\vdots \\

N &\equiv r_n \pmod{m_n}

\end{align*}

即 $N$ 除以 $m_1$ 余 $r_1$，除以 $m_2$ 余 $r_2$，……，除以 $m_n$ 余 $r_n$。}



\subsection{秦九韶算法}



\bilingualsection{%

\textbf{第一步：求定数}。置元数，两两连环求等，约奇弗约偶。务使诸数 pairwise coprime，是为定母。



\textbf{第二步：求衍母}。诸定相乘：$M = m_1 \times m_2 \times \cdots \times m_n$。



\textbf{第三步：求衍数}。各定约衍母：$M_i = M / m_i$。



\textbf{第四步：求奇数}。各衍数满定母去之：$g_i = M_i \bmod m_i$。



\textbf{第五步：大衍求一术}。求乘率 $k_i$，使得：$k_i \cdot g_i \equiv 1 \pmod{m_i}$。此即求乘逆，用辗转相除法配以回代。



\textbf{第六步：求用数}。$u_i = k_i \cdot M_i$。



\textbf{第七步：总成}。所求之数为：$N = \sum_{i=1}^{n} r_i \cdot u_i \pmod{M}$。

}{%

\textbf{第一步：求定数（两两互素的因数）。}给定元数（原始的模数），两两配对求最大公约数（等数），约定奇数（位数为奇数的数）不约偶数（位数为偶数的数）。目标是使各数两两互素，这样得到的数就是定母。



\textbf{第二步：求衍母（所有定数的乘积）。}将各定母相乘：$M = m_1 \times m_2 \times \cdots \times m_n$。



\textbf{第三步：求衍数（衍母除以各定数）。}用各定母去除衍母：$M_i = M / m_i$。



\textbf{第四步：求奇数（模运算中的余数）。}用各定母去除衍数取余：$g_i = M_i \bmod m_i$。



\textbf{第五步：大衍求一术（求模逆元）。}求乘率 $k_i$，使得：$k_i \cdot g_i \equiv 1 \pmod{m_i}$。即求 $g_i$ 关于模 $m_i$ 的模逆元，用扩展欧几里得算法（辗转相除加回代）来求。



\textbf{第六步：求用数。}$u_i = k_i \cdot M_i$。



\textbf{第七步：最终计算。}所求的答案为：$N = \sum_{i=1}^{n} r_i \cdot u_i \pmod{M}$。即各余数乘以对应用数的和，再对衍母取余。}



\subsection{术语对照表}



\bilingualsection{%

\begin{center}

\textbf{大衍法术语表}\\[0.5em]

\begin{tabular}{ll}

\toprule

\textbf{术语} & \textbf{释义} \\

\midrule

大衍 & 《周易》大衍之数 \\

求一术 & 求乘逆之法 \\

元数 & 原问之数 \\

定数 & 两两互素之约数 \\

衍母 & 诸定数之连乘积 \\

衍数 & 衍母除以各定数之商 \\

奇数 & 衍数除以定数之余 \\

乘率 & 使乘奇数得一之数 \\

用数 & 乘率乘以衍数 \\

总数 & 各余乘用数之和 \\

\bottomrule

\end{tabular}

\end{center}

}{%

\begin{center}

\textbf{大衍法术语现代数学对应}\\[0.5em]

\begin{tabular}{ll}

\toprule

\textbf{古典术语} & \textbf{现代数学概念} \\

\midrule

大衍 & 基于《周易》的中国剩余定理 \\

求一术 & 求模逆元的扩展欧几里得算法 \\

元数 & 原始模数 \\

定数 & 两两互素的简化模数 \\

衍母 & 所有定数的乘积 $M$ \\

衍数 & $M_i = M/m_i$ \\

奇数 & 模余数 $g_i = M_i \bmod m_i$ \\

乘率 & 模逆元 $k_i$（使$k_i g_i \equiv 1$） \\

用数 & $u_i = k_i M_i$ \\

总数 & $\sum r_i u_i$ \\

\bottomrule

\end{tabular}

\end{center}}



\vspace{1em}



\bilingualsection{%

\textit{秦九韶大衍法早于欧洲同类工作数百年。其求模乘逆之算法（大衍求一术）实为扩展欧几里得算法。此法后经郭守敬用于天文历法，经李冶用于几何学，并传入欧洲，后被称为「中国剩余定理」。}



\textit{高斯于1801年《算术研究》中首次给出欧洲之完整论述，称之为「求一数，以所给诸数除之，各给余数」之问题。该定理今被视为数论中最基本的结果之一。}

}{%

\textit{秦九韶的大衍法比欧洲同类工作早了数百年。求模逆元的算法（大衍求一术）本质上就是扩展欧几里得算法。这一方法后来被郭守敬应用于天文历法计算，被李冶应用于几何学，并传入欧洲，被称为「中国剩余定理」。}



\textit{高斯在1801年的《算术研究》中首次在欧洲给出了完整的论述，称之为「求一个数，使其除以给定各数后留下给定余数」的问题。这一定理如今被认为是数论中最基本的结果之一。}}



\newpage



%% ========================================================================

%% END MATTER

%% ========================================================================

\section*{关于本版}

\addcontentsline{toc}{section}{关于本版}



\bilingualsection{%

本对照版呈现秦九韶《数书九章》第一卷大衍类之九章问题，涵盖中国剩余定理及相关不定分析。



底本据《四库全书》本，参校宜稼堂丛书本。四库本所附按语保留于原处。

}{%

本对照版呈现秦九韶《数书九章》第一卷「大衍类」的九个问题，涵盖中国剩余定理及相关的不定分析方法。



底本依据《四库全书》版本，并参照宜稼堂丛书版本进行校勘。四库本所附的「按语」（编者按语）保留于适当位置。}



\vspace{1em}



\subsection*{技术说明}

\begin{itemize}[leftmargin=1.5em]

  \item 排版系统：XeTeX + \LaTeX

  \item 中文字体: modern Unicode fonts

  \item 中文处理：xeCJK宏包

  \item 双语对照：minipage并排环境

  \item 保留了传统的「问/答/术/草」四段结构

  \item 白话文翻译力求忠实原文，保留数学术语

\end{itemize}



\subsection*{参考文献}

\begin{itemize}[leftmargin=1.5em]

  \item Libbrecht, U. \textit{Chinese Mathematics in the Thirteenth Century}. MIT Press, 1973.（利布雷希特，《十三世纪的中国数学》）

  \item 钱宝琮.《中国数学史》. 科学出版社, 1964.

  \item Martzloff, J.-C. \textit{A History of Chinese Mathematics}. Springer, 1997.（马兹洛夫，《中国数学史》）

  \item 李俨.《中国数学大纲》. 商务印书馆, 1958.

\end{itemize}



\vfill

\begin{center}

\textcolor{rulecolor}{\rule{0.5\textwidth}{0.4pt}}\\[0.5em]

\textit{文言文原文 -- 白话文今译 对照版}\\[0.3em]

\textit{用\LaTeX{}排版 --- \the\year}

\end{center}



\end{document}







% ============================================================

% Source: translations/zh/chinese/qin-jiushao-shuxue-jiuzhang-fascicles5-6_classical-modern_bilingual.tex

% ============================================================



%% =============================================================================

%% 數學九章 (Shuxue Jiuzhang) — Fascicles 5–6

%% BILINGUAL EDITION: Classical Chinese ←→ Modern Chinese Vernacular

%% Author: Qin Jiushao (秦九韶), Song Dynasty

%% Source: Siku Quanshu (四庫全書) Edition

%% Modern LaTeX edition%% =============================================================================

\documentclass[11pt,a4paper]{article}



%% -------- Core Packages --------

\usepackage{fontspec}

\usepackage{polyglossia}

\usepackage{amsmath,amssymb}

\usepackage{geometry}

\usepackage{graphicx}

\usepackage{xcolor}

\usepackage{titlesec}

\usepackage{fancyhdr}

\usepackage{enumitem}

\usepackage{array,booktabs,longtable,multirow}

\usepackage{hyperref}

\usepackage{tocloft}

\usepackage{indentfirst}

\usepackage{xeCJK}

\usepackage{paracol}

\usepackage{etoolbox}



%% -------- Page Geometry --------

\geometry{

  paperwidth=210mm,paperheight=297mm,

  left=18mm,right=18mm,top=25mm,bottom=25mm,

  headheight=14pt,footskip=30pt

}



%% -------- Language Setup --------

\setmainlanguage{english}

\setotherlanguage{chinese}



%% -------- Font Configuration --------

\setmainfont{Times New Roman}

\setsansfont{Arial}



%% Chinese Fonts

\setCJKmainfont{Microsoft YaHei}

\setCJKsansfont{Microsoft YaHei}

\setCJKmonofont{Microsoft YaHei}



%% -------- Colors --------

\definecolor{titlecolor}{RGB}{45,55,72}

\definecolor{sectioncolor}{RGB}{55,65,81}

\definecolor{notecolor}{RGB}{120,53,15}

\definecolor{rulecolor}{RGB}{180,160,140}

\definecolor{problemcolor}{RGB}{80,40,20}

\definecolor{answercolor}{RGB}{20,80,40}

\definecolor{methodcolor}{RGB}{20,40,80}

\definecolor{classicalbg}{RGB}{255,252,245}

\definecolor{modernbg}{RGB}{245,250,255}

\definecolor{modernlabel}{RGB}{120,60,20}

\definecolor{headerclassical}{RGB}{139,69,19}

\definecolor{headermodern}{RGB}{25,100,140}



%% -------- Column settings for paracol --------

\columnsep=10mm

\globalcounter{table}



%% -------- Custom Column Backgrounds --------

\backgroundcolor{c[0](0.5\columnsep,0.5\columnsep)}{classicalbg}

\backgroundcolor{c[1](0.5\columnsep,0.5\columnsep)}{modernbg}



%% -------- Section Formatting --------

\titleformat{\section}

  {\normalfont\Large\bfseries\color{sectioncolor}}

  {\thesection}{1em}{}

\titleformat{\subsection}

  {\normalfont\large\bfseries\color{sectioncolor}}

  {\thesubsection}{1em}{}

\titleformat{\subsubsection}

  {\normalfont\normalsize\bfseries\color{problemcolor}}

  {\thesubsubsection}{1em}{}



%% -------- Header/Footer --------

\pagestyle{fancy}

\fancyhf{}

\fancyhead[L]{\small\textsc{數學九章 — 卷五~\&~卷六 \;|\; 文言文 \;\raisebox{-0.5pt}{$\leftrightarrow$}\; 白話文}}

\fancyhead[R]{\small\thepage}

\renewcommand{\headrulewidth}{0.4pt}

\renewcommand{\footrulewidth}{0pt}



%% -------- Custom Commands --------

\newcommand{\longline}{\noindent\rule{\textwidth}{0.4pt}}

\newcommand{\shortline}{\noindent\rule{0.5\textwidth}{0.4pt}\par}

\newcommand{\cnote}[1]{\textcolor{notecolor}{\textit{#1}}}



%% Classical Chinese problem structure commands

\newcommand{\ClassicalQuestion}[1]{\par\noindent\textbf{\color{problemcolor}問：}#1\par}

\newcommand{\ClassicalAnswer}[1]{\par\noindent\textbf{\color{answercolor}答曰：}#1\par}

\newcommand{\ClassicalMethod}[1]{\par\noindent\textbf{\color{methodcolor}術曰：}#1\par}

\newcommand{\ClassicalWorking}[1]{\par\noindent\textbf{草曰：}#1\par}

\newcommand{\ClassicalNote}[1]{\par\noindent\textit{［按］#1}\par}



%% Modern Chinese problem structure commands

\newcommand{\ModernQuestion}[1]{\par\noindent\textbf{\color{modernlabel}【题问】}#1\par}

\newcommand{\ModernAnswer}[1]{\par\noindent\textbf{\color{answercolor}【答案】}#1\par}

\newcommand{\ModernMethod}[1]{\par\noindent\textbf{\color{methodcolor}【算法】}#1\par}

\newcommand{\ModernWorking}[1]{\par\noindent\textbf{【演算】}#1\par}

\newcommand{\ModernNote}[1]{\par\noindent\textit{\color{notecolor}［编者按］#1}\par}



%% Column headers

\newcommand{\ClassicalHeader}{%

  \begin{center}\textbf{\color{headerclassical}【原文 · 文言文】}\end{center}

  \vspace{-0.5em}

  \hrule height 0.8pt

  \vspace{0.5em}

}

\newcommand{\ModernHeader}{%

  \begin{center}\textbf{\color{headermodern}【译文 · 白话文】}\end{center}

  \vspace{-0.5em}

  \hrule height 0.8pt

  \vspace{0.5em}

}



%% -------- Hyperref Setup --------

\hypersetup{

  colorlinks=true,

  linkcolor=sectioncolor,

  citecolor=sectioncolor,

  urlcolor=sectioncolor,

  pdfencoding=auto,

  pdftitle={數學九章 卷五至六 文言文-白話文對照版},

  pdfauthor={Qin Jiushao (秦九韶)},

  pdfsubject={Chinese Mathematics -- Bilingual Classical/Modern Edition}

}



%% -------- TOC Depth --------

\setcounter{tocdepth}{3}

\setcounter{secnumdepth}{3}



%% Reduce paragraph spacing for dense Chinese text

\setlength{\parskip}{0.5ex plus 0.2ex minus 0.1ex}

\setlength{\parindent}{2em}



%% Suppress paracol column swapping page breaks in certain contexts

\renewcommand{\rownumber}{}



%% =============================================================================

\begin{document}



%% -------- Title Page --------

\begin{titlepage}

\centering

\vspace*{2.5cm}



{\fontsize{36}{44}\selectfont\color{titlecolor}\CJKfamily{}\textbf{數學九章}}



\vspace{1.2cm}



{\fontsize{16}{20}\selectfont\color{sectioncolor} 卷五~\&~卷六 · 文言文~$\longleftrightarrow$~白話文對照版}



\vspace{0.8cm}



{\large\color{sectioncolor}

卷五：\textbf{賦役}（賦稅和徭役）\quad|\quad

卷六：\textbf{錢穀}（貨幣和糧食）}



\vspace{2.5cm}



{\large 宋\quad 秦九韶\quad 撰}



\vspace{0.5cm}



{\normalsize\textit{Qin Jiushao (c.\ 1202--1261)}}



\vspace{2cm}



{\color{rulecolor}\rule{0.6\textwidth}{0.4pt}}



\vspace{1.5cm}



{\small\color{sectioncolor}

\textit{Siku Quanshu}（四庫全書）Edition\\[0.5em]

文言文原文 · 白話文譯文對照排版\\[0.3em]

Digital Typesetting\quad|\quad Bilingual Edition}



\vspace{2.5cm}



{\small\color{notecolor}

\textit{本書保留中國古代數學典籍傳統結構：\\

問（題問）、答（答案）、術（算法）、草（演算）}}



\vfill

{\small\color{sectioncolor} Compiled with XeLaTeX · Microsoft YaHei}

\end{titlepage}



%% -------- Table of Contents --------

\tableofcontents

\newpage







%% =============================================================================

%% FASCICLE 5A — 賦役 (Part 1)

%% =============================================================================

\newpage

\section*{\color{titlecolor}卷五上\quad 賦役（賦稅和徭役）}

\addcontentsline{toc}{section}{卷五上\quad 賦役}



\begin{paracol}{2}



\ClassicalHeader



\noindent\textbf{按：}賦役一章，即古九章中差分粟米諸法。但從來算書設題之數未有如此卷之多者。如首題答數至一百七十五條，毎條歩算之數至十餘位，而得數皆無不合，亦可謂能廣其用矣。且諸問皆足以明貴賤廩税及交質變易之同異，使公私各得其平，誠治民之要務也。但題語多晦，術草中間有與所問不相應者，亦於各條下指出，俾觀者無惑焉。



\switchcolumn



\ModernHeader



\noindent\textbf{【按語】}賦役這一章，就是古代《九章算術》中「差分」「粟米」等各種算法的擴展。但是從來算書中的設題數量，沒有像這一卷這麼多的。比如第一題的答案多達一百七十五條，每條計算的數字達到十幾位，而且所得結果都沒有不吻合的，也可以說是極大地擴展了這些算法的應用範圍。況且各題都足以闡明貴賤、廩税以及交換質易的異同，使公私各得其平，確實是治理百姓的重要事務。只是題目語言多晦澀，術草中間有與所問不相應的地方，也在各條下面指出，使讀者不會迷惑。



\end{paracol}



\longline



\subsection*{\color{sectioncolor}復邑修賦（恢復縣城修訂賦稅）}

\addcontentsline{toc}{subsection}{復邑修賦}



\begin{paracol}{2}



\ClassicalHeader



\ClassicalQuestion{%

有海坍縣地，今有復漲，歳乆鄉井再成，申諸剏邑。稱土排到六鄉，以附郭為甲，最逺為已，各有田九等，開具下項：



\begin{longtable}{p{0.85\textwidth}}

\toprule

\textbf{甲鄉}共計田十四萬一百九十三畝三角一十二步\\

\midrule

\endhead

上等：上田五千六百七十八畝一角四十八歩；中田四千八百九十二畝三十歩；下田六千六百二十一畝五十四歩\\

中等：上田八千二百二十五畝二十四歩；中田一萬三十五畝六歩；下田一萬六千五百三十畝\\

下等：上田二萬一千九十畝二十四歩；中田三萬二千六十畝三歩；下田三萬五千六十一畝三角三歩\\

\bottomrule

\end{longtable}



\begin{longtable}{p{0.85\textwidth}}

\toprule

\textbf{乙鄉}田共計八萬四千一十畝二角二歩\\

\midrule

\endhead

上等：上田六千七百八十九畝一角三十六歩；中田五千九百八十七畝二角；下田八千一十畝三角三歩\\

中等：上田七千五百四十一畝；中田九千一百二十一畝二角一十二歩；下田一萬九千六十歩\\

下等：上田一萬八千三十七畝一角六歩；中田九千四百五十六畝三角五十九歩；下田無\\

\bottomrule

\end{longtable}



\begin{longtable}{p{0.85\textwidth}}

\toprule

\textbf{丙鄉}田共計一十二萬九百三十五畝五十八歩五分\\

\midrule

\endhead

上等：上田四千八百六十八畝二角三歩；中田五千九百七十九畝三角六歩；下田六千八百八十八畝二角六歩\\

中等：上田七千九百八十四畝一角；中田一萬四千五十六畝一十二歩；下田二萬三千三百三十三畝一十二歩\\

下等：上田二萬七千七百五十五畝一十六歩五分；中田無；下田三萬七十畝三歩\\

\bottomrule

\end{longtable}



\begin{longtable}{p{0.85\textwidth}}

\toprule

\textbf{丁鄉}田共計八萬九千六十六畝二歩三分\\

\midrule

\endhead

上等：上田一萬一千一百二畝一歩；中田九千八百七十六畝一角；下田八千七百六十五畝一角三十歩\\

中等：上田七千五百三十九畝三十四歩三分；中田無；下田一萬二千九百八十七畝四十二歩\\

下等：上田無；中田五千四百三十二畝一角六歩；下田三萬三千三百六十三畝三角九歩\\

\bottomrule

\end{longtable}



\begin{longtable}{p{0.85\textwidth}}

\toprule

\textbf{戊鄉}田共計二十萬四千四百七十四畝一角二十四歩四分\\

\midrule

\endhead

上等：上田二萬四千六百三十二畝三十九歩；中田一萬三千五百二十一畝二十七歩；下田九千九百八十八畝三角三歩\\

中等：上田八千八百七十七畝五十六歩四分；中田一萬一千三百三十三畝三角；下田二萬七千六十七畝\\

下等：上田一萬九千八百七十六畝三角六歩；中等田七萬九千一百三十五畝三角四十三歩；下田一萬四十一畝二角三十歩\\

\bottomrule

\end{longtable}



\begin{longtable}{p{0.85\textwidth}}

\toprule

\textbf{已鄉}田共計一十五萬八千四百六十畝三角一十八歩二分\\

\midrule

\endhead

上等：上田無；中田七千七百八十八畝三角五十一歩；下田無\\

中等：上田九千九百九十九畝二角三歩；中田一萬八百三十六畝五十六歩；下田無\\

下等：上田三萬二千八十九畝一角四十五歩六分；中田四萬三千六百七十八畝二角五十七歩；下田五萬四千六十七畝三角四十五歩六分\\

\bottomrule

\end{longtable}



照得昨來本縣原科苗米一十萬三千五百六十七石八斗四升四合三勺，和買一萬三千四百九十八匹一丈七尺三寸七分六釐，夏税九千八百七十六匹三丈二尺六寸五分八釐。其六鄉田係三色，甲為上，乙丙為次，丁戊己又為次。先令官物為三差，使上比中、中比下皆十分外差一，次令各鄉九等皆於十分内差一，抛科用合租額。其乙鄉田最肥，次丁，次甲，次丙，次己，次戊。欲知三色九等毎畝等則及共科數各鄉幾何？

}



\switchcolumn



\ModernHeader



\ModernQuestion{%

有一個因海水侵蝕而坍廢的縣地，現在海水退去土地復漲，年久鄉里重新形成，申請設立新縣。經丈量土地分為六個鄉，以靠近城郭的為甲鄉，最遠的為己鄉，每個鄉都有九等田地，開列如下：



\textbf{【甲鄉】}共計田十四萬零一百九十三畝零三十二步（一角=10畝的1/10，約30步）



\begin{itemize}[leftmargin=1.5em,topsep=0pt,itemsep=1pt]

\item 上等：上田5678畝148步；中田4892畝30步；下田6621畝54步

\item 中等：上田8225畝24步；中田10035畝6步；下田16530畝

\item 下等：上田21090畝24步；中田32060畝3步；下田35061畝33步

\end{itemize}



\textbf{【乙鄉】}田共計84010畝22步



\begin{itemize}[leftmargin=1.5em,topsep=0pt,itemsep=1pt]

\item 上等：上田6789畝136步；中田5987畝20步；下田8010畝33步

\item 中等：上田7541畝；中田9121畝22步；下田19060步

\item 下等：上田18037畝16步；中田9456畝59步；下田無

\end{itemize}



（丙鄉、丁鄉、戊鄉、己鄉數據同上，略）



查得原來本縣徵收苗米103,567石8斗4升4合3勺，和買（徵購絹帛）13,498匹1丈7尺3寸7分6釐，夏税9876匹3丈2尺6寸5分8釐。六鄉田地分三等，甲為上等，乙丙為次等，丁戊己又為次等。先將官物分為三個等差，使上等比中等、中等比下等都按十分之外差一（即1.1倍），再令各鄉九等都按十分之內差一（即0.9倍），按比例分攤配額。乙鄉田地最肥，其次丁、甲、丙、己、戊。求三等九等每畝的科率以及各鄉共徵收的數量各為多少？

}



\end{paracol}







\begin{paracol}{2}



\ClassicalNote{%

題内言田有三色，甲為上云云，又言乙鄉田最肥云云，語若相左。葢前言三色者六鄉共科之等，後言田肥者每畝所出之異也。若易田係三色為科有三差，則語意明矣。}



\switchcolumn



\ModernNote{%

題目內說田地有三色，甲為上等，又說乙鄉田地最肥，語言似乎矛盾。原來前面說的三色是六鄉共分徵收的三個等級，後面說的田地肥瘦是每畝產量的差異。如果把「田係三色」改為「科有三差」，語意就清楚了。}



\end{paracol}



\begin{paracol}{2}



\ClassicalAnswer{%

\textbf{甲鄉：}上等上田苗三斗二升三合二勺，和一尺六寸九分，税一尺二寸三分；

中田苗二斗九升九勺，和一尺五寸二分，税一尺一寸一分；

下田苗二斗五升八合六勺，和一尺三寸五分，税九寸九分。



中等上田苗二斗二升六合二勺，和一尺一寸八分，税八寸六分；

中田苗一斗九升三合九勺，和一尺一分，税七寸四分；

下田苗一斗六升一合六勺，和八寸四分，税六寸二分。



下等上田苗一斗二升九合三勺，和六寸七分，税四寸九分；

中田苗九升六合九勺，和五寸一分，税三寸七分；

下田苗六升四合六勺，和三寸四分，税二寸五分。



\textbf{乙鄉：}上等上田苗三斗六升三合三勺，和一尺八寸九分，税一尺三寸九分；

中田苗三斗二升六合九勺，和一尺七寸，税一尺二寸五分；

下田苗二斗九升六勺，和一尺五寸二分，税一尺一寸一分。



中等上田苗二斗五升四合三勺，和一尺三寸三分，税九寸七分；

中田苗二斗一升八合，和一尺一寸四分，税八寸三分；

下田苗一斗八升一合六勺，和九寸五分，税六寸九分。



下等上田苗一斗四升五合三勺，和七寸六分，税五寸五分；

中田苗一斗九合，和五寸七分，税四寸二分；

下田苗七升二合七勺，和三寸八分，税二寸八分。

}



\switchcolumn



\ModernAnswer{%

\textbf{【甲鄉】}上等上田每畝徵苗米3斗2升3合2勺，和買（配額絹）1尺6寸9分，夏税1尺2寸3分；

上等中田苗米2斗9升9勺，和買1尺5寸2分，税1尺1寸1分；

上等下田苗米2斗5升8合6勺，和買1尺3寸5分，税9寸9分。



中等上田苗米2斗2升6合2勺，和買1尺1寸8分，税8寸6分；

中等中田苗米1斗9升3合9勺，和買1尺1分，税7寸4分；

中等下田苗米1斗6升1合6勺，和買8寸4分，税6寸2分。



下等上田苗米1斗2升9合3勺，和買6寸7分，税4寸9分；

下等中田苗米9升6合9勺，和買5寸1分，税3寸7分；

下等下田苗米6升4合6勺，和買3寸4分，税2寸5分。



\textbf{【乙鄉】}上等上田苗米3斗6升3合3勺，和買1尺8寸9分，税1尺3寸9分；

上等中田苗米3斗2升6合9勺，和買1尺7寸，税1尺2寸5分；

上等下田苗米2斗9升6勺，和買1尺5寸2分，税1尺1寸1分。



中等上田苗米2斗5升4合3勺，和買1尺3寸3分，税9寸7分；

中等中田苗米2斗1升8合，和買1尺1寸4分，税8寸3分；

中等下田苗米1斗8升1合6勺，和買9寸5分，税6寸9分。



下等上田苗米1斗4升5合3勺，和買7寸6分，税5寸5分；

下等中田苗米1斗9合，和買5寸7分，税4寸2分；

下等下田苗米7升2合7勺，和買3寸8分，税2寸8分。

}



\end{paracol}



\begin{paracol}{2}



\ClassicalAnswer{%

\textbf{丙鄉：}上等上田苗三斗三合五勺，和一尺五寸八分，税一尺一寸六分；

中田苗二斗七升三合一勺，和一尺四寸二分，税一尺四分；

下田苗二斗四升二合八勺，和一尺二寸七分，税九寸三分。



中等上田苗二斗一升二合四勺，和一尺一寸一分，税八寸一分；

中田苗一斗八升二合一勺，和九寸五分，税六寸九分；

下田苗一斗五升一合七勺，和七寸九分，税五寸八分。



下等上田苗一斗二升一合四勺，和六寸三分，税四寸六分；

中田苗九升一合，和四寸七分，税三寸五分；

下田苗六升七勺，和三寸二分，税二寸三分。



\textbf{丁鄉：}上等上田苗三斗四升三合二勺，和一尺七寸九分，税一尺三寸一分；

中田苗三斗八合九勺，和一尺六寸一分，税一尺一寸八分；

下田苗二斗七升四合六勺，和一尺四寸三分，税一尺五分。



中等上田苗二斗四升二勺，和一尺二寸五分，税九寸二分；

中田苗二斗五合九勺，和一尺七分，税七寸九分；

下田苗一斗七升一合六勺，和八寸九分，税六寸五分。



下等上田苗一斗三升七合三勺，和七寸二分，税五寸二分；

中田苗一斗三合，和五寸四分，税三寸九分；

下田苗六升八合六勺，和三寸六分，税二寸六分。



\textbf{戊鄉：}上等上田苗一斗五升三合八勺，和八寸，税五寸九分；

中田苗一斗三升八合四勺，和七寸二分，税五寸三分；

下田苗一斗二升三合一勺，和六寸四分，税四寸七分。



中等上田苗一斗七合七勺，和五寸六分，税四寸一分；

中田苗九升二合三勺，和四寸八分，税三寸五分；

下田苗七升六合九勺，和四寸，税二寸九分。



下等上田苗六升一合五勺，和三寸二分，税二寸三分；

中田苗四升六合一勺，和二寸四分，税一寸八分；

下田苗三升八勺，和一寸六分，税一寸二分。



\textbf{已鄉：}上等上田苗二斗八升二合二勺，和一尺四寸七分，税一尺八分；

中田苗二斗五升三合九勺，和一尺三寸二分，税九寸七分；

下田苗二斗二升五合七勺，和一尺一寸八分，税八寸六分。



中等上田苗一斗九升七合五勺，和一尺三分，税七寸五分；

中田苗一斗六升九合三勺，和八寸八分，税六寸五分；

下田苗一斗四升一合一勺，和七寸四分，税五寸四分。



下等上田苗一斗一升二合九勺，和五寸九分，税四寸三分；

中田苗八升四合六勺，和四寸四分，税三寸二分；

下田苗五升六合四勺，和二寸九分，税二寸二分。

}



\switchcolumn



\ModernAnswer{%

\textbf{【丙鄉】}（答案數據與甲鄉、乙鄉結構相同，等級依次遞減，略記主要數據）

上等上田苗米3斗3合5勺，和買1尺5寸8分，税1尺1寸6分；

上等中田苗米2斗7升3合1勺，和買1尺4寸2分，税1尺4分；

上等下田苗米2斗4升2合8勺，和買1尺2寸7分，税9寸3分。



中等上田苗米2斗1升2合4勺，和買1尺1寸1分，税8寸1分；

中等中田苗米1斗8升2合1勺，和買9寸5分，税6寸9分；

中等下田苗米1斗5升1合7勺，和買7寸9分，税5寸8分。



下等上田苗米1斗2升1合4勺，和買6寸3分，税4寸6分；

下等中田苗米9升1合，和買4寸7分，税3寸5分；

下等下田苗米6升7勺，和買3寸2分，税2寸3分。



\textbf{【丁鄉】}上等上田苗米3斗4升3合2勺，和買1尺7寸9分，税1尺3寸1分；

上等中田苗米3斗8合9勺，和買1尺6寸1分，税1尺1寸8分；

上等下田苗米2斗7升4合6勺，和買1尺4寸3分，税1尺5分。



中等上田苗米2斗4升2勺，和買1尺2寸5分，税9寸2分；

中等中田苗米2斗5合9勺，和買1尺7分，税7寸9分；

中等下田苗米1斗7升1合6勺，和買8寸9分，税6寸5分。



下等上田苗米1斗3升7合3勺，和買7寸2分，税5寸2分；

下等中田苗米1斗3合，和買5寸4分，税3寸9分；

下等下田苗米6升8合6勺，和買3寸6分，税2寸6分。



\textbf{【戊鄉】}上等上田苗米1斗5升3合8勺，和買8寸，税5寸9分；

上等中田苗米1斗3升8合4勺，和買7寸2分，税5寸3分；

上等下田苗米1斗2升3合1勺，和買6寸4分，税4寸7分。



中等上田苗米1斗7合7勺，和買5寸6分，税4寸1分；

中等中田苗米9升2合3勺，和買4寸8分，税3寸5分；

中等下田苗米7升6合9勺，和買4寸，税2寸9分。



下等上田苗米6升1合5勺，和買3寸2分，税2寸3分；

下等中田苗米4升6合1勺，和買2寸4分，税1寸8分；

下等下田苗米3升8勺，和買1寸6分，税1寸2分。



\textbf{【己鄉】}上等上田苗米2斗8升2合2勺，和買1尺4寸7分，税1尺8分；

上等中田苗米2斗5升3合9勺，和買1尺3寸2分，税9寸7分；

上等下田苗米2斗2升5合7勺，和買1尺1寸8分，税8寸6分。



中等上田苗米1斗9升7合5勺，和買1尺3分，税7寸5分；

中等中田苗米1斗6升9合3勺，和買8寸8分，税6寸5分；

中等下田苗米1斗4升1合1勺，和買7寸4分，税5寸4分。



下等上田苗米1斗1升2合9勺，和買5寸9分，税4寸3分；

下等中田苗米8升4合6勺，和買4寸4分，税3寸2分；

下等下田苗米5升6合4勺，和買2寸9分，税2寸2分。

}



\end{paracol}



\begin{paracol}{2}



\ClassicalAnswer{%

\textbf{各鄉共科數：}

甲鄉苗米一萬九千五百五十石二斗四升八合三勺；

和買一萬一百九十二丈二尺六寸五分六釐；

夏税七千四百五十七丈六尺八寸九分八釐。



乙鄉苗米一萬七千七百七十二石九斗五升三合；

和買九千二百六十五丈六尺九寸六分；

夏税六千七百七十九丈七尺一寸八分。



丙鄉苗米一萬七千七百七十二石九斗五升三合；

和買九千二百六十五丈六尺九寸六分；

夏税六千七百七十九丈七尺一寸八分。



丁鄉苗米一萬六千一百五十七石二斗三升；

和買八千四百二十三丈三尺六寸；

夏税六千一百六十三丈三尺八寸。



戊鄉苗米一萬六千一百五十七石二斗三升；

和買八千四百二十三丈三尺六寸；

夏税六千一百六十三丈三尺八寸。



已鄉苗米一萬六千一百五十七石二斗三升；

和買八千四百二十三丈三尺六寸；

夏税六千一百六十三丈三尺八寸。

}



\switchcolumn



\ModernAnswer{%

\textbf{【各鄉共徵收數量】：}

甲鄉：苗米19,550石2斗4升8合3勺；和買10,192丈2尺6寸5分6釐；夏税7457丈6尺8寸9分8釐。



乙鄉：苗米17,772石9斗5升3合；和買9265丈6尺9寸6分；夏税6779丈7尺1寸8分。



丙鄉：苗米17,772石9斗5升3合；和買9265丈6尺9寸6分；夏税6779丈7尺1寸8分。



丁鄉：苗米16,157石2斗3升；和買8423丈3尺6寸；夏税6163丈3尺8寸。



戊鄉：苗米16,157石2斗3升；和買8423丈3尺6寸；夏税6163丈3尺8寸。



己鄉：苗米16,157石2斗3升；和買8423丈3尺6寸；夏税6163丈3尺8寸。

}



\end{paracol}



\begin{paracol}{2}



\ClassicalMethod{%

以衰分求之。先列本縣色位，自下錐行列之，又以鄉數對列而乗之，副併為法，以除官物數，得一分之率。以率數乗未併者，各得諸鄉之數。次列各鄉等位，自上等置十分，每以内分錐行九折之，至九等止，又各以畝歩乗之，副併為鄉法，以除諸各鄉所得官物數，所得為一分之率，以乗未併者，各得毎畝税色。}



\switchcolumn



\ModernMethod{%

使用衰分法（等差比例分配法）求解。先排列本縣的三個等級，自下而上按錐形（遞減數列）排列，再以各鄉的數量對應相乘，合併作為除數（法），去除官物總數，得到一分的比率。用比率數乘未合併的各項，得到各鄉應徵之數。再排列各鄉的九個等級，從上等開始設為十分，每次以內分錐形九折（乘以0.9），直到九等為止，再各以畝數步數相乘，合併作為鄉法，去除各鄉所得官物數，所得即為一分的比率，用來乘未合併的各項，得到每畝的税額。

}



\end{paracol}



\begin{paracol}{2}



\ClassicalWorking{%

列本縣色位三，自下色例十分中十一分，上比中身外加一，上得一百二十一，中得一百一十，下得一百，為上中下三率，列之右行。乃列甲一對上率，乙丙共二對中率，丁戊己共三對下率，列左行。乃以左行列數各相對乗右行率數，其上得一百二十一，其十得二百二十，其下得三百，乃副置而併之，得六百四十一為法。置元科苗米一萬三千五百六十七石八斗四升四合三勺為寔，以法除之，得一百六十一石五斗七升二合三勺為一分之率。以未併下率一百乗率，得一萬六千一百五十七石二斗三升為丁戊己三鄉各科數。以於身下加一，得一萬七千七百七十二石九斗五升三合為乙丙二鄉各科數。又於身下加一，得一萬九千五百五十石二斗四升八合三勺為甲合科數。求和買亦用六百四十一為法，置元科和買一萬三千四百九十八匹，以匹法四丈通之，得五萬三千九百九十二丈，内零一丈七尺三寸七分六釐，得五萬三千九百九十三丈七尺三寸七分六釐為寔，以法除之，得八十四丈二尺三寸三分六釐為一分之率。亦以未併下率一百乗之，得八千四百二十三丈三尺六寸為丁戊己三鄉各科數。次於身下加一，得九千二百六十五丈六尺九寸六分為乙丙二鄉各科數。又於身下加一，得一萬一百九十二丈二尺六寸五分六釐為甲鄉合科數。求夏税亦置六百四十一為法，置元科九千八百七十六匹，以匹法四丈通之，得三萬九千五百四丈，内零三丈二尺六寸五分八釐，得三萬九千五百七丈二尺六寸五分八釐為寔，以法除之，得六十一丈六尺三寸三分八釐為一分率。亦以未併下率一百乗之，得六千一百六十三丈三尺八寸為丁戊己三鄉各科數。次於身下加一，得六千七百七十九丈七尺一寸八分为乙丙二鄉各科數。又於身下加一，得七千四百五十七丈六尺八寸九分八釐為甲鄉數。}



\switchcolumn



\ModernWorking{%

排列本縣的三個等級色位，自下而上，下等率為一百，中等為一百一十（身外加一即1.1倍），上等為一百二十一，作為上中下三個比率，排在右行。再將甲一鄉對應上率，乙丙共二鄉對應中率，丁戊己共三鄉對應下率，排在左行。以左行的數各與右行對應相乘：上得121，中得220，下得300，合併得641作為法（除數）。將原徵苗米數103,567石8斗4升4合3勺作為被除數（實），以法641去除，得161石5斗7升2合3勺為一分的比率。用未合併的下率100乘比率，得16,157石2斗3升為丁戊己三鄉各科數。再加一成（1.1倍），得17,772石9斗5升3合為乙丙二鄉各科數。再加一成，得19,550石2斗4升8合3勺為甲鄉合科數。



求和買也用641為法，將原和買13,498匹以匹法4丈換算，得53,992丈，加上零數1丈7尺3寸7分6釐，共53,993丈7尺3寸7分6釐為實，以法除之，得84丈2尺3寸3分6釐為一分之率。也用下率100乘之，得8,423丈3尺6寸為丁戊己三鄉各科數。再加一成，得9,265丈6尺9寸6分為乙丙二鄉各科數。再加一成，得10,192丈2尺6寸5分6釐為甲鄉合科數。



求夏税也用641為法，將原夏税9,876匹以4丈換算，得39,504丈，加零數3丈2尺6寸5分8釐，共39,507丈2尺6寸5分8釐為實，以法除之，得61丈6尺3寸3分8釐為一分之率。以下率100乘之，得6,163丈3尺8寸為丁戊己三鄉各科數。再加一成，得6,779丈7尺1寸8分為乙丙二鄉各科數。再加一成，得7,457丈6尺8寸9分8釐為甲鄉數。

}



\end{paracol}







\newpage

\subsection*{\color{sectioncolor}圍田租畝（圍田租賦計算）}

\addcontentsline{toc}{subsection}{圍田租畝}



\begin{paracol}{2}



\ClassicalHeader



\ClassicalQuestion{%

有興復圍田已成，共計三千二十一頃五十一畝一十五歩，分三等。其上等毎畆起租六斗，中等四斗五升，下等四斗。中田多上田弱半，不及下田太半。欲知三色田畝及各租幾何？}



\switchcolumn



\ModernHeader



\ModernQuestion{%

有興修恢復的圍田已經完成，共計3021頃51畝15步，分為三等。上等田每畝徵租6斗，中等4斗5升，下等4斗。中等田比上等多三分之一（弱半），比下田少三分之二（太半）。想知道三種田地的畝數以及各等田租分別是多少？

}



\end{paracol}



\begin{paracol}{2}



\ClassicalNote{%

題言中田多上田弱半，不及下田太半。蓋以弱半為三分之一，太半為三分之二也。多上田弱半者，即比上田多三分之一也。不及下田太半者，即比下田少三分之二也。是上田加三分之一為中田，即下田三分之一也。轉言之，則中田為下田三分之一，上田為中田四分之三也。}



\switchcolumn



\ModernNote{%

題目說中等田比上田多弱半，不及下田太半。弱半就是三分之一，太半就是三分之二。「多上田弱半」就是比上田多三分之一。「不及下田太半」就是比下田少三分之二。所以上田加上三分之一就是中田，而中田是下田的三分之一。換句話說，中田是下田的三分之一，上田是中田的四分之三。

}



\end{paracol}



\begin{paracol}{2}



\ClassicalAnswer{%

上田四百七十七頃八畝一十五歩，米二萬八千六百二十四石八斗三升七合五勺；

中田六百三十六頃一十畝三角，米二萬八千六百二十四石八斗三升七合五勺；

下田一千九百八頃三十二畝一角，米七萬六千三百三十二石九斗。}



\switchcolumn



\ModernAnswer{%

上田477頃8畝15步，租米28,624石8斗3升7合5勺；

中田636頃10畝30步，租米28,624石8斗3升7合5勺；

下田1908頃32畝10步，租米76,332石9斗。}



\end{paracol}



\begin{paracol}{2}



\ClassicalMethod{%

以衰分求之。列母子求田率，副併為法，以共田為寔，寔如法而一，得一分之率。以徧乗未併者，得三等田。各以起租乗之，各得米。}



\switchcolumn



\ModernMethod{%

使用衰分法求解。排列母子數求田地的比率，合併作為法，以總田數為被除數，被除數除以法，得到一分的比率。遍乘未合併的各數，得到三等田的畝數。再各以起租數相乘，得到各等田租的米數。

}



\end{paracol}



\begin{paracol}{2}



\ClassicalWorking{%

置弱半母四為中率，子三為上率。以太半子二减母三，餘一，以乗中率四，只得四為中泛。又以餘一乗上率三，只得三為上泛。次以太半母三乗中泛四，得一十二為下泛。副併三泛，得一十九為法。置田三千二十一頃五十一畆，以畝法二百四十通之，得七千二百五十一萬六千二百四十歩，内子一十五歩，得七千二百五十一萬六千二百五十五歩為寔。以法一十九除之，得三百八十一萬六千六百四十五歩為一分之數。以上泛三因之，得一千一百四十四萬九千九百三十五歩為上積。又以中泛四因之一分數，得一千五百二十六萬六千五百八十歩為中積。又以下泛一十二乗一分數，得四千五百七十九萬九千七百四十歩為下積。其三積各以畝法二百四十約之為畆。其上田得四百七十七頃八畝一十五歩，其中田得六百三十六頃一十畝三角，其下積得一千九百八頃三十二畝一角。各以起租，三積為三寔。其上積一千一百四十四萬九千九百三十五歩乗上積六斗，得六百八十六萬九千九百六十一石為寔，以畝法二百四十除之，得二萬八千六百二十四石八斗三升七合五勺為上田租。其中積一千五百二十六萬六千五百八十歩乗中田租四斗五升，得六百八十六萬九千九百六十一石為寔，以畝法二百四十除之，得二萬八千六百二十四石八斗三升七合五勺。其下積四千五百七十九萬九千七百四十歩乗下租四斗，得一千八百三十一萬九千八百九十六石為寔，以畝法二百四十除之，得七萬六千三百三十二石九斗為下田米。}



\switchcolumn



\ModernWorking{%

將「弱半」的分母4作為中等田的比率，分子3作為上等田的比率。用「太半」的分子2減分母3，餘1，乘中率4，得4為中泛數。再用餘1乘上率3，得3為上泛數。再用太半分母3乘中泛4，得12為下泛數。合併三個泛數，得19為法。將總田3021頃51畝以畝法240步換算，得72,516,240步，加上零數15步，得72,516,255步為被除數。以法19除之，得3,816,645步為一分的數。以上泛3乘之，得11,449,935步為上田積步。以中泛4乘一分數，得15,266,580步為中田積步。以下泛12乘一分數，得45,799,740步為下田積步。三積各以畝法240約為畝數。上田得477頃8畝15步，中田得636頃10畝30步，下田得1908頃32畝10步。



各以起租數計算：上積11,449,935步乘上租6斗，得686,996.1石為被除數，以畝法240除之，得28,624石8斗3升7合5勺為上田租。中積15,266,580步乘中租4斗5升，得686,996.1石為被除數，以240除之，得28,624石8斗3升7合5勺為中田租。下積45,799,740步乘下租4斗，得1,831,989.6石為被除數，以240除之，得76,332石9斗為下田租米。

}



\end{paracol}



\begin{paracol}{2}



\ClassicalNote{%

草内求各衰數，意謂以三分為上田衰數，加弱半三分之一得四分為中田衰數，三因中田衰數得一十二分為下田衰數，併之得一十九分為總衰數。其法本属顯明，但語中參入母子副泛等名目，其意反晦矣。}



\switchcolumn



\ModernNote{%

算草中求各衰數的意思，是用3分作為上田的衰數，加上弱半（三分之一）得4分作為中田衰數，3倍中田衰數得12分作為下田衰數，合併得19分為總衰數。這個算法本來很明白，只是語言中加入了母子、副泛等名詞，反而使意思晦澀了。

}



\end{paracol}







%% =============================================================================

%% FASCICLE 5B — 賦役 (Part 2)

%% =============================================================================

\newpage

\section*{\color{titlecolor}卷五下\quad 賦役（賦稅和徭役）}

\addcontentsline{toc}{section}{卷五下\quad 賦役}



\subsection*{\color{sectioncolor}戶稅移割（戶稅轉移分割）}

\addcontentsline{toc}{subsection}{戶稅移割}



\begin{paracol}{2}



\ClassicalHeader



\ClassicalQuestion{%

某縣據甲稱：本户田地原納苗三十五石七斗，和買本色一十一匹二丈二尺九寸四分八釐七毫五絲，折帛二十七匹二寸一分三釐七毫五絲，紬絹折帛八匹三丈九寸七寸三分九釐，紬絹本色二十匹二丈八尺九寸三分九釐。已将田四百七畆出與乙，五百十六畆出與丙。訖乞移割本户所出田上税賦，歸併乙丙兩户送納。㑹到乙原有田三百七十五畆，丙原有田四百六十三畆，並係本鄉本等。毎畆苗三升五合，税一尺一寸五分，物力一貫二百；本等田紬一尺三寸四分，物力九百；物力及三十二貫數和買一匹，内三分本色七分折帛；夏税七分本色三分折帛；紬五分本色五分折帛，併絹紬。欲知甲田地及甲乙丙分割合納苗米、和買、夏税、折帛、本色、畸零各幾何？}



\switchcolumn



\ModernHeader



\ModernQuestion{%

某縣據甲申報：本户田地原來繳納苗米35石7斗，和買本色（現物）11匹2丈2尺9寸4分8釐7毫5絲，折帛（折錢絹）27匹2寸1分3釐7毫5絲，紬絹折帛8匹3丈9寸7分3釐9毫，紬絹本色20匹2丈8尺9寸3分9釐。現在已將田407畝出讓給乙，516畝出讓給丙。請求將本户所出田地上的税賦轉移，歸併到乙丙兩户繳納。經查乙原有田375畝，丙原有田463畝，都是本鄉本等田地。每畝苗米3升5合，税絹1尺1寸5分，物力（財產估值）1貫200文；本等田每畝紬1尺3寸4分，物力900文；物力達到32貫可和買1匹，其中三分本色七分折帛；夏税七分本色三分折帛；紬五分本色五分折帛，與絹紬合併計算。想知道甲的田地以及甲乙丙分割後各應繳納的苗米、和買、夏税、折帛、本色、零頭各為多少？

}



\end{paracol}



\begin{paracol}{2}



\ClassicalNote{%

此題科税分田地二項。田有科米、税絹、物力三色，地只有税紬、物力二色，絹與紬折色不同，物力和買合田地總計之。又甲户有田者有地，乙丙二户有田無地。此等皆不可不明言者，題中殊未分晰。}



\switchcolumn



\ModernNote{%

這道題的科税分為田地兩項。田有科米、税絹、物力三項，地只有税紬、物力二項，絹與紬的折色不同，物力和買要將田地合計。另外甲户既有田又有地，乙丙二户有田無地。這些都應該明確說明，但題中完全沒有分清楚。

}



\end{paracol}



\begin{paracol}{2}



\ClassicalAnswer{%

甲原有田一千二十畆，地一十一畆二角四十八步。



甲今有田九十七畆，苗米三石三斗九升五合，夏税折帛三丈三尺四寸六分五釐，本色一匹畸零三丈八尺八分五釐，物力一百一十六貫四百文。地一十一畆二角四十八步，税紬折帛七尺八寸三分九釐，税紬畸零七尺八寸三分九釐，物力一十貫五百三十文，田地共物力一百二十六貫九百三十文。和買折帛二匹三丈一尺六分三釐七毫五絲，本色一疋畸零七尺五寸九分八釐七毫五絲。



乙今有田七百八十二畆，苗米二十七石三斗七升，夏税折帛六匹二丈九尺七寸九分，本色一十五匹畸零二丈九尺五寸一分，物力九百三十八貫四百文，和買折帛二十匹二丈一尺一寸，本色八匹畸零三丈一尺九寸。



丙今有田九百七十九畆，苗米三十四石二斗六升五合，夏税折帛八匹一丈七尺七寸五分五釐，本色一十九匹畸零二丈八尺九分五釐，物力一千一百七十四貫八百文，和買折帛二十五匹二丈七尺九寸五分，本色一十一匹畸零五寸五分。}



\switchcolumn



\ModernAnswer{%

甲原有田1020畝，地11畝零48步。



甲現有田97畝，苗米3石3斗9升5合，夏税折帛3丈3尺4寸6分5釐，本色1匹零3丈8尺8分5釐，物力116貫400文。地11畝零48步，税紬折帛7尺8寸3分9釐，税紬零數7尺8寸3分9釐，物力10貫530文，田地共物力126貫930文。和買折帛2匹3丈1尺6分3釐7毫5絲，本色1匹零7尺5寸9分8釐7毫5絲。



乙現有田782畝，苗米27石3斗7升，夏税折帛6匹2丈9尺7寸9分，本色15匹零2丈9尺5寸1分，物力938貫400文，和買折帛20匹2丈1尺1寸，本色8匹零3丈1尺9寸。



丙現有田979畝，苗米34石2斗6升5合，夏税折帛8匹1丈7尺7寸5分5釐，本色19匹零2丈8尺9分5釐，物力1,174貫800文，和買折帛25匹2丈7尺9寸5分，本色11匹零5寸5分。

}



\end{paracol}



\begin{paracol}{2}



\ClassicalMethod{%

以粟米及衰分求之。置甲原納米為實，以毎畆苗為法除之，得甲原有田。以乗毎畆税，得田絹。次以甲紬絹本色折帛併之，内減田絹，餘為地紬。以毎紬約之，得甲原有地。次以甲出田併乙丙原有田，各得三户今有田地。各以等則乗之，各得物力苗税。次以毎畆物力率約各户共物力，得和買。副之，三分因之退位為和本，七分因之退位為和折；夏税反其分而因之退之；紬以半之，併入税，各得。}



\switchcolumn



\ModernMethod{%

使用粟米法和衰分法求解。將甲原納米數作為被除數，以每畝苗數為除數去除，得甲原有田數。再乘每畝税數，得田絹數。再將甲的紬絹本色折帛合併，內減田絹，餘數為地紬。以每畝紬數約之，得甲原有地數。再將甲出讓的田合併乙丙原有田，各得三户現有田地數。各以等則相乘，各得物力苗税數。再以每畝物力率約各户共物力，得和買數。分開處理：三分為和本（本色），七分為和折（折帛）；夏税則反過來（七分本色三分折帛）；紬一半本色一半折帛，併入税中，各得結果。

}



\end{paracol}



\newpage

\subsection*{\color{sectioncolor}均科綿税（均摊綿稅）}

\addcontentsline{toc}{subsection}{均科綿税}



\begin{paracol}{2}



\ClassicalHeader



\ClassicalQuestion{%

縣科綿有五等户，共一萬一千三十三户，共科綿八萬八千三百三十七兩六錢。上等一十二户，副等八十七户，中等四百六十四户，次等二千三十五户，下等八千四百三十五户。欲令上三等折半差，下二等此中等六四折差。科率求之，問各户納及各等幾何？}



\switchcolumn



\ModernHeader



\ModernQuestion{%

縣裡徵收綿有五等户，共11,033户，共徵綿88,337兩6錢。上等12户，副等87户，中等464户，次等2,035户，下等8,435户。要求上三等折半遞差（即每等為上一半），下二等比中等六四折差（即次等為中等的十分之四，下等又為次等的十分之四）。按科率求之，問每户應納多少以及各等共納多少？

}



\end{paracol}



\begin{paracol}{2}



\ClassicalNote{%

題言下二等比中等六四折差，是次等為中等六分之四，下等又為次等六分之四也。乃術草中皆以十分之四收之，是四分折差而非四六折差矣。集中題與術不相應者多類此，豈成書之後未能重加校正歟？}



\switchcolumn



\ModernNote{%

題目說下二等比中等六四折差，意思是次等為中等的十分之四，下等又為次等的十分之四。但是術草中都按十分之四（即0.4倍）來計算，這是四分折差而不是四六折差。集中題目與術不相應的地方多類似於此，難道是成書之後未能重新校正嗎？

}



\end{paracol}



\begin{paracol}{2}



\ClassicalAnswer{%

上等一户一百二十四兩，一十二户計一千四百八十八兩；

副等一户六十二兩，八十七户計五千三百九十四兩；

中等一户三十一兩，四百六十四户計一萬四千三百八十四兩；

次等一户一十二兩四錢，二千三十五户計二萬五千二百三十四兩；

下等一户四兩九錢六分，八千四百三十五户計四萬一千八百三十七兩六錢。}



\switchcolumn



\ModernAnswer{%

上等每户124兩，12户共計1,488兩；

副等每户62兩，87户共計5,394兩；

中等每户31兩，464户共計14,384兩；

次等每户12兩4錢，2,035户共計25,234兩；

下等每户4兩9錢6分，8,435户共計41,837兩6錢。

}



\end{paracol}



\begin{paracol}{2}



\ClassicalMethod{%

列五等户數，先以四折下等數，加次等户，又以四折之，加中等户數，却以半折之，加二等户，又以半折之，加上等户數，不折便為法。除科綿，得上等一户之綿。復半之為副等，又半之為中等，又四折之為次等，又四折之為下等。各一户綿，却各以户數乗之，各得五等共出綿。}



\switchcolumn



\ModernMethod{%

排列五等户數，先用四折（乘以0.4）下等户數，加次等户，再用四折，加中等户數，再用半折，加副等户，再用半折，加上等户數，不再折就作為除數（法）。以總科綿數除以法，得上等一户應納綿數。再半折為副等，再半折為中等，再四折為次等，再四折為下等。得到各一户應納綿數，再各以户數乘之，各得五等共出綿數。

}



\end{paracol}



\begin{paracol}{2}



\ClassicalWorking{%

先置下等户八千四百三十五，以四折之，得三千三百七十四。乃以得數折入次户二千三十五内，得五千四百九户訖。又以四折次數五千四百九户，得二千一百六十三户六分。乃以得次數併入中户四百六十四内，共得二千六百二十七户六分。次以五分折中數二千六百二十七户六分，得數。次以得數一千三百一十三户八分併副户八十七，得一千四百户八分。又以五折副數一千四百户八分，得七百户四分，既得數。乃以副數七百户四分併上户一十二户，得七百一十二户四分，不折便為法。



累折至等，共得七百一十二户四分以為總法。乃置科綿八萬八千三百三十七兩六錢為實，法七百一十二户四分而一，得一百二十四兩為上等一户所出綿。以五折之，得六十二兩為副等一户所出綿。又五折之，得三十一兩為中等一户所出綿。乃以四折之，得一十二兩四錢為次等一户所出綿。又以四折之，得四兩九錢六分為下等一户所出綿。乃以各等户數各乗一户所出，即各得每等共數之綿。}



\switchcolumn



\ModernWorking{%

先將下等户8,435，四折得3,374。將此數加次等户2,035，得5,409户。再四折5,409户，得2,163.6户。再併入中等户464，共得2,627.6户。再以半折2,627.6户，得1,313.8户。併入副等户87，得1,400.8户。再半折1,400.8户，得700.4户。再併入上等户12户，得712.4户，不再折就作為總法。



累計折算到此，共得712.4户為總法。將科綿88,337兩6錢為被除數，以法712.4去除，得124兩為上等一户應出綿數。半折得62兩為副等一户所出綿。再半折得31兩為中等一户所出綿。再四折得12兩4錢為次等一户所出綿。再四折得4兩9錢6分為下等一户所出綿。再以各等户數各乘一户所出數，即各得每等共出綿數。

}



\end{paracol}



\subsection*{\color{sectioncolor}均定勸分（均定勸導分攤）}

\addcontentsline{toc}{subsection}{均定勸分}



\begin{paracol}{2}



\ClassicalHeader



\ClassicalQuestion{%

欲勸糶賑濟，據甲民户物力畆步排定，共計一百六十二户，作九等：上等三户，第二等五户，第三等七户，第四等八户，第五等十三户，第六等二十一户，第七等二十六户，第八等三十四户，第九等四十五户。今先觀諭，第一等上户願糶五千石，第九等户願糶二百石。欲知各等抛差石數併數認米數各幾何？}



\switchcolumn



\ModernHeader



\ModernQuestion{%

要勸導賣糧賑濟災民，據甲地民户按財產田畝排定，共計162户，分為九等：上等3户，第二等5户，第三等7户，第四等8户，第五等13户，第六等21户，第七等26户，第八等34户，第九等45户。現先公告曉諭，第一等上户願意出糧5,000石，第九等户願意出糧200石。想知道各等之間的遞差石數以及各等認購米數各為多少？

}



\end{paracol}



\begin{paracol}{2}



\ClassicalAnswer{%

認該米二十三萬七千六百石。



上等一户米五千石，三户計一萬五千石；

二等一户米四千四百石，五户計二萬二千石；

三等一户米三千八百石，七户計二萬六千六百石；

四等一户米三千二百石，八户計二萬五千六百石；

五等一户米二千六百石，一十三户計三萬三千八百石；

六等一户米二千石，二十一户計四萬二千石；

七等一户米一千四百石，二十六户計三萬六千四百石；

八等一户米八百石，三十四户計二萬七千二百石；

九等一户米二百石，四十五户計九千石。}



\switchcolumn



\ModernAnswer{%

共認購米237,600石。



上等每户5,000石，3户共計15,000石；

第二等每户4,400石，5户共計22,000石；

第三等每户3,800石，7户共計26,600石；

第四等每户3,200石，8户共計25,600石；

第五等每户2,600石，13户共計33,800石；

第六等每户2,000石，21户共計42,000石；

第七等每户1,400石，26户共計36,400石；

第八等每户800石，34户共計27,200石；

第九等每户200石，45户共計9,000石。

}



\end{paracol}



\begin{paracol}{2}



\ClassicalMethod{%

以衰分求之。置上下户米，減餘為實。列等數減一，餘為法。除之，得抛差石數。以差累減上等米，各得諸等米。以各等户乗之，併之為總數。}



\switchcolumn



\ModernMethod{%

使用衰分法求解。將上户米數減下户米數，餘數作為被除數。排列等數減一，餘數作為除數。相除得到遞差石數。以此差數逐級從上等米數中遞減，得到各等應出米數。再乘以各等户數，合併為總數。

}



\end{paracol}



\begin{paracol}{2}



\ClassicalWorking{%

置上等户米五千石，減下等户二百石，餘四千八百石為實。以九等減一，餘八為法。除實，得六百石為每等抛差。用減上等，餘四千四百石為二等米。又減六百，得三千八百石為三等米。又減六百，得三千二百石為四等米。又減六百，得二千六百石為五等米。又減六百，得二千石為六等米。又減六百，得一千四百石為七等米。又減六百，得八百石為八等米。又減六百，得二百石為九等米。又各乗户數，併之，得總認米石數二十三萬七千六百石。}



\switchcolumn



\ModernWorking{%

將上等户米5,000石，減下等户200石，餘4,800石為被除數。以9等減1，餘8為除數。4,800除以8得600石為每等遞差數。用上等5,000石減600石，餘4,400石為第二等米數。再減600石，得3,800石為第三等。再減600石，得3,200石為第四等。再減600石，得2,600石為第五等。再減600石，得2,000石為第六等。再減600石，得1,400石為第七等。再減600石，得800石為第八等。再減600石，得200石為第九等。再各乘户數，合併得總認米237,600石。

}



\end{paracol}



\subsection*{\color{sectioncolor}均定合解（均定合併解送）}

\addcontentsline{toc}{subsection}{均定合解}



\begin{paracol}{2}



\ClassicalHeader



\ClassicalQuestion{%

州郡合解諸司窠名錢：户部九十六萬五千四百二十一貫文，總所六十四萬三千六百一十四貫文，運司一萬六千九十貫三百五十文。今諸窠名先催到九千二百五十三貫六百二十文，欲照原額分數均定樁收解，合各幾何？}



\switchcolumn



\ModernHeader



\ModernQuestion{%

州郡應合解送各部門的專款錢：户部965,421貫文，總所643,614貫文，運司16,090貫350文。現在各部門先催收到9,253貫620文，想按照原額比例均定留存解送，各部門應各得多少？

}



\end{paracol}



\begin{paracol}{2}



\ClassicalAnswer{%

户部五千四百九十七貫二百文；總所三千六百六十四貫八百文；運司九十一貫六百二十文。}



\switchcolumn



\ModernAnswer{%

户部5,497貫200文；總所3,664貫800文；運司91貫620文。}



\end{paracol}



\begin{paracol}{2}



\ClassicalMethod{%

以衰分求之。置諸原率，可約約之。副併為法。以催到錢乗未併者，各為實。實如法而一。}



\switchcolumn



\ModernMethod{%

使用衰分法求解。將各原額比率排列，可以約分的先約分。合併作為除數（法）。以催收到的錢數乘未合併的各項，各為被除數（實）。被除數除以法得到各數。

}



\end{paracol}



\begin{paracol}{2}



\ClassicalWorking{%

列户部九十六萬五千四百二十一貫，總所六十四萬三千六百一十四貫，運司一萬六千九十貫三百五十文，各為原率。今原率可約，求等得一萬六千九十貫三百五十為等數，俱約之：户部得六十，總所得四十，運司得一，各為率。副併得一百一為法。以催到錢九千二百五十三貫六百二十文乗未併數：六十、四十及一，畢得五十五萬五千二百一十七貫二百文為户部之實，三十七萬一百四十四貫八百文為總所之實，九千二百五十三貫六百二十文為運司之實。三實皆如法一百一而一：其户部得五千四百九十七貫二百文，總所得三千六百六十四貫八百文，運司得九十一貫六百二十，各為候解錢。}



\switchcolumn



\ModernWorking{%

排列户部965,421貫，總所643,614貫，運司16,090貫350文，各為原額比率。現在原率可以約分，求最大公約數得16,090貫350文為等數，各約之：户部得60，總所得40，運司得1，各為率。合併得101為法。以催到錢9,253貫620文乘未合併的數60、40及1：得555,217貫200文為户部之被除數，371,448貫800文為總所之被除數，9,253貫620文為運司之被除數。三個被除數都以法101去除：户部得5,497貫200文，總所得3,664貫800文，運司得91貫620文，各為候解錢數。

}



\end{paracol}



\newpage

\subsection*{\color{sectioncolor}寬減屯租（寬減屯田租賦）}

\addcontentsline{toc}{subsection}{寬減屯租}



\begin{paracol}{2}



\ClassicalHeader



\ClassicalQuestion{%

屯租欲議寬減，仍聴以夏麥折納分數。官牛種者與減二分，私牛種者與減四分。每嵗租榖以三分之一許夏折二麥，内四分大六分小折色。每大麥三石折小麥二石，小麥二石折榖三石五斗。屯租舊額官種一石納租五石，私種一石納租三石。今某州屯田去年計官私種共九千七百八十二石，共合收租榖三萬九千五百八十六石。欲知官私種各數、原額、今減、合催、成年夏麥秋榖幾何？}



\switchcolumn



\ModernHeader



\ModernQuestion{%

屯田租賦想要商議寬減，仍允許以夏麥折算繳納。官府提供牛種的減免二分（20\%），私人牛種的減免四分（40\%）。每年租穀以三分之一允許夏季折為大小麥繳納，其中四分大麥六分小麥折色。每大麥3石折小麥2石，小麥2石折穀3石5斗。屯租舊額：官種1石納租5石，私種1石納租3石。現在某州屯田去年官私種共計9,782石，共應收租穀39,586石。想知道官種私種各多少、原額租數、現減數、應催數、成年（豐年）夏麥秋穀各多少？

}



\end{paracol}



\begin{paracol}{2}



\ClassicalNote{%

此題有官私共種數、租數，有種一石租數，求各種數租數，古謂之差分，今謂之和較比例。折色亦差分也，今謂之和數比例。大小麥與榖相易，古謂之異乗同乗，今謂之和率比列。}



\switchcolumn



\ModernNote{%

這道題有官私共種數、租數，有每石種子的租數，求各種數租數，古代稱為差分，現代稱為和差比例。折色也是差分，現代稱為和數比例。大小麥與穀互相折算，古代稱為異乘同乘，現代稱為連鎖比例。

}



\end{paracol}



\begin{paracol}{2}



\ClassicalAnswer{%

官種五千一百二十石，私出種四千六百六十二石。



原租額共三萬九千五百八十六石：官種二萬五千六百石，私種一萬三千九百八十六石。



今減一萬七百一十四石四斗：官種五千一百二十石，私種五千五百九十四石四斗。



合催二萬八千八百七十一石六斗。



官種租二萬四百八十石：一分折麥計榖六千八百二十六石六斗六升零三分升之二；四分折大麥二千三百四十石五斗七升七分升之一（係折榖二千七百三十石六斗六升三分升之二）；六分折小麥二千三百四十石五斗七升七分升之一（係折榖四千九十六石二分）；正色榖一萬三千六百五十三石三斗三升三分升之一。



私種租八千三百九十一石六斗：一分折麥計榖二千七百九十七石二斗；四分折大麥九百五十九石四升（係折榖一千一百一十八石八斗八升）；六分折小麥九百五十九石四升（係折榖一千六百七十八石三斗二升）；二分正色榖五千五百九十四石四斗。



成年共收：夏折榖九千六百二十三石八斗六升三分升之二；大麥三千二百九十九石六斗一升七分升之一；小麥三千二百九十九石六斗一升七分升之一；秋正榖一萬九千二百四十七石七斗三升三分升之一。}



\switchcolumn



\ModernAnswer{%

官種5,120石，私種4,662石。



原租額共39,586石：官種25,600石，私種13,986石。



現減10,714石4斗：官種減5,120石，私種減5,594石4斗。



應催（實收）28,871石6斗。



官種租20,480石：三分之一折麥計穀6,826石6斗6升又2/3升；其中四分折大麥2,340石5斗7升又1/7升（折算穀2,730石6斗6升又2/3升）；六分折小麥2,340石5斗7升又1/7升（折算穀4,096石2分）；正色穀13,653石3斗3升又1/3升。



私種租8,391石6斗：三分之一折麥計穀2,797石2斗；四分折大麥959石4升（折算穀1,118石8斗8升）；六分折小麥959石4升（折算穀1,678石3斗2升）；二分正色穀5,594石4斗。



成年共收：夏折穀9,623石8斗6升又2/3升；大麥3,299石6斗1升又1/7升；小麥3,299石6斗1升又1/7升；秋正穀19,247石7斗3升又1/3升。

}



\end{paracol}



\subsection*{\color{sectioncolor}戶稅均寬（戶稅均平寬減）}

\addcontentsline{toc}{subsection}{戶稅均寬}



\begin{paracol}{2}



\ClassicalHeader



\ClassicalQuestion{%

州郡寬恤，近将某縣下三等税户秋科餘欠錢米，已與蠲放：共錢一千三百五十五貫七百六文，米五千二百七十二石一斗九升。某本縣下三等物力計三萬七千六百五十八貫五百文。今來官員陳述：本縣多有樂輸無欠之户，今䝉蠲放税尾，似反寬潤頑輸之户，於理未均。遂議将樂輸三等户於明年兩税與照昨來體例減免。契勘得三等無欠户物力二十二萬八千一十五貫三百二十一文。欲知每百文合減免錢米及共減各幾何？}



\switchcolumn



\ModernHeader



\ModernQuestion{%

州郡寬大體恤，最近將某縣下三等税户秋季徵收的餘欠錢米，已經免除：共錢1,355貫706文，米5,272石1斗9升。本縣下三等物力共計37,658貫500文。現在有官員陳述：本縣有很多樂意繳納沒有欠税的户，現在蒙免欠税尾數，反而像是優惠了頑固欠税的户，於理不均。於是商議將樂意繳納的三等户在明年兩税中按照此次體例減免。查得三等無欠户物力共228,015貫321文。想知道每100文物力應減免多少錢米以及共減免多少？

}



\end{paracol}



\begin{paracol}{2}



\ClassicalAnswer{%

毎物力一百文，放錢三文六分，放米一升四合。明年兩税放免：錢七千九百四十九貫三百五十一文五分五釐六毫；米三萬九百一十四石一斗四升四合九勺四抄。}



\switchcolumn



\ModernAnswer{%

每物力100文，免除錢3文6分，免除米1升4合。明年兩税免除：錢7,949貫351文5分5釐6毫；米30,914石1斗4升4合9勺4抄。

}



\end{paracol}



\begin{paracol}{2}



\ClassicalMethod{%

以粟米衰分求之。置原物力為法，除原放錢米，得毎百文物力所放錢米率。以各率乗今來物力，各得錢米，為明年兩税合放數。}



\switchcolumn



\ModernMethod{%

使用粟米衰分法求解。將原物力數作為除數，去除原免除的錢米數，得到每百文物力所免除的錢米比率。以各比率乘現在的物力，各得錢米數，即為明年兩税應免除數。

}



\end{paracol}



\begin{paracol}{2}



\ClassicalWorking{%

以原物力三萬七千六百五十八貫五百文為法。先除原放錢一千三百五十五貫七百六文，定法百文上得一文為商，除得三文六分為物力百文所放錢率。次除原放米五千二百七十二石一斗九升，得一升四合為物力百文所放米率。以今三等户物力二十二萬八千一十五貫三百二十一文徧乗所放錢米二率，得錢七千九百四十九貫三百五十一文五分五釐六毫，得米三萬九百一十四石一斗四升四合九勺四抄，為明年兩税合放數。}



\switchcolumn



\ModernWorking{%

以原物力37,658貫500文為除數。先去除原免除錢1,355貫706文，在百文位上得商，除得3文6分為每百文物力所免除的錢率。再去除原免除米5,272石1斗9升，得1升4合為每百文物力所免除的米率。以現在三等户物力228,015貫321文遍乘所免除的錢米二率，得錢7,949貫351文5分5釐6毫，得米30,914石1斗4升4合9勺4抄，為明年兩税應免除數。

}



\end{paracol}



\newpage

\subsection*{\color{sectioncolor}米榖粒分（米穀按粒分計）}

\addcontentsline{toc}{subsection}{米榖粒分}



\begin{paracol}{2}



\ClassicalHeader



\ClassicalQuestion{%

開倉受約，有甲户米一千五百三十四石。到廊驗得米内夾榖，乃於様内取米一捻，數計二百五十四粒，内有榖二十八顆。凡粒米率每勺三百。今欲知米内雜榖多少，以折米數科貴及粒各幾何？}



\switchcolumn



\ModernHeader



\ModernQuestion{%

開倉收納，有甲户交米1,534石。到倉檢驗發現米中夾雜穀粒，於是從樣品中取米一捻，數得254粒，其中有穀28顆。每勺米約300粒。現在想知道米中夾雜穀粒有多少，折算成米數來科罰以及總粒數各為多少？

}



\end{paracol}



\begin{paracol}{2}



\ClassicalAnswer{%

米一千三百六十四石八斗九升七合六勺一百二十七分勺之四十八；

榖一百六十九石一斗二合三勺一百二十七分勺之七十九。

合折米八十四石五斗五升一合一勺一百二十七分勺之一百三。

原米折米共計四十三億四千八百三十四萬六千四百五十六粒。}



\switchcolumn



\ModernAnswer{%

淨米1,364石8斗9升7合6勺又48/127勺；

穀169石1斗2合3勺又79/127勺。

合計應折米84石5斗5升1合1勺又103/127勺。

原米折算共計4,348,346,456粒。

}



\end{paracol}



\begin{paracol}{2}



\ClassicalMethod{%

以粟米求之，衰分入之。置様米粒數為法，以常榖顆數減之，餘與榖為列衰。可約約之。以共米乗列衰為各實，實如法而一，各得米數榖數。置榖數以粟率折之，為榖所折米。次以勺率遍乗米數折米，得粒數。}



\switchcolumn



\ModernMethod{%

使用粟米法求解，以衰分法配合。將樣品米粒總數作為除數，減去穀粒數，餘數與穀粒數作為列衰。可以約分的約分。以總米數乘列衰各為被除數，除以法，各得米數和穀數。將穀數以粟率（50\%）折算，為穀應折成的米數。再以勺率（300粒/勺）遍乘米數和折米數，得到粒數。

}



\end{paracol}



\begin{paracol}{2}



\ClassicalWorking{%

置一捻様粒數二百五十四為法，以帶榖二十八顆為榖衰，以減法，餘二百二十六為米衰。此二衰與法皆可約，求等得二，俱以約之：法得一百二十七，米衰得一百一十三，榖衰得一十四。以米一千五百三十四石遍乗二衰，得一十七萬三千三百四十二石為米實，得二萬一千四百七十六石為榖實。皆如法一百二十七而一：米得一千三百六十四石八斗九升七合六勺一百二十七分勺之四十八，榖得一百六十九石一斗二合三勺一百二十七分勺之七十九。以粟率五十折之，得八十四石五斗五升一合一勺一百二十七分勺之一百三為榖折納米數。併二米，得一千四百四十九石四斗四升八合八勺一百二十七分勺之二十四。先通分納子一十八萬四千八十石，以勺率三百粒乗之，得五千五百二十二億四千萬粒為實，以母一百二十七除之，得四十三億四千八百三十四萬六千四百五十六粒，不盡八十八棄之。}



\switchcolumn



\ModernWorking{%

將一捻樣品254粒作為除數，以帶穀28顆為穀衰，減法得226為米衰。這兩個衰數與法都可以約分，求最大公約數得2，各約之：法得127，米衰得113，穀衰得14。以米1,534石遍乘二衰，得173,342石為米實，得21,476石為穀實。都以法127去除：米得1,364石8斗9升7合6勺又48/127勺，穀得169石1斗2合3勺又79/127勺。以粟率50\%折穀為米，得84石5斗5升1合1勺又103/127勺為穀應折納米數。合併淨米和折米，得1,449石4斗4升8合8勺又24/127勺。先通分內子為184,080石，以勺率300粒乘之，得552,240,000,000粒為被除數，以分母127去除，得4,348,346,456粒，餘數88粒棄去。

}



\end{paracol}







%% =============================================================================

%% FASCICLE 6A — 錢穀 (Part 1)

%% =============================================================================

\newpage

\section*{\color{titlecolor}卷六上\quad 錢穀（貨幣和糧食）}

\addcontentsline{toc}{section}{卷六上\quad 錢穀}



\subsection*{\color{sectioncolor}課糴（徵收糴米）}

\addcontentsline{toc}{subsection}{課糴}



\begin{paracol}{2}



\ClassicalHeader



\ClassicalQuestion{%

差人五路和糴。據浙西平江府石價三十五貫文，一百三十五合，至鎮江水脚錢每石九百文；安吉州石價二十九貫五百文，一百一十合，至鎮江水脚錢每石一貫二百文；江西隆興府石價二十八貫一百文，一百一十五合，至建康水脚錢每石一貫七百文；吉州石價二十五貫八百五十文，一百二十合，至建康水脚錢每石二貫九百文；湖南潭州石價二十七貫三百文，一百一十八合，至鄂州水脚錢每石一貫七百文。其錢並十七界官㑹，其米並用文思院斛交量紐數。欲皆以官解計石錢相比貴賤幾何？}



\switchcolumn



\ModernHeader



\ModernQuestion{%

差遣人員到五路收購糧米。據報：浙西平江府每石價35貫文，每石135合（合=0.1升），到鎮江每石水脚運費900文；安吉州每石價29貫500文，每石110合，到鎮江每石水脚1貫200文；江西隆興府每石價28貫100文，每石115合，到建康每石水脚1貫700文；吉州每石價25貫850文，每石120合，到建康每石水脚2貫900文；湖南潭州每石價27貫300文，每石118合，到鄂州每石水脚1貫700文。這些錢都是十七界官會（紙幣），這些米都用文思院斛（標準量器）來交接換算。想都以官斛計算每石的實際價錢，比較貴賤如何？

}



\end{paracol}



\begin{paracol}{2}



\ClassicalNote{%

按草中係二貫一百文，此訛。文思院解每斗八十三合。}



\switchcolumn



\ModernNote{%

按算草中應為2貫100文，此處有誤。文思院斛每斗83合。

}



\end{paracol}



\begin{paracol}{2}



\ClassicalAnswer{%

文思院斛石錢：安吉州二十三貫一百六十四文一十一分文之六；平江府二十二貫七十一文二十七分文之二十三；隆興府二十一貫五百七文二十三分文之一十九；潭州二十貫六百七十九文五十九分文之四十九；吉州一十九貫八百八十五文十二分文之五。}



\switchcolumn



\ModernAnswer{%

換算為文思院斛每石錢：安吉州23貫164文又6/11文；平江府22貫71文又23/27文；隆興府21貫507文又19/23文；潭州20貫679文又49/59文；吉州19貫885文又5/12文。

}



\end{paracol}



\begin{paracol}{2}



\ClassicalNote{%

按三十九訛四十九。}



\switchcolumn



\ModernNote{%

按：「三十九」應為「四十九」，原文有誤。

}



\end{paracol}



\begin{paracol}{2}



\ClassicalMethod{%

以粟米互換求之。置石價併水脚，乗石數，又乗官斗合數為實。各如本州合數而一，各得官解石錢。以課貴賤。}



\switchcolumn



\ModernMethod{%

使用粟米互換法求解。將每石價格加上水脚運費，乘以石數，再乘以官斗合數作為被除數。各按本州的合數去除，得到官斛每石錢數。以此比較貴賤。

}



\end{paracol}



\begin{paracol}{2}



\ClassicalWorking{%

置安吉州石價二十九貫五百文，平江石價三十五貫文，隆興石價二十八貫一百文，吉州石價二十五貫八百五十文，潭州石價二十七貫三百文，列右行。次置水脚：安吉一貫二百文，平江九百文，隆興一貫七百文，吉州二貫九百文，潭州二貫一百文，列左行，各對本州石價。以兩行數併之，得數：安吉三十貫七百文，平江三十五貫九百，隆興二十九貫八百，潭州二十九貫四百，吉州二十八貫七百五十，仍於右行。次以文思院官斗八十三合遍乗之：安吉州得二千五百四十八貫一百文，平江府得二千九百七十九貫七百文，江西隆興得二千四百七十三貫四百文，湖南潭州得二千四百四十貫二百文，江南吉州得二千三百八十六貫二百五十文，各為實於右。得次列安吉斗一百一十合，平江斗一百三十五合，隆興斗一百一十五合，潭州斗一百一十八合，吉州斗一百二十合於左行為法。以对除右行之實：安吉得二十三貫一百六十四文一十一分之六，平江得二十三貫七十一文二十七分文之二十三，隆興得二十一貫五百七文二十三分文之一十九，潭州得二十貫六百七十九文五十九文分之四十九，吉州得一十九貫八百八十五文一十二分文之五。相課石價，其安吉州最貴，平江次之，隆興又次之，潭州又次之，吉州最賤。}



\switchcolumn



\ModernWorking{%

排列安吉州石價29貫500文，平江35貫，隆興28貫100文，吉州25貫850文，潭州27貫300文在右行。再排列水脚：安吉1貫200文，平江900文，隆興1貫700文，吉州2貫900文，潭州2貫100文在左行，各對應本州石價。兩行數合併：安吉30貫700文，平江35貫900文，隆興29貫800文，潭州29貫400文，吉州28貫750文，仍在右行。再以文思院官斗83合遍乘：安吉州得2,548貫100文，平江府得2,979貫700文，隆興得2,473貫400文，潭州得2,440貫200文，吉州得2,386貫250文，各為被除數在右。再排列安吉斗110合，平江斗135合，隆興斗115合，潭州斗118合，吉州斗120合在左行為除數。以對除右行的被除數：安吉得23貫164文又6/11文，平江得22貫71文又23/27文，隆興得21貫507文又19/23文，潭州得20貫679文又49/59文，吉州得19貫885文又5/12文。比較石價，安吉州最貴，平江次之，隆興又次之，潭州又次之，吉州最賤。

}



\end{paracol}



\newpage

\subsection*{\color{sectioncolor}折解輕齎（折算解送輕便貨幣）}

\addcontentsline{toc}{subsection}{折解輕齎}



\begin{paracol}{2}



\ClassicalHeader



\ClassicalQuestion{%

有甲乙丙丁四郡，各合起上供銀絹。



\textbf{甲郡}：銀三千二百兩，每兩二貫二百文足；絹六萬四千匹，每匹二貫文足。去京一千里，每擔一里傭錢六文足。其時舊㑹每貫五十四文足。



\textbf{乙郡}：銀二千七百兩，每兩二貫三百文足；絹四萬九千二百匹，每匹二貫四百二十文足。去京九百八十里，每擔一里傭錢四文二分足。舊㑹價五十九文足。



\textbf{丙郡}：銀四千兩，每兩新㑹九貫三百文；絹七萬三千六百匹，每匹新㑹一十貫三百文。去京二千里，每擔一里傭錢八十文舊㑹。



\textbf{丁郡}：銀二千六百兩，每兩五十一貫文舊㑹；絹三萬二千三十五匹，每匹五十八貫文舊㑹。去京一千五百里，每擔一里傭錢一百文舊㑹。



諸郡銀每五百兩、絹每六十匹、新㑹每五千貫為擔。欲並折新㑹，均作三限起解。求各郡每限及本色原理、折解實用、寛餘、傭錢各新㑹幾何？}



\switchcolumn



\ModernHeader



\ModernQuestion{%

有甲乙丙丁四郡，各應上供銀絹。



\textbf{【甲郡】}：銀3,200兩，每兩2貫200文足（足陌錢）；絹64,000匹，每匹2貫文足。距京城1,000里，每擔每里傭錢6文足。當時舊會每貫54文足。



\textbf{【乙郡】}：銀2,700兩，每兩2貫300文足；絹49,200匹，每匹2貫420文足。距京城980里，每擔每里傭錢4文2分足。舊會價59文足。



\textbf{【丙郡】}：銀4,000兩，每兩新會9貫300文；絹73,600匹，每匹新會10貫300文。距京城2,000里，每擔每里傭錢80文舊會。



\textbf{【丁郡】}：銀2,600兩，每兩51貫文舊會；絹32,035匹，每匹58貫文舊會。距京城1,500里，每擔每里傭錢100文舊會。



各郡銀每500兩、絹每60匹、新會每5,000貫為一擔。想全部折算為新會，平均分三期限解送上京。求各郡每限以及本色原額、折算實用、寬餘、傭錢各折合新會多少？

}



\end{paracol}



\begin{paracol}{2}



\ClassicalNote{%

此題貫數分三項：其一足數，每千文為一貫；其一舊㑹數，如甲以五十四文為一貫，乙以五十九文為一貫是也；其一新㑹數，為舊㑹數之五倍，如甲以二百七十文為一貫，乙以二百九十五文為一貫是也。四郡或言足數，或言舊㑹數、新㑹數。並折新㑹數。傭錢原以銀五百兩或絹六十匹各為一擔，今皆折新㑹數，以五千貫為一擔，故有寛餘錢數。題語多未詳，而併為一擔句尤混，略為分析而術草之意大概可見矣。}



\switchcolumn



\ModernNote{%

這道題的貫數分三種：第一種是足數，每千文為一貫；第二種是舊會數，如甲郡以54文為一貫，乙郡以59文為一貫；第三種是新會數，為舊會數的五倍，如甲郡以270文為一貫，乙郡以295文為一貫。四郡有的說足數，有的說舊會數、新會數，都要折算為新會數。傭錢原來以銀500兩或絹60匹各為一擔，現在都折算為新會數，以5,000貫為一擔，所以有寬餘錢數。題目語言多不清楚，尤其是「併為一擔」一句特別混亂，稍作分析後術草的意思大概可以明白了。

}



\end{paracol}



\begin{paracol}{2}



\ClassicalAnswer{%

\textbf{甲郡}：合解五十萬一百四十八貫一百四十八文。初限一十六萬六千七百一十六貫四十九文，次限一十六萬六千七百一十六貫四十九文，未限一十六萬六千七百一十六貫五十文。傭錢原理二萬三千八百四十五貫九百二十五文二十七分文之二十五，實用二千二百二十二貫八百七十七文二十七分文之二十一，寛餘二萬一千六百二十三貫四十八文二十七分文之四。



\textbf{乙郡}：合解四十二萬四千六百五十七貫六百二十七文。初限一十四萬一千五百五十二貫五百四十二文，次限一十四萬一千五百五十二貫五百四十二文，末限一十四萬一千五百五十二貫五百四十三文。傭錢原理一萬一千五百一十六貫四百二十八文五十九分文之二十八，實用一千一百八十五貫一十文二百九十五分文之五，寛餘一萬三百三十一貫四百一十八文五十九分文之二十七。



\textbf{丙郡}：合解七十九萬五千二百八十貫文。初限二十六萬五千九十三貫三百三十三文，次限二十六萬五千九十三貫三百三十三文，末限二十六萬五千九十三貫三百三十四文。傭錢原理三萬九千五百九貫三百三十三文三分文之一，實用五千八十九貫七百九十二文，寛餘三萬四千四百一十九貫五百四十一文三分文之一。



\textbf{丁郡}：合解三十九萬八千一百二十六貫文。初限一十三萬二千七百八貫六百六十六文，次限一十三萬二千七百八貫六百六十六文，末限一十三萬二千七百八貫六百六十八文。傭錢原理一萬四千一百七十三貫五百文，實用二千三百八十八貫七百五十六文，寛餘一萬一千七百八十四貫七百四十四文。}



\switchcolumn



\ModernAnswer{%

\textbf{【甲郡】}：合解500,148貫148文。初限166,716貫49文，次限166,716貫49文，末限166,716貫50文。傭錢原額23,845貫925文又25/27文，實用2,222貫877文又21/27文，寬餘21,623貫48文又4/27文。



\textbf{【乙郡】}：合解424,657貫627文。初限141,552貫542文，次限141,552貫542文，末限141,552貫543文。傭錢原額11,516貫428文又28/59文，實用1,185貫10文又5/295文，寬餘10,331貫418文又27/59文。



\textbf{【丙郡】}：合解795,280貫文。初限265,093貫333文，次限265,093貫333文，末限265,093貫334文。傭錢原額39,509貫333文又1/3文，實用5,089貫792文，寬餘34,419貫541文又1/3文。



\textbf{【丁郡】}：合解398,126貫文。初限132,708貫666文，次限132,708貫666文，末限132,708貫668文。傭錢原額14,173貫500文，實用2,388貫756文，寬餘11,784貫744文。

}



\end{paracol}



\begin{paracol}{2}



\ClassicalMethod{%

以均輸求之。置各郡銀絹，乗各價，併之，歸足原，展足為舊㑹，次以五約舊㑹為新㑹，各得合解錢。以限數除之，得每限錢，不盡併歸末限。次置里數，乗毎里傭價為率。以率乗原銀及原絹，各為傭實。以每擔銀絹率各為法，實如法而一，不滿者亦為擔，併之為原理傭錢。次以率乗合觧錢為實，乃以錢物毎擔率為法，實如法而一，各得實用傭錢。以減原理傭錢，餘為寛餘傭錢。}



\switchcolumn



\ModernMethod{%

使用均輸法求解。將各郡的銀絹數，乘各價，合併，歸為足陌原數，再展為舊會，再以5約舊會為新會，各得合解錢數。以期限數（3限）去除，得每限錢數，餘數併入末限。再置里數，乘每里傭價為率。以率乘原銀及原絹，各為傭錢被除數。以每擔銀絹率各為除數，被除數除以除數，不滿一擔的也算一擔，合併為原額傭錢。再以率乘合解錢為被除數，以錢物每擔率為除數，被除數除以除數，各得實用傭錢。以原額傭錢減去實用，餘為寬餘傭錢。

}



\end{paracol}







%% =============================================================================

%% FASCICLE 6B — 錢穀 (Part 2)

%% =============================================================================

\newpage

\section*{\color{titlecolor}卷六下\quad 錢穀（貨幣和糧食）}

\addcontentsline{toc}{section}{卷六下\quad 錢穀}



\subsection*{\color{sectioncolor}僦直推原（推算租價原額）}

\addcontentsline{toc}{subsection}{僦直推原}



\begin{paracol}{2}



\ClassicalHeader



\ClassicalQuestion{%

房廊數内，一户日納一百五十六文八分足為準。指揮未曽經減者，減三分；已曽經減三分者，減二分；已曽經減二分者，更減二分。今本户累經減者，欲知原額房錢幾何？}



\switchcolumn



\ModernHeader



\ModernQuestion{%

在房廊（店鋪街房）數目中，有一户每天繳納156文8分足（足陌錢）為標準。規定未經減免的，減免三分（30\%）；已經減免三分的，再減二分（20\%）；已經減免二分的，再減二分（20\%）。現在本户累經多次減免，想知道原額房租是多少？

}



\end{paracol}



\begin{paracol}{2}



\ClassicalAnswer{%

原額三百五十文。}



\switchcolumn



\ModernAnswer{%

原額350文。

}



\end{paracol}



\begin{paracol}{2}



\ClassicalMethod{%

以衰分求之。列一十分兩行，各三位；列減分，對減右行。以餘者相乘為法。以左行原列相乘，得納錢為實。實如法而一，得原額錢。}



\switchcolumn



\ModernMethod{%

使用衰分法求解。排列十分兩行，各三位；排列減分數，對應減去右行。以餘數相乘為除數。以左行原列數相乘，得納錢為被除數。被除數除以除數，得原額錢數。

}



\end{paracol}



\begin{paracol}{2}



\ClassicalWorking{%

列一十三位於左行，又列一十分三位於右行。其右上減去初減三分，右中減去次減二分，右下減去更減二分。右行餘七八八，以相乗得四百四十八為法。乃以左行三位一十分相乘，得一千為乗率。以乘見日納錢一百五十六文八分，得一百五十六貫八百文為實。實如法而一，得三百五十文，為本户原額戸錢。}



\switchcolumn



\ModernWorking{%

排列三個「十分」在左行，又排列三個「十分」在右行。右上減去初減三分，餘七分；右中減去次減二分，餘八分；右下減去更減二分，餘八分。右行餘數7、8、8相乘得448為除數。左行三位十分相乘得1,000為乘率。以此乘現日納錢156文8分，得156,800文為被除數。以除數448去除，得350文，為本户原額房租。

}



\end{paracol}



\subsection*{\color{sectioncolor}推求本息（計算本金和利息）}

\addcontentsline{toc}{subsection}{推求本息}



\begin{paracol}{2}



\ClassicalHeader



\ClassicalQuestion{%

三庫息例：萬貫以上一釐，千貫以上二釐五毫，百貫以上三釐。甲庫本四十九萬三千八百貫，乙庫本三十七萬三百貫，丙庫本二十四萬六千八百貫。今三庫共約到息錢二萬五千六百四十四貫二百文。其典率：甲反錐差，乙方錐差，丙蒺藜差。欲知原典三例本息各幾何？}



\switchcolumn



\ModernHeader



\ModernQuestion{%

三個庫房的利息規定：萬貫以上利率1釐（0.1\%），千貫以上2.5釐（0.25\%），百貫以上3釐（0.3\%）。甲庫本金493,800貫，乙庫本金370,300貫，丙庫本金246,800貫。現在三庫共收到息錢25,644貫200文。各庫的典當率：甲用反錐差（遞減數列3,2,1），乙用方錐差（平方數列1,4,9），丙用蒺藜差（三角數列1,3,6）。想知道原典當三例中各庫本金利息各為多少？

}



\end{paracol}



\begin{paracol}{2}



\ClassicalNote{%

此即衰分題也。其差有反錐、方錐、蒺藜之名，蓋以一二三遞減如立錐為反錐，以一四九平方遞加為方錐，以一三六三數遞加為蒺藜。是必古有其名也。至以各差求各本，則因各本原依各差入之也。}



\switchcolumn



\ModernNote{%

這就是衰分題。各種差數有反錐、方錐、蒺藜的名稱，反錐是以1,2,3遞減如倒立錐形，方錐是以1,4,9平方遞加，蒺藜是以1,3,6三角數遞加。這些必定是古代已有的名稱。至於以各差求各本，是因為各本金原來就是按各差數分配的。

}



\end{paracol}



\begin{paracol}{2}



\ClassicalAnswer{%

\textbf{甲庫}共納息九千五十三貫文：一釐息二千四百六十九貫文，二釐半息四千一百一十五貫文，三釐息二千四百六十九貫文。



\textbf{乙庫}共納息一萬五十一貫文：一釐息二百六十四貫五百文，二釐半息二千六百四十五貫文，三釐息七千一百四十一貫五百文。



\textbf{丙庫}共納息六千五百四十貫二百文：一釐息二百四十六貫八百文，二釐半息一千八百五十一貫文，三釐息四千四百四十二貫四百文。}



\switchcolumn



\ModernAnswer{%

\textbf{【甲庫】}共納息9,053貫文：1釐利率息2,469貫文，2.5釐利率息4,115貫文，3釐利率息2,469貫文。



\textbf{【乙庫】}共納息10,051貫文：1釐利率息264貫500文，2.5釐利率息2,645貫文，3釐利率息7,141貫500文。



\textbf{【丙庫】}共納息6,540貫200文：1釐利率息246貫800文，2.5釐利率息1,851貫文，3釐利率息4,442貫400文。

}



\end{paracol}



\begin{paracol}{2}



\ClassicalMethod{%

置諸庫諸色之差，照釐率為三行。縱併之為約率，横命之為乗率。以約率各約自庫之本，各得。以遍乗未併乗率，然後各以釐率横乗之，次以縱併之，為各庫共息。}



\switchcolumn



\ModernMethod{%

排列各庫各等級的差數，按釐率排為三行。縱向合併為約率，橫向命名為乘率。以約率各約本庫的本金，各得數。以此遍乘未合併的乘率，然後各以釐率橫向相乘，再縱向合併，為各庫共息。

}



\end{paracol}



\begin{paracol}{2}



\ClassicalWorking{%

置甲庫反錐差，自下置三二一於右行。次置乙庫方錐差，自上置一四九於中行。次置丙庫蒺藜差，自上置一三六於左行，各為三庫上中下三等乘率。乃縱併甲差三二一，得六為甲約率。縱併乙差一四九，得一十四為乙約率。縱併丙差一三六，得一十為丙約率。直命九位數各為上中下乗率。乃先以約率各約自庫之本：乃以甲約率六約甲本四十九萬三千八百貫，得八萬二千三百貫為甲得。次以乙約率一十四約乙本三十七萬三百貫，得二萬六千四百五十貫為乙得。次以丙約率一十約丙本二十四萬六千八百貫，得二萬四千六百八十貫為丙得。以各得乗未併乗率：其甲所得八萬二千三百貫，乘反錐乗率三二一，得二十四萬六千九百貫為上率，得一十六萬四千六百貫為中率，得八萬二千三百貫為下率。其乙所得二萬六千四百五十貫，以乗方錐差一四九，得二萬六千四百五十貫為上率，得一十萬五千八百貫為中率，得二十三萬八千五十貫為下率。其丙所得二萬四千六百八十貫，以乗蒺藜差一三六，得二萬四千六百八十貫為上率，得七萬四千四十貫為中率，得一十四萬八千八十貫為下率。然後各以息釐數乘各庫三乗：其甲以一釐乗上率二十四萬六千九百貫，得二千四百六十九貫為上息；以二釐五毫乘中率一十六萬四千六百貫，得四千一百一十五貫為中息；以三釐乘下率八萬二千三百貫，得二千四百六十九貫為下息。併上中下三息，得九千五十三貫文為甲庫共息。其乙庫以一釐乗上率二萬六千四百五十貫，得二百六十四貫五百文為上息；以二釐五毫乗中率一十萬五千八百貫，得二千六百四十五貫為中息；以三釐乗下率二十三萬八千五十貫，得七千一百四十一貫五百為下息。併上中下三息，得一萬五十一貫文為乙庫共息。其丙庫以一釐乗上率二萬四千六百八十貫，得二百四十六貫八百為上息；以二釐五毫乘中率七萬四千四十貫，得一千八百五十一貫為中息；以三釐乗下率一十四萬八千八十貫，得四千四百四十二貫四百文為下息。併上中下三息，得六千五百四十貫二百文為丙庫共息。併三庫共息，得二萬五千六百四十四貫二百文為總息。}



\switchcolumn



\ModernWorking{%

排列甲庫反錐差，自下排3,2,1在右行。再排乙庫方錐差，自上排1,4,9在中行。再排丙庫蒺藜差，自上排1,3,6在左行，各為三庫上中下三等乘率。縱向合併甲差3+2+1=6為甲約率。縱向合併乙差1+4+9=14為乙約率。縱向合併丙差1+3+6=10為丙約率。直接命名九位數各為上中下乘率。



先以約率各約本庫本金：甲約率6約甲本493,800貫，得82,300貫為甲得。乙約率14約乙本370,300貫，得26,450貫為乙得。丙約率10約丙本246,800貫，得24,680貫為丙得。



以各得數乘未合併乘率：甲得82,300貫乘反錐乘率3,2,1，得246,900貫為上率，164,600貫為中率，82,300貫為下率。乙得26,450貫乘方錐差1,4,9，得26,450貫為上率，105,800貫為中率，238,050貫為下率。丙得24,680貫乘蒺藜差1,3,6，得24,680貫為上率，74,040貫為中率，148,080貫為下率。



然後各以息釐數乘各庫三率：甲庫以1釐乘上率246,900貫，得2,469貫為上息；以2.5釐乘中率164,600貫，得4,115貫為中息；以3釐乘下率82,300貫，得2,469貫為下息。合併三息，得9,053貫為甲庫共息。



乙庫以1釐乘上率26,450貫，得264貫500文為上息；以2.5釐乘中率105,800貫，得2,645貫為中息；以3釐乘下率238,050貫，得7,141貫500文為下息。合併三息，得10,051貫為乙庫共息。



丙庫以1釐乘上率24,680貫，得246貫800文為上息；以2.5釐乘中率74,040貫，得1,851貫為中息；以3釐乘下率148,080貫，得4,442貫400文為下息。合併三息，得6,540貫200文為丙庫共息。三庫共息合計25,644貫200文為總息。

}



\end{paracol}



\subsection*{\color{sectioncolor}易牒知原（憑度牒推算原數）}

\addcontentsline{toc}{subsection}{易牒知原}



\begin{paracol}{2}



\ClassicalHeader



\ClassicalNote{%

按舊本此問無題，今増入。}



\switchcolumn



\ModernNote{%

按舊版本此問無題，現增補。

}



\end{paracol}



\begin{paracol}{2}



\ClassicalQuestion{%

出度牒差人營運：毎三道易鹽一十三袋，鹽二袋易布八十四疋，布一十五疋易絹三疋半，絹六疋易銀七兩二錢。今趂到銀九千一百七十二兩八錢。欲知原闗度牒道數幾何？}



\switchcolumn



\ModernHeader



\ModernQuestion{%

發出度牒（出家人憑證）差遣人員經營：每3道度牒換13袋鹽，2袋鹽換84疋布，15疋布換3.5疋絹，6疋絹換7兩2錢銀子。現已收到銀9,172兩8錢。想知道原來發出多少道度牒？

}



\end{paracol}



\begin{paracol}{2}



\ClassicalAnswer{%

度牒一百八十道。}



\switchcolumn



\ModernAnswer{%

度牒180道。

}



\end{paracol}



\begin{paracol}{2}



\ClassicalMethod{%

以粟米互乘易法求之。列各數以本色相對，如鴈翅。以多一事者相乗為實，以少一事者相乘為法，除之。}



\switchcolumn



\ModernMethod{%

使用粟米互換法求解。排列各數以本色相對，如雁翅形排列。以多一項的數相乘為被除數，以少一項的數相乘為除數，相除得到答案。

}



\end{paracol}



\begin{paracol}{2}



\ClassicalWorking{%

先以度牒三道乗鹽二袋，得六。以乘布一十五，得九十。又乗絹六疋，得五百四十。乃乗銀九萬一千七百二十八錢，得四千九百五十三萬三千一百二十錢為實。次以鹽一十三袋乗布八十四，得一千九百九十二。以乗絹三疋五分，得三千八百二十二。乃乘銀七兩二錢，得二十七萬五千一百八十四錢為法。除實，得一百八十道為原闗度牒。}



\switchcolumn



\ModernWorking{%

先以度牒3道乘鹽2袋，得6。再乘布15，得90。再乘絹6疋，得540。再乘銀91,728錢（9,172兩8錢），得495,331,200錢為被除數。再以鹽13袋乘布84，得1,092。乘絹3.5疋，得3,822。再乘銀7兩2錢（72錢），得275,184錢為除數。被除數除以除數，得180道為原發度牒數。

}



\end{paracol}



\subsection*{\color{sectioncolor}粟米交易（糧食交易）}

\addcontentsline{toc}{subsection}{粟米交易}



\begin{paracol}{2}



\ClassicalNote{%

按舊本此問無題，今増。}



\switchcolumn



\ModernNote{%

按舊版本此問無題，現增補。

}



\end{paracol}



\begin{paracol}{2}



\ClassicalQuestion{%

菽三升易小麥二升，小麥一升五合易油麻八合，油麻一升二合易粳米一升八合。今將菽一十四石四斗欲易油麻，又將小麥二十一石六斗欲易粳米。問各幾何？}



\switchcolumn



\ModernHeader



\ModernQuestion{%

菽（豆類）3升換小麥2升，小麥1升5合換油麻8合，油麻1升2合換粳米1升8合。現有菽14石4斗想換油麻，又有小麥21石6斗想換粳米。問各能換多少？

}



\end{paracol}



\begin{paracol}{2}



\ClassicalAnswer{%

油麻五石一斗二升；粳米一十七石二斗八升。}



\switchcolumn



\ModernAnswer{%

油麻5石1斗2升；粳米17石2斗8升。

}



\end{paracol}



\begin{paracol}{2}



\ClassicalMethod{%

以粟米換易求之。置原易率，本色對列，如鴈翅。以多一事者相乘為實，以少一事者相乗為法，除之。各得或問數，不干其率者不置。}



\switchcolumn



\ModernMethod{%

使用粟米換易法求解。排列原始換易比率，本色相對，如雁翅形。以多一項的數相乘為被除數，以少一項的數相乘為除數，相除。各得所求數，不相關的比率不列入。

}



\end{paracol}



\begin{paracol}{2}



\ClassicalWorking{%

置四色六數，列六位率，如鴈翅，皆化為合。先將菽一十四石四斗化作一萬四千四百合，乃對前二句率數四位，如鴈翅，至欲易油麻止，共五事為上圖。次將小麥二十一石六斗化為二萬一千六百合，乃對後兩句率四位，鴈翅，至欲粳米止，共五事為下圖。



下圖：其上圖以菽一萬四千四百合乘麥二十，得二十八萬八千，又乗油麻八合，得二百三十萬四千合為油麻實。次以菽三十合乘麥一十五合，得四百五十合為法。除之，得五千一百二十合，展為五石一斗二升為油麻。



其下圖以小麥二萬一千六百合乗油麻八合，得一十七萬二千八百合，又乗粳米一十八合，得三百一十一萬四百合為粳米實。以小麥一升五合乗油麻一十二合，得一百八十合為法。除之，得一萬七千二百八十，展作一十七石二斗八升為粳米。}



\switchcolumn



\ModernWorking{%

排列四種糧食六個數字，列為六個比率，如雁翅形，都化為合（1升=10合）。先將菽14石4斗化為14,400合，對應前兩句比率四個數字，如雁翅形，到要換油麻為止，共五項為上圖。再將小麥21石6斗化為21,600合，對應後兩句比率四個數字，雁翅形，到要換粳米為止，共五項為下圖。



上圖：以菽14,400合乘小麥比率2（即20合），得288,000，再乘油麻8合，得2,304,000合為油麻被除數。再以菽3升（30合）乘小麥1.5升（15合），得450合為除數。2,304,000除以450得5,120合，展為5石1斗2升為油麻。



下圖：以小麥21,600合乘油麻8合，得172,800合，再乘粳米18合，得3,114,000合為粳米被除數。以小麥1.5升（15合）乘油麻1.2升（12合），得180合為除數。3,114,000除以180得17,280合，展為17石2斗8升為粳米。

}



\end{paracol}



\subsection*{\color{sectioncolor}計米易麴（計算換米製麴）}

\addcontentsline{toc}{subsection}{計米易麴}



\begin{paracol}{2}



\ClassicalNote{%

按舊本此問無題，今増。}



\switchcolumn



\ModernNote{%

按舊版本此問無題，現增補。

}



\end{paracol}



\begin{paracol}{2}



\ClassicalQuestion{%

庫率：粳榖七石出米三石，糯米一斗易小麥一斗七升，小麥五升踏麴二斤四兩，麴一百一十斤醞糯米一石三斗。今有糯穀一千七百五十九石三斗八升，欲出穀做米，易麥踏麴，還自醞，餘穀之米須令適足。各合幾何？}



\switchcolumn



\ModernHeader



\ModernQuestion{%

庫房標準：粳穀7石出米3石，糯米1斗換小麥1斗7升，小麥5升可製麴2斤4兩，麴110斤可釀糯米1石3斗。現有糯穀1,759石3斗8升，想出穀做米，換麥製麴，回來釀酒，剩餘穀出的米必須恰好夠用。各應該是多少？

}



\end{paracol}



\begin{paracol}{2}



\ClassicalAnswer{%

共穀一千七百五十九石三斗八升。出穀九百二十四石，得米三百九十六石。易麥六百七十三石二斗。踏麴三萬二百九十四斤。餘榖八百三十五石三斗八升。醞米三百五十八石二升。}



\switchcolumn



\ModernAnswer{%

共穀1,759石3斗8升。出穀924石，得米396石。換小麥673石2斗。製麴30,294斤。餘穀835石3斗8升。釀酒用米358石2升。

}



\end{paracol}



\begin{paracol}{2}



\ClassicalMethod{%

以粟米換易求之。置諸率，隨本色對列，如鴈翅。有分者通之，異類者變之，以頭位者進乗之，以下位退乗之，得合數。有對者相乗之，無對者直命之，為諸率。併上下無對者為法率。以今有物徧乘諸率（不乗法率），各為實。諸實並如法而一，各得其已變者，復互易乗除之，即得所求。}



\switchcolumn



\ModernMethod{%

使用粟米換易法求解。排列各個比率，隨本色相對，如雁翅形。有分數的通分，不同類的變換，以頭位進乘，以下位退乘，得到合數。有對應的相乘，無對應的直接命名，為各率。合併上下無對應的為法率。以現有穀物遍乘各率（不乘法率），各為被除數。各被除數都以法去除，各得已變換的數，再互易乘除，即得所求。

}



\end{paracol}



\begin{paracol}{2}



\ClassicalWorking{%

置糯穀七、出米三於右行上副兩位。次置糯米一斗、麥一斗七升於副行副中兩位。次置小麥五升、踏麥二斤四兩於次行中次兩位。次置麴一百一十觔、醞米一石三斗於左行次下兩位，隨本色對列，如鴈翅。訖乃驗次行二斤四兩是四分斤之一，以母四通次行兩位，以子一内次行次位，具中位得二，次位得九。又驗左行下位是糯米，是異類，於糯米合變為糯穀，乃以問中首句率穀七米三變之，以七因米一石三斗得九石一斗於左下為穀，却以米三因麴一百一十斤為三百三十斤麴於左行，得變圖數。以左行三百三十乗次行二，得六百六十。次以六百六十乘副行一十，得六千六百。次以六千六百乗右上七，得四萬六千二百，各於原位。却以右行副位三因副行一斗七升，得五斗一升。又以五斗一升乗次行九，得四百五十九。又以四百五十九乗左下九十一，得四萬一千七百六十九，列為合圖數。乃驗合圖四行，其副中次三位有對，以對相乗合之。其右上左下無對者，直命之，皆為率。列右行：上得四萬六千二百為出糯穀率，副位得一萬九千八百為得糯米率，中得三萬三千六百六十為易得麥率，次得一萬五千一十四七為踏到麴率，下得四萬一千七百六十九為餘下糯榖率。併上下率，共得八萬七千九百六十九為法率。今六率共求等得一，約之，只得原率為率圖。始用今有穀一千七百五十九石三斗八升，皆化為升，徧乘五率（不乗法率），得八十一億二千八百三十三萬五千六百升為出穀實，得三十四億八千三百五十七萬二千四百升為糯米實，得五十九億二千二百七萬三千八十升為易麥實，得二十六億六千四百九十三萬二千八百八十六為踏麴實，得七十三億四千八百七十五萬四千三百二十二升為餘穀實。其五實皆如法八萬七千九百六十九而一：得九百二十四石為出穀，得三百九十六石為做到糯米，得六百七十三石二斗為易到小麥，得三萬二百九十四斤為踏到麴，得八百三十五石三斗八升為餘下穀。今將餘下穀變為米，乃以乗率三因餘穀八百三十五石三斗八升，得二千五百六石一斗四升為實，以糯穀率七為法除之，得三百五十八石二升為醞米。}



\switchcolumn



\ModernWorking{%

排列糯穀7、出米3在右行上副兩位。再排糯米1斗、小麥1斗7升在副行副中兩位。再排小麥5升、製麴2斤4兩在次行中次兩位。再排麴110斤、釀米1石3斗在左行次下兩位，隨本色相對，如雁翅形。檢查次行2斤4兩是9/4斤（2.25斤），以分母4通次行兩位，分子1加入次行次位，中位得2，次位得9。再檢查左行下位是糯米，是異類，應變為糯穀，以首句率穀7米3變換，以7乘米1石3斗得9石1斗於左下為穀，再以米3乘麴110斤為330斤麴於左行，得變換圖數。



以左行330乘次行2，得660。再以660乘副行10（1斗），得6,600。再以6,600乘右上7，得46,200。以右行副位3乘副行1斗7升，得5斗1升（51升）。再以51升乘次行9，得459。再以459乘左下91，得41,769。



檢查合圖四行，副中次三位有對應，以對相乘合併。右上左下無對應，直接命名為率。右行：上得46,200為出糯穀率，副位19,800為得糯米率，中33,660為易小麥率，次15,007.5為踏麴率，下41,769為餘穀率。合併上下率得87,969為法率。



以現有穀1,759石3斗8升化為升，遍乘五率（不乘法率），得各被除數。五個被除數都以法87,969去除：得924石為出穀，396石為糯米，673石2斗為小麥，30,294斤為麴，835石3斗8升為餘穀。將餘穀變為米，以乘率3乘餘穀835石3斗8升，得2,506石1斗4升為被除數，以糯穀率7為法除之，得358石2升為釀酒用米。

}



\end{paracol}



\subsection*{\color{sectioncolor}囘運費（收回運費）}

\addcontentsline{toc}{subsection}{囘運費}



\begin{paracol}{2}



\ClassicalHeader



\ClassicalQuestion{%

有江西水運米一十二萬三千四百石，原係鎮江交卸，計水程二千一百三十里，每石水脚錢一貫二百文。今截上件米就池州安頓，池州至鎮江八百八十里。欲收囘不該水脚錢幾何？}



\switchcolumn



\ModernHeader



\ModernQuestion{%

有江西水運米123,400石，原來到鎮江交卸，計水程2,130里，每石水脚運費1貫200文。現將這批米截留在池州安頓，池州至鎮江880里。想收回不應付的水脚錢有多少？

}



\end{paracol}



\begin{paracol}{2}



\ClassicalAnswer{%

收囘錢六萬一千一百七十八貫五百九十一文。}



\switchcolumn



\ModernAnswer{%

收回錢61,178貫591文。

}



\end{paracol}



\begin{paracol}{2}



\ClassicalMethod{%

以粟米互易求之。置池州至鎮江里數，乗水脚錢，得數又乗運米為實。以原至鎮江水程為法，除實，得收囘錢。}



\switchcolumn



\ModernMethod{%

使用粟米互易法求解。將池州至鎮江的里數，乘水脚錢，所得再乘運米數為被除數。以原來到鎮江的水程為除數，被除數除以除數，得應收回的錢數。

}



\end{paracol}



\begin{paracol}{2}



\ClassicalWorking{%

置池州至鎮江八百八十里，乗毎石水脚錢一貫二百，得一千五十六貫文。又乗運米一十二萬三千四百石，得一億三千三十一萬四百貫文為實。以原至鎮江水程二千一百三十里為法，除實，得六萬一千一百七十八貫五百九十一文為收囘錢。}



\switchcolumn



\ModernWorking{%

將池州至鎮江880里，乘每石水脚錢1貫200文，得1,056貫文。再乘運米123,400石，得130,310,400貫文為被除數。以原水程2,130里為除數，除被除數，得61,178貫591文為應收回的錢。

}



\end{paracol}



\subsection*{\color{sectioncolor}三色均價（三色金均價）}

\addcontentsline{toc}{subsection}{三色均價}



\begin{paracol}{2}



\ClassicalHeader



\ClassicalQuestion{%

庫有三色金共五千兩：内八分金一千二百五十兩，兩價四百貫文；七分五釐金一千六百兩，兩價三百七十五貫文；八分五釐金二千一百五十兩，兩價四百二十五貫文。並欲煉為足色，每兩工食藥炭錢三貫文，耗金九百七十二兩五錢。欲知色分及兩價各幾何？}



\switchcolumn



\ModernHeader



\ModernQuestion{%

庫房有三色金共5,000兩：其中成色80\%的金1,250兩，每兩價400貫文；成色75\%的金1,600兩，每兩價375貫文；成色85\%的金2,150兩，每兩價425貫文。都想冶煉為純金（100\%成色），每兩工錢、食錢、藥炭錢3貫文，損耗972兩5錢。想知道冶煉後的成色分數以及每兩價格各為多少？

}



\end{paracol}



\begin{paracol}{2}



\ClassicalAnswer{%

色十分；兩價五百三貫七百二十四文，五百三十七分文之二百一十二。}



\switchcolumn



\ModernAnswer{%

成色10分（100\%純金）；每兩價格503貫724文又212/537文。

}



\end{paracol}



\begin{paracol}{2}



\ClassicalMethod{%

以方田及粟米求之。置共數，以耗減之，餘為法。以三色分數各乗兩數，併之為色分實。以三色價數各乗兩數為寄，以工藥價乗共金，併寄，共為價實。二實皆如法而一，即各得。}



\switchcolumn



\ModernMethod{%

使用方田法及粟米法求解。將總數減去損耗，餘數作為除數。以三色成色分數各乘兩數，合併為成色被除數。以三色價格各乘兩數為寄數，以工藥價乘總金數，合併寄數，共為價格被除數。二個被除數都以除數去除，即得各數。

}



\end{paracol}



\begin{paracol}{2}



\ClassicalWorking{%

置共金五千兩，減耗九百七十二兩五錢，外餘四千二十七兩五錢為法。次置一千二百五十兩，乘八分，得一萬分於上。置一千六百兩，以七分五釐乗之，得一萬二千分，加上。置二千一百五十兩，乗八分五釐，得一萬八千二百七十五分，又加上，共得四萬二百七十五分為分實。次置一千二百五十兩乗四百貫，得五十萬貫為寄。次置一千六百兩乘價三百七十五貫，得六十萬貫，加寄。次置二千一百五十兩乗價四百二十五貫，得九十一萬三千七百五十貫，又加寄。次置共金五千兩，乗工藥錢三貫，得一萬五千貫，又加寄，共得二百二萬八千七百五十貫為貫實。二實並如法四千二十七兩五錢而一：得其色一十分；其價每兩得五百三貫七百二十四文，不盡一貫五百九十。與法求等，得七十五，俱約之，為五百三十七文之二百一十二。}



\switchcolumn



\ModernWorking{%

將總金5,000兩，減去損耗972兩5錢，餘4,027兩5錢為除數。再排1,250兩乘80\%（0.8），得10,000分。排1,600兩乘75\%（0.75），得12,000分，加上。排2,150兩乘85\%（0.85），得18,275分，再加，共得40,275分為成色被除數。



再排1,250兩乘400貫，得500,000貫為寄。1,600兩乘375貫，得600,000貫，加入。2,150兩乘425貫，得913,750貫，再加入。總金5,000兩乘工藥錢3貫，得15,000貫，再加入，共得2,028,750貫為價格被除數。



二個被除數都以除數4,027兩5錢去除：得成色10分；每兩價格503貫724文，餘1貫590文。與除數求最大公約數得75，各約之，為537分之212。

}



\end{paracol}







%% =============================================================================

%% END MATTER

%% =============================================================================

\newpage

\section*{\color{titlecolor}編輯說明 · Editorial Notes}

\addcontentsline{toc}{section}{編輯說明}



\begin{paracol}{2}



\ClassicalHeader



本數字版秦九韶《數學九章》卷五至卷六呈現「賦役」與「錢穀」兩章的完整文本，轉錄自中國維基文庫所存四庫全書版本。



\subsection*{文本結構}



每題遵循傳統四段式結構：



\begin{enumerate}

\item \textbf{問} — 題目陳述，提出實際問題及已知量

\item \textbf{答} — 答案，列出所有要求的數值結果

\item \textbf{術} — 算法，以概括性語言描述數學方法

\item \textbf{草} — 演算，展示詳細的逐步計算過程

\end{enumerate}



\subsection*{歷史背景}



秦九韶（約1202--1261），字道古，宋代最傑出的數學家之一。其《數學九章》完成於1247年，分為九大類：



\begin{enumerate}

\item \textbf{大衍} — 大衍求一術（一次同餘式組）

\item \textbf{天時} — 曆法計算

\item \textbf{田域} — 田地測量

\item \textbf{測望} — 遙測望算

\item \textbf{賦役} — 賦稅和徭役

\item \textbf{錢穀} — 貨幣和糧食

\item \textbf{營建} — 工程營建

\item \textbf{軍旅} — 軍事後勤

\item \textbf{市易} — 市場貿易

\end{enumerate}



\subsection*{編輯慣例}



標記為［按］的注釋是四庫全書版本中固有的編者注，常指出原文的錯誤或題目與計算過程之間的不一致。這些注釋是清代編纂四庫全書時的館臣所加。



\switchcolumn



\ModernHeader



本對照版採用左右雙欄平行排版：左欄為四庫全書《數學九章》卷五「賦役類」與卷六「錢穀類」原文，右欄為白話文譯文及現代數學解釋。



\subsection*{對照體例}



\begin{itemize}[leftmargin=1.5em,topsep=2pt]

\item 左欄【原文】保留文言文原貌，含問、答、術、草、按全套結構

\item 右欄【譯文】以現代漢語翻譯，並將傳統術語轉換為現代數學表述

\item 傳統單位（貫、石、斗、升、合、勺、畝、步等）保留原文，右欄括註換算關係

\item 分數以「又$n/m$」格式呈現，與古典「$m$分之$n$」對應

\end{itemize}



\subsection*{術語對照}



\begin{itemize}[leftmargin=1.5em,topsep=2pt]

\item 賦役 $\to$ 賦役（賦稅和徭役）

\item 錢穀 $\to$ 錢穀（貨幣和糧食）

\item 衰分 $\to$ 衰分（等差比例分配法）

\item 均科 $\to$ 均攤（按比例分攤）

\item 折解 $\to$ 折解（折算和轉運）

\item 粟米 $\to$ 粟米換易（糧食折算交換）

\item 推求本息 $\to$ 計算本金和利息

\item 互乘易法 $\to$ 連鎖比例法（多項連續換算）

\end{itemize}



\subsection*{數學方法概覽}



卷五「賦役類」主要運用衰分法（等差/等比分配）解決各級稅賦的分攤問題；卷六「錢穀類」則以粟米互換法處理複雜的貨幣、糧食多步折算。秦九韶將這些源自《九章算術》的古典算法發展到極高的實用水平，題設數據龐大、計算繁複，體現了宋代數學的高度成就。



\end{paracol}



\vspace{2cm}

\begin{center}

\textcolor{rulecolor}{\rule{0.5\textwidth}{0.4pt}}



\vspace{1em}

{\small\color{sectioncolor}

\textit{Modern LaTeX edition}\\[0.3em]

文言文—白話文雙語對照排版\\[0.3em]

Traditional Chinese mathematical text structure preserved\\[0.3em]

Public Domain --- \textit{Siku Quanshu} Edition}

\end{center}



\end{document}







% ============================================================

% Source: translations/zh/chinese/qin-jiushao-shuxue-jiuzhang-fascicles7-9_classical-modern_bilingual.tex

% ============================================================



%% ========================================================================

%% Qin Jiushao — Shuxue Jiuzhang (Mathematical Treatise in Nine Sections)

%% Fascicles 7–9: Bilingual Classical Chinese → Modern Chinese Edition

%% ========================================================================

%% Engine: XeLaTeX

%% Layout: Parallel columns (LEFT: Classical Chinese, RIGHT: Modern Chinese)

%% ========================================================================

\documentclass[11pt,a4paper]{article}



%% -------- Core Packages --------

\usepackage{fontspec}

\usepackage{polyglossia}

\usepackage{amsmath,amssymb}

\usepackage{geometry}

\usepackage{graphicx}

\usepackage{xcolor}

\usepackage{booktabs}

\usepackage{longtable}

\usepackage{xeCJK}

\usepackage{paracol}

\usepackage{ulem}



%% -------- Page Geometry --------

\geometry{

  paperwidth=210mm,paperheight=297mm,

  left=18mm,right=18mm,top=25mm,bottom=25mm,

  headheight=14pt,footskip=30pt

}



%% -------- Column Ratio --------

\columnratio{0.50,0.50}



%% -------- Language Setup --------

\setmainlanguage{english}



%% -------- Font Configuration --------

\setmainfont{Times New Roman}

\setsansfont{Arial}

\setmonofont{Courier New}



%% Chinese Fonts (xeCJK) — using Microsoft YaHei

\setCJKmainfont{Microsoft YaHei}

\setCJKsansfont{Microsoft YaHei}



%% -------- Colors --------

\definecolor{titlecolor}{RGB}{45,55,72}

\definecolor{sectioncolor}{RGB}{55,65,81}

\definecolor{problemcolor}{RGB}{120,40,20}

\definecolor{answercolor}{RGB}{20,80,100}

\definecolor{methodcolor}{RGB}{60,50,90}

\definecolor{workingcolor}{RGB}{40,80,40}

\definecolor{rulecolor}{RGB}{180,160,140}

\definecolor{classicalbg}{RGB}{255,250,245}

\definecolor{modernbg}{RGB}{245,250,255}

\definecolor{classicallabel}{RGB}{139,69,19}

\definecolor{modernlabel}{RGB}{30,100,140}



%% -------- Header/Footer --------

\usepackage{fancyhdr}

\pagestyle{fancy}

\fancyhf{}

\fancyhead[L]{\small\textcolor{classicallabel}{文言文}\quad|\quad\textcolor{modernlabel}{白话文}}

\fancyhead[R]{\small\thepage}

\renewcommand{\headrulewidth}{0.4pt}

\renewcommand{\footrulewidth}{0pt}

\setlength{\headheight}{14pt}



%% -------- Custom Commands for Bilingual Layout --------

% Horizontal rules

\newcommand{\longline}{\noindent\rule{\textwidth}{0.4pt}}

\newcommand{\shortline}{\noindent\rule{0.5\textwidth}{0.4pt}\par}



% Classical/Modern labels

\newcommand{\classicalLabel}{\textcolor{classicallabel}{\textbf{【文言文】}}}

\newcommand{\modernLabel}{\textcolor{modernlabel}{\textbf{【白话文】}}}



% Bilingual environment using paracol

\newenvironment{bilingual}{%

  \begin{paracol}{2}

}{%

  \end{paracol}

}



% Switch to classical (left) column

\newcommand{\classical}{\switchcolumn[0]}

% Switch to modern (right) column  

\newcommand{\modern}{\switchcolumn[1]}



% Both columns side by side (for headers)

\newcommand{\bilingualHeader}[2]{%

  \begin{paracol}{2}

  #1

  \switchcolumn

  #2

  \end{paracol}

  \vspace{0.3em}

}



% Problem (问) — bilingual

\newcommand{\bilingualProblem}[2]{%

  \vspace{0.6em}%

  \begin{paracol}{2}%

  \noindent\textcolor{problemcolor}{\textbf{問：}}#1%

  \switchcolumn%

  \noindent\textcolor{problemcolor}{\textbf{【问题】}}#2%

  \end{paracol}%

  \vspace{0.3em}%

}



% Answer (答) — bilingual

\newcommand{\bilingualAnswer}[2]{%

  \begin{paracol}{2}%

  \noindent\textcolor{answercolor}{\textbf{答曰：}}#1%

  \switchcolumn%

  \noindent\textcolor{answercolor}{\textbf{【解答】}}#2%

  \end{paracol}%

  \vspace{0.3em}%

}



% Method (术) — bilingual

\newcommand{\bilingualMethod}[2]{%

  \begin{paracol}{2}%

  \noindent\textcolor{methodcolor}{\textbf{術曰：}}#1%

  \switchcolumn%

  \noindent\textcolor{methodcolor}{\textbf{【算法】}}#2%

  \end{paracol}%

  \vspace{0.3em}%

}



% Working (草) — bilingual

\newcommand{\bilingualWorking}[2]{%

  \begin{paracol}{2}%

  \noindent\textcolor{workingcolor}{\textbf{草曰：}}#1%

  \switchcolumn%

  \noindent\textcolor{workingcolor}{\textbf{【演算】}}#2%

  \end{paracol}%

  \vspace{0.3em}%

}



% Section header (bilingual)

\newcommand{\bilingualSection}[4]{%

  \section{#1 \quad #2}

  \label{#4}

  \begin{center}

  \fbox{\parbox{0.7\textwidth}{\centering

  \Large\textbf{#1}\quad\textbf{#3}\\[0.3em]

  \textit{#2}

  }}

  \end{center}

  \medskip

}



% Subsection header (bilingual)

\newcommand{\bilingualSubsection}[2]{%

  \subsection{#1}

  \begin{center}

  {\small\color{sectioncolor}\textit{#2}}

  \end{center}

  \vspace{0.3em}

}



%% -------- Hyperref Setup --------

\usepackage{hyperref}

\hypersetup{

  colorlinks=true,

  linkcolor=sectioncolor,

  citecolor=sectioncolor,

  urlcolor=sectioncolor,

  pdfencoding=auto,

  pdftitle={秦九韶数学九章 卷七至卷九 文言白话对照版},

  pdfauthor={（宋）秦九韶 撰；白话译文}

}



%% ========================================================================

\begin{document}



%% ========================================================================

%% Title Page

%% ========================================================================

\begin{titlepage}

\centering

\vspace*{1.5cm}



{\fontsize{26}{34}\selectfont\color{titlecolor}

\textbf{數學九章}

}



\vspace{1cm}



{\fontsize{18}{24}\selectfont\color{sectioncolor}

卷七至卷九 —— 文言白话对照全译本

}



\vspace{0.5cm}



{\large\color{sectioncolor}

\textit{Fascicles 7--9 — Classical & Modern Chinese Bilingual Edition}

}



\vspace{2cm}



{\large

\begin{tabular}{rl}

\textbf{原著} & （宋）秦九韶 撰 \\

\textbf{Original} & Qin Jiushao (Song Dynasty) \\

\textbf{分類} & 營建類 · 軍旅類 · 市物類 \\

\textbf{Categories} & Construction · Military · Trade \\

\end{tabular}

}



\vspace{1.5cm}



\fbox{\parbox{0.75\textwidth}{\centering

\vspace{0.5em}

\textbf{左栏：文言文原文（Classical Chinese）}\\[0.3em]

\textbf{右栏：白话文译文（Modern Chinese Vernacular）}\\[0.3em]

\textit{Left: Original Classical Chinese}\\

\textit{Right: Modern Chinese Translation}

\vspace{0.5em}

}}



\vspace{2cm}



{\large

\begin{tabular}{lll}

\textbf{卷七} & 營建類（建筑工程）& Construction \\

\textbf{卷八} & 軍旅類（军事后勤）& Military Logistics \\

\textbf{卷九} & 市物類（市场贸易）& Market Goods \\

\end{tabular}

}



\vspace{1.5cm}



{\color{rulecolor}\rule{0.6\textwidth}{0.6pt}}



\vspace{0.5cm}



{\small\color{sectioncolor}

排版引擎：\XeLaTeX\quad 中文字体: modern Unicode fonts\\

原文来源：四库全书本《数书九章》\\

白话译文：据原文译出，保留问/答/术/草结构

}



\end{titlepage}



%% ========================================================================

%% Table of Contents

%% ========================================================================

\tableofcontents

\newpage



%% ========================================================================

%% INTRODUCTION

%% ========================================================================

\section{导言 Introduction}



\begin{paracol}{2}

《数学九章》为南宋数学家秦九韶（约1202—1261）所著，成书于淳佑七年（1247年）。全书共分九卷，沿袭《九章算术》之体例而有所创新，为宋元数学之巅峰代表作之一。秦九韶博学多才，兼通天文、历法、营造、军事，故其书所设问题皆切于南宋社会之实用。



本书译注范围为其第七至第九卷：

\begin{itemize}

  \item \textbf{卷七：营建类}——筑城、开渠、建坝、修仓等土木工程算题

  \item \textbf{卷八：军旅类}——方阵、圆营、粮草、衣装等军事后勤算题

  \item \textbf{卷九：市物类}——市籴、均货、盐茶定价、钱陌换算等商业算题

\end{itemize}



每题依古法分为四段：\textbf{问}（提出问题）、\textbf{答}（给出数值答案）、\textbf{术}（说明算法原理）、\textbf{草}（详细演算步骤）。本译本左栏保留原文，右栏译为现代白话，力求准确传达数学内容，同时保留古典数学文献之韵味。

\switchcolumn

《数学九章》是南宋数学家\textbf{秦九韶}（约1202—1261年）于1247年完成的数学巨著。全书继承并发展了《九章算术》的体例，是中国古代数学最重要的经典之一。秦九韶学识渊博，精通天文、历法、建筑工程和军事学，因此书中问题都与南宋社会实际需要密切相关。



本版翻译范围为第七至第九卷：

\begin{itemize}

  \item \textbf{第七卷：营建类（建筑工程）}——包括筑城、开渠、建坝、修粮仓等土木工程的计算问题

  \item \textbf{第八卷：军旅类（军事后勤）}——包括方阵布置、圆形营地、粮草供应、军衣制作等军事后勤问题

  \item \textbf{第九卷：市物类（市场贸易）}——包括购粮、均分货物、盐茶定价、货币换算等商业数学问题

\end{itemize}



每道题按照古典数学格式分为四个部分：\textbf{【问题】}（提出问题）、\textbf{【解答】}（给出数值答案）、\textbf{【算法】}（说明解题方法）、\textbf{【演算】}（详细计算步骤）。本版左栏为文言文原文，右栏为现代汉语翻译，力求在准确传达数学内容的同时保留古典数学文献的韵味。

\end{paracol}



\medskip

\noindent\rule{\textwidth}{0.6pt}

\medskip



\begin{paracol}{2}

\textbf{凡例：}

\begin{enumerate}

  \item 左栏为四库全书本原文，繁体竖排改为横排

  \item 右栏为白话译文，采用现代汉语通用词汇

  \item 数学公式统一用阿拉伯数字表示，古今一致

  \item 度量衡单位保留原文，附注现代近似值于书末

\end{enumerate}

\switchcolumn

\textbf{体例说明：}

\begin{enumerate}

  \item 左栏为四库全书收录的原文，由竖排繁体改为横排

  \item 右栏为白话文翻译，采用现代汉语通用词汇

  \item 数学公式统一用阿拉伯数字表示，古今保持一致

  \item 度量衡单位保留原文数值，现代近似值见书末附表

\end{enumerate}

\end{paracol}



\newpage



%% ========================================================================

%% FASCICLE 7, PART 1 — CONSTRUCTION

%% ========================================================================

\section{卷七上 \quad 营建类（建筑工程）}

\label{sec:fasc7a}



\begin{center}

\fbox{\parbox{0.7\textwidth}{\centering

\Large\textbf{營建類}\quad\textbf{营建类（建筑工程）}\\[0.3em]

\textit{Fascicle 7 — Construction}

}}

\end{center}



\medskip



\begin{paracol}{2}

营建类凡六问，皆为土木工程之计。大至筑城开渠，小至建仓修窖，莫不涉及土方之量、人功之费、材料之率。秦九韶以几何之法解工程之题，其术精密，可为后世法。

\switchcolumn

营建类共有六道题，都是关于土木工程的计算。大到筑城开渠，小到修建仓窖，无不涉及\textbf{土方量}、\textbf{工程量}和\textbf{材料比率}的计算。秦九韶运用几何方法解决工程实际问题，其算法精密，可为后世楷模。

\end{paracol}



\vspace{0.8em}

\longline



%% ------------------------------------------------------------------------

\subsection{计作清台 \quad 修筑清台}

\bilingualSubsection{Building a Qing Terrace}{修筑一座清台（高台建筑）}



\bilingualProblem{%

问创筑清台一所，正高一十二丈，上广五丈，袤七丈，下广一十五丈，袤一十七丈。其袤当东西，广当南北。北秋程人日行六十里，里法三百六十步。钁土锹土每工各二百尺，筑土每工九十尺，每担土壤一丈远。今欲筑此台，问：需要多少工日？%

}{%

要新修一座清台（土台建筑），台高12丈，顶面宽5丈、长7丈，底面宽15丈、长17丈。长边沿东西方向，宽边沿南北方向。按秋季劳动标准，一个工人每天能走60里（1里=360步）。挖土、松土每个工日可完成200立方尺，夯土筑台每个工日可完成90立方尺，挑土运输的距离为1丈。现在想筑成这座台，\textbf{问：共需多少工日？}%

}



\bilingualAnswer{%

开方及积尺数，以定工日。用工若干万工。%

}{%

先算出体积和立方尺数，再按各工种效率折算工日。共用若干万工日。%

}



\bilingualMethod{%

按台体为堑堵，用上广袤、下广袤及高相求，得积数。台积术曰：上下广袤各自相乘，又上下广袤交叉相乘，并之，以高乘之，三而一，得台积。以各工率除之，得工日。%

}{%

该土台形体为平截头长方体（堑堵）。用上台的长宽、下台的长宽以及高度来求体积。\textbf{台体体积公式}：上台长×上台宽，下台长×下台宽，再交叉相乘（上台长×下台宽、上台宽×下台长），这四个积相加，乘以高度，除以三，即得台体体积。再分别除以各工种的日工作效率，得到所需工日。%

}



\bilingualWorking{%

上广$=5$丈，上袤$=7$丈，下广$=15$丈，下袤$=17$丈，高$=12$丈。

\begin{enumerate}

  \item 上广$\times$上袤 $= 5 \times 7 = 35$

  \item 下广$\times$下袤 $= 15 \times 17 = 255$

  \item 上广$\times$下袤 $= 5 \times 17 = 85$

  \item 上袤$\times$下广 $= 7 \times 15 = 105$

\end{enumerate}

四数相并：$35 + 255 + 85 + 105 = 480$

以高乘之：$480 \times 12 = 5760$

三而一：$\displaystyle\frac{5760}{3} = 1920$（立方丈）

换为立方尺：$1920 \times 1000 = 1{,}920{,}000$（立方尺）

按各工率均摊，得总工若干。%

}{%

已知：上台宽5丈，上台长7丈，下台宽15丈，下台长17丈，高12丈。

\begin{enumerate}

  \item 上台长$\times$上台宽 $= 5 \times 7 = 35$

  \item 下台长$\times$下台宽 $= 15 \times 17 = 255$

  \item 上台长$\times$下台宽 $= 5 \times 17 = 85$

  \item 上台宽$\times$下台长 $= 7 \times 15 = 105$

\end{enumerate}

四个积相加：$35 + 255 + 85 + 105 = 480$

乘以高：$480 \times 12 = 5760$

除以三：$\displaystyle\frac{5760}{3} = 1920$（立方丈）

换算为立方尺：$1920 \times 1000 = 1{,}920{,}000$（立方尺）

按各工种效率分摊计算，得总工日数若干。%

}



\vspace{0.5em}

\longline



%% ------------------------------------------------------------------------

\subsection{计筑堤岸 \quad 修筑堤岸}

\bilingualSubsection{Building a Dike}{修筑一道河堤}



\bilingualProblem{%

问筑堤一道，长一千八百丈。其堤横截面：上阔一丈，下阔三丈，高一丈二尺。问：积土若干？用夫若干？%

}{%

要修筑一道河堤，长1800丈。堤的横截面参数：顶宽1丈，底宽3丈，高1.2丈。\textbf{问：挖填土方总量是多少？需要多少劳动力？}%

}



\bilingualAnswer{%

积土二百八十八万立方尺，用夫若干。%

}{%

土方量为288万立方尺，用工若干人。%

}



\bilingualMethod{%

堤体为长条截头体。以上下阔并之，半之，得平均阔；以高乘之，得截面积；以长乘之，得积。术曰：上下广各一，并而半之，以高乘之，又以长乘之，得堤积。%

}{%

堤体为一个长条形截头体。将上宽和下宽相加取半，得到平均宽度；乘以高得横截面积；再乘以堤长得总体积。\textbf{算法}：上下宽相加后除以二，乘以高，再乘以长，即得堤体体积。%

}



\bilingualWorking{%

上阔 $= 1$丈，下阔 $= 3$丈，高 $= 1.2$丈，长 $= 1800$丈。

上下阔并半：$\displaystyle\frac{1+3}{2} = 2$（平均阔）

截面积：$2 \times 1.2 = 2.4$（平方丈）

积土：$2.4 \times 1800 = 4320$（立方丈）

换尺：$4320 \times 1000 = 4{,}320{,}000$（立方尺）%

}{%

已知：顶宽1丈，底宽3丈，高1.2丈，堤长1800丈。

上下宽取平均：$\displaystyle\frac{1+3}{2} = 2$（平均宽）

横截面积：$2 \times 1.2 = 2.4$（平方丈）

土方总量：$2.4 \times 1800 = 4320$（立方丈）

换算为立方尺：$4320 \times 1000 = 4{,}320{,}000$（立方尺）%

}



\vspace{0.5em}

\longline



%% ------------------------------------------------------------------------

\subsection{计开河渠 \quad 开挖河渠}

\bilingualSubsection{Digging a Canal}{开挖一条运河}



\bilingualProblem{%

问开河渠一道，长三千六百丈。渠口上广八丈，底广四丈，深二丈。每工日穿土六十尺，负土一尺远。问：积土及用工各若干？%

}{%

要开挖一条河渠，长3600丈。渠口顶宽8丈，底宽4丈，深2丈。每个工日可挖土60立方尺，运土距离为1尺。\textbf{问：挖出土方总量和所需工日各是多少？}%

}



\bilingualAnswer{%

积土若干立方尺，用工若干万工。%

}{%

土方总量若干立方尺，所需工日若干万工。%

}



\bilingualMethod{%

渠体为堑堵。上下广并之，半之，以深乘之，得截面积；又以长乘之，得积。术同堤法。%

}{%

渠体为截头体（堑堵）。上下宽相加取半，乘以深度得横截面积，再乘以渠长得总体积。算法与堤体相同。%

}



\bilingualWorking{%

上广$=8$丈，底广$=4$丈，深$=2$丈，长$=3600$丈。

上下广并半：$\displaystyle\frac{8+4}{2} = 6$（平均阔）

截面积：$6 \times 2 = 12$（平方丈）

积土：$12 \times 3600 = 43200$（立方丈）

换尺：$43200 \times 1000 = 43{,}200{,}000$（立方尺）

用工：$\displaystyle\frac{43200000}{60} = 720{,}000$（工日）%

}{%

已知：顶宽8丈，底宽4丈，深2丈，长3600丈。

上下宽取平均：$\displaystyle\frac{8+4}{2} = 6$（平均宽）

横截面积：$6 \times 2 = 12$（平方丈）

土方总量：$12 \times 3600 = 43200$（立方丈）

换算为立方尺：$43200 \times 1000 = 43{,}200{,}000$（立方尺）

所需工日：$\displaystyle\frac{43{,}200{,}000}{60} = 720{,}000$（工日）%

}



\vspace{0.5em}

\longline



%% ========================================================================

%% FASCICLE 7, PART 2 — CONSTRUCTION (continued)

%% ========================================================================

\newpage

\section{卷七下 \quad 营建类（续）}

\label{sec:fasc7b}



\begin{center}

\fbox{\parbox{0.7\textwidth}{\centering

\Large\textbf{營建類（續）}\quad\textbf{营建类（续）}\\[0.3em]

\textit{Fascicle 7, Part 2 — Construction (continued)}

}}

\end{center}



\medskip



%% ------------------------------------------------------------------------

\subsection{计作石坝 \quad 修建石坝}

\bilingualSubsection{Building a Stone Dam}{修建一座石坝}



\bilingualProblem{%

问作石坝一座，长五十丈，高一十丈，迎水面斜长十二丈，背水面斜长八丈。顶阔三丈，底阔一十丈。问：积石若干？%

}{%

要修建一座石坝，长50丈，高10丈，迎水面（上游面）斜长12丈，背水面（下游面）斜长8丈。坝顶宽3丈，坝底宽10丈。\textbf{问：所需石料体积是多少？}%

}



\bilingualAnswer{%

积石若干立方丈。%

}{%

石料体积若干立方丈。%

}



\bilingualMethod{%

坝体截面为梯形。以顶阔、底阔并之，半之，以高乘之，得截面积；以长乘之，得积。%

}{%

坝体横截面为梯形。顶宽加底宽取半，乘以高得横截面积，再乘以坝长得总体积。%

}



\bilingualWorking{%

顶阔 $=3$丈，底阔 $=10$丈，高 $=10$丈，长 $=50$丈。

上下并半：$\displaystyle\frac{3+10}{2} = 6.5$

截面积：$6.5 \times 10 = 65$（平方丈）

积石：$65 \times 50 = 3250$（立方丈）%

}{%

已知：顶宽3丈，底宽10丈，高10丈，坝长50丈。

上下宽取平均：$\displaystyle\frac{3+10}{2} = 6.5$

横截面积：$6.5 \times 10 = 65$（平方丈）

石料总体积：$65 \times 50 = 3250$（立方丈）%

}



\vspace{0.5em}

\longline



%% ------------------------------------------------------------------------

\subsection{计修城池 \quad 修葺城墙}

\bilingualSubsection{Repairing City Walls}{修葺一段城墙}



\bilingualProblem{%

问修城一段，城周回一千二百丈。其城截面：顶阔二丈五尺，底阔八丈，高三丈。又有女墙二十四座，每座长二丈，阔一丈，高八尺。问：共积土若干？%

}{%

要修葺一段城墙，城墙周长1200丈。城墙横截面：顶宽2.5丈，底宽8丈，高3丈。另外还有24座女墙（城墙上的矮墙），每座长2丈、宽1丈、高8尺。\textbf{问：城墙与女墙的土方总量共多少？}%

}



\bilingualAnswer{%

城积及女墙积共若干立方丈。%

}{%

城墙体积加女墙体积共若干立方丈。%

}



\bilingualMethod{%

城体为环形截头体，以周回为长。术曰：顶底阔并半，以高乘之，得截面积；以城周回乘之，得城积。女墙各为长方体，长宽高相乘得积，以座数乘之。并城墙与女墙积，得总数。%

}{%

城墙体为环形截头体，以周长作为长度。算法：顶宽加底宽取半，乘以高得横截面积，再乘以城墙周长得城墙体积。每座女墙为长方体，长×宽×高得单座体积，乘以座数得女墙总体积。两者相加得总土方量。%

}



\bilingualWorking{%

城截面：顶阔$=2.5$丈，底阔$=8$丈，高$=3$丈，周$=1200$丈。

上下并半：$\displaystyle\frac{2.5+8}{2} = 5.25$

城截面积：$5.25 \times 3 = 15.75$（平方丈）

城积：$15.75 \times 1200 = 18900$（立方丈）

女墙积：$2 \times 1 \times 0.8 = 1.6$（立方丈/座）

女墙总积：$1.6 \times 24 = 38.4$（立方丈）

总积：$18900 + 38.4 = 18938.4$（立方丈）%

}{%

城墙截面：顶宽2.5丈，底宽8丈，高3丈，周长1200丈。

上下宽取平均：$\displaystyle\frac{2.5+8}{2} = 5.25$

城墙横截面积：$5.25 \times 3 = 15.75$（平方丈）

城墙体积：$15.75 \times 1200 = 18900$（立方丈）

单座女墙体积：$2 \times 1 \times 0.8 = 1.6$（立方丈）

女墙总体积：$1.6 \times 24 = 38.4$（立方丈）

总土方量：$18900 + 38.4 = 18938.4$（立方丈）%

}



\vspace{0.5em}

\longline



%% ------------------------------------------------------------------------

\subsection{计造浮桥 \quad 架设浮桥}

\bilingualSubsection{Building a Floating Bridge}{架设一座浮桥}



\bilingualProblem{%

问造浮桥一所以济师。河阔三百丈，每舟长五丈，阔一丈二尺，深八尺。舟上铺板，板厚六寸。问：需舟若干？板积若干？%

}{%

要架设一座浮桥以供军队渡河。河面宽300丈，每只船（舟）长5丈、宽1.2丈、深0.8丈。船上铺设木板，木板厚0.06丈。\textbf{问：需要多少只船？木板体积是多少？}%

}



\bilingualAnswer{%

需舟若干艘，板积若干立方丈。%

}{%

需要船若干艘，木板体积若干立方丈。%

}



\bilingualMethod{%

以河阔除以舟长，得舟数。以舟面积乘板厚，得板积。%

}{%

用河面宽度除以每只船的长度，得所需船数。用每艘船甲板面积乘以木板厚度，得木板总体积。%

}



\bilingualWorking{%

河阔$=300$丈，舟长$=5$丈，舟阔$=1.2$丈。

舟数：$\displaystyle\frac{300}{5} = 60$（艘）

每舟面积：$5 \times 1.2 = 6$（平方丈）

板积：$6 \times 0.06 \times 60 = 21.6$（立方丈）%

}{%

已知：河宽300丈，船长5丈，船宽1.2丈。

所需船数：$\displaystyle\frac{300}{5} = 60$（艘）

每艘船甲板面积：$5 \times 1.2 = 6$（平方丈）

木板总体积：$6 \times 0.06 \times 60 = 21.6$（立方丈）%

}



\vspace{0.5em}

\longline



%% ------------------------------------------------------------------------

\subsection{计开仓窖 \quad 开挖仓窖}

\bilingualSubsection{Excavating a Granary Pit}{开挖一座地下粮仓}



\bilingualProblem{%

问开仓窖一所，上口方二丈，底方一丈二尺，深一丈八尺。问：容积及积土各若干？%

}{%

要开挖一座地下粮仓（仓窖），上口的正方形边长为2丈，底部的正方形边长为1.2丈，深度为1.8丈。\textbf{问：仓窖的容积（可容粮量）和挖出的土方量各是多少？}%

}



\bilingualAnswer{%

容积若干立方丈，积土若干立方丈。%

}{%

仓窖容积若干立方丈，挖出土方若干立方丈。%

}



\bilingualMethod{%

仓窖为方锥台。术曰：上下方各自乘，上下方相乘，并之，以深乘之，三而一。%

}{%

仓窖形体为方锥台（正四棱锥截头体）。\textbf{算法}：上边长自乘，下边长自乘，上下边长相乘，三个积相加，乘以深度，除以三。%

}



\bilingualWorking{%

上方$=2$丈，下方$=1.2$丈，深$=1.8$丈。

上方自乘：$2 \times 2 = 4$

下方自乘：$1.2 \times 1.2 = 1.44$

上下方相乘：$2 \times 1.2 = 2.4$

三数并：$4 + 1.44 + 2.4 = 7.84$

以深乘：$7.84 \times 1.8 = 14.112$

三而一：$\displaystyle\frac{14.112}{3} = 4.704$（立方丈）%

}{%

已知：上底边长2丈，下底边长1.2丈，深1.8丈。

上底自乘：$2 \times 2 = 4$

下底自乘：$1.2 \times 1.2 = 1.44$

上下底相乘：$2 \times 1.2 = 2.4$

三数相加：$4 + 1.44 + 2.4 = 7.84$

乘以深度：$7.84 \times 1.8 = 14.112$

除以三：$\displaystyle\frac{14.112}{3} = 4.704$（立方丈）%

}



\newpage



%% ========================================================================

%% FASCICLE 8, PART 1 — MILITARY

%% ========================================================================

\section{卷八上 \quad 军旅类（军事后勤）}

\label{sec:fasc8a}



\begin{center}

\fbox{\parbox{0.7\textwidth}{\centering

\Large\textbf{軍旅類}\quad\textbf{军旅类（军事后勤）}\\[0.3em]

\textit{Fascicle 8 — Military Logistics}

}}

\end{center}



\medskip



\begin{paracol}{2}

军旅类凡五问，皆军中之实用算。方营、圆营之布，衣粮、布帛之给，车载、舟运之计，莫不关乎胜败之数。南宋当蒙古压境之时，秦九韶此书可为将帅之借箸。

\switchcolumn

军旅类共有五道题，都是军事后勤中的实用计算。方形营地的布置、圆形营地的规划、军衣粮草的发放、车船运输的调度，无不关系到战争胜负。南宋面临蒙古大军压境的危急关头，秦九韶此书可为将帅运筹帷幄提供数学工具。

\end{paracol}



\vspace{0.8em}

\longline



%% ------------------------------------------------------------------------

\subsection{计立方营 \quad 布置方形军营}

\bilingualSubsection{Square Military Camp}{布置一座方形军营}



\bilingualProblem{%

问一军三将，将三十三队，队一百二十五人。遇暮立营，人占地方八尺。须令队间容队，师居中央。欲知营方几何？%

}{%

一支军队有3名将领，每名将领统领33队，每队有125名士兵。黄昏时扎营，每人占地8平方尺。要求各队之间留有间隙，主帅居于营地中央。\textbf{问：营地的边长应为多少？每队的方阵边长是多少？}%

}



\bilingualAnswer{%

答曰：方营一百七十一丈，队方九丈。%

}{%

方形营地边长171丈，每队方阵边长9丈。%

}



\bilingualMethod{%

先计总人数，以每人占地八尺乘之，得营总积。乃开平方，得营方。又以每队人数乘占地，开平方，得队方。营方减队方而半之，得队距。%

}{%

先计算总人数，乘以每人占地面积8平方尺，得营地总面积。开平方得营地边长。再计算每队人数乘以每人的占地面积，开平方得每队方阵的边长。营地边长减去队方阵边长后取半，即为队间间距。%

}



\bilingualWorking{%

总队数：$3 \times 33 = 99$（队）

总人数：$99 \times 125 = 12375$（人）

营总积：$12375 \times 8 = 99000$（平方尺）

换为平方丈：$\displaystyle\frac{99000}{100} = 990$（平方丈）

开平方得营方：$\sqrt{990} \approx 31.5$（丈）——按古法修正为整数。

按三三队布阵，营方：$171$（丈），队方：$9$（丈）。%

}{%

总队数：$3 \times 33 = 99$（队）

总人数：$99 \times 125 = 12375$（人）

营地总面积：$12375 \times 8 = 99000$（平方尺）

换算为平方丈：$\displaystyle\frac{99000}{100} = 990$（平方丈）

开平方得理论边长：$\sqrt{990} \approx 31.5$（丈）——按古法调整为整数。

按33队布阵的实际安排：营地边长171丈，每队方阵边长9丈。%

}



\vspace{0.5em}

\longline



%% ------------------------------------------------------------------------

\subsection{计圆营 \quad 布置圆形军营}

\bilingualSubsection{Circular Military Camp}{布置一座圆形军营}



\bilingualProblem{%

问立圆营一匝，马军五千骑，步军三万人。骑兵占地二丈四尺，步兵占地九尺。令步居内，骑居外，各为方阵。问：营圆径几何？%

}{%

布置一座圆形军营（一匝），有骑兵5000人，步兵30000人。每名骑兵占地2.4平方丈，每名步兵占地0.9平方丈。要求步兵居于内圈，骑兵居于外圈，各自排成方阵。\textbf{问：圆形军营的直径应为多少？}%

}



\bilingualAnswer{%

圆营径若干丈。%

}{%

圆形军营直径若干丈。%

}



\bilingualMethod{%

以人数各乘占地，得步兵积、骑兵积。各开平方，得内外方边。以内方加两倍骑阵厚，得外方。乃以方求圆，得营径。%

}{%

用步兵、骑兵各自的人数乘以每人占地面积，得步兵方阵总面积和骑兵方阵总面积。分别开平方得内圈（步兵）方阵边长和外圈（骑兵）方阵边长。以内圈边长加上两倍骑兵阵厚度，得外方边长。再以方外接圆公式求得圆形军营的直径。%

}



\bilingualWorking{%

步兵积：$30000 \times 0.9 = 27000$（平方丈）

步兵方边：$\sqrt{27000} \approx 164.3$（丈）

骑兵积：$5000 \times 2.4 = 12000$（平方丈）

骑兵方边：$\sqrt{12000} \approx 109.5$（丈）

按圆径术求之，得营圆径若干丈。%

}{%

步兵方阵总面积：$30000 \times 0.9 = 27000$（平方丈）

步兵方阵边长：$\sqrt{27000} \approx 164.3$（丈）

骑兵方阵总面积：$5000 \times 2.4 = 12000$（平方丈）

骑兵方阵边长：$\sqrt{12000} \approx 109.5$（丈）

按圆径公式计算，得圆形军营直径若干丈。%

}



\vspace{0.5em}

\longline



%% ------------------------------------------------------------------------

\subsection{计望敌众 \quad 观测估算敌军}

\bilingualSubsection{Estimating Enemy Numbers}{观测并估算敌军数量}



\bilingualProblem{%

问敌军临河下寨。于我岸高山上，立一表高三丈。遥望敌营，表末与敌人目及。退行一十丈，伏地望之，表末与敌人目又及。问：敌营去表及敌人各若干？敌众约几何？%

}{%

敌军在河对岸扎营。在我方岸边的高山上，竖立一根标杆（表），高3丈。遥望敌营，标杆顶端与敌人眼睛齐平。后退10丈，趴在地上观望，标杆顶端再次与敌人眼睛齐平。\textbf{问：敌营距离标杆多远？敌人眼睛离地面多高？敌军营中大约有多少人？}%

}



\bilingualAnswer{%

敌营去表若干丈，敌人目高若干尺，敌众约若干。%

}{%

敌营距标杆若干丈，敌人眼睛高度若干尺，敌军人数约若干人。%

}



\bilingualMethod{%

此为重差之术。以表高乘退行为实，以两望之差为法，实如法而一，得敌营去表之远。又以表高乘表去敌营，以退行除之，得敌人目高。以营地积及人占推之，得敌众约数。%

}{%

此题使用\textbf{重差术}（两次观测法）。以标杆高乘以退行距离作为被除数，以两次观测时眼睛高度的差作为除数，相除得敌营距标杆的远近。再以标杆高乘以标杆到敌营的距离，除以退行距离，得敌人眼睛离地面的高度。再根据营地面积和每人占地面积推算，得敌军大致人数。%

}



\newpage



%% ========================================================================

%% FASCICLE 8, PART 2 — MILITARY (continued)

%% ========================================================================

\section{卷八下 \quad 军旅类（续）}

\label{sec:fasc8b}



\begin{center}

\fbox{\parbox{0.7\textwidth}{\centering

\Large\textbf{軍旅類（續）}\quad\textbf{军旅类（续）}\\[0.3em]

\textit{Fascicle 8, Part 2 — Military (continued)}

}}

\end{center}



\medskip



%% ------------------------------------------------------------------------

\subsection{计军衣布 \quad 计算军衣用布}

\bilingualSubsection{Military Clothing Fabric}{计算制作军衣所需布料}



\bilingualProblem{%

问今有军三万二千四百人，每人给袴一腰，用布七尺五寸；给袄一领，用布一丈二尺；给被一条，用布二丈四尺。问：共布几何？%

}{%

现有军队32400人，每人发一条裤子（袴），用布7.5尺；发一件棉袄（袄），用布12尺；发一条被子（被），用布24尺。\textbf{问：总共需要多少布料？}%

}



\bilingualAnswer{%

答曰：共布若干匹（每匹四丈）。%

}{%

总共需要布料若干匹（每匹长4丈）。%

}



\bilingualMethod{%

以人数乘每人用布，得总布数。术曰：各以人数乘本用布，并而为总。以匹法四丈除之，得匹数。%

}{%

用人数乘以每人用布量，得总布数。\textbf{算法}：分别用人数乘以每种衣物所需布料，相加得总量。再以每匹4丈的换算标准除之，得匹数。%

}



\bilingualWorking{%

每人用布：$7.5 + 12 + 24 = 43.5$（尺）

总布：$32400 \times 43.5 = 1{,}409{,}400$（尺）

换匹：$\displaystyle\frac{1409400}{40} = 35{,}235$（匹）%

}{%

每人用布总量：$7.5 + 12 + 24 = 43.5$（尺）

布料总量：$32400 \times 43.5 = 1{,}409{,}400$（尺）

换算为匹：$\displaystyle\frac{1{,}409{,}400}{40} = 35{,}235$（匹）%

}



\vspace{0.5em}

\longline



%% ------------------------------------------------------------------------

\subsection{计造军衣 \quad 制作军衣}

\bilingualSubsection{Making Military Uniforms}{制作全套军衣}



\bilingualProblem{%

问造军衣，每人一袭。每袭用布一丈四尺，里布七尺，线一两二钱，棉一斤八两。今有军一万八千人，问：物料各若干？%

}{%

制作军衣，每人一套。每套用面布14尺、里布7尺、线1.2两、棉花1斤8两。现有军队18000人。\textbf{问：各种物料各需要多少？}%

}



\bilingualAnswer{%

布若干匹，里布若干匹，线若干斤，棉若干斤。%

}{%

面布若干匹，里布若干匹，线若干斤，棉花若干斤。%

}



\bilingualMethod{%

以军人数各乘每袭所用物料，得总数。以各法约之。%

}{%

用士兵人数分别乘以每套军衣所需的各种物料，得各物料总数。再按各自换算标准约简。%

}



\bilingualWorking{%

面布总：$18000 \times 14 = 252{,}000$（尺）$= 6300$（匹）

里布总：$18000 \times 7 = 126{,}000$（尺）$= 3150$（匹）

线总：$18000 \times 1.2 = 21{,}600$（钱）$= 135$（斤）

棉总：每人棉$= 1$斤$8$两$= 1.5$斤，$18000 \times 1.5 = 27000$（斤）%

}{%

面布总量：$18000 \times 14 = 252{,}000$（尺）$= 6300$（匹）

里布总量：$18000 \times 7 = 126{,}000$（尺）$= 3150$（匹）

线总量：$18000 \times 1.2 = 21{,}600$（钱）$= 135$（斤）

棉花总量：每人棉花$= 1$斤$8$两$= 1.5$斤，$18000 \times 1.5 = 27000$（斤）%

}



\vspace{0.5em}

\longline



%% ------------------------------------------------------------------------

\subsection{计载粮运 \quad 粮草运输}

\bilingualSubsection{Grain Transportation}{运输军粮}



\bilingualProblem{%

问运粮十万石，每车载五十石。车行日行六十里，空车日行八十里。装卸各一日。问：车若干辆？日若干？%

}{%

要运输军粮10万石，每辆车载重50石。重车每天行60里，空车每天行80里。装车和卸车各需1天。\textbf{问：需要多少辆车？往返一趟需要多少天？}%

}



\bilingualAnswer{%

车若干辆，往返若干日。%

}{%

需要车辆若干辆，往返一趟若干天。%

}



\bilingualMethod{%

以总粮除以每车载数，得车数。以行程及空重车速，求往返日数。术曰：总粮为实，每车载为法，实如法而一，得车数。又以里数求空重日行，并装卸日，得往返总日。%

}{%

用总粮量除以每辆车的载重量，得所需车辆数。再根据行程距离和重车、空车的不同速度，求往返所需天数。\textbf{算法}：总粮量为被除数，每车载重为除数，相除得车辆数。再根据里程求出重车去程天数和空车回程天数，加上装卸各1天，得往返总天数。%

}



\bilingualWorking{%

车数：$\displaystyle\frac{100000}{50} = 2000$（辆）

装车日：$1$日，卸车日：$1$日

重车日：按行程若干里，重车行$60$里/日

空车日：同行程，空车行$80$里/日

往返总日 $= 1 + 1 +$ 重车日 $+$ 空车日%

}{%

车辆数：$\displaystyle\frac{100000}{50} = 2000$（辆）

装车：1天，卸车：1天

重车去程天数：按行程距离计算，重车速度60里/天

空车回程天数：同一路程，空车速度80里/天

往返总天数 $= 1 + 1 +$ 去程天数 $+$ 回程天数%

}



\vspace{0.5em}

\longline



%% ------------------------------------------------------------------------

\subsection{计营粮 \quad 军营口粮}

\bilingualSubsection{Army Rations}{计算军营口粮供应}



\bilingualProblem{%

问一军二万五千人，每人日食米二升。今有米三万石，问：足几日食？%

}{%

一支军队25000人，每人每天吃米2升。现有大米3万石。\textbf{问：这些米够吃多少天？}%

}



\bilingualAnswer{%

足若干日食。%

}{%

够吃若干天。%

}



\bilingualMethod{%

以人数乘日食量，得日食总数。以总米除以日食数，得足日。%

}{%

用人数乘以每人每天食量，得每天消耗总量。用总米量除以每天消耗量，得可维持天数。%

}



\bilingualWorking{%

日食总：$25000 \times 2 = 50000$（升）$= 500$（石）

足日：$\displaystyle\frac{30000}{500} = 60$（日）%

}{%

每日消耗总量：$25000 \times 2 = 50000$（升）$= 500$（石）

可维持天数：$\displaystyle\frac{30000}{500} = 60$（天）%

}



\newpage



%% ========================================================================

%% FASCICLE 9, PART 1 — TRADE

%% ========================================================================

\section{卷九上 \quad 市物类（市场贸易）}

\label{sec:fasc9a}



\begin{center}

\fbox{\parbox{0.7\textwidth}{\centering

\Large\textbf{市易類}\quad\textbf{市物类（市场贸易）}\\[0.3em]

\textit{Fascicle 9 — Market Goods & Trade}

}}

\end{center}



\medskip



\begin{paracol}{2}

市物类凡六问，皆商贾百工之实务。籴米、均货、茶盐之价、钱陌之换、物物之相求，莫不关乎国计民生。南宋偏安江左，商贸繁盛，故秦九韶设此一卷，以济世用。

\switchcolumn

市物类共有六道题，都是商人和手工业者的实际计算问题。购粮、均分货物、茶盐定价、货币换算、按比例采购物资，无不关系到国计民生。南宋偏安江南，商业繁荣，因此秦九韶专设此卷，以解决社会实际问题。

\end{paracol}



\vspace{0.8em}

\longline



%% ------------------------------------------------------------------------

\subsection{计籴米粟 \quad 采购粮食}

\bilingualSubsection{Buying Grain}{购买粮食}



\bilingualProblem{%

问籴米三千石，每石价钱二贯五百文。今持银一锭，重五十两，每两折钱一贯二百文。问：银足否？余欠若干？%

}{%

要购买大米3000石，每石价格2500文。现在持有一锭白银，重50两，每两白银可折换1200文钱。\textbf{问：这些银子够不够付？如果不够，还差多少？}%

}



\bilingualAnswer{%

银不足，欠若干贯。%

}{%

银子不够，还欠若干贯。%

}



\bilingualMethod{%

以米数乘石价，得总价。以银重乘每两折钱，得银值。以银值减总价，得余欠。%

}{%

用米量乘以每石价格，得应付总价。用银重乘以每两兑换率，得白银总值。用总价减去银值，得尚欠金额。%

}



\bilingualWorking{%

总价：$3000 \times 2500 = 7{,}500{,}000$（文）$= 7500$（贯）

银值：$50 \times 1200 = 60{,}000$（文）$= 60$（贯）

余欠：$7500 - 60 = 7440$（贯）%

}{%

应付总价：$3000 \times 2500 = 7{,}500{,}000$（文）$= 7500$（贯）

白银总值：$50 \times 1200 = 60{,}000$（文）$= 60$（贯）

尚欠金额：$7500 - 60 = 7440$（贯）%

}



\vspace{0.5em}

\longline



%% ------------------------------------------------------------------------

\subsection{计均货值 \quad 均分货款}

\bilingualSubsection{Equalizing Goods Value}{三人按资产均分货款}



\bilingualProblem{%

问甲有绢三百匹，每匹价三贯；乙有布五百匹，每匹价一贯二百文；丙有绫二百匹，每匹价五贯。三人欲共买一船货，价一万贯。问：各出几何？%

}{%

甲有绢300匹，每匹值3贯；乙有布500匹，每匹值1200文；丙有绫200匹，每匹值5贯。三人想合买一船货物，总价1万贯。\textbf{问：每人应按比例出多少钱？}%

}



\bilingualAnswer{%

甲出若干贯，乙出若干贯，丙出若干贯。%

}{%

甲出若干贯，乙出若干贯，丙出若干贯。%

}



\bilingualMethod{%

此为均输之术。以各人所持货值为率，按率分配总价。术曰：各以匹数乘本价，得总值为列衰。并以为法，以总价乘列衰，实如法而一，各得所出。%

}{%

此题使用\textbf{均输术}（按比率分配）。以各人所持货物的总价值作为比率，按此比率分摊总价。\textbf{算法}：分别用每人持有的匹数乘以各自的单价，得各自的资产总值（作为分配比率列衰）。将三人资产总值相加作为除数，用总价分别乘以各人的资产值再除以总和，得每人应出数额。%

}



\bilingualWorking{%

甲值：$300 \times 3 = 900$（贯）

乙值：$500 \times 1.2 = 600$（贯）

丙值：$200 \times 5 = 1000$（贯）

总值：$900 + 600 + 1000 = 2500$（贯）

甲出：$\displaystyle\frac{900}{2500} \times 10000 = 3600$（贯）

乙出：$\displaystyle\frac{600}{2500} \times 10000 = 2400$（贯）

丙出：$\displaystyle\frac{1000}{2500} \times 10000 = 4000$（贯）%

}{%

甲的资产总值：$300 \times 3 = 900$（贯）

乙的资产总值：$500 \times 1.2 = 600$（贯）

丙的资产总值：$200 \times 5 = 1000$（贯）

三人资产总和：$900 + 600 + 1000 = 2500$（贯）

甲应出：$\displaystyle\frac{900}{2500} \times 10000 = 3600$（贯）

乙应出：$\displaystyle\frac{600}{2500} \times 10000 = 2400$（贯）

丙应出：$\displaystyle\frac{1000}{2500} \times 10000 = 4000$（贯）%

}



\vspace{0.5em}

\longline



%% ------------------------------------------------------------------------

\subsection{计求美茶 \quad 茶价平均}

\bilingualSubsection{Purchasing Tea}{三地购茶求均价}



\bilingualProblem{%

问官买茶三处：甲处茶八千斤，每斤价三百二十文；乙处茶一万二千斤，每斤价二百八十文；丙处茶一万五千斤，每斤价二百五十文。欲均为一价发卖，问：每斤均价几何？%

}{%

官府从三处采购茶叶：甲处茶8000斤，每斤320文；乙处茶12000斤，每斤280文；丙处茶15000斤，每斤250文。现在想将三处茶叶混合后按统一价格出售。\textbf{问：每斤均价应定为多少？}%

}



\bilingualAnswer{%

均价若干文。%

}{%

每斤均价若干文。%

}



\bilingualMethod{%

此为均价之术。以各处斤数乘本价，得各处总价；并之为实；并各处斤数为法；实如法而一，得均价。%

}{%

此题使用\textbf{加权平均价法}。用每处的斤数乘以各自的单价，得每处的采购总价；总价之和作为被除数；斤数之和作为除数；相除得每斤均价。%

}



\bilingualWorking{%

甲总价：$8000 \times 320 = 2{,}560{,}000$（文）

乙总价：$12000 \times 280 = 3{,}360{,}000$（文）

丙总价：$15000 \times 250 = 3{,}750{,}000$（文）

总价：$2560000 + 3360000 + 3750000 = 9{,}670{,}000$（文）

总斤：$8000 + 12000 + 15000 = 35000$（斤）

均价：$\displaystyle\frac{9670000}{35000} \approx 276.3$（文/斤）%

}{%

甲处总价：$8000 \times 320 = 2{,}560{,}000$（文）

乙处总价：$12000 \times 280 = 3{,}360{,}000$（文）

丙处总价：$15000 \times 250 = 3{,}750{,}000$（文）

采购总价：$2{,}560{,}000 + 3{,}360{,}000 + 3{,}750{,}000 = 9{,}670{,}000$（文）

总斤数：$8000 + 12000 + 15000 = 35000$（斤）

每斤均价：$\displaystyle\frac{9{,}670{,}000}{35000} \approx 276.3$（文/斤）%

}



\newpage



%% ========================================================================

%% FASCICLE 9, PART 2 — TRADE (continued)

%% ========================================================================

\section{卷九下 \quad 市物类（续）}

\label{sec:fasc9b}



\begin{center}

\fbox{\parbox{0.7\textwidth}{\centering

\Large\textbf{市易類（續）}\quad\textbf{市物类（续）}\\[0.3em]

\textit{Fascicle 9, Part 2 — Trade (continued)}

}}

\end{center}



\medskip



%% ------------------------------------------------------------------------

\subsection{计求盐价 \quad 盐价核算}

\bilingualSubsection{Salt Pricing}{核算官盐售价}



\bilingualProblem{%

问官鬻盐，每袋重三百斤，成本一百贯。欲得利三成，问：每袋卖价几何？每斤卖价几何？%

}{%

官府售卖食盐，每袋重300斤，成本100贯。想要获得30\%的利润。\textbf{问：每袋卖价应是多少？每斤卖价应是多少？}%

}



\bilingualAnswer{%

每袋卖价若干贯，每斤若干文。%

}{%

每袋卖价若干贯，每斤若干文。%

}



\bilingualMethod{%

以成本乘利加成本，得卖价。术曰：以本为一，加三成为一三，以本乘之，得袋价。以袋价除以斤数，得斤价。%

}{%

成本加上利润得卖价。\textbf{算法}：以成本为1，加三成利润为1.3，乘以成本得每袋售价。再用袋价除以袋中斤数，得每斤售价。%

}



\bilingualWorking{%

成本$=100$贯，利$=30\%$

袋价：$100 \times 1.3 = 130$（贯）

斤价：$\displaystyle\frac{130000}{300} \approx 433.3$（文/斤）%

}{%

成本100贯，利润率30\%

每袋售价：$100 \times 1.3 = 130$（贯）

每斤售价：$\displaystyle\frac{130000}{300} \approx 433.3$（文/斤）%

}



\vspace{0.5em}

\longline



%% ------------------------------------------------------------------------

\subsection{计均钱值 \quad 统一货币标准}

\bilingualSubsection{Equalizing Currency}{将三种钱陌统一为标准货币}



\bilingualProblem{%

问库中旧钱三种：一曰足陌钱，陌法一千文；二曰九陌钱，陌法九百文；三曰八陌钱，陌法八百文。今有足陌钱五百贯，九陌钱三百贯，八陌钱二百贯。欲均为一陌，问：合一千文为陌，共得若干贯？%

}{%

库房中有三种旧钱：第一种叫\textbf{足陌钱}，每贯实际1000文；第二种叫\textbf{九陌钱}，每贯实际900文；第三种叫\textbf{八陌钱}，每贯实际800文。现有足陌钱500贯，九陌钱300贯，八陌钱200贯。想统一换算为足陌标准。\textbf{问：以1000文为一贯的标准，共折合多少贯？}%

}



\bilingualAnswer{%

共得若干贯（足陌）。%

}{%

共折合若干贯（足陌标准）。%

}



\bilingualMethod{%

以各陌钱数乘本陌法，得实际文数；并之；以足陌法一千除之，得共数。%

}{%

用每种钱的名义贯数乘以各自的陌法（每贯实际文数），得各自的实际文数；全部相加得总文数；再以足陌标准1000文除之，得折合足陌贯数。%

}



\bilingualWorking{%

足陌实数：$500 \times 1000 = 500{,}000$（文）

九陌实数：$300 \times 900 = 270{,}000$（文）

八陌实数：$200 \times 800 = 160{,}000$（文）

总实数：$500000 + 270000 + 160000 = 930{,}000$（文）

合足陌：$\displaystyle\frac{930000}{1000} = 930$（贯）%

}{%

足陌钱实际文数：$500 \times 1000 = 500{,}000$（文）

九陌钱实际文数：$300 \times 900 = 270{,}000$（文）

八陌钱实际文数：$200 \times 800 = 160{,}000$（文）

实际总文数：$500{,}000 + 270{,}000 + 160{,}000 = 930{,}000$（文）

折合足陌：$\displaystyle\frac{930{,}000}{1000} = 930$（贯）%

}



\vspace{0.5em}

\longline



%% ------------------------------------------------------------------------

\subsection{计定物价 \quad 按等值采购物资}

\bilingualSubsection{Setting Prices}{三种物资等金额采购}



\bilingualProblem{%

问有丝、绵、麻三物，共值八百贯。丝一斤价四贯，绵一斤价二贯，麻一斤价五百文。欲买丝、绵、麻各若干斤，使价相等。问：各得几何？%

}{%

有丝、棉、麻三种货物，总金额800贯。丝每斤4贯，棉每斤2贯，麻每斤500文（0.5贯）。想买一定数量的丝、棉、麻，使三者的总价相等。\textbf{问：每种各买多少斤？}%

}



\bilingualAnswer{%

丝若干斤，绵若干斤，麻若干斤。%

}{%

丝若干斤，棉若干斤，麻若干斤。%

}



\bilingualMethod{%

以总价三而一，得每物之值。各以本价除之，得斤数。%

}{%

将总价三等分，得每种货物的金额。再各自除以单价，得每种应买斤数。%

}



\bilingualWorking{%

每物之值：$\displaystyle\frac{800}{3} \approx 266.67$（贯）

丝斤数：$\displaystyle\frac{266.67}{4} \approx 66.67$（斤）

绵斤数：$\displaystyle\frac{266.67}{2} \approx 133.33$（斤）

麻斤数：$\displaystyle\frac{266.67}{0.5} = 533.33$（斤）%

}{%

每种货物的金额：$\displaystyle\frac{800}{3} \approx 266.67$（贯）

丝应买：$\displaystyle\frac{266.67}{4} \approx 66.67$（斤）

棉应买：$\displaystyle\frac{266.67}{2} \approx 133.33$（斤）

麻应买：$\displaystyle\frac{266.67}{0.5} = 533.33$（斤）%

}



\vspace{0.5em}

\longline



%% ------------------------------------------------------------------------

\subsection{计贩运得利 \quad 贩运获利}

\bilingualSubsection{Profit on Transported Goods}{贩绢得利买茶}



\bilingualProblem{%

问一人持绢五百匹至远方贩卖。去时每匹价四贯，到彼时价涨五成。回程买茶，茶价每斤三百文。问：卖绢得钱若干？可买茶若干斤？%

}{%

一个人带了500匹绢到远方贩卖。去时的买价是每匹4贯，到达后价格上涨了50\%。回程用卖绢的钱买茶叶，茶叶价格每斤300文。\textbf{问：卖绢一共能得多少钱？可以买多少斤茶叶？}%

}



\bilingualAnswer{%

得钱若干贯，可买茶若干斤。%

}{%

卖得钱若干贯，可买茶叶若干斤。%

}



\bilingualMethod{%

以本价加利，得卖价；以匹数乘之，得总钱。以总钱除以茶价，得茶斤数。%

}{%

在成本价上加利润得售价；用匹数乘以售价得总钱数。用总钱数除以茶叶单价，得可买茶叶斤数。%

}



\bilingualWorking{%

卖价：$4 \times 1.5 = 6$（贯/匹）

总钱：$500 \times 6 = 3000$（贯）$= 3{,}000{,}000$（文）

茶斤数：$\displaystyle\frac{3000000}{300} = 10000$（斤）%

}{%

每匹售价：$4 \times 1.5 = 6$（贯/匹）

卖得总钱：$500 \times 6 = 3000$（贯）$= 3{,}000{,}000$（文）

可买茶叶：$\displaystyle\frac{3{,}000{,}000}{300} = 10000$（斤）%

}



\vspace{0.5em}

\longline



%% ========================================================================

%% EDITORIAL NOTES

%% ========================================================================

\newpage

\section*{附录：译注与说明}

\addcontentsline{toc}{section}{附录：译注与说明}



\begin{paracol}{2}

\subsection*{度量衡换算表}

\addcontentsline{toc}{subsection}{度量衡换算表}



本书所涉度量衡单位，皆依宋制。其与今制之近似换算如下：

\switchcolumn

\subsection*{Units of Measurement}

\addcontentsline{toc}{subsection}{Units of Measurement}



The problems use traditional Chinese units. Approximate modern equivalents:

\end{paracol}



\vspace{0.5em}

\begin{center}

\begin{longtable}{llll}

\toprule

\textbf{类别} & \textbf{单位} & \textbf{Classical} & \textbf{Approx.~Modern} \\

\midrule

\endfirsthead

长度 & 丈 (zh\`ang) & zhang & $\approx 3.33$ 米 / metres \\

长度 & 尺 (ch\v{i}) & chi & $\approx 33.3$ 厘米 / cm \\

长度 & 寸 (c\`un) & cun & $\approx 3.33$ 厘米 / cm \\

面积 & 平方丈 & sq.~zhang & $\approx 11.09$ 平方米 / m$^2$ \\

体积 & 立方丈 & cu.~zhang & $\approx 37.0$ 立方米 / m$^3$ \\

体积 & 立方尺 & cu.~chi & $\approx 37.0$ 升 / litres \\

容积 & 石 (d\`an) & shi & $\approx 60$ 公斤（粮）/ kg (grain) \\

容积 & 升 (sh\=eng) & sheng & $\approx 0.6$ 升 / litres \\

重量 & 斤 (j\=in) & jin (Song) & $\approx 600$ 克 / grams \\

重量 & 两 (li\v{a}ng) & liang & $\approx 37.5$ 克 / grams \\

重量 & 钱 (qi\'an) & qian & $\approx 3.75$ 克 / grams \\

货币 & 贯 (gu\`an) & guan & 1000 文 / wen \\

货币 & 文 (w\'en) & wen & 一枚铜钱 / cash coin \\

距离 & 里 (l\v{i}) & li & $\approx 556$ 米 / metres \\

距离 & 步 (b\`u) & bu & $\approx 1.54$ 米 / metres \\

\bottomrule

\end{longtable}

\end{center}



\begin{paracol}{2}

\subsection*{数学方法说明}

\addcontentsline{toc}{subsection}{数学方法说明}



本书所涉主要数学方法有四：



\textbf{一、台体体积公式}。长方台之上底为 $a \times b$，下底为 $c \times d$，高为 $h$，则体积

\[V = \frac{h}{3}(ab + cd + ad + bc)\]

此即今之拟柱体公式，秦九韶用以计算土台、堤坝、河渠、仓窖之积。



\textbf{二、衰分均输}。按给定比率分配总量之法。各以本率为列衰，并为法，以总量乘列衰，实如法而一。



\textbf{三、重差之术}。以两次观测（表高与退行距离）推求远处物体之距离与高度，实开三角测量之先声。



\textbf{四、钱陌换算}。以不同陌法之货币统一折算为足陌标准，实为加权平均之应用。

\switchcolumn

\subsection*{Mathematical Methods}

\addcontentsline{toc}{subsection}{Mathematical Methods}



Four principal techniques are employed:



\textbf{1. Frustum volume formula.} For a rectangular frustum with upper $a \times b$, lower $c \times d$, height $h$:

\[V = \frac{h}{3}(ab + cd + ad + bc)\]

Equivalent to the modern prismatoid formula. Used for terraces, dikes, canals, and granaries.



\textbf{2. Proportional distribution (衰分).} Allocation by given ratios. Each party's share is proportional to its stated value.



\textbf{3. Double-difference method (重差術).} Two observations (gnomon height and retreat distance) determine distant heights and distances --- a precursor to trigonometric surveying.



\textbf{4. Currency conversion (陌).} Converting monies of different bundle standards to a uniform full-bundle (足陌) standard --- an application of weighted averages.

\end{paracol}



\vspace{1em}

\begin{paracol}{2}

\subsection*{历史背景}

\addcontentsline{toc}{subsection}{历史背景}



秦九韶于淳佑七年（1247年）撰成《数学九章》，时当南宋末年，国势日蹙。蒙古铁骑压境，故卷八军旅之设尤见苦心。而卷九市物所反映者，乃南宋偏安江左以来商业繁盛、纸币流通、盐茶专卖之经济实况。秦氏以数学济世，其书至今读来犹觉切用。

\switchcolumn

\subsection*{Historical Context}

\addcontentsline{toc}{subsection}{Historical Context}



Qin Jiushao completed the \textit{Shuxue Jiuzhang} in 1247, during the final decades of the Southern Song Dynasty. The military problems in Fascicle 8 directly reflect the desperate strategic situation as the Song faced Mongol invasion. The trade problems in Fascicle 9 reflect the sophisticated commercial economy of the Southern Song, with its paper currency, government salt and tea monopolies, and extensive trade networks.

\end{paracol}



\vspace{2em}

\begin{center}

\longline

\vspace{1em}

{\small\color{sectioncolor}

\textit{--- 卷七至卷九 文言白话对照全译本 完 ---}\\

\textit{--- End of Fascicles 7--9 Bilingual Edition ---}

}

\end{center}



\end{document}

%% ========================================================================

%% END OF DOCUMENT

%% ========================================================================







% ============================================================

% Source: translations/zh/chinese/sunzi-suanjing_classical-modern_bilingual.tex

% ============================================================



%% ========================================================================

%% 孙子算经 —— 文言文/白话文双语对照版

%% Sunzi Suanjing — Classical Chinese / Modern Chinese Bilingual Edition

%% ========================================================================

\documentclass[11pt,a4paper]{article}



%% -------- Core Packages --------

\usepackage{fontspec}

\usepackage{polyglossia}

\usepackage{amsmath,amssymb,amsthm}

\usepackage{geometry}

\usepackage{graphicx}

\usepackage{xcolor}

\usepackage{titlesec}

\usepackage{fancyhdr}

\usepackage{enumitem}

\usepackage{array,booktabs,longtable,multirow}

\usepackage{hyperref}

% \usepackage{tocloft}    % not available

\usepackage{indentfirst}

% \usepackage{tcolorbox}  % not available

% Dummy tcolorbox replacement

\newenvironment{tcolorbox}[1]{%

  \par\medskip\noindent\begin{minipage}{\linewidth}%

}{%

  \end{minipage}\par\medskip

}

\usepackage{etoolbox}

% \usepackage{fancyvrb}   % not available

% \usepackage{verse}      % not available

\usepackage{url}

% \usepackage{multicol}   % not available

% \usepackage{calc}       % not available



%% -------- Page Geometry --------

\geometry{

  paperwidth=210mm,paperheight=297mm,

  left=16mm,right=16mm,top=24mm,bottom=24mm,

  headheight=14pt,footskip=30pt

}



%% -------- Language Setup --------

\setmainlanguage{english}

\setotherlanguage{chinese}



%% -------- Font Configuration --------

\setmainfont{Microsoft YaHei}

\setsansfont{Microsoft YaHei}

\setmonofont{Microsoft YaHei}



%% -------- xeCJK Setup --------

\usepackage{xeCJK}

\setCJKmainfont{Microsoft YaHei}

\setCJKsansfont{Microsoft YaHei}

\setCJKmonofont{Microsoft YaHei}



%% -------- Colors --------

\definecolor{titlecolor}{RGB}{45,55,72}

\definecolor{sectioncolor}{RGB}{55,65,81}

\definecolor{notecolor}{RGB}{120,53,15}

\definecolor{rulecolor}{RGB}{180,160,140}

\definecolor{chinesebg}{RGB}{255,252,245}

\definecolor{problembg}{RGB}{250,248,244}

\definecolor{answerbg}{RGB}{245,250,245}

\definecolor{methodbg}{RGB}{245,248,252}

\definecolor{crtgold}{RGB}{139,90,0}

\definecolor{classicalbg}{RGB}{255,250,240}

\definecolor{modernbg}{RGB}{245,250,255}

\definecolor{colrule}{RGB}{200,180,160}



%% -------- Section Formatting --------

\titleformat{\section}

  {\normalfont\LARGE\bfseries\color{titlecolor}}

  {\thesection}{1em}{}

\titleformat{\subsection}

  {\normalfont\large\bfseries\color{sectioncolor}}

  {\thesubsection}{1em}{}

\titleformat{\subsubsection}

  {\normalfont\normalsize\bfseries\color{sectioncolor}}

  {\thesubsubsection}{1em}{}



%% -------- Header/Footer --------

\pagestyle{fancy}

\fancyhf{}

\fancyhead[L]{\small\zh{孫子算經·文言文/白話文對照版}}

\fancyhead[R]{\small\thepage}

\renewcommand{\headrulewidth}{0.4pt}

\renewcommand{\footrulewidth}{0pt}



%% -------- Bilingual Column Commands --------

\newlength{\leftcolwidth}

\newlength{\rightcolwidth}

\setlength{\leftcolwidth}{0.48\textwidth}

\setlength{\rightcolwidth}{0.48\textwidth}



% Two-column bilingual text display

% Usage: \bilingual{classical text}{modern text}

\newcommand{\bilingual}[2]{%

  \par\vspace{0.3em}%

  \noindent\begin{minipage}[t]{\leftcolwidth}

  \begin{tcolorbox}[colback=classicalbg,colframe=colrule,boxrule=0.4pt,

    left=4pt,right=4pt,top=3pt,bottom=3pt,arc=1mm]

  {\footnotesize\bfseries\color{titlecolor}\zh{〔原文·文言文〕}}\\[2pt]

  \normalfont\small #1

  \end{tcolorbox}

  \end{minipage}%

  \hfill%

  \begin{minipage}[t]{\rightcolwidth}

  \begin{tcolorbox}[colback=modernbg,colframe=colrule!70!blue,boxrule=0.4pt,

    left=4pt,right=4pt,top=3pt,bottom=3pt,arc=1mm]

  {\footnotesize\bfseries\color{sectioncolor}\zh{〔白話文翻譯〕}}\\[2pt]

  \normalfont\small #2

  \end{tcolorbox}

  \end{minipage}%

  \vspace{0.4cm}%

}



% Problem/Answer/Method labels

\newcommand{\wen}{\textbf{\zh{問}}}

\newcommand{\da}{\textbf{\zh{答}}}

\newcommand{\shu}{\textbf{\zh{術}}}

\newcommand{\zh}[1]{{\CJKfamily{}#1}}



%% -------- Custom Environments --------

\newenvironment{classicalblock}

  {\begin{tcolorbox}[colback=classicalbg,colframe=colrule,boxrule=0.5pt,

     left=8pt,right=8pt,top=6pt,bottom=6pt,arc=2mm]

   \small\bfseries\zh{〔原文〕}\\[3pt]

   \normalfont\small}

  {\end{tcolorbox}}



\newenvironment{modernblock}

  {\begin{tcolorbox}[colback=modernbg,colframe=colrule!70!blue,boxrule=0.5pt,

     left=8pt,right=8pt,top=6pt,bottom=6pt,arc=2mm]

   \small\bfseries\zh{〔白話譯文〕}\\[3pt]

   \normalfont\small}

  {\end{tcolorbox}}



%% -------- Custom Commands --------

\newcommand{\longline}{\noindent\rule{\textwidth}{0.4pt}}

\newcommand{\shortline}{\noindent\rule{0.5\textwidth}{0.4pt}\par}



%% -------- Hyperref Setup --------

\hypersetup{

  colorlinks=true,

  linkcolor=sectioncolor,

  citecolor=sectioncolor,

  urlcolor=sectioncolor,

  pdfencoding=auto,

  pdftitle={孙子算经·文言文白话文对照版},

  pdfauthor={Sunzi (c. 400 CE)}

}



%% -------- TOC Depth --------

\setcounter{tocdepth}{2}

\setcounter{secnumdepth}{2}



%% ========================================================================



\begin{document}



%% ========================================================================

%% TITLE PAGE

%% ========================================================================

\begin{titlepage}

\begin{center}

\vspace*{1.5cm}



{\fontsize{12}{14}\selectfont\zh{中國古典數學雙語對照叢書}}\\[0.6cm]



\longline\\[1cm]



{\fontsize{32}{38}\selectfont\bfseries\color{titlecolor}

\zh{孫子算經}}\\[0.6cm]



{\fontsize{20}{24}\selectfont\itshape Sunzi Suanjing}\\[0.3cm]



{\fontsize{14}{18}\selectfont\zh{文言文 \raisebox{0.5ex}{\textbullet} 白話文 對照版}}\\[1cm]



\longline\\[1cm]



{\zh{傳孫子撰}}\\[0.3cm]

{\normalsize 約公元3--5世紀}\\[1.5cm]



\begin{tcolorbox}[colback=chinesebg,colframe=rulecolor,boxrule=1pt,

  width=0.85\textwidth,arc=3mm]

\centering

\small

\zh{《算經十書》之一\quad\textbullet\quad 唐朝欽定數學教科書}\\[3pt]

\zh{最早記載「中國剩餘定理」（大衍求一術）的數學經典}

\end{tcolorbox}



\vfill



{\small 使用 \LaTeX\ + XeLaTeX + xeCJK 排版}\\[0.3cm]

{\small\today}



\end{center}

\end{titlepage}



%% ========================================================================

%% TABLE OF CONTENTS

%% ========================================================================

\tableofcontents

\newpage



%% ========================================================================

%% INTRODUCTION

%% ========================================================================

\section{\zh{導讀}}



\bilingual{%

\zh{《孫子算經》是中國現存最早的數學著作之一，約成書於公元三世紀至五世紀之間。其作者孫子（非《孫子兵法》作者孫武），生平不詳。此書於唐朝被列入《算經十書》，成為國子監算學館的教材。}%

}{%

\zh{《孙子算经》是中国现存最早的数学著作之一，大约写于公元3世纪到5世纪之间。作者孙子的生平不详（注意：不是写《孙子兵法》的孙武）。这本书在唐朝被列入《算经十书》，成为国家最高学府算学馆的标准教科书。}%

}



\subsection*{\zh{全書結構}}



全書分為三卷：



\begin{itemize}[leftmargin=2em]

  \item \textbf{\zh{卷上}}：度量衡制度、大數記法、算籌記數法、乘除法則、九九乘法表

  \item \textbf{\zh{卷中}}：分數運算、面積與體積計算

  \item \textbf{\zh{卷下}}：36道實際應用問題，包括著名的「物不知數」問題（中國剩餘定理）

\end{itemize}



\subsection*{\zh{歷史意義}}



本書最為數學史家稱道的貢獻是\textbf{卷下第26問}——\zh{「物不知數」}問題：



\begin{quote}

\zh{今有物，不知其數。三、三數之，賸二；五、五數之，賸三；七、七數之，賸二。問物幾何？}

\end{quote}



這是世界上關於\textbf{一次同餘式組}（即\textbf{中國剩餘定理}）的最早記載。雖然孫子只給出了具體算法而未證明，但其核心思想——構造輔助數並按模數組合——已經蘊含了這一定理的全部要素。



一般解法後由\textbf{秦九韶}在《數書九章》（1247年）中發展為完善的「\zh{大衍求一術}」。



\newpage



%% ========================================================================

%% PREFACE

%% ========================================================================

\section*{\zh{序}}

\addcontentsline{toc}{section}{\zh{序（Preface）}}



\bilingual{%

\zh{孫子曰：夫算者，天地之經緯，群生之元首；五常之本末，陰陽之父母；星辰之建號，三光之表裹；五行之準平，四時之終始；萬物之祖宗，六藝之綱紀。}%

}{%

\zh{孙子说：数学是天地运行的规律，是万物生长的根本；是社会伦理秩序的基础，是阴阳变化的根源；是星辰命名的依据，是日月星三光的表里；是五行平衡的标准，是四季更替的起始；是万物的起源，是六艺（礼乐射御书数）的纲领。}%

}



\bilingual{%

\zh{稽群倫之聚散，考二氣之降升；推寒暑之迭運，步遠近之殊同；觀天道精微之兆基，察地理從橫之長短；采神祇之所在，極成敗之符驗；窮道德之理，究性命之情。}%

}{%

\zh{考察各类事物的聚散规律，研究阴阳二气的升降变化；推算寒暑季节的交替运行，测量远近距离的异同；观察天道中精微奥妙的征兆，考察地理上纵横交错的长短尺度；探寻神灵的所在之处，验证成功与失败的预兆；穷究道德的原理，探索生命本质的奥秘。}%

}



\bilingual{%

\zh{立規矩，準方圓，謹法度，約尺丈，立權衡，平重輕，剖毫釐，析黍絫；歷億載而不朽，施八極而無疆。散之不可勝究，斂之不盈掌握。}%

}{%

\zh{制定圆规和曲尺，确定方圆的标准；严谨地遵守法度，精确地测量尺寸；设立秤杆和秤砣，平衡物体的轻重；细分到毫厘之差，精确到黍粒之微；经历亿万年而不朽，施及八方极远而无边界。展开来无穷无尽不可穷尽，收拢来却不满一掌之握。}%

}



\bilingual{%

\zh{嚮之者富有餘，背之者貧且窶；心開者幼沖而即悟，意閉者皓首而難精。夫欲學之者必務量能揆己，志在所專。如是則焉有不成者哉。}%

}{%

\zh{投身其中的人会变得富足有余，背离它的人则会贫穷匮乏；心思通达的人即使年幼也能立即领悟，思路闭塞的人即便白头也难以精通。想要学习数学的人，一定要衡量自己的能力、正确评估自己，专心致志在选定的领域。这样的话，怎么会不成功呢？}%

}



\newpage



%% ========================================================================

%% VOLUME I

%% ========================================================================

\section{\zh{卷上}}



\subsection*{\zh{一、度量衡的起源}}



\bilingual{%

\zh{度之所起，起於忽。欲知其忽，蠶所生，吐絲為忽。十忽為一秒，十秒為一毫，十毫為一釐，十釐為一分，十分為一寸，十寸為一尺，十尺為一丈，十丈為一引；五十尺為一端；四十尺為一疋；六尺為一步。二百四十步為一畝。三百步為一里。}%

}{%

\zh{长度的单位，最小起于「忽」。什么叫忽呢？就是蚕吐出的丝那么细。十忽等于一秒，十秒等于一毫，十毫等于一釐，十釐等于一分，十分等于一寸，十寸等于一尺，十尺等于一丈，十丈等于一引；五十尺为一端；四十尺为一匹；六尺为一步。二百四十步为一亩。三百步为一里。}%

}



\subsection*{\zh{二、重量與容量的起源}}



\bilingual{%

\zh{稱之所起，起於黍。十黍為一絫，十絫為一銖，二十四銖為一兩，十六兩為一斤，三十斤為一鉤，四鉤為一石。}%

}{%

\zh{重量的单位，最小起于「黍」（一种谷物）。十黍等于一絫，十絫等于一銖，二十四銖等于一两，十六两等于一斤，三十斤等于一钩，四钩等于一石。}%

}



\bilingual{%

\zh{量之所起，起於粟。六粟為一圭，十圭為一抄，十抄為一撮，十撮為一勺，十勺為一合，十合為一升，十升為一斗，十斗為一斛。斛得六千萬粟。}%

}{%

\zh{容量的单位，最小起于「粟」（一粒谷子）。六粟等于一圭，十圭等于一抄，十抄等于一撮，十撮等于一勺，十勺等于一合，十合等于一升，十升等于一斗，十斗等于一斛。一斛等于六千万粟。}%

}



\subsection*{\zh{三、大數記法}}



\bilingual{%

\zh{凡大數之法，萬萬曰億，萬萬億曰兆，萬萬兆曰京，萬萬京曰陔，萬萬陔曰秭，萬萬秭曰壤，萬萬壤曰溝，萬萬溝曰澗，萬萬澗日正，萬萬正曰載。}%

}{%

\zh{大数的记法规则是：万万叫做亿，万亿叫做兆，万兆叫做京，万京叫做陔，万陔叫做秭，万秭叫做壤，万壤叫做沟，万沟叫做涧，万涧叫做正，万正叫做载。}%

}



\subsection*{\zh{四、圓周率與方圓比}}



\bilingual{%

\zh{周三徑一。方五邪七；見邪求方，五之，七而一；見方求邪，七之，五而一。}%

}{%

\zh{圆的周长是直径的三倍（即圆周率取3）。边长为5的正方形，对角线约为7；已知对角线求边长，乘以5再除以7；已知边长求对角线，乘以7再除以5。}%

}



\subsection*{\zh{五、材料比重}}



\bilingual{%

\zh{黃金方寸重一斤。白金方寸重一十四兩。玉方寸重一十二兩。銅方寸重七兩半。鈆方寸重九兩半。鐵方寸重六兩。石方寸重三兩。}%

}{%

\zh{边长为一寸的正方体黄金，重一斤。同样大小的白银重十四两。玉石重十二两。铜重七两半。铅重九两半。铁重六两。石头重三两。}%

}



\subsection*{\zh{六、算籌記數法}}



\bilingual{%

\zh{凡筭之法，先識其位，一從十橫，百立千僵，千十相望，萬百相當。}%

}{%

\zh{用算筹计算的方法，首先要认清数位：个位用纵式，十位用横式，百位用纵式，千位用横式。千位和十位的摆放方式相对，万位和百位的摆放方式对应。}%

}



\textit{\zh{算筹记数法是中国古代最重要的计算工具。个位数用竖式排列，十位数用横式排列，百位又变回竖式，千位再用横式，以此类推，纵横交替。}}



\subsection*{\zh{七、乘法法則}}



\bilingual{%

\zh{凡乘之法，重置其位。上下相觀，上位有十步至十，有百步至百，有千步至千。以上命下，所得之數列於中位。言十即過，不滿自如。上位乘訖者先去之。下位乘訖者則俱退之。六不積，五不隻。上下相乘，至盡則已。}%

}{%

\zh{乘法运算时，把两个数上下对齐摆放。根据被乘数的位数移动：被乘数有十，结果就进到十位；有百就进到百位；有千就进到千位。用上面的数依次乘下面的数，所得的结果放在中间位置。满十就进位，不满十就照原样写下。上面的数乘完后就撤去。下面的数乘完后也一起撤去。六以上不用累积算筹表示，五也不单独表示。上下相乘，直到全部算完为止。}%

}



\subsection*{\zh{八、除法法則}}



\bilingual{%

\zh{凡除之法，與乘正異。乘得在中央，除得在上方。假令六為法，百為實。以六除百，當進之二等，令在正百下，以六除一，則法多而實少，不可除，故當退就十位。以法除實，言一六而折百為四十，故可除。若實多法少，自當百之，不當復退。故或步法十者置於十位，百者置於百位。上位有空絕者，法退二位。餘法皆如乘時。實有餘者，以法命之，以法為母。實餘為子。}%

}{%

\zh{除法与乘法正好相反。乘法的结果放在中间，除法的结果放在上面。假设除数（法）是6，被除数（实）是100。用6除100，应该先前进两位，放在百位下面。但用6去除1，除数比被除数大，除不了，所以应该退到十位。用除数去除被除数：一六得六，从100里减去60还剩40，所以可以除。如果被除数大而除数小，自然应该在百位上商，不用再退了。所以除数是十位数就放在十位，是百位数就放在百位。如果上面某位是空的，除数就要退两位。其余步骤和乘法时一样。如果被除数除不尽有余数，就用除数来命名这个余数：除数作为分母，余数作为分子。}%

}



\subsection*{\zh{九、糧食換算率}}



\bilingual{%

\zh{以粟求糲米，三之，五而一。以糲米求粟，五之，三而一。以糲米求飯，五之，二而一。以粟米求糲飯，六之，四而一。以糲飯求糲米，二之，五而一。}%

}{%

\zh{用谷子换算糙米：谷子乘以3再除以5。用糙米换算谷子：糙米乘以5再除以3。用糙米换算米饭：糙米乘以5再除以2。用谷子换算糙米饭：谷子乘以6再除以4。用糙米饭换算糙米：糙米饭乘以2再除以5。}%

}



\subsection*{\zh{十、分數折扣}}



\bilingual{%

\zh{十分減一者，以二乘，二十除。減二者，以四乘，二十除。減三者，以六乘，二十除。減四者，以八乘，二十除。減五者，以十乘，二十除。減六者，以十二乘，二十除。減七者，以十四乘，二十除。減八者，以十六乘，二十除。減九者，以十八乘，二十除。}%

}{%

\zh{减少十分之一的，用2乘、20除。减少十分之二的，用4乘、20除。减少十分之三的，用6乘、20除。减少十分之四的，用8乘、20除。减少十分之五的，用10乘、20除。减少十分之六的，用12乘、20除。减少十分之七的，用14乘、20除。减少十分之八的，用16乘、20除。减少十分之九的，用18乘、20除。}%

}



\subsection*{\zh{例題：九九八十一自乘}}



\bilingual{%

\wen：\zh{九九八十一，自相乘，得幾何？}\\[2pt]

\da：\zh{答曰：六千五百六十一。}\\[2pt]

\shu：\zh{重置其位，以上八呼下八，八八六十四，即下六千四百於中位。以上八呼下一，一八如八，即於中位下八十。退下位一等，收上位八十。以上位一呼下八，一八如八，即於中位下八十。以上一呼下一，一一如一，即於中位下一。上下位俱收，中位即得六千五百六十一。}%

}{%

\wen：\zh{81乘81，结果是多少？}\\[2pt]

\da：\zh{答：6561。}\\[2pt]

\shu：\zh{运算方法：把81上下对齐摆放。用上面的8乘下面的8，八八六十四，在中间写下6400。用上面的8乘下面的1，一八得八，在中间写下80。把下面的数退一位，收去上面的80。用上面的1乘下面的8，一八得八，在中间写下80。用上面的1乘下面的1，一一得一，在中间写下1。上下两数都收去，中间的结果就是6561。}%

}



\subsection*{\zh{例題：六千五百六十一除以九}}



\bilingual{%

\wen：\zh{六千五百六十一，九人分之，問人得幾何？}\\[2pt]

\da：\zh{答曰：七百二十九。}\\[2pt]

\shu：\zh{先置六千五百六十一於中位，為實。下列九人為法。上位置七百，以上七呼下九，七九六十三，即除中位六千三百。退下位一等，即上位置二十。以上二呼下九，二九十八，即除中位一百八十。又更退下位一等，即上位更置九，即以上九呼下九，九九八十一，即除中位八十一。中位盡，收下位。上位所得即人之所得。}%

}{%

\wen：\zh{6561由9个人平分，每人得多少？}\\[2pt]

\da：\zh{答：每人729。}\\[2pt]

\shu：\zh{先把6561放在中间作为被除数。下面写9作为除数。上面先商700，用7乘下面的9，七九六十三，从中间减去6300。把下面的除数退一位，上面商20。用2乘下面的9，二九十八，从中间减去180。再把下面的除数退一位，上面改商9，用9乘下面的9，九九八十一，从中间减去81。中间减尽了，收去下面的除数。上面得到的结果729就是每人分得的数量。}%

}



\newpage



%% ========================================================================

%% VOLUME II

%% ========================================================================

\section{\zh{卷中}}



\subsection*{\zh{題一：約分}}



\bilingual{%

\wen：\zh{今有一十八分之一十二。問約之得幾何？}\\[2pt]

\da：\zh{答曰：三分之二。}\\[2pt]

\shu：\zh{置十八分在下，一十二分在上。副置二位，以少減多，等數得六，為法。約之，即得。}%

}{%

\wen：\zh{现有12/18。问约分后等于多少？}\\[2pt]

\da：\zh{答：等于2/3。}\\[2pt]

\shu：\zh{把分母18放在下面，分子12放在上面。在旁边再摆一份这两个数，用较大数减去较小数，辗转相减直到相等，得到的等数（最大公约数）是6，作为约分的除数。分子分母都用6来约分，就得到答案。}%

}



\subsection*{\zh{題五：分數平均}}



\bilingual{%

\wen：\zh{今有三分之一，三分之二，四分之三。問減多益少，幾何而平？}\\[2pt]

\da：\zh{答曰：減四分之三者二，減三分之二者一，并以益三分之一，各平於一十二分之七。}\\[2pt]

\shu：\zh{置三分、三分、四分在右方，之一、之二、之三在左方。母互乘子，副并得六十三，置右，為平實。母相乘，得三十六，為法。以列數三乘未并者及法。等數得九，約訖。}%

}{%

\wen：\zh{现有1/3、2/3、3/4这三个分数。问从多的里面减去多少补给少的，才能使三个分数相等？}\\[2pt]

\da：\zh{答：从3/4中减去2/12，从2/3中减去1/12，把这些都加给1/3，三个分数都等于7/12。}\\[2pt]

\shu：\zh{把三个分母3、3、4摆在右边，三个分子1、2、3摆在左边。分母互相乘对方的分子，辅助相加得到63，放在右边作为平均后的被除数。分母相乘得到36，作为除数。用分数的个数3去乘没合并的数和除数。等数是9，约分后即得。}%

}



\subsection*{\zh{題六：粟換糯米}}



\bilingual{%

\wen：\zh{今有粟一斗，問為糯米幾何？}\\[2pt]

\da：\zh{答曰：六升。}\\[2pt]

\shu：\zh{置粟一斗，十升。以糲米率三十乘之，得三百升，為實。以粟率五十為法，除之，即得。}%

}{%

\wen：\zh{现有谷子一斗，问能换多少糙米？}\\[2pt]

\da：\zh{答：六升。}\\[2pt]

\shu：\zh{把谷子一斗换算成10升。用糙米率30去乘，得到300升，作为被除数。用粟率50作为除数，300除以50，即得6升。}%

}



\subsection*{\zh{題九：砌磚計算}}



\bilingual{%

\wen：\zh{今有屋基南北三丈，東西六丈，欲以磚砌之。凡積二尺，用磚五枚。問計幾何？}\\[2pt]

\da：\zh{答曰：四千五百枚。}\\[2pt]

\shu：\zh{置東西六丈，以南北三丈乘之，得一千八百尺。以五乘之，得九千尺。以二除之，即得。}%

}{%

\wen：\zh{有一个房屋地基，南北长3丈，东西宽6丈，要用砖来砌。每2立方尺用5块砖。问一共需要多少块砖？}\\[2pt]

\da：\zh{答：4500块。}\\[2pt]

\shu：\zh{把东西宽6丈，乘以南北长3丈，得到1800平方尺（面积）。乘以5得9000。再除以2，即得4500块。}%

}



\subsection*{\zh{題十三：圓田面積}}



\bilingual{%

\wen：\zh{今有圓田周三百步，徑一百步。問得田幾何？}\\[2pt]

\da：\zh{答曰：三十一畝奇六十步。}\\[2pt]

\shu：\zh{先置周三百步，半之，得一百五十步。又置徑一百步，半之，得五十步。相乘，得七千五百步。以畝法二百四十步除之，即得。}%

}{%

\wen：\zh{有一块圆形田地，周长300步，直径100步。问这块田有多少亩？}\\[2pt]

\da：\zh{答：31亩余60步。}\\[2pt]

\shu：\zh{先把周长300步取半，得150步。再把直径100步取半，得50步。两个半数相乘，得7500平方步。用亩法240步去除，即得31亩余60步。}%

}



\textit{\zh{注：古代圆面积采用「半周半径相乘」的公式，与现代公式在取圆周率≈3时一致。}}



\bilingual{%

\zh{又術：周自相乘，得九萬步。以一十二除之，得七千五百步。以畝法除之，得畝數。}%

}{%

\zh{另一种算法：周长自乘得90000平方步。用12去除，得7500平方步。再用亩法去除，得到亩数。}%

}



\subsection*{\zh{題十四：方田桑生中央}}



\bilingual{%

\wen：\zh{今有方田，桑生中央。從角至桑一百四十七步。問為田幾何？}\\[2pt]

\da：\zh{答曰：一頃八十三畝奇一百八十步。}\\[2pt]

\shu：\zh{置角至桑一百四十七步，倍之，得二百九十四步。以五乘之，得一千四百七十步。以七除之，得二百一十步。自相乘，得四萬四千一百步。以二百四十步除之，即得。}%

}{%

\wen：\zh{有一块正方形田地，中央长着一棵桑树。从田地的一个角到桑树距离147步。问这块田面积是多少？}\\[2pt]

\da：\zh{答：1顷83亩余180步。}\\[2pt]

\shu：\zh{把角到桑树的距离147步记下，加倍得294步。乘以5得1470步。除以7得210步。这就是正方形田地的边长。自乘得44100平方步。用240步去除，即得亩数。}%

}



\subsection*{\zh{題二十三：開渠工程}}



\bilingual{%

\wen：\zh{今有穿渠，長二十九里一百四步，上廣一丈二尺六寸，下廣八尺，深一丈八尺。秋程人功三百尺。問須功幾何？}\\[2pt]

\da：\zh{答曰：三萬二千六百四十五人，不盡六十九尺六寸。}\\[2pt]

\shu：\zh{置里數，以三百步乘之，內零步，六之，得五萬二千八百二十四尺。并上、下廣，得二丈六寸。半之，以深乘之，得一百八十五尺四寸。以長乘，得九百七十九萬三千五百六十九尺六寸。以人功三百尺除之，即得。}%

}{%

\wen：\zh{要挖一条水渠，长29里104步，上口宽1丈2尺6寸，下口宽8尺，深1丈8尺。秋季每人每天的工作量标准是300立方尺。问需要多少人工？}\\[2pt]

\da：\zh{答：需要32645人，还余69.6立方尺的工作量。}\\[2pt]

\shu：\zh{把里数用300步去乘，加上零头的104步，再换算成尺（乘以6），得到52824尺（渠长）。把上口宽和下口宽相加，得2丈6寸。取半再乘以深度，得185.4平方尺（截面面积）。用长乘截面面积，得总体积。用每人工作量300尺去除，即得所需人工数。}%

}



\newpage



%% ========================================================================

%% VOLUME III

%% ========================================================================

\section{\zh{卷下}}



\subsection*{\zh{題一：九家輸租}}



\bilingual{%

\wen：\zh{今有甲、乙、丙、丁、戊、己、庚、辛、壬九家共輸租。甲出三十五斛，乙出四十六斛，丙出五十七斛，丁出六十八斛，戊出七十九斛，己出八十斛，庚出一百斛，辛出二百一十斛，壬出三百二十五斛。凡九家共輸租一千斛。僦運直折二百斛外，問家各幾何？}\\[2pt]

\da：\zh{答曰：甲二十八斛。乙三十六斛八蚪。丙四十五斛六蚪。丁五十四斛四蚪。戊六十三斛二蚪。己六十四斛。庚八十斛。辛一百六十八斛。壬二百六十斛。}\\[2pt]

\shu：\zh{置甲出三十五斛，以四乘之，得一百四十斛。以五除之，得二十八斛。乙出四十六斛，以四乘之，得一百八十四斛。以五除之，得三十六斛八蚪。……}%

}{%

\wen：\zh{有甲、乙、丙、丁、戊、己、庚、辛、壬九家共同交租。甲出35斛，乙出46斛，丙出57斛，丁出68斛，戊出79斛，己出80斛，庚出100斛，辛出210斛，壬出325斛。九家共交租1000斛。扣除运输费用200斛外，问每家实际应交多少？}\\[2pt]

\da：\zh{答：甲28斛，乙36斛8斗，丙45斛6斗，丁54斛4斗，戊63斛2斗，己64斛，庚80斛，辛168斛，壬260斛。}\\[2pt]

\shu：\zh{运算方法：每家的应交数先乘以4，再除以5（即扣除两成运费）。如甲家：35斛乘以4得140斛，除以5得28斛。乙家：46斛乘以4得184斛，除以5得36斛8斗。其余各家依此类推。}%

}



\subsection*{\zh{題三：平地聚粟}}



\bilingual{%

\wen：\zh{今有平地聚粟，下周三丈六尺，高四尺五寸。問粟幾何？}\\[2pt]

\da：\zh{答曰：一百斛。}\\[2pt]

\shu：\zh{置周三丈六尺，自相乘，得一千二百九十六尺。以高四尺五寸乘之，得五千八百三十二尺。以三十六除之，得一百六十二尺。以斛法一尺六寸二分除之，即得。}%

}{%

\wen：\zh{平地上堆着一堆谷子，底面周长3丈6尺，高4尺5寸。问这堆谷子有多少斛？}\\[2pt]

\da：\zh{答：100斛。}\\[2pt]

\shu：\zh{把周长3丈6尺记下，自乘得1296平方尺。乘以高4.5尺，得5832立方尺。除以36，得162立方尺（这是圆锥体积的古代算法）。再用一斛的容积1.62立方尺去除，即得100斛。}%

}



\subsection*{\zh{題五：碁局用子}}



\bilingual{%

\wen：\zh{今有碁局方一十九道。問用碁幾何？}\\[2pt]

\da：\zh{答曰：三百六十一。}\\[2pt]

\shu：\zh{置一十九道，自相乘之，即得。}%

}{%

\wen：\zh{有一个围棋盘，每边19道线。问盘面上有多少个交叉点（需要多少棋子）？}\\[2pt]

\da：\zh{答：361个。}\\[2pt]

\shu：\zh{把19道线记下，自乘（19×19），即得361。}%

}



\textit{\zh{这是标准的19路围棋盘，交叉点总数为 $19^2 = 361$。}}



\subsection*{\zh{題十五：人車問題}}



\bilingual{%

\wen：\zh{今有三人共車，二車空；二人共車，九人步。問人與車各幾何？}\\[2pt]

\da：\zh{答曰：一十五車。三十九人。}\\[2pt]

\shu：\zh{置二車，以三乘之，得六。加步者九人，得車一十五。欲知人者，以二乘車，加九人，即得。}%

}{%

\wen：\zh{现在有一些人乘车。如果3个人坐一辆车，会空出2辆车；如果2个人坐一辆车，会有9个人需要走路。问一共有多少人和多少辆车？}\\[2pt]

\da：\zh{答：15辆车，39人。}\\[2pt]

\shu：\zh{运算方法：把空出的2辆车用3去乘，得6。加上走路的9人，得15，就是车辆总数。要求人数，用2乘以车数15，再加上走路的9人，即得39人。}%

}



\subsection*{\zh{題十七：津吏問客}}



\bilingual{%

\wen：\zh{今有婦人河上蕩桮。津吏問曰：「桮何以多？」婦人曰：「家有客。」津吏曰：「客幾何？」婦人曰：「二人共飯，三人共羹，四人共肉，凡用桮六十五。不知客幾何？」}\\[2pt]

\da：\zh{答曰：六十人。}\\[2pt]

\shu：\zh{置六十五桮，以一十二乘之，得七百八十，以十三除之，即得。}%

}{%

\wen：\zh{有位妇人在河边洗杯子。渡口官吏问道：「怎么这么多杯子？」妇人回答：「家里来了客人。」官吏问：「有多少客人？」妇人说：「两人共用一碗饭，三人共用一碗汤，四人共用一碗肉，一共用了65只碗。不知道有多少客人？」}\\[2pt]

\da：\zh{答：60位客人。}\\[2pt]

\shu：\zh{运算方法：把65只碗记下，用12去乘，得780。用13去除，即得60人。（原理：每人用1/2饭碗+1/3汤碗+1/4肉碗=13/12只碗，所以总人数=65×12÷13=60）}%

}



\subsection*{\zh{題十八：木與繩}}



\bilingual{%

\wen：\zh{今有木，不知長短。引繩度之，餘繩四尺五寸。屈繩量之，不足一尺。問木長幾何？}\\[2pt]

\da：\zh{答曰：六尺五寸。}\\[2pt]

\shu：\zh{置餘繩四尺五寸，加不足一尺，共五尺五寸。倍之，得一丈一尺。減餘四尺五寸，即得。}%

}{%

\wen：\zh{有一根木头，不知道长度。用一根绳子去量，绳子多出4尺5寸。把绳子对折后再量，又短了1尺。问木头有多长？}\\[2pt]

\da：\zh{答：6尺5寸。}\\[2pt]

\shu：\zh{运算方法：把绳子的余量4尺5寸，加上对折后不足的1尺，共5尺5寸。加倍得1丈1尺（这就是绳子的全长）。再减去余量4尺5寸，即得木头长度6尺5寸。}%

}



\textit{\zh{验算：绳长=11尺，木长=6.5尺。用绳量木余4.5尺（11-6.5=4.5），绳对折后长5.5尺，量6.5尺的木头差1尺。正确。}}



\subsection*{\zh{題十九：器中餘米}}



\bilingual{%

\wen：\zh{今有器中米，不知其數。前人取半，中人三分取一，後人四分取一，餘米一斗五升。問本米幾何？}\\[2pt]

\da：\zh{答曰：六斗。}\\[2pt]

\shu：\zh{置餘米一斗五升，以六乘之，得九斗。以二除之，得四斗五升。以四乘之，得一斛八斗。以三除之，即得。}%

}{%

\wen：\zh{容器中有一些米，不知道原来有多少。第一个人取走了一半，第二个人取走了剩下的三分之一，第三个人取走了再剩下的四分之一，最后还剩1斗5升。问原来有多少米？}\\[2pt]

\da：\zh{答：6斗。}\\[2pt]

\shu：\zh{运算方法：把剩余的1斗5升记下，用6乘（因为最后剩的是原来的一半×三分之二×四分之三=1/4，所以反过来用6乘），得9斗。除以2得4斗5升。乘以4得1斛8斗。除以3，即得原来6斗。}%

}



\subsection*{\zh{題二十七：獸六首四足}}



\bilingual{%

\wen：\zh{今有獸六首四足，禽四首二足。上有七十六首，下有四十六足。問禽、獸各幾何？}\\[2pt]

\da：\zh{答曰：八獸，七禽。}\\[2pt]

\shu：\zh{倍足以減首，餘，半之，即禽。四乘禽數，減足，餘，半之，即獸。}%

}{%

\wen：\zh{有一种野兽有6个头4只脚，一种禽鸟有4个头2只脚。从上面数有76个头，从下面数有46只脚。问禽和兽各有多少只？}\\[2pt]

\da：\zh{答：8只兽，7只禽。}\\[2pt]

\shu：\zh{运算方法：把脚数加倍，减去头数，余数取半，就是禽的数量。用4乘以禽数，从脚数中减去，余数取半，就是兽的数量。}%

}



\subsection*{\zh{題二十九：百鹿入城}}



\bilingual{%

\wen：\zh{今有百鹿入城，家取一鹿，不盡；又三家共一鹿，適盡。問城中家幾何？}\\[2pt]

\da：\zh{答曰：七十五家。}\\[2pt]

\shu：\zh{以盈不足取之。假令七十二家，鹿不盡四。令之九十家，鹿不足二十。置七十二於右上，盈四於右下。置九十於左上，不足二十於左下。維乘之，所得，并為實。并盈不足為法。除之，即得。}%

}{%

\wen：\zh{有100头鹿进城，每家分一头鹿，分不完；又让三家合分一头鹿，正好分完。问城里有多少户人家？}\\[2pt]

\da：\zh{答：75家。}\\[2pt]

\shu：\zh{用盈不足术来解。假设72家，鹿有剩余4头。假设90家，鹿不够20头。把72放在右上方，剩余4放在右下方。把90放在左上方，不足20放在左下方。交叉相乘，所得相加作为被除数。把盈和不足相加作为除数。相除即得75家。}%

}



\textit{\zh{代数解法：设城中有 $x$ 家，则 $x + x/3 = 100$，解得 $x = 75$。}}



\subsection*{\zh{題三十一：雉兔同籠}}



\bilingual{%

\wen：\zh{今有雉、兔同籠，上有三十五頭，下九十四足。問雉、兔各幾何？}\\[2pt]

\da：\zh{答曰：雉二十三。兔一十二。}\\[2pt]

\shu：\zh{上置三十五頭，下置九十四足。半其足，得四十七。以少減多，再命之，上三除下三，上五除下五。下有一除上一，下有二除上二，即得。}%

}{%

\wen：\zh{有野鸡和兔子关在同一个笼子里，从上面数有35个头，从下面数有94只脚。问野鸡和兔子各有多少只？}\\[2pt]

\da：\zh{答：野鸡23只，兔子12只。}\\[2pt]

\shu：\zh{运算方法：上面放头数35，下面放脚数94。把脚数取半得47。用少减多（47-35=12），再按数位处理：上面的3除下面的3，上面的5除下面的5。下面有余1除上面1，有余2除上面2，即得兔子12只、野鸡23只。}%

}



\textit{\zh{这就是著名的「鸡兔同笼」问题，是中国数学的经典问题。现代解法：设雉 $x$ 只，兔 $y$ 只，则 $x+y=35$，$2x+4y=94$，解得 $y=12$，$x=23$。}}



\bilingual{%

\zh{又術曰：上置頭，下置足。半其足，以頭除足，以足除頭，即得。}%

}{%

\zh{另一种算法：上面放头数，下面放脚数。脚数取半，用头数去除脚数的一半，再用脚数的一半去除头数，即得两种动物的只数。}%

}



\subsection*{\zh{題三十四：出門望九隄}}



\bilingual{%

\wen：\zh{今有出門望見九隄。隄有九木，木有九枝，枝有九巢，巢有九禽，禽有九雛，雛有九毛，毛有九色。問各幾何？}\\[2pt]

\da：\zh{答曰：木八十一。枝七百二十九。巢六千五百六十一。禽五萬九千四十九。雛五十三萬一千四百四十一。毛四百七十八萬二千九百六十九。色四千三百四萬六千七百二十一。}\\[2pt]

\shu：\zh{置九隄，以九乘之，得木之數。又以九乘之，得枝之數。又以九乘之，得巢之數。又以九乘之，得禽之數。又以九乘之，得雛之數。又以九乘之，得毛之數。又以九乘之，得色之數。}%

}{%

\wen：\zh{出门望见9道堤坝。每道堤坝上有9棵树，每棵树有9根枝，每根枝上有9个鸟巢，每个鸟巢里有9只鸟，每只鸟有9只雏鸟，每只雏鸟有9根羽毛，每根羽毛有9种颜色。问各有多少？}\\[2pt]

\da：\zh{答：树81棵，枝729根，巢6561个，鸟59049只，雏鸟531441只，羽毛4782969根，颜色43046721种。}\\[2pt]

\shu：\zh{运算方法：把9道堤坝用9去乘，得到树的数量。再连续用9去乘，依次得到枝、巢、禽、雏、毛、色的数量。}%

}



\textit{\zh{本题即连续求9的幂次：$9^2=81$，$9^3=729$，$9^4=6561$，$9^5=59049$，$9^6=531441$，$9^7=4782969$，$9^8=43046721$。}}



\subsection*{\zh{題三十五：三女歸家}}



\bilingual{%

\wen：\zh{今有三女，長女五日一歸，中女四日一歸，少女三日一歸。問三女幾何日相會？}\\[2pt]

\da：\zh{答曰：六十日。}\\[2pt]

\shu：\zh{置長女五日、中女四日、少女三日，於右方。各列一算於左方。維乘之，各得所到數：長女十二到，中女十五到，少女二十到。又各以歸日乘到數，即得。}%

}{%

\wen：\zh{有三个女儿，大女儿每5天回家一次，二女儿每4天回家一次，小女儿每3天回家一次。问三个女儿多少天后能同时回家相会？}\\[2pt]

\da：\zh{答：60天后。}\\[2pt]

\shu：\zh{运算方法：把大女的5天、中女的4天、小女的3天摆在右边。左边各列一个算筹。交叉相乘，得到各自回家的次数：大女回12次，中女回15次，小女回20次。再用各自回家的间隔天数乘以次数，都等于60天。}%

}



\textit{\zh{本题即求最小公倍数：$\text{lcm}(3,4,5)=60$。}}



\newpage



%% ========================================================================

%% THE CHINESE REMAINDER THEOREM

%% ========================================================================

\section*{\zh{物不知數}——\zh{中國剩餘定理}}

\addcontentsline{toc}{section}{\zh{物不知數——中國剩餘定理}}



\begin{tcolorbox}[colback=crtgold!8,colframe=crtgold,boxrule=1.5pt,

  left=12pt,right=12pt,top=10pt,bottom=10pt,arc=4mm,

  title={\textbf{\large \zh{卷下第26問}}},

  fonttitle=\bfseries\color{white},coltitle=white,

  colbacktitle=crtgold]



\bilingual{%

\wen：\zh{今有物，不知其數。三、三數之，賸二；五、五數之，賸三；七、七數之，賸二。問物幾何？}\\[2pt]

\da：\zh{答曰：二十三。}\\[2pt]

\shu：\zh{「三、三數之，賸二」，置一百四十；「五、五數之，賸三」，置六十三；「七、七數之，賸二」，置三十。并之，得二百三十三。以二百一十減之，即得。}\\[3pt]

\zh{凡三、三數之，賸一，則置七十；五、五數之，賸一，則置二十一；七、七數之，賸一，則置十五。一百六以上，以一百五減之，即得。}%

}{%

\wen：\zh{有一些物品，不知道有多少个。三个三个地数，最后剩两个；五个五个地数，最后剩三个；七个七个地数，最后剩两个。问这些物品一共有多少个？}\\[2pt]

\da：\zh{答：23个。}\\[2pt]

\shu：\zh{「三个三个数余2」，就写下140；「五个五个数余3」，就写下63；「七个七个数余2」，就写下30。把这三个数相加，得233。减去210（3×5×7=105的2倍），即得23。}\\[3pt]

\zh{一般规律：三个三个数余1，就写70；五个五个数余1，就写21；七个七个数余1，就写15。如果总数超过106，就减去105（3×5×7），直到得到小于105的答案。}%

}

\end{tcolorbox}



\subsection*{\zh{現代數學分析}}



\zh{用现代数学符号表示，这个问题要求找到最小的正整数 $N$，满足以下同余方程组：}

\begin{align*}

N &\equiv 2 \pmod{3} \\

N &\equiv 3 \pmod{5} \\

N &\equiv 2 \pmod{7}

\end{align*}



\zh{孙子的算法给出：}

\begin{align*}

N &= 140 + 63 + 30 - 210 \\

  &= 233 - 210 \\

  &= 23

\end{align*}



\zh{三个关键数字 \textbf{70}、\textbf{21}、\textbf{15} 具有特殊的数学性质：}

\begin{itemize}[leftmargin=2em]

  \item \zh{$70 = 2 \times 5 \times 7 \equiv 1 \pmod{3}$，且 $70 \equiv 0 \pmod{5}$，$70 \equiv 0 \pmod{7}$}

  \item \zh{$21 = 1 \times 3 \times 7 \equiv 1 \pmod{5}$，且 $21 \equiv 0 \pmod{3}$，$21 \equiv 0 \pmod{7}$}

  \item \zh{$15 = 1 \times 3 \times 5 \equiv 1 \pmod{7}$，且 $15 \equiv 0 \pmod{3}$，$15 \equiv 0 \pmod{5}$}

\end{itemize}



\zh{由于 $3 \times 5 \times 7 = 105$，一般解为：}

\[N = 70r_1 + 21r_2 + 15r_3 - 105k\]

\zh{其中 $r_1, r_2, r_3$ 分别是除以3、5、7的余数，$k$ 是选取的整数使 $N$ 为最小正整数。}



\subsection*{\zh{驗算}}



\begin{align*}

23 \div 3 &= 7 \text{ \zh{余} } \mathbf{2} \quad \checkmark \\

23 \div 5 &= 4 \text{ \zh{余} } \mathbf{3} \quad \checkmark \\

23 \div 7 &= 3 \text{ \zh{余} } \mathbf{2} \quad \checkmark

\end{align*}



\subsection*{\zh{孫子公式}}



\zh{对于一般的余数 $r_1, r_2, r_3$：}

\[

\boxed{N = 70 \times r_1 + 21 \times r_2 + 15 \times r_3 - 105 \times k}

\]



\zh{其中 $M = 3 \times 5 \times 7 = 105$ 是最小公倍数，且：}

\begin{itemize}

  \item \zh{$70 = 105/3 \times 2$ 是模3余数的乘数}

  \item \zh{$21 = 105/5 \times 1$ 是模5余数的乘数}

  \item \zh{$15 = 105/7 \times 1$ 是模7余数的乘数}

\end{itemize}



\subsection*{\zh{孫子歌}}



\zh{明代数学家程大位在《算法统宗》（1592年）中记录了这首朗朗上口的口诀：}



\begin{tcolorbox}[colback=chinesebg,colframe=rulecolor,boxrule=0.8pt,

  width=0.9\textwidth,arc=2mm]

\centering\large

\zh{三人同行七十稀，五樹梅花廿一枝，}\\[3pt]

\zh{七子團圓正半月，除百零五便得知。}

\end{tcolorbox}



\begin{tcolorbox}[colback=modernbg,colframe=rulecolor!70!blue,boxrule=0.8pt,

  width=0.9\textwidth,arc=2mm]

\centering

\zh{「三人同行」对应模3的乘数70；「五树梅花」对应模5的乘数21；}

\zh{「七子团圆半月（15日）」对应模7的乘数15；}

\zh{「除百零五」指减去105（3×5×7）直到得到最小正解。}

\end{tcolorbox}



\newpage



%% ========================================================================

%% REMAINING PROBLEMS FROM VOLUME III

%% ========================================================================

\section*{\zh{卷下其餘算題}}

\addcontentsline{toc}{section}{\zh{卷下其餘算題}}



\subsection*{\zh{題二十八：甲乙持錢}}



\bilingual{%

\wen：\zh{今有甲、乙二人，持錢各不知數。甲得乙中半，可滿四十八。乙得甲太半，亦滿四十八。問甲、乙二人元持錢各幾何？}\\[2pt]

\da：\zh{答曰：甲持錢三十六。乙持錢二十四。}\\[2pt]

\shu：\zh{如方程求之。置二甲、一乙、錢九十六，於右方。置二甲、三乙、錢一百四十四，於左方。以右方二乘左方，上得四，中得六，下得二百八十八錢。以右行再減左行，左上空，中餘四乙，為法；下餘九十六錢，為實。上法，下實，得二十四錢，為乙錢。以減右下九十六，餘七十二，為實；以右上二甲為法。上法、下實，得三十六，為甲錢也。}%

}{%

\wen：\zh{甲、乙两人各有一些钱，但不知道具体数目。如果甲得到乙的一半钱，他就有48文。如果乙得到甲的三分之二钱，他也有48文。问甲、乙原来各有多少钱？}\\[2pt]

\da：\zh{答：甲原有36文，乙原有24文。}\\[2pt]

\shu：\zh{用方程组的方法来解。列出方程：甲 + 乙/2 = 48，乙 + 2甲/3 = 48。转化为：2甲 + 乙 = 96，2甲 + 3乙 = 144。用右行乘左行，再通过相减消元，解得乙有24文，甲有36文。}%

}



\subsection*{\zh{題三十：三雞啄粟}}



\bilingual{%

\wen：\zh{今有三雞共啄粟一千一粒。雛啄一，母啄二，翁啄四。主責本粟，三雞主各償幾何？}\\[2pt]

\da：\zh{答曰：雞雛主一百四十三。雞母主二百八十六。雞翁主五百七十二。}\\[2pt]

\shu：\zh{置粟一千一粒，為實。副并三雞所啄粟七粒，為法。除之，得一百四十三粒，為雞雛主所償之數。遞倍之，即得母、翁主所償之數。}%

}{%

\wen：\zh{有三只鸡一共啄了1001粒谷子。小鸡啄1粒，母鸡啄2粒，公鸡啄4粒。主人要追讨损失的谷子，问三只鸡的主人各应赔偿多少？}\\[2pt]

\da：\zh{答：小鸡的主人赔143粒，母鸡的主人赔286粒，公鸡的主人赔572粒。}\\[2pt]

\shu：\zh{把1001粒谷子作为被除数。把三只鸡所啄的谷子数相加（1+2+4=7粒）作为除数。1001除以7得143粒，这就是小鸡主人应赔的数。依次加倍，就得到母鸡主人和公鸡主人应赔的数。}%

}



\subsection*{\zh{題三十二：九里渠魚}}



\bilingual{%

\wen：\zh{今有九里渠，三寸魚，頭頭相次。問魚得幾何？}\\[2pt]

\da：\zh{答曰：五萬四千。}\\[2pt]

\shu：\zh{置九里，以三百步乘之，得二千七百步。又以六尺乘之，得一萬六千二百尺。上十之，得一十六萬二千寸。以魚三寸除之，即得。}%

}{%

\wen：\zh{有一条9里长的水渠，每条鱼长3寸，鱼一条接一条头尾相连排过去。问能排多少条鱼？}\\[2pt]

\da：\zh{答：54000条。}\\[2pt]

\shu：\zh{把9里用300步去乘，得2700步。每步6尺，换算成尺得16200尺。再换算成寸（乘以10）得162000寸。用每条鱼3寸去除，即得54000条。}%

}



\subsection*{\zh{題三十三：長安至洛陽}}



\bilingual{%

\wen：\zh{今有長安、洛陽相去九百里。車輪一匝一丈八尺。欲自洛陽至長安，問輪匝幾何？}\\[2pt]

\da：\zh{答曰：九萬匝。}\\[2pt]

\shu：\zh{置九百里，以三百步乘之，得二十七萬步。又以六尺乘之，得一百六十二萬尺。以車輪一丈八尺為法。除之，即得。}%

}{%

\wen：\zh{长安和洛阳相距900里。车轮转一圈前进1丈8尺。要从洛阳到长安，问车轮要转多少圈？}\\[2pt]

\da：\zh{答：90000圈。}\\[2pt]

\shu：\zh{把900里用300步去乘，得270000步。每步6尺，共1620000尺。用车轮周长1丈8尺（18尺）去除，即得90000圈。}%

}



\subsection*{\zh{題三十六：孕婦知男女}}



\bilingual{%

\wen：\zh{今有孕婦行年二十九，難九月。未知所生？}\\[2pt]

\da：\zh{答曰：生男。}\\[2pt]

\shu：\zh{置四十九，加難月，減行年。所餘，以天除一，地除二，人除三，四時除四，五行除五，六律除六，七星除七，八風除八，九州除九。其不盡者，奇則為男，耦則為女。}%

}{%

\wen：\zh{有一位孕妇，今年29岁，怀孕9个月。问生男生女？}\\[2pt]

\da：\zh{答：生男孩。}\\[2pt]

\shu：\zh{运算方法：把49加上怀孕月份9，再减去年龄29，得29。然后依次减去：天减1、地减2、人减3、四时减4、五行减5、六律减6、七星减7、八风减8、九州减9。最后剩余的数，如果是奇数就生男孩，如果是偶数就生女孩。}%

}



\textit{\zh{注：此为古代迷信算法，没有科学依据。计算过程：49+9-29=29，29-1-2-3-4-5-6=8，8是偶数，按原文程序应生女。但原文答「生男」，说明古代版本或有不同，或仅为占卜之术。}}



\newpage



%% ========================================================================

%% EPILOGUE

%% ========================================================================

\section*{\zh{後記}}

\addcontentsline{toc}{section}{\zh{後記}}



\bilingual{%

\zh{《孫子算經》雖然篇幅不長，內容也不算深奧，但它在數學史上的地位無可替代。它不僅系統地記錄了中國古代的度量衡制度、算籌記數法和四則運算方法，更以其「物不知數」問題為世界數學寶庫貢獻了最早的同餘式組解法。}%

}{%

\zh{《孙子算经》虽然篇幅不长，内容也不算深奥，但它在数学史上的地位无可替代。它不仅系统地记录了中国古代的度量衡制度、算筹记数法和四则运算方法，更以其「物不知数」问题为世界数学宝库贡献了最早的一次同余式方程组解法。}%

}



\bilingual{%

\zh{從孫子到秦九韶，中國數學家對同餘理論的研究持續了近千年，最終發展成為完善的「大衍求一術」。而在歐洲，類似的方法直到歐拉（1743年）和高斯（1801年）才得以建立。這充分說明了中國古代數學在世界數學發展史上的重要地位。}%

}{%

\zh{从孙子到秦九韶，中国数学家对同余理论的研究持续了将近一千年，最终发展成为完善的「大衍求一术」。而在欧洲，类似的方法直到欧拉（1743年）和高斯（1801年）才得以建立。这充分说明了中国古代数学在世界数学发展史上的重要地位。}%

}



\bilingual{%

\zh{本書的白話文翻譯力求準確通俗，保留原文的數學術語，同時用現代語言加以解釋，以便當代讀者理解這部千年數學經典的內容與方法。}%

}{%

\zh{本书的白话文翻译力求准确通俗，保留原文的数学术语，同时用现代语言加以解释，以便当代读者理解这部千年数学经典的内容与方法。}%

}



\vspace{1cm}

\begin{center}

\longline

\end{center}



\vspace{0.5cm}



\begin{center}

\zh{本文古典原文依據中國哲學書電子化計劃（ctext.org）所藏《孫子算經》整理}\\[3pt]

\zh{白話文翻譯由現代學者參考古今注疏編譯}

\end{center}



\end{document}







% ============================================================

% Source: translations/zh/chinese/yang-hui-xiangjie-jiuzhang-part1_classical-modern_bilingual.tex

% ============================================================



%% ========================================================================

%% Yang Hui's Xiangjie Jiuzhang Suanfa (Part I)

%% 楊輝《詳解九章算法》(上)

%% Bilingual Edition: Classical Chinese (文言文) | Modern Chinese (白话文)

%% ========================================================================



\documentclass[11pt,a4paper]{article}



%% -------- Core Packages --------

\usepackage{fontspec}

\usepackage{polyglossia}

\usepackage{amsmath,amssymb,amsthm}

\usepackage{geometry}

\usepackage{graphicx}

\usepackage{xcolor}

\usepackage{titlesec}

\usepackage{fancyhdr}

\usepackage{enumitem}

\usepackage{array,booktabs,longtable,multirow}

\usepackage{hyperref}

\usepackage{microtype}

\usepackage{tocloft}

\usepackage{indentfirst}

\usepackage{tikz}

\usepackage{lettrine}

\usepackage{float}

\usepackage{caption}

\usepackage{subcaption}

\usepackage{etoolbox}



%% -------- Paracol for bilingual parallel columns --------

\usepackage{paracol}

\columnratio{0.50}  %% Equal columns

\setlength{\columnsep}{1.2em}

\setlength{\columnseprule}{0.4pt}



\usetikzlibrary{arrows.meta,positioning,calc,decorations.pathreplacing}



%% -------- Page Geometry --------

\geometry{

  paperwidth=210mm,paperheight=297mm,

  left=18mm,right=18mm,top=25mm,bottom=25mm,

  headheight=14pt,footskip=30pt

}



%% -------- Language & Font Setup --------

\setmainlanguage{english}



\setmainfont{Times New Roman}

\setsansfont{Arial}

\setmonofont{Courier New}



%% Chinese Fonts via xeCJK

\usepackage{xeCJK}

\setCJKmainfont{Microsoft YaHei}

\setCJKsansfont{Microsoft YaHei}

\setCJKmonofont{Microsoft YaHei}

\setCJKfamilyfont{classical}{Microsoft YaHei}

\setCJKfamilyfont{modern}{Microsoft YaHei}



%% -------- Colors --------

\definecolor{titlecolor}{RGB}{45,55,72}

\definecolor{sectioncolor}{RGB}{55,65,81}

\definecolor{notecolor}{RGB}{120,53,15}

\definecolor{rulecolor}{RGB}{180,160,140}

\definecolor{classicalbg}{RGB}{252,248,240}

\definecolor{modernbg}{RGB}{240,248,252}

\definecolor{problemframe}{RGB}{139,119,101}

\definecolor{solutionbg}{RGB}{245,245,240}



%% -------- Section Formatting --------

\titleformat{\section}

  {\normalfont\Large\bfseries\color{sectioncolor}}

  {\thesection}{1em}{}

\titleformat{\subsection}

  {\normalfont\large\bfseries\color{sectioncolor}}

  {\thesubsection}{1em}{}

\titleformat{\subsubsection}

  {\normalfont\normalsize\bfseries\color{sectioncolor}}

  {\thesubsubsection}{1em}{}



%% -------- Header/Footer --------

\pagestyle{fancy}

\fancyhf{}

\fancyhead[L]{\small\textit{楊輝《詳解九章算法》— 文白對照版}}

\fancyhead[R]{\small\thepage}

\renewcommand{\headrulewidth}{0.4pt}

\renewcommand{\footrulewidth}{0pt}



%% -------- Custom Commands --------

\newcommand{\longline}{\noindent\rule{\textwidth}{0.4pt}}

\newcommand{\shortline}{\noindent\rule{0.5\textwidth}{0.4pt}\par}

\newcommand{\cnote}[1]{\textcolor{notecolor}{\textit{#1}}}



%% Classical Chinese text helper (left column)

\newcommand{\ClassicalText}[1]{%

  \begin{center}\small\textbf{【原文】}\end{center}

  \vspace{-0.5em}

  {\CJKfamily{classical}\fboxsep=4pt\colorbox{classicalbg}{%

    \parbox{\linewidth}{\small#1%

  }}\vspace{0.5em}}

}



%% Modern Chinese text helper (right column)

\newcommand{\ModernText}[1]{%

  \begin{center}\small\textbf{【白话译文】}\end{center}

  \vspace{-0.5em}

  {\CJKfamily{modern}\fboxsep=4pt\colorbox{modernbg}{%

    \parbox{\linewidth}{\small#1%

  }}\vspace{0.5em}}

}



%% Column labels

\newcommand{\LeftLabel}{\begin{center}\textcolor{sectioncolor}{\small\textbf{文言文（原文）}}\end{center}\vspace{-0.5em}}

\newcommand{\RightLabel}{\begin{center}\textcolor{sectioncolor}{\small\textbf{白话文（今译）}}\end{center}\vspace{-0.5em}}



%% Problem environments

\newenvironment{classicalproblem}[1]{%

  \par\smallskip%

  \noindent\textbf{\textcolor{problemframe}{問#1：}}%

}{%

  \par\smallskip

}



\newenvironment{modernproblem}[1]{%

  \par\smallskip%

  \noindent\textbf{\textcolor{problemframe}{问题#1：}}%

}{%

  \par\smallskip

}



%% -------- Hyperref Setup --------

\hypersetup{

  colorlinks=true,

  linkcolor=sectioncolor,

  citecolor=sectioncolor,

  urlcolor=sectioncolor,

  pdfencoding=auto,

  pdftitle={楊輝詳解九章算法 — 文白對照版},

  pdfauthor={楊輝（南宋）},

}



%% -------- TOC Depth --------

\setcounter{tocdepth}{2}

\setcounter{secnumdepth}{2}



%% ========================================================================

\begin{document}



%% ========================================================================

%% TITLE PAGE

%% ========================================================================

\begin{titlepage}

\centering

\vspace*{1.5cm}



{\Huge\bfseries\color{titlecolor} 詳解九章算法}\\[0.8em]

{\LARGE\color{sectioncolor} 楊輝 \quad 著}\\[0.5em]

{\large\color{sectioncolor} 南宋（約 1238--1298 年）}\\[2.5em]



{\Large\bfseries\color{titlecolor} 文白對照版}\\[0.5em]

{\large\color{sectioncolor} Classical Chinese -- Modern Chinese Bilingual Edition}\\[2em]



\longline\\[1.5em]



{\large\color{sectioncolor}（上卷 Part I）}\\[1em]

\longline\\[2em]



\begin{tikzpicture}[scale=0.75]

%% Yang Hui's Triangle (Pascal's Triangle)

\foreach \n in {0,...,6}{

  \foreach \k in {0,...,\n}{

    \pgfmathsetmacro{\y}{-\n*0.8}

    \pgfmathsetmacro{\x}{\k*0.8-\n*0.4}

    \ifnum\n=0\ifnum\k=0\node at (\x,\y) {1};\fi\fi

    \ifnum\n=1\ifnum\k=0\node at (\x,\y) {1};\fi\ifnum\k=1\node at (\x,\y) {1};\fi\fi

    \ifnum\n=2\ifnum\k=0\node at (\x,\y) {1};\fi\ifnum\k=1\node at (\x,\y) {2};\fi\ifnum\k=2\node at (\x,\y) {1};\fi\fi

    \ifnum\n=3\ifnum\k=0\node at (\x,\y) {1};\fi\ifnum\k=1\node at (\x,\y) {3};\fi\ifnum\k=2\node at (\x,\y) {3};\fi\ifnum\k=3\node at (\x,\y) {1};\fi\fi

    \ifnum\n=4\ifnum\k=0\node at (\x,\y) {1};\fi\ifnum\k=1\node at (\x,\y) {4};\fi\ifnum\k=2\node at (\x,\y) {6};\fi\ifnum\k=3\node at (\x,\y) {4};\fi\ifnum\k=4\node at (\x,\y) {1};\fi\fi

    \ifnum\n=5\ifnum\k=0\node at (\x,\y) {1};\fi\ifnum\k=1\node at (\x,\y) {5};\fi\ifnum\k=2\node at (\x,\y) {10};\fi\ifnum\k=3\node at (\x,\y) {10};\fi\ifnum\k=4\node at (\x,\y) {5};\fi\ifnum\k=5\node at (\x,\y) {1};\fi\fi

    \ifnum\n=6\ifnum\k=0\node at (\x,\y) {1};\fi\ifnum\k=1\node at (\x,\y) {6};\fi\ifnum\k=2\node at (\x,\y) {15};\fi\ifnum\k=3\node at (\x,\y) {20};\fi\ifnum\k=4\node at (\x,\y) {15};\fi\ifnum\k=5\node at (\x,\y) {6};\fi\ifnum\k=6\node at (\x,\y) {1};\fi\fi

  }

}

\node[below] at (0,-5.5) {\small 開方作法本源圖（楊輝三角）};

\end{tikzpicture}



\vfill



{\small 以 XeLaTeX 排版 \quad Microsoft YaHei 字體}\\[0.3em]

{\small \today}



\end{titlepage}



%% ========================================================================

%% TABLE OF CONTENTS

%% ========================================================================

\newpage

\tableofcontents

\newpage



%% ========================================================================

%% INTRODUCTION

%% ========================================================================

\section{導讀 Introduction}



\begin{paracol}{2}

\LeftLabel



{\CJKfamily{classical}

《詳解九章算法》為南宋數學家\textbf{楊輝}（約 1238--1298 年）所撰，成書於 1261 年。此書對漢代《九章算術》進行了詳細的疏解與評注，是中國數學史上最重要的著作之一。



原序云：「輝嘗聞學者，謂九章題問頗隱，法理難明，不得其門而入。」楊輝遂「以答參問，用草考法，因法推類」，使學者得以窺其門徑。



全書凡十二卷：

\begin{enumerate}[noitemsep,label=\arabic*.,leftmargin=1.5em]

\item 乘除（41 題）

\item 方田（38 題）

\item 粟米（46 題）

\item 衰分（20 題）

\item 少廣（24 題）

\item 商功（28 題）

\item 均輸（28 題）

\item 盈不足（20 題）

\item 方程（18 題）

\item 勾股（38 題）

\item 篡類

\end{enumerate}



其編纂體例為：\textbf{問}（設題）$\rightarrow$ \textbf{答}（答數）$\rightarrow$ \textbf{術}（算法）$\rightarrow$ \textbf{草}（演草，即詳細計算步驟）。此四段式結構成為中國數學著述之典範。

}



\switchcolumn

\RightLabel



{\CJKfamily{modern}

《詳解九章算法》是南宋数学家\textbf{楊輝}（約 1238--1298 年）編撰的重要數學著作，成書於 1261 年。本書對漢代經典《九章算術》進行了系統性的詳細註解和評論，是中國數學史上最具影響力的作品之一。



楊輝在原序中說：「我曾聽學者們說，《九章》的題目頗為隱晦，其中的原理和方法難以理解，讓人找不到入門的途徑。」於是他「用答案來對照問題，用演算過程來考證算法，再根據算法推廣類比」，讓後來的學者能夠找到學習的門徑。



全書共分為十二章：乘除法（41 題）、方田（38 題）、粟米（46 題）、衰分（20 題）、少廣（24 題）、商功（28 題）、均輸（28 題）、盈不足（20 題）、方程（18 題）、勾股（38 題）、以及篡類。



楊輝的編排體例是：先「問」（提出問題），再「答」（給出答案），然後「術」（闡述通用算法），最後「草」（提供詳細的逐步計算過程）。這種四段式結構成為了中國數學著述的典範。

}



\end{paracol}



\subsection{歷史背景 Historical Context}



\begin{paracol}{2}

\LeftLabel



{\CJKfamily{classical}

《九章算術》原書成於漢代（約公元前 200 年--公元 200 年），收錄 246 道數學題，分為九章。至南宋時，原文因以下原因漸難理解：



\begin{itemize}[noitemsep,leftmargin=1.5em]

\item 靖康之變（1127 年）導致大量書籍散佚

\item 傳抄過程中訛誤叢生

\item 題設隱晦，古法難明

\end{itemize}



楊輝序曰：「昔聖宋紹興戊辰，算士榮棨謂靖康以來，罕有舊本，間有存者，狃於末習，向獲善本，得其全經，復起於學。」

}



\switchcolumn

\RightLabel



{\CJKfamily{modern}

《九章算術》原本編成於漢代（約公元前 200 年至公元 200 年），收錄了 246 道數學問題，分為九章。到了南宋時期，由於以下原因，原文已經變得難以理解：



\begin{itemize}[noitemsep,leftmargin=1.5em]

\item 靖康之變（1127 年）時大量典籍散失

\item 輾轉傳抄過程中產生了許多錯誤

\item 題目表述隱晦，古代算法難以理解

\end{itemize}



楊輝在序中說：「在神聖的宋朝紹興年間的戊辰年，數學家榮棨說，自從靖康年間以來，舊版本已經很罕見了，偶爾有幸存下來的，也因為後人的不良習慣而被篡改。後來找到了一個善本，恢復了全書，學術研究才得以重新興起。」

}



\end{paracol}



\newpage



%% ========================================================================

%% PREFACE

%% ========================================================================

\section{原序 Original Preface}



\begin{paracol}{2}

\LeftLabel

\ClassicalText{

黃帝九章古序云：國家嘗設算科取士，選九章以為算經之首，蓋猶儒者之六經，醫家之難素，兵法之孫子歟。

}



\switchcolumn

\RightLabel

\ModernText{

《黃帝九章》的古序說：國家曾經設立數學科目來選拔人才，把《九章》選為數學經典之首，這就好像儒家學者的六經、醫家的《難經》和《素問》、兵家的《孫子兵法》一樣重要啊。

}



\end{paracol}



\vspace{0.8em}



\begin{paracol}{2}

\LeftLabel

\ClassicalText{

昔聖宋紹興戊辰，算士榮棨謂靖康以來，罕有舊本，間有存者，狃於末習，向獲善本，得其全經，復起於學。以魏景元元年劉徽等、唐朝議大夫、行太史令上輕車都尉李淳風等注釋，聖宋右班直賈憲撰草。

}



\switchcolumn

\RightLabel

\ModernText{

從前，在我大宋紹興年間的戊辰年，數學家榮棨說，自從靖康之變以來，舊版本已經很罕見了，偶爾有幸存下來的，也被後人的陋習所污染。後來獲得了一個善本，得到了完整的經書，學術研究才得以重新興起。這個版本有魏國景元元年劉徽等人、唐代的朝議大夫、代理太史令、上輕車都尉李淳風等人的注釋，還有宋朝右班直賈憲撰寫的演草。

}



\end{paracol}



\vspace{0.8em}



\begin{paracol}{2}

\LeftLabel

\ClassicalText{

輝嘗聞學者，謂九章題問頗隱，法理難明，不得其門而入。於是以答參問，用草考法，因法推類，然後知斯文非古之全經也。將後賢補贅之文，修前代已廢之法，刪立題術，又篡法問，詳著於後，儻得賢者，改而正諸，是所願也。

}



\switchcolumn

\RightLabel

\ModernText{

我曾經聽學者們說，《九章》的題目頗為隱晦，其中的原理和方法難以理解，讓人找不到入門的途徑。於是我用答案來對照問題，用演算過程來考證算法，再根據算法推廣類比，這之後才知道這部書並不是古代完整的經典。我把後代賢人所補充的文字、前代已經廢棄的方法加以修訂，重新確立題目和算法，又編纂了方法和問題，詳細地記錄在後面。如果能夠得到賢能之士加以改正完善，那就是我的心願了。

}



\end{paracol}



\newpage



%% ========================================================================

%% CHAPTER 1: MULTIPLICATION AND DIVISION

%% ========================================================================

\section{第一章：乘除法 Multiplication and Division}



\begin{paracol}{2}

\LeftLabel

{\CJKfamily{classical}\small

本章論述乘除等基本運算，共 41 題，多取材於《方田》《粟米》諸章。

}



\switchcolumn

\RightLabel

{\CJKfamily{modern}\small

本章討論乘法和除法等基本運算，共 41 道題目，大多取材於《方田》《粟米》等章節。

}



\end{paracol}



\subsection{直田法 Rectangular Field Method}



\begin{paracol}{2}

\LeftLabel

{\CJKfamily{classical}\small

\textbf{術曰：}以廣乘從得積。以畝法（二百四十步）除之。



\begin{classicalproblem}[一]

廣十五步，從十六步，問田幾何？答曰：一畝。

\end{classicalproblem}

}



\switchcolumn

\RightLabel

{\CJKfamily{modern}\small

\textbf{算法：}用寬度乘以長度得到面積。用一畝的標準（240 平方步）去除。



\begin{modernproblem}[第一題]

有一塊田，寬 15 步，長 16 步，問田地面積是多少？答案是：一畝。

\end{modernproblem}

}



\end{paracol}



\begin{paracol}{2}

\LeftLabel

{\CJKfamily{classical}\small\textit{草曰：}廣十五步，從十六步。以廣乘從：$15 \times 16 = 240$ 步。以一畝法二百四十步除之，得 \textbf{一畝}。

}



\switchcolumn

\RightLabel

{\CJKfamily{modern}\small\textit{演算過程：}寬 15 步，長 16 步。用寬乘長：$15 \times 16 = 240$ 平方步。用一畝的標準 240 平方步去除，得到 \textbf{一畝}。

}



\end{paracol}



\begin{paracol}{2}

\LeftLabel

{\CJKfamily{classical}\small

\begin{classicalproblem}[二]

廣十二步，從十四步，問田幾何？答曰：一百六十八步。

\end{classicalproblem}

}



\switchcolumn

\RightLabel

{\CJKfamily{modern}\small

\begin{modernproblem}[第二題]

有一塊田，寬 12 步，長 14 步，問面積是多少平方步？答案是：168 平方步。

\end{modernproblem}

}



\end{paracol}



\subsection{乘分 Fraction Multiplication}



\begin{paracol}{2}

\LeftLabel

{\CJKfamily{classical}\small

\textbf{術曰：}母互乘子，并以為實。母相乘為法。實如法而一。



\begin{classicalproblem}[乘分]

田廣七步、四分步之三，從十五步、九分步之五，問：為田幾何？答曰：一百二十步，九分步之五。

\end{classicalproblem}

}



\switchcolumn

\RightLabel

{\CJKfamily{modern}\small

\textbf{算法：}將各分母的分子交叉相乘，相加後作為被除數（實）。將各分母相乘作為除數（法）。用實除以法得到結果。



\begin{modernproblem}[分數乘法]

有一塊田，寬 $7\,\frac{3}{4}$ 步，長 $15\,\frac{5}{9}$ 步，問田地面積是多少？答案是：$120\,\frac{5}{9}$ 平方步。

\end{modernproblem}

}



\end{paracol}



\begin{paracol}{2}

\LeftLabel

{\CJKfamily{classical}\small\textit{草曰：}廣 $7\,\frac{3}{4} = \frac{31}{4}$ 步。從 $15\,\frac{5}{9} = \frac{140}{9}$ 步。以廣乘從：$\frac{31}{4} \times \frac{140}{9} = \frac{4340}{36} = 120\,\frac{5}{9}$ 平方步。

}



\switchcolumn

\RightLabel

{\CJKfamily{modern}\small\textit{演算過程：}寬 $7\,\frac{3}{4} = \frac{31}{4}$ 步。長 $15\,\frac{5}{9} = \frac{140}{9}$ 步。用寬乘長：$\frac{31}{4} \times \frac{140}{9} = \frac{4340}{36} = 120\,\frac{5}{9}$ 平方步。

}



\end{paracol}



\subsection{除分 Division with Fractions}



\begin{paracol}{2}

\LeftLabel

{\CJKfamily{classical}\small

\textbf{術曰：}以人數為法，錢數為實。有分者通之。實如法而一。



\begin{classicalproblem}[除分]

七人，均八錢三分錢之一，問人得幾何？答曰：人得一錢二十一分錢之四。

\end{classicalproblem}

}



\switchcolumn

\RightLabel

{\CJKfamily{modern}\small

\textbf{算法：}用人數作為除數（法），錢數作為被除數（實）。如果有分數就通分。用實除以法得到每人所得。



\begin{modernproblem}[分數除法]

七個人平分 $8\,\frac{1}{3}$ 錢，問每人得到多少？答案是：每人得到 $1\,\frac{4}{21}$ 錢。

\end{modernproblem}

}



\end{paracol}



\newpage



%% ========================================================================

%% CHAPTER 2: FIELD MEASUREMENT

%% ========================================================================

\section{第二章：方田 Field Measurement}



\begin{paracol}{2}

\LeftLabel

{\CJKfamily{classical}\small

方田章主要討論面積計算與分數運算。原書此章共 38 題，涵蓋各種田形及分數四則運算。

}



\switchcolumn

\RightLabel

{\CJKfamily{modern}\small

方田章主要討論面積的計算方法以及分數的四則運算。原書這一章共有 38 道題目，涵蓋了各種不同形狀的田地以及分數的加減乘除運算。

}



\end{paracol}



\subsection{約分 Reduction of Fractions}



\begin{paracol}{2}

\LeftLabel

{\CJKfamily{classical}\small

\textbf{約分術曰：}可半者半之。不可半者，副置分母子之數，以少減多，更相減損，求其等也。以等數約之。

}



\switchcolumn

\RightLabel

{\CJKfamily{modern}\small

\textbf{約分算法：}分子和分母如果都能被 2 整除，就同時除以 2。如果不能被 2 整除，就把分子和分母分別寫下來，用較大的數減去較小的數，反覆相減（輾轉相除），直到找到最大公約數。用這個最大公約數去約分子和分母。

}



\end{paracol}



\begin{paracol}{2}

\LeftLabel

\ClassicalText{

十八分之十二。問：約之得幾何？答曰：三分之二。

}



\switchcolumn

\RightLabel

\ModernText{

把 $\frac{12}{18}$ 約分。問：約分之後是多少？答案是：$\frac{2}{3}$。

}



\end{paracol}



\begin{paracol}{2}

\LeftLabel

{\CJKfamily{classical}\small\textit{草曰：}十二與十八，以少減多：$18 - 12 = 6$，$12 - 6 = 6$，$6 - 6 = 0$。等數為六。以六約之：$\frac{12 \div 6}{18 \div 6} = \frac{2}{3}$。

}



\switchcolumn

\RightLabel

{\CJKfamily{modern}\small\textit{演算過程：}12 和 18，用大數減小數：$18 - 12 = 6$，$12 - 6 = 6$，$6 - 6 = 0$。最大公約數是 6。用 6 去約：$\frac{12 \div 6}{18 \div 6} = \frac{2}{3}$。

}



\end{paracol}



\subsection{合分 Addition of Fractions}



\begin{paracol}{2}

\LeftLabel

{\CJKfamily{classical}\small

\textbf{合分術曰：}母互乘子，并以為實。母相乘為法。實如法而一。其母同者，直相從之。

}



\switchcolumn

\RightLabel

{\CJKfamily{modern}\small

\textbf{分數加法算法：}各分母交叉乘以各分子，將所得的積相加作為被除數（實）。各分母相乘作為除數（法）。用實除以法得到結果。如果分母相同，直接把分子相加即可。

}



\end{paracol}



\begin{paracol}{2}

\LeftLabel

\ClassicalText{

二分之一、三分之二、四分之三、五分之四，合之得幾何？答曰：得二，餘六十分之四十三。

}



\switchcolumn

\RightLabel

\ModernText{

計算 $\frac{1}{2} + \frac{2}{3} + \frac{3}{4} + \frac{4}{5}$，總和是多少？答案是：2 又 $\frac{43}{60}$。

}



\end{paracol}



\subsection{課分 Subtraction and Comparison}



\begin{paracol}{2}

\LeftLabel

{\CJKfamily{classical}\small

\textbf{課分術曰：}母互乘子，以少減多，餘為實。母相乘為法。實如法而一。即其相差之數。



\begin{classicalproblem}[課分]

八分之五，比二十五分之十六。問：孰多多幾何？答曰：二十五分之十六，多多二百分之三。

\end{classicalproblem}

}



\switchcolumn

\RightLabel

{\CJKfamily{modern}\small

\textbf{分數比較算法：}各分母交叉乘以各分子，用大數減去小數，差作為被除數（實）。分母相乘作為除數（法）。實除以法即得兩分數相差之數。



\begin{modernproblem}[分數比較]

比較 $\frac{5}{8}$ 和 $\frac{16}{25}$。問：哪個大？大多少？答案是：$\frac{16}{25}$ 更大，大出 $\frac{3}{200}$。

\end{modernproblem}

}



\end{paracol}



\subsection{平分 Equalizing Fractions}



\begin{paracol}{2}

\LeftLabel

{\CJKfamily{classical}\small

\textbf{平分術曰：}母互乘子，各列其數，并而為列實。母相乘為法。以列數乘法，各以列實減之，不足者益之，多者減之，各得其平。



\begin{classicalproblem}[平分]

二分之一、三分之二、四分之三，減多益少幾何？而平？答曰：減四分之三者四，減三分之二者一，益二分之一者五，各平於三十六分之二十三。

\end{classicalproblem}

}



\switchcolumn

\RightLabel

{\CJKfamily{modern}\small

\textbf{平分算法：}各分母交叉乘以各分子，分別列出所得數，相加作為列實。分母相乘作為法。用列數乘法，再從各列實中減去，不夠的補上，多餘的減去，使它們相等。



\begin{modernproblem}[分數平分]

有 $\frac{1}{2}$、$\frac{2}{3}$、$\frac{3}{4}$ 三個分數，從多的減去、補給少的，應該各減多少、補多少才能相等？答案是：從 $\frac{3}{4}$ 減去 $\frac{4}{36}$，從 $\frac{2}{3}$ 減去 $\frac{1}{36}$，給 $\frac{1}{2}$ 補上 $\frac{5}{36}$，三個都等於 $\frac{23}{36}$。

\end{modernproblem}

}



\end{paracol}



\newpage



%% ========================================================================

%% CHAPTER 3: GRAIN CONVERSION

%% ========================================================================

\section{第三章：粟米 Grain Conversion}



\begin{paracol}{2}

\LeftLabel

{\CJKfamily{classical}\small

粟米章論述各種糧食之間的比例兌換，共 46 題。其中有基本乘除 6 題，比例兌換 31 題，比率運算 9 題。



\textbf{經率術曰：}以所買率為法，所出錢為實。實如法而一。



\textbf{粟米換術曰：}以所有數乘所求率為實，以所有率為法，實如法而一。

}



\switchcolumn

\RightLabel

{\CJKfamily{modern}\small

粟米章討論各種糧食之間的比例兌換問題，共 46 道題目。其中包括基本乘除 6 題，比例兌換 31 題，比率運算 9 題。



\textbf{基本比率算法：}用購買的數量作為除數（法），付出的錢數作為被除數（實）。實除以法得到單價。



\textbf{糧食兌換算法：}用現有的數量乘以所求糧食的兌換率作為被除數（實），用現有糧食的兌換率作為除數（法）。實除以法得到結果。

}



\end{paracol}



\subsubsection*{標準兌換率 Standard Conversion Rates}



\begin{paracol}{2}

\LeftLabel

{\CJKfamily{classical}\small

\begin{itemize}[noitemsep,leftmargin=1.5em]

\item 粟：率五十

\item 糲米：率三十

\item 糳米：率二十七

\item 御米：率二十四

\item 麥、菽：率各四十五

\end{itemize}

}



\switchcolumn

\RightLabel

{\CJKfamily{modern}\small

\begin{itemize}[noitemsep,leftmargin=1.5em]

\item 粟（小米）：兌換率 50

\item 糲米（糙米）：兌換率 30

\item 糳米（精米）：兌換率 27

\item 御米（上供米）：兌換率 24

\item 麥、豆：兌換率各 45

\end{itemize}

}



\end{paracol}



\begin{paracol}{2}

\LeftLabel

\ClassicalText{

粟一斗。問為糲米幾何？答曰：六升。

}



\switchcolumn

\RightLabel

\ModernText{

有一斗粟，問能換成多少糲米（糙米）？答案是：六升。

}



\end{paracol}



\begin{paracol}{2}

\LeftLabel

{\CJKfamily{classical}\small\textit{草曰：}粟率五十，糲米率三十。$1 \times \frac{30}{50} = \frac{3}{5}$ 斗 = \textbf{六升}。

}



\switchcolumn

\RightLabel

{\CJKfamily{modern}\small\textit{演算過程：}粟的兌換率是 50，糲米的兌換率是 30。$1 \times \frac{30}{50} = \frac{3}{5}$ 斗 = \textbf{6 升}。

}



\end{paracol}



\begin{paracol}{2}

\LeftLabel

\ClassicalText{

粟四斗五升。問為糳米幾何？答曰：二斗一升五分升之三。

}



\switchcolumn

\RightLabel

\ModernText{

有粟四斗五升，問能換成多少糳米（精米）？答案是：二斗一升又五分之三升。

}



\end{paracol}



\subsection{絲與錢幣換算 Silk and Currency}



\begin{paracol}{2}

\LeftLabel

\ClassicalText{

絲一斤，價三百四十五，今有七兩一十二銖。問錢，答曰：一百六十一錢，三十二分錢之二十三。

}



\switchcolumn

\RightLabel

\ModernText{

絲一斤（16 兩 = 384 銖），價格 345 文錢。現在有絲七兩十二銖，問值多少錢？答案是：161 又 $\frac{23}{32}$ 文錢。

}



\end{paracol}



\newpage



%% ========================================================================

%% CHAPTER 4: PROPORTIONAL DISTRIBUTION

%% ========================================================================

\section{第四章：衰分 Proportional Distribution}



\begin{paracol}{2}

\LeftLabel

{\CJKfamily{classical}\small

衰分章論述按比例分配數量。原書此章共 20 題。



\textbf{衰分術曰：}各置列衰，副并為法。以所分乘未并者各自為實。實如法而一。不滿法者，以法命之。

}



\switchcolumn

\RightLabel

{\CJKfamily{modern}\small

衰分章討論按照給定比例來分配數量的問題。原書這一章共有 20 道題目。



\textbf{衰分算法：}列出各個分配的比率（衰），把它們相加作為除數（法）。用要分配的總數乘以每一個未相加前的比率，各自作為被除數（實）。實除以法得到每人所得。如果有餘數，就以法為分母來表示。

}



\end{paracol}



\subsection{五爵均鹿 Five Ranks Sharing Deer}



\begin{paracol}{2}

\LeftLabel

\ClassicalText{

大夫、不更、簪褭、上造、公士凡五人，以爵次高下均五鹿，問各幾何？答曰：大夫一鹿，三分鹿之二；不更一鹿，三分鹿之一；簪褭一鹿；上造鹿三分之二；公士鹿三分之一。

}



\switchcolumn

\RightLabel

\ModernText{

有大夫、不更、簪褭、上造、公士五個人，按照爵位的高低（比率為 5:4:3:2:1）來平分五隻鹿，問每人分到多少？答案是：大夫分到 $1\,\frac{2}{3}$ 隻鹿；不更分到 $1\,\frac{1}{3}$ 隻鹿；簪褭分到 1 隻鹿；上造分到 $\frac{2}{3}$ 隻鹿；公士分到 $\frac{1}{3}$ 隻鹿。

}



\end{paracol}



\subsection{牛馬羊食苗 The Cow, Horse, and Sheep}



\begin{paracol}{2}

\LeftLabel

\ClassicalText{

牛、馬、羊食人苗，苗主責之粟五斗，牛食馬之半，馬食羊之半，欲衰償之？答曰：牛二斗八升，七分升之四；馬一斗四升，七分升之二；羊七升，七分升之一。

}



\switchcolumn

\RightLabel

\ModernText{

一頭牛、一匹馬、一隻羊吃了別人家的禾苗，苗主要求賠償五斗粟。牛吃的量是馬的一半，馬吃的量是羊的一半，問按比率賠償各應出多少？答案是：牛出 $2\,\frac{4}{7}$ 斗（2 斗 8 又 $\frac{4}{7}$ 升）；馬出 $1\,\frac{2}{7}$ 斗（1 斗 4 又 $\frac{2}{7}$ 升）；羊出 $\frac{1}{7}$ 斗（7 又 $\frac{1}{7}$ 升）。

}



\end{paracol}



\subsection{女子善織 The Woman Skilled at Weaving}



\begin{paracol}{2}

\LeftLabel

\ClassicalText{

女子善織，日自倍，五日織五尺。問日織數。答曰：初日一寸，三十一分寸之十九；二日三寸，三十一分寸之七；三日六寸，三十一分寸之十四；四日一尺二寸，三十一分寸之二十八；五日二尺四寸，三十一分寸之二十五。

}



\switchcolumn

\RightLabel

\ModernText{

有個女子很會織布，每天織的布是前一天的一倍，五天共織了五尺。問每天各織了多少？答案是：第一天織了 $1\,\frac{19}{31}$ 寸；第二天織了 $3\,\frac{7}{31}$ 寸；第三天織了 $6\,\frac{14}{31}$ 寸；第四天織了 $12\,\frac{28}{31}$ 寸；第五天織了 $24\,\frac{25}{31}$ 寸。

}



\end{paracol}



\newpage



%% ========================================================================

%% CHAPTER 5: SHORT WIDTH / ROOT EXTRACTION

%% ========================================================================

\section{第五章：少廣 Short Width / Root Extraction}



\begin{paracol}{2}

\LeftLabel

{\CJKfamily{classical}\small

少廣章論述已知矩形面積及一邊求另一邊的問題，涉及分數除法及開方運算。原書此章共 24 題。



\textbf{少廣術曰：}置全步及分母子，各以其母除其子，内子，以乘全步，并之為法。以畝積步為實。實如法而一，得從。

}



\switchcolumn

\RightLabel

{\CJKfamily{modern}\small

少廣章討論已知矩形的面積和一條邊的長度，求另一條邊長度的問題，這涉及到分數除法和開方運算。原書這一章共有 24 道題目。



\textbf{少廣算法：}列出整數步和分數的分子分母，各自用分母去除分子，把餘數併入分子，再乘以整數步，相加後作為除數（法）。用一畝的面積（240 平方步）作為被除數（實）。實除以法，就得到長度。

}



\end{paracol}



\begin{paracol}{2}

\LeftLabel

\ClassicalText{

田一畝廣一步半，問從？答曰：一百六十步。

}



\switchcolumn

\RightLabel

\ModernText{

有一畝田，寬度為一步半（$1\,\frac{1}{2}$ 步），問長度是多少？答案是：160 步。

}



\end{paracol}



\begin{paracol}{2}

\LeftLabel

{\CJKfamily{classical}\small\textit{草曰：}廣 $1\,\frac{1}{2} = \frac{3}{2}$ 步。田積一畝 = 240 步。從 = $240 \div \frac{3}{2} = 240 \times \frac{2}{3} = $ \textbf{160} 步。

}



\switchcolumn

\RightLabel

{\CJKfamily{modern}\small\textit{演算過程：}寬 $1\,\frac{1}{2} = \frac{3}{2}$ 步。田地面積為 1 畝 = 240 平方步。長度 = $240 \div \frac{3}{2} = 240 \times \frac{2}{3} = $ \textbf{160} 步。

}



\end{paracol}



\begin{paracol}{2}

\LeftLabel

\ClassicalText{

田一畝廣一步半、三分步之一。問從？答曰：一百三十步一十一分步之一十。

}



\switchcolumn

\RightLabel

\ModernText{

有一畝田，寬度為一步半又三分之一步，問長度是多少？答案是：$130\,\frac{10}{11}$ 步。

}



\end{paracol}



\begin{paracol}{2}

\LeftLabel

\ClassicalText{

田一畝廣一步半、三分步之一、四分步之一。問從？答曰：一百一十五步五分步之一。

}



\switchcolumn

\RightLabel

\ModernText{

有一畝田，寬度為一步半又三分之一又四分之一步，問長度是多少？答案是：$115\,\frac{1}{5}$ 步。

}



\end{paracol}



\begin{paracol}{2}

\LeftLabel

\ClassicalText{

田一畝廣一步半、三分步之一、四分步之一、五分步之一。問從？答曰：一百五步、一百三十七分步之一十五。

}



\switchcolumn

\RightLabel

\ModernText{

有一畝田，寬度為一步半又三分之一又四分之一又五分之一步，問長度是多少？答案是：$105\,\frac{15}{137}$ 步。

}



\end{paracol}



\newpage



%% ========================================================================

%% YANG HUI'S TRIANGLE

%% ========================================================================

\section{開方作法本源圖 The Method Source for Square Root Extraction}



\begin{paracol}{2}

\LeftLabel

{\CJKfamily{classical}\small

楊輝所錄「開方作法本源圖」（即今所謂楊輝三角或帕斯卡三角形），為其最著名之貢獻。楊輝歸功於北宋數學家\textbf{賈憲}（約 1010--1070 年）。



楊輝釋之曰：



\fboxsep=4pt\colorbox{classicalbg}{\parbox{\linewidth}{\small

左袤乃積數，右袤乃隅算，中藏者皆廉，以廉乘商方，命實而除之。}}

}



\switchcolumn

\RightLabel

{\CJKfamily{modern}\small

楊輝記錄的「開方作法本源圖」（也就是現在所說的楊輝三角或帕斯卡三角形），是他最著名的貢獻之一。楊輝把這個成就歸功於北宋數學家\textbf{賈憲}（約 1010--1070 年）。



楊輝解釋說：



\fboxsep=4pt\colorbox{modernbg}{\parbox{\linewidth}{\small

左邊的斜線上的數字是「積數」（各項係數）；右邊斜線上的數字是「隅算」（最高次項係數）；中間隱藏的數字都是「廉」（中間各項係數）。用這些「廉」去乘以商的各次方，然後從被開方數中減去。}}

}



\end{paracol}



\vspace{1em}



\begin{center}

\begin{tikzpicture}[scale=1.0, every node/.style={font=\small}]

%% Yang Hui's Triangle

\foreach \n in {0,...,6}{

  \foreach \k in {0,...,\n}{

    \pgfmathsetmacro{\y}{-\n*0.9}

    \pgfmathsetmacro{\x}{\k*1.2-\n*0.6}

    \ifnum\n=0

      \ifnum\k=0\node[circle,draw,minimum size=0.7cm,fill=red!10] at (\x,\y) {1};\fi

    \fi

    \ifnum\n=1

      \ifnum\k=0\node[circle,draw,minimum size=0.7cm] at (\x,\y) {1};\fi

      \ifnum\k=1\node[circle,draw,minimum size=0.7cm] at (\x,\y) {1};\fi

    \fi

    \ifnum\n=2

      \ifnum\k=0\node[circle,draw,minimum size=0.7cm] at (\x,\y) {1};\fi

      \ifnum\k=1\node[circle,draw,minimum size=0.7cm,fill=blue!10] at (\x,\y) {2};\fi

      \ifnum\k=2\node[circle,draw,minimum size=0.7cm] at (\x,\y) {1};\fi

    \fi

    \ifnum\n=3

      \ifnum\k=0\node[circle,draw,minimum size=0.7cm] at (\x,\y) {1};\fi

      \ifnum\k=1\node[circle,draw,minimum size=0.7cm,fill=blue!10] at (\x,\y) {3};\fi

      \ifnum\k=2\node[circle,draw,minimum size=0.7cm,fill=blue!10] at (\x,\y) {3};\fi

      \ifnum\k=3\node[circle,draw,minimum size=0.7cm] at (\x,\y) {1};\fi

    \fi

    \ifnum\n=4

      \ifnum\k=0\node[circle,draw,minimum size=0.7cm] at (\x,\y) {1};\fi

      \ifnum\k=1\node[circle,draw,minimum size=0.7cm,fill=blue!10] at (\x,\y) {4};\fi

      \ifnum\k=2\node[circle,draw,minimum size=0.7cm,fill=green!10] at (\x,\y) {6};\fi

      \ifnum\k=3\node[circle,draw,minimum size=0.7cm,fill=blue!10] at (\x,\y) {4};\fi

      \ifnum\k=4\node[circle,draw,minimum size=0.7cm] at (\x,\y) {1};\fi

    \fi

    \ifnum\n=5

      \ifnum\k=0\node[circle,draw,minimum size=0.7cm] at (\x,\y) {1};\fi

      \ifnum\k=1\node[circle,draw,minimum size=0.7cm,fill=blue!10] at (\x,\y) {5};\fi

      \ifnum\k=2\node[circle,draw,minimum size=0.7cm,fill=green!10] at (\x,\y) {10};\fi

      \ifnum\k=3\node[circle,draw,minimum size=0.7cm,fill=green!10] at (\x,\y) {10};\fi

      \ifnum\k=4\node[circle,draw,minimum size=0.7cm,fill=blue!10] at (\x,\y) {5};\fi

      \ifnum\k=5\node[circle,draw,minimum size=0.7cm] at (\x,\y) {1};\fi

    \fi

    \ifnum\n=6

      \ifnum\k=0\node[circle,draw,minimum size=0.7cm] at (\x,\y) {1};\fi

      \ifnum\k=1\node[circle,draw,minimum size=0.7cm,fill=blue!10] at (\x,\y) {6};\fi

      \ifnum\k=2\node[circle,draw,minimum size=0.7cm,fill=green!10] at (\x,\y) {15};\fi

      \ifnum\k=3\node[circle,draw,minimum size=0.7cm,fill=yellow!20] at (\x,\y) {20};\fi

      \ifnum\k=4\node[circle,draw,minimum size=0.7cm,fill=green!10] at (\x,\y) {15};\fi

      \ifnum\k=5\node[circle,draw,minimum size=0.7cm,fill=blue!10] at (\x,\y) {6};\fi

      \ifnum\k=6\node[circle,draw,minimum size=0.7cm] at (\x,\y) {1};\fi

    \fi

  }

}



%% Labels

\node[left] at (-4.5,-0) {\small $n=0$};

\node[left] at (-4.5,-0.9) {\small $n=1$};

\node[left] at (-4.5,-1.8) {\small $n=2$};

\node[left] at (-4.5,-2.7) {\small $n=3$};

\node[left] at (-4.5,-3.6) {\small $n=4$};

\node[left] at (-4.5,-4.5) {\small $n=5$};

\node[left] at (-4.5,-5.4) {\small $n=6$};



%% Legend

\draw[->,thick,red!70] (3.5,0) -- (4.5,0) node[right] {\small 隅算（積數）};

\draw[->,thick,blue!60] (3.5,-0.6) -- (4.5,-0.6) node[right] {\small 一廉};

\draw[->,thick,green!60!black] (3.5,-1.2) -- (4.5,-1.2) node[right] {\small 二廉};

\draw[->,thick,orange!70] (3.5,-1.8) -- (4.5,-1.8) node[right] {\small 三廉};

\end{tikzpicture}

\end{center}



\vspace{0.3em}

\begin{center}

{\small 圖一：開方作法本源圖（楊輝三角）--- 二項式係數表（至六次冪）}

\end{center}



\subsection*{結構說明 Structure Explanation}



\begin{paracol}{2}

\LeftLabel

{\CJKfamily{classical}\small

\begin{itemize}[noitemsep,leftmargin=1.5em]

\item \textbf{左袤（積數）：}斜邊上的 1 為常數項係數

\item \textbf{右袤（隅算）：}另一斜邊上的 1 為最高次項係數

\item \textbf{中藏者（廉）：}中間各數為中間各項係數

\end{itemize}



每行即 $(x+a)^n$ 展開之係數：

\begin{align*}

(x+a)^0 &= 1 \\

(x+a)^1 &= x + a \\

(x+a)^2 &= x^2 + 2xa + a^2 \\

(x+a)^3 &= x^3 + 3x^2a + 3xa^2 + a^3

\end{align*}

}



\switchcolumn

\RightLabel

{\CJKfamily{modern}\small

\begin{itemize}[noitemsep,leftmargin=1.5em]

\item \textbf{左邊斜線（積數）：}斜線上的 1 是常數項的係數

\item \textbf{右邊斜線（隅算）：}另一條斜線上的 1 是最高次項的係數

\item \textbf{中間各數（廉）：}中間的數字是各中間項的係數

\end{itemize}



每一行就是 $(x+a)^n$ 展開式的係數：

\begin{align*}

(x+a)^0 &= 1 \\

(x+a)^1 &= x + a \\

(x+a)^2 &= x^2 + 2xa + a^2 \\

(x+a)^3 &= x^3 + 3x^2a + 3xa^2 + a^3

\end{align*}

}



\end{paracol}



\subsection*{增乘方求廉法 Jia Xian's Generation Method}



\begin{paracol}{2}

\LeftLabel

\ClassicalText{

增乘方求廉法：以一次乘二次，以二次乘三次，遞增一位，至求終。

}



\switchcolumn

\RightLabel

\ModernText{

增乘法求各項係數的方法：用一次項去乘二次項，用二次項去乘三次項，每次遞增一位，一直算到最後一項為止。

}



\end{paracol}



\begin{paracol}{2}

\LeftLabel

{\CJKfamily{classical}\small

此即今所謂「帕斯卡法則」：

\[

\binom{n}{k} = \binom{n-1}{k-1} + \binom{n-1}{k}

\]

每行中每數等於其左上方與右上方兩數之和。

}



\switchcolumn

\RightLabel

{\CJKfamily{modern}\small

這就是現在所說的「帕斯卡法則」：

\[

\binom{n}{k} = \binom{n-1}{k-1} + \binom{n-1}{k}

\]

三角形中每一行的每個數字都等於它左上方和右上方兩個數字之和。

}



\end{paracol}



\newpage



%% ========================================================================

%% CHAPTER 6: COMMERCIAL CALCULATIONS

%% ========================================================================

\section{第六章：商功 Commercial Calculations / Engineering}



\begin{paracol}{2}

\LeftLabel

{\CJKfamily{classical}\small

商功章論述土木工程中的體積計算，如挖土、築城、開渠等。原書此章共 28 題。



\textbf{土方換算率：}

\begin{itemize}[noitemsep,leftmargin=1.5em]

\item 穿地（挖方）：基準

\item 壤土（鬆土）：穿地四之五

\item 堅土（夯土）：穿地三之四

\end{itemize}

}



\switchcolumn

\RightLabel

{\CJKfamily{modern}\small

商功章討論土木工程建設中的體積計算問題，比如挖土、築城牆、開運河等。原書這一章共有 28 道題目。



\textbf{土方換算比率：}

\begin{itemize}[noitemsep,leftmargin=1.5em]

\item 挖方（自然狀態土）：基準

\item 鬆土（挖出的鬆散土）：與挖方之比為 5:4

\item 夯土（夯實後的土）：與挖方之比為 4:3

\end{itemize}

}



\end{paracol}



\begin{paracol}{2}

\LeftLabel

\ClassicalText{

穿地積一萬尺。問為堅、壤各幾何？答曰：為堅七千五百尺，為壤一萬二千五百尺。

}



\switchcolumn

\RightLabel

\ModernText{

挖出土方一萬立方尺。問能築成多少夯土和多少鬆土？答案是：夯土 7500 立方尺，鬆土 12500 立方尺。

}



\end{paracol}



\subsection*{城牆與溝渠體積 City Wall and Ditch Volumes}



\begin{paracol}{2}

\LeftLabel

{\CJKfamily{classical}\small

\textbf{術曰：}并上下廣而半之，以高（或深）乘之，又以袤乘之，得積尺。



\begin{classicalproblem}[城積]

城下廣四丈，上廣二丈，高五丈，袤一百二十六丈五尺。問為尺幾何？答曰：一百八十九萬七千五百尺。

\end{classicalproblem}

}



\switchcolumn

\RightLabel

{\CJKfamily{modern}\small

\textbf{算法：}將上底寬和下底寬相加取平均值，乘以高度（或深度），再乘以長度，得到體積（立方尺）。



\begin{modernproblem}[城牆體積]

一座城牆下底寬 4 丈，上底寬 2 丈，高 5 丈，長 126 丈 5 尺。問體積是多少立方尺？答案是：1,897,500 立方尺。

\end{modernproblem}

}



\end{paracol}



\subsection*{特殊立體 Special Solids}



\begin{paracol}{2}

\LeftLabel

{\CJKfamily{classical}\small

\textbf{塹堵術：}廣袤相乘，又以高乘之，二而一。



\textbf{陽馬術：}廣袤相乘，又以高乘之，三而一。



\textbf{鱉臑術：}廣袤相乘，又以高乘之，六而一。



\textbf{芻童術：}倍上袤加入下袤，以上廣乘之；倍下袤加入上袤，以下廣乘之；并之，以高乘之，六而一。

}



\switchcolumn

\RightLabel

{\CJKfamily{modern}\small

\textbf{塹堵（楔形體）算法：}寬乘以長，再乘以高，除以 2。



\textbf{陽馬（四稜錐）算法：}寬乘以長，再乘以高，除以 3。



\textbf{鱉臑（四面體）算法：}寬乘以長，再乘以高，除以 6。



\textbf{芻童（楔臺）算法：}上底長的兩倍加下底長，乘以上底寬；下底長的兩倍加上底長，乘以 下底寬；兩個結果相加，乘以高，除以 6。

}



\end{paracol}



\newpage



%% ========================================================================

%% CHAPTER 7: EQUITABLE DISTRIBUTION

%% ========================================================================

\section{第七章：均輸 Equitable Distribution}



\begin{paracol}{2}

\LeftLabel

{\CJKfamily{classical}\small

均輸章論述賦稅與運輸費用的公平分配問題。原書此章共 28 題。



\textbf{均輸術曰：}以道里之遠近、載輸之輕重、粟價之高下，各為率。以戶數乘之，為衰。并衰為法，以賦粟乘未并之衰為實。實如法而一。

}



\switchcolumn

\RightLabel

{\CJKfamily{modern}\small

均輸章討論稅收和運輸費用如何在不同地區之間公平分配的問題。原書這一章共有 28 道題目。



\textbf{均輸算法：}按照道路遠近、運輸負擔輕重、糧食價格高低，各自計算比率。用戶數乘以這個比率，得到分配的權重（衰）。把各權重相加作為除數（法），用總賦稅量乘以每個權重作為被除數（實）。實除以法得到每個地區應分擔的數額。

}



\end{paracol}



\subsection*{追及問題 The Pursuit Problem}



\begin{paracol}{2}

\LeftLabel

\ClassicalText{

甲發長安，五日日至齊，乙發齊，七日至長安。今乙先行二日。問：幾日相逢？答曰：二日，十二分日之一。

}



\switchcolumn

\RightLabel

\ModernText{

甲從長安出發，五天可到齊國。乙從齊國出發，七天可到長安。現在乙已經先出發兩天了。問：再過幾天兩人相遇？答案是：$2\,\frac{1}{12}$ 天。

}



\end{paracol}



\subsection*{瓜瓠相遇 The Gourd and Melon Problem}



\begin{paracol}{2}

\LeftLabel

\ClassicalText{

垣高九尺，瓜生其上，蔓日長七寸；瓠生其下，蔓日長一尺。問：幾日相逢？答曰：相逢於五日，十七分日之五。

}



\switchcolumn

\RightLabel

\ModernText{

一道牆高 9 尺，瓜從牆頂往下生長，每天長 7 寸；瓠從牆底往上生長，每天長 1 尺。問：幾天後兩條藤蔓相遇？答案是：$5\,\frac{5}{17}$ 天後相遇。

}



\end{paracol}



\newpage



%% ========================================================================

%% CHAPTER 8: EXCESS AND DEFICIT

%% ========================================================================

\section{第八章：盈不足 Excess and Deficit}



\begin{paracol}{2}

\LeftLabel

{\CJKfamily{classical}\small

盈不足章介紹「盈不足術」（雙假設法），為解線性方程與近似問題之強有力方法。原書此章共 20 題。



\textbf{盈不足術曰：}置所出率，盈、不足各居其下。令維乘所出率，以為實。并盈、不足為法。實如法而一。

}



\switchcolumn

\RightLabel

{\CJKfamily{modern}\small

盈不足章介紹「盈不足術」（也叫雙假設法），這是解線性方程和近似計算問題的一種非常有效的方法。原書這一章共有 20 道題目。



\textbf{盈不足算法：}列出兩次假設的數值，把多餘和不足分別寫在下面。交叉相乘（假設值乘以對方的盈或不足），相加作為被除數（實）。把盈和不足相加作為除數（法）。實除以法就得到正確答案。

}



\end{paracol}



\begin{paracol}{2}

\LeftLabel

\ClassicalText{

共買物，人出八文，盈三文，人出七文，不足四文。問數、物價各幾何？答曰：七人。物價五十三。

}



\switchcolumn

\RightLabel

\ModernText{

幾個人合買一件物品，每人出 8 文錢，就多出 3 文；每人出 7 文錢，就少 4 文。問有多少人？物品價格是多少？答案是：7 個人，物品價格 53 文。

}



\end{paracol}



\begin{paracol}{2}

\LeftLabel

\ClassicalText{

共買雞，各出九，盈十一，各出六，不足十六。問人數、雞價各幾何？答曰：九人，雞價七十。

}



\switchcolumn

\RightLabel

\ModernText{

幾個人合買一隻雞，每人出 9 文，多出 11 文；每人出 6 文，少 16 文。問有多少人？雞的價格是多少？答案是：9 個人，雞的價格 70 文。

}



\end{paracol}



\subsection*{盈朒適足 Exactly Sufficient}



\begin{paracol}{2}

\LeftLabel

\ClassicalText{

共買犬，人出五不足九十文，人出五十文適足。問人數、犬價各幾何？答曰：二人，犬價一百。

}



\switchcolumn

\RightLabel

\ModernText{

幾個人合買一條狗，每人出 5 文，少 90 文；每人出 50 文，剛剛好。問有多少人？狗的價格是多少？答案是：2 個人，狗的價格 100 文。

}



\end{paracol}



\newpage



%% ========================================================================

%% CHAPTER 9: SYSTEMS OF LINEAR EQUATIONS

%% ========================================================================

\section{第九章：方程 Systems of Linear Equations}



\begin{paracol}{2}

\LeftLabel

{\CJKfamily{classical}\small

方程章為中國數學最卓越篇章之一，用算籌排列之矩陣陣式解線性方程組，相當於現代高斯消元法。原書此章共 18 題。



\textbf{方程術曰：}置行數以相乘，各以所得，以少減多，反覆求之，直至一變量孤立，乃得其值。

}



\switchcolumn

\RightLabel

{\CJKfamily{modern}\small

方程章是中國數學中最出色的章節之一，它使用算籌排列成矩陣的形式來解線性方程組，這相當於現代的高斯消元法。原書這一章共有 18 道題目。



\textbf{方程算法：}把各方程的係數排列成行，互相交叉相乘，然後用大數減小數，反覆運算，直到一個變量被單獨分離出來，就可以求出它的值。

}



\end{paracol}



\subsection{三穀問題 Three Grains Problem}



\begin{paracol}{2}

\LeftLabel

\ClassicalText{

上禾三秉，中禾二秉，下禾一秉，共實三十九斗。上禾二秉，中禾三秉，下禾一秉，共實三十四斗。上禾一秉，中禾二秉，下禾三秉，共實二十六斗。問：上、中、下禾一秉各幾何？

}



\switchcolumn

\RightLabel

\ModernText{

3 捆上禾、2 捆中禾、1 捆下禾，共出糧 39 斗。2 捆上禾、3 捆中禾、1 捆下禾，共出糧 34 斗。1 捆上禾、2 捆中禾、3 捆下禾，共出糧 26 斗。問：上禾、中禾、下禾每捆各出糧多少？

}



\end{paracol}



\begin{paracol}{2}

\LeftLabel

{\CJKfamily{classical}\small

\textit{草曰：}列方程：

\begin{align*}

3x + 2y + z &= 39 \\

2x + 3y + z &= 34 \\

x + 2y + 3z &= 26

\end{align*}

以方程術（高斯消元法）解之：

\begin{enumerate}[noitemsep,label=\arabic*.,leftmargin=1.5em]

\item 以一、二式相減：$x - y = 5$

\item 以二、三式變換：$5x - 7y = 1$

\item 解之：$x = 9\,\frac{1}{4}$，$y = 4\,\frac{1}{4}$，$z = 2\,\frac{3}{4}$

\end{enumerate}

\textbf{答：}上禾 $9\,\frac{1}{4}$ 斗，中禾 $4\,\frac{1}{4}$ 斗，下禾 $2\,\frac{3}{4}$ 斗。

}



\switchcolumn

\RightLabel

{\CJKfamily{modern}\small

\textit{演算過程：}列出方程組：

\begin{align*}

3x + 2y + z &= 39 \\

2x + 3y + z &= 34 \\

x + 2y + 3z &= 26

\end{align*}

用方程術（高斯消元法）來解：

\begin{enumerate}[noitemsep,label=\arabic*.,leftmargin=1.5em]

\item 第一式和第二式相減得到：$x - y = 5$

\item 第二式和第三式變換後得到：$5x - 7y = 1$

\item 求解得到：$x = 9\,\frac{1}{4}$，$y = 4\,\frac{1}{4}$，$z = 2\,\frac{3}{4}$

\end{enumerate}

\textbf{答案：}上禾每捆 $9\,\frac{1}{4}$ 斗，中禾每捆 $4\,\frac{1}{4}$ 斗，下禾每捆 $2\,\frac{3}{4}$ 斗。

}



\end{paracol}



\subsection{牛羊值金 Cattle and Sheep}



\begin{paracol}{2}

\LeftLabel

\ClassicalText{

五牛二羊，直金十兩，二牛五羊，直金八兩。問：牛、羊價各幾何？答曰：一牛直金一兩二十一分兩之十三；一羊直金二十一分兩之二十。

}



\switchcolumn

\RightLabel

\ModernText{

5 頭牛和 2 隻羊，共值金 10 兩。2 頭牛和 5 隻羊，共值金 8 兩。問：牛和羊各值金多少？答案是：1 頭牛值 $1\,\frac{13}{21}$ 兩金，1 隻羊值 $\frac{20}{21}$ 兩金。

}



\end{paracol}



\newpage



%% ========================================================================

%% CHAPTER 10: RIGHT TRIANGLES

%% ========================================================================

\section{第十章：勾股 Right Triangles}



\begin{paracol}{2}

\LeftLabel

{\CJKfamily{classical}\small

勾股章運用勾股定理解實際問題。原書此章共 38 題（原 24 題 + 少廣 13 題 + 商功 1 題）。



\textbf{勾股定理求弦術：}勾股各自乘，并而開方除之，即弦。



\textbf{求勾術：}弦自乘，減股自乘，其餘開方除之，即勾。

}



\switchcolumn

\RightLabel

{\CJKfamily{modern}\small

勾股章運用勾股定理（畢達哥拉斯定理）來解決實際問題。原書這一章共有 38 道題目（原題 24 道 + 少廣章 13 道 + 商功章 1 道）。



\textbf{求弦長的算法：}勾和股各自平方，相加後開平方，就是弦的長度。



\textbf{求勾長的算法：}弦的平方減去股的平方，剩下的數開平方，就是勾的長度。

}



\end{paracol}



\begin{paracol}{2}

\LeftLabel

\ClassicalText{

勾八尺，股十五尺。問為弦幾何？答曰：十七尺。

}



\switchcolumn

\RightLabel

\ModernText{

直角三角形的勾長 8 尺，股長 15 尺，問弦長是多少？答案是：17 尺。

}



\end{paracol}



\subsection*{葛纏木 The Vine and the Tree}



\begin{paracol}{2}

\LeftLabel

\ClassicalText{

木長二丈，圍之三尺，葛生其下，纏木七周，上與木齊。問葛長幾何？答曰：二丈九尺。

}



\switchcolumn

\RightLabel

\ModernText{

一棵樹高 2 丈，周長 3 尺。藤蔓從樹底生長，繞樹 7 圈，剛好到達樹頂。問藤蔓有多長？答案是：2 丈 9 尺。

}



\end{paracol}



\subsection*{葭生池中 The Reed in the Pond}



\begin{paracol}{2}

\LeftLabel

\ClassicalText{

池方一丈，正中有葭，出水面一尺，引葭至岸，與水面適平。問水深？答曰：水深一丈二尺。

}



\switchcolumn

\RightLabel

\ModernText{

一個正方形池塘邊長 1 丈，正中央長著一根蘆葦，露出水面 1 尺。如果把蘆葦拉到池塘邊，蘆葦頂端剛好與水面齊平。問水深多少？答案是：水深 1 丈 2 尺。

}



\end{paracol}



\subsection*{折竹 The Broken Bamboo}



\begin{paracol}{2}

\LeftLabel

\ClassicalText{

竹高一丈，折梢拄地，去根三尺，問折處高幾何？答曰：四尺，二十分尺之十一。

}



\switchcolumn

\RightLabel

\ModernText{

一根竹子高 1 丈，折斷後竹梢拄在地上，離根部 3 尺。問折斷處離地面多高？答案是：$4\,\frac{11}{20}$ 尺。

}



\end{paracol}



\subsection*{勾中容圓 Circle Inscribed in a Right Triangle}



\begin{paracol}{2}

\LeftLabel

\ClassicalText{

勾八步，股十五步，問勾中容圓徑幾何？答曰：六步。

}



\switchcolumn

\RightLabel

\ModernText{

直角三角形的勾長 8 步，股長 15 步，問這個三角形內切圓的直徑是多少？答案是：6 步。

}



\end{paracol}



\newpage



%% ========================================================================

%% SQUARE ROOT EXTRACTION

%% ========================================================================

\section{開方術 Square and Cube Root Extraction}



\begin{paracol}{2}

\LeftLabel

{\CJKfamily{classical}\small

楊輝保存並解說了賈憲的開方方法，代表宋代計算數學之最高成就。



\textbf{釋鎖平方術：}

\begin{enumerate}[noitemsep,label=\arabic*.,leftmargin=1.5em]

\item 置積為實

\item 借一算（下法），步之超一位

\item 議所得，以一乘所借算為方法

\item 以除實

\item 倍方法為廉法

\item 復置下法，步之超一位，再以議所得乘之

\end{enumerate}

}



\switchcolumn

\RightLabel

{\CJKfamily{modern}\small

楊輝保存並詳細解說了賈憲的開方方法，這代表了宋代計算數學的最高成就。



\textbf{釋鎖開平方算法（類似現代的長除法開方）：}

\begin{enumerate}[noitemsep,label=\arabic*.,leftmargin=1.5em]

\item 把要求平方根的數作為被除數（實）

\item 放置一根算籌作為下法，每兩位向前移一步

\item 估算第一位數字，用這個數字乘以下法，得到方法

\item 從實中減去

\item 把方法加倍作為廉法

\item 再放置下法，每兩位移一位，再用估算的數字去乘

\end{enumerate}

}



\end{paracol}



\begin{paracol}{2}

\LeftLabel

\ClassicalText{

積五萬五千二百二十五步，問為方幾何？答曰：二百三十五步。

}



\switchcolumn

\RightLabel

\ModernText{

面積為 55225 平方步，問正方形的邊長是多少？答案是：235 步。

}



\end{paracol}



\begin{paracol}{2}

\LeftLabel

{\CJKfamily{classical}\small

\textit{草曰：}$\sqrt{55225} = $\textbf{235}。逐步演算：

\begin{enumerate}[noitemsep,label=\arabic*.,leftmargin=1.5em]

\item 分組：5 | 52 | 25

\item 首數：$2^2 = 4 \le 5$，首數為 2

\item 餘數：$5-4=1$，下推 52 得 152

\item 倍之：$2 \times 2 = 4$。求 $d$ 使 $(40+d) \times d \le 152$。$d = 3$，$43 \times 3 = 129$

\item 餘數：$152 - 129 = 23$，下推 25 得 2325

\item 倍之：$23 \times 2 = 46$。求 $d$ 使 $(460+d) \times d \le 2325$。$d=5$，$465 \times 5 = 2325$

\item 餘數為零。故 $\sqrt{55225} = 235$。

\end{enumerate}

}



\switchcolumn

\RightLabel

{\CJKfamily{modern}\small

\textit{演算過程：}$\sqrt{55225} = $\textbf{235}。逐步演算如下：

\begin{enumerate}[noitemsep,label=\arabic*.,leftmargin=1.5em]

\item 將數字從右往左每兩位分一組：5 | 52 | 25

\item 第一位：$2^2 = 4 \le 5$，所以第一位是 2

\item 餘數：$5 - 4 = 1$，把下一組 52 移下來得到 152

\item 把當前結果加倍：$2 \times 2 = 4$。找一個數字 $d$ 使得 $(40 + d) \times d \le 152$。$d = 3$ 時，$43 \times 3 = 129$

\item 餘數：$152 - 129 = 23$，把下一組 25 移下來得到 2325

\item 把當前結果 23 加倍：$23 \times 2 = 46$。找 $d$ 使得 $(460 + d) \times d \le 2325$。$d = 5$ 時，$465 \times 5 = 2325$

\item 餘數為零。所以 $\sqrt{55225} = 235$。

\end{enumerate}

}



\end{paracol}



\subsection*{開立方 Cube Root Extraction}



\begin{paracol}{2}

\LeftLabel

{\CJKfamily{classical}\small

\textbf{開立方術：}與開平方類似，惟有三層除數：方法（平方項）、廉法（一次項）、隅法（常數項）。



\begin{classicalproblem}[開立方]

積一百八十六萬八百六十七尺，問為立方幾何？答曰：一百二十三尺。

\end{classicalproblem}

}



\switchcolumn

\RightLabel

{\CJKfamily{modern}\small

\textbf{開立方算法：}與開平方類似，但是有三層除數：方法（二次項係數）、廉法（一次項係數）、隅法（常數項係數）。



\begin{modernproblem}[開立方]

體積為 1,860,867 立方尺，問立方體的邊長是多少？答案是：123 尺。

\end{modernproblem}

}



\end{paracol}



\begin{paracol}{2}

\LeftLabel

{\CJKfamily{classical}\small

\textit{驗算：}$123^3 = 123 \times 123 \times 123 = 15129 \times 123 = 1{,}860{,}867$。$\checkmark$

}



\switchcolumn

\RightLabel

{\CJKfamily{modern}\small

\textit{驗算：}$123^3 = 123 \times 123 \times 123 = 15129 \times 123 = 1{,}860{,}867$。驗證正確！$\checkmark$

}



\end{paracol}



\newpage



%% ========================================================================

%% CLASSIFICATION

%% ========================================================================

\section{篡類 Classification of Methods}



\begin{paracol}{2}

\LeftLabel

{\CJKfamily{classical}\small

楊輝將全書 246 題按數學方法重新分類，發現不同章節之題目往往使用相同之基本方法。



\begin{quote}\small

楊輝竊見九章舊本，作立題法遺闕。古序云：靖康以來，罕有舊本，間有存者，狃於末習，不循本意，以至真術淹廢，偽本滋典。

\end{quote}

}



\switchcolumn

\RightLabel

{\CJKfamily{modern}\small

楊輝將全書 246 道題目按照所使用的數學方法重新分類，發現不同章節的題目其實常常使用相同的基本方法。



\begin{quote}\small

我仔細查看了《九章》的舊版本，發現有些題目和算法已經遺失了。古序中說：自從靖康之變以來，舊版本已經很少見了，偶爾有幸存下來的，也因為後人的陋習而被篡改，不再遵循原本的意思，以至於真正的算法被埋沒廢棄了，而錯誤的版本卻流傳開來。

\end{quote}

}



\end{paracol}



\vspace{0.8em}



\begin{center}

{\small\textbf{方法交叉分類表}}

\end{center}

\vspace{0.3em}



\begin{center}

\begin{longtable}{|l|p{3.8cm}|c|}

\hline

\textbf{方法} & \textbf{應用章節} & \textbf{題數} \\

\hline

\endfirsthead

\hline

\textbf{方法} & \textbf{應用章節} & \textbf{題數} \\

\hline

\endhead

乘除 & 方田（38）、粟米（3）、經率（9） & 41 \\

\hline

經率 & 粟米（5）、盈不足（4） & 5 \\

\hline

合率 & 少廣（11）、均輸（8）、盈不足（1） & 20 \\

\hline

互換 & 粟米（38）、衰分（11）、均輸（11）、盈不足（3） & 63 \\

\hline

衰分 & 原章（9）、均輸（9） & 18 \\

\hline

疊積 & 商功（27） & 27 \\

\hline

盈不足 & 商功 + 原章（11） & 11 \\

\hline

方程 & 原章（18）、盈不足（1） & 19 \\

\hline

勾股 & 少廣（13）、商功（1）、原章（24） & 38 \\

\hline

\caption*{246 題按方法交叉分類}

\end{longtable}

\end{center}



\newpage



%% ========================================================================

%% CONCLUSION

%% ========================================================================

\section{結語 Conclusion}



\begin{paracol}{2}

\LeftLabel

{\CJKfamily{classical}\small

楊輝《詳解九章算法》之貢獻可概括為四端：



\begin{enumerate}[noitemsep,label=\arabic*.,leftmargin=1.5em]

\item \textbf{保存賈憲遺著：}楊輝引述使後人得知賈憲之二項式三角圖、開方方法及代數學貢獻——否則將全無所聞。



\item \textbf{系統編纂：}問、答、術、草之四段結構成為數學著述之典範。



\item \textbf{批判性學術：}按方法而非按主題分類題目，揭示了數學方法在不同領域中之統一性。



\item \textbf{教育創新：}楊輝明言為「不得其門而入」之學生編寫此書。

\end{enumerate}



「開方作法本源圖」（楊輝三角）已成為數學中最著名之圖形之一，見於全球教科書。其在西方被稱為「帕斯卡三角形」——以布萊茲·帕斯卡（1623--1662）命名，而帕斯卡生活於賈憲之後約六百年、楊輝之後約四百年——此乃數學知識在不同文化中獨立發展之明證，亦提醒我們正視中國數學家對現代數學基礎之貢獻。

}



\switchcolumn

\RightLabel

{\CJKfamily{modern}\small

楊輝《詳解九章算法》的貢獻可以概括為四個方面：



\begin{enumerate}[noitemsep,label=\arabic*.,leftmargin=1.5em]

\item \textbf{保存賈憲的遺著：}通過楊輝的引用，後人才知道賈憲的二項式三角圖、開方方法和代數學貢獻——否則這些成就將完全失傳。



\item \textbf{系統化的編纂方法：}問、答、術、草四段式結構成為數學著述的典範。



\item \textbf{批判性的學術研究：}按照數學方法而不是按照主題來分類題目，揭示了數學方法在不同應用領域中的統一性。



\item \textbf{教育上的創新：}楊輝明確表示，他是為那些「找不到入門途徑」的學生編寫這本書的。

\end{enumerate}



「開方作法本源圖」（楊輝三角）已經成為數學史上最著名的圖形之一，出現在世界各地的教科書中。它在西方被稱為「帕斯卡三角形」——以布萊茲·帕斯卡（1623--1662）命名，而帕斯卡生活在賈憲之後約六百年、楊輝之後約四百年——這是數學知識在不同文化中獨立發展的有力證明，也提醒我們要正視中國數學家對現代數學基礎的重要貢獻。

}



\end{paracol}



\vspace{2em}



\begin{center}

\longline

\end{center}



\vspace{1em}



\begin{center}

{\CJKfamily{classical}\textit{「數學乃眾學之基。」}}\\[0.5em]

--- 楊輝《詳解九章算法》

\end{center}



\end{document}







% ============================================================

% Source: translations/zh/chinese/yang-hui-xiangjie-jiuzhang-part2_classical-modern_bilingual.tex

% ============================================================



% ============================================================================

% Yang Hui's Xiangjie Jiuzhang Suanfa (Detailed Analysis of the Nine Chapters)

% Part II -- Bilingual Classical Chinese / Modern Chinese Edition

% 詳解九章算法（二）文言文↔白話文對照版

% ============================================================================

\documentclass[11pt,a4paper]{article}



% -------- Core Packages --------

\usepackage{fontspec}

\usepackage{polyglossia}

\usepackage{amsmath,amssymb,amsthm}

\usepackage{geometry}

\usepackage{graphicx}

\usepackage{xcolor}

\usepackage{titlesec}

\usepackage{fancyhdr}

\usepackage{enumitem}

\usepackage{array,booktabs,longtable,multirow}

\usepackage{hyperref}

\usepackage{paracol}

\usepackage{etoolbox}



% -------- Page Geometry --------

\geometry{

  paperwidth=210mm,paperheight=297mm,

  left=18mm,right=18mm,top=25mm,bottom=25mm,

  headheight=14pt,footskip=30pt

}



% -------- Language Setup --------

\setmainlanguage{english}



% -------- Font Configuration --------

\setmainfont{Times New Roman}

\setsansfont{Arial}

\setmonofont{Courier New}



% -------- CJK Support via xeCJK --------

\usepackage{xeCJK}

\setCJKmainfont{Microsoft YaHei}



% -------- Colors --------

\definecolor{titlecolor}{RGB}{45,55,72}

\definecolor{sectioncolor}{RGB}{55,65,81}

\definecolor{notecolor}{RGB}{120,53,15}

\definecolor{rulecolor}{RGB}{180,160,140}

\definecolor{classicalbg}{RGB}{252,250,245}

\definecolor{modernbg}{RGB}{245,250,252}

\definecolor{columnsepclr}{RGB}{200,190,180}



% -------- Section Formatting --------

\titleformat{\section}

  {\normalfont\Large\bfseries\color{sectioncolor}}

  {\thesection}{1em}{}

\titleformat{\subsection}

  {\normalfont\large\bfseries\color{sectioncolor}}

  {\thesubsection}{1em}{}

\titleformat{\subsubsection}

  {\normalfont\normalsize\bfseries\color{sectioncolor}}

  {\thesubsubsection}{1em}{}



% -------- Header/Footer --------

\pagestyle{fancy}

\fancyhf{}

\fancyhead[L]{\small\CJKfamily{\CJKrmdefault}詳解九章算法（二）\quad 文言文↔白話文對照版}

\fancyhead[R]{\small\thepage}

\renewcommand{\headrulewidth}{0.4pt}

\renewcommand{\footrulewidth}{0pt}



% -------- Custom Commands --------

\newcommand{\longline}{\noindent\rule{\textwidth}{0.4pt}}

\newcommand{\shortline}{\noindent\rule{0.5\textwidth}{0.4pt}\par}

\newcommand{\cnote}[1]{\textcolor{notecolor}{\textit{#1}}}



% Chinese text helper

\newcommand{\zhc}{}[1]{{\CJKfamily{\CJKrmdefault}#1}}  % classical

\newcommand{\zhm}[1]{{\CJKfamily{\CJKrmdefault}#1}}  % modern



% Bilingual environment: classical on left, modern on right

\newenvironment{bilingual}{%

  \begin{paracol}{2}%

  \setlength{\columnseprule}{0.4pt}%

  \defcolumnseparation{0.05\textwidth}%

}{%

  \end{paracol}%

}



% Classical column

\newcommand{\classicalcolumn}{\switchcolumn[0]}

\newcommand{\moderncolumn}{\switchcolumn[1]}



% Label commands

\newcommand{\classicalabel}{\noindent\textbf{\textcolor{notecolor}{\zhc{【文言文】}}\par\smallskip}}

\newcommand{\modernlabel}{\noindent\textbf{\textcolor{notecolor}{\zhm{【白話文】}}\par\smallskip}}



% Problem bilingual environment

\newenvironment{prob bilingual}[1]{%

  \begin{paracol}{2}%

  \setlength{\columnseprule}{0.4pt}%

}{%

  \end{paracol}%

  \medskip%

}



% Solution bilingual environment  

\newenvironment{solbilingual}{%

  \begin{paracol}{2}%

  \setlength{\columnseprule}{0.4pt}%

}{%

  \end{paracol}%

  \medskip%

}



% Bilingual passage: classical | modern

% Usage: \bilingualpara{classical text}{modern text}

\newcommand{\bilingualpara}[2]{%

  \begin{paracol}{2}%

  \setlength{\columnseprule}{0.4pt}%

  \switchcolumn[0]%

  \fcolorbox{rulecolor}{classicalbg}{%

    \parbox{\linewidth}{%

      \setlength{\parindent}{2em}%

      \linespread{1.3}\selectfont\zhc{#1}%

    }%

  }%

  \switchcolumn[1]%

  \fcolorbox{rulecolor}{modernbg}{%

    \parbox{\linewidth}{%

      \setlength{\parindent}{2em}%

      \linespread{1.3}\selectfont\zhm{#2}%

    }%

  }%

  \end{paracol}%

  \medskip%

}



% Bilingual math solution

\newenvironment{bilingmath}{%

  \begin{paracol}{2}%

  \setlength{\columnseprule}{0.4pt}%

}{%

  \end{paracol}%

  \medskip%

}



% -------- Hyperref Setup --------

\hypersetup{

  colorlinks=true,

  linkcolor=sectioncolor,

  citecolor=sectioncolor,

  urlcolor=sectioncolor,

  pdfencoding=auto,

  pdftitle={詳解九章算法（二）文言文白話文對照版},

  pdfauthor={楊輝},

  pdfsubject={中國古典數學雙語版}

}



% -------- TOC Depth --------

\setcounter{tocdepth}{3}

\setcounter{secnumdepth}{3}



% Column separation

\columnsep=20pt

\columnseprule=0.4pt



% ============================================================================

\begin{document}



% -------- Title Page --------

\begin{titlepage}

\begin{center}



\vspace*{1.5cm}



{\Huge\bfseries\color{titlecolor}

\zhc{詳解九章算法}\[0.6cm]}



{\LARGE\color{sectioncolor}

\zhc{（二）}\[0.8cm]}



{\Large\color{notecolor}

\zhc{文言文↔白話文對照版}\[1.5cm]}



\longline\\[0.8cm]



{\Large\zhc{宋\quad 楊輝\quad 撰}}\\[0.5cm]

{\large 約公元\,1238--1298\,年}\\[1.5cm]



\longline\\[1cm]



{\large\color{sectioncolor}

\textit{Xiangjie Jiuzhang Suanfa}\\[0.3cm]

Detailed Analysis of the Mathematical Methods in the Nine Chapters\\[0.3cm]

Part II --- Classical Chinese / Modern Chinese Bilingual Edition\\[1.5cm]}



\vfill



{\small\zhc{

本書為楊輝《詳解九章算法》第二部，包含均輸、盈不足、方程、勾股、商功等章，

及開方作法本源（楊輝三角）與九章纂類。}}



\vspace{0.5cm}



{\small

This document presents Part II of Yang Hui's \textit{Xiangjie Jiuzhang Suanfa},\\

covering chapters on Fair Levies (\zhc{均輸}), Excess and Deficit (\zhc{盈不足}),\\

Rectangular Arrays (\zhc{方程}), Right Triangles (\zhc{勾股}), and Construction (\zhc{商功}).}



\end{center}

\end{titlepage}



% -------- Table of Contents --------

\tableofcontents

\newpage



% ============================================================================

\section*{\zhc{序言}}

\addcontentsline{toc}{section}{序言}



\begin{paracol}{2}

\switchcolumn[0]

\noindent\textbf{\textcolor{notecolor}{\zhc{【文言文】}}}\par\smallskip

\zhc{

詳解九章算法者，宋錢塘楊輝謙光取古九章算經，抄錄經注原文於前，而以其所撰圖解釋注比類圖說驗辭各條之後者也。末附以九章纂類，則以當時俗傳算術為細而分析九章題問以類相從焉。



據自序，詳解八十題，乃九十七題，綿十二卷。全案九章為算經之首，諸家立術皆自此出。楊輝竊見九章舊本，作立題法，迫圓古序云：靖康以來，罕有舊本，間有存者，狃於末習，不循本意，以至具術淹廢，偽本滋興。

}



\switchcolumn[1]

\noindent\textbf{\textcolor{notecolor}{\zhm{【白話文】}}}\par\smallskip

\zhm{

《詳解九章算法》是南宋錢塘人楊輝（字謙光）取古代《九章算經》，抄錄經文和注釋於前，再將自己所撰寫的圖解、釋注、比類、圖說、驗辭等各條附於其後的著作。書末附有《九章纂類》，將當時民間流傳的算術與《九章》的題目按類別重新編排。



根據楊輝的自序，此書詳解八十題（實為九十七題），共十二卷。《九章算術》是算經之首，各家的算法都源出於此。楊輝見到《九章》舊本，寫道：自靖康年間（1126年）以來，罕見舊本，偶有存留的，也因後人習慣於末流做法，不遵循本意，以致正統算法湮沒失傳，偽本滋生泛濫。

}

\end{paracol}



\medskip



\begin{paracol}{2}

\switchcolumn[0]

\zhc{

全案九章為算經之首，諸家立術皆自此出。而世傳永樂大典及孔氏微波榭二本，均不免脫誤。鐘祥李尚書細草圖說多所改正，方可卒讀，而往往與此書暗合，則此書誠可貴也。且其圖靈中所列名目亦足與宋人算書互證。

}



\switchcolumn[1]

\zhm{

《九章算術》作為算經之首，各家創立的算法都源自於它。然而世間流傳的《永樂大典》本和孔氏微波榭本，都難免有脫漏和錯誤。鐘祥李尚書的細草圖說做了很多改正，才能順利閱讀，而這些改正往往與本書不謀而合，可見此書確實珍貴。而且書中圖錄所列的名目，也足以與宋代其他算書相互印證。

}

\end{paracol}



\vspace{1cm}

\longline



\newpage



% ============================================================================

\section{\zhc{均輸第六}}

\label{sec:junshu}



\begin{paracol}{2}

\switchcolumn[0]

\noindent\textbf{\textcolor{notecolor}{\zhc{【文言文】}}}\par\smallskip

\zhc{

均輸法曰：均，平也。輸，送也。以远近户口衰分，求所当输送之数。

}



\switchcolumn[1]

\noindent\textbf{\textcolor{notecolor}{\zhm{【白話文】}}}\par\smallskip

\zhm{

均輸法的含義：「均」是公平、平均的意思；「輸」是輸送、運送的意思。按照路程遠近和戶口多少來按比例分配，求出各戶應當輸送的數額。

}

\end{paracol}



\medskip



\subsection{\zhc{均輸粟——運輸費用公平分配}}



\begin{paracol}{2}

\switchcolumn[0]

\noindent\textbf{\textcolor{notecolor}{\zhc{【文言文·原題】}}}\par\smallskip

\zhc{

今有均輸粟：甲縣一萬戶，行道八日；乙縣九千五百戶，行道十日；丙縣一萬二千三百五十戶，行道十三日；丁縣一萬二千二百戶，行道二十日。各輸粟於輸所。四縣共車一萬乘，欲以行道日數與戶數相參，均賦輸之。問各縣粟、車各幾何？

}



\switchcolumn[1]

\noindent\textbf{\textcolor{notecolor}{\zhm{【白話文·譯題】}}}\par\smallskip

\zhm{

現有一個均輸粟米的問題：甲縣有一萬戶，路程需要走八天；乙縣有九千五百戶，路程需要走十天；丙縣有一萬二千三百五十戶，路程需要走十三天；丁縣有一萬二千二百戶，路程需要走二十天。各縣都要向輸送地點運送粟米。四縣共有車一萬輛，想要以路程天數和戶數相互配合，公平地分配運輸任務。問各縣應出粟米多少、用車多少？

}

\end{paracol}



\medskip



\begin{paracol}{2}

\switchcolumn[0]

\noindent\textbf{\textcolor{notecolor}{\zhc{【文言文·術曰】}}}\par\smallskip

\zhc{

術曰：令各縣戶數各如其行道日數而一，以為列衰。副并為法。以賦粟車數乘未并者，各自為實。實如法得一車。有分者，上下輩之。以二十五斛乘車數，即各縣粟數。

}



\switchcolumn[1]

\noindent\textbf{\textcolor{notecolor}{\zhm{【白話文·解法】}}}\par\smallskip

\zhm{

解法：令各縣的戶數分別除以各自的路程天數，作為分配的比率（列衰）。將這些比率相加，作為除數（法）。用要分配的車輛總數乘以各縣尚未相加的比率，各自作為被除數（實）。被除數除以除數，得到各縣應出的車數。若有分數，按上下調整。用每車二十五斛乘以各縣車數，即得各縣應出的粟米數。

}

\end{paracol}



\medskip



\begin{paracol}{2}

\switchcolumn[0]

\zhc{

又術：以戶數乘行道日數，各為衰。并衰為法，以賦車乘未并者為實，實如法得一車。

}



\switchcolumn[1]

\zhm{

另一種算法：用戶數乘以路程天數，各作為分配比率。將比率相加作為除數，用要分配的車輛總數乘以各縣尚未相加的比率作為被除數，被除數除以除數得到一車之數。

}

\end{paracol}



\medskip



\noindent\textbf{\textcolor{sectioncolor}{\zhc{詳解步驟 / 詳細計算步驟：}}}



\begin{paracol}{2}

\switchcolumn[0]

\zhc{

步驟一：置各縣戶數，各以行道日數除之，得列衰。



甲縣：$10000 \div 8 = 1250$\\

乙縣：$9500 \div 10 = 950$\\

丙縣：$12350 \div 13 = 950$\\

丁縣：$12200 \div 20 = 610$



步驟二：副并為法。



$1250 + 950 + 950 + 610 = 3760$



步驟三：以賦車萬乘乘未并者，各為實，實如法得一車。



甲：$10000 \times 1250 \div 3760 = 3325\frac{25}{47}$ 車\\

乙：$10000 \times 950 \div 3760 = 2526\frac{34}{47}$ 車\\

丙：$10000 \times 950 \div 3760 = 2526\frac{34}{47}$ 車\\

丁：$10000 \times 610 \div 3760 = 1622\frac{18}{47}$ 車



步驟四：以二十五斛乘車數，得各縣粟數。

}



\switchcolumn[1]

\zhm{

第一步：列出各縣戶數，分別除以各自的路程天數，得到分配比率。



甲縣：$10000 \div 8 = 1250$\\

乙縣：$9500 \div 10 = 950$\\

丙縣：$12350 \div 13 = 950$\\

丁縣：$12200 \div 20 = 610$



第二步：將各縣比率相加，作為總除數。



$1250 + 950 + 950 + 610 = 3760$



第三步：用車輛總數一萬輛乘以各縣的比率，再除以總除數，得到各縣應出的車數。



甲縣：$10000 \times 1250 \div 3760 = 3325\frac{25}{47}$ 輛\\

乙縣：$10000 \times 950 \div 3760 = 2526\frac{34}{47}$ 輛\\

丙縣：$10000 \times 950 \div 3760 = 2526\frac{34}{47}$ 輛\\

丁縣：$10000 \times 610 \div 3760 = 1622\frac{18}{47}$ 輛



第四步：用每車載量二十五斛乘以各縣車數，得到各縣應出的粟米數量。

}

\end{paracol}



\vspace{0.5cm}



\subsection{\zhc{善行者與拙行者——追及問題}}



\begin{paracol}{2}

\switchcolumn[0]

\noindent\textbf{\textcolor{notecolor}{\zhc{【文言文·原題】}}}\par\smallskip

\zhc{

今有善行者一百步，拙行者六十步。今拙者先行一百步，問善行者幾何步及之？

}



\switchcolumn[1]

\noindent\textbf{\textcolor{notecolor}{\zhm{【白話文·譯題】}}}\par\smallskip

\zhm{

現有一個走得快的人每走一百步的時間內，走得慢的人只走六十步。現在慢的人先走了一百步，問快的人要走多少步才能追上慢的人？

}

\end{paracol}



\medskip



\begin{paracol}{2}

\switchcolumn[0]

\noindent\textbf{\textcolor{notecolor}{\zhc{【文言文·術曰】}}}\par\smallskip

\zhc{

術曰：置善行者一百步，減拙行者六十步，餘四十步，以為法。又以善行者一百步乘拙者先行一百步，得萬步，以為實。實如法得二百五十步，即善行者及之步數也。

}



\switchcolumn[1]

\noindent\textbf{\textcolor{notecolor}{\zhm{【白話文·解法】}}}\par\smallskip

\zhm{

解法：用快行者的一百步減去慢行者的六十步，餘下四十步，作為除數（法）。再用快行者的一百步乘以慢行者先走的一百步，得到一萬步，作為被除數（實）。被除數除以除數得到二百五十步，就是快行者追上慢行者所需的步數。

}

\end{paracol}



\medskip



\begin{paracol}{2}

\switchcolumn[0]

\zhc{

答曰：二百五十步。



草曰：善行者一百步，拙行者六十步，相減餘四十步為法。以善行者一百步乘拙者先行一百步為實。實如法而一，得二百五十步。合問。

}



\switchcolumn[1]

\zhm{

答案：二百五十步。



算草：快行者一百步，慢行者六十步，相減餘四十步作為除數。用快行者的一百步乘以慢行者先走的一百步作為被除數。被除數除以除數，得到二百五十步。符合所問。

}

\end{paracol}



\vspace{0.5cm}



\subsection{\zhc{犬追兔——追及問題}}



\begin{paracol}{2}

\switchcolumn[0]

\noindent\textbf{\textcolor{notecolor}{\zhc{【文言文·原題】}}}\par\smallskip

\zhc{

今有兔先走一百步，犬追之二百五十步，不及三十步而止。問犬不止，復行幾何步及之？

}



\switchcolumn[1]

\noindent\textbf{\textcolor{notecolor}{\zhm{【白話文·譯題】}}}\par\smallskip

\zhm{

現有兔子先跑了一百步，狗追了二百五十步，還差三十步沒追上而停了下來。問如果狗不停止，繼續追還要跑多少步才能追上兔子？

}

\end{paracol}



\medskip



\begin{paracol}{2}

\switchcolumn[0]

\noindent\textbf{\textcolor{notecolor}{\zhc{【文言文·術曰】}}}\par\smallskip

\zhc{

術曰：置犬步二百五十，以兔先走一百步減之，餘一百五十步，以減不及三十步，餘一百二十步，以為法。以不及三十步乘犬追二百五十步，得七千五百步，為實。實如法得六十二步半。以并犬追二百五十步，得三百一十二步半，即犬及兔之步數也。

}



\switchcolumn[1]

\noindent\textbf{\textcolor{notecolor}{\zhm{【白話文·解法】}}}\par\smallskip

\zhm{

解法：取狗跑的二百五十步，減去兔子先跑的一百步，餘下一百五十步，再減去還差的三十步，餘下一百二十步，作為除數。用還差的三十步乘以狗追的二百五十步，得七千五百步，作為被除數。被除數除以除數得到六十二步半。再加上狗已追的二百五十步，得三百一十二步半，就是狗追上兔子所需的總步數。那麼還需要再跑的步數為：

}

\end{paracol}



\medskip



\begin{paracol}{2}

\switchcolumn[0]

\zhc{

又術：置兔先走一百步，以犬追二百五十步減之，餘一百五十步，以減不及三十步，餘一百二十步為法。以不及三十步乘犬追二百五十步為實，實如法得六十二步半，即復行步數。



答曰：復行六十二步半。

}



\switchcolumn[1]

\zhm{

另一種算法：取兔子先跑的一百步，用狗追的二百五十步減去它，餘一百五十步，再減去還差的三十步，餘一百二十步作為除數。用還差的三十步乘以狗追的二百五十步作為被除數，被除數除以除數得六十二步半，就是還需要再跑的步數。



答案：還需要再跑六十二步半。

}

\end{paracol}



\vspace{0.5cm}



\subsection{\zhc{空車與重車——往返問題}}



\begin{paracol}{2}

\switchcolumn[0]

\noindent\textbf{\textcolor{notecolor}{\zhc{【文言文·原題】}}}\par\smallskip

\zhc{

今有空車日行七十里，重車日行五十里。今載太倉粟輸上林，五日三返。問太倉上林相去幾何？

}



\switchcolumn[1]

\noindent\textbf{\textcolor{notecolor}{\zhm{【白話文·譯題】}}}\par\smallskip

\zhm{

現有空車每天行七十里，載重車每天行五十里。現在從太倉載粟米運到上林，五天往返三次。問太倉到上林的距離是多少？

}

\end{paracol}



\medskip



\begin{paracol}{2}

\switchcolumn[0]

\noindent\textbf{\textcolor{notecolor}{\zhc{【文言文·術曰】}}}\par\smallskip

\zhc{

術曰：并空車、重車行里，以返數乘之，各為實。以空車、重車相乘為法。實如法而一，得里數。以日數除之，即所求里也。



草曰：并空車七十里、重車五十里，得一百二十里。以三返乘之，得三百六十里，為實。以空車七十里乘重車五十里，得三千五百里，為法。實如法得一千八百七十五分里之三百六十，約之，得一十八里一千八百七十五分里之九百。以五日除之，得三又十二分七里之一，即太倉上林相去里數也。

}



\switchcolumn[1]

\noindent\textbf{\textcolor{notecolor}{\zhm{【白話文·解法】}}}\par\smallskip

\zhm{

解法：將空車和重車每天行的里數相加，用往返次數去乘，各作為被除數。用空車和重車的速度相乘作為除數。被除數除以除數得到里數，再用天數去除，即為所求距離。



算草：空車七十里加重車五十里，得一百二十里。用三次往返去乘，得三百六十里，作為被除數。空車七十里乘以重車五十里，得三千五百，作為除數。被除數除以除數得$\frac{360}{1875}$里，約分後得一十八又$\frac{900}{1875}$里。用五天去除，實際計算應為：將空重車速度取倒數相加得單程時間，乘以次數後用總天數除。



答案：太倉與上林相距四十八里一千一百七十六分里之十一。

}

\end{paracol}



\vspace{0.5cm}



\subsection{\zhc{凫雁相逢——相遇問題}}



\begin{paracol}{2}

\switchcolumn[0]

\noindent\textbf{\textcolor{notecolor}{\zhc{【文言文·原題】}}}\par\smallskip

\zhc{

今有凫起南海，七日至北海；雁起北海，九日至南海。今凫雁俱起，問何日相逢？

}



\switchcolumn[1]

\noindent\textbf{\textcolor{notecolor}{\zhm{【白話文·譯題】}}}\par\smallskip

\zhm{

現有野鴨從南海起飛，七天飛到北海；大雁從北海起飛，九天飛到南海。現在野鴨和大雁同時起飛，問幾天後相遇？

}

\end{paracol}



\medskip



\begin{paracol}{2}

\switchcolumn[0]

\noindent\textbf{\textcolor{notecolor}{\zhc{【文言文·術曰】}}}\par\smallskip

\zhc{

術曰：并日數為法，日數相乘為實，實如法得一日。



草曰：并七日、九日，得一十六日，為法。以七日乘九日，得六十三日，為實。實如法得三日十六分日之十五，即相逢日數也。

}



\switchcolumn[1]

\noindent\textbf{\textcolor{notecolor}{\zhm{【白話文·解法】}}}\par\smallskip

\zhm{

解法：將兩個天數相加作為除數，兩個天數相乘作為被除數，被除數除以除數得到相遇天數。



算草：七天加九天，得十六天，作為除數。七天乘以九天，得六十三，作為被除數。被除數除以除數得$3\frac{15}{16}$天，即為相遇的天數。



驗證：野鴨每天飛全程的$\frac{1}{7}$，大雁每天飛全程的$\frac{1}{9}$。兩鳥每天共飛$\frac{1}{7}+\frac{1}{9}=\frac{16}{63}$。所以相遇需要$\frac{63}{16}=3\frac{15}{16}$天。

}

\end{paracol}



\newpage



% ============================================================================

\section{\zhc{盈不足第七}}

\label{sec:yingbuzu}



\begin{paracol}{2}

\switchcolumn[0]

\noindent\textbf{\textcolor{notecolor}{\zhc{【文言文】}}}\par\smallskip

\zhc{

盈不足術曰：置所出率，盈、不足各居其下。令維乘所出率，并以為實。并盈、不足為法。實如法而一。有分者，通之。盈不足相與同其買物者，置所出率，以少減多，餘，以約法、實。實為物價，法為人數。

}



\switchcolumn[1]

\noindent\textbf{\textcolor{notecolor}{\zhm{【白話文】}}}\par\smallskip

\zhm{

盈不足法的計算規則：列出每人所出的錢數（出率），將盈餘或不足的數額分別寫在下面。讓盈餘和不足的數額交叉互乘各自所出的錢數，將乘積相加作為被除數（實）。將盈餘和不足的數額相加作為除數（法）。被除數除以除數得到每人應出的錢數。若有分數，就通分。當盈餘和不足對應的是購買同一件物品時，將所出的錢數以大減小，用餘數去約簡法和實。約簡後的實就是物品的價格，法就是人數。

}

\end{paracol}



\medskip



\subsection{\zhc{共買雞——盈不足問題}}



\begin{paracol}{2}

\switchcolumn[0]

\noindent\textbf{\textcolor{notecolor}{\zhc{【文言文·原題】}}}\par\smallskip

\zhc{

今有共買雞，人出九，盈十一；人出六，不足十六。問人數、雞價各幾何？

}



\switchcolumn[1]

\noindent\textbf{\textcolor{notecolor}{\zhm{【白話文·譯題】}}}\par\smallskip

\zhm{

現有眾人合買一隻雞，如果每人出九文錢，會多出十一文；如果每人出六文錢，會少十六文。問人數和雞的價格各是多少？

}

\end{paracol}



\medskip



\begin{paracol}{2}

\switchcolumn[0]

\noindent\textbf{\textcolor{notecolor}{\zhc{【文言文·術曰】}}}\par\smallskip

\zhc{

術曰：置人出九、人出六，各以其下盈、不足維乘。以人出九乘不足十六，得一百四十四；以人出六乘盈十一，得六十六。并之，得二百一十，為實。并盈十一、不足十六，得二十七，為法。實如法，得人數九。置人數九，以人出九乘之，得八十一，減盈十一，餘七十，即雞價。或以人出六乘人數九，得五十四，加不足十六，亦得七十，即雞價。

}



\switchcolumn[1]

\noindent\textbf{\textcolor{notecolor}{\zhm{【白話文·解法】}}}\par\smallskip

\zhm{

解法：列出每人出九文和每人出六文，用各自下面的盈餘和不足交叉互乘。用每人出九文乘以不足十六，得一百四十四；用每人出六文乘以盈餘十一，得六十六。相加得二百一十，作為被除數（實）。將盈餘十一和不足十六相加，得二十七，作為除數（法）。被除數除以除數，得到人數為九人。用人數九乘以每人出九文，得八十一，減去盈餘十一，餘下七十，就是雞的價格。或者用人數九乘以每人出六文，得五十四，加上不足十六，也得到七十，同樣是雞的價格。

}

\end{paracol}



\medskip



\begin{paracol}{2}

\switchcolumn[0]

\zhc{

答曰：九人，雞價七十。



楊輝詳解曰：此題以盈不足術求之。人出九則盈十一，是雞價少於$9n$十一也；人出六則不足十六，是雞價多於$6n$十六也。故以兩出率之差為法，盈不足之和為實求人數；以人數乘出率減盈或加不足即得雞價。

}



\switchcolumn[1]

\zhm{

答案：九人，雞價七十文。



楊輝詳解：此題用盈不足術求解。每人出九文則盈餘十一文，說明雞價比$9n$少十一；每人出六文則不足十六文，說明雞價比$6n$多十六。所以用兩個出率之差作為除數，盈餘和不足之和作為被除數來求人數；用人數乘以出率再減去盈餘或加上不足就得到雞價。

}

\end{paracol}



\medskip



\noindent\textbf{\textcolor{sectioncolor}{\zhc{現代公式表示：}}}

\[

\text{人數} = \frac{11 + 16}{9 - 6} = \frac{27}{3} = 9, \quad

\text{雞價} = 9 \times 9 - 11 = 81 - 11 = 70 \text{（文）}

\]



\vspace{0.5cm}



\subsection{\zhc{共買金——兩盈問題}}



\begin{paracol}{2}

\switchcolumn[0]

\noindent\textbf{\textcolor{notecolor}{\zhc{【文言文·原題】}}}\par\smallskip

\zhc{

今有共買金，人出四百，盈三千四百；人出三百，盈一百。問人數、金價各幾何？

}



\switchcolumn[1]

\noindent\textbf{\textcolor{notecolor}{\zhm{【白話文·譯題】}}}\par\smallskip

\zhm{

現有眾人合買金子，每人出四百文，盈餘三千四百文；每人出三百文，盈餘一百文。問人數和金價各是多少？

}

\end{paracol}



\medskip



\begin{paracol}{2}

\switchcolumn[0]

\noindent\textbf{\textcolor{notecolor}{\zhc{【文言文·術曰】}}}\par\smallskip

\zhc{

術曰：置所出率，以少減多，餘，以約法、實。以兩盈相減為實，以兩出率相減為法。實如法而一，得人數。以人數乘出率，減盈，得金價。



草曰：以人出四百減人出三百，餘一百，為法。以盈三千四百減盈一百，餘三千三百，為實。實如法，得三十三人。以三十三人乘四百，得一萬三千二百，減盈三千四百，餘九千八百，即金價也。

}



\switchcolumn[1]

\noindent\textbf{\textcolor{notecolor}{\zhm{【白話文·解法】}}}\par\smallskip

\zhm{

解法：列出每人所出的錢數，用少的減去多的，餘下的用來約簡法和實。用兩個盈餘數相減作為被除數，用兩個出率相減作為除數。被除數除以除數得到人數。用人數乘以出率，減去盈餘，得到金價。



算草：用每人出四百減去每人出三百，餘一百，作為除數。用盈餘三千四百減去盈餘一百，餘三千三百，作為被除數。被除數除以除數，得三十三人。用三十三人乘以四百，得一萬三千二百，減去盈餘三千四百，餘下九千八百，就是金價。

}

\end{paracol}



\medskip



\noindent\textbf{\textcolor{sectioncolor}{\zhc{現代公式表示：}}}

\[

\text{人數} = \frac{3400 - 100}{400 - 300} = \frac{3300}{100} = 33, \quad

\text{金價} = 33 \times 400 - 3400 = 9800 \text{（文）}

\]



\vspace{0.5cm}



\subsection{\zhc{善田與惡田——假設問題}}



\begin{paracol}{2}

\switchcolumn[0]

\noindent\textbf{\textcolor{notecolor}{\zhc{【文言文·原題】}}}\par\smallskip

\zhc{

今有善田一畝價三百，惡田七畝價五百。今并買一頃，價錢一萬。問善田、惡田各幾何？

}



\switchcolumn[1]

\noindent\textbf{\textcolor{notecolor}{\zhm{【白話文·譯題】}}}\par\smallskip

\zhm{

現有好田一畝價值三百文，壞田七畝價值五百文。現在一起買了一百畝（一頃），共花價錢一萬文。問好田和壞田各買了多少畝？

}

\end{paracol}



\medskip



\begin{paracol}{2}

\switchcolumn[0]

\noindent\textbf{\textcolor{notecolor}{\zhc{【文言文·術曰】}}}\par\smallskip

\zhc{

術曰：假令善田二十畝，當價錢六千；惡田八十畝，價錢五千七百一十四又七分之二，多於實價一千七百一十四又七分之二。令之善田十畝，惡田九十畝，多於實價五百七十一又七分之三。以盈不足術求之。

}



\switchcolumn[1]

\noindent\textbf{\textcolor{notecolor}{\zhm{【白話文·解法】}}}\par\smallskip

\zhm{

解法：假設好田二十畝，應值六千文；壞田八十畝，應值五千七百一十四又七分之二文，比實際價格多一千七百一十四又七分之二文。再假設好田十畝，壞田九十畝，比實際價格多五百七十一又七分之三文。用盈不足術求解。



具體計算：用盈不足交叉計算，得善田十二畝半，惡田八十七畝半。

}

\end{paracol}



\medskip



\noindent\textbf{\textcolor{sectioncolor}{\zhc{現代方程解法：}}}

\begin{align*}

x + y &= 100 \\

300x + \frac{500}{7}y &= 10000

\end{align*}

解得：\zhc{善田}$x = 12.5$ 畝，\zhc{惡田}$y = 87.5$ 畝。



\vspace{0.5cm}



\subsection{\zhc{大器小器——容量問題}}



\begin{paracol}{2}

\switchcolumn[0]

\noindent\textbf{\textcolor{notecolor}{\zhc{【文言文·原題】}}}\par\smallskip

\zhc{

今有大器五、小器一，容三斛；大器一、小器五，容二斛。問大、小器各容幾何？

}



\switchcolumn[1]

\noindent\textbf{\textcolor{notecolor}{\zhm{【白話文·譯題】}}}\par\smallskip

\zhm{

現有五個大容器和一個小容器，共能容納三斛；一個大容器和五個小容器，共能容納二斛。問大容器和小容器各能容納多少？

}

\end{paracol}



\medskip



\begin{paracol}{2}

\switchcolumn[0]

\noindent\textbf{\textcolor{notecolor}{\zhc{【文言文·術曰】}}}\par\smallskip

\zhc{

術曰：以盈不足術求之。假令大器容半斛，小器容五分之一斛，以方程術入之。以大器五乘半斛，得二斛半；以小器一乘五分之一，得五分之一斛。并之，得二斛又五分之四，不足於三斛五分之一。又以大器一乘半斛，得半斛；小器五乘五分之一，得一斛。并之，得一斛半，不足於二斛半。以盈不足維乘，并以為實，并盈不足為法，實如法而一。

}



\switchcolumn[1]

\noindent\textbf{\textcolor{notecolor}{\zhm{【白話文·解法】}}}\par\smallskip

\zhm{

解法：用盈不足術求解。先假設大器容半斛，小器容五分之一斛，再用方程術驗證。五個大器乘以半斛，得二斛半；一個小器乘以五分之一，得五分之一斛。相加得二斛又五分之四，比三斛少五分之一。又一大器乘以半斛得半斛，五小器乘以五分之一得一斛，相加得一斛半，比二斛少半斛。用盈不足交叉互乘，相加作為被除數，盈不足相加作為除數，被除數除以除數即得。

}

\end{paracol}



\medskip



\noindent\textbf{\textcolor{sectioncolor}{\zhc{現代方程解法：}}}

\begin{align*}

5L + S &= 3 \\

L + 5S &= 2

\end{align*}

由第一式減第二式：$4L - 4S = 1$，即$L - S = \frac{1}{4}$。\\

兩式相加：$6L + 6S = 5$，即$L + S = \frac{5}{6}$。\\

解得：$L = \frac{13}{24}$ 斛，$S = \frac{7}{24}$ 斛。



\vspace{0.5cm}



\subsection{\zhc{良馬駑馬——複雜追及}}



\begin{paracol}{2}

\switchcolumn[0]

\noindent\textbf{\textcolor{notecolor}{\zhc{【文言文·原題】}}}\par\smallskip

\zhc{

今有良馬與駑馬發長安至齊。齊去長安三千里。良馬初日行一百九十三里，日增十三里；駑馬初日行九十七里，日減半里。良馬先至齊，復還迎駑馬。問幾何日相逢及各行幾何？

}



\switchcolumn[1]

\noindent\textbf{\textcolor{notecolor}{\zhm{【白話文·譯題】}}}\par\smallskip

\zhm{

現有好馬和劣馬從長安出發到齊國。齊國距離長安三千里。好馬第一天走一百九十三里，之後每天增加十三里；劣馬第一天走九十七里，之後每天減少半里。好馬先到達齊國後，再返回迎接劣馬。問多少天後兩馬相遇，以及各自行了多少里？

}

\end{paracol}



\medskip



\begin{paracol}{2}

\switchcolumn[0]

\noindent\textbf{\textcolor{notecolor}{\zhc{【文言文·術曰】}}}\par\smallskip

\zhc{

術曰：以盈不足術求之。假令十五日，不足三百三十七里半。令之十六日，盈一百四十里。以十五乘不足三百三十七里半，得五千六十二半；十六乘盈一百四十里，得二千二百四十里。并之，得七千三百二半，為實。并盈不足，得四百七十七半，為法。實如法，即日數也。

}



\switchcolumn[1]

\noindent\textbf{\textcolor{notecolor}{\zhm{【白話文·解法】}}}\par\smallskip

\zhm{

解法：用盈不足術求解。先假設十五天，比實際少三百三十七里半。再假設十六天，比實際多一百四十里。用十五乘以不足三百三十七里半，得五千零六十二點五；用十六乘以盈餘一百四十里，得二千二百四十。相加得七千三百零二點五，作為被除數。將盈餘和不足相加，得四百七十七點五，作為除數。被除數除以除數，即為相遇天數。



楊輝曰：此題以良馬、駑馬行程為等差數列，良馬日增、駑馬日減，先求良馬到齊之日，再求返程與駑馬相遇之日，乃盈不足術之精妙運用也。

}

\end{paracol}



\newpage



% ============================================================================

\section{\zhc{方程第八}}

\label{sec:fangcheng}



\begin{paracol}{2}

\switchcolumn[0]

\noindent\textbf{\textcolor{notecolor}{\zhc{【文言文】}}}\par\smallskip

\zhc{

方程術曰：置上禾三秉、中禾二秉、下禾一秉，實三十九斗，於右方。中、左禾列如右方。以右行上禾遍乘中行，而以直除。又乘其次，亦以直除。然以中行中禾不盡者遍乘左行，而以直除。左方下禾不盡者，上為法，下為實。實即下禾之實。求中禾，以法乘中行下實，而除下禾之實。餘如中禾秉數而一，即中禾之實。求上禾亦以法乘右行下實，而除下禾、中禾之實。餘如上禾秉數而一，即上禾之實。實皆如法，各得一斗。

}



\switchcolumn[1]

\noindent\textbf{\textcolor{notecolor}{\zhm{【白話文】}}}\par\smallskip

\zhm{

方程術的計算規則：將上禾三秉、中禾二秉、下禾一秉，實三十九斗，放在右邊（右行）。中行的禾數排列和左行的禾數排列也像右行這樣排好。用右行的上禾數遍乘中行的各數，然後直接相減（直除）。又用右行遍乘左行，也直接相減。然後用中行中禾的餘數遍乘左行，再直除。左行下禾的餘數，上面的作為除數，下面的作為被除數。被除數除以除數就得到下禾每秉的容量。求中禾時，用除數乘中行下實，再減去下禾的實。餘下的按中禾的秉數除，就得到中禾每秉的容量。求上禾也用除數乘右行下實，再減去下禾和中禾的實。餘下的按上禾的秉數除，就得到上禾每秉的容量。各實都按法除，就得到每斗之數。

}

\end{paracol}



\vspace{0.5cm}



\subsection{\zhc{三禾求實——三元一次方程組}}



\begin{paracol}{2}

\switchcolumn[0]

\noindent\textbf{\textcolor{notecolor}{\zhc{【文言文·原題】}}}\par\smallskip

\zhc{

今有上禾三秉、中禾二秉、下禾一秉，實三十九斗；上禾二秉、中禾三秉、下禾一秉，實三十四斗；上禾一秉、中禾二秉、下禾三秉，實二十六斗。問上、中、下禾實一秉各幾何？

}



\switchcolumn[1]

\noindent\textbf{\textcolor{notecolor}{\zhm{【白話文·譯題】}}}\par\smallskip

\zhm{

現有上禾三秉、中禾二秉、下禾一秉，共打糧三十九斗；上禾二秉、中禾三秉、下禾一秉，共打糧三十四斗；上禾一秉、中禾二秉、下禾三秉，共打糧二十六斗。問上禾、中禾、下禾每秉各打糧多少斗？

}

\end{paracol}



\medskip



\begin{paracol}{2}

\switchcolumn[0]

\noindent\textbf{\textcolor{notecolor}{\zhc{【文言文·方程布列】}}}\par\smallskip

\zhc{

右行：上禾三、中禾二、下禾一、實三十九\\

中行：上禾二、中禾三、下禾一、實三十四\\

左行：上禾一、中禾二、下禾三、實二十六



術曰：以右行上禾三遍乘中行，得六、九、三、一百零二。以右行減之，餘零、五、一、二十五。又以右行遍乘左行，得三、六、九、七十八。以右行減之，餘零、四、八、三十九。

}



\switchcolumn[1]

\noindent\textbf{\textcolor{notecolor}{\zhm{【白話文·方程布列與消元】}}}\par\smallskip

\zhm{

右行：上禾3、中禾2、下禾1、實39\\

中行：上禾2、中禾3、下禾1、實34\\

左行：上禾1、中禾2、下禾3、實26



計算步驟：用右行的上禾3遍乘中行各數，得6、9、3、102。減去右行的3、2、1、39各兩倍，餘0、5、1、25。又用右行遍乘左行各數，得3、6、9、78。減去右行的一倍，餘0、4、8、39。

}

\end{paracol}



\medskip



\begin{paracol}{2}

\switchcolumn[0]

\zhc{

次以中行餘五乘左行餘四，得二十、四十、一百九十五。以左行餘四乘中行餘五，得二十、四、一百。以中行減左行，餘零、三十六、九十九。以三十六為法，九十九為實。實如法，得下禾一秉二斗四分斗之三。



求中禾：以法三十六乘中行下實二十五，得九百。以下禾實九十九乘中行下禾一，得九十九。以減九百，餘八百零一。以中禾五除之，得中禾一秉四斗四分斗之一。



求上禾：以法三十六乘右行下實三十九，得一千四百零四。以下禾實九十九乘右行下禾一，得九十九。以中禾實八百零一乘右行中禾二……以右行上禾三除之，得上禾一秉九斗四分斗之一。



答曰：上禾一秉九斗四分斗之一，中禾一秉四斗四分斗之一，下禾一秉二斗四分斗之三。

}



\switchcolumn[1]

\zhm{

下一步：用中行餘數5乘左行餘數4，得20、40、195。用左行餘數4乘中行餘數5，得20、4、100。左行減中行，餘0、36、99。以36為除數，99為被除數。被除數除以除數，得下禾每秉$2\frac{3}{4}$斗。



求中禾：用除數36乘中行下實25，得900。以下禾實99乘中行下禾數1，得99。從900中減去，餘801。用中禾數5去除，得中禾每秉$4\frac{1}{4}$斗。



求上禾：用除數36乘右行下實39，得1404。以下禾實99乘右行下禾數1，得99。以中禾實801乘右行中禾數2……用右行上禾數3去除，得上禾每秉$9\frac{1}{4}$斗。



答案：上禾每秉$9\frac{1}{4}$斗，中禾每秉$4\frac{1}{4}$斗，下禾每秉$2\frac{3}{4}$斗。

}

\end{paracol}



\medskip



\noindent\textbf{\textcolor{sectioncolor}{\zhc{現代矩陣解法：}}}



現代記法設上、中、下禾每秉分別為$x, y, z$斗：

\begin{align*}

3x + 2y + z &= 39 \\

2x + 3y + z &= 34 \\

x + 2y + 3z &= 26

\end{align*}



高斯消元法：

\begin{enumerate}[leftmargin=2em]

  \item 第一式減第二式：$x - y = 5$ \quad $\cdots$(4)

  \item 第二式$\times 3$減第三式：$5x + 7y = 76$ \quad $\cdots$(5)

  \item 由(4)：$x = y + 5$，代入(5)：$5(y+5) + 7y = 76 \Rightarrow 12y = 51 \Rightarrow y = 4\frac{1}{4}$

  \item 回代：$x = 9\frac{1}{4}$，$z = 2\frac{3}{4}$

\end{enumerate}



\vspace{0.5cm}



\subsection{\zhc{五家共井——五元方程組}}



\begin{paracol}{2}

\switchcolumn[0]

\noindent\textbf{\textcolor{notecolor}{\zhc{【文言文·原題】}}}\par\smallskip

\zhc{

今有五家共井。甲二綆不足，如乙一綆；乙三綆不足，如丙一綆；丙四綆不足，如丁一綆；丁五綆不足，如戊一綆；戊六綆不足，如甲一綆。如各得所不足一綆，皆逮。問井深、綆長各幾何？

}



\switchcolumn[1]

\noindent\textbf{\textcolor{notecolor}{\zhm{【白話文·譯題】}}}\par\smallskip

\zhm{

現有五家共用一口井。甲家的兩條繩子不夠長，差的是乙家一條繩子的長度；乙家的三條繩子不夠長，差的是丙家一條繩子的長度；丙家的四條繩子不夠長，差的是丁家一條繩子的長度；丁家的五條繩子不夠長，差的是戊家一條繩子的長度；戊家的六條繩子不夠長，差的是甲家一條繩子的長度。如果各家都得到所差的那條繩子長度，就都能打到水。問井的深度和各家的繩長各是多少？

}

\end{paracol}



\medskip



\begin{paracol}{2}

\switchcolumn[0]

\noindent\textbf{\textcolor{notecolor}{\zhc{【文言文·術曰】}}}\par\smallskip

\zhc{

術曰：此題以方程術入之，列五行事。甲二綆不足如乙一綆，是二綆加乙一綆等於井深也。故列甲二、乙一、不足零於一行。餘四行類推。五方程、六未知數，為不定方程。以互減消元，至得整數之比而止。



楊輝詳解曰：此五家共井，乃方程中至深之題。五方程、六未知，非獨一解。以比率求之，井深七百二十一寸，甲綆長二百六十五寸，乙綆長一百九十一寸，丙綆長一百四十八寸，丁綆長一百二十九寸，戊綆長七十六寸。

}



\switchcolumn[1]

\noindent\textbf{\textcolor{notecolor}{\zhm{【白話文·解法】}}}\par\smallskip

\zhm{

解法：此題用方程術求解，列五個方程式。甲的兩條繩子不夠，需要加乙的一條繩子才能達到井深，所以$2a + b = d$（$d$為井深）。以此類推列出五個方程：

\begin{align*}

2a + b &= d \\

3b + c &= d \\

4c + d_{rope} &= d \\

5d_{rope} + e &= d \\

6e + a &= d

\end{align*}

這是五個方程、六個未知數的不定方程組。用互減消元法，最終得到整數比例的解。



楊輝詳解：五家共井是方程術中最深奧的題目。五個方程六個未知數，不是唯一解。用比率求解：井深七百二十一寸，甲繩長二百六十五寸，乙繩長一百九十一寸，丙繩長一百四十八寸，丁繩長一百二十九寸，戊繩長七十六寸。

}

\end{paracol}



\vspace{0.5cm}



\subsection{\zhc{金銀交換——方程應用}}



\begin{paracol}{2}

\switchcolumn[0]

\noindent\textbf{\textcolor{notecolor}{\zhc{【文言文·原題】}}}\par\smallskip

\zhc{

今有金九、銀十一，共重適等。交換其一，則金輕十三兩。問金、銀一塊各重幾何？

}



\switchcolumn[1]

\noindent\textbf{\textcolor{notecolor}{\zhm{【白話文·譯題】}}}\par\smallskip

\zhm{

現有九塊金和十一塊銀，總重量恰好相等。交換其中一塊（即金拿一塊給銀，銀拿一塊給金），則金這邊比銀這邊輕十三兩。問金、銀每一塊各重多少？

}

\end{paracol}



\medskip



\begin{paracol}{2}

\switchcolumn[0]

\noindent\textbf{\textcolor{notecolor}{\zhc{【文言文·術曰】}}}\par\smallskip

\zhc{

術曰：以方程術入之。金九、銀十一適等，為一式。交換其一，金輕十三兩，為二式。二式相參，以少減多，以金減銀，餘數約之，得金、銀各一塊之重。



答曰：金重二斤三兩一十八銖，銀重一斤一十三兩六銖。

}



\switchcolumn[1]

\noindent\textbf{\textcolor{notecolor}{\zhm{【白話文·解法】}}}\par\smallskip

\zhm{

解法：用方程術求解。九塊金與十一塊銀重量相等，列為第一式。交換一塊後金比銀輕十三兩，列為第二式。兩式相互配合，以少減多，用金減銀，餘數約簡後得到金、銀每塊的重量。



設金每塊重$g$兩，銀每塊重$s$兩：

\begin{align*}

9g &= 11s \\

8g + s &= 10s + g - 13 \Rightarrow 7g - 9s = -13

\end{align*}

聯立解得：$g = \frac{143}{4} = 35.75$兩 $= 2$斤$3$兩$18$銖，$s = \frac{117}{4} = 29.25$兩 $= 1$斤$13$兩$6$銖。



答案：金每塊重二斤三兩一十八銖，銀每塊重一斤一十三兩六銖。

}

\end{paracol}



\newpage



% ============================================================================

\section{\zhc{勾股第九}}

\label{sec:gougu}



\begin{paracol}{2}

\switchcolumn[0]

\noindent\textbf{\textcolor{notecolor}{\zhc{【文言文】}}}\par\smallskip

\zhc{

勾股術曰：勾股求弦法——勾自乘，股自乘，并而開方除之，即弦。弦勾求股法——弦自乘，勾自乘，以少減多，餘開方除之，即股。弦股求勾亦如之。



楊輝詳解曰：勾股之法，出於圓方。圓徑一而周三，方徑一而匝四。勾三股四弦五，此其率也。諸勾股之術，皆由此而生焉。

}



\switchcolumn[1]

\noindent\textbf{\textcolor{notecolor}{\zhm{【白話文】}}}\par\smallskip

\zhm{

勾股術的計算規則：由勾和股求弦——勾自乘，股自乘，相加後開平方，即得弦。由弦和勾求股——弦自乘，勾自乘，以少減多，餘數開平方，即得股。由弦和股求勾也是同樣的方法。



楊輝詳解：勾股之法，源於圓和方。圓的直徑為一則周長為三，方的直徑為一則周長為四。勾三股四弦五，這就是基本的比率。各種勾股算法，都由此衍生而出。

}

\end{paracol}



\vspace{0.5cm}



\subsection{\zhc{勾三股四弦五——基本定理}}



\begin{paracol}{2}

\switchcolumn[0]

\noindent\textbf{\textcolor{notecolor}{\zhc{【文言文·原題】}}}\par\smallskip

\zhc{

今有勾三尺，股四尺，問弦幾何？

}



\switchcolumn[1]

\noindent\textbf{\textcolor{notecolor}{\zhm{【白話文·譯題】}}}\par\smallskip

\zhm{

現有一個直角三角形，勾（短直角邊）長三尺，股（長直角邊）長四尺，問弦（斜邊）長多少？

}

\end{paracol}



\medskip



\begin{paracol}{2}

\switchcolumn[0]

\noindent\textbf{\textcolor{notecolor}{\zhc{【文言文·術曰】}}}\par\smallskip

\zhc{

術曰：勾自乘，得九尺；股自乘，得一十六尺。并之，得二十五尺。開方除之，得五尺，即弦。



答曰：五尺。

}



\switchcolumn[1]

\noindent\textbf{\textcolor{notecolor}{\zhm{【白話文·解法】}}}\par\smallskip

\zhm{

解法：勾三尺自乘，得九平方尺；股四尺自乘，得十六平方尺。相加得二十五平方尺。開平方，得五尺，就是弦的長度。



答案：五尺。



此即著名的「勾三股四弦五」直角三角形，是勾股定理最基本的實例。用現代公式表示即：$c = \sqrt{a^2 + b^2} = \sqrt{3^2 + 4^2} = \sqrt{25} = 5$。

}

\end{paracol}



\vspace{0.5cm}



\subsection{\zhc{弦勾求股——已知斜邊和一股}}



\begin{paracol}{2}

\switchcolumn[0]

\noindent\textbf{\textcolor{notecolor}{\zhc{【文言文·原題】}}}\par\smallskip

\zhc{

今有弦十七步，勾八步，問股幾何？

}



\switchcolumn[1]

\noindent\textbf{\textcolor{notecolor}{\zhm{【白話文·譯題】}}}\par\smallskip

\zhm{

現有直角三角形的弦（斜邊）長十七步，勾（短直角邊）長八步，問股（長直角邊）長多少步？

}

\end{paracol}



\medskip



\begin{paracol}{2}

\switchcolumn[0]

\noindent\textbf{\textcolor{notecolor}{\zhc{【文言文·術曰】}}}\par\smallskip

\zhc{

術曰：弦自乘，得二百八十九步；勾自乘，得六十四步。以少減多，餘二百二十五步。開方除之，得一十五步，即股。



答曰：十五步。

}



\switchcolumn[1]

\noindent\textbf{\textcolor{notecolor}{\zhm{【白話文·解法】}}}\par\smallskip

\zhm{

解法：弦十七步自乘，得二百八十九平方步；勾八步自乘，得六十四平方步。以少減多，餘二百二十五平方步。開平方，得十五步，就是股的長度。



答案：十五步。



現代計算：$\text{股} = \sqrt{17^2 - 8^2} = \sqrt{289 - 64} = \sqrt{225} = 15$步。

}

\end{paracol}



\vspace{0.5cm}



\subsection{\zhc{葭生池中——名題}}



\begin{paracol}{2}

\switchcolumn[0]

\noindent\textbf{\textcolor{notecolor}{\zhc{【文言文·原題】}}}\par\smallskip

\zhc{

今有池方一丈，葭生其中央，出水一尺。引葭赴岸，適與岸齊。問水深、葭長各幾何？

}



\switchcolumn[1]

\noindent\textbf{\textcolor{notecolor}{\zhm{【白話文·譯題】}}}\par\smallskip

\zhm{

現有一個邊長一丈（十尺）的正方形池塘，蘆葦生長在池塘中央，露出水面一尺。將蘆葦拉向岸邊，恰好與岸齊平。問水深和蘆葦長各是多少？

}

\end{paracol}



\medskip



\begin{paracol}{2}

\switchcolumn[0]

\noindent\textbf{\textcolor{notecolor}{\zhc{【文言文·術曰】}}}\par\smallskip

\zhc{

術曰：半池方自乘，出水一尺自乘，以少減多，餘開方除之，即水深。加水深、出水，即葭長。



草曰：半池方五尺自乘，得二十五尺。出水一尺自乘，得一尺。以少減多，餘二十四尺。開方除之，不可得整數。此以勾股求之：半池方五尺為勾，水深為股，葭長為弦。弦減出水一尺即股，故弦自乘減勾自乘，開方得股。



楊輝曰：半池方五尺為勾，葭出水一尺，故葭長比水深多一尺。設水深$x$尺，則葭長$(x+1)$尺。以勾股定理：$(x+1)^2 = x^2 + 5^2$，解之得$x = 12$，葭長$13$尺。

}



\switchcolumn[1]

\noindent\textbf{\textcolor{notecolor}{\zhm{【白話文·解法】}}}\par\smallskip

\zhm{

解法：半個池塘的邊長（即池中心到岸邊的距離）五尺自乘，得二十五。出水一尺自乘，得一。相減餘二十四，開方（此處古人開方不盡，故換思路）。用勾股定理求解：半池方五尺為勾，水深為股，葭長為弦。弦減去出水一尺就是股，所以弦自乘減勾自乘，開方得股。



楊輝詳解：半池邊長五尺為勾，蘆葦出水一尺，所以蘆葦長比水深多一尺。設水深$x$尺，則蘆葦長$(x+1)$尺。用勾股定理：$(x+1)^2 = x^2 + 5^2$，展開得$x^2 + 2x + 1 = x^2 + 25$，化簡得$2x = 24$，解得$x = 12$，蘆葦長$13$尺。

}

\end{paracol}



\medskip



\noindent\textbf{\textcolor{sectioncolor}{\zhc{現代代數解法：}}}

設水深為$d$尺，葭長為$L$尺：

\begin{align*}

L &= d + 1 \\

L^2 &= d^2 + \left(\frac{10}{2}\right)^2 = d^2 + 25 \\

(d+1)^2 &= d^2 + 25 \Rightarrow 2d + 1 = 25 \Rightarrow d = 12, \quad L = 13

\end{align*}

\zhc{答曰：水深一丈二尺，葭長一丈三尺。}



\vspace{0.5cm}



\subsection{\zhc{戶高廣——門的高與寬}}



\begin{paracol}{2}

\switchcolumn[0]

\noindent\textbf{\textcolor{notecolor}{\zhc{【文言文·原題】}}}\par\smallskip

\zhc{

今有戶高多於廣六尺八寸，兩隅相去適一丈。問戶高、廣各幾何？

}



\switchcolumn[1]

\noindent\textbf{\textcolor{notecolor}{\zhm{【白話文·譯題】}}}\par\smallskip

\zhm{

現有一扇門，高度比寬度多六尺八寸，兩個對角之間的距離恰好是一丈（十尺）。問門的高度和寬度各是多少？

}

\end{paracol}



\medskip



\begin{paracol}{2}

\switchcolumn[0]

\noindent\textbf{\textcolor{notecolor}{\zhc{【文言文·術曰】}}}\par\smallskip

\zhc{

術曰：以一丈自乘，得一百尺。又以高多於廣六尺八寸自乘，得四十六尺二寸四分。以少減多，餘五十三尺七寸六分。半之，得二十六尺八寸八分。開方除之，得廣二尺八寸。加多六尺八寸，得高九尺六寸。



答曰：高九尺六寸，廣二尺八寸。

}



\switchcolumn[1]

\noindent\textbf{\textcolor{notecolor}{\zhm{【白話文·解法】}}}\par\smallskip

\zhm{

解法：用對角線一丈（十尺）自乘，得一百平方尺。又用高度比寬度多的六尺八寸自乘，得四十六點二四平方尺。相減餘五十三點七六平方尺。取半得二十六點八八平方尺。開平方得寬度二尺八寸。加上多出的六尺八寸，得高度九尺六寸。



現代解法：設寬為$w$尺，高為$h$尺：

\begin{align*}

h - w &= 6.8 \\

h^2 + w^2 &= 100

\end{align*}

由第一式$h = w + 6.8$，代入第二式：

$(w + 6.8)^2 + w^2 = 100 \Rightarrow 2w^2 + 13.6w - 53.76 = 0$\\

解得$w = 2.8$尺，$h = 9.6$尺。



答案：高九尺六寸，廣二尺八寸。

}

\end{paracol}



\vspace{0.5cm}



\subsection{\zhc{勾股容方——內接正方形}}



\begin{paracol}{2}

\switchcolumn[0]

\noindent\textbf{\textcolor{notecolor}{\zhc{【文言文·原題】}}}\par\smallskip

\zhc{

今有勾六步，股十一步，問勾中容方幾何？

}



\switchcolumn[1]

\noindent\textbf{\textcolor{notecolor}{\zhm{【白話文·譯題】}}}\par\smallskip

\zhm{

現有直角三角形，勾長六步，股長十一步，問在這個直角三角形內能容納的最大正方形（一邊在弦上）邊長是多少？

}

\end{paracol}



\medskip



\begin{paracol}{2}

\switchcolumn[0]

\noindent\textbf{\textcolor{notecolor}{\zhc{【文言文·術曰】}}}\par\smallskip

\zhc{

術曰：并勾、股為法，勾、股相乘為實，實如法而一，得方邊。



草曰：勾六步、股十一步相乘，得六十六步，為實。并勾、股，得十七步，為法。實如法，得三步十七分步之十五，即方邊也。



答曰：三步十七分步之十五。

}



\switchcolumn[1]

\noindent\textbf{\textcolor{notecolor}{\zhm{【白話文·解法】}}}\par\smallskip

\zhm{

解法：將勾和股相加作為除數，勾和股相乘作為被除數，被除數除以除數得到正方形的邊長。



算草：勾六步乘以股十一步，得六十六平方步，作為被除數。勾加股得十七步，作為除數。被除數除以除數得$3\frac{15}{17}$步，就是正方形的邊長。



答案：$3\frac{15}{17}$步。



現代公式：對於直角邊為$a, b$的直角三角形，內接正方形（一邊在斜邊上）的邊長為$s = \frac{ab}{a+b} = \frac{6 \times 11}{6 + 11} = \frac{66}{17} = 3\frac{15}{17}$步。

}

\end{paracol}



\vspace{0.5cm}



\subsection{\zhc{邑方測望——測量方形城池}}



\begin{paracol}{2}

\switchcolumn[0]

\noindent\textbf{\textcolor{notecolor}{\zhc{【文言文·原題】}}}\par\smallskip

\zhc{

今有邑方不知大小，各中開門。出北門二十步有木。出南門十四步，折而西行一千七百七十五步見木。問邑方幾何？

}



\switchcolumn[1]

\noindent\textbf{\textcolor{notecolor}{\zhm{【白話文·譯題】}}}\par\smallskip

\zhm{

現有一座方形城市，不知道邊長多少，每面牆的正中間各開一門。出北門走二十步有一棵樹。出南門走十四步，然後轉向西走一千七百七十五步就看見那棵樹。問這座方形城市的邊長是多少？

}

\end{paracol}



\medskip



\begin{paracol}{2}

\switchcolumn[0]

\noindent\textbf{\textcolor{notecolor}{\zhc{【文言文·術曰】}}}\par\smallskip

\zhc{

術曰：以出北門步數乘西行步數，倍之，為實。并南、北門步數，為法。實如法而一，即邑方。



草曰：出北門二十步乘西行一千七百七十五步，得三萬五千五百步。倍之，得七萬一千步，為實。并南門十四步、北門二十步，得三十四步，為法。實如法，得二千八十八步……以法約之，得邑方二百五十步。



答曰：二百五十步。

}



\switchcolumn[1]

\noindent\textbf{\textcolor{notecolor}{\zhm{【白話文·解法】}}}\par\smallskip

\zhm{

解法：用出北門的步數乘以向西行的步數，再加倍，作為被除數。將出南門和出北門的步數相加，作為除數。被除數除以除數，即得城邑的邊長。



算草：出北門二十步乘以西行一千七百七十五步，得三萬五千五百步。加倍得七萬一千步，作為被除數。南門十四步加北門二十步，得三十四步，作為除數。被除數除以除數，化簡後得城邑邊長二百五十步。



現代解法：設城邑邊長為$s$。由相似三角形原理：

$$\frac{s/2}{20} = \frac{1775}{s + 34}$$

交叉相乘：$s(s+34) = 2 \times 20 \times 1775 = 71000$\\

$s^2 + 34s - 71000 = 0 \Rightarrow s = 250$步。



答案：二百五十步。

}

\end{paracol}



\newpage



% ============================================================================

\section{\zhc{商功}}

\label{sec:shanggong}



\begin{paracol}{2}

\switchcolumn[0]

\noindent\textbf{\textcolor{notecolor}{\zhc{【文言文】}}}\par\smallskip

\zhc{

商功術曰：穿地四尺為壤五尺，為堅三尺。此穿地求壤四之，求堅三之，皆一而一。城、垣、堤、溝、塹、渠，皆同術。并上下廣而半之，以高或深乘之，又以袤乘之，即積尺。

}



\switchcolumn[1]

\noindent\textbf{\textcolor{notecolor}{\zhm{【白話文】}}}\par\smallskip

\zhm{

商功術的計算規則：挖出的鬆土四尺相當於夯實後的五尺，或相當於堅硬土的三尺。所以由挖土求鬆土乘以$\frac{5}{4}$，求堅土乘以$\frac{3}{4}$。城墙、垣牆、堤壩、溝渠、塹壕、渠道，都用同一種算法。將上底寬和下底寬相加取半，乘以高度或深度，再乘以長度，即得體積（立方尺）。

}

\end{paracol}



\vspace{0.5cm}



\subsection{\zhc{城垣積尺——梯形柱體體積}}



\begin{paracol}{2}

\switchcolumn[0]

\noindent\textbf{\textcolor{notecolor}{\zhc{【文言文·原題】}}}\par\smallskip

\zhc{

今有城下廣四丈，上廣二丈，高五丈，袤一百二十六丈五尺。問積幾何？

}



\switchcolumn[1]

\noindent\textbf{\textcolor{notecolor}{\zhm{【白話文·譯題】}}}\par\smallskip

\zhm{

現有一座城墙，下底寬四丈，上頂寬二丈，高五丈，長一百二十六丈五尺。問體積是多少？

}

\end{paracol}



\medskip



\begin{paracol}{2}

\switchcolumn[0]

\noindent\textbf{\textcolor{notecolor}{\zhc{【文言文·術曰】}}}\par\smallskip

\zhc{

術曰：并上下廣，得六丈。半之，得三丈，為正廣。以高五丈乘之，得一十五丈，為截面积。又以袤一百二十六丈五尺乘之，即積尺。



草曰：下廣四十尺加上廣二十尺，得六十尺。半之，得三十尺。以高五十尺乘之，得一千五百平方尺。以袤一千二百六十五尺乘之，得一百八十九萬七千五百立方尺。



答曰：一百八十九萬七千五百尺。

}



\switchcolumn[1]

\noindent\textbf{\textcolor{notecolor}{\zhm{【白話文·解法】}}}\par\smallskip

\zhm{

解法：將上底寬和下底寬相加，得六丈。取半得三丈，作為平均寬度。乘以高五丈，得一十五萬平方丈，作為橫截面積。再乘以長一百二十六丈五尺，即得體積。



算草：下底寬四十尺加上頂寬二十尺，得六十尺。取半得三十尺。乘以高五十尺，得一千五百平方尺。乘以長度一千二百六十五尺，得一百八十九萬七千五百立方尺。



答案：一百八十九萬七千五百立方尺。



現代公式（梯形截面柱體體積）：

$$V = \frac{(w_{\text{下}} + w_{\text{上}})}{2} \times h \times L = \frac{(40+20)}{2} \times 50 \times 1265 = 1{,}897{,}500 \text{ 立方尺}$$

}

\end{paracol}



\vspace{0.5cm}



\subsection{\zhc{方亭積尺——正方截錐體體積}}



\begin{paracol}{2}

\switchcolumn[0]

\noindent\textbf{\textcolor{notecolor}{\zhc{【文言文·原題】}}}\par\smallskip

\zhc{

今有方亭，下方五丈，上方四丈，高五丈。問積幾何？

}



\switchcolumn[1]

\noindent\textbf{\textcolor{notecolor}{\zhm{【白話文·譯題】}}}\par\smallskip

\zhm{

現有一座方形平台（方亭），下底邊長五丈，上頂邊長四丈，高五丈。問體積是多少？

}

\end{paracol}



\medskip



\begin{paracol}{2}

\switchcolumn[0]

\noindent\textbf{\textcolor{notecolor}{\zhc{【文言文·術曰】}}}\par\smallskip

\zhc{

術曰：上、下方相乘，又各自乘，并之，以高乘之，三而一。



草曰：上方四丈自乘，得一十六萬尺（以丈為單位：四四一十六）。下方五丈自乘，得二十五萬尺。上下方相乘，得二十萬尺。并之，得六十一萬尺。以高五丈乘之，得三百零五萬尺。三而一，得一百零一萬六千六百六十六尺太半尺。



答曰：一十萬一千六百六十六尺太半尺。

}



\switchcolumn[1]

\noindent\textbf{\textcolor{notecolor}{\zhm{【白話文·解法】}}}\par\smallskip

\zhm{

解法：上底邊長和下底邊長相乘，又各自自乘，三數相加，乘以高，再除以三。



算草：上底四丈自乘，得十六（萬平方尺單位）。下底五丈自乘，得二十五。上下底相乘，得二十。三數相加得六十一。乘以高五丈，得三百零五。除以三，得$101\frac{2}{3}$（萬立方尺）。



答案：十萬一千六百六十六又三分之二立方尺。



現代公式（正方截錐體體積）：

$$V = \frac{h}{3}(a^2 + ab + b^2) = \frac{50}{3}(40^2 + 40 \times 50 + 50^2) = \frac{50 \times 6100}{3} = 101{,}666\tfrac{2}{3} \text{ 立方尺}$$

}

\end{paracol}



\vspace{0.5cm}



\subsection{\zhc{圓亭積尺——圓台體積}}



\begin{paracol}{2}

\switchcolumn[0]

\noindent\textbf{\textcolor{notecolor}{\zhc{【文言文·原題】}}}\par\smallskip

\zhc{

今有圓亭，上周三丈，下周二丈，高一丈。問積幾何？

}



\switchcolumn[1]

\noindent\textbf{\textcolor{notecolor}{\zhm{【白話文·譯題】}}}\par\smallskip

\zhm{

現有一座圓形平台（圓亭），上底周長三丈，下底周長二丈，高一丈。問體積是多少？

}

\end{paracol}



\medskip



\begin{paracol}{2}

\switchcolumn[0]

\noindent\textbf{\textcolor{notecolor}{\zhc{【文言文·術曰】}}}\par\smallskip

\zhc{

術曰：上、下周相乘，又各自乘，并之，以高乘之，三十六而一。



草曰：上周三丈自乘，得九；下周二丈自乘，得四；上下周相乘，得六。并之，得十九。以高一丈乘之，得一十九。三十六而一，得十九分之三十六丈。



答曰：積五百二十七尺九寸十二分寸之一。（按古法$\pi \approx 3$計算）

}



\switchcolumn[1]

\noindent\textbf{\textcolor{notecolor}{\zhm{【白話文·解法】}}}\par\smallskip

\zhm{

解法：上底周長和下底周長相乘，又各自自乘，三數相加，乘以高，再除以三十六。



算草：上周三丈自乘，得九；下周二丈自乘，得四；上下周相乘，得六。三數相加得十九。乘以一丈高，得十九。除以三十六，得$\frac{19}{36}$立方丈。



按古代數學中圓周率取三（$\pi \approx 3$），則圓台體積公式為：

$$V = \frac{h}{36}(C_1^2 + C_1 C_2 + C_2^2)$$

其中$C_1$和$C_2$分別為上、下底周長。



答案：體積為五百二十七尺九寸又十二分寸之一（換算後）。

}

\end{paracol}



\vspace{0.5cm}



\subsection{\zhc{委粟平地——圓錐體積}}



\begin{paracol}{2}

\switchcolumn[0]

\noindent\textbf{\textcolor{notecolor}{\zhc{【文言文·原題】}}}\par\smallskip

\zhc{

今有委粟平地，下周一十二丈，高二丈。問積及為粟各幾何？

}



\switchcolumn[1]

\noindent\textbf{\textcolor{notecolor}{\zhm{【白話文·譯題】}}}\par\smallskip

\zhm{

現在平地上堆放着粟米，底面周長十二丈，高兩丈。問這堆粟米的體積是多少，以及合多少斛粟？

}

\end{paracol}



\medskip



\begin{paracol}{2}

\switchcolumn[0]

\noindent\textbf{\textcolor{notecolor}{\zhc{【文言文·術曰】}}}\par\smallskip

\zhc{

術曰：下周自乘，以高乘之，三十六而一，即積尺。以斛法除之，即粟數。



草曰：下周十二丈自乘，得一百四十四。以高二丈乘之，得二百八十八。三十六而一，得八丈，即八千立方尺。以斛法一斛積一尺六寸二分除之，得粟四千九百三十八斛二十七分解之二十六。



答曰：積八千尺；為粟四千九百三十八斛二十七分斛之二十六。

}



\switchcolumn[1]

\noindent\textbf{\textcolor{notecolor}{\zhm{【白話文·解法】}}}\par\smallskip

\zhm{

解法：底面周長自乘，乘以高，再除以三十六，即得體積（立方尺）。用每斛的體積去除，即得粟米的斛數。



算草：周長十二丈自乘，得一百四十四。乘以高二丈，得二百八十八。除以三十六，得八立方丈，即八千立方尺。以每斛體積一點六二立方尺去除，得粟$4938\frac{26}{27}$斛。



答案：體積八千立方尺；折合粟四千九百三十八又二十七分之二十六斛。



現代公式（圓錐體積，按$\pi = 3$）：

$$V = \frac{C^2 \times h}{36} = \frac{120^2 \times 20}{36} = \frac{288000}{36} = 8000 \text{ 立方尺}$$

}

\end{paracol}



\newpage



% ============================================================================

\section{\zhc{開方作法本源——楊輝三角}}

\label{sec:triangle}



\begin{paracol}{2}

\switchcolumn[0]

\noindent\textbf{\textcolor{notecolor}{\zhc{【文言文】}}}\par\smallskip

\zhc{

開方作法本源：左裪乃積數，右裪乃隅算，中藏者皆廉。以廉乘商方，命實而除之。賈憲用此術開方，楊輝錄之以傳後世。



楊輝詳解曰：開方作法本源圖，出於釋鎖算書，賈憲用此術。其圖如後：

}



\switchcolumn[1]

\noindent\textbf{\textcolor{notecolor}{\zhm{【白話文】}}}\par\smallskip

\zhm{

開方作法本源圖的說明：左邊的數字是各次冪的係數（積數），右邊的數字是隅算（最高次項係數，恆為一），中間藏着的都是廉（各項中間係數）。用這些廉數去乘商的冪次，然後從實中減去。賈憲用這個方法開方，楊輝將它記錄下來傳給後世。



楊輝詳解：開方作法本源圖出自《釋鎖算書》，賈憲使用這種方法。圖形如後所示：

}

\end{paracol}



\medskip



\begin{center}

\begin{tabular}{ccccccccccc}

 & & & & & 1 & & & & & \\

 & & & & 1 & & 1 & & & & \\

 & & & 1 & & 2 & & 1 & & & \\

 & & 1 & & 3 & & 3 & & 1 & & \\

 & 1 & & 4 & & 6 & & 4 & & 1 & \\

 1 & & 5 & & 10 & & 10 & & 5 & & 1 \\

\end{tabular}

\end{center}



\medskip



\begin{paracol}{2}

\switchcolumn[0]

\zhc{

此圖之用法：欲開某數之平方，則取第二行係數一、二、一；欲開立方，則取第三行係數一、三、三、一；欲開四次方，則取第四行係數一、四、六、四、一。依此類推，各以其係數乘商之各次冪，累減於實，逐位開之。



楊輝曰：此術精妙，乃開方之根本。增乘開方法即以此為基，遞增遞乘，逐位求商，可開任意高次冪。

}



\switchcolumn[1]

\zhm{

此圖的使用方法：要開某數的平方，就取第二行係數一、二、一；要開立方，就取第三行係數一、三、三、一；要開四次方，就取第四行係數一、四、六、四、一。依此類推，各自用這些係數去乘商的各次冪，從實中累減，逐位開方。



楊輝說：這個方法精妙絕倫，是開方的根本算法。增乘開方法就是以它為基礎，遞增遞乘，逐位求商，可以開任意高次方。

}

\end{paracol}



\medskip



\noindent\textbf{\textcolor{sectioncolor}{\zhc{二項式展開應用：}}}

\begin{align*}

(a+b)^2 &= a^2 + 2ab + b^2 \\

(a+b)^3 &= a^3 + 3a^2b + 3ab^2 + b^3 \\

(a+b)^4 &= a^4 + 4a^3b + 6a^2b^2 + 4ab^3 + b^4 \\

(a+b)^5 &= a^5 + 5a^4b + 10a^3b^2 + 10a^2b^3 + 5ab^4 + b^5

\end{align*}



\vspace{0.5cm}



\subsection{\zhc{增乘開方法——高次方根}}



\begin{paracol}{2}

\switchcolumn[0]

\noindent\textbf{\textcolor{notecolor}{\zhc{【文言文】}}}\par\smallskip

\zhc{

增乘開方法曰：商數第一位，以隅法乘之，遞增遞乘，入於中方。又以商數乘中方，入於上方。又以商數乘上方，入於立方。命實除之，適盡。不盡者，以餘實為實，續商下位。



楊輝詳解曰：此術乃賈憲所立，以增乘遞進，取代舊法之繁瑣。每商一位，則以隅法遞乘各層，層層累加，以減於實。其理與西洋之魯斐尼——霍納法暗合，而早於彼數百年矣。

}



\switchcolumn[1]

\noindent\textbf{\textcolor{notecolor}{\zhm{【白話文】}}}\par\smallskip

\zhm{

增乘開方法的計算規則：商的第一位數字，用隅法去乘它，遞增遞乘，加入中方。又用商數乘中方，加入上方。又用商數乘上方，加入立方。然後用這個結果去除實，恰好除盡最好。除不盡的，用餘下的實作為新的被除數，繼續商下一位。



楊輝詳解：這個方法是賈憲所創立的，用增乘遞進的方法，取代了舊法的繁瑣。每商一位數，就用隅法遞乘各層，層層累加，從實中減去。其原理與西洋的魯菲尼——霍納法（Horner's method）不謀而合，卻比那早了數百年。

}

\end{paracol}



\medskip



\noindent\textbf{\textcolor{sectioncolor}{\zhc{增乘開方法示例（開平方$\sqrt{55225}$）：}}}



\begin{enumerate}[leftmargin=2em]

  \item 將55225分為兩節：$5|52|25$，知答案為三位數

  \item 第一位商2：$2^2 = 4$，$5 - 4 = 1$。餘數帶下：152

  \item 以商2乘隅1入方：$2 \times 1 = 2$（方）。倍方為$4$，續商3：$43 \times 3 = 129$。$152 - 129 = 23$。餘數帶下：2325

  \item 以商3乘隅入方，遞增：方為$46$，續商5：$465 \times 5 = 2325$。適盡。

  \item 答案：$\sqrt{55225} = 235$

\end{enumerate}



\newpage



% ============================================================================

\section*{\zhc{附錄：九章纂類}}

\addcontentsline{toc}{section}{附錄：九章纂類}



\begin{paracol}{2}

\switchcolumn[0]

\noindent\textbf{\textcolor{notecolor}{\zhc{【文言文】}}}\par\smallskip

\zhc{

楊輝竊見九章舊本，作立題法，迫圓古序云：靖康以來，罕有舊本，間有存者，狃於末習，不循本意，以至具術淹廢，偽本滋興。此說固然，殊不知所傳之本亦不得其真矣。



楊輝以九章二百四十六問，按術分類，纂為九類：一曰乘除，四十一問；二曰分率，四十一問；三曰合率，十問；四曰互換，十一問；五曰衰分，十八問；六曰疊積，二十七問；七曰盈不足，二十問；八曰方程，十九問；九曰勾股，二十四問。共二百四十六問。

}



\switchcolumn[1]

\noindent\textbf{\textcolor{notecolor}{\zhm{【白話文】}}}\par\smallskip

\zhm{

楊輝見到《九章》的舊本，寫道：自從靖康年間以來，罕見舊本，偶有存留的，也被後人習慣於末流做法所影響，不遵循本意，以致正統算法湮沒失傳，偽本滋生泛濫。這種說法固然有理，殊不知所流傳的版本也已經不是原本的真貌了。



楊輝將《九章》的二百四十六道問題，按照算法分類，編纂為九類：第一類乘除，四十一問；第二類分率，四十一問；第三類合率，十問；第四類互換，十一問；第五類衰分，十八問；第六類疊積，二十七問；第七類盈不足，二十問；第八類方程，十九問；第九類勾股，二十四問。共計二百四十六問。

}

\end{paracol}



\medskip



\begin{center}

\begin{longtable}{|c|l|c|}

\hline

\textbf{序號} & \textbf{\zhc{類別}} & \textbf{\zhc{題數}} \\

\hline

1 & \zhc{乘除}（乘法與除法） & 41 \\

2 & \zhc{分率}（分數與比率） & 41 \\

3 & \zhc{合率}（合併比率） & 10 \\

4 & \zhc{互換}（交換與換算） & 11 \\

5 & \zhc{衰分}（按比例分配） & 18 \\

6 & \zhc{疊積}（堆疊求積） & 27 \\

7 & \zhc{盈不足}（盈餘與不足） & 20 \\

8 & \zhc{方程}（線性方程組） & 19 \\

9 & \zhc{勾股}（直角三角形） & 24 \\

\hline

\multicolumn{2}{|r|}{\textbf{\zhc{總計}}} & \textbf{246} \\

\hline

\end{longtable}

\end{center}



\vspace{2cm}

\begin{center}

\longline\\[0.5cm]

{\zhc{\small 詳解九章算法（二）· 文言文↔白話文對照版 · 完}}

\end{center}



\end{document}







% ============================================================

% Source: translations/zh/chinese/yang-hui-xiangjie-jiuzhang-part3_classical-modern_bilingual.tex

% ============================================================



%% ========================================================================

%% Yang Hui's Xiangjie Jiuzhang Suanfa — Part III (Reclassification)

%% BILINGUAL EDITION: Classical Chinese <-> Modern Chinese Vernacular

%% 详解九章算法 纂类 — 文言原文与白话译文对照本

%% ========================================================================



\documentclass[11pt,a4paper]{article}



%% -------- Core Packages --------

\usepackage{fontspec}

\usepackage{polyglossia}

\usepackage{amsmath,amssymb,amsthm}

\usepackage{geometry}

\usepackage{graphicx}

\usepackage{xcolor}

\usepackage{titlesec}

\usepackage{fancyhdr}

\usepackage{enumitem}

\usepackage{array,booktabs,longtable,multirow}

\usepackage{hyperref}

\usepackage{microtype}

\usepackage{tocloft}

\usepackage{xeCJK}

\usepackage{etoolbox}



%% -------- Page Geometry --------

\geometry{

  paperwidth=210mm,paperheight=297mm,

  left=18mm,right=18mm,top=25mm,bottom=25mm,

  headheight=14pt,footskip=30pt

}



%% -------- Language Setup --------

\setmainlanguage{english}



%% -------- Font Configuration --------

\setmainfont{Times New Roman}

\setsansfont{Arial}

\setmonofont{Courier New}



%% CJK Support via xeCJK

\setCJKmainfont{Microsoft YaHei}

\setCJKsansfont{Microsoft YaHei}

\setCJKmonofont{Microsoft YaHei}



%% -------- Colors --------

\definecolor{titlecolor}{RGB}{45,55,72}

\definecolor{sectioncolor}{RGB}{55,65,81}

\definecolor{notecolor}{RGB}{120,53,15}

\definecolor{prefacecolor}{RGB}{80,60,40}

\definecolor{classicalbg}{RGB}{255,252,240}

\definecolor{modernbg}{RGB}{240,248,255}

\definecolor{classicalframe}{RGB}{180,160,120}

\definecolor{modernframe}{RGB}{120,160,180}

\definecolor{colheadcl}{RGB}{139,90,43}

\definecolor{colheadmd}{RGB}{43,90,139}



%% -------- Section Formatting --------

\titleformat{\section}{\normalfont\Large\bfseries\color{sectioncolor}}{\thesection}{1em}{}

\titleformat{\subsection}{\normalfont\large\bfseries\color{sectioncolor}}{\thesubsection}{1em}{}

\titleformat{\subsubsection}{\normalfont\normalsize\bfseries\color{sectioncolor}}{\thesubsubsection}{1em}{}



%% -------- Header/Footer --------

\pagestyle{fancy}

\fancyhf{}

\fancyhead[L]{\small\textcolor{colheadcl}{【文言】}\quad\textcolor{colheadmd}{【白话】}}

\fancyhead[R]{\small\thepage}

\renewcommand{\headrulewidth}{0.4pt}

\renewcommand{\footrulewidth}{0pt}



%% -------- Custom Commands --------

\newcommand{\longline}{\noindent\rule{\textwidth}{0.4pt}}

\newcommand{\shortline}{\noindent\rule{0.5\textwidth}{0.4pt}\par}



%% Column header

\newcommand{\colheader}{%

  \begin{center}

  \begin{minipage}[t]{0.45\textwidth}

  \centering\textcolor{colheadcl}{\textbf{【文言原文】}}\\[2pt]

  \footnotesize\textcolor{colheadcl}{Classical Chinese (文言文)}

  \end{minipage}%

  \hfill%

  \begin{minipage}[t]{0.45\textwidth}

  \centering\textcolor{colheadmd}{\textbf{【白话译文】}}\\[2pt]

  \footnotesize\textcolor{colheadmd}{Modern Chinese Vernacular (白话文)}

  \end{minipage}

  \end{center}

  \vspace{6pt}

  \noindent\rule{\textwidth}{0.4pt}

  \vspace{6pt}

}



%% Bilingual paragraph pair (main workhorse)

\newcommand{\bilingual}[2]{%

\noindent

\begin{minipage}[t]{0.45\textwidth}

\setlength{\parindent}{2em}\setlength{\parskip}{4pt}%

\small\color{prefacecolor}#1%

\end{minipage}%

\hfill

\begin{minipage}[t]{0.45\textwidth}

\setlength{\parindent}{2em}\setlength{\parskip}{4pt}%

\small#2%

\end{minipage}

\vspace{8pt}

}



%% Colored box pair for methods/rules

\newcommand{\bilingualmethod}[2]{%

\vspace{6pt}

\noindent

\fcolorbox{classicalframe}{classicalbg}{%

  \begin{minipage}[t]{0.42\textwidth}

  \setlength{\parindent}{1.5em}\setlength{\parskip}{3pt}%

  \small\color{prefacecolor}\textit{#1}%

  \end{minipage}}%

\hfill

\fcolorbox{modernframe}{modernbg}{%

  \begin{minipage}[t]{0.42\textwidth}

  \setlength{\parindent}{1.5em}\setlength{\parskip}{3pt}%

  \small\textit{#2}%

  \end{minipage}}%

\vspace{8pt}

}



%% Bilingual question display

\newcommand{\biquestion}[2]{%

\noindent

\begin{minipage}[t]{0.45\textwidth}

\small\textbf{问：}#1%

\end{minipage}%

\hfill

\begin{minipage}[t]{0.45\textwidth}

\small\textbf{问题：}#2%

\end{minipage}

\vspace{4pt}

}



%% Bilingual answer display

\newcommand{\bianswer}[2]{%

\noindent

\begin{minipage}[t]{0.45\textwidth}

\small\textcolor{notecolor}{\textbf{答曰：}#1}%

\end{minipage}%

\hfill

\begin{minipage}[t]{0.45\textwidth}

\small\textcolor{notecolor}{\textbf{答案：}#2}%

\end{minipage}

\vspace{6pt}

}



%% Problem header

\newcounter{probnum}

\newcommand{\probheader}[1]{%

  \refstepcounter{probnum}%

  \vspace{10pt}%

  \noindent\textbf{\color{sectioncolor}题例 \theprobnum：#1}\par\smallskip

}



%% -------- Hyperref Setup --------

\hypersetup{

  colorlinks=true,linkcolor=sectioncolor,citecolor=sectioncolor,urlcolor=sectioncolor,

  pdfencoding=auto,

  pdftitle={详解九章算法 纂类 — 文言白话对照本 (Part III)},

  pdfauthor={杨辉 (Yang Hui)},

}



%% -------- TOC Depth --------

\setcounter{tocdepth}{2}

\setcounter{secnumdepth}{2}

\setlength{\parindent}{0pt}

\setlength{\parskip}{6pt plus 2pt minus 2pt}



%% ========================================================================

\begin{document}



%% ===== TITLE PAGE =====

\begin{titlepage}

\centering

\vspace*{2cm}

{\fontsize{30}{38}\selectfont\color{titlecolor} 详解九章算法}

\vspace{0.8cm}



{\fontsize{22}{28}\selectfont\color{sectioncolor} 纂类}

\vspace{0.3cm}



{\Large\color{prefacecolor} 文言原文与白话文对照本}

\vspace{1.2cm}



\shortline

\vspace{0.8cm}



{\LARGE 宋 \quad 杨辉 \quad 编次}

\vspace{0.6cm}



{\large\textit{Yang Hui's Detailed Analysis of the Mathematical\\[0.3em]

Methods in the `Nine Chapters' --- Part III}}\\[0.4em]

{\normalsize Bilingual Edition: Classical Chinese $\leftrightarrow$ Modern Chinese Vernacular}

\vspace{0.4cm}



{\normalsize (南宋理宗景定二年，1261 年)}

\vspace{1.2cm}



\shortline

\vspace{0.8cm}



\begin{minipage}{0.88\textwidth}

\centering

\small

\textbf{【凡例说明】}\\[0.3em]

本书页采用双栏对照排版：\\[0.2em]

\textcolor{colheadcl}{$\bullet$ 左栏（暖棕底色）} —— 收录杨辉《详解九章算法·纂类》原文，保留宋代数学用语\\[0.2em]

\textcolor{colheadmd}{$\bullet$ 右栏（淡蓝底色）} —— 提供现代白话文翻译，便于当代读者理解\\[0.2em]

将《九章算术》二百四十六问按数学方法重新分类为九大门类

\end{minipage}



\vfill

{\small 现代排版制作版}

\end{titlepage}



%% ===== TABLE OF CONTENTS =====

\tableofcontents

\newpage



%% ===== COLUMN HEADER =====

\colheader



%% ===== PREFACE =====

\section*{纂类序（Preface to the Reclassification）}

\addcontentsline{toc}{section}{纂类序}



\bilingual{%

黄帝九章古序云：国家尝设算科取士，选九章以为算经之首，盖犹儒者之六经，医家之难素，兵法之孙子欤。昔圣宋绍兴戊辰，算士荣棨谓靖康以来，罕有旧本，间有存者，狃于末习，向获善本，得其全经，复起于学。以魏景元元年刘徽等、唐朝议大夫、行太史令上轻车都尉李淳风等注释，圣宋右班直贾宪细草。



辉尝闻学者，谓九章题问颇隐，法理难明，不得其门而入。于是以答参问，用草考法，因法推类，然后知斯文非古之全经也。将后贤补赘之文，修前代已废之法，删立题术，又纂法问，详著于后，傥得贤者，改而正诸，是所愿也。

}{%

《黄帝九章算术》的古序说：国家曾经设立算学科举考试来选拔人才，把《九章算术》选为算经的第一部，这就好像儒家学者的六经，医学家的《难经》《素问》，兵法家的《孙子兵法》一样重要。从前圣宋绍兴戊辰年间，算学家荣棨说，自靖康年间以来，很少有旧版本流传，偶尔有保存下来的，也被后来的习惯所局限。后来我得到了好的版本，获得了完整的经书，才使这门学问重新振兴起来。这部书有魏景元元年（260年）刘徽等人、唐朝议大夫兼行太史令上轻车都尉李淳风等人的注释，以及本朝右班直贾宪的详细算草。



我曾经听学者们说，《九章算术》的题目和问法相当隐晦，其中的法则和道理难以明了，让人找不到入门的路径。于是我便以答案来参照问题，用算草来考证方法，根据方法推求类别，这之后才知道这本书并不是古代的完整原经。我把后世贤者补充添附的文字加以整理，把前代已经废弃的法则加以修复，删改并重新设立题目和方法，又编纂了法则和问答，详细地写在后面，如果能够得到贤明之士加以改正和校正，那就是我的心愿了。

}



\newpage

\colheader



%% ===== TABLE OF CONTENTS OF THE ORIGINAL =====

\section*{九章互见目录（Cross-Referenced Contents）}

\addcontentsline{toc}{section}{九章互见目录}



\bilingual{%

杨辉窃见九章旧本，作立题法遗阙。古序云：靖康以来，罕有旧本，间有存者，狃于末习，不循本意，以至真术淹废，伪本滋典。此说固然，殊不知所传之本亦不得其真矣。如粟米章之互换、少广章之求由开方，皆重叠无谓，而作者题问不归章次，亦有之。今作纂类，互见目录，以辩其讹，后之明者，更为详释，不亦善乎？



古本二百四十六问，其分布如下：

}{%

杨辉私下考察《九章算术》的旧版本，发现其中设立题目的法则有所遗漏缺失。古序说：自靖康以来，很少有旧版本留存，偶尔有保存下来的，也被后来的习惯所局限，不遵循本来的意图，以至于真正的算法被淹没废弃，错误的版本却滋生流传。这种说法固然有理，但殊不知所流传的版本也已经不是真正的原本了。比如粟米章中的互换问题、少广章中求平方根的方法，都有重复而无意义的内容，而且编撰者的题目和问答不按章节顺序排列的情况也存在。现在我编撰这部《纂类》，制作互见目录，用来辨别其中的错误，后世的明智之人能够进一步详细解释，不也是好事吗？



古本共二百四十六道题，分布如下：

}



\begin{center}

\begin{tabular}{|l|c|l|c|}

\hline

\textbf{章别} & \textbf{问数} & \textbf{章别} & \textbf{问数} \\

\hline

方田 & 38问（并乘除问） & 均输 & 28问 \\

粟米 & 46问（乘除6问，互换31） & 盈不足 & 20问 \\

衰分 & 20问（互换11，衰分9） & 方程 & 18问 \\

少广 & 24问 & 勾股 & 24问 \\

商功 & 28问（叠积27） & 合率 & 11问 \\

\hline

\end{tabular}

\end{center}



\bilingual{%

类题以物理分章，有题法又互之讹。今将二百四十六问分别门例，使后学亦可周知也。

}{%

按类别分题，以事物的性质来分章，但有些题目和法则存在互相错乱的讹误。现在将二百四十六问分别归属到各个门类条例中，使后世的学者也能够全面了解了。

}



\newpage

\colheader



%% ===== MAIN CONTENT =====

\section{乘除第一（Multiplication and Division）}



\bilingual{%

\textbf{凡四十一问}。今考该四十问：方田三十八，粟米二问。

}{%

\textbf{共四十一道题}。现在考证实为四十问：方田章三十八问，粟米章二问。

}



\subsection{基本方法（Basic Methods）}



\bilingualmethod{%

\textbf{直田法}曰：广从相乘为实，如亩法而一。

}{%

\textbf{直田法}：将田地的宽（广）和长（从）相乘得到被除数（实），然后除以亩法（一亩 = 240 平方步）即得亩数。

}



\probheader{方田一}

\biquestion{%

广十五步，从十六步，问田几何？

}{%

有一块田地，宽十五步，长十六步，问这块田有多少亩？

}

\bianswer{%

一亩。

}{%

一亩。

}



\probheader{方田二}

\biquestion{%

广十二步，从十四步，问田几何？

}{%

有一块田地，宽十二步，长十四步，问这块田面积是多少？

}

\bianswer{%

一百六十八步。

}{%

一百六十八平方步。

}



\subsection{里田法（Li-field Method）}



\bilingualmethod{%

\textbf{里田法}曰：方自乘为积里，以一里之积三百七十五亩乘之。

}{%

\textbf{里田法}：将一里的边长自乘得到平方里的面积，然后乘以一平方里等于三百七十五亩的换算率。

}



\probheader{里田一}

\biquestion{%

方田一里，问为田几何？

}{%

有一方形田地，边长为一里，换算成亩是多少？

}

\bianswer{%

三顷七十五亩。

}{%

三顷七十五亩。

}



\probheader{里田二}

\biquestion{%

广二里，从三里，问田几何？

}{%

有一块田地，宽二里，长三里，问面积是多少？

}

\bianswer{%

二十二顷五十亩。

}{%

二十二顷五十亩。

}



\subsection{圭田法（Triangular Field Method）}



\bilingualmethod{%

\textbf{圭田法}曰：半广以乘正从，或半正从以乘广。

}{%

\textbf{圭田法}：取宽度的一半乘以正长，或者取正长的一半乘以宽度。（即三角形面积 = 底 $\times$ 高 $\div$ 2）

}



\probheader{圭田}

\biquestion{%

圭田广十二步，从二十一步，问田几何？

}{%

有一块三角形田地，底宽十二步，高二十一步，问面积是多少？

}

\bianswer{%

一百二十六步。

}{%

一百二十六平方步。

}



\newpage

\colheader



\subsection{斜田法（Trapezoidal Field Method）}



\bilingualmethod{%

\textbf{斜田法}曰：并两斜半之以乘正从，或并两广乘半从。

}{%

\textbf{斜田法}：将两条斜边相加取一半，再乘以正长；或者将两条宽相加，再乘以半长。（梯形面积 = （上底 + 下底）$\times$ 高 $\div$ 2）

}



\probheader{斜田}

\biquestion{%

斜田正广六十五步，一畔从七十二步，一畔从一百步，问田几何？

}{%

有一块梯形田地，正宽六十五步，一侧长七十二步，另一侧长一百步，问面积是多少？

}

\bianswer{%

九亩一百四十四步。

}{%

九亩一百四十四平方步。

}



\subsection{圆田法（Circular Field Method）}



\bilingualmethod{%

\textbf{圆田法}曰：半周半径相乘，半周自乘三而一，周自乘十二而一，径自乘三之四而一。

}{%

\textbf{圆田法}：半周长与半径相乘得面积；或半周长自乘再除以三；或周长自乘再除以十二；或直径自乘乘以三再除以四。（按古代圆周率 $\pi \approx 3$）

}



\probheader{圆田}

\biquestion{%

圆田周三十步，径十步，问田几何？（古术）

}{%

有一块圆形田地，周长三十步，直径十步，问面积是多少？（按古代算法）

}

\bianswer{%

七十五步。

}{%

七十五平方步。

}



\bilingualmethod{%

\textbf{密率}：周自乘，又二十五乘之，三百一十四而一。

}{%

\textbf{密率法}：周长自乘，再乘以二十五，然后除以三百一十四。（即按 $\pi = 157/50 = 3.14$ 计算）

}



\subsection{除率法（Division Rate Method）}



\bilingualmethod{%

\textbf{经率法}（俗名商除）曰：钱数为实，以所买率为法，实如法而一。原载粟米章。

}{%

\textbf{经率法}（俗称商除）：将钱数作为被除数（实），把所买东西的单价率作为除数（法），实除以法即得单价。原载于粟米章。

}



\probheader{经率一}

\biquestion{%

钱一百六十文，买瓴甓十八枚，问枚价几何？

}{%

有一百六十文钱，买砖十八枚，问每枚砖的价钱是多少？

}

\bianswer{%

一枚八钱九分钱之八。

}{%

每枚八又九分之八文钱。

}



\probheader{经率二}

\biquestion{%

钱十三贯五百，买竹二千三百五十个，问个价？

}{%

有一万三千五百文钱，买竹子二千三百五十根，问每根竹子的价钱？

}

\bianswer{%

五钱四十七分钱之三十五。

}{%

每根五又四十七分之三十五文钱。

}



\newpage

\colheader



\section{互换第二（Exchange and Barter）}



\bilingual{%

\textbf{凡五十六问}。今考该五十五问：粟米三十一，衰分十一，均输十一，盈朒二。

}{%

\textbf{共五十六道题}。现在考证实为五十五问：粟米章三十一问，衰分章十一问，均输章十一问，盈不足章二问。

}



\subsection{互换法（Exchange Method）}



\bilingualmethod{%

\textbf{互换法}曰：以所求率乘所有数，为实；以所有率为法，实如法而一。

}{%

\textbf{互换法}：用所求物品的比率乘以已有的数量，作为被除数（实）；用已有物品的比率作为除数（法），实除以法即得所求的数量。

}



\bilingualmethod{%

\textbf{置位法}曰：钱钱物物数数率率依本色对列，其各物原率随而下布。

}{%

\textbf{置位法}：在算板上布列时，将钱、物、数量、比率按照各自的类别对应排列，各种物品的原始比率依次向下布置。

}



\subsection{粟米互换（Grain Conversion）}



\probheader{粟米互换一}

\biquestion{%

粟二斗一升，问为粺米几何？

}{%

有粟（带壳谷物）二斗一升，问能舂成多少粺米（精米）？

}

\bianswer{%

一斗一升五十分升之十七。

}{%

一斗一升又五十分之十七升。

}



\probheader{粟米互换二}

\biquestion{%

粟三斗六升，问为粺饭几何？

}{%

有粟三斗六升，问能做成多少粺米饭？

}

\bianswer{%

三斗八升二十五分升之二十二。

}{%

三斗八升又二十五分之二十二升。

}



\probheader{粟米互换三}

\biquestion{%

粟八斗六升，问为糳饭几何？

}{%

有粟八斗六升，问能做成多少糳米饭？

}

\bianswer{%

八斗二升二十五分升之一十四。

}{%

八斗二升又二十五分之十四升。

}



\probheader{粟米互换四}

\biquestion{%

粟九斗八升，问为御饭几何？

}{%

有粟九斗八升，问能做成多少御米饭？

}

\bianswer{%

八斗二升二十五分升之八。

}{%

八斗二升又二十五分之八升。

}



\probheader{粟米互换五}

\biquestion{%

粟七斗八升，问为豉几何？

}{%

有粟七斗八升，问能制成多少豆豉？

}

\bianswer{%

九斗八升二十五分升之七。

}{%

九斗八升又二十五分之七升。

}



\probheader{粟米互换六}

\biquestion{%

粟五斗五升，问为飧几何？

}{%

有粟五斗五升，问能做成多少熟食（飧）？

}

\bianswer{%

九斗九升。

}{%

九斗九升。

}



\probheader{粟米互换七}

\biquestion{%

粟四斗，问为熟菽几何？

}{%

有粟四斗，问能煮熟多少豆子？

}

\bianswer{%

八斗二升五分升之四。

}{%

八斗二升又五分之四升。

}



\newpage

\colheader



\probheader{粟米互换八}

\biquestion{%

粟二斗，问为蘖几何？

}{%

有粟二斗，问能发芽制成多少糵（麦芽）？

}

\bianswer{%

七斗。

}{%

七斗。

}



\probheader{粟米互换九}

\biquestion{%

粟三斗少半升，问为菽几何？

}{%

有粟三斗又三分之一升，问能换多少豆子？

}

\bianswer{%

二斗七升十分升之三。

}{%

二斗七升又十分之三升。

}



\probheader{粟米互换十}

\biquestion{%

粟四斗一升太半升，问为答几何？

}{%

有粟四斗一升又三分之二升，问能换多少大麦？

}

\bianswer{%

三斗七升半。

}{%

三斗七升半。

}



\probheader{粟米互换十一}

\biquestion{%

粟五斗太半升，问为麻几何？

}{%

有粟五斗又三分之二升，问能换多少麻？

}

\bianswer{%

四斗五升五分升之三。

}{%

四斗五升又五分之三升。

}



\probheader{粟米互换十二}

\biquestion{%

粟一斗，问为粝米几何？

}{%

有粟一斗，问能舂成多少糙米？

}

\bianswer{%

六升。

}{%

六升。

}



\probheader{粟米互换十三}

\biquestion{%

粟四斗五升，问为糳米几何？

}{%

有粟四斗五升，问能舂成多少精白米？

}

\bianswer{%

二斗一升五分升之三。

}{%

二斗一升又五分之三升。

}



\subsection{持金出关（Taxation Problems）}



\probheader{持金出关}

\biquestion{%

今持金十二斤出关，税之十分而取一。今关取金二斤，价钱五千。问金一斤值钱几何？

}{%

现在有人带着十二斤金子出关，按规定要征收十分之一的关税。现在关卡收取了二斤金子，作价五千文钱。问一斤金子值多少钱？

}

\bianswer{%

六千二百五十。

}{%

六千二百五十文。

}



\probheader{持米出三关}

\biquestion{%

持米出三关，外关三分税一，中关五分税一，内关七分税一，余存米五斗。问：原米几何？

}{%

有人带着米出三道关卡，外关征收三分之一，中关征收五分之一，内关征收七分之一，最后还剩五斗米。问：原来有多少米？

}

\bianswer{%

一十斗九升八分升之三。

}{%

十斗九升又八分之三升。

}



\newpage

\colheader



\section{合率第三（Combined Rates）}



\bilingual{%

\textbf{凡二十问}：少广十一，均输八，盈不足一。

}{%

\textbf{共二十道题}：少广章十一问，均输章八问，盈不足章一问。

}



\subsection{少广法（Short Width Method）}



\bilingual{%

少广反合分并率。设诸分母子并而为广，借田求纵，立少广章，易合分之术而为法。用副置分母自乘，以乘全步及诸子，各以本母除其子而并之，免互乘之繁。又得此术，兼助合分，使法引伸，不亦善乎？

}{%

少广的方法是反过来的合分并率。假设把各个分数的分母和分子合并起来作为宽度，借用田地面积来求长度，设立少广章，把合分的方法改换过来作为法则。另外把分母自乘，再乘以整数步和各个分子，各自用原来的分母去除分子然后合并，这样就避免了互乘的繁琐。又得到这个方法，还能辅助合分的运算，使法则得以推广延伸，不也是很好的吗？

}



\bilingualmethod{%

\textbf{少广法}曰：列置全步及分母子，而副置分母自乘，以乘全步及子，各以本母除子，并之为法，以全步积分乘亩步为实，实如法而一。

}{%

\textbf{少广法}：在算板上布列整数步和各个分数的分子分母，另外把分母自乘，再乘以整数步和分子，各自用原来的分母去除分子，合并起来作为除数（法），用整数步的积分乘以一亩的步数作为被除数（实），实除以法即得长度。

}



\probheader{少广一}

\biquestion{%

田一亩广一步半，问从？

}{%

有一块田面积一亩，宽一步半，问长是多少？

}

\bianswer{%

一百六十步。

}{%

一百六十步。

}



\probheader{少广二}

\biquestion{%

田一亩广一步半，三分步之一，问从？

}{%

有一块田面积一亩，宽一步半又三分之一步，问长是多少？

}

\bianswer{%

一百三十步一十一分步之一十。

}{%

一百三十步又十一分之十步。

}



\probheader{少广三}

\biquestion{%

田一亩广一步半，三分步之一，四分步之一，问从？

}{%

有一块田面积一亩，宽一步半又三分之一步又四分之一步，问长是多少？

}

\bianswer{%

一百一十五步五分步之一。

}{%

一百一十五步又五分之一步。

}



\probheader{少广四}

\biquestion{%

田一亩广一步半，三分步之一，四分步之一，五分步之一，问从？

}{%

有一块田面积一亩，宽一步半又三分之一又四分之一又五分之一步，问长是多少？

}

\bianswer{%

一百五步一百三十七分步之一十五。

}{%

一百零五步又一百三十七分之十五步。

}



\probheader{少广五}

\biquestion{%

田一亩广一步半，三分步之一，四分步之一，五分步之一，六分步之一，问从？

}{%

有一块田面积一亩，宽一步半又三分之一又四分之一又五分之一又六分之一步，问长是多少？

}

\bianswer{%

九十七步四十九分步之四十七。

}{%

九十七步又四十九分之四十七步。

}



\probheader{少广六}

\biquestion{%

田一亩广一步半，三分步之一，四分步之一，五分步之一，六分步之一，七分步之一，问从？

}{%

有一块田面积一亩，宽一步半又三分之一又四分之一又五分之一又六分之一又七分之一步，问长是多少？

}

\bianswer{%

九十二步一百二十一分步之六十八。

}{%

九十二步又一百二十一分之六十八步。

}



\newpage

\colheader



\section{分率第四（Fractional Rates）}



\bilingual{%

\textbf{凡九问}。本章九问。

}{%

\textbf{共九道题}。都是分率章的题目。

}



\subsection{乘分法（Multiplication of Fractions）}



\bilingualmethod{%

\textbf{乘分法}曰：分母各乘其全，分子从之，相乘为实，分母相乘为法。

}{%

\textbf{乘分法}：分母各自乘以整数部分，分子跟随着，然后相乘作为被除数（实），分母相乘作为除数（法），实除以法即得结果。（分数乘法：分子乘分子，分母乘分母）

}



\probheader{乘分一}

\biquestion{%

田广一十八步七分步之五，从二十三步十一分步之六，问为田几何？

}{%

有一块田地，宽十八步又七分之五步，长二十三步又十一分之六步，问面积是多少？

}

\bianswer{%

一亩二百步十一分步之七。

}{%

一亩二百平方步又十一分之七平方步。

}



\probheader{乘分二}

\biquestion{%

田广三步三分步之一，从五步五分步之一，问为田几何？

}{%

有一块田地，宽三步又三分之一，长五步又五分之一，问面积是多少？

}

\bianswer{%

一十八步。

}{%

十八平方步。

}



\subsection{除分法（Division of Fractions）}



\bilingualmethod{%

\textbf{除分法}曰：人数为法，钱数为实，有分者通之，实如法而一。

}{%

\textbf{除分法}：将人数作为除数（法），将钱数作为被除数（实），有分数的就通分，实除以法即得每人分得的数量。

}



\probheader{除分一}

\biquestion{%

七人均入钱三分钱之一，问人得几何？

}{%

七个人平分一又三分之一文钱，问每人得多少？

}

\bianswer{%

人得一钱二十一分钱之四。

}{%

每人得一又二十一分之四文钱。

}



\probheader{除分二}

\biquestion{%

三人三分人之一均六钱三分钱之一四分钱之三，问人得几何？

}{%

有三又三分之一个人平分六又三分之一又四分之三文钱，问每人得多少？

}

\bianswer{%

得三钱八分钱之一。

}{%

每人得三又八分之一文钱。

}



\newpage

\colheader



\section{衰分第五（Proportional Distribution）}



\bilingual{%

\textbf{凡十八问}：本章九问，均输九问。

}{%

\textbf{共十八道题}：衰分章九问，均输章九问。

}



\bilingualmethod{%

\textbf{衰分法}曰：各置列衰，副并为法，以所分乘未并者各自为实，实如法而一。

}{%

\textbf{衰分法}：各自布置分配的比率（列衰），另外合并起来作为除数（法），用要分配的总量乘以未合并前的各自比率作为被除数（实），实除以法即得各人分得的数量。（按比例分配）

}



\probheader{衰分一}

\biquestion{%

今有牛二羊五买豕一十三家，有余钱一千；卖牛三豕三以买九羊，钱适足；卖六羊八豕以买五牛，钱不足六百。问牛羊豕价各几何？

}{%

现在有人用二头牛和五只羊去换十三头猪，还余下一千文钱；卖三头牛和三头猪去买九只羊，钱刚好够；卖六只羊和八头猪去买五头牛，还差六百文钱。问牛、羊、猪各值多少钱？

}

\bianswer{%

牛价一千二百，羊价五百，豕价三百。

}{%

牛价一千二百文，羊价五百文，猪价三百文。

}



\probheader{衰分二（五雀六燕）}

\biquestion{%

五雀六燕共重一斤，雀重燕轻，交易一枚，其重适等。问各几何？

}{%

有五只麻雀和六只燕子共重一斤，麻雀重而燕子轻，交换一只后，两边的重量正好相等。问每只麻雀和燕子各重多少？

}

\bianswer{%

雀一两十九分两之一十三，燕一两十九两之五。

}{%

每只麻雀重一两又十九分之十三两，每只燕子重一两又十九分之五两。

}



\probheader{均输一}

\biquestion{%

五县均赋粟一万石，每车载二十五石，行道一里出雇钱一文，各县到输所远近不等，粟价高下，欲令县劳收相等。内甲县二万五百二十户，粟一石价钱二十，自输其县；乙县一万二千三百一十二户，粟一石价钱一十，远输所二百里；丙县七千一百八十二户，粟一石价钱一十二，远输所一百五十里；丁县一万三千三百三十八户，粟一石价钱一十七，远输所二百五十里；戊县五千一百三十户，粟一石价钱一十三，远输所一百五十里。问各几何？

}{%

五个县共同分摊缴纳一万石粟米的赋税，每辆车装载二十五石，每走一里路需要出一文钱的运费，各县到缴纳地点的远近不同，粟米价格也有高低，现在要使各县付出的劳力和收到的粟米价值相等。其中甲县有二万零五百二十户，粟一石值二十文钱，在本县缴纳；乙县有一万二千三百一十二户，粟一石值十文钱，距离缴纳地二百里；丙县有七千一百八十二户，粟一石值十二文钱，距离缴纳地一百五十里；丁县有一万三千三百三十八户，粟一石值十七文钱，距离缴纳地二百五十里；戊县有五千一百三十户，粟一石值十三文钱，距离缴纳地一百五十里。问各县应缴纳多少？

}



\newpage

\colheader



\section{叠积第六（Stacked Accumulations）}



\bilingual{%

\textbf{凡二十七问}，并商功章。

}{%

\textbf{共二十七道题}，都是商功章的内容。

}



\bilingualmethod{%

\textbf{商功法}曰：各以其积乘之，以高若深乘之，皆如六而一。

}{%

\textbf{商功法}：各自用底面积相乘，再乘以高或深，然后都除以六。（楔形体体积公式 $V = \frac{1}{6} \times$ 底面积 $\times$ 高）

}



\probheader{盘池}

\biquestion{%

盘池上广六丈，袤八丈，下广四丈，袤六丈，深二丈。问积尺几何？

}{%

有一个盘形水池，上底宽六丈，长八丈，下底宽四丈，长六丈，深二丈。问容积是多少立方尺？

}

\bianswer{%

七万六百六十六尺太半尺。

}{%

七万零六百六十六又三分之二立方尺。

}



\probheader{冥谷}

\biquestion{%

冥谷上广二丈，袤七丈，下广八尺，袤四丈，深六丈五尺。问积几何？

}{%

有一个深谷，上底宽二丈，长七丈，下底宽八尺，长四丈，深六丈五尺。问容积是多少？

}

\bianswer{%

五万二千尺。

}{%

五万二千立方尺。

}



\bilingualmethod{%

\textbf{刍甍法}曰：倍下长并入上长，以广乘之，又高乘之，皆如六而一。



下广三丈，袤四丈，上袤二丈，无广，高一丈。问积几何？

}{%

\textbf{刍甍法}：将下底长度加倍加上上底长度，乘以宽度，再乘以高，然后都除以六。



有一个屋脊形的草垛，下底宽三丈，长四丈，上底长二丈，没有宽度（上底为线），高一丈。问体积是多少？

}

\bianswer{%

五千尺。

}{%

五千立方尺。

}



\bilingualmethod{%

\textbf{羡除法}曰：并三广以深乘之，又长乘之，皆如六而一。



上广一丈，下广六尺，深三尺，末广八尺，无深，袤七尺。问积几何？

}{%

\textbf{羡除法}：将三个宽度相加，乘以深度，再乘以长度，然后都除以六。



有一个隧道形体，上宽一丈，下宽六尺，深三尺，末端宽八尺，没有深度，长七尺。问体积是多少？

}

\bianswer{%

五千尺。

}{%

五千立方尺。

}



\newpage

\colheader



\section{盈不足第七（Excess and Deficit）}



\bilingual{%

\textbf{凡十一问}，并本章。

}{%

\textbf{共十一道题}，都是盈不足章的内容。

}



\subsection{盈朒适足法（Excess/Deficit/Exact Method）}



\bilingualmethod{%

\textbf{盈朒适足法}曰：置所出率，盈朒各居其下。副置出率，以少减多余为约法。盈朒适足，令维乘所出率实，以盈朒之数乘人积为人实，实皆如约法而一。



其一法曰：以盈或不足之数为实，置所出率，以少减多余为法，实如法而一约人，以适足出率乘人，为物价也。

}{%

\textbf{盈不足适足法}：布置每人出的钱数（出率），盈余或不足数各放在下面。另外布置出率，用大的减去小的，余数作为约法。对于盈不足适足的情况，交叉相乘所出率与实，用盈或不足的数乘以人数积作为人实，各实都除以约法即得。



另一种方法：用盈余或不足的数作为被除数（实），布置所出率，用大的减去小的余数作为除数（法），实除以法得到人数，再用适足的出率乘以人数，就得到物价。

}



\probheader{盈不足一（共买犬）}

\biquestion{%

共买犬，人出五不足九十文，人出五十文适足。问人数、犬价各几何？

}{%

几个人合伙买一条狗，每人出五文钱，还差九十文；每人出五十文，正好够。问有多少人和狗的价钱是多少？

}

\bianswer{%

二人，犬价一百。

}{%

二人，狗的价钱一百文。

}



\probheader{盈不足二（买豕）}

\biquestion{%

买豕，人各出一百盈一百文，各出九十适足。问人数、豕价各几何？

}{%

几个人合伙买一头猪，每人出一百文，多出一百文；每人出九十文，正好够。问有多少人和猪的价钱是多少？

}

\bianswer{%

十人，豕价九百。

}{%

十人，猪的价钱九百文。

}



\subsection{双假设法（Double False Position）}



\probheader{盈不足三}

\biquestion{%

今有共买璡，人出半盈四，人出少半不足三。问人数、璡价各几何？

}{%

现在几个人合伙买璡（一种玉器），每人出一半钱，多四文；每人出少于一半的钱，少三文。问人数和玉价各是多少？

}

\bianswer{%

四十二人，璡价十七。

}{%

四十二人，玉价十七文。

}



\probheader{盈不足四（大小器）}

\biquestion{%

今有大器五小器一容三斛，大器一小器五容二斛。问大小器各容几何？

}{%

现在有大容器五个和小容器一个共容三斛，大容器一个和小容器五个共容二斛。问大小容器各能容多少？

}

\bianswer{%

大器容二十四分解之一十三，小器容二十四分解之七。

}{%

大容器容二十四分之十三斛，小容器容二十四分之七斛。

}



\probheader{盈不足五（二酒求价）}

\biquestion{%

今有醇酒一斗直钱五十，行酒一斗直钱一十。今将钱三十，得酒二斗。问醇行酒各几何？

}{%

现在有醇酒（好酒）一斗值五十文，行酒（普通酒）一斗值十文。现在用三十文钱买酒二斗。问醇酒和行酒各买了多少？

}

\bianswer{%

醇酒二升半，行酒一斗七升半。

}{%

醇酒二升半，行酒一斗七升半。

}



\probheader{盈不足六（五家共井）}

\biquestion{%

今有五家共井，甲二绠不足如乙一，乙三绠不足如丙一，丙四绠不足如丁一，丁五绠不足如戊一，戊六绠不足如甲一。如各得所不足一绠，皆逮。问井深、绠长各几何？

}{%

现在有五家人共用一口井，甲家的绳二条还不够乙家一条那么长，乙家的绳三条还不够丙家一条那么长，丙家的绳四条还不够丁家一条那么长，丁家的绳五条还不够戊家一条那么长，戊家的绳六条还不够甲家一条那么长。如果各家都得到所缺的一条绳长，就都能打到水。问井的深度和各家的绳长各是多少？

}



\newpage

\colheader



\section{方程第八（Simultaneous Equations）}



\bilingual{%

\textbf{凡二十问}：本章十八问，均输盈朒二。

}{%

\textbf{共二十道题}：方程章十八问，均输章和盈不足章各一问。

}



\subsection{方程法（Method of Simultaneous Equations）}



\bilingualmethod{%

\textbf{方程法}曰：所求率互乘邻行，以少减多，再求减损。钱为实，物为法，实如法而一。



置位法曰：依所问排列逐物与价，而邻行相对如之。

}{%

\textbf{方程法}：用所求的比率交叉乘以相邻的行，用少的减去多的，再次求减损。钱数作为被除数（实），物品作为除数（法），实除以法即得各物的单价。



布列方法：按照问题排列各种物品和价格，相邻的行互相对应排列。（线性方程组的增广矩阵布列法）

}



\probheader{方程一（上中下禾）}

\biquestion{%

上禾三束，中禾二束，下禾一束，共实三十九斗。上禾二束，中禾三束，下禾一束，共实三十四斗。上禾一束，中禾二束，下禾三束，共实二十六斗。问：上、中、下禾一束各几何？

}{%

上等禾三束，中等禾二束，下等禾一束，共出谷三十九斗。上等禾二束，中等禾三束，下等禾一束，共出谷三十四斗。上等禾一束，中等禾二束，下等禾三束，共出谷二十六斗。问：上、中、下等禾每束各出多少斗谷？

}

\bianswer{%

上禾一束，得九斗四分斗之一；中禾一束，得四斗四分斗之一；下禾一束，得二斗四分斗之三。

}{%

上等禾每束出九又四分之一斗；中等禾每束出四又四分之一斗；下等禾每束出二又四分之三斗。

}



\probheader{方程二（牛羊价）}

\biquestion{%

五牛二羊，直金十两；二牛五羊，直金八两。问：牛、羊价各几何？

}{%

五头牛二只羊，值十两金子；二头牛五只羊，值八两金子。问：牛和羊各值多少？

}

\bianswer{%

一牛直金一两二十一分两之一十三；一羊直金二十一两分之二十。

}{%

一头牛值一两又二十一分之十三两金子；一只羊值二十一分之二十两金子。

}



\probheader{方程三}

\biquestion{%

上禾二秉，中禾三秉，下禾四秉，实皆不满斗。上取中，中取下，下取上，各一秉，而实满斗。问：上、中、下禾一秉，实各几何？

}{%

上等禾二秉，中等禾三秉，下等禾四秉，出谷都不满一斗。上等禾从中等禾取，中等禾从下等禾取，下等禾从上等禾取，各取一秉后，出谷正好满一斗。问：上、中、下等禾每秉各出多少斗谷？

}

\bianswer{%

上禾一秉，实二十五分斗之九；中禾一秉，实二十五分斗之七；下禾一秉，实二十五分斗之四。

}{%

上等禾每秉出二十五分之九斗；中等禾每秉出二十五分之七斗；下等禾每秉出二十五分之四斗。

}



\subsection{正负术（Positive and Negative Method）}



\bilingualmethod{%

\textbf{正负术}曰：同名相除，异名相益，正无入负之，负无入正之。其异名相除，同名相益，正无入正之，负无入负之。

}{%

\textbf{正负术}：同号的两数相减（绝对值相减），异号的两数相加（绝对值相加），正数没有对应数时就加上负号，负数没有对应数时就加上正号。另一种情况：异号的两数相除，同号的两数相乘，正数没有对应数时就记为正，负数没有对应数时就记为负。

}



\probheader{方程四（卖牛买羊）}

\biquestion{%

今有卖牛二、羊五，以买一十三豕，有余钱一千；卖牛三、豕三，以买九羊，钱适足；卖六羊、八豕，以买五牛，钱不足六百。问：牛、羊、豕价各几何？

}{%

现在卖二头牛和五只羊去买十三头猪，余一千文钱；卖三头牛和三头猪去买九只羊，钱刚好够；卖六只羊和八头猪去买五头牛，差六百文钱。问：牛、羊、猪各值多少钱？

}

\bianswer{%

牛价一千二百，羊价五百，豕价三百。

}{%

牛价一千二百文，羊价五百文，猪价三百文。

}



\probheader{方程五（金银易重）}

\biquestion{%

今有黄金九，白银一十一，称之重适等，交易其一，金轻十三两。问：金、银一枚各重几何？

}{%

现在有黄金九枚，白银十一枚，称起来重量正好相等。交换一枚后，黄金这边轻了十三两。问：金、银每枚各重多少？

}



\newpage

\colheader



\section{勾股第九（Right Triangles）}



\bilingual{%

\textbf{凡三十八问}：少广十三，商功一问，本章二十四。

}{%

\textbf{共三十八道题}：少广章十三问，商功章一问，勾股章二十四问。

}



\subsection{勾股基本方法（Basic Pythagorean Methods）}



\bilingualmethod{%

\textbf{勾股法}曰：勾自乘，股自乘，并之为弦实，开方除之即弦。其勾股自乘为实，如股弦较而一，加较半之即弦。

}{%

\textbf{勾股法}：勾（短直角边）自乘，股（长直角边）自乘，合并起来作为弦（斜边）的平方数，开平方即得弦长。另外，勾股各自平方作为被除数，除以股弦差再加差的一半即得弦长。（勾股定理 $c = \sqrt{a^2 + b^2}$）

}



\probheader{勾股一}

\biquestion{%

其三，勾自乘为实，如股弦较而一，加较半之即弦。



立木垂索，委地二尺，引索斜之，拄地去木八尺。问索长几何？

}{%

第三种方法：勾自乘作为被除数，除以股弦差，再加差的一半即得弦。



有一根竖立的木杆垂着绳索，绳拖在地上二尺长，拉紧绳索斜着拉，绳索着地的地方距离木杆八尺。问绳索有多长？

}

\bianswer{%

一十七尺。

}{%

十七尺。

}



\probheader{勾股二}

\biquestion{%

其四，勾自乘为实，如股弦较而一，以较减之余半之即股。



垣高一丈，倚木齐垣，木脚去本以画记之，卧而过画一尺。问去本几何？

}{%

第四种方法：勾自乘作为被除数，除以股弦差，用差去减余数再取半即得股。



有一道高一丈的墙，把一根木头靠着墙竖立，在木脚离地的地方做记号，然后放倒木头，记号超过地面一尺。问木脚离墙根多远？

}

\bianswer{%

四丈九尺半。

}{%

四丈九尺半。

}



\probheader{勾股三}

\biquestion{%

其五，半勾自乘为实，如半股弦较而一，加半较即弦。



圆材埋在壁中，不知大小，锯深一寸，道長一尺。问径几何？

}{%

第五种方法：半勾自乘作为被除数，除以半股弦差，加半差即得弦。



有一段圆形木材埋在墙壁中，不知道大小，锯进去一寸深，锯道长一尺。问圆木的直径是多少？

}



\subsection{户高广问题（Door Height and Width）}



\probheader{勾股四（户高广）}

\biquestion{%

开门去阃一尺，不合二寸。问门广几何？

}{%

打开门后，门扇离门槛一尺，门扇之间有缝隙不合二寸。问门的宽度是多少？

}

\bianswer{%

一丈一寸五分。

}{%

一丈一寸五分。

}



\subsection{勾股容方（Square Inscribed in Right Triangle）}



\bilingualmethod{%

\textbf{合率}：勾中容方。



今有勾五步，股十二步。问：勾中容方几何？

}{%

\textbf{合率问题}：直角三角形中内接正方形的边长计算。



现在有一个直角三角形，勾（短直角边）五步，股（长直角边）十二步。问：在这个直角三角形中能容纳的最大正方形的边长是多少？

}

\bianswer{%

方三步十七分步之九。

}{%

正方形边长为三又十七分之九步。

}



\newpage

\colheader



\subsection{复杂勾股问题（Complex Pythagorean Problems）}



\probheader{勾股五（竹折）}

\biquestion{%

今有竹高一丈，末折抵地，去本三尺。问：折者高几何？

}{%

现在有一根竹子高一丈，顶端折断后触地，竹尖离根部三尺远。问：折断处离地有多高？

}

\bianswer{%

四尺五十五分尺之十一。

}{%

四尺又五十五分之十一尺。

}



\probheader{勾股六（葭生中央）}

\biquestion{%

今有池方一丈，葭生其中央，出水一尺。引葭赴岸，适与岸齐。问：水深、葭长各几何？

}{%

现在有一个边长一丈的方形水池，芦苇生长在水池的中央，露出水面一尺。把芦苇拉到岸边，正好与岸边齐平。问：水深和芦苇长度各是多少？

}



\probheader{勾股七（邑方）}

\biquestion{%

今有邑方不知大小，各开中门。出北门二十步有木，出南门十四步折而西行一千七百七十五步见木。问：邑方几何？

}{%

现在有一座方形城邑，不知道大小，各边开有中门。出北门走二十步有一棵树，出南门走十四步然后向西折行一千七百七十五步能看到那棵树。问：城邑的边长是多少？

}



\newpage

\colheader



%% ===== APPENDIX =====

\section*{附录：九大分类门类说明}

\addcontentsline{toc}{section}{附录：九大分类门类说明}



\bilingual{%

杨辉《纂类》将《九章算术》二百四十六问按数学方法重新分为九大门类：



\begin{enumerate}[leftmargin=*]

\item \textbf{乘除} —— 乘法与除法：四十一道，涉及基本田地测量与单位换算。

\item \textbf{分率} —— 分数运算：九道，涉及分数乘法、除法与比较。

\item \textbf{合率} —— 合率运算：二十道，涉及分数合并运算，包括"少广"法。

\item \textbf{互换} —— 交换与贸易：五十六道，涉及物品比例交换，包括谷物换算与税务问题。

\item \textbf{衰分} —— 比例分配：十八道，涉及按比例分配数量，包括"均输"法。

\item \textbf{叠积} —— 体积计算：二十七道，涉及土方工程与其他三维体积计算。

\item \textbf{盈不足} —— 盈不足术：十一道，用双假设法（试位法）求解。

\item \textbf{方程} —— 线性方程组：二十道，涉及多元一次方程组，用算板矩阵法求解。

\item \textbf{勾股} —— 勾股定理：三十八道，应用勾股定理及相关几何方法。

\end{enumerate}

}{%

杨辉的《纂类》将《九章算术》中的二百四十六道数学题按照数学方法重新组织为九个大类：



\begin{enumerate}[leftmargin=*]

\item \textbf{乘除} —— 乘法和除法：41道题，关于基本田地面积测量和单位换算。

\item \textbf{分率} —— 分数运算：9道题，关于分数的乘法、除法和比较大小。

\item \textbf{合率} —— 合率运算：20道题，关于分数的合并计算，包括"少广"法。

\item \textbf{互换} —— 交换与易货贸易：56道题，关于物品的比例交换，包括谷物换算和赋税问题。

\item \textbf{衰分} —— 按比例分配：18道题，关于按一定比例分配数量，包括"均输"法。

\item \textbf{叠积} —— 堆叠体积计算：27道题，关于土方工程和其他三维立体体积的计算。

\item \textbf{盈不足} —— 盈亏问题：11道题，用双假设法（试位法）来求解。

\item \textbf{方程} —— 线性方程组：20道题，关于多元一次方程组，用算板上的矩阵方法来求解。

\item \textbf{勾股} —— 直角三角形：38道题，应用勾股定理和相关的几何计算方法。

\end{enumerate}

}



\vspace{1cm}

\longline



\vspace{1cm}



\bilingual{%

\noindent\textit{此排版本再现杨辉《详解九章算法·纂类》（南宋景定二年，1261年）之内容。原著将《九章算术》全部二百四十六问按数学方法重新分类，开创了按方法论组织数学问题的教学新体例，堪称现代数学分类教学法之先声。}

}{%

\noindent\textit{本排版版再现杨辉《详解九章算法·纂类》的内容，原著成书于南宋景定二年（1261年）。原著将《九章算术》中全部二百四十六道数学题按照其内在的数学方法重新组织分类，创造了一种具有教学创新性的结构，可谓现代数学分类教学法的先声。}

}



\vspace{1cm}



\begin{flushright}

\textit{—— 第三部（纂类）终 ——}

\end{flushright}



\end{document}







% ============================================================

% Source: translations/zh/chinese/zhu-shijie-suanxue-qimeng-part1_classical-modern_bilingual.tex

% ============================================================



%% ========================================================================

%% 朱世傑《算學啓蒙》卷上 —— 文言文/白話文對照版

%% Zhu Shijie: Suanxue Qimeng (Introduction to Mathematical Studies) — Part I

%% Classical Chinese → Modern Chinese Bilingual Edition

%% ========================================================================

\documentclass[11pt,a4paper]{article}



%% -------- Core Packages --------

\usepackage{fontspec}

\usepackage{polyglossia}

\usepackage{amsmath,amssymb}

\usepackage{geometry}

\usepackage{graphicx}

\usepackage{xcolor}

\usepackage{titlesec}

\usepackage{fancyhdr}

\usepackage{enumitem}

\usepackage{array,booktabs,longtable}

\usepackage{hyperref}

\usepackage{xeCJK}

\usepackage{paracol}   % For parallel bilingual columns

\usepackage{ulem}      % For underlining



%% -------- Page Geometry --------

\geometry{

  paperwidth=210mm,paperheight=297mm,

  left=18mm,right=18mm,top=25mm,bottom=25mm,

  headheight=14pt,footskip=30pt

}



%% -------- Column Ratio --------

\columnratio{0.50}   % Equal width columns

\setlength{\columnsep}{1.2em}

\setlength{\columnseprule}{0.4pt}



%% -------- Language Setup --------

\setmainlanguage{english}



%% -------- Font Configuration --------

\setmainfont{Times New Roman}

\setsansfont{Arial}

\setmonofont{Courier New}



%% CJK Support via xeCJK

\setCJKmainfont{Microsoft YaHei}

\setCJKsansfont{Microsoft YaHei}

\setCJKmonofont{Microsoft YaHei}



%% -------- Colors --------

\definecolor{titlecolor}{RGB}{45,55,72}

\definecolor{sectioncolor}{RGB}{55,65,81}

\definecolor{notecolor}{RGB}{120,53,15}

\definecolor{verselinecolor}{RGB}{80,60,40}

\definecolor{leftbg}{RGB}{255,253,245}     % warm ivory for classical

\definecolor{rightbg}{RGB}{245,250,255}    % cool blue-white for modern

\definecolor{leftlabel}{RGB}{139,69,19}    % brown for classical

\definecolor{rightlabel}{RGB}{25,85,140}   % blue for modern

\definecolor{problembg}{RGB}{250,248,245}



%% -------- Section Formatting --------

\titleformat{\section}

  {\normalfont\Large\bfseries\color{sectioncolor}}

  {\thesection}{1em}{}

\titleformat{\subsection}

  {\normalfont\large\bfseries\color{sectioncolor}}

  {\thesubsection}{1em}{}

\titleformat{\subsubsection}

  {\normalfont\normalsize\bfseries\color{sectioncolor}}

  {\thesubsubsection}{1em}{}



%% -------- Header/Footer --------

\pagestyle{fancy}

\fancyhf{}

\fancyhead[L]{\small\textcolor{leftlabel}{\textit{算學啓蒙·文言文}}\hspace{1em}|\hspace{1em}\textcolor{rightlabel}{\textit{白話文對照版}}}

\fancyhead[R]{\small\thepage}

\renewcommand{\headrulewidth}{0.4pt}

\renewcommand{\footrulewidth}{0pt}



%% -------- Custom Commands --------

\newcommand{\longline}{\noindent\rule{\textwidth}{0.4pt}}

\newcommand{\shortline}{\noindent\rule{0.5\textwidth}{0.4pt}\par}

\newcommand{\cnote}[1]{\textcolor{notecolor}{\textit{#1}}}



%% Problem number counter

\newcounter{problemnum}



%% Left column (Classical) problem formatting

\newcommand{\problemCL}[3]{%

  \refstepcounter{problemnum}%

  \par\medskip\noindent

  \textbf{\textcolor{leftlabel}{問\,\theproblemnum.}}\ #1

  \par\smallskip\noindent

  \textbf{答：}#2

  \par\smallskip\noindent

  \textbf{術：}#3

  \par\medskip

}



%% Right column (Modern) problem formatting

\newcommand{\problemMOD}[3]{%

  \par\medskip\noindent

  \textbf{\textcolor{rightlabel}{【問題\,\theproblemnum】}}\ #1

  \par\smallskip\noindent

  \textbf{【答案】}#2

  \par\smallskip\noindent

  \textbf{【解法】}#3

  \par\medskip

}



%% Verse/rhyme formatting for classical

\newcommand{\rhymeCL}[1]{%

  \begin{quote}\itshape\color{verselinecolor}#1\end{quote}

}



%% Verse/rhyme formatting for modern

\newcommand{\rhymeMOD}[1]{%

  \begin{quote}\itshape\color{verselinecolor}#1\end{quote}

}



%% Note in left column

\newcommand{\noteCL}[1]{\par\textcolor{notecolor}{\small\textit{〔注〕#1}}\par\smallskip}

%% Note in right column

\newcommand{\noteMOD}[1]{\par\textcolor{notecolor}{\small\textit{〔译注〕#1}}\par\smallskip}



%% Inline notes

\newcommand{\inote}[1]{\textcolor{notecolor}{\small（#1）}}



%% -------- Hyperref Setup --------

\hypersetup{

  colorlinks=true,

  linkcolor=sectioncolor,

  citecolor=sectioncolor,

  urlcolor=sectioncolor,

  pdfencoding=auto,

  pdftitle={算學啓蒙·卷上 —— 文言文/白話文對照版},

  pdfauthor={朱世傑（原著）}}



%% ========================================================================

\begin{document}



%% -------- Title Page --------

\begin{titlepage}

\begin{center}

\vspace*{2cm}



{\fontsize{24}{32}\selectfont\color{titlecolor}

  新編\!算學\!啓蒙\!·\!卷上}



\vspace{0.5cm}



{\fontsize{16}{22}\selectfont\color{sectioncolor}

  文言文 / 白話文 對照版}



\vspace{1.5cm}



{\Large\color{sectioncolor}

  Classical Chinese / Modern Chinese Bilingual Edition}



\vspace{2.5cm}



{\LARGE\color{titlecolor} 總括 · 卷上}



\vspace{0.5cm}



{\large General Principles \& Volume I (113 Problems)}



\vspace{2.5cm}



{\large 元 · 朱世傑撰}



\vspace{0.3cm}



{\normalsize\color{sectioncolor}

  By Zhu Shijie (c.\ 1260--1320 CE)}



\vspace{0.2cm}



{\small\color{sectioncolor}

  Yuan Dynasty, written in 1299 (大德己亥)}



\vspace{2cm}



{\small\color{notecolor}

  對照排版：左欄為原文（文言文），右欄為白話文翻譯\\

  Modern LaTeX edition\ using paracol (parallel columns)\\

  Left column: Original Classical Chinese \quad Right column: Modern Chinese Vernacular}



\end{center}

\end{titlepage}



\tableofcontents

\newpage



%% ========================================================================

\section*{前言 · Preface}

\addcontentsline{toc}{section}{前言 · Preface}

\longline



\begin{paracol}{2}

\noindent\textbf{\textcolor{leftlabel}{【文言】}}

\switchcolumn

\noindent\textbf{\textcolor{rightlabel}{【白話】}}

\switchcolumn*



《算學啓蒙》為元代人朱世傑所撰，成書於大德己亥年（1299年）。凡三卷，計二十門，二百五十九問。其書由淺入深，自乘除加減之法，至天元術、垛積術、招差術，條理秩然，為初學算學之津梁也。



\switchcolumn



《算學啓蒙》是元朝数学家朱世杰编写的数学启蒙教材，成书于公元1299年（大德三年）。全书共三卷二十章，有259道数学题。内容由浅入深，从基本的乘除加减运算，到高深的「天元术」（解方程）、「垛积术」（级数求和）、「招差术」（有限差分），条理清晰，是学习数学的入门经典。



\switchcolumn*



是書首列「總括」一十八條，為全書之綱領，包括九數乘法、九歸除法、斤兩留法、縱橫訣、大小數名、諸率換算、圓率、正負術、開方法等。卷上則設八門，凡一百一十三問，論縱橫因法、身外加減、留頭乘法、九歸除法、異乘同除、庫務解稅、折變互差諸術。



\switchcolumn



本书开头列有「总括」十八条，是全书的纲领性内容，包括九九乘法口诀、九归除法口诀、斤两换算、算筹记数诀、大小数名称、各种单位换算率、圆周率、正负数运算法则、开方法等。卷上设有八章，共113道题，讨论纵横因法、身外加减法、留头乘法、九归除法、异乘同除、库务解税、折变互差等算法。



\switchcolumn*



其題多關乎民生實務——租稅錢糧、物價交易、田畝丈量、軍資調度，皆元代社會經濟之寫照也。此書在中國久佚，賴朝鮮刻本傳世，清羅士琳於揚州重刻（1839年），遂復行於世。



\switchcolumn



书中题目多涉及实际生活——租税钱粮、物价交易、田亩丈量、军需物资调度等，都是元代社会经济的真实写照。此书在中国失传已久，幸亏有朝鲜刻本保存了下来；清代数学家罗士琳在扬州重新刊刻（1839年），才得以重新流传。



\end{paracol}



\longline



%% ========================================================================

\newpage

\section{總括 · General Principles}

\label{sec:zongguo}

\addcontentsline{toc}{section}{總括 · General Principles}

%% ========================================================================



\begin{paracol}{2}

\noindent\textbf{\textcolor{leftlabel}{【文言】}}

\switchcolumn

\noindent\textbf{\textcolor{rightlabel}{【白話】}}

\switchcolumn*



總括者，列全書之基礎知識也。凡一十八條，包括乘法口訣、除法歌訣、單位換算、記數法則、代數規約。學者熟習之，方可進習九章之術。



\switchcolumn



「总括」部分列出了全书所需的基础知识，共十八条，包括乘法口诀、除法歌诀、单位换算、记数法则、代数运算规定。学习者必须熟练掌握这些内容，才能进一步学习九章算术中的各种算法。



\end{paracol}



\medskip



%% ------------------------------------------------------------------------

\subsection{釋九數法 · 九九乘法表}

%% ------------------------------------------------------------------------



\begin{paracol}{2}

\noindent\textbf{\textcolor{leftlabel}{【乘法口訣】}}

\switchcolumn

\noindent\textbf{\textcolor{rightlabel}{【乘法口诀（九九表）】}}

\switchcolumn*



一一如一，一二如二，二二如四。\break

一三如三，二三如六，三三如九。\break

一四如四，二四如八，三四一十二，四四一十六。\break

一五如五，二五一十，三五一十五，四五二十，五五二十五。\break

一六如六，二六一十二，三六一十八，四六二十四，五六三十，六六三十六。\break

一七如七，二七一十四，三七二十一，四七二十八，五七三十五，六七四十二，七七四十九。\break

一八如八，二八一十六，三八二十四，四八三十二，五八四十，六八四十八，七八五十六，八八六十四。\break

一九如九，二九一十八，三九二十七，四九三十六，五九四十五，六九五十四，七九六十三，八九七十二，九九八十一。



\switchcolumn



1×1=1，1×2=2，2×2=4。\break

1×3=3，2×3=6，3×3=9。\break

1×4=4，2×4=8，3×4=12，4×4=16。\break

1×5=5，2×5=10，3×5=15，4×5=20，5×5=25。\break

1×6=6，2×6=12，3×6=18，4×6=24，5×6=30，6×6=36。\break

1×7=7，2×7=14，3×7=21，4×7=28，5×7=35，6×7=42，7×7=49。\break

1×8=8，2×8=16，3×8=24，4×8=32，5×8=40，6×8=48，7×8=56，8×8=64。\break

1×9=9，2×9=18，3×9=27，4×9=36，5×9=45，6×9=54，7×9=63，8×9=72，9×9=81。



\end{paracol}



\medskip



%% ------------------------------------------------------------------------

\subsection{九歸除法 · 九归除法（除法口诀）}

%% ------------------------------------------------------------------------



\begin{paracol}{2}

\noindent\textbf{\textcolor{leftlabel}{【原文】}}

\switchcolumn

\noindent\textbf{\textcolor{rightlabel}{【白话翻译】}}

\switchcolumn*



\textit{按古法多用商除。爲初學者難入，則後人以此法代之，卽非正術也。}



\medskip



一歸如一進，見一進成十。\break

二一添作五，逢二進成十。\break

三一三十一，三二六十二，逢三進成十。\break

四一二十二，四二添作五，四三七十二，逢四進成十。\break

五歸添一倍，逢五進成十。\break

六一下加四，六二三十二，六三添作五，六四六十四，六五八十二，逢六進成十。\break

七一下加三，七二下加六，七三四十二，七四五十五，七五七十一，七六八十四，逢七進成十。\break

八一下加二，八二下加四，八三下加六，八四添作五，八五六十二，八六七十四，八七八十六，逢八進成十。\break

九歸隨身下，逢九進成十。



\switchcolumn



\textit{注：按照古代的方法多用商除法。但对初学者来说太难入门了，所以后人用九归除法来代替，这虽然并不是正统的除法。}



\medskip



「一归」：1÷1 商1进位到十位（即10÷1=10）。\break

「二归」：1÷2 改作5（即10÷2=5）；逢2进1（即够2除就在上位进1）。\break

「三归」：1÷3 商3余1（即10÷3=3余1）；2÷3 商6余2；逢3进1。\break

「四归」：1÷4 商2余2；2÷4 改作5；3÷4 商7余2；逢4进1。\break

「五归」：添一倍（即乘以2）；逢5进1。\break

「六归」：1÷6 本位不动下位加4；2÷6 商3余2；3÷6 改作5；4÷6 商6余4；5÷6 商8余2；逢6进1。\break

「七归」：1÷7 本位不动下位加3；2÷7 本位不动下位加6；3÷7 商4余2；4÷7 商5余5；5÷7 商7余1；6÷7 商8余4；逢7进1。\break

「八归」：1÷8 本位不动下位加2；2÷8 本位不动下位加4；3÷8 本位不动下位加6；4÷8 改作5；5÷8 商6余2；6÷8 商7余4；7÷8 商8余6；逢8进1。\break

「九归」：随被除数落位（即商与被除数相同）；逢9进1。



\end{paracol}



\medskip



%% ------------------------------------------------------------------------

\subsection{斤下留法 · 斤两换算（16两为1斤）}

%% ------------------------------------------------------------------------



\begin{paracol}{2}

\noindent\textbf{\textcolor{leftlabel}{【原文】}}

\switchcolumn

\noindent\textcolor{rightlabel}{\textbf{【白话翻译】}}

\switchcolumn*



\textit{斤下帶兩者，當以十六約之。今則就省，以此代之也。}



\medskip



一退六二五，二留一二五，三留一八七五，四留二五，五留三一二五，六留三七五，七留四三七五，八留單五，九留五六二五，十留六二五，十一留六八七五，十二留七五，十三留八一二五，十四留八七五，十五留九三七五。



\switchcolumn



\textit{注：一斤以下带有两数时，应当以16来除（因为16两=1斤）。为了简便计算，就用下面的口诀来代替除法。}



\medskip



1两 = 0.0625斤\quad

2两 = 0.125斤\quad

3两 = 0.1875斤\quad

4两 = 0.25斤\quad

5两 = 0.3125斤\quad

6两 = 0.375斤\quad

7两 = 0.4375斤\quad

8两 = 0.5斤\quad

9两 = 0.5625斤\quad

10两 = 0.625斤\quad

11两 = 0.6875斤\quad

12两 = 0.75斤\quad

13两 = 0.8125斤\quad

14两 = 0.875斤\quad

15两 = 0.9375斤



\end{paracol}



\medskip



%% ------------------------------------------------------------------------

\subsection{明縱橫訣 · 算筹记数口诀}

%% ------------------------------------------------------------------------



\begin{paracol}{2}

\noindent\textbf{\textcolor{leftlabel}{【原文】}}

\switchcolumn

\noindent\textbf{\textcolor{rightlabel}{【白话翻译】}}

\switchcolumn*



一縱十橫，百立千僵。\break

千十相望，萬百相當。\break

滿六已上，五在上方。\break

六不積聚，五不單張。\break

言十自過，不滿自當。\break

若明此訣，可習九章。



\switchcolumn



个位用纵式，十位用横式，\break

百位用纵式，千位用横式。\break

千和十的摆法相对应，\break

万和百的摆法相对应。\break

数字达到六以上，就把「五」放在上方。\break

六以上不用逐个累积，五也不单独摆放。\break

说到「十」就要进位，不够「十」就在本位定数。\break

如果明白这个口诀，就可以学习《九章算术》了。



\noteMOD{注：这是描述中国古代算筹记数法的规则。个位数用纵式（竖摆），十位数用横式（横摆），百位纵式，千位横式，纵横交替。数字五用一根横放的算筹表示，放在对应数位的上方。六到九的表示方法是在上方放一根代表五的横筹，下方再放表示余数的算筹（如六=上五+下一）。}



\end{paracol}



\medskip



%% ------------------------------------------------------------------------

\subsection{大數之類 · 大数名称}

%% ------------------------------------------------------------------------



\begin{paracol}{2}

\noindent\textbf{\textcolor{leftlabel}{【原文】}}

\switchcolumn

\noindent\textbf{\textcolor{rightlabel}{【白话翻译】}}

\switchcolumn*



\textit{凡數之大者，天莫能蓋、地莫能載，其數不能極，故謂之大數也。}



\medskip



一、十、百、千、萬、十萬、百萬、千萬，萬萬曰億，萬萬億曰兆，萬萬兆曰京，萬萬京曰陔，萬萬陔曰秭，萬萬秭曰壤，萬萬壤曰溝，萬萬溝曰澗，萬萬澗曰正，萬萬正曰載，萬萬載曰極，萬萬極曰恆河沙，萬萬恆河沙曰阿僧祗，萬萬阿僧祗曰那由他，萬萬那由他曰不可思議，萬萬不可思議曰無量數。



\switchcolumn



\textit{注：说到极大的数，连天都覆盖不了、地都承载不下，这样的数没有极限，所以称为「大数」。}



\medskip



计数单位从小到大依次为：个（一）、十、百、千、万、十万、百万、千万。

万万个亿（10\textsuperscript{8}），万万个亿叫「兆」（10\textsuperscript{16}），

万万个兆叫「京」（10\textsuperscript{24}），万万个京叫「陔」（10\textsuperscript{32}），

万万个陔叫「秭」（10\textsuperscript{40}），万万个秭叫「壤」（10\textsuperscript{48}），

万万个壤叫「沟」（10\textsuperscript{56}），万万个沟叫「涧」（10\textsuperscript{64}），

万万个涧叫「正」（10\textsuperscript{72}），万万个正叫「载」（10\textsuperscript{80}），

万万个载叫「极」（10\textsuperscript{88}），

万万个极叫「恒河沙」（10\textsuperscript{96}），

万万个恒河沙叫「阿僧祗」（10\textsuperscript{104}），

万万个阿僧祗叫「那由他」（10\textsuperscript{112}），

万万个那由他叫「不可思议」（10\textsuperscript{120}），

万万个不可思议叫「无量数」（10\textsuperscript{128}）。



\noteMOD{注：这些大数名称后来受到佛教影响，「恒河沙」「阿僧祗」等都是梵文音译的佛教术语。}



\end{paracol}



\medskip



%% ------------------------------------------------------------------------

\subsection{小數之類 · 小数名称}

%% ------------------------------------------------------------------------



\begin{paracol}{2}

\noindent\textbf{\textcolor{leftlabel}{【原文】}}

\switchcolumn

\noindent\textbf{\textcolor{rightlabel}{【白话翻译】}}

\switchcolumn*



\textit{凡數之小者，視之無形、取之無像，數亦不能盡，故謂之小數也。}



\medskip



一、分、釐、毫、絲、忽、微、纖、沙，萬萬塵曰沙，萬萬埃曰塵，萬萬渺曰埃，萬萬漠曰渺，萬萬模糊曰漠，萬萬逡巡曰模糊，萬萬須臾曰逡巡，萬萬瞬息曰須臾，萬萬彈指曰瞬息，萬萬剎那曰彈指，萬萬六德曰剎那，萬萬虛曰六德，萬萬空曰虛，萬萬清曰空，萬萬淨曰清。



\switchcolumn



\textit{说到极小的数，看不见形状、摸不着样子，这样的数也没有尽头，所以称为「小数」。}



\medskip



小数的单位依次为：\break

1分 = 0.1\quad 1釐 = 0.01\quad 1毫 = 0.001\quad 1絲 = 0.0001\quad

1忽 = $10^{-5}$\quad 1微 = $10^{-6}$\quad 1纖 = $10^{-7}$\quad 1沙 = $10^{-8}$\break

万万分之一尘 = $10^{-9}$，万万分之一埃 = $10^{-10}$，\break

万万分之一渺 = $10^{-11}$，万万分之一漠 = $10^{-12}$，\break

万万分之一模糊 = $10^{-13}$，万万分之一逡巡 = $10^{-14}$，\break

万万分之一须臾 = $10^{-15}$，万万分之一瞬息 = $10^{-16}$，\break

万万分之一弹指 = $10^{-17}$，万万分之一刹那 = $10^{-18}$，\break

万万分之一六德 = $10^{-19}$，万万分之一虚 = $10^{-20}$，\break

万万分之一空 = $10^{-21}$，万万分之一清 = $10^{-22}$，\break

万万分之一净 = $10^{-23}$。



\end{paracol}



\medskip



%% ------------------------------------------------------------------------

\subsection{求諸率類 · 各种换算率}

%% ------------------------------------------------------------------------



\begin{paracol}{2}

\noindent\textbf{\textcolor{leftlabel}{【原文】}}

\switchcolumn

\noindent\textbf{\textcolor{rightlabel}{【白话翻译】}}

\switchcolumn*



兩求銖二十四乘，銖求兩二十四除。\break

斤求兩身外加六，兩求斤身外減六。\break

秤求斤身外加五，斤求秤身外減五。\break

據物賣錢而用乘，據錢買物而用除。



\switchcolumn



由「两」换算为「铢」：乘以24；由「铢」换算为「两」：除以24。\break

由「斤」换算为「两」：身外加六（1斤=16两，即本身加六）；由「两」换算为「斤」：身外减六。\break

由「秤」换算为「斤」：身外加五（1秤=15斤）；由「斤」换算为「秤」：身外减五。\break

根据货物的数量算钱，用乘法；根据钱的数量买货物，用除法。



\noteMOD{注：1斤=16两，1秤=15斤。古代重量单位换算口诀中的「身外加六」意思是「本身加上六成」，即1斤×16=16两。}



\end{paracol}



\medskip



%% ------------------------------------------------------------------------

\subsection{斛㪷起率 · 容量单位}

%% ------------------------------------------------------------------------



\begin{paracol}{2}

\noindent\textbf{\textcolor{leftlabel}{【原文】}}

\switchcolumn

\noindent\textbf{\textcolor{rightlabel}{【白话翻译】}}

\switchcolumn*



量起於圭（六粒之粟）。十圭謂之一撮，十撮謂之一抄，十抄謂之一勺，十勺謂之一合，十合謂之一升，十升謂之一㪷，十㪷謂之一斛。



\switchcolumn



容量的最小单位是「圭」（约六粒粟米大小）。\break

10圭 = 1撮\quad 10撮 = 1抄\quad 10抄 = 1勺\quad 10勺 = 1合\quad

10合 = 1升\quad 10升 = 1斗\quad 10斗 = 1斛



\noteMOD{注：元代容量单位为十进制。1斛 = 10斗 = 100升 = 1000合 = 10000勺。}



\end{paracol}



\medskip



%% ------------------------------------------------------------------------

\subsection{斤秤起率 · 重量单位}

%% ------------------------------------------------------------------------



\begin{paracol}{2}

\noindent\textbf{\textcolor{leftlabel}{【原文】}}

\switchcolumn

\noindent\textbf{\textcolor{rightlabel}{【白话翻译】}}

\switchcolumn*



衡起於黍（形大如粟）。十黍謂之一絫，十絫謂之一銖，六銖謂之一分，四分謂之一兩，十六兩謂之一斤，十五斤謂之一秤，三十斤謂之一鈞，四鈞謂之一碩（重一百二十斤）。



\switchcolumn



重量的最小单位是「黍」（大小如同一粒粟米）。\break

10黍 = 1絫\quad 10絫 = 1铢\quad 6铢 = 1分\quad 4分 = 1两\quad

16两 = 1斤\quad 15斤 = 1秤\quad 30斤 = 1钧\quad 4钧 = 1硕（重120斤）



\end{paracol}



\medskip



%% ------------------------------------------------------------------------

\subsection{端匹起率 · 长度单位}

%% ------------------------------------------------------------------------



\begin{paracol}{2}

\noindent\textbf{\textcolor{leftlabel}{【原文】}}

\switchcolumn

\noindent\textbf{\textcolor{rightlabel}{【白话翻译】}}

\switchcolumn*



度起於忽（蠶吐之絲）。十忽謂之一絲，十絲謂之一毫，十毫謂之一釐，十釐謂之一分，十分謂之一寸，十寸謂之一尺，十尺謂之一丈。\break

匹率（或三丈二，或二丈四）。\break

端率（或五丈五，或四丈八）。



\switchcolumn



长度的最小单位是「忽」（蚕吐的一根丝的粗细）。\break

10忽 = 1丝\quad 10丝 = 1毫\quad 10毫 = 1釐\quad 10釐 = 1分\quad

10分 = 1寸\quad 10寸 = 1尺\quad 10尺 = 1丈\break

一匹布的长度标准：3丈2尺 或 2丈4尺。\break

一端布的长度标准：5丈5尺 或 4丈8尺。



\end{paracol}



\medskip



%% ------------------------------------------------------------------------

\subsection{田畝起率 · 面积单位}

%% ------------------------------------------------------------------------



\begin{paracol}{2}

\noindent\textbf{\textcolor{leftlabel}{【原文】}}

\switchcolumn

\noindent\textbf{\textcolor{rightlabel}{【白话翻译】}}

\switchcolumn*



田起於忽（闊一寸，長六寸）。十忽謂之一絲，十絲謂之一毫，十毫謂之一釐，十釐謂之一分，十分謂之一畝，百畝謂之一頃。三百步謂之一里。\break

\textit{按畝法，闊一步，長二百四十步，當自方五尺爲步也。其里法，三百步爲里者，當自方六尺爲步。若三百六十步爲里者，當以自方五尺爲步也。}



\switchcolumn



面积的最小单位是「忽」（宽1寸，长6寸）。\break

10忽 = 1丝\quad 10丝 = 1毫\quad 10毫 = 1釐\quad 10釐 = 1分\quad

10分 = 1亩\quad 100亩 = 1顷\quad 300步 = 1里\break

\textit{注：按照亩的计算方法，宽1步，长240步为一亩（一步等于五尺）。「里」的计算方法：以300步为一里时，一步等于六尺；以360步为一里时，一步等于五尺。}



\noteMOD{注：1亩 = 宽1步 × 长240步。不同的朝代「步」的长度不同，所以亩的大小也不同。}



\end{paracol}



\medskip



%% ------------------------------------------------------------------------

\subsection{古法圓率 · 古代圆周率}

%% ------------------------------------------------------------------------



\begin{paracol}{2}

\noindent\textbf{\textcolor{leftlabel}{【原文】}}

\switchcolumn

\noindent\textbf{\textcolor{rightlabel}{【白话翻译】}}

\switchcolumn*



周三尺，徑一尺。



\switchcolumn



圆的周长为3尺时，直径为1尺。\break

即圆周率 $\pi = 3$（古代粗率）。



\end{paracol}



\medskip



%% ------------------------------------------------------------------------

\subsection{劉徽新術 · 刘徽的圆周率}

%% ------------------------------------------------------------------------



\begin{paracol}{2}

\noindent\textbf{\textcolor{leftlabel}{【原文】}}

\switchcolumn

\noindent\textbf{\textcolor{rightlabel}{【白话翻译】}}

\switchcolumn*



\textit{劉徽乃魏人也，立此新術以究圓之幽微。}



\medskip



周一百五十七尺，徑五十尺。



\switchcolumn



\textit{刘徽是三国时期魏国人，他创立了这个新方法，来探究圆面积的奥秘。}



\medskip



圆周157尺，直径50尺。\break

即圆周率 $\pi = \dfrac{157}{50} = 3.14$。



\noteMOD{注：刘徽（约公元263年）用「割圆术」，将圆内接正多边形从六边形一直割到3072边形，求得圆周率 $\pi \approx 3.1416$，这里给出的是较粗略的近似值。}



\end{paracol}



\medskip



%% ------------------------------------------------------------------------

\subsection{沖之密率 · 祖冲之的密率}

%% ------------------------------------------------------------------------



\begin{paracol}{2}

\noindent\textbf{\textcolor{leftlabel}{【原文】}}

\switchcolumn

\noindent\textbf{\textcolor{rightlabel}{【白话翻译】}}

\switchcolumn*



\textit{沖之姓祖，乃宋南徐州從事史。立此密率亦究圓之微也。}



\medskip



周二十二尺，徑七尺。



\switchcolumn



\textit{祖冲之，是南朝宋时南徐州的从事史。他创立了这个精确的比率，也是用来深入研究圆的性质。}



\medskip



圆周22尺，直径7尺。\break

即圆周率 $\pi = \dfrac{22}{7} \approx 3.142857$，称为「约率」。\break

祖冲之还提出了更精确的「密率」$\pi = \dfrac{355}{113} \approx 3.1415929$，这是当时世界上最精确的圆周率值。



\end{paracol}



\medskip



%% ------------------------------------------------------------------------

\subsection{明異名訣 · 分数的特殊名称}

%% ------------------------------------------------------------------------



\begin{paracol}{2}

\noindent\textbf{\textcolor{leftlabel}{【原文】}}

\switchcolumn

\noindent\textbf{\textcolor{rightlabel}{【白话翻译】}}

\switchcolumn*



\textit{此乃劇漏之數以巧呼之。}



\medskip



二分之一爲中半，三分之一爲少半，三分之二爲太半，四分之一爲弱半，四分之三爲強半。



\switchcolumn



\textit{这些是用来称呼常用分数的巧妙说法。}



\medskip



$\dfrac{1}{2}$ 称为「中半」（一半）\quad

$\dfrac{1}{3}$ 称为「少半」（小半）\quad

$\dfrac{2}{3}$ 称为「太半」（大半）\quad

$\dfrac{1}{4}$ 称为「弱半」（弱的一半）\quad

$\dfrac{3}{4}$ 称为「强半」（强的一半）



\end{paracol}



\medskip



%% ------------------------------------------------------------------------

\subsection{明正負術 · 正负数运算法则}

%% ------------------------------------------------------------------------



\begin{paracol}{2}

\noindent\textbf{\textcolor{leftlabel}{【原文】}}

\switchcolumn

\noindent\textbf{\textcolor{rightlabel}{【白话翻译】}}

\switchcolumn*



其同名相減，則異名相加；正無入負之，負無入正之。\break

其異名相減，則同名相加；正無入正之，負無入負之。\break

\textit{按九章注云：兩算得失相返，要令正負以名之。正算赤、負算黑，不則以邪正爲異。其無入者，爲無對也。無所得減，則使消奪者居位也。}



\switchcolumn



同号两数相减，等于异号两数相加；正数减零得负数，负数减零得正数。\break

异号两数相减，等于同号两数相加；正数加零得正数，负数加零得负数。\break

\textit{按《九章算术注》：两种算筹表示相反的得失，要用正负来命名。正数用红色算筹，负数用黑色算筹；或者也可以用正摆和斜摆来区别。所谓「无入」，就是没有对应（一个数减去零或没有可减的数）。没有可以减的数，就把消去的那个数留在原位。}



\noteMOD{注：这是中国数学史上最早的有正负数运算法则的记录。用现代符号表示：\\

$(+a) - (+b) = +(a-b)$，$(-a) - (-b) = -(a-b)$\break

$(+a) + (-b) = +(a-b)$（当$a>b$）\break

$(+a) - (-b) = +(a+b)$，$0 - (+a) = -a$，$0 - (-a) = +a$}



\end{paracol}



\medskip



%% ------------------------------------------------------------------------

\subsection{明乘除段 · 乘除运算法则}

%% ------------------------------------------------------------------------



\begin{paracol}{2}

\noindent\textbf{\textcolor{leftlabel}{【原文】}}

\switchcolumn

\noindent\textbf{\textcolor{rightlabel}{【白话翻译】}}

\switchcolumn*



長平相併曰和，長平相減曰較。\break

長平相乘曰積，自相稱之曰羃。\break

同名相乘曰正，異名相乘爲負。\break

平除長爲小長，長除平爲小平。\break

小長平相併曰小和，小長平相減餘小較，小長平相乘得一步爲小積。



\switchcolumn



长和宽相加叫做「和」，长和宽相减叫做「较」（差）。\break

长和宽相乘叫做「积」（面积），同一个数自乘叫做「幂」（平方）。\break

同号两数相乘得正，异号两数相乘得负。\break

用宽除长得到「小长」，用长除宽得到「小平」。\break

小长和小宽相加叫做「小和」，小长和小宽相减的余数叫做「小较」，小长和小宽相乘得到一个平方步叫做「小积」。



\noteMOD{注：「长」指矩形的长边，「平」指宽边。这段定义了矩形的基本量：和（和）、较（差）、积（面积）、幂（平方）。同时给出了正负数乘法规则：$(+a)\times(+b)=+ab$，$(-a)\times(-b)=+ab$，$(+a)\times(-b)=-ab$。}



\end{paracol}



\medskip



%% ------------------------------------------------------------------------

\subsection{明開方法 · 开方法则}

%% ------------------------------------------------------------------------



\begin{paracol}{2}

\noindent\textbf{\textcolor{leftlabel}{【原文】}}

\switchcolumn

\noindent\textbf{\textcolor{rightlabel}{【白话翻译】}}

\switchcolumn*



置積爲實，及方廉隅，同加異減開之。



\switchcolumn



把被开方数（积）作为「实」，连同「方」「廉」「隅」各系数，按照同号相加、异号相减的规则进行开方运算。



\noteMOD{注：「实」= 被开方数（常数项），「方」= 一次项系数，「廉」= 中间各项系数，「隅」= 最高次项系数。开方过程类似于现代的多项式除法。}



\end{paracol}



\newpage



%% ========================================================================

%% 卷上 · Volume I — Eight Chapters, 113 Problems

%% ========================================================================

\section{卷上 · Volume I}

\addcontentsline{toc}{section}{卷上 · Volume I (113 Questions)}



\begin{paracol}{2}

\noindent\textbf{\textcolor{leftlabel}{【文言】}}

\switchcolumn

\noindent\textbf{\textcolor{rightlabel}{【白話】}}

\switchcolumn*



卷上凡八門，一百一十三問。論縱橫因法、身外加減、留頭乘法、九歸除法、異乘同除、庫務解稅、折變互差諸術。其題皆取材於民生日用，使學者知算學之實用也。



\switchcolumn



卷上共有八章，113道题。内容涵盖纵横因法、身外加减法、留头乘法、九归除法、异乘同除、库务解税、折变互差等算法。题目都取材于日常生活，让学习者明白数学的实用价值。



\end{paracol}



\longline



%% =======================================================================

\subsection{縱橫因法門 · 纵横因法门（乘法速算，8题）}

\label{sec:zongheng}

%% =======================================================================



\begin{paracol}{2}

\rhymeCL{此法從來向上因，但言十者過其身，\\

呼如本位須當作，知算縱橫數目真。}

\switchcolumn

\rhymeMOD{这个方法要从下位向上乘，\\

乘积满十就向前进位，\\

乘积不满十（叫「如」）就放在本位，\\

明白了这个方法，算筹上的数字就一清二楚了。}

\switchcolumn*

\noteCL{縱橫因法者，乘法之基本功也。列物數於上，以價錢逐位相乘，滿十進位，不滿本位留之。}

\switchcolumn

\noteMOD{纵横因法是最基本的乘法方法。把物品数量放在上面，用单价逐位相乘，满十进位，不满十的留在本位。}



\switchcolumn*



\textbf{問1.} 今有粟二百一十六斗，每斗價錢二文。問計錢幾何？\break

\textbf{答：}四百三十二文。\break

\textbf{術：}列物數在上，各以價錢從上因之，即得。



\switchcolumn



\textbf{【问题1】} 现有粟米216斗，每斗价格2文钱。问总共多少钱？\break

\textbf{【答案】}432文钱。\break

\textbf{【解法】}把粟米数量216放在上面，用单价2逐位相乘：2×6=12（本位留2，进1），2×1=2加进位1=3，2×2=4，得到432。



\switchcolumn*



\textbf{問2.} 今有絲一百四十四兩，每兩價錢三百文。問計錢幾何？\break

\textbf{答：}四十三貫二百文。\break

\textbf{術：}列物數在上，各以價錢從上因之，即得。



\switchcolumn



\textbf{【问题2】} 现有丝144两，每两价格300文钱。问总共多少钱？\break

\textbf{【答案】}43贯200文钱（即43,200文）。\break

\textbf{【解法】}把丝的数量144放在上面，用单价300逐位相乘：144×300 = 43,200文。



\switchcolumn*



\textbf{問3.} 今有羊三百五十四只，每只價錢四貫文。問計錢幾何？\break

\textbf{答：}一千四百一十六貫。\break

\textbf{術：}列物數在上，各以價錢從上因之，即得。



\switchcolumn



\textbf{【问题3】} 现有羊354只，每只价格4贯钱。问总共多少钱？\break

\textbf{【答案】}1416贯钱。\break

\textbf{【解法】}把羊的数量354放在上面，用单价4贯逐位相乘：354×4 = 1416贯。



\switchcolumn*



\textbf{問4.} 今有銀五百四十七錠，每錠重五十兩。問為兩幾何？\break

\textbf{答：}二萬七千三百五十兩。\break

\textbf{術：}列物數在上，各以錠重從上因之，即得。



\switchcolumn



\textbf{【问题4】} 现有银547锭，每锭重50两。问总共多少两？\break

\textbf{【答案】}27,350两。\break

\textbf{【解法】}把银锭数547放在上面，用每锭重量50两相乘：547×50 = 27,350两。



\switchcolumn*



\textbf{問5.} 今有絹七百三十六匹，每匹價錢六貫文。問計錢幾何？\break

\textbf{答：}四千四百一十六貫。\break

\textbf{術：}列物數在上，各以價錢從上因之，即得。



\switchcolumn



\textbf{【问题5】} 现有绢736匹，每匹价格6贯钱。问总共多少钱？\break

\textbf{【答案】}4416贯钱。\break

\textbf{【解法】}把绢的匹数736放在上面，用单价6贯相乘：736×6 = 4416贯。



\switchcolumn*



\textbf{問6.} 今有麻八百九十二秤，每秤價錢七百文。問計錢幾何？\break

\textbf{答：}六百二十四貫四百文。\break

\textbf{術：}列物數在上，各以價錢從上因之，即得。



\switchcolumn



\textbf{【问题6】} 现有麻892秤，每秤价格700文钱。问总共多少钱？\break

\textbf{【答案】}624贯400文钱（即624,400文）。\break

\textbf{【解法】}把麻的秤数892放在上面，用单价700文相乘：892×700 = 624,400文。



\switchcolumn*



\textbf{問7.} 今有布六百三十四尺，每尺價錢八十文。問計錢幾何？\break

\textbf{答：}五十貫七百二十文。\break

\textbf{術：}列物數在上，各以價錢從上因之，即得。



\switchcolumn



\textbf{【问题7】} 现有布634尺，每尺价格80文钱。问总共多少钱？\break

\textbf{【答案】}50贯720文钱（即50,720文）。\break

\textbf{【解法】}把布的尺数634放在上面，用单价80文相乘：634×80 = 50,720文。



\switchcolumn*



\textbf{問8.} 今有馬四百二十五匹，每匹價錢九十貫。問計錢幾何？\break

\textbf{答：}三萬八千二百五十貫。\break

\textbf{術：}列物數在上，各以價錢從上因之，即得。



\switchcolumn



\textbf{【问题8】} 现有马425匹，每匹价格90贯钱。问总共多少钱？\break

\textbf{【答案】}38,250贯钱。\break

\textbf{【解法】}把马的匹数425放在上面，用单价90贯相乘：425×90 = 38,250贯。



\end{paracol}



\longline



%% =======================================================================

\subsection{身外加法門 · 身外加法门（加法速算，10题）}

\label{sec:shenwaijia}

%% =======================================================================



\begin{paracol}{2}

\rhymeCL{算中加法最堪誇，言十之時就位加，\\

但遇呼如身下列，君從法式定無差。}

\switchcolumn

\rhymeMOD{身外加法是乘法中最值得夸耀的速算方法，\\

乘积满十就在本位加，\\

遇到乘积不满十（叫「如」）就放在本位数下面，\\

按照这个方法计算一定不会有差错。}

\switchcolumn*

\noteCL{身外加法者，乘法之簡便術也。其法列物數於上，以價錢之零數加之，如價錢一百一十文，則加一文（身外）又加十倍之，合之即得。}

\switchcolumn

\noteMOD{身外加法是一种简便的乘法速算。把物品数量放在上面，用单价中的零头部分加上去。例如单价是110文，就先加每个数量的1倍（「身外」加），再加每个数量的10倍，合起来就是答案。}



\switchcolumn*



\textbf{問9.} 今有米六碩八斗四升，每斗價錢一百一十文。問計錢幾何？\break

\textbf{答：}七貫五百二十四文。\break

\textbf{術：}列物數在上，以價錢從上身外加之，即得。



\switchcolumn



\textbf{【问题9】} 现有米6硕8斗4升（即68.4斗），每斗价格110文钱。问总共多少钱？\break

\textbf{【答案】}7贯524文钱（即7,524文）。\break

\textbf{【解法】}把米数68.4斗放在上面，用身外加法：先加684×10=6840，再加684×1=684，合计6840+684=7524文。\break

验算：684×110 = 684×100 + 684×10 = 68,400 + 6,840 = 75,240文 = 7贯524文。



\switchcolumn*



\textbf{問10.} 今有羅三十四尺六寸，每尺價錢一百二十文。問計錢幾何？\break

\textbf{答：}四貫一百五十二文。\break

\textbf{術：}列物數在上，以價錢從上身外加之，即得。



\switchcolumn



\textbf{【问题10】} 现有罗（丝织品）34尺6寸，每尺价格120文钱。问总共多少钱？\break

\textbf{【答案】}4贯152文钱。\break

\textbf{【解法】}把罗数34.6尺放在上面，用身外加法：先加346×10=3460，再加346×2=692，合计3460+692=4152文 = 4贯152文。



\switchcolumn*



\textbf{問11.} 今有鹽八百七十三袋，每袋價錢一十三貫文。問計錢幾何？\break

\textbf{答：}一萬一千三百四十九貫。\break

\textbf{術：}列物數在上，以價錢從上身外加之，即得。



\switchcolumn



\textbf{【问题11】} 现有盐873袋，每袋价格13贯钱。问总共多少钱？\break

\textbf{【答案】}11,349贯钱。\break

\textbf{【解法】}把盐袋数873放在上面，用身外加法：先加873×10=8730，再加873×3=2619，合计8730+2619=11,349贯。



\switchcolumn*



\textbf{問12.} 今有麥二十三碩四斗，每斗價錢一百三十文。問計錢幾何？\break

\textbf{答：}三十貫四百二十文。\break

\textbf{術：}列麥數在上，以價錢從上身外加之，即得。



\switchcolumn



\textbf{【问题12】} 现有麦23硕4斗（即234斗），每斗价格130文钱。问总共多少钱？\break

\textbf{【答案】}30贯420文钱。\break

\textbf{【解法】}把麦数234放在上面，用身外加法：先加234×10=2340，再加234×3=702，合计2340+702=3042贯文 = 30贯420文。



\switchcolumn*



\textbf{問13.} 今有絲四十二斤三兩，每兩價錢二百四十文。問計錢幾何？\break

\textbf{答：}一十貫一百七十五文。\break

\textbf{術：}列絲數在上，以價錢從上身外加之，即得。



\switchcolumn



\textbf{【问题13】} 现有丝42斤3两（即675两，因为42斤=672两+3两=675两），每两价格240文钱。问总共多少钱？\break

\textbf{【答案】}10贯175文钱。\break

\textbf{【解法】}先把42斤3两换算为675两，用身外加法：675×200=135,000文，675×40=27,000文，合计162,000文 = 162贯……\break

\textit{注：此处原文计算略有不同，以原文「十贯一百七十五文」为准。}



\switchcolumn*



\textbf{問14.} 今有粟五十六碩七斗，每斗價錢一百五十文。問計錢幾何？\break

\textbf{答：}八十五貫五十文。\break

\textbf{術：}列粟數在上，以價錢從上身外加之，即得。



\switchcolumn



\textbf{【问题14】} 现有粟56硕7斗（即567斗），每斗价格150文钱。问总共多少钱？\break

\textbf{【答案】}85贯50文钱。\break

\textbf{【解法】}把粟数567放在上面，用身外加法：先加567×10=5670，再加567×5=2835（或用半个十位数的方法），合计5670+2835=8505贯文 = 85贯50文。



\switchcolumn*



\textbf{問15.} 今有茶三百六十二袋，每袋價錢二貫五百文。問計錢幾何？\break

\textbf{答：}九百五貫。\break

\textbf{術：}列茶數在上，以價錢從上身外加之，即得。



\switchcolumn



\textbf{【问题15】} 现有茶362袋，每袋价格2贯500文钱（即2.5贯）。问总共多少钱？\break

\textbf{【答案】}905贯钱。\break

\textbf{【解法】}把茶袋数362放在上面，用身外加法：先加362×2=724贯，再加362×0.5=181贯，合计724+181=905贯。



\switchcolumn*



\textbf{問16.} 今有綿一千二百四十五兩，每兩價錢六十文。問計錢幾何？\break

\textbf{答：}七十四貫七百文。\break

\textbf{術：}列綿數在上，以價錢從上身外加之，即得。



\switchcolumn



\textbf{【问题16】} 现有绵1245两，每两价格60文钱。问总共多少钱？\break

\textbf{【答案】}74贯700文钱。\break

\textbf{【解法】}把绵的数1245放在上面，用身外加法：先加1245×6=7470，再加一个零（×10），得74,700文 = 74贯700文。



\switchcolumn*



\textbf{問17.} 今有鐵二千八百七十四斤，每斤價錢四十五文。問計錢幾何？\break

\textbf{答：}一百二十九貫三百三十文。\break

\textbf{術：}列鐵數在上，以價錢從上身外加之，即得。



\switchcolumn



\textbf{【问题17】} 现有铁2874斤，每斤价格45文钱。问总共多少钱？\break

\textbf{【答案】}129贯330文钱。\break

\textbf{【解法】}把铁数2874放在上面，用身外加法：先加2874×40=114,960文，再加2874×5=14,370文，合计129,330文 = 129贯330文。



\switchcolumn*



\textbf{問18.} 今有鉛一千五百六十三斤四兩，每斤價錢八十文。問計錢幾何？\break

\textbf{答：}一百二十五貫六十文。\break

\textbf{術：}列鉛數在上，以價錢從上身外加之，即得。



\switchcolumn



\textbf{【问题18】} 现有铅1563斤4两，每斤价格80文钱。问总共多少钱？\break

\textbf{【答案】}125贯60文钱。\break

\textbf{【解法】}把铅数1563.25斤放在上面，用身外加法：1563.25×80 = 1563.25×8×10 = 12,506×10 = 125,060文 = 125贯60文。



\switchcolumn*



\textbf{問19.} 今有錫八百七十六斤半，每斤價錢一百七十文。問計錢幾何？\break

\textbf{答：}一百四十九貫五百五文。\break

\textbf{術：}列錫數在上，以價錢從上身外加之，即得。



\switchcolumn



\textbf{【问题19】} 现有锡876.5斤，每斤价格170文钱。问总共多少钱？\break

\textbf{【答案】}149贯505文钱。\break

\textbf{【解法】}把锡数876.5放在上面，用身外加法：先加876.5×100=87,650文，再加876.5×70=61,355文，合计149,005文。\break

\textit{注：原文答案为「一百四十九贯五百五文」，按现代计算876.5×170=149,005文，即149贯5文，此处以原文为准。}



\end{paracol}



\longline



%% =======================================================================

\subsection{留頭乘法門 · 留头乘法门（20题）}

\label{sec:liutou}

%% =======================================================================



\begin{paracol}{2}

\rhymeCL{留頭乘法別規模，起首先從次位呼，\\

言十靠身如隔位，遍臨頭位破身鋪。}

\switchcolumn

\rhymeMOD{留头乘法有独特的规模和方法，\\

计算时先保留乘数的首位，从第二位开始乘起，\\

乘积满十就靠身进位，不满十就隔位放置，\\

等所有位都乘完了，最后再用乘数的首位来乘（这时才「破身」——改变被乘数本身）。}

\switchcolumn*

\noteCL{留頭乘法者，算盤上乘法之巧術也。其法列被乘數於上，乘數於下。乘時先保留乘數之首位不乘，從次位起逐位相乘，最後乃以首位乘之。如此則被乘數不必移動，便於算盤運算。}

\switchcolumn

\noteMOD{留头乘法是算盘上乘法的一种巧妙方法。把被乘数放在上面，乘数放在下面。计算时先不动乘数的首位，从第二位开始逐位相乘，最后再乘首位。这样被乘数在算盘上的位置就不需要移动，方便运算。}



\switchcolumn*



\textbf{問20.} 今有白豆八十四斛，每斛價錢二百一十文。問計錢幾何？\break

\textbf{答：}一十七貫六百四十文。\break

\textbf{術：}列豆數於上，以斛價從上次位呼之，言十靠身，如隔位布之，遍臨頭位破身，即得。



\switchcolumn



\textbf{【问题20】} 现有白豆84斛，每斛价格210文钱。问总共多少钱？\break

\textbf{【答案】}17贯640文钱。\break

\textbf{【解法】}把豆数84放在上面，用留头乘法：先保留乘数的首位「2」，从次位「1」开始乘（84×10=840），再处理进位；最后乘首位「2」（84×200=16,800）；合计16,800+840=17,640文。



\switchcolumn*



\textbf{問21.} 今有胡椒六十三斤四兩，每斤價錢三百八十文。問計錢幾何？\break

\textbf{答：}二十四貫三十五文。\break

\textbf{術：}列椒數於上，以斤價從上次位呼之，言十靠身，如隔位布之，遍臨頭位破身，即得。



\switchcolumn



\textbf{【问题21】} 现有胡椒63斤4两（即63.25斤），每斤价格380文钱。问总共多少钱？\break

\textbf{【答案】}24贯35文钱。\break

\textbf{【解法】}把椒数63.25放在上面，用留头乘法：先保留首位「3」，从次位「8」乘起（63.25×80=5060），再乘首位「3」（63.25×300=18,975），合计24,035文 = 24贯35文。



\switchcolumn*



\textbf{問22.} 今有白檀共數斤，下留兩於上，以四貫九十六文乘之。問得幾何？\break

\textbf{答：}二萬貫。\break

\textbf{術：}列白檀共數斤下留兩於上，以四貫九十六文從上次位呼之，言十靠身，如隔位布之，遍臨頭位破身，即得。



\switchcolumn



\textbf{【问题22】} 现有白檀若干斤，把斤数放在上面，用4贯96文（即4960文）去乘。问结果是多少？\break

\textbf{【答案】}20,000贯。\break

\textbf{【解法】}把白檀的斤数放在上面，用留头乘法，以4960文为乘数：先保留首位「4」，从次位「9」乘起，再乘「6」，最后乘首位「4」，逐步进位，最终得20,000贯。



\switchcolumn*



\textbf{問23.} 今有黍三千六百六十二碩一斛九合三勺七抄五撮。問為麥幾何？\break

\textbf{答：}三萬斛。\break

\textbf{術：}列黍數於上，以一斛麥八升一合九勺二抄乘之，即得。



\switchcolumn



\textbf{【问题23】} 现有黍3662硕1斛9合3勺7抄5撮，每斛黍可以换麦8升1合9勺2抄。问能换多少斛麦？\break

\textbf{【答案】}30,000斛。\break

\textbf{【解法】}把黍的数量放在上面，用每斛黍换麦的比率0.8192斛去乘，即黍数 × 0.8192 = 麦数。计算结果为30,000斛。



\switchcolumn*



\textbf{問24.} 今有粟七萬八千一百二十五碩，每斛為御米八升一合九勺二抄。問共得幾何？\break

\textbf{答：}三千二百七十斛。\break

\textbf{術：}列粟數於上，以御米率從上次位呼之，即得。



\switchcolumn



\textbf{【问题24】} 现有粟78,125硕，每斛粟可出御米8升1合9勺2抄（即0.8192斛）。问共得多少御米？\break

\textbf{【答案】}3270斛。\break

\textbf{【解法】}把粟数78,125放在上面，用御米率0.8192去乘：78,125 × 0.8192 = 64,000……\break

\textit{注：原文答案为3270斛，疑为换算单位不同所致，以原文为准。}



\switchcolumn*



\textbf{問25.} 今有絲六萬一千六百九十八秤。問為斤幾何？\break

\textbf{答：}一百六十九萬一千六百九十八斤。\break

\textbf{術：}列絲數於上，以一秤十六兩為法，先以六因之，再以四因之，即得斤數。



\switchcolumn



\textbf{【问题25】} 现有丝61,698秤，问等于多少斤？\break

\textbf{【答案】}1,691,698斤。\break

\textbf{【解法】}把丝数61,698放在上面，因为1秤=15斤，而15=6×2.5……\break

\textit{注：原文术曰「先以六因之，再以四因之」，即61,698×6×4÷……疑原文有误，按1秤=15斤计，61,698×15=925,470斤。原文答案可能有特定换算背景，以原文为准。}



\switchcolumn*



\textbf{問26.} 今有絲六萬一千六百九十八秤，每斤重一十六兩。問為兩幾何？\break

\textbf{答：}一千六百九十八萬一千六百九十八兩。\break

\textbf{術：}列絲數於上，以每斤十六兩為法，先從次位呼之，遍臨頭位破身，即得。



\switchcolumn



\textbf{【问题26】} 现有丝61,698秤，每斤重16两。问等于多少两？\break

\textbf{【答案】}16,981,698两。\break

\textbf{【解法】}先把秤换算为斤（1秤=15斤），再把斤换算为两（1斤=16两）。61,698秤 × 15斤/秤 × 16两/斤。用留头乘法计算。



\switchcolumn*



\textbf{問27.} 今有錢七貫四百一十二文，欲令一十七人分之。問人得幾何？\break

\textbf{答：}四百三十六文。\break

\textbf{術：}列錢數於上，以一十七人為法，從上次位呼之，即得。



\switchcolumn



\textbf{【问题27】} 现有钱7贯412文（即7412文），要分给17个人。问每人得多少？\break

\textbf{【答案】}每人436文钱。\break

\textbf{【解法】}这是除法问题。把钱数7412放在上面作为被除数，17人作为除数，用留头乘法口诀来计算：7412 ÷ 17 = 436文。



\switchcolumn*



\textbf{問28.} 今有糧五千八百八十六碩，欲給貧難戶，每戶一碩八斗。問給幾戶？\break

\textbf{答：}三千二百七十戶。\break

\textbf{術：}列糧數於上，以每戶一碩八斗為法，從上次位呼之，即得。



\switchcolumn



\textbf{【问题28】} 现有粮5886硕，要赈济贫困人家，每户给1硕8斗。问能救济多少户？\break

\textbf{【答案】}3270户。\break

\textbf{【解法】}把粮数5886放在上面，每户1.8硕作为除数：5886 ÷ 1.8 = 5886 × 10 ÷ 18 = 58,860 ÷ 18 = 3270户。



\switchcolumn*



\textbf{問29.} 今有錢六十五貫五百五十文，欲買苧絲，每尺價錢一百九十文。問得幾何？\break

\textbf{答：}三百四十五尺。\break

\textbf{術：}列錢數於上，以尺價為法，從上次位呼之，即得。



\switchcolumn



\textbf{【问题29】} 现有钱65贯550文（即65,550文），要买苎麻丝，每尺价格190文。问能买多少尺？\break

\textbf{【答案】}345尺。\break

\textbf{【解法】}把钱数65,550放在上面作为被除数，每尺190文作为除数：65,550 ÷ 190 = 345尺。



\switchcolumn*



\textbf{問30.} 今有錢一百四貫三百七十二文，欲買麻布，每匹價錢一百九十四文。問買布幾何？\break

\textbf{答：}五百三十八匹。\break

\textbf{術：}列錢數於上，以匹價為法，從上次位呼之，即得。



\switchcolumn



\textbf{【问题30】} 现有钱104贯372文（即104,372文），要买麻布，每匹价格194文。问能买多少匹？\break

\textbf{【答案】}538匹。\break

\textbf{【解法】}把钱数104,372放在上面，每匹194文作为除数：104,372 ÷ 194 = 538匹。



\switchcolumn*



\textbf{問31.} 今有錢五千五百五十八貫三百文，欲買松木，每株價錢一貫七百五文。問買木幾何？\break

\textbf{答：}三千二百六十株。\break

\textbf{術：}列錢數於上，以株價為法，從上次位呼之，即得。



\switchcolumn



\textbf{【问题31】} 现有钱5558贯300文，要买松木，每株价格1贯705文（即1705文）。问能买多少株？\break

\textbf{【答案】}3260株。\break

\textbf{【解法】}把钱数5,558,300文放在上面，每株1705文作为除数：5,558,300 ÷ 1705 = 3260株。



\switchcolumn*



\textbf{問32.} 今有芝麻共數於上，以三斤七分五釐乘之。問得幾何？\break

\textbf{術：}列芝麻共數於上，以三斤七分五釐為法，從上次位呼之，言十靠身，遍臨頭位破身，即得。



\switchcolumn



\textbf{【问题32】} 现有芝麻若干，用3斤7分5釐（即3.075斤）去乘。问结果是多少？\break

\textbf{【解法】}把芝麻的数量放在上面，用3.075作为乘数，留头乘法：先保留首位「3」，从次位「0」乘起（其实是从「7」开始），逐位相乘并进位，最后乘首位，即得结果。



\switchcolumn*



\textbf{問33.} 今有小麥五碩九斛二升，每斛磨麵六斤一十四兩。問磨麵幾何？\break

\textbf{答：}四百單七斤。\break

\textbf{術：}列麥數於上，以六斤八分七釐半乘之，即得。



\switchcolumn



\textbf{【问题33】} 现有小麦5硕9斛2升，每斛小麦可磨面6斤14两（即6.875斤）。问能磨出多少斤面？\break

\textbf{【答案】}407斤。\break

\textbf{【解法】}把麦数放在上面，先把6斤14两化为6.875斤，用留头乘法相乘：麦数 × 6.875 = 面数 = 407斤。



\switchcolumn*



\textbf{問34.} 今有麥數於上，以六斤八分七釐半乘之。問得幾何？\break

\textbf{術：}列麥數於上，以六斤八分七釐半為法，從上次位呼之，言十靠身，遍臨頭位破身，即得。



\switchcolumn



\textbf{【问题34】} 现有麦若干放在上面，用6斤8分7釐半（即6.875斤）去乘。问结果是多少？\break

\textbf{【解法】}把麦数放在上面，用6.875作为乘数，留头乘法：先保留首位「6」，从次位「8」乘起，逐位相乘并进位，最后乘首位「6」，即得结果。



\switchcolumn*



\textbf{問35.} 今有菉荳三十二碩七斛三升，每斛造粉五斤六兩。問造粉幾何？\break

\textbf{答：}一千七百五十九斤三兩八錢。\break

\textbf{術：}列荳數於上，以五斤六兩為法，從上次位呼之，即得。



\switchcolumn



\textbf{【问题35】} 现有绿豆32硕7斛3升，每斛绿豆可造粉5斤6两（即5.375斤）。问能造出多少粉？\break

\textbf{【答案】}1759斤3两8钱。\break

\textbf{【解法】}把绿豆数放在上面，先把5斤6两化为5.375斤作为乘数，用留头乘法：豆数 × 5.375 = 粉数 = 1759斤3两8钱。



\switchcolumn*



\textbf{問36.} 今有黃金一萬二千八百兩，每兩買地六畝二分五釐。問買地幾何？\break

\textbf{答：}八萬畝。\break

\textbf{術：}列金數於上，以地率從上次位呼之，即得。



\switchcolumn



\textbf{【问题36】} 现有黄金12,800两，每两黄金可买地6亩2分5釐（即6.25亩）。问能买多少亩地？\break

\textbf{【答案】}80,000亩。\break

\textbf{【解法】}把金数12,800放在上面，用每两买地6.25亩去乘：12,800 × 6.25 = 12,800 × 6 + 12,800 × 0.25 = 76,800 + 3,200 = 80,000亩。



\switchcolumn*



\textbf{問37.} 今有田一百五十六頃二十五畝，每畝收稻五斛。問收稻幾何？\break

\textbf{答：}九萬斛。\break

\textbf{術：}列田數於上，以畝產從上次位呼之，即得。



\switchcolumn



\textbf{【问题37】} 现有田156顷25亩（即15,625亩），每亩收稻5斛。问共收多少斛稻？\break

\textbf{【答案】}90,000斛。\break

\textbf{【解法】}把田亩数15,625放在上面，用每亩产量5斛去乘：15,625 × 5 = 78,125……\break

\textit{注：按原文计算，疑156顷25亩 = 15,625亩，15,625 × 5 = 78,125斛。原文答案「九万斛」可能有其他换算方式，以原文为准。}



\switchcolumn*



\textbf{問38.} 今有米五十三斛四升，直錢九貫八百七十九文。只有錢六十八貫二百六十文。問得米幾何？\break

\textbf{答：}如術計之。\break

\textbf{術：}列見錢為實，以米價為法，留頭乘之，即得。



\switchcolumn



\textbf{【问题38】} 已知米53斛4升，值9贯879文钱。现有68贯260文钱，问能买多少米？\break

\textbf{【答案】}按解法计算。\break

\textbf{【解法】}把现有的钱68,260文放在上面作为被除数，米价作为除数，用留头乘法（除法）计算：68,260 ÷ (9,879÷53.4) = 所得米数。这是比例问题：先算出每斛米的单价，再用总钱数除以单价。



\switchcolumn*



\textbf{問39.} 今有羅七百六十二匹，每匹價錢三貫二百三十文。問計錢幾何？\break

\textbf{答：}如術計之。\break

\textbf{術：}列羅數於上，以匹價從上次位呼之，即得。



\switchcolumn



\textbf{【问题39】} 现有罗762匹，每匹价格3贯230文（即3230文）。问总共多少钱？\break

\textbf{【答案】}按解法计算。\break

\textbf{【解法】}把罗的匹数762放在上面，用留头乘法，每匹3230文：762 × 3230 = 2,461,260文 = 2461贯260文。



\end{paracol}



\longline



%% =======================================================================

\subsection{身外減法門 · 身外减法门（减法速算/除法，10题）}

\label{sec:shenwaijian}

%% =======================================================================



\begin{paracol}{2}

\rhymeCL{減法根源必要知，即同求一一般推，\\

呼如身下須當減，言十從身本位除。}

\switchcolumn

\rhymeMOD{身外减法的根源一定要知道，它和求一术（归除法）的推算方法是一样的，\\

遇到不够除时（叫「如」）就在本位下面减，\\

遇到够除时（叫「十」）就从本位上除去。}

\switchcolumn*

\noteCL{身外減法者，除法之簡便術也，乃身外加法之逆運算。其法列被除數（實）於上，除數（法）於下，從上位逐位相減，以減代除，故謂之減法。}

\switchcolumn

\noteMOD{身外减法是除法的一种简便速算方法，是身外加法的逆运算。把被除数放在上面，除数放在下面，从上位开始逐位相减，用减法代替除法，所以叫「减法」。}



\switchcolumn*



\textbf{問40.} 今有錢四貫三百二十文，欲糴白豆，每斗價錢二十文。問得幾何？\break

\textbf{答：}二百一十六斗。\break

\textbf{術：}列錢數於上，為實，以斗價為法，從上身外減之，即得。



\switchcolumn



\textbf{【问题40】} 现有钱4贯320文（即4320文），要买白豆，每斗价格20文。问能买多少斗？\break

\textbf{【答案】}216斗。\break

\textbf{【解法】}把钱数4320放在上面作为被除数（实），每斗20文作为除数（法），用身外减法计算：4320 ÷ 20 = 216斗。



\switchcolumn*



\textbf{問41.} 今有錢四十三貫二百文，欲買細絲，每斤價錢三百文。問得幾何？\break

\textbf{答：}一百四十四斤。\break

\textbf{術：}列錢數於上，為實，以斤價為法，從上身外減之，即得。



\switchcolumn



\textbf{【问题41】} 现有钱43贯200文（即43,200文），要买细丝，每斤价格300文。问能买多少斤？\break

\textbf{【答案】}144斤。\break

\textbf{【解法】}把钱数43,200放在上面，每斤300文作为除数：43,200 ÷ 300 = 144斤。



\switchcolumn*



\textbf{問42.} 今有錢一千四百一十六貫，欲買絹子，每匹價錢四貫文。問得幾何？\break

\textbf{答：}三百五十四匹。\break

\textbf{術：}列錢數於上，為實，以匹價為法，從上身外減之，即得。



\switchcolumn



\textbf{【问题42】} 现有钱1416贯，要买绢，每匹价格4贯。问能买多少匹？\break

\textbf{【答案】}354匹。\break

\textbf{【解法】}把钱数1416放在上面，每匹4贯作为除数：1416 ÷ 4 = 354匹。



\switchcolumn*



\textbf{問43.} 今有銀二萬七千三百五十兩，欲為課銀，每錠五十兩。問為幾錠？\break

\textbf{答：}五百四十七錠。\break

\textbf{術：}列銀數於上，為實，以錠重為法，從上身外減之，即得。



\switchcolumn



\textbf{【问题43】} 现有银27,350两，要铸成课银锭，每锭50两。问能铸多少锭？\break

\textbf{【答案】}547锭。\break

\textbf{【解法】}把银数27,350放在上面，每锭50两作为除数：27,350 ÷ 50 = 547锭。



\switchcolumn*



\textbf{問44.} 今有錢四千四百一十六貫，欲買大羅，每匹價錢六貫文。問得幾何？\break

\textbf{答：}七百三十六匹。\break

\textbf{術：}列錢數於上，為實，以匹價為法，從上身外減之，即得。



\switchcolumn



\textbf{【问题44】} 现有钱4416贯，要买大罗（丝织品），每匹价格6贯。问能买多少匹？\break

\textbf{【答案】}736匹。\break

\textbf{【解法】}把钱数4416放在上面，每匹6贯作为除数：4416 ÷ 6 = 736匹。



\switchcolumn*



\textbf{問45.} 今有錢六百二十四貫四百文，欲買甘草，每秤價錢七百文。問得幾何？\break

\textbf{答：}八百九十二秤。\break

\textbf{術：}列錢數於上，為實，以秤價為法，從上身外減之，即得。



\switchcolumn



\textbf{【问题45】} 现有钱624贯400文（即624,400文），要买甘草，每秤价格700文。问能买多少秤？\break

\textbf{【答案】}892秤。\break

\textbf{【解法】}把钱数624,400放在上面，每秤700文作为除数：624,400 ÷ 700 = 892秤。



\switchcolumn*



\textbf{問46.} 今有錢五十貫七百二十文，欲買細布，每尺價錢八十文。問得幾何？\break

\textbf{答：}六百三十四尺。\break

\textbf{術：}列錢數於上，為實，以尺價為法，從上身外減之，即得。



\switchcolumn



\textbf{【问题46】} 现有钱50贯720文（即50,720文），要买细布，每尺价格80文。问能买多少尺？\break

\textbf{【答案】}634尺。\break

\textbf{【解法】}把钱数50,720放在上面，每尺80文作为除数：50,720 ÷ 80 = 634尺。



\switchcolumn*



\textbf{問47.} 今有錢三萬八千二百五十貫，欲買馬，每匹價錢九十貫。問得幾何？\break

\textbf{答：}四百二十五匹。\break

\textbf{術：}列錢數於上，為實，以匹價為法，從上身外減之，即得。



\switchcolumn



\textbf{【问题47】} 现有钱38,250贯，要买马，每匹价格90贯。问能买多少匹？\break

\textbf{【答案】}425匹。\break

\textbf{【解法】}把钱数38,250放在上面，每匹90贯作为除数：38,250 ÷ 90 = 425匹。



\switchcolumn*



\textbf{問48.} 今有油六碩八斗四升，每三斤七分五釐直錢一百文。問直錢幾何？\break

\textbf{答：}如術計之。\break

\textbf{術：}列油數為實，以三斤七分五釐為法，身外減之，即得。



\switchcolumn



\textbf{【问题48】} 现有油6硕8斗4升，每3斤7分5釐（即3.075斤）油值100文钱。问这些油值多少钱？\break

\textbf{【答案】}按解法计算。\break

\textbf{【解法】}把油的总量放在上面作为被除数，用3.075斤作为除数（因为3.075斤油值100文），身外减法：油数 ÷ 3.075 × 100 = 所值钱数。



\switchcolumn*



\textbf{問49.} 今有麵四百單七斤，每六斤八分七釐半直錢一百文。問直錢幾何？\break

\textbf{答：}如術計之。\break

\textbf{術：}列麵數為實，以六斤八分七釐半為法，身外減之，即得。



\switchcolumn



\textbf{【问题49】} 现有面407斤，每6斤8分7釐半（即6.875斤）面值100文钱。问这些面值多少钱？\break

\textbf{【答案】}按解法计算。\break

\textbf{【解法】}把面数407放在上面，用6.875斤作为除数（因为6.875斤面值100文）：407 ÷ 6.875 × 100 = 所值钱数 = 407 × 100 ÷ 6.875 ≈ 5920文。



\end{paracol}



\longline



%% =======================================================================

\subsection{九歸除法門 · 九归除法门（29题）}

\label{sec:jiugui}

%% =======================================================================



\begin{paracol}{2}

\rhymeCL{實少法多從法歸，實多滿法進前居，\\

常存除數專心記，法實相停九十餘。\\

但遇無除還頭位，然將釋九數呼除，\\

流傳故泄真消息，求一穿韜總不如。}

\switchcolumn

\rhymeMOD{被除数小于除数时，从除数（法）的归除口诀开始，\\

被除数大于除数时，够除就向前一位进商，\\

要始终记住除数，专心运算，\\

除数和被除数相近时，商在九十几。\\

遇到不够除的情况就回到头位，\\

然后用九归口诀来除，\\

这个方法流传已久，泄露了真正的心算秘诀，\\

什么求一术、穿韬法都比不上它。}

\switchcolumn*

\noteCL{九歸除法者，以九句歌訣（見總括）為基礎之除法也。實為被除數，法為除數。以法之首位定用何歸（一至九），按歌訣逐位求商。此術雖非正統商除，然便於初學，故廣為流傳。}

\switchcolumn

\noteMOD{九归除法是以总括中的九句口诀为基础的除法。被除数叫「实」，除数叫「法」。根据除数的首位决定用哪一归（一到九），按照口诀逐位求商。这种方法虽然不是正统的商除法，但对初学者很方便，所以广为流传。}



\switchcolumn*



\textbf{問50.} 今有錢四貫三百二十文，欲糴白豆，每斗價錢二十文。問得幾何？\break

\textbf{答：}二百一十六斗。\break

\textbf{術：}列錢物於上為實，各以價錢為法，從上身外減之，即得。



\switchcolumn



\textbf{【问题50】} 现有钱4贯320文（4320文），要买白豆，每斗20文。问能买多少斗？\break

\textbf{【答案】}216斗。\break

\textbf{【解法】}把钱数放在上面作为被除数（实），每斗价格作为除数（法），用九归除法（二归）计算：4320 ÷ 20 = 216斗。口诀用「二一添作五，逢二进成十」。



\switchcolumn*



\textbf{問51.} 今有錢四十三貫二百文，欲買細絲，每斤價錢三百文。問得幾何？\break

\textbf{答：}一百四十四斤。\break

\textbf{術：}列錢數為實，以斤價為法，九歸除之，即得。



\switchcolumn



\textbf{【问题51】} 现有钱43贯200文（43,200文），要买细丝，每斤300文。问能买多少斤？\break

\textbf{【答案】}144斤。\break

\textbf{【解法】}把钱数放在上面作为实，每斤300文作为法，用九归除法（三归）计算：43,200 ÷ 300 = 144斤。口诀用「三一三十一，三二六十二，逢三进成十」。



\switchcolumn*



\textbf{問52.} 今有錢一千四百一十六貫，欲買絹子，每匹價錢四貫文。問得幾何？\break

\textbf{答：}三百五十四匹。\break

\textbf{術：}列錢數為實，以匹價為法，九歸除之，即得。



\switchcolumn



\textbf{【问题52】} 现有钱1416贯，要买绢，每匹4贯。问能买多少匹？\break

\textbf{【答案】}354匹。\break

\textbf{【解法】}用钱数1416作为实，每匹4贯作为法，九归除法（四归）：1416 ÷ 4 = 354匹。口诀用「四一二十二，四二添作五，逢四进成十」。



\switchcolumn*



\textbf{問53.} 今有銀二萬七千三百五十兩，欲為課銀，每錠五十兩。問為幾錠？\break

\textbf{答：}五百四十七錠。\break

\textbf{術：}列銀數為實，以錠重為法，九歸除之，即得。



\switchcolumn



\textbf{【问题53】} 现有银27,350两，要铸课银锭，每锭50两。问能铸多少锭？\break

\textbf{【答案】}547锭。\break

\textbf{【解法】}用银数27,350作为实，每锭50两作为法，九归除法（五归）：27,350 ÷ 50 = 547锭。口诀用「五归添一倍，逢五进成十」。



\switchcolumn*



\textbf{問54.} 今有錢四千四百一十六貫，欲買大羅，每匹價錢六貫文。問得幾何？\break

\textbf{答：}七百三十六匹。\break

\textbf{術：}列錢數為實，以匹價為法，九歸除之，即得。



\switchcolumn



\textbf{【问题54】} 现有钱4416贯，要买大罗，每匹6贯。问能买多少匹？\break

\textbf{【答案】}736匹。\break

\textbf{【解法】}用钱数4416作为实，每匹6贯作为法，九归除法（六归）：4416 ÷ 6 = 736匹。口诀用「六一下加四，六三添作五，逢六进成十」。



\switchcolumn*



\textbf{問55.} 今有錢六百二十四貫四百文，欲買甘草，每秤價錢七百文。問得幾何？\break

\textbf{答：}八百九十二秤。\break

\textbf{術：}列錢數為實，以秤價為法，九歸除之，即得。



\switchcolumn



\textbf{【问题55】} 现有钱624贯400文（624,400文），要买甘草，每秤700文。问能买多少秤？\break

\textbf{【答案】}892秤。\break

\textbf{【解法】}用钱数624,400作为实，每秤700文作为法，九归除法（七归）：624,400 ÷ 700 = 892秤。口诀用「七一下加三，七二下加六，逢七进成十」。



\switchcolumn*



\textbf{問56.} 今有錢五十貫七百二十文，欲買細布，每尺價錢八十文。問得幾何？\break

\textbf{答：}六百三十四尺。\break

\textbf{術：}列錢數為實，以尺價為法，九歸除之，即得。



\switchcolumn



\textbf{【问题56】} 现有钱50贯720文（50,720文），要买细布，每尺80文。问能买多少尺？\break

\textbf{【答案】}634尺。\break

\textbf{【解法】}用钱数50,720作为实，每尺80文作为法，九归除法（八归）：50,720 ÷ 80 = 634尺。口诀用「八一下加二，八四添作五，逢八进成十」。



\switchcolumn*



\textbf{問57.} 今有木幾何，直錢三千二百六十株。問每株價錢幾何？\break

\textbf{答：}三貫二百六十文。\break

\textbf{術：}列錢物於上為實，各以價錢為法，從上身外減之，即得。



\switchcolumn



\textbf{【问题57】} 已知买木共花了能买3260株木的钱，每株价格多少？\break

\textbf{【答案】}3贯260文。\break

\textbf{【解法】}把总钱数放在上面作为实，3260株作为法，用九归除法计算每株的价格。这是上一题的逆运算。



\switchcolumn*



\textbf{問58.} 今有絹幾何，每匹四貫，共錢一千四百一十六貫。問匹數幾何？\break

\textbf{答：}三百五十四匹。\break

\textbf{術：}列錢數為實，以匹價為法，九歸除之，即得。



\switchcolumn



\textbf{【问题58】} 已知每匹绢4贯，共钱1416贯。问买了多少匹？\break

\textbf{【答案】}354匹。\break

\textbf{【解法】}用钱数1416作为实，每匹4贯作为法，九归除法：1416 ÷ 4 = 354匹。



\switchcolumn*



\textbf{問59.} 今有銀幾何，每錠五十兩，共二萬七千三百五十兩。問錠數幾何？\break

\textbf{答：}五百四十七錠。\break

\textbf{術：}列銀數為實，以錠重為法，九歸除之，即得。



\switchcolumn



\textbf{【问题59】} 已知每锭银50两，共有27,350两。问有多少锭？\break

\textbf{【答案】}547锭。\break

\textbf{【解法】}用银数27,350作为实，每锭50两作为法：27,350 ÷ 50 = 547锭。



\switchcolumn*



\textbf{問60.} 今有甘草幾何，每秤七百文，共錢六百二十四貫四百文。問秤數幾何？\break

\textbf{答：}八百九十二秤。\break

\textbf{術：}列錢數為實，以秤價為法，九歸除之，即得。



\switchcolumn



\textbf{【问题60】} 已知每秤甘草700文，共钱624贯400文。问有多少秤？\break

\textbf{【答案】}892秤。\break

\textbf{【解法】}用钱数624,400作为实，每秤700文作为法：624,400 ÷ 700 = 892秤。



\switchcolumn*



\textbf{問61.} 今有麻布幾何，每匹一百九十四文，共錢一百四貫三百七十二文。問匹數幾何？\break

\textbf{答：}五百三十八匹。\break

\textbf{術：}列錢數為實，以匹價為法，九歸除之，即得。



\switchcolumn



\textbf{【问题61】} 已知每匹麻布194文，共钱104贯372文（104,372文）。问有多少匹？\break

\textbf{【答案】}538匹。\break

\textbf{【解法】}用钱数104,372作为实，每匹194文作为法：104,372 ÷ 194 = 538匹。



\switchcolumn*



\textbf{問62.} 今有細布幾何，每尺八十文，共錢五十貫七百二十文。問尺數幾何？\break

\textbf{答：}六百三十四尺。\break

\textbf{術：}列錢數為實，以尺價為法，九歸除之，即得。



\switchcolumn



\textbf{【问题62】} 已知每尺细布80文，共钱50贯720文（50,720文）。问有多少尺？\break

\textbf{【答案】}634尺。\break

\textbf{【解法】}用钱数50,720作为实，每尺80文作为法：50,720 ÷ 80 = 634尺。



\switchcolumn*



\textbf{問63.} 今有松木幾何，每株一貫七百五文，共錢五千五百五十八貫三百文。問株數幾何？\break

\textbf{答：}三千二百六十株。\break

\textbf{術：}列錢數為實，以株價為法，九歸除之，即得。



\switchcolumn



\textbf{【问题63】} 已知每株松木1贯705文（1705文），共钱5558贯300文（5,558,300文）。问有多少株？\break

\textbf{【答案】}3260株。\break

\textbf{【解法】}用钱数5,558,300作为实，每株1705文作为法：5,558,300 ÷ 1705 = 3260株。



\switchcolumn*



\textbf{問64.} 今有白豆幾何，每斗二十文，共錢四貫三百二十文。問斗數幾何？\break

\textbf{答：}二百一十六斗。\break

\textbf{術：}列錢數為實，以斗價為法，九歸除之，即得。



\switchcolumn



\textbf{【问题64】} 已知每斗白豆20文，共钱4贯320文（4320文）。问有多少斗？\break

\textbf{【答案】}216斗。\break

\textbf{【解法】}用钱数4320作为实，每斗20文作为法：4320 ÷ 20 = 216斗。



\switchcolumn*



\textbf{問65.} 今有絲幾何，每斤三百文，共錢四十三貫二百文。問斤數幾何？\break

\textbf{答：}一百四十四斤。\break

\textbf{術：}列錢數為實，以斤價為法，九歸除之，即得。



\switchcolumn



\textbf{【问题65】} 已知每斤丝300文，共钱43贯200文（43,200文）。问有多少斤？\break

\textbf{【答案】}144斤。\break

\textbf{【解法】}用钱数43,200作为实，每斤300文作为法：43,200 ÷ 300 = 144斤。



\switchcolumn*



\textbf{問66.} 今有馬幾何，每匹九十貫，共錢三萬八千二百五十貫。問匹數幾何？\break

\textbf{答：}四百二十五匹。\break

\textbf{術：}列錢數為實，以匹價為法，九歸除之，即得。



\switchcolumn



\textbf{【问题66】} 已知每匹马90贯，共钱38,250贯。问有多少匹？\break

\textbf{【答案】}425匹。\break

\textbf{【解法】}用钱数38,250作为实，每匹90贯作为法：38,250 ÷ 90 = 425匹。



\switchcolumn*



\textbf{問67.} 今有羅幾何，每匹價錢三貫二百三十文，共錢一千七百一十一貫四百文。問匹數幾何？\break

\textbf{答：}如術計之。\break

\textbf{術：}列錢數為實，以匹價為法，九歸除之，即得。



\switchcolumn



\textbf{【问题67】} 已知每匹罗3贯230文（3230文），共钱1711贯400文。问有多少匹？\break

\textbf{【答案】}按解法计算。\break

\textbf{【解法】}用钱数1711贯400文（1,711,400文）作为实，每匹3230文作为法：1,711,400 ÷ 3230 = 530匹。



\switchcolumn*



\textbf{問68.} 今有大羅幾何，每匹六貫，共錢四千四百一十六貫。問匹數幾何？\break

\textbf{答：}七百三十六匹。\break

\textbf{術：}列錢數為實，以匹價為法，九歸除之，即得。



\switchcolumn



\textbf{【问题68】} 已知每匹大罗6贯，共钱4416贯。问有多少匹？\break

\textbf{【答案】}736匹。\break

\textbf{【解法】}用钱数4416作为实，每匹6贯作为法：4416 ÷ 6 = 736匹。



\switchcolumn*



\textbf{問69.} 今有苧絲幾何，每尺一百九十文，共錢六十五貫五百五十文。問尺數幾何？\break

\textbf{答：}三百四十五尺。\break

\textbf{術：}列錢數為實，以尺價為法，九歸除之，即得。



\switchcolumn



\textbf{【问题69】} 已知每尺苎丝190文，共钱65贯550文（65,550文）。问有多少尺？\break

\textbf{【答案】}345尺。\break

\textbf{【解法】}用钱数65,550作为实，每尺190文作为法：65,550 ÷ 190 = 345尺。



\switchcolumn*



\textbf{問70.} 今有米幾何，每斗一百一十文，共錢七貫五百二十四文。問斗數幾何？\break

\textbf{答：}六碩八斗四升。\break

\textbf{術：}列錢數為實，以斗價為法，九歸除之，即得。



\switchcolumn



\textbf{【问题70】} 已知每斗米110文，共钱7贯524文（7524文）。问有多少斗？\break

\textbf{【答案】}6硕8斗4升（即68.4斗）。\break

\textbf{【解法】}用钱数7524作为实，每斗110文作为法：7524 ÷ 110 = 68.4斗 = 6硕8斗4升。



\switchcolumn*



\textbf{問71.} 今有鹽幾何，每袋一十三貫，共錢一萬一千三百四十九貫。問袋數幾何？\break

\textbf{答：}八百七十三袋。\break

\textbf{術：}列錢數為實，以袋價為法，九歸除之，即得。



\switchcolumn



\textbf{【问题71】} 已知每袋盐13贯，共钱11,349贯。问有多少袋？\break

\textbf{【答案】}873袋。\break

\textbf{【解法】}用钱数11,349作为实，每袋13贯作为法：11,349 ÷ 13 = 873袋。



\switchcolumn*



\textbf{問72.} 今有羊幾何，每只四貫，共錢一千四百一十六貫。問只數幾何？\break

\textbf{答：}三百五十四只。\break

\textbf{術：}列錢數為實，以只價為法，九歸除之，即得。



\switchcolumn



\textbf{【问题72】} 已知每只羊4贯，共钱1416贯。问有多少只？\break

\textbf{【答案】}354只。\break

\textbf{【解法】}用钱数1416作为实，每只4贯作为法：1416 ÷ 4 = 354只。



\switchcolumn*



\textbf{問73.} 今有油六碩八斗四升，每三斤七分五釐直錢一百文。問直錢幾何？\break

\textbf{答：}如術計之。\break

\textbf{術：}列油數為實，以三斤七分五釐為法，九歸除之，即得。



\switchcolumn



\textbf{【问题73】} 现有油6硕8斗4升，每3斤7分5釐（3.075斤）油值100文。问这些油值多少钱？\break

\textbf{【答案】}按解法计算。\break

\textbf{【解法】}把油的总数放在上面作为实，用3.075作为法（因为3.075斤油值100文），九归除法计算：油数 ÷ 3.075 × 100 = 所值钱数。



\switchcolumn*



\textbf{問74.} 今有麵四百單七斤，每六斤八分七釐半直錢一百文。問直錢幾何？\break

\textbf{答：}如術計之。\break

\textbf{術：}列麵數為實，以六斤八分七釐半為法，九歸除之，即得。



\switchcolumn



\textbf{【问题74】} 现有面407斤，每6斤8分7釐半（6.875斤）面值100文。问这些面值多少钱？\break

\textbf{【答案】}按解法计算。\break

\textbf{【解法】}把面数407放在上面作为实，用6.875作为法：407 ÷ 6.875 × 100 = 所值钱数。



\switchcolumn*



\textbf{問75.} 今有羅三十四尺六寸，每尺一百二十文。問計錢幾何？\break

\textbf{答：}四貫一百五十二文。\break

\textbf{術：}列羅數於上，以尺價為法，九歸除之，即得。



\switchcolumn



\textbf{【问题75】} 现有罗34尺6寸（34.6尺），每尺120文。问总共多少钱？\break

\textbf{【答案】}4贯152文。\break

\textbf{【解法】}把罗数34.6放在上面，每尺120文去乘：34.6 × 120 = 4152文 = 4贯152文。



\switchcolumn*



\textbf{問76.} 今有胡椒六十三斤四兩，每斤三百八十文。問計錢幾何？\break

\textbf{答：}二十四貫三十五文。\break

\textbf{術：}列椒數於上，以斤價為法，九歸除之，即得。



\switchcolumn



\textbf{【问题76】} 现有胡椒63斤4两（63.25斤），每斤380文。问总共多少钱？\break

\textbf{【答案】}24贯35文（24,035文）。\break

\textbf{【解法】}把椒数63.25放在上面，每斤380文去乘：63.25 × 380 = 24,035文。



\switchcolumn*



\textbf{問77.} 今有馬幾何，每匹價錢九十貫。問計錢幾何？\break

\textbf{答：}如術計之。\break

\textbf{術：}列馬數於上，以匹價為法，九歸除之，即得。



\switchcolumn



\textbf{【问题77】} 现有马若干，每匹90贯。问总共多少钱？\break

\textbf{【答案】}按解法计算。\break

\textbf{【解法】}把马匹数放在上面，用每匹90贯去乘，即得总钱数。



\switchcolumn*



\textbf{問78.} 今有綿三百四十五斤，欲作綿被，每床用綿七斤。問作幾床？\break

\textbf{答：}四十九床，餘二斤。\break

\textbf{術：}列綿數為實，以每床用綿為法，九歸除之，即得。



\switchcolumn



\textbf{【问题78】} 现有绵345斤，要做棉被，每床用绵7斤。问能做多少床？\break

\textbf{【答案】}49床，余2斤。\break

\textbf{【解法】}把绵数345放在上面作为实，每床7斤作为法：345 ÷ 7 = 49床余2斤。口诀用「七三四十二，七四五十五，逢七进成十」。



\end{paracol}



\longline



%% =======================================================================

\subsection{異乘同除門 · 异乘同除门（比例运算，8题）}

\label{sec:yicheng}

%% =======================================================================



\begin{paracol}{2}

\noindent\textbf{\textcolor{leftlabel}{【章旨】}}

\switchcolumn

\noindent\textbf{\textcolor{rightlabel}{【章节说明】}}

\switchcolumn*



此門論異乘同除之術，即比例換算之法也。以物易物、據錢買物、據物賣錢，皆以此術求之。其法：以原物乘今錢為實，以原錢為法，除之即得。



\switchcolumn



这一章讨论异乘同除的算法，也就是比例换算的方法。用物品交换物品、根据钱数买物品、根据物品算价钱，都用这个方法来求解。算法是：用原来的物品数量乘以现在的钱数作为被除数，用原来的钱数作为除数，除之即得。



\switchcolumn*



\textbf{問79.} 今有錢九貫八百七十九文，糴米五碩三斛四升。只有錢六十八貫二百六十文。問得米幾何？\break

\textbf{答：}六十八貫二百六十五文。\break

\textbf{術：}列只有米數，以乘今有錢為實，以原錢為法，除之，即得。



\switchcolumn



\textbf{【问题79】} 已知9贯879文钱能买米5硕3斛4升。现有68贯260文钱，问能买多少米？\break

\textbf{【答案】}68贯265文。\break

\textbf{【解法】}用异乘同除法：现在的钱数68,260 × 原来买的米数5.34 ÷ 原来的钱数9,879 = 现在能买的米数。\break

即 (68,260 × 5.34) ÷ 9,879 = 所求米数。



\switchcolumn*



\textbf{問80.} 今有絲七匹，通尺內子得二千三十一尺，以乘上位得六百三十七萬九千三百七十一。又以實以斤銖法三百八十四銖除之，得三萬六千五百四十二銖，為實。\break

\textbf{答：}九千七百六十五斤一十兩四銖。\break

\textbf{術：}列七匹通尺內子，以乘上位，又以斤銖法除之，即得。



\switchcolumn



\textbf{【问题80】} 现有丝7匹，换算成尺并加上零数共得2031尺，再乘以上位的数得到6,379,371。然后用斤铢法384铢去除，得到36,542铢作为被除数。\break

\textbf{【答案】}9765斤10两4铢。\break

\textbf{【解法】}把7匹丝换算为尺数（含零数），再乘以上位给定的数，然后用384铢（1斤=16两=384铢）去除，得到斤数：36,542 ÷ 384 = 95斤余62铢，再换算……\break

\textit{注：原文答案为「九千七百六十五斤」，疑有更复杂的换算过程。}



\switchcolumn*



\textbf{問81.} 今有絲八斤二兩二十一銖，織錦七匹六尺五寸。只有絲九十五斤三兩六銖。問織錦幾何？\break

\textbf{答：}如術計之。\break

\textbf{術：}（匹法同前）列只有絲數，以乘錦數為實，以原絲為法，除之，即得。



\switchcolumn



\textbf{【问题81】} 已知丝8斤2两21铢可以织锦7匹6尺5寸。现有丝95斤3两6铢。问能织多少锦？\break

\textbf{【答案】}按解法计算。\break

\textbf{【解法】}用异乘同除法：现有的丝数 × 原来织的锦数 ÷ 原来的丝数 = 现在能织的锦数。\break

先把所有重量统一换算为铢，然后计算：(95斤3两6铢 × 7匹6尺5寸) ÷ (8斤2两21铢) = 所求锦数。



\switchcolumn*



\textbf{問82.} 今有錢五貫八百三十文，買麩八斗五升。只有錢三十七貫七百六十文。問買麩幾何？\break

\textbf{答：}五碩五斗二升。\break

\textbf{術：}列只有錢數，以乘原麩為實，以原錢為法，除之，即得。



\switchcolumn



\textbf{【问题82】} 已知5贯830文钱能买麸8斗5升。现有37贯760文钱。问能买多少麸？\break

\textbf{【答案】}5硕5斗2升（即55.2斗）。\break

\textbf{【解法】}用异乘同除法：现在的钱数37,760 × 原来买的麸数8.5 ÷ 原来的钱数5,830 = 现在能买的麸数。\break

(37,760 × 8.5) ÷ 5,830 = 55.2斗 = 5硕5斗2升。



\switchcolumn*



\textbf{問83.} 今有錢一貫一百七十五文，買麥五斗四升。只有錢二十三貫五百文。問買麥幾何？\break

\textbf{答：}十碩八斗。\break

\textbf{術：}列只有錢數，以乘原麥為實，以原錢為法，除之，即得。



\switchcolumn



\textbf{【问题83】} 已知1贯175文钱能买麦5斗4升。现有23贯500文钱。问能买多少麦？\break

\textbf{【答案】}10硕8斗（即108斗）。\break

\textbf{【解法】}用异乘同除法：23,500 × 5.4 ÷ 1,175 = 108斗 = 10硕8斗。



\switchcolumn*



\textbf{問84.} 今有錢二貫六百文，買麵三斤四兩。只有錢四十七貫四百五十文。問買麵幾何？\break

\textbf{答：}六十三斤四兩。\break

\textbf{術：}列只有錢數，以乘原麵為實，以原錢為法，除之，即得。



\switchcolumn



\textbf{【问题84】} 已知2贯600文钱能买面3斤4两。现有47贯450文钱。问能买多少面？\break

\textbf{【答案】}63斤4两。\break

\textbf{【解法】}用异乘同除法：47,450 × 3.25 ÷ 2,600 = 63.25斤 = 63斤4两。



\switchcolumn*



\textbf{問85.} 今有錢一貫二百文，買鹽二碩四斗。只有錢三十二貫三百文。問買鹽幾何？\break

\textbf{答：}六十四碩八斗。\break

\textbf{術：}列只有錢數，以乘原鹽為實，以原錢為法，除之，即得。



\switchcolumn



\textbf{【问题85】} 已知1贯200文钱能买盐2硕4斗。现有32贯300文钱。问能买多少盐？\break

\textbf{【答案】}64硕8斗（即648斗）。\break

\textbf{【解法】}用异乘同除法：32,300 × 24 ÷ 1,200 = 646斗……\break

按原文计算：32,300 × 24 ÷ 1,200 = 646斗。原文答案为「六十四硕八斗」= 648斗。



\switchcolumn*



\textbf{問86.} 今有銀二十四銖，直錢三貫六百文。只有銀五十九兩八銖。問直錢幾何？\break

\textbf{答：}八十九貫七百文。\break

\textbf{術：}列只有銀數，以乘原直錢為實，以原銀為法，除之，即得。



\switchcolumn



\textbf{【问题86】} 已知24铢银值3贯600文钱。现有银59两8铢。问值多少钱？\break

\textbf{【答案】}89贯700文。\break

\textbf{【解法】}先把59两8铢换算为铢：59×24 + 8 = 1416 + 8 = 1424铢。\break

用异乘同除法：1424 × 3600 ÷ 24 = 213,600文 = 213贯600文……\break

\textit{注：原文答案为89贯700文，疑有不同的换算方式。}



\end{paracol}



\longline



%% =======================================================================

\subsection{庫務解稅門 · 库务解税门（财政税收，11题）}

\label{sec:kuwu}

%% =======================================================================



\begin{paracol}{2}

\noindent\textbf{\textcolor{leftlabel}{【章旨】}}

\switchcolumn

\noindent\textbf{\textcolor{rightlabel}{【章节说明】}}

\switchcolumn*



此門論庫務、解稅、利息、錢鈔換算之事，皆官府財政之實務也。其題涉及本利計算、錢鈔兌換、合金成色、課稅徵收，反映了元代財政管理之制度。



\switchcolumn



这一章讨论国库管理、税收解送、利息计算、钱币兑换等问题，都是政府财政管理的实际事务。题目涉及本金利息计算、纸币银两兑换、合金成色、课税征收等内容，反映了元代的财政管理制度。



\switchcolumn*



\textbf{問87.} 今有本銀四千五百六兩，經三箇月一十二日，月利銀三百六兩。問本利共幾何？\break

\textbf{答：}本銀四千五十兩，利銀三百六兩。\break

\textbf{術：}置本銀以月數乘之，又以日數乘之，以三十日除之，得利息，加入本銀，即得。



\switchcolumn



\textbf{【问题87】} 现有本金4050两，经过3个月12天，每月利息按比例计算。问本金和利息共多少？\break

\textbf{【答案】}本金4050两，利息306两。\break

\textbf{【解法】}把本金放在上面，用月数去乘，再用日数去乘，除以30天，得到利息，加入本金，即得本利合计。\break

计算：利息 = 4050 × 3 × 12 ÷ 30 = 4050 × 36 ÷ 30 = 4860两……\break

\textit{注：原文「月利銀三百六兩」疑为已知条件，即每月利息按比例相当于306两。实际计算方法为：本金 × 月利率 × 时间 = 利息。}



\switchcolumn*



\textbf{問88.} 今有本銀四千五十兩，月利一分二釐。問三箇月一十二日，利銀幾何？\break

\textbf{答：}一百四十五兩八錢。\break

\textbf{術：}置本銀以月利乘之，又以日數乘之，以三十日除之，即得。



\switchcolumn



\textbf{【问题88】} 现有本金4050两，月利率为1分2釐（即1.2\%）。问3个月12天的利息是多少？\break

\textbf{【答案】}145两8钱。\break

\textbf{【解法】}把本金4050放在上面，乘以月利率1.2\%：4050 × 0.012 = 48.6两（每月利息）。\break

3个月利息 = 48.6 × 3 = 145.8两 = 145两8钱。\break

12天的利息 = 48.6 × 12 ÷ 30 = 19.44两，合计165.24两……\break

\textit{注：原文答案为145两8钱，疑只计算了3个月（不含12天），4050 × 0.012 × 3 = 145.8两。}



\switchcolumn*



\textbf{問89.} 今有課銀三萬六千兩，每五十兩為一錠。問為幾錠？\break

\textbf{答：}七百二十錠。\break

\textbf{術：}置課銀以錠重除之，即得。



\switchcolumn



\textbf{【问题89】} 现有课银36,000两，每50两铸成1锭。问能铸多少锭？\break

\textbf{【答案】}720锭。\break

\textbf{【解法】}把课银36,000放在上面，除以每锭50两：36,000 ÷ 50 = 720锭。



\switchcolumn*



\textbf{問90.} 今有鈔二千四百貫，欲換銀，每銀一兩直鈔八貫。問得銀幾何？\break

\textbf{答：}三百兩。\break

\textbf{術：}置鈔數為實，以銀價為法，除之，即得。



\switchcolumn



\textbf{【问题90】} 现有纸币2400贯，要换成白银，每两白银值纸币8贯。问能换多少两？\break

\textbf{【答案】}300两。\break

\textbf{【解法】}把钞数2400贯放在上面作为被除数，每两银价8贯作为除数：2400 ÷ 8 = 300两。



\switchcolumn*



\textbf{問91.} 今有銀三百兩，欲換鈔，每兩直鈔八貫。問得鈔幾何？\break

\textbf{答：}二千四百貫。\break

\textbf{術：}置銀數為實，以鈔價為法，乘之，即得。



\switchcolumn



\textbf{【问题91】} 现有银300两，要换成纸币，每两银值纸币8贯。问能换多少贯？\break

\textbf{【答案】}2400贯。\break

\textbf{【解法】}把银数300放在上面，乘以每两银价8贯：300 × 8 = 2400贯。这是上一题的逆运算。



\switchcolumn*



\textbf{問92.} 今有金一十二兩五錢，內減五十兩，餘即入銀，合同為法。實如法而一，得本銀，反減共還銀數，餘即利銀也。\break

\textbf{答：}一十二兩五錢。\break

\textbf{術：}列金數為實，以八分為法，除之，得本銀，反減共還銀數，餘即利銀也。



\switchcolumn



\textbf{【问题92】} 现有金12两5钱，题目涉及合金成色的计算。用特定的方法算出本金银，再用总还款数减去本金，余数就是利息。\break

\textbf{【答案】}12两5钱。\break

\textbf{【解法】}把金数放在上面作为被除数，以8分（0.8）为除数相除，得到本银数，再用总还款数反减本银数，余数就是利息。这是一道涉及金融计算的题目。



\switchcolumn*



\textbf{問93.} 今有銀一十二兩五錢，內減五十兩，餘即入銀，合同為法。實如法而一，得六兩，問為顏色分數幾何？\break

\textbf{答：}八分色。\break

\textbf{術：}列金數為實，以本銀為法，除之，即得顏色分數。



\switchcolumn



\textbf{【问题93】} 现有银12两5钱，经过某种合金计算得到6两，问合金的成色（纯度）是多少？\break

\textbf{【答案】}八分色（即80\%的纯度）。\break

\textbf{【解法】}把金（银）数放在上面作为被除数，以本银为除数：纯度 = 实际银含量 ÷ 总重量。如 10 ÷ 12.5 = 0.8 = 八分色（80\%纯度）。



\switchcolumn*



\textbf{問94.} 今有絲六十二斤半，換金一十兩，內八分半金五兩七分半金五兩。問二色金各得幾何？\break

\textbf{術：}列絲數為實，以金價為法，除之，即得各金數。



\switchcolumn



\textbf{【问题94】} 现有丝62.5斤，要换成10两金，其中一种是成色8.5分的金5两，另一种是成色7.5分的金5两。问两种金子各得多少？\break

\textbf{【解法】}把丝数62.5斤放在上面作为被除数，以金价为除数来计算。这是关于不同成色金子的兑换问题：两种金子的平均成色决定了兑换比率。\break

8.5分金5两 + 7.5分金5两 = 平均成色8分金10两。用62.5斤丝去换，按金价计算各得多少。



\switchcolumn*



\textbf{問95.} 今有銀三千五百六兩，直錢三十五貫六百文。問每兩直錢幾何？\break

\textbf{答：}一十文。\break

\textbf{術：}置銀數為實，以直錢為法，除之，即得。



\switchcolumn



\textbf{【问题95】} 现有银3506两，总共值35贯600文。问每两值多少钱？\break

\textbf{【答案】}每两10文。\break

\textbf{【解法】}把银数3506放在上面，总钱数35,600文为除数：35,600 ÷ 3506 ≈ 10.15文……\break

\textit{注：原文答案为「一十文」，疑银数应为3560两，则35,600 ÷ 3560 = 10文。}



\switchcolumn*



\textbf{問96.} 今有鈔一百六十五貫五百文，買絲三十七斤八兩。問每斤價鈔幾何？\break

\textbf{答：}四貫四百文。\break

\textbf{術：}置鈔數為實，以絲數為法，除之，即得。



\switchcolumn



\textbf{【问题96】} 现有钞165贯500文（165,500文），买丝37斤8两（即37.5斤）。问每斤丝价格多少？\break

\textbf{【答案】}每斤4贯400文。\break

\textbf{【解法】}把钞数165,500文放在上面作为被除数，丝数37.5斤作为除数：165,500 ÷ 37.5 = 4413.33文……\break

\textit{注：原文答案为每斤4贯400文，即4400文。若165,000 ÷ 37.5 = 4400文，则原文「五百文」应为零数。}



\switchcolumn*



\textbf{問97.} 今有課鈔一千貫，納官每百貫收工墨錢一貫。問工墨錢幾何？\break

\textbf{答：}十貫。\break

\textbf{術：}置課鈔為實，以百貫為法，除之，即得。



\switchcolumn



\textbf{【问题97】} 现有课钞1000贯要缴纳给官府，每100贯收取工墨钱（手续费）1贯。问工墨钱是多少？\break

\textbf{【答案】}10贯。\break

\textbf{【解法】}把课钞1000贯放在上面，每100贯为单位：1000 ÷ 100 = 10贯工墨钱。\break

这是元代税收制度中的手续费计算。



\end{paracol}



\longline



%% =======================================================================

\subsection{折變互差門 · 折变互差门（比例换算，15题）}

\label{sec:zhebian}

%% =======================================================================



\begin{paracol}{2}

\rhymeCL{折變互差要詳知，多少相乘作實推，\\

以少求多因作母，除來便見兩無虧。}

\switchcolumn

\rhymeMOD{折变互差的方法要详细了解，\\

把已知的数量和比率相乘作为被除数，\\

以少求多时就把「少」作为除数，\\

除出来的结果双方都不亏。}

\switchcolumn*

\noteCL{折變互差者，論物價折算、合金成色、田地測量之術也。以已知之數推未知之數，用比例之法求解。}

\switchcolumn

\noteMOD{折变互差是讨论物价折算、合金成色、田地测量等算法。用已知数来推算未知数，用比例的方法来求解。}



\switchcolumn*



\textbf{問98.} 今有香油三兩，折菜油四兩，只云香油斤價四百文。問菜油斤價幾何？\break

\textbf{答：}二十五貫四百二十五文。\break

\textbf{術：}列香油數為實，以折率乘之，又以香油價為法，除之，即得。



\switchcolumn



\textbf{【问题98】} 已知3两香油可以折算为4两菜油，又知香油每斤价格400文。问菜油每斤价格多少？\break

\textbf{【答案】}25贯425文。\break

\textbf{【解法】}用比例法：香油3两 : 菜油4两 = 香油价 : 菜油价。\break

菜油价 = 香油价 × 香油量 ÷ 菜油量 = 400 × 3 ÷ 4 = 300文/斤……\break

\textit{注：原文答案为「二十五贯四百二十五文」，疑有不同的换算方式。}



\switchcolumn*



\textbf{問99.} 今有香油三兩，折菜油四兩，只云菜油八十四斤一十二兩，問香油幾何？\break

\textbf{答：}如術計之。\break

\textbf{術：}列菜油數為實，以折率除之，即得。



\switchcolumn



\textbf{【问题99】} 已知3两香油折算为4两菜油，现有菜油84斤12两（84.75斤）。问相当于多少香油？\break

\textbf{【答案】}按解法计算。\break

\textbf{【解法】}用比例法：香油 = 菜油 × 3 ÷ 4 = 84.75 × 3 ÷ 4 = 63.5625斤 = 63斤9两。



\switchcolumn*



\textbf{問100.} 今有麵四百單七斤，每六斤八分七釐半直錢一百文。問直錢幾何？\break

\textbf{答：}如術計之。\break

\textbf{術：}列麵數為實，以六斤八分七釐半為法，除之，即得。



\switchcolumn



\textbf{【问题100】} 现有面407斤，每6斤8分7釐半（6.875斤）面值100文。问这些面值多少钱？\break

\textbf{【答案】}按解法计算。\break

\textbf{【解法】}把面数407放在上面，6.875斤作为除数：407 ÷ 6.875 × 100 = 407 × 100 ÷ 6.875 = 5920文 = 5贯920文。



\switchcolumn*



\textbf{問101.} 今有油六碩八斗四升，每三斤七分五釐直錢一百文。問直錢幾何？\break

\textbf{答：}如術計之。\break

\textbf{術：}列油數為實，以三斤七分五釐為法，除之，即得。



\switchcolumn



\textbf{【问题101】} 现有油6硕8斗4升，每3斤7分5釐（3.075斤）油值100文。问这些油值多少钱？\break

\textbf{【答案】}按解法计算。\break

\textbf{【解法】}把油的总数量放在上面，3.075作为除数：油数 ÷ 3.075 × 100 = 所值钱数。



\switchcolumn*



\textbf{問102.} 今有麵數為實，以六斤八分七釐半除之。問得幾何？\break

\textbf{術：}列麵數為實，以六斤八分七釐半為法，除之，即得。



\switchcolumn



\textbf{【问题102】} 现有面若干作为被除数，用6斤8分7釐半（6.875斤）去除。问得多少？\break

\textbf{【解法】}把面数放在上面作为被除数（实），用6.875斤作为除数（法），除之即得。这是面数和单价之间的换算。



\switchcolumn*



\textbf{問103.} 今有羅三十四匹六尺，每尺一百二十文。問計錢幾何？\break

\textbf{答：}四貫一百五十二文。\break

\textbf{術：}列羅數通為尺，以尺價乘之，即得。



\switchcolumn



\textbf{【问题103】} 现有罗34匹6尺，每尺120文。问总共多少钱？\break

\textbf{【答案】}4贯152文。\break

\textbf{【解法】}先把34匹6尺全部换算为尺数（按1匹=某尺数），再乘以每尺120文：总尺数 × 120 = 总钱数 = 4152文。



\switchcolumn*



\textbf{問104.} 今有絹七匹，每匹四貫。問計錢幾何？\break

\textbf{答：}二十八貫。\break

\textbf{術：}列絹數，以匹價乘之，即得。



\switchcolumn



\textbf{【问题104】} 现有绢7匹，每匹4贯。问总共多少钱？\break

\textbf{【答案】}28贯。\break

\textbf{【解法】}把绢的匹数7放在上面，乘以每匹4贯：7 × 4 = 28贯。



\switchcolumn*



\textbf{問105.} 今有絲一斤價錢二貫文，問一兩價錢幾何？\break

\textbf{答：}一百二十五文。\break

\textbf{術：}置絲斤價為實，以十六兩為法，除之，即得。



\switchcolumn



\textbf{【问题105】} 现有丝每斤价格2贯（2000文），问每两价格多少？\break

\textbf{【答案】}125文。\break

\textbf{【解法】}把丝的斤价2000文放在上面，除以16两：2000 ÷ 16 = 125文/两。



\switchcolumn*



\textbf{問106.} 今有田一畝，闊十五步，長十六步。問為頃畝幾何？\break

\textbf{答：}一畝。\break

\textbf{術：}列闊步、長步相乘得積步，以畝法二百四十步除之，即得。



\switchcolumn



\textbf{【问题106】} 现有田1亩，宽15步，长16步。问合多少顷亩？\break

\textbf{【答案】}1亩。\break

\textbf{【解法】}把宽15步和长16步相乘得到积步（面积步数）：15 × 16 = 240平方步。\break

用亩法240平方步去除：240 ÷ 240 = 1亩。\break

\textit{注：古代1亩 = 宽1步 × 长240步 = 240平方步。这里15×16=240平方步，正好等于1亩。}



\switchcolumn*



\textbf{問107.} 今有粟九斗，每斛直錢一貫二百文。問直錢幾何？\break

\textbf{答：}一貫八十文。\break

\textbf{術：}置粟數為實，以斛價為法，乘之，即得。



\switchcolumn



\textbf{【问题107】} 现有粟9斗，每斛价格1贯200文（1200文）。问值多少钱？\break

\textbf{【答案】}1贯80文（1080文）。\break

\textbf{【解法】}先把9斗换算为斛：9斗 = 0.9斛。\break

再用0.9斛 × 每斛1200文 = 1080文 = 1贯80文。



\switchcolumn*



\textbf{問108.} 今有金五十兩，令二八共造，問各幾何？\break

\textbf{答：}金四十兩，銀一十兩。\break

\textbf{術：}置共金為實，以分母為法，除之，各以分子乘之，即得。



\switchcolumn



\textbf{【问题108】} 现有金50两，要按「二八」的比例来铸造合金，问金和银各用多少？\break

\textbf{【答案】}金40两，银10两。\break

\textbf{【解法】}「二八」即十分之八为金，十分之二为银。\break

金 = 50 × 8 ÷ 10 = 40两；银 = 50 × 2 ÷ 10 = 10两。



\switchcolumn*



\textbf{問109.} 今有銀二十四兩，直錢三貫六百文。只有銀五十九兩八銖。問直錢幾何？\break

\textbf{答：}八十九貫七百文。\break

\textbf{術：}列只有銀數，以乘原直錢為實，以原銀為法，除之，即得。



\switchcolumn



\textbf{【问题109】} 已知24两银值3贯600文（3600文）。现有银59两8铢。问值多少钱？\break

\textbf{【答案】}89贯700文。\break

\textbf{【解法】}先把59两8铢换算：59两8铢 = 59×24 + 8 = 1416 + 8 = 1424铢。\break

24两 = 24 × 24 = 576铢。\break

用异乘同除法：1424 × 3600 ÷ 576 = 8900文……\break

\textit{注：原文答案为89贯700文（89,700文）。}



\switchcolumn*



\textbf{問110.} 今有絲八斤二兩二十一銖，織錦七匹六尺五寸。只有絲九十五斤三兩六銖。問織錦幾何？\break

\textbf{術：}（匹法同前）列只有絲數，以乘錦數為實，以原絲為法，除之，即得。



\switchcolumn



\textbf{【问题110】} 已知丝8斤2两21铢可以织锦7匹6尺5寸。现有丝95斤3两6铢。问能织多少锦？\break

\textbf{【解法】}用异乘同除法：先把所有重量统一换算为铢，然后\break

现有丝数 × 原来织的锦数 ÷ 原来的丝数 = 现在能织的锦数。



\switchcolumn*



\textbf{問111.} 今有錢五貫八百三十文，買麩八斗五升。只有錢三十七貫七百六十文。問買麩幾何？\break

\textbf{答：}五碩五斗二升。\break

\textbf{術：}列只有錢數，以乘原麩為實，以原錢為法，除之，即得。



\switchcolumn



\textbf{【问题111】} 已知5贯830文能买麸8斗5升。现有37贯760文。问能买多少麸？\break

\textbf{【答案】}5硕5斗2升（55.2斗）。\break

\textbf{【解法】}异乘同除法：37,760 × 8.5 ÷ 5,830 = 55.2斗 = 5硕5斗2升。



\switchcolumn*



\textbf{問112.} 今有錢一貫一百七十五文，買麥五斗四升。只有錢二十三貫五百文。問買麥幾何？\break

\textbf{答：}十碩八斗。\break

\textbf{術：}列只有錢數，以乘原麥為實，以原錢為法，除之，即得。



\switchcolumn



\textbf{【问题112】} 已知1贯175文能买麦5斗4升。现有23贯500文。问能买多少麦？\break

\textbf{【答案】}10硕8斗（108斗）。\break

\textbf{【解法】}异乘同除法：23,500 × 5.4 ÷ 1,175 = 108斗 = 10硕8斗。



\end{paracol}



\vspace{1cm}

\longline



%% ========================================================================

\section*{跋 · 卷上终}

\addcontentsline{toc}{section}{跋 · 卷上终}

%% ========================================================================



\begin{paracol}{2}

\noindent\textbf{\textcolor{leftlabel}{【文言】}}

\switchcolumn

\noindent\textbf{\textcolor{rightlabel}{【白話】}}

\switchcolumn*



右卷上凡八門，計一百一十三問。論縱橫因法、身外加減、留頭乘法、九歸除法、異乘同除、庫務解稅、折變互差之術。由淺入深，循序漸進，使學者習之，庶幾可以進窺卷中、卷下之奧也。



\switchcolumn



以上卷上共八章，计113道题。讨论了纵横因法、身外加减法、留头乘法、九归除法、异乘同除、库务解税、折变互差等算法。由浅入深、循序渐进，使学习者通过练习这些题目，可以为进一步学习卷中、卷下的深奥内容打下基础。



\switchcolumn*



朱氏之書，條理秩然。首列總括一十八條，立全書之綱紀；次分三卷，為初學之階梯。其術雖云啓蒙，實具大匠之規模。後之習算者，當由是而進窺天元、垛積之精微，則算學之能事畢矣。



\switchcolumn



朱世杰的这本书，条理清晰。开头列出总括十八条，确立了全书的纲领；然后分为三卷，作为初学者的阶梯。此书虽说是启蒙，实际上具有大师的格局。后世学习数学的人，应当由此进而探究天元术、垛积术的精妙，那就算数学的精髓都掌握了。



\end{paracol}



\vspace{2cm}

\begin{center}

\textcolor{sectioncolor}{\rule{0.3\textwidth}{0.4pt}}\\[0.5cm]

{\small\color{notecolor} 卷上終 · End of Volume I\\

總計八門 \quad 一百一十三問 \quad 18 General Principles + 113 Problems = 131 Items}

\end{center}



\end{document}







% ============================================================

% Source: translations/zh/chinese/zhu-shijie-suanxue-qimeng-part2_classical-modern_bilingual.tex

% ============================================================



\documentclass[11pt,a4paper]{article}



% Core packages

\usepackage{fontspec}

\usepackage{polyglossia}

\usepackage{amsmath,amssymb,amsthm}

\usepackage{geometry}

\usepackage{graphicx}

\usepackage{xcolor}

\usepackage{fancyhdr}

\usepackage{enumitem}

\usepackage{array,booktabs,longtable}

\usepackage{hyperref}

\usepackage{paracol}

\usepackage{indentfirst}

\usepackage{float}

\usepackage{caption}

\usepackage{subcaption}

\usepackage{tikz}



% Page geometry - wider for bilingual columns

\geometry{

  paperwidth=210mm,paperheight=297mm,

  left=18mm,right=18mm,top=25mm,bottom=25mm,

  headheight=14pt,footskip=30pt

}



% Language setup

\setmainlanguage{english}



% Font configuration

\setmainfont{Times New Roman}

\setsansfont{Arial}

\setmonofont{Courier New}



% xeCJK for Chinese

\usepackage{xeCJK}

\setCJKmainfont{Microsoft YaHei}

\setCJKsansfont{Microsoft YaHei}

\setCJKmonofont{Microsoft YaHei}



% Colors

\definecolor{titlecolor}{RGB}{60,30,10}

\definecolor{sectioncolor}{RGB}{80,40,15}

\definecolor{chaptercolor}{RGB}{100,50,20}

\definecolor{notecolor}{RGB}{120,53,15}

\definecolor{rulecolor}{RGB}{180,160,140}

\definecolor{problemcolor}{RGB}{30,60,90}

\definecolor{classicalbg}{RGB}{255,250,245}

\definecolor{modernbg}{RGB}{245,250,255}

\definecolor{classicalborder}{RGB}{160,120,80}

\definecolor{modernborder}{RGB}{80,120,160}

\definecolor{headerclassical}{RGB}{100,60,30}

\definecolor{headermodern}{RGB}{30,60,100}



% Section formatting

\usepackage{titlesec}

\titleformat{\section}

  {\normalfont\Large\bfseries\color{sectioncolor}}

  {\thesection}{1em}{}

\titleformat{\subsection}

  {\normalfont\large\bfseries\color{sectioncolor}}

  {\thesubsection}{1em}{}

\titleformat{\subsubsection}

  {\normalfont\normalsize\bfseries\color{sectioncolor}}

  {\thesubsubsection}{1em}{}



% Header/Footer

\pagestyle{fancy}

\fancyhf{}

\fancyhead[L]{\small\textcolor{headerclassical}{文言文}\hspace{0.3em}|\hspace{0.3em}\textcolor{headermodern}{白话文} \textit{算學啟蒙·卷中}}

\fancyhead[R]{\small\thepage}

\renewcommand{\headrulewidth}{0.4pt}

\renewcommand{\footrulewidth}{0pt}



% Column ratio for paracol (left=classical, right=modern)

\columnratio{0.50,0.50}



% Custom commands

\newcommand{\longline}{\noindent\rule{\textwidth}{0.4pt}}

\newcommand{\shortline}{\noindent\rule{0.5\textwidth}{0.4pt}\par}

\newcommand{\cnote}[1]{\textcolor{notecolor}{\textit{#1}}}



% Classical text environment (left column)

\newenvironment{classicalbox}{%

  \begin{tikzpicture}

  \node[draw=classicalborder, fill=classicalbg, rounded corners=3pt,

        inner sep=6pt, text width=\linewidth] \bgroup

}{%

  \egroup;

  \end{tikzpicture}

}



% Modern text environment (right column)

\newenvironment{modernbox}{%

  \begin{tikzpicture}

  \node[draw=modernborder, fill=modernbg, rounded corners=3pt,

        inner sep=6pt, text width=\linewidth] \bgroup

}{%

  \egroup;

  \end{tikzpicture}

}



% Bilingual problem environment

\newenvironment{bilingualproblem}[2]{%

  \vspace{0.5em}

  \begin{paracol}{2}

  \begin{leftcolumn}

  \begin{tikzpicture}

  \node[draw=classicalborder, fill=classicalbg, rounded corners=3pt,

        inner sep=6pt, text width=\linewidth] \bgroup

  \textbf{\textcolor{headerclassical}{#1}}\\[0.3em]

}{%

  \egroup;

  \end{tikzpicture}

  \end{leftcolumn}

  \end{paracol}

  \vspace{0.5em}

}



% Problem in classical Chinese

\newcommand{\classicalproblem}[1]{%

  \begin{tikzpicture}

  \node[draw=classicalborder, fill=classicalbg, rounded corners=3pt,

        inner sep=6pt, text width=\linewidth] {#1};

  \end{tikzpicture}

}



% Problem in modern Chinese

\newcommand{\modernproblem}[1]{%

  \begin{tikzpicture}

  \node[draw=modernborder, fill=modernbg, rounded corners=3pt,

        inner sep=6pt, text width=\linewidth] {#1};

  \end{tikzpicture}

}



% Column headers

\newcommand{\classicalheader}{%

  \begin{center}

  \textbf{\textcolor{headerclassical}{\large 文言文（原文）}}

  \end{center}

  \vspace{-0.3em}

  \hrule height 0.5pt \color{classicalborder}

  \vspace{0.5em}

}

\newcommand{\modernheader}{%

  \begin{center}

  \textbf{\textcolor{headermodern}{\large 白话文（译文）}}

  \end{center}

  \vspace{-0.3em}

  \hrule height 0.5pt \color{modernborder}

  \vspace{0.5em}

}



% Hyperref

\hypersetup{

  colorlinks=true,

  linkcolor=sectioncolor,

  citecolor=sectioncolor,

  urlcolor=sectioncolor,

  pdfencoding=auto,

  pdftitle={算學啟蒙·卷中 --- 文言文/白话文 对照版},

  pdfauthor={朱世傑（原著）},

}



% TOC depth

\setcounter{tocdepth}{2}

\setcounter{secnumdepth}{2}



% Synchronized columns

\switchcolumn*\newcommand{\syncmark}{\switchcolumn*}



\begin{document}



% ============================================

% TITLE PAGE

% ============================================

\begin{titlepage}

\centering

\vspace*{1.5cm}



{\fontsize{12}{14}\selectfont\textsc{中国古代数学经典双语对照丛书}}

\vspace{0.5cm}



{\fontsize{10}{12}\selectfont\textit{Chinese Classical Mathematical Texts --- Bilingual Edition}}

\vspace{0.8cm}



\longline

\vspace{1.2cm}



{\fontsize{26}{32}\selectfont\bfseries\color{titlecolor}

算學啟蒙·卷中\\[0.5cm]

}

\vspace{0.3cm}



{\fontsize{16}{20}\selectfont\color{sectioncolor}

\textit{Suanxue Qimeng} --- Volume II (Middle)\\[0.3cm]

}



\vspace{0.5cm}



{\fontsize{22}{26}\selectfont\color{chaptercolor}

文言文 / 白话文 对照版\\[0.2cm]

{\large Classical Chinese / Modern Chinese Bilingual Edition}

}



\vspace{1.5cm}

\longline

\vspace{0.8cm}



{\large 元·松庭朱世傑 編撰\\[0.3cm]

\textit{Compiled by Zhu Shijie (Songting), Yuan Dynasty}\\[0.5cm]

}



\vspace{0.3cm}



{\large 1299年刊于扬州\\[0.3cm]

\textit{Published in Yangzhou, 1299 CE}\\[0.5cm]

}



\vspace{0.8cm}



\begin{minipage}{0.85\textwidth}

\centering

\textbf{卷中七章七十一问：}\\[0.3cm]

田畝形段门 · 仓屯积粟门 · 双据互换门 · 求差分和门 · 差分均配门 · 商功修筑门 · 贵贱反率门

\end{minipage}



\vspace{1cm}



\begin{tikzpicture}

\node[draw=classicalborder, fill=classicalbg, rounded corners=5pt, inner sep=12pt, text width=0.75\textwidth, align=center] {

\textbf{双语对照说明}\\[0.3em]

\textbf{左栏：}文言文原文（古典数学术语保留原貌）\\

\textbf{右栏：}白话文译文（现代汉语解释与注释）

};

\end{tikzpicture}



\vfill



{\small\textit{以 \LaTeX\ 及 \textsf{XeLaTeX} 排版，使用 Arial CJK 字体}\\

\today}



\end{titlepage}



% ============================================

% TABLE OF CONTENTS

% ============================================

\tableofcontents

\newpage



% ============================================

% INTRODUCTION

% ============================================

\section{简介 · Introduction}



\begin{paracol}{2}

\begin{leftcolumn}

\classicalheader

\end{leftcolumn}

\begin{rightcolumn}

\modernheader

\end{rightcolumn}

\end{paracol}



\vspace{0.3em}



\begin{paracol}{2}

\begin{leftcolumn}

《算學啓蒙》者，元燕山朱世傑松庭之所撰也。書成于大德三年（1299年），刻于揚州。凡三卷，計二十門，二百五十九問。自淺入深，由加減乘除而至天元術，誠算學之津梁也。

\end{leftcolumn}

\begin{rightcolumn}

《算學啓蒙》是元代燕山（今北京附近）人朱世傑（字松庭）所撰写的数学教科书。此书成书于元大德三年（1299年），在扬州刊刻。全书共三卷、二十门、二百五十九道算题。内容由浅入深，从加减乘除四则运算一直讲到天元术（一元高次方程），真可谓数学入门的桥梁。

\end{rightcolumn}

\end{paracol}



\vspace{0.5em}



\begin{paracol}{2}

\begin{leftcolumn}

世傑周遊四方二十余年，學徒雲集。其學尤精于天元、四元、垛積、招差諸術。所著之書，今存者唯《算學啓蒙》與《四元玉鑑》二種而已。

\end{leftcolumn}

\begin{rightcolumn}

朱世傑周游四方二十余年，前来求学的人从各地云集而来。他的学问尤其精通天元术、四元术、垛積术和招差术。他所写的著作，如今保存下来的只有《算學啓蒙》和《四元玉鑑》两种。

\end{rightcolumn}

\end{paracol}



\subsection{卷中结构 · Structure of Volume II}



\begin{paracol}{2}

\begin{leftcolumn}

\classicalheader

\end{leftcolumn}

\begin{rightcolumn}

\modernheader

\end{rightcolumn}

\end{paracol}



\begin{paracol}{2}

\begin{leftcolumn}

卷中凡七門，七十一問：

\begin{enumerate}[noitemsep]

\item 田畝形段門（十六問）

\item 倉屯積粟門（九問）

\item 雙據互換門（六問）

\item 求差分和門（九問）

\item 差分均配門（十問）

\item 商功修築門（十三問）

\item 貴賤反率門（八問）

\end{enumerate}

\end{leftcolumn}

\begin{rightcolumn}

第二卷共分七章，包含七十一道算题：

\begin{enumerate}[noitemsep]

\item 田畝形段门（16题）--- 田地面积计算

\item 仓屯积粟门（9题）--- 粮仓容积计算

\item 双据互换门（6题）--- 复合比例换算

\item 求差分和门（9题）--- 和差问题

\item 差分均配门（10题）--- 按比例分配

\item 商功修筑门（13题）--- 工程土方计算

\item 贵贱反率门（8题）--- 反比例（价格与比率）问题

\end{enumerate}

\end{rightcolumn}

\end{paracol}



\newpage



% ============================================

% CHAPTER 1: 田畝形段門

% ============================================

\section{第一章：田畝形段門 · Field Area Computations}



\begin{paracol}{2}

\begin{leftcolumn}

\classicalheader

\end{leftcolumn}

\begin{rightcolumn}

\modernheader

\end{rightcolumn}

\end{paracol}



\begin{paracol}{2}

\begin{leftcolumn}

田畝形段門，計一十六問。此門專論方田、直田、勾股田、梯田、圭田、圓田諸形之面積算法。古法以二百四十平方步為一畝。

\end{leftcolumn}

\begin{rightcolumn}

田畝形段门共十六道题。这一章专门讨论方形田、矩形田、直角三角形田、梯形田、三角形田、圆形田等各种田地形状的面积计算方法。按照古代标准，二百四十平方步等于一畝。

\end{rightcolumn}

\end{paracol}



\vspace{0.5em}



\subsection{第一問：方田 · Square Field}



\begin{paracol}{2}

\begin{leftcolumn}

\classicalproblem{

\textbf{今有}方田一段，自方九十六步，\textbf{問為}田幾何？



\textbf{答曰：}三十八畝四分。



\textbf{術曰：}列九十六步自乘，得九千二百一十六步，為田積也，以畝法二百四十步除之，\textbf{合問}。}

\end{leftcolumn}

\begin{rightcolumn}

\modernproblem{

\textbf{题目：}现有一块正方形田地，每边长九十六步，\textbf{问}这块田的面积是多少？



\textbf{答案：}三十八畝四分（即38.4畝）。



\textbf{解法：}把九十六步自乘（平方），得到九千二百一十六平方步，这就是田地的面积；再用畝法二百四十步去除，\textbf{就得到了答案}。}

\end{rightcolumn}

\end{paracol}



\vspace{0.3em}



\begin{paracol}{2}

\begin{leftcolumn}

\textit{今按：}方田之術，以邊自乘得積，以畝法二百四十除之，即得畝數。此術與《九章算術》方田章同。

\end{leftcolumn}

\begin{rightcolumn}

\textit{现代解释：}正方形田地的计算方法是：边长自乘得到面积，再用二百四十（每畝的平方步数）去除，就得到畝数。这个方法与《九章算术》方田章的算法相同。



\textbf{现代公式：}$A = \dfrac{96^2}{240} = \dfrac{9216}{240} = 38.4$ 畝

\end{rightcolumn}

\end{paracol}



\vspace{0.5em}



\subsection{第二問：直田 · Rectangular Field}



\begin{paracol}{2}

\begin{leftcolumn}

\classicalproblem{

\textbf{今有}直田一段，長四十九步，闊二十四步，\textbf{問為}田幾何？



\textbf{答曰：}四畝九分。



\textbf{術曰：}列長四十九步，以闊二十四步乘之，得一千一百七十六步，為田積也，以畝法二百四十步除之，\textbf{合問}。}

\end{leftcolumn}

\begin{rightcolumn}

\modernproblem{

\textbf{题目：}现有一块矩形田地，长四十九步，宽二十四步，\textbf{问}这块田的面积是多少？



\textbf{答案：}四畝九分（即4.9畝）。



\textbf{解法：}把长四十九步与宽二十四步相乘，得到一千一百七十六平方步，这就是田地的面积；再用畝法二百四十步去除，\textbf{就得到了答案}。}

\end{rightcolumn}

\end{paracol}



\begin{paracol}{2}

\begin{leftcolumn}

\textit{今按：}直田者，長闊相乘得積。其理與方田同，惟邊長不等耳。

\end{leftcolumn}

\begin{rightcolumn}

\textit{现代解释：}矩形田地就是将长和宽相乘得到面积。原理与正方形田地相同，只是边长不相等。



\textbf{现代公式：}$A = \dfrac{49 \times 24}{240} = \dfrac{1176}{240} = 4.9$ 畝

\end{rightcolumn}

\end{paracol}



\vspace{0.5em}



\subsection{第三問：勾股田 · Right-Triangle Field}



\begin{paracol}{2}

\begin{leftcolumn}

\classicalproblem{

\textbf{今有}勾股田一段，勾三十六步，股六十二步，\textbf{問為}田幾何？



\textbf{答曰：}四畝六分五厘。



\textbf{術曰：}列股六十二步，以勾三十六步乘之，折半，得一千一百一十六步，為田積也，以畝法二百四十步除之，\textbf{合問}。}

\end{leftcolumn}

\begin{rightcolumn}

\modernproblem{

\textbf{题目：}现有一块直角三角形田地，勾（底边）三十六步，股（高）六十二步，\textbf{问}这块田的面积是多少？



\textbf{答案：}四畝六分五厘（即4.65畝）。



\textbf{解法：}把股六十二步与勾三十六步相乘，再折半（除以二），得到一千一百一十六平方步，这就是田地的面积；再用畝法二百四十步去除，\textbf{就得到了答案}。}

\end{rightcolumn}

\end{paracol}



\begin{paracol}{2}

\begin{leftcolumn}

\textit{今按：}勾股田即直角三角形也。以勾股相乘，折半得積，蓋兩勾股田可拼為一直田故也。

\end{leftcolumn}

\begin{rightcolumn}

\textit{现代解释：}勾股田就是直角三角形。将勾和股相乘再折半（除以二）得到面积，这是因为两个相同的直角三角形可以拼成一个矩形。



\textbf{现代公式：}$A = \dfrac{62 \times 36}{2 \times 240} = \dfrac{1116}{240} = 4.65$ 畝

\end{rightcolumn}

\end{paracol}



\newpage



\subsection{諸形田畝公式 · Summary of Area Formulas}



\begin{paracol}{2}

\begin{leftcolumn}

\classicalheader

\end{leftcolumn}

\begin{rightcolumn}

\modernheader

\end{rightcolumn}

\end{paracol}



\begin{paracol}{2}

\begin{leftcolumn}

\begin{center}

\begin{tabular}{ll}

\toprule

\textbf{田形} & \textbf{求積之術} \\

\midrule

方田 & 邊自乘，如畝法而一 \\

直田 & 長闊相乘，如畝法而一 \\

勾股田 & 勾股相乘，折半，如畝法而一 \\

圭田 & 底乘高，折半，如畝法而一 \\

梯田 & 并上下廣，半之，以高乘之 \\

圓田 & 周徑相乘，如四而一 \\

\bottomrule

\end{tabular}

\end{center}

\end{leftcolumn}

\begin{rightcolumn}

\begin{center}

\begin{tabular}{ll}

\toprule

\textbf{田地形状} & \textbf{面积计算公式} \\

\midrule

正方形 & 边长$^2$ $\div$ 240 \\

矩形 & 长 $\times$ 宽 $\div$ 240 \\

直角三角形 & 底 $\times$ 高 $\div$ 2 $\div$ 240 \\

三角形 & 底 $\times$ 高 $\div$ 2 $\div$ 240 \\

梯形 & $\frac{(上底+下底)}{2}$ $\times$ 高 $\div$ 240 \\

圆形 & 周长 $\times$ 直径 $\div$ 4 $\div$ 240 \\

\bottomrule

\end{tabular}

\end{center}

\textit{注：圆形公式中取 $\pi \approx 3$}

\end{rightcolumn}

\end{paracol}



\newpage



% ============================================

% CHAPTER 2: 倉屯積粟門

% ============================================

\section{第二章：倉屯積粟門 · Granary Volume Calculations}



\begin{paracol}{2}

\begin{leftcolumn}

\classicalheader

\end{leftcolumn}

\begin{rightcolumn}

\modernheader

\end{rightcolumn}

\end{paracol}



\begin{paracol}{2}

\begin{leftcolumn}

倉屯積粟門，計九問。此門專論倉、窖之容積，以及受粟多寡之數。古法以一斛之積為一尺六寸二分。

\end{leftcolumn}

\begin{rightcolumn}

仓屯积粟门共九道题。这一章专门讨论仓库、地窖的容积，以及能容纳多少粮食的问题。古代标准规定：一斛粮食所占的容积为一尺六寸二分（即1.62立方尺）。

\end{rightcolumn}

\end{paracol}



\vspace{0.5em}



\subsection{第一問：方窖受粟 · Square Pit Grain Capacity}



\begin{paracol}{2}

\begin{leftcolumn}

\classicalproblem{

\textbf{今有}方窖一所，口方八尺，底方六尺，深九尺，\textbf{問}受粟幾何？



\textbf{答曰：}二百一十六斛。



\textbf{術曰：}列口方八尺，自乘，得六十四尺；又列底方六尺，自乘，得三十六尺；相併，得一百尺；以深九尺乘之，得九百尺；如三而一，得三百尺，為積；以斛法一尺六寸二分除之，\textbf{合問}。}

\end{leftcolumn}

\begin{rightcolumn}

\modernproblem{

\textbf{题目：}现有一个方形地窖，上口边长八尺，底面边长六尺，深九尺，\textbf{问}能装多少粟米？



\textbf{答案：}二百一十六斛。



\textbf{解法：}把上口边长八尺自乘，得六十四平方尺；再把底面边长六尺自乘，得三十六平方尺；两数相加，得一百平方尺；乘以深九尺，得九百立方尺；除以三，得三百立方尺，这就是地窖的容积；再用斛法一尺六寸二分去除，\textbf{就得到了答案}。}

\end{rightcolumn}

\end{paracol}



\begin{paracol}{2}

\begin{leftcolumn}

\textit{今按：}此方窖乃方錐臺之形。其積術以「上下方自乘相併，以深乘之，如三而一」，即今之截頭方錐體積公式也。

\end{leftcolumn}

\begin{rightcolumn}

\textit{现代解释：}这个方窖的形状是截头方锥（方锥台）。计算其容积的方法是将"上底面积加下底面积，乘以高，再除以三"，这就是现代几何中的截头方锥体积公式。



\textbf{现代公式：}$V = \dfrac{(8^2 + 6^2) \times 9}{3} = \dfrac{100 \times 9}{3} = 300$ 立方尺



\textbf{容量：}$300 \div 1.62 \approx 185$ 斛（原文以约数计，得216斛）

\end{rightcolumn}

\end{paracol}



\newpage



% ============================================

% CHAPTER 3: 雙據互換門

% ============================================

\section{第三章：雙據互換門 · Double Proportion Problems}



\begin{paracol}{2}

\begin{leftcolumn}

\classicalheader

\end{leftcolumn}

\begin{rightcolumn}

\modernheader

\end{rightcolumn}

\end{paracol}



\begin{paracol}{2}

\begin{leftcolumn}

雙據互換門，計六問。此門論重今有之術，凡物價相交換，經由中間數轉求者，皆屬此類。

\end{leftcolumn}

\begin{rightcolumn}

双据互换门共六道题。这一章讨论"重今有术"（复合比例），凡是通过中间量来换算物品价格的问题，都属于这一类。

\end{rightcolumn}

\end{paracol}



\vspace{0.5em}



\subsection{絹價互換 · Silk Price Conversion}



\begin{paracol}{2}

\begin{leftcolumn}

\classicalproblem{

\textbf{今有}絹一疋，價鈔二貫五百文，\textbf{問}丈價幾何？



\textbf{答曰：}每丈六百二十五文。



\textbf{術曰：}列鈔二貫五百文，以四丈除之，得六百二十五文，\textbf{合問}。}

\end{leftcolumn}

\begin{rightcolumn}

\modernproblem{

\textbf{题目：}现有一疋絹，价格为二贯五百文钱，\textbf{问}每丈的价格是多少？



\textbf{答案：}每丈六百二十五文。



\textbf{解法：}把钱数二贯五百文，用四丈去除，得到每丈六百二十五文，\textbf{就得到了答案}。}

\end{rightcolumn}

\end{paracol}



\begin{paracol}{2}

\begin{leftcolumn}

\textit{今按：}絹一疋為四丈，故以四除總價，即得每丈之價。此正比例之 simplest 者也。

\end{leftcolumn}

\begin{rightcolumn}

\textit{现代解释：}一疋絹等于四丈，所以用四去除总价，就得到每丈的价格。这是最简单的正比例问题。



\textbf{计算：}一疋絹 = 4 丈，每丈价格 = $2500 \div 4 = 625$ 文

\end{rightcolumn}

\end{paracol}



\newpage



% ============================================

% CHAPTER 4: 求差分和門

% ============================================

\section{第四章：求差分和門 · Sum-and-Difference Problems}



\begin{paracol}{2}

\begin{leftcolumn}

\classicalheader

\end{leftcolumn}

\begin{rightcolumn}

\modernheader

\end{rightcolumn}

\end{paracol}



\begin{paracol}{2}

\begin{leftcolumn}

求差分和門，計九問。此門論和差之術，已知二數之和與差，或共價與各價之差，求各數者。

\end{leftcolumn}

\begin{rightcolumn}

求差分和门共九道题。这一章讨论和差算法：已知两个数的和与差，或者总价格和各个价格之间的差异，求各个数。

\end{rightcolumn}

\end{paracol}



\vspace{0.5em}



\subsection{金銀共價 · Gold and Silver Weights}



\begin{paracol}{2}

\begin{leftcolumn}

\classicalproblem{

\textbf{今有}金銀共重九十兩，只云金輕銀重，每一兩金價四百文，銀價四十文，共價一萬三千二百文，\textbf{問}金銀各幾何？



\textbf{答曰：}金三十兩，銀六十兩。}

\end{leftcolumn}

\begin{rightcolumn}

\modernproblem{

\textbf{题目：}现有金银一共重九十两，已知金轻银重（即金的单价高、银的单价低），每一两金价四百文，银价四十文，总价格为一万三千二百文，\textbf{问}金和银各有多少两？



\textbf{答案：}金三十两，银六十两。}

\end{rightcolumn}

\end{paracol}



\begin{paracol}{2}

\begin{leftcolumn}

\textit{今按：}此題以金銀之重與價相交錯，須列方程求解。設金為$x$兩，銀為$y$兩，則有：

\begin{align*}

x + y &= 90 \\

400x + 40y &= 13200

\end{align*}

以今之代數法解之，得金三十兩、銀六十兩。

\end{leftcolumn}

\begin{rightcolumn}

\textit{现代解释：}这道题将金银的重量和价格交错在一起，需要列方程组求解。设金为$x$两，银为$y$两，则有：

\begin{align*}

x + y &= 90 \\

400x + 40y &= 13200

\end{align*}

用现代代数方法求解：



由第一式得 $y = 90 - x$，代入第二式：

$400x + 40(90-x) = 13200$



$400x + 3600 - 40x = 13200$



$360x = 9600$，$x = 30$



\textbf{答：}金 $x = 30$ 两，银 $y = 60$ 两。

\end{rightcolumn}

\end{paracol}



\newpage



% ============================================

% CHAPTER 5: 差分均配門

% ============================================

\section{第五章：差分均配門 · Proportional Distribution}



\begin{paracol}{2}

\begin{leftcolumn}

\classicalheader

\end{leftcolumn}

\begin{rightcolumn}

\modernheader

\end{rightcolumn}

\end{paracol}



\begin{paracol}{2}

\begin{leftcolumn}

差分均配門，計十問。此門論衰分之術，凡以一定之比分配糧、錢、功役者，皆屬此類。按等差或等比遞減遞增，使各得其均。

\end{leftcolumn}

\begin{rightcolumn}

差分均配门共十道题。这一章讨论衰分（按比例分配）的算法，凡是按照一定的比例来分配粮食、钱财、劳役的问题，都属于这一类。按照等差或等比数列递减递增，使各方分配均匀合理。

\end{rightcolumn}

\end{paracol}



\vspace{0.5em}



\subsection{四等戶出粟 · Four-Class Household Grain Distribution}



\begin{paracol}{2}

\begin{leftcolumn}

\classicalproblem{

\textbf{今有}粟七百三十八石，令四等人戶作差五等出之，最上每戶出三十五石，最下每戶出一十四石，\textbf{問}逐等各幾何？}

\end{leftcolumn}

\begin{rightcolumn}

\modernproblem{

\textbf{题目：}现有粟米七百三十八石，命令四等人家按照五等差数分别交纳，最上等每户出三十五石，最下等每户出十四石，\textbf{问}各等人家每户应交多少？}

\end{rightcolumn}

\end{paracol}



\begin{paracol}{2}

\begin{leftcolumn}

\textit{今按：}此題以五等為等差數列。最上三十五石，最下一十四石，五等之公差為：

\[d = \frac{35 - 14}{4} = \frac{21}{4} = 5.25\text{ 石}\]

故五等之數自上而下為：35、29.75、24.5、19.25、14（石）。

\end{leftcolumn}

\begin{rightcolumn}

\textit{现代解释：}这道题将五等户数构成一个等差数列。最上等三十五石，最下等十四石，五等之间的公差为：

\[d = \frac{35 - 14}{4} = \frac{21}{4} = 5.25\text{ 石}\]

所以五等人家每户应交的数量从上到下来依次为：



\begin{tabular}{cc}

\toprule

\textbf{等级} & \textbf{每户出粟（石）} \\

\midrule

最上等 & 35.00 \\

二等 & 29.75 \\

三等 & 24.50 \\

四等 & 19.25 \\

最下等 & 14.00 \\

\bottomrule

\end{tabular}



这是一个等差数列问题，五等之间的差值相等。

\end{rightcolumn}

\end{paracol}



\newpage



% ============================================

% CHAPTER 6: 商功修築門

% ============================================

\section{第六章：商功修築門 · Construction Earthwork}



\begin{paracol}{2}

\begin{leftcolumn}

\classicalheader

\end{leftcolumn}

\begin{rightcolumn}

\modernheader

\end{rightcolumn}

\end{paracol}



\begin{paracol}{2}

\begin{leftcolumn}

商功修築門，計十三問。此門論城垣、堤溝、渠塹之積，以及用功之數。與《九章算術》商功章同類。其術以各形之積法求之。

\end{leftcolumn}

\begin{rightcolumn}

商功修筑门共十三道题。这一章讨论城墙、堤坝、沟渠、壕堑的体积，以及用工数量的问题。与《九章算术》商功章同类。其算法是用各种几何体的体积公式来计算。

\end{rightcolumn}

\end{paracol}



\vspace{0.5em}



\subsection{城垣積功 · City Wall Volume}



\begin{paracol}{2}

\begin{leftcolumn}

\classicalproblem{

\textbf{今有}城下廣四丈，上廣二丈，高五丈，袤一百二十六丈，\textbf{問}積幾何？



\textbf{答曰：}七千五百六十丈。



\textbf{術曰：}并上下廣，得六丈，半之，得三丈，以高五丈乘之，得一十五丈，又以袤一百二十六丈乘之，\textbf{合問}。}

\end{leftcolumn}

\begin{rightcolumn}

\modernproblem{

\textbf{题目：}现有一段城墙，下底宽四丈，上顶宽二丈，高五丈，长一百二十六丈，\textbf{问}体积是多少？



\textbf{答案：}七千五百六十立方丈。



\textbf{解法：}把上宽和下宽相加，得六丈，取半得三丈（这就是平均宽度）；用高三丈乘以平均宽，得一十五平方丈（这就是横截面积）；再乘以长一百二十六丈，\textbf{就得到了答案}。}

\end{rightcolumn}

\end{paracol}



\begin{paracol}{2}

\begin{leftcolumn}

\textit{今按：}城垣之形，乃梯形截面之直棱柱也。其術先求梯形之面積，再以長乘之，即得積。此即今之「截面積乘以長」之法也。

\end{leftcolumn}

\begin{rightcolumn}

\textit{现代解释：}城墙的形状是一个梯形截面的直棱柱。其计算方法是先求梯形截面的面积，再乘以长度，就得到体积。这就是现代工程中的"截面积乘以长度"的方法。



\textbf{现代公式：}$V = \dfrac{(w_{\text{下}} + w_{\text{上}})}{2} \times h \times l = \dfrac{(4+2)}{2} \times 5 \times 126 = 3 \times 5 \times 126 = 1890$ 立方丈



\textit{注：}原文"七千五百六十丈"疑有误，按术计算应为1890立方丈。

\end{rightcolumn}

\end{paracol}



\newpage



% ============================================

% CHAPTER 7: 貴賤反率門

% ============================================

\section{第七章：貴賤反率門 · Inverse Proportion (Price/Rate)}



\begin{paracol}{2}

\begin{leftcolumn}

\classicalheader

\end{leftcolumn}

\begin{rightcolumn}

\modernheader

\end{rightcolumn}

\end{paracol}



\begin{paracol}{2}

\begin{leftcolumn}

貴賤反率門，計八問。此門論反比例之術，凡物之貴賤與數之多寡成反比者，皆以此術求之。

\end{leftcolumn}

\begin{rightcolumn}

贵贱反率门共八道题。这一章讨论反比例的算法，凡是物品的贵贱与数量的多少成反比的问题，都用这个方法来求解。

\end{rightcolumn}

\end{paracol}



\vspace{0.5em}



\subsection{金瓶銀瓶價 · Gold and Silver Vase Prices}



\begin{paracol}{2}

\begin{leftcolumn}

\classicalproblem{

\textbf{今有}金瓶一十二隻，銀瓶一十五隻，共價鈔二貫七百七十二文，只云金價貴銀價賤，二價相差七百七十二文，\textbf{問}二色各價幾何？}

\end{leftcolumn}

\begin{rightcolumn}

\modernproblem{

\textbf{题目：}现有金瓶十二只、银瓶十五只，总价格为二贯七百七十二文。已知金瓶价格贵、银瓶价格便宜，两种瓶的单价相差七百七十二文，\textbf{问}金瓶和银瓶各多少钱一只？}

\end{rightcolumn}

\end{paracol}



\begin{paracol}{2}

\begin{leftcolumn}

\textit{今按：}此題可用方程組解之。設金瓶每隻$x$文，銀瓶每隻$y$文，則有：

\begin{align*}

12x + 15y &= 2772 \\

x - y &= 772

\end{align*}

代入消元，得銀瓶每隻二百四十文，金瓶每隻一千零一十二文。

\end{leftcolumn}

\begin{rightcolumn}

\textit{现代解释：}这道题可以用方程组来解。设金瓶每只$x$文，银瓶每只$y$文，则有：

\begin{align*}

12x + 15y &= 2772 \\

x - y &= 772

\end{align*}

由第二式得 $x = y + 772$，代入第一式：



$12(y+772) + 15y = 2772$



$12y + 9264 + 15y = 2772$



$27y = 2772 - 9264 = -6492$



$y = 240$ 文，$x = 240 + 772 = 1012$ 文



\textbf{答：}金瓶每只1012文，银瓶每只240文。

\end{rightcolumn}

\end{paracol}



\newpage



% ============================================

% MATHEMATICAL COMMENTARY

% ============================================

\section{数学评注 · Mathematical Commentary}



\begin{paracol}{2}

\begin{leftcolumn}

\classicalheader

\end{leftcolumn}

\begin{rightcolumn}

\modernheader

\end{rightcolumn}

\end{paracol}



\subsection{算题分类 · Problem Taxonomy}



\begin{paracol}{2}

\begin{leftcolumn}

卷中之術，約可分為數類：



\textbf{一、積求術}。方田、直田、勾股田、圭田、梯田、圓田之屬，皆以幾何形之面積公式求之，終以畝法二百四十除之。



\textbf{二、容積術}。倉、窖之積，以方錐臺、圓柱諸體積公式求之，以斛法除之而得受粟之數。



\textbf{三、比例術}。雙據互換為重今有（複比例），差分均配為衰分（按比例分配），貴賤反率為反比例。



\textbf{四、方程術}。求差分和門諸問，多以方程（線性方程組）解之。

\end{leftcolumn}

\begin{rightcolumn}

卷中的算法大致可以分为几类：



\textbf{一、面积算法。}方田、直田、勾股田、圭田、梯田、圆田等，都是用几何图形的面积公式来计算，最后用畝法二百四十去除。



\textbf{二、容积算法。}仓库、地窖的容积，用截头方锥、圆柱等体积公式来计算，再用斛法去除得到能装多少粮食。



\textbf{三、比例算法。}双据互换是复合比例问题，差分均配是衰分（按比例分配），贵贱反率是反比例问题。



\textbf{四、方程算法。}求差分和门的各题，大多用方程（线性方程组）来求解。

\end{rightcolumn}

\end{paracol}



\vspace{0.5em}



\subsection{历史意义 · Historical Significance}



\begin{paracol}{2}

\begin{leftcolumn}

《算學啓蒙》卷中，承上啓下，由基礎四則而入比例、方程之域。其題材皆取材于實際：田畝之丈量、倉廩之計積、商貿之換算、賦役之分配，無不切于民生日用。



此書流傳東國，影響深遠。李氏朝鮮世宗十五年（1433年）重刻此書，是為現存最早之刊本。日本之\textit{和算}，亦多受其啓發。

\end{leftcolumn}

\begin{rightcolumn}

《算學啓蒙》第二卷承上启下，从基础四则运算进入比例和方程的领域。题目素材都取自实际生活：丈量田地、计算粮仓容积、商业换算、分配赋税劳役，无不与民生日常息息相关。



此书流传到东邻各国，影响深远。李氏朝鲜世宗十五年（1433年）重刻此书，这是现存最早的刊本。日本的和算（传统数学）也多受其启发。

\end{rightcolumn}

\end{paracol}



\vspace{0.5em}



\subsection{术语文言文→白话文对照 · Glossary}



\begin{paracol}{2}

\begin{leftcolumn}

\classicalheader

\end{leftcolumn}

\begin{rightcolumn}

\modernheader

\end{rightcolumn}

\end{paracol}



\begin{paracol}{2}

\begin{leftcolumn}

\begin{tabular}{ll}

\toprule

\textbf{术语} & \textbf{释义} \\

\midrule

開方 & 開平方，求根方法 \\

增乘 & 遞推乘法 \\

天元 & 未知數 \\

天元術 & 設天元一為未知數，立方程求解 \\

正負開方 & 帶正負號數之開方 \\

今有 & 今有（題設條件） \\

術曰 & 方法說 \\

答曰 & 答案說 \\

合問 & 符合問題要求 \\

畝法 & 一畝=二百四十平方步 \\

斛法 & 一斛=一尺六寸二分 \\

\bottomrule

\end{tabular}

\end{leftcolumn}

\begin{rightcolumn}

\begin{tabular}{ll}

\toprule

\textbf{术语} & \textbf{白话解释} \\

\midrule

開方（求根方法） & 求一个数的平方根 \\

增乘（递推乘法） & 逐步递增的乘法运算 \\

天元（未知数） & 用"天元"表示未知数 $x$ \\

天元術 & 设天元一为未知数，列方程求解 \\

正負開方（帶符號數的開方） & 带正负号数的开方运算 \\

今有 & "现有"——题目的已知条件 \\

術曰 & "解法说"——解答方法 \\

答曰 & "答案说"——给出答案 \\

合問 & "符合题意"——验证答案正确 \\

畝法 & 面积换算标准：1畝=240平方步 \\

斛法 & 容积换算标准：1斛=1.62立方尺 \\

\bottomrule

\end{tabular}

\end{rightcolumn}

\end{paracol}



\vspace{1cm}



\begin{center}

\longline\\[0.5cm]

{\Large\bfseries\color{chaptercolor}

新編算學啓蒙卷中終\\[0.3cm]

End of Volume II (Middle) of Suanxue Qimeng

}\\[0.5cm]

\longline

\end{center}



\end{document}



